\documentclass[11pt,letterpaper]{article}
\usepackage[T1]{fontenc}
\usepackage[utf8]{inputenc}
\usepackage{lmodern}
\usepackage[margin=1in,headheight=15pt]{geometry}
\usepackage{amsmath,amssymb,amsthm,mathtools,mathrsfs}
\usepackage{microtype,booktabs,longtable,tabularx,array,enumitem}
\usepackage{xcolor,fancyhdr,needspace}
\definecolor{Ink}{HTML}{1F3648}
\definecolor{Accent}{HTML}{2A5F73}
\usepackage[colorlinks=true,linkcolor=Ink,citecolor=Accent,urlcolor=Accent]{hyperref}
\usepackage{xurl}
\usepackage[nameinlink,noabbrev]{cleveref}
\usepackage{aliascnt}
\hypersetup{%
  pdftitle={Differential Algebra and Differential Equations over the Surreal and Surcomplex Numbers},
  pdfsubject={Merged research report: the finite-primitive criterion, the phase obstruction, and coherent coordinate systems},
  pdfkeywords={surreal numbers, surcomplex numbers, Berarducci--Mantova derivation, H-field, Hahn series, differential Galois theory, regular-singular systems, autonomous differential equations, abelian varieties, log-atomic numbers, critical potentials}}
\pagestyle{fancy}\fancyhf{}
\fancyhead[L]{\small Differential algebra over the surreal and surcomplex numbers}
\fancyhead[R]{\small\thepage}
\renewcommand{\headrulewidth}{0.3pt}
\setlength{\parskip}{0.25em}
\setlength{\emergencystretch}{3em}
\widowpenalty=10000 \clubpenalty=10000 \displaywidowpenalty=10000
\setlist{itemsep=0.2em,topsep=0.3em}
\setcounter{tocdepth}{2}
\setcounter{secnumdepth}{3}
\numberwithin{equation}{section}

% Every theorem-like environment shares the theorem counter through an alias
% counter, so that \cref prints the environment's own name (Corollary,
% Warning, ...) and not the name of the shared counter.
\newtheorem{theorem}{Theorem}[section]
\newcommand{\diffnewthm}[3]{\newaliascnt{#1}{theorem}%
  \newtheorem{#1}[#1]{#2}\aliascntresetthe{#1}\crefname{#1}{#2}{#3}}
\diffnewthm{proposition}{Proposition}{Propositions}
\diffnewthm{lemma}{Lemma}{Lemmas}
\diffnewthm{corollary}{Corollary}{Corollaries}
\theoremstyle{definition}
\diffnewthm{definition}{Definition}{Definitions}
\diffnewthm{example}{Example}{Examples}
\diffnewthm{convention}{Convention}{Conventions}
\theoremstyle{remark}
\diffnewthm{remark}{Remark}{Remarks}
\diffnewthm{warning}{Warning}{Warnings}
\crefname{theorem}{Theorem}{Theorems}
\crefname{section}{Section}{Sections}
\crefname{table}{Table}{Tables}
\crefname{part}{Part}{Parts}

\newcommand{\No}{\mathbf{No}}
\newcommand{\SC}{\mathbf{No}[i]}
\newcommand{\R}{\mathbb R}
\newcommand{\C}{\mathbb C}
\newcommand{\Q}{\mathbb Q}
\newcommand{\Z}{\mathbb Z}
\newcommand{\N}{\mathbb N}
\newcommand{\der}{\partial}
\newcommand{\Ofin}{\mathcal O}
\newcommand{\mm}{\mathfrak m}
\newcommand{\Pcl}{\mathcal P}
\newcommand{\Acl}{\mathcal A}
\newcommand{\Izero}{\mathcal I_{0}}
\newcommand{\Obs}{\mathfrak o}
\newcommand{\pp}{\operatorname{pp}}
\newcommand{\ct}{\operatorname{ct}}
\newcommand{\st}{\operatorname{st}}
\newcommand{\fin}{\operatorname{fin}}
\newcommand{\supp}{\operatorname{supp}}
\newcommand{\Ree}{\operatorname{Re}}
\newcommand{\Imm}{\operatorname{Im}}
\newcommand{\cis}{\operatorname{cis}}
\newcommand{\Emm}{\mathrm{E}}
\newcommand{\Lmm}{\mathrm{L}}
\newcommand{\Tor}{\mathbb T}
\newcommand{\Strip}{\mathcal S}
\newcommand{\Estrip}{E_{\Strip}}
\newcommand{\Hring}{\mathscr H}
\newcommand{\Hol}{\mathscr O}
\newcommand{\dercoef}{\der_{\mathrm H}}
\newcommand{\dertil}{\widetilde{\der}}
\newcommand{\derZ}{\der_{Z}}
\newcommand{\Dz}{D_{z}}
\newcommand{\Mat}{\operatorname{Mat}}
\newcommand{\GL}{\operatorname{GL}}
\newcommand{\tr}{\operatorname{tr}}
\newcommand{\Span}{\operatorname{span}}
\newcommand{\Ker}{\operatorname{ker}}
\newcommand{\gm}{\mathbf G_{\mathrm m}}
\newcommand{\SL}{\operatorname{SL}}
\newcommand{\Un}{\operatorname{U}}
\newcommand{\Herm}{\operatorname{Herm}}
\newcommand{\Symm}{\operatorname{Sym}}
\newcommand{\Hom}{\operatorname{Hom}}
\newcommand{\Homd}{\operatorname{Hom}_{\der}}
\newcommand{\Ph}{\operatorname{Ph}}
\newcommand{\Rot}{\mathcal R}
\newcommand{\Tser}{\mathcal T}
\newcommand{\abs}[1]{\lvert#1\rvert}
\newcommand{\Repo}{\texttt{VladimirReshetnikov/Surreal}}
\newcommand{\Commit}{\nolinkurl{e260237db9b71da8b74a0c13c8e6355119091100}}
\newcommand{\CommitB}{\nolinkurl{aa846271b4dcae2c055b216126a87210292ec19b}}
\newcommand{\CommitC}{\nolinkurl{a3124af79f66b8b9c196d76b4cbc5ac3938907c4}}
\newcommand{\CommitD}{\nolinkurl{048b72cf7cbfc8ab246e4f73788c10460cb3f6e0}}
% Notation of the transfinite critical-potential part (member 13); see
% diff:cp:conv:alphabet.
\newcommand{\Sch}{\operatorname{Sch}}
\newcommand{\Bor}{\operatorname{Bor}}
% Notation of the regular-singular part (member 12) and the autonomous part
% (member 11); see diff:rs:conv:alphabet and diff:aut:conv:alphabet.
\newcommand{\derT}{\der_{\tau}}
\newcommand{\KR}{K_{\R}}
\newcommand{\Eig}{\operatorname{Eig}}
\newcommand{\Lev}{\operatorname{Lev}}
\newcommand{\hob}{\operatorname{hob}}
\newcommand{\pr}{\operatorname{pr}}
\newcommand{\Sol}{\operatorname{Sol}}
\newcommand{\crit}{\mathrm{crit}}
\newcommand{\Dn}{\eth}
\newcommand{\Lie}{\operatorname{Lie}}
\newcommand{\dlog}{\operatorname{dlog}}
\newcommand{\Log}{\operatorname{Log}}
\newcommand{\src}[1]{{\small\textsf{[#1]}}}
\newcolumntype{L}[1]{>{\raggedright\arraybackslash}p{#1}}
\newcolumntype{Y}{>{\raggedright\arraybackslash}X}

\begin{document}
\hypersetup{pageanchor=false}
\begin{titlepage}
\centering
\vspace*{0.4in}
{\large\scshape A merged research report\par}
\vspace{0.45in}
{\LARGE\bfseries\color{Ink} Differential Algebra and\\[0.2em]
Differential Equations over the\\[0.2em]
Surreal and Surcomplex Numbers\par}
\vspace{0.3in}
{\Large The finite-primitive criterion, the phase obstruction,\\[0.15em]
and support-certified coordinate systems\par}
\vspace{0.45in}
{\large September 2026\par}
\vspace{0.45in}
\begin{minipage}{0.93\textwidth}
\begin{abstract}
\noindent
Eleven independently written manuscripts developed the theory of the
Berarducci--Mantova derivation $\der$ on the surreal class field $\No$,
normalized by $\der\omega=1$, and of its unique extension to the surcomplex
field $\SC$. This report is their union. Seven of them share a twelve-item
spine of scalar and workspace theory; four joined later and carry the
matrix, regular-singular, autonomous and transfinite second-order strata.

The organizing identity is exact: the image of the logarithmic derivative is
\[
 \{\der y/y : y\in\SC^{\times}\}=\No+i\Acl,
 \qquad \Acl:=\der\mm=\der\Ofin,
\]
where $\Ofin$ is the class of finite surreals and $\mm$ its ideal of
infinitesimals. Hence $\der y=(a+ib)y$ has a nonzero surcomplex solution
exactly when $b$ has a \emph{finite} surreal primitive. The two printed forms
$\No+i\der\mm$ and $\No+i\der\Ofin$ agree only through a proved lemma, and
$\Acl$ is a \emph{strict} subclass of $\mm$: the coefficient $\omega^{-1}$ is
infinitesimal while its primitive $\log\omega$ is infinite. The purely infinite
part of the accumulated phase is an $\R$-linear surjection onto the class of
purely infinite surreals, a complete gauge invariant, and the middle term of a
four-term exact sequence. From it follow sharp support and iterated-logarithm
thresholds, a decision procedure for rational coefficients, the pointwise
maximality of the strip $\No+i\Ofin$ as the domain of any
derivation-compatible exponential, the nonexistence of a nonzero solution of
$\der y=iy$ or $\der^{2}y+y=0$, a complete constant-coefficient and
constant-matrix classification in which only real characteristic roots
contribute, a triangular fundamental-matrix criterion, exact resolvents with
factorially divergent Hahn solutions, differentially stable set-sized
workspaces, and oscillatory Picard--Vessiot extensions with differential
Galois tori --- at the cost, in the real form, of the $H$-field ordering rule.

A later part classifies \emph{finite matrix} systems. Every $\der y=Ay$ with
$A$ an ordinary finite matrix over $\SC$ is gauge equivalent to
$i\operatorname{diag}(\der P_{1},\dots,\der P_{n})$ with each $P_{j}$ purely
infinite, and the multiset $\{P_{j}\}$ is a complete invariant of differential
equivalence; for a Hermitian $H$ it is computed by integrating the ordinary
algebraic eigenvalues of $H$ and keeping the purely infinite parts of the
primitives, with neither commutation nor a spectral gap assumed. From that
follow a gap-free finite-primitive perturbation theorem, an exact kernel
count, normalized non-Abelian integration, the complete cone of positive
invariant metrics, a real rotation-block descent, a positive Ermakov--Pinney
dichotomy, an exact Hahn expansion for the oscillatory Airy equation, and the
phase-lattice Picard--Vessiot torus of an arbitrary finite system. A
deliberate counterexample --- a nonnormal coefficient with instantaneous
eigenvalues $\pm i$ and a fundamental matrix already in the base field ---
marks the boundary of the eigenvalue formula, not of the finite-system classification. That part also
invokes a separately imported linear-splitting and first-order surjectivity
theorem, not a consequence of the scalar criterion and flagged wherever used.

Three later parts classify regular-singular systems $\der y=-t\,Ay$ with an
ordinary residual matrix and a positive real-power Hahn perturbation, for which
$\Ph(-tA)=\{(\Imm\lambda_{j})\log t\}$ is an explicit non-Hermitian sector reached
without the import; first-order autonomous equations with constant
coefficients, for \emph{every} normalized surreal derivation, with the
logarithmic-derivative image of a constant abelian variety; and second-order
equations on the log-atomic tower, with a sharp threshold $1/4$ at every
ordinal depth.

A deliberately separate half treats \emph{coordinate} differential equations
in fixed-common-domain Hahn rings $\Hol(U)((t^{\Gamma}))$: existence,
uniqueness, support certificates, variation of constants, a Liouville
determinant identity, monodromy valued in $\GL_{d}(K_{\Gamma})$, and a
polynomially nonlinear initial-value theorem, all built on strong summability
and an ordered-word support lemma proved here in full. The two halves share
notation and almost nothing else; several theorems are true for one operator
and false for the other on the same printed equation, and the report keeps
their type distinctions intact rather than flattening them.

The Berarducci--Mantova construction is imported, not reproved. Nothing here is
claimed as new, refereed, or machine-checked.
\end{abstract}
\end{minipage}
\vfill
\begin{minipage}{0.93\textwidth}\small
\textbf{Status.} A merged exposition assembled from eleven independent
manuscripts: seven pinned to \Repo\ at commit \Commit; the eighth,
which supplies \cref{diff:part:matrix}, pinned to the later revision
\CommitB; the ninth and tenth, which supply
\cref{diff:aut:part,diff:rs:part}, pinned to \CommitC; and the eleventh,
which supplies \cref{diff:cp:part}, pinned to \CommitD. Deep structural inputs
are attributed; the connecting deductions are proved in the text. No priority
claim, no solution of a named published open problem, no independent
refereeing, and no proof-assistant verification is asserted.
\Cref{diff:sec:nonclaims} is the consolidated register of what is
\emph{not} claimed; it is part of the result, not a disclaimer appended to it.
\end{minipage}
\end{titlepage}
\hypersetup{pageanchor=true}
\pagenumbering{roman}
\begingroup\small\setlength{\parskip}{0pt}
\tableofcontents
\endgroup
\clearpage
\pagenumbering{arabic}
\part{Foundations, and the distinctions that must survive}
\label{diff:part:foundations}

\section{What this report is}
\label{diff:sec:intro}

\subsection{Eleven manuscripts, one theory}
\label{diff:sub:seven}

This report is the union of eleven independently written manuscripts on the
differential algebra of the surreal and surcomplex numbers. Seven of them form
the original core: they were prepared separately, they are all pinned to the
same documentation snapshot, and they prove --- in seven notations --- the same
twelve-item spine:

\begin{enumerate}[label=(\roman*),leftmargin=2.6em]
\item the Berarducci--Mantova derivation $\der$ extends uniquely to
$\SC=\No[i]$, because differentiating $i^{2}+1=0$ forces $\der i=0$; its
constant field is exactly $\C$;
\item derivatives of finite elements are infinitesimal;
\item an infinitesimal exponential and logarithm, defined as Hahn sums, with
their chain rules;
\item a branch-free polar decomposition with a finite angle;
\item a primitive normalized by a vanishing $\omega^{0}$ coefficient;
\item the rank-one criterion --- $\der y=(a+ib)y$ has a nonzero surcomplex
solution if and only if $b$ has a finite primitive --- and the exact
logarithmic-derivative image;
\item the purely infinite obstruction, its exact sequence, and a gauge normal
form;
\item $\der y=iy$ and $\der^{2}y+y=0$ have only the zero solution;
\item there is no globally derivation-compatible complex exponential;
\item the constant-coefficient classification, in which only real
characteristic roots contribute, with its constant-matrix corollary;
\item a set-sized differential localization;
\item an oscillatory Picard--Vessiot extension with no new constants.
\end{enumerate}

Every one of the seven proves items (i)--(vi) and (viii)--(ix).  The
remaining four are proved by most but not all: the gauge normal form of (vii)
is absent from three members, its exact sequence from a fourth; the
constant-coefficient classification (x) is absent from one member, which
explicitly declines it, and its constant-matrix corollary from two others; and
at least five prove each of (xi) and (xii).  \Cref{diff:app:provenance} gives
the member-by-member record, and no claim below rests on a member that does not
prove what it is cited for. Keeping them apart would mean printing one theory seven times,
including seven proofs of $\der i=0$ and seven proofs that $\der y=iy$ has no
nonzero solution. Each is nevertheless the unique source of some genuinely
distinct payload, and those payloads are all built on the shared spine and
cannot be extracted without carrying it along. \Cref{diff:app:provenance}
records, member by member, what each contributed.

The eighth manuscript joined later and is different in kind. It does not
reprove the spine; it cites it, names the gap this report's own scope ledger
recorded --- no classification of arbitrary nontriangular matrix systems ---
and closes the Hermitian half of it. Its mathematics is
\cref{diff:part:matrix}, and it is pinned to a later revision, \CommitB, of
the same repository. Several of the twelve spine items are reproved inside it by
its own routes, and where that happens the routes are recorded as agreement
(\cref{diff:rem:matagreement,diff:rem:matrankone}) rather than printed twice.
It also brings one \emph{new imported input}, the linear-splitting and
first-order surjectivity theorem of \cref{diff:thm:matsurjectivity}, which the
first seven do not use and which the register of non-claims accordingly
records as an import rather than as a proved consequence.

The ninth and tenth manuscripts --- members~11 and~12 in the numbering of
\cref{diff:app:provenance} --- joined next, arrived together, and are pinned to
a third revision, \CommitC{} \cite{RepoSnapshotC}. Each compared itself with the repository's root
README, its catalogue and this report's README at that revision, not with this
article, and neither contradicts the report or the other. Between that revision
and this merge the article changed only by the addition of
\cref{diff:rem:tailspans}, and the two members' descriptions of what the report
already contained were checked against the article and hold. They share no
theorem and are printed as two new parts.
\begin{itemize}[leftmargin=1.6em]
\item Member~12 supplies \cref{diff:rs:part}: regular-singular matrix systems
over real-exponent Hahn workspaces. It is the matrix version of the rank-one
resonant normal form of \cref{diff:part:workspaces}, and over $\SC$ it computes
one explicit non-Hermitian sector of the phase spectrum of
\cref{diff:part:matrix} without that part's imported input. It partly answers
items N24 and N57 and the computable-filtration question of
\cref{diff:warn:matlimits}, and it does not bear on item~N41.
\item Member~11 supplies \cref{diff:aut:part}: all first-order autonomous
equations with constant coefficients, on curves and on abelian varieties. It
lifts items N21 and N25 and the remark after \cref{diff:cor:riccati} for that
class, and it is proved for every normalized derivation, which required a
scoped amendment of the one-derivation rule (\cref{diff:aut:conv:scope}).
\end{itemize}
Their reproofs of results already here --- smallness, the missing oscillator,
the ordered-word lemma, the strip logarithm, the logarithmic separator, the
missing logarithmic primitive --- are recorded as agreements and cited, not
printed twice.

The eleventh manuscript --- member~13 --- joined after them and is pinned to a
fourth revision, \CommitD{} \cite{RepoSnapshotD}. At that revision this article
was byte-identical to the one it was merged into, and the member's description
of what the report contains --- the imported derivative, the type
distinctions, the summability conventions, the finite-primitive criterion and
the normalized primitive --- holds. It supplies \cref{diff:cp:part}: the
second-order equations $\der^{2}y+(Q_{\alpha}+c(\ell_{\alpha}^{\dagger})^{2})y=0$
along the ordinal log-atomic tower, with the sharp threshold $c=\frac14$ at
every ordinal depth, the failure of every set of earlier scales to decide it,
and a limit-stage filtration by differential Hahn fields with a Borel
Picard--Vessiot group. It amends item~N21 for one explicit family of
second-order equations, and it is adjacent to, without answering, member~12's
layered logarithmic continuation (\cref{diff:sub:nc:open}). Its reproof of the
obstruction for nonreal powers of an infinite scale is this report's criterion,
and is recorded as agreement (\cref{diff:cp:cor:nonreal}); its connection with
the missing oscillator is made precise in \cref{diff:cp:rem:oscillator}.

Merging carries a specific risk, which \cref{diff:sec:hazard} addresses before
any theorem is stated: several different objects in these manuscripts are
called ``phase'', and two nearly identical classes are both called the
exponential domain. A merged sentence that unifies either word, or either
symbol, proves false statements. This report fixes one name and one symbol per
object, keeps the two classes visibly distinct, and states the distinction
wherever a reader might carry one sense into another.

\subsection{The documented gap}
\label{diff:sub:gap}

The seven audits agree on the gap and on its character. All are pinned to the
same snapshot of \Repo:
\begin{center}\small\Commit\end{center}
That snapshot describes fifteen research packages covering surreal normal
forms and constructions, surcomplex analysis and geometry, polynomial algebra,
trigonometry, foundations and symbolic computation \cite{RepoMap}.
Differentiation is not absent from that collection. What is absent is the step from the collection's careful
\emph{distinctions} between kinds of derivative to existence, obstruction and
support-control \emph{theorems} for differential equations.

Three boundaries in the existing documentation delimit the gap precisely.
The trigonometry report constructs finite phases and states explicitly that it
chooses no derivation on $\No$, and that its classification of global phase
extensions is not a classification of differential-equation solutions
\cite{RepoTrig}. The computer-algebra report separates an external-variable
derivative, a formal monomial derivative and a derivation of the surreal
scalar field, and its subsection ``Differential equations must name their
category'' issues a uniqueness warning rather than developing a solution
theory \cite{RepoCAS}. The foundations report insists that a fine derivative,
a formal-series germ, a coherent Hahn datum and a scalar-field derivation stay
four different interfaces \cite{RepoFound}. The analysis report goes further
still: its derivation-obstruction theorem already shows that one chosen
canonical global surcomplex exponential is incompatible with an internal scalar
derivation normalized by $\der\omega=1$ \cite{RepoAnalysis}. That existing
theorem is prior coverage, acknowledged as such; \cref{diff:cor:noglobalPsi}
below is strictly stronger in its phase-independence and is compatible with it.

Every one of the seven audits was targeted: the catalogue, the specialist
READMEs, and selected source passages, at one pinned revision. None was a
line-by-line certification of the fifteen reports or of their archived
sources, and none used the absence of a code-search hit as evidence of
absence. \Cref{diff:sec:nonclaims} records those audit limits in full.

\subsection{What is imported and what is proved}
\label{diff:sub:imported}

The deep input is the existence of the derivation, together with its strong
additivity, its constant field, its compatibility with the surreal
exponential, its surjectivity and its $H$-field positivity. These are imported
from Berarducci and Mantova \cite{BM} and are not reproved; their construction
through log-atomic numbers, path derivatives and nested truncation ranks is
not reproduced. The positive-support summability lemma is imported in its
standard form from Neumann \cite{Neumann}, and its ordered-word core is
additionally proved in full in \cref{diff:app:support} because the original
paper was not retrieved. Restricted composition and the global
composition-incompatibility theorem for $\omega$-series are separately imported
from \cite{BMGerms} and are flagged wherever used. The universal $H$-field
results of \cite{ADH} are background for
\cref{diff:part:foundations,diff:part:scalar,diff:part:workspaces,diff:rs:part,diff:aut:part,diff:part:coordinate},
not premises; \cref{diff:cp:part} uses a different part of \cite{ADH}, the
ordinal log-atomic formulas of its Section~2, as a premise, and not its
elementary embedding; in \cref{diff:part:matrix} alone they become a premise, because
\cref{diff:thm:matsurjectivity} transfers along the elementary embedding of
\cite{ADH} two further published facts --- splitting of monic scalar operators
into first-order factors and surjectivity of nonzero scalar operators
\cite[Lemma 2.4.19]{ADHNorm}. That import is a genuine addition to the
hypothesis base of this report and is flagged at every use
(\cref{diff:warn:matimport}, item~N53). Ordinary complex
initial-value theory \cite{Teschl} supplies the analytic ingredient of the
coordinate half, and Picard--Vessiot terminology is used in the standard sense
of \cite{PS,Singer}.

\Cref{diff:rs:part,diff:aut:part} add no new import of that kind. \Cref{diff:rs:part} uses
the Berarducci--Mantova clauses \emph{without} surjectivity and without
\cref{diff:thm:matsurjectivity} (\cref{diff:rs:rem:imports}); classical
regular-singular theory \cite{Levelt,KW,vdH} is its acknowledged antecedent.
\Cref{diff:aut:part} is stated for the whole class of normalized derivations
(\cref{diff:aut:def:normalized}), whose one member needed for existence is the
Berarducci--Mantova derivation. Surjectivity is assumed only for the
global primitive-obstruction construction and its finite-primitive
criterion; the explicit speed threshold needs no such assumption.
The part imports the classical local normal forms and Rosenlicht's
algebraic-independence theorem as presented by Noordman, van der Put and Top
\cite{NPT}, with the standard setting of abelian varieties \cite{Milne} and the
valuative reading of properness \cite{Stacks}.
\Cref{diff:cp:part} imports one input about the derivation that
\cref{diff:thm:BM} does not list: the ordinal log-atomic tower and its
derivatives, limit stages included \cite[Section~2]{ADH}
(\cref{diff:cp:prop:tower}). It does not use surjectivity in its main
construction, and it takes the classical Euler and
Kneser--Weber--Hartman--Hille thresholds as its acknowledged antecedents
\cite{KT}.

Everything else is proved here. None of it is claimed to be new.

\begin{remark}[Reading order]
\label{diff:rem:reading}
\Cref{diff:part:foundations} fixes conventions, imports the derivation,
derives smallness, and installs the vocabulary discipline. Readers who want
only the headline should then read \cref{diff:sec:phasecrit}.
\Cref{diff:part:scalar} is the scalar theory over $\SC$;
\cref{diff:part:matrix} raises it to finite matrix systems and contains the
spectral--integral theorem \cref{diff:thm:matspectral};
\cref{diff:part:workspaces} is its set-sized and Hahn workspaces;
\cref{diff:rs:part} classifies regular-singular matrix systems over those
workspaces and over $\SC$; \cref{diff:aut:part} classifies first-order
autonomous equations with constant coefficients for every normalized
derivation; \cref{diff:cp:part} classifies one family of second-order
equations along the transfinite log-atomic tower and continues the missing
oscillator there;
\cref{diff:part:coordinate} is the separate coordinate theory, whose only
dependence on \cref{diff:part:scalar} is negative --- it explains what the
scalar obstruction does \emph{not} forbid. \Cref{diff:part:limits} collects the
computational contracts and the register of non-claims.
\end{remark}

\section{Conventions, classes, and normal forms}
\label{diff:sec:conventions}

\subsection{Class-theoretic ground rules}
\label{diff:sub:classes}

We reason classically: ZFC with definable classes, or NBG with choice, or any
conservative class-theoretic presentation with enough replacement for the
constructions displayed below. The symbols $\No$ and $\SC$ denote proper
classes; a ``class field'' assertion means that the field operations are class
functions satisfying the usual identities, and a class-level statement is a
statement about elements.

Every individual number, every normal-form support, every family that is
summed, every polynomial, and every matrix below is a \emph{set}. No
proper-class family of nonzero terms is summed and no sum is indexed by $\No$.

\begin{warning}[No class of proper-class cosets]
\label{diff:warn:cosets}
The suggestive notation $\No/\Ofin$ must never be implemented as a collection
whose members are cosets of $\Ofin$: a coset of a proper class is not an
object that can belong to another class. Throughout this report it denotes the
\emph{represented} quotient, namely the class $\Pcl$ of purely infinite normal
forms together with the projection $\pp$ of \cref{diff:eq:split}. Every
obstruction value is an actual surreal number, hence a set. Likewise,
automorphism groups over a proper-class base are encoded by constant
multipliers of finitely many generators, not by forming a class whose members
are proper-class maps; and rank-one equations are classified by their scalar
coefficients together with an explicitly stated gauge relation, not by a
class-sized collection of modules. Inside a fixed set-sized field, ordinary
set quotients are of course legitimate; \cref{diff:sec:localization} supplies
such fields.
\end{warning}

\subsection{Normal forms and the three-part splitting}
\label{diff:sub:normalforms}

$\No$ is real closed and $\SC=\No[i]$ is algebraically closed. We use the
Conway normal form, the real surreal (Gonshor) exponential and its inverse
$\log$ on positive elements as background \cite{BM,Gonshor}. A real surreal
has a unique expansion
\begin{equation}\label{diff:eq:normal}
 x=\sum_{\alpha\in S}r_{\alpha}\omega^{\alpha},
 \qquad r_{\alpha}\in\R\setminus\{0\},
\end{equation}
with $S$ a reverse well-ordered set of surreal exponents. For surcomplex
numbers the coefficients lie in $\C$. It is often convenient to reverse the
orientation and write
\begin{equation}\label{diff:eq:tmonomial}
 t^{\gamma}:=\omega^{-\gamma},
\end{equation}
so that a support is well ordered in increasing $\gamma$ and positive
exponents describe infinitesimal monomials. We abbreviate $t=t^{1}=\omega^{-1}$.

\begin{warning}[Two different exponent constructions]
\label{diff:warn:monomialmap}
In \cref{diff:eq:normal} and \cref{diff:eq:tmonomial},
$\gamma\mapsto\omega^{-\gamma}$ is the Conway \emph{monomial map}. For an
arbitrary surreal $\gamma$ it must not be silently replaced by
$\exp(-\gamma\log\omega)$: the omega-map and general surreal exponentiation are
different constructions. The identity
$\der(\omega^{r})=r\omega^{r-1}$, equivalently
$\der(t^{r})=-r\,t^{r+1}$, is asserted for \emph{ordinary real} $r$ only, where
the two interpretations agree. The safe universal identity is
$\der(\omega^{\gamma})=\omega^{\gamma}\,\der\log(\omega^{\gamma})$ with the
actual Gonshor logarithm of the monomial. A symbolic implementation needs
different constructors for the two operations. For the record, the expression
$\exp(\gamma\log\omega)$ has derivative
$\exp(\gamma\log\omega)\bigl((\der\gamma)\log\omega+\gamma/\omega\bigr)$, which
is not a general rule for the monomial map.
\end{warning}

Write
\[
 \Ofin=\{x\in\No:\abs x<n\text{ for some }n\in\N\},
 \qquad
 \mm=\{x\in\No:\abs x<1/n\text{ for every }n\geq1\}.
\]
Thus $\Ofin$ is the convex valuation ring of finite elements and $\mm$ is its
maximal ideal, with $\Ofin=\R\oplus\mm$. Let $\Pcl$ be the additive class of
\emph{purely infinite} normal forms: those all of whose exponents in
\cref{diff:eq:normal} are strictly positive, together with $0$. It is an
additive real vector class. It is \emph{not} the class of all infinite
surreals. Splitting the support of \cref{diff:eq:normal} at exponent $0$ gives
the unique additive decomposition
\begin{equation}\label{diff:eq:split}
 x=\pp(x)+\ct(x)+\varepsilon,
 \qquad \pp(x)\in\Pcl,\quad\ct(x)\in\R,\quad\varepsilon\in\mm,
\end{equation}
and we write $\fin(x)=\ct(x)+\varepsilon=x-\pp(x)$. All of $\pp$, $\ct$ and
$\fin$ are additive and $\R$-linear, and extend $\C$-linearly to $\SC$
componentwise. So $\No=\Pcl\oplus\R\oplus\mm$ and $\Ofin=\R\oplus\mm$, and $x$
is finite exactly when $\pp(x)=0$.

\begin{warning}[$\ct$ is not a standard part and is not multiplicative]
\label{diff:warn:ct}
$\ct(x)$ is the coefficient of the monomial $\omega^{0}=1$, defined for
\emph{every} surreal, including infinite ones. The standard part $\st$ is
defined only on $\Ofin$, where it is a ring homomorphism with kernel $\mm$;
the two agree only there. $\ct$ is additive and real-linear but is not
multiplicative: $\ct(\omega)=\ct(\omega^{-1})=0$ while
$\ct(\omega\cdot\omega^{-1})=1$. None of $\pp$, $\ct$, $\fin$ is a field
homomorphism.
\end{warning}

For complex values, conjugation fixes $\No$ and sends $i$ to $-i$, and
$\abs{u+iv}=(u^{2}+v^{2})^{1/2}$. Put $\Ofin_{\C}=\Ofin+i\Ofin$ and
$\mm_{\C}=\mm+i\mm$; standard part is taken coordinatewise and lands in $\C$.
An element of modulus one is necessarily finite, whatever the birthdays of its
coordinates.

\subsection{Summability is a support condition}
\label{diff:sub:summability}

\begin{definition}[Strong summability]
\label{diff:def:summable}
A \emph{set}-indexed family $(x_{j})_{j\in J}$ of Hahn normal forms is
\emph{summable} (equivalently \emph{strongly summable}) when
$\bigcup_{j}\supp x_{j}$ is well ordered in the increasing-exponent convention
of \cref{diff:eq:tmonomial} and, at every exponent, only finitely many $x_{j}$
have a nonzero coefficient. Its sum is formed coefficientwise, by those finite
coefficient sums.
\end{definition}

\begin{warning}[What summability is not]
\label{diff:warn:summability}
Strong summability is a statement about supports. Increasing valuations alone
do not establish it. It licenses coefficientwise rearrangement of a summable
family; exchanging two infinite sums requires strong summability of the joint
family, not merely of each row. And \emph{writing} a formal geometric or
exponential series is not by itself a summability proof. None of the Hahn
identities below is asserted to be a limit of partial sums in the fine surreal
topology, in which set-indexed convergence is degenerate. Conversely, a
fixed-rank set-sized Hahn field does have nontrivial valuation-topological
convergence, and full-class topological statements must not simply be
transferred to it. The three notions --- fine topology, fixed-rank valuation
topology, strong summability --- are never identified here.

A concrete separator: with $\Gamma=\Z^{2}$ lexicographically ordered,
$v(t)=(0,1)$ and $v(u)=(1,0)$, one has $v(t^{n})<v(u)$ for every ordinary $n$.
The identity $\sum_{n\geq0}t^{n}=(1-t)^{-1}$ is strongly summable, yet its
partial sums do not converge in that valuation topology. At higher rank
$n\sigma$ need not become cofinal, and \cref{diff:app:support} establishes no
such cofinality; no argument in this report uses one.
\end{warning}

The combinatorial engine is the positive-support lemma.

\begin{lemma}[Positive-support calculus]
\label{diff:lem:neumann}
Let $\Gamma$ be an ordered abelian group and let $S\subseteq\Gamma_{>0}$ be
well ordered. Then the additive monoid
\[
 S^{*}=\{0\}\cup S\cup(S+S)\cup(S+S+S)\cup\cdots
\]
is well ordered, and for each $\gamma\in\Gamma$ only finitely many finite
\emph{ordered} words $(s_{1},\dots,s_{n})$ in $S$ satisfy
$s_{1}+\dots+s_{n}=\gamma$. Consequently, if $h$ has positive support, then
for ordinary real or complex coefficients $a_{n}$ the family $(a_{n}h^{n})_{n\geq0}$
is summable.
\end{lemma}

This is the lemma underlying the classical Hahn--Neumann construction
\cite{Neumann}, used in the summability preliminaries of \cite{BM,BMGerms}.
\Cref{diff:app:support} proves it in full, by way of Higman's finite-word
lemma \cite{Higman}, and stresses that the count is over \emph{ordered} words:
that stronger form is what noncommuting matrix products require in
\cref{diff:sec:coherentlinear,diff:sec:twoscale}, and an unordered or multiset
count will not carry those theorems.

\begin{warning}[Splitting a differentiated sum incurs two obligations]
\label{diff:warn:splitting}
If a family $(a_{k}x^{k})_{k}$ is summable, strong additivity gives
$\der\sum_{k}a_{k}x^{k}=\sum_{k}\der(a_{k}x^{k})$. To split the right-hand
side into a coefficient-derivative sum and a variable-derivative sum, both
separated families must \emph{themselves} be summable. Summability of a family
of sums does not imply summability of its summand families, because
cancellation can hide a violation. Every Taylor evaluation used below carries
the required positive-support certificate;
\cref{diff:prop:chain} is the one place where the summability of the two split
families is \emph{derived} rather than assumed, and it says exactly which
extra structure makes that possible.
\end{warning}

\section{The imported derivation, and smallness as a theorem}
\label{diff:sec:BM}

\subsection{The imported package}
\label{diff:sub:BMpackage}

\begin{theorem}[Berarducci--Mantova: imported input]
\label{diff:thm:BM}
There is a distinguished derivation $\der:\No\to\No$ --- the simplest
Berarducci--Mantova derivation --- with the following properties.
\begin{enumerate}[label=\textup{(BM\arabic*)},leftmargin=3.4em]
\item\label{diff:BM:leibniz} Leibniz: $\der(xy)=x\der y+y\der x$.
\item\label{diff:BM:const} Constant field: $\Ker\der=\R$, exactly.
\item\label{diff:BM:additive} Strong additivity:
$\der\bigl(\sum_{j\in J}x_{j}\bigr)=\sum_{j\in J}\der x_{j}$ for every
set-indexed summable family, \emph{and} the family of derivatives is itself
summable.
\item\label{diff:BM:exp} Exponential compatibility:
$\der(\exp x)=(\exp x)\der x$, for the real surreal exponential.
\item\label{diff:BM:norm} Normalization: $\der\omega=1$.
\item\label{diff:BM:surj} Surjectivity: $\der\No=\No$.
\item\label{diff:BM:positivity} $H$-field positivity: $x>\R$ implies
$\der x>0$.
\end{enumerate}
\end{theorem}

This is the package supplied by \cite[Theorems A/6.30 and B/7.7]{BM}; the
theorem numbering is that of the linked preprint. Its proof is deep and is not
reproduced here. Assigning a derivative to a few familiar monomials is not a
construction of a strongly additive derivation on every surreal: log-atomic
numbers require extra initial data, the construction specifies a
pre-derivation there and extends it through summable path derivatives, and both
the summability proof and the surjectivity proof are substantial inputs rather
than consequences of the product rule. This report uses their conclusions, not
a pretended implementation.

\begin{warning}[Two clauses that are easy to drop, and one that is not in the list]
\label{diff:warn:BMclauses}
Clause \ref{diff:BM:additive} is load-bearing for every Hahn-sum computation
below, including the infinitesimal exponential and logarithm
(\cref{diff:lem:smallchain}), the Laurent resolvents
(\cref{diff:thm:resolvent}), and the coefficient derivation on common-domain
families (\cref{diff:thm:coefderivation}); it is \emph{not} supplied by the
Leibniz rule. Clause \ref{diff:BM:positivity} is load-bearing for
\cref{diff:lem:smallnegative} and hence for the entire obstruction theory, and
for the ordering obstruction of \cref{diff:thm:orderobstruction}. Clause
\ref{diff:BM:norm} is indispensable for discussing a \emph{named} equation such
as $\der y=iy$: every named-equation statement below is a statement about this
derivation with this normalization, and none of them is derivation-agnostic.
(Scope amendment, added with member~11: the sentence just written holds for
every part except \cref{diff:aut:part}, whose statements are proved for every
derivation satisfying all of the clauses above except \ref{diff:BM:surj}, and
which says so at each statement; see \cref{diff:aut:conv:scope}. No statement
outside that part is thereby made derivation-agnostic.)

What is deliberately \emph{not} in the list is \emph{smallness}, the assertion
$\der\mm\subseteq\mm$. Three of the seven source manuscripts import it as a
clause of the Berarducci--Mantova package, and one of those flags it as
specific to the chosen derivation rather than a general fact. The other four
derive it from \ref{diff:BM:positivity}. \Cref{diff:sub:smallness} derives it,
for two reasons: advertising smallness as part of an imported theorem overstates
the import when it follows from an axiom already assumed; and the derived
version additionally yields strict order reversal on the infinitesimals
(\cref{diff:lem:smallnegative}), which the imported version does not give and
which \cref{diff:thm:idealstructure} uses.
\end{warning}

\begin{warning}[No uniqueness, and no other derivation is classified]
\label{diff:warn:nouniqueness}
Other surreal derivations exist. \Cref{diff:thm:BM} is a usable interface, not
a uniqueness characterization, and the elementary properties listed do not
characterize every possible surreal derivation. Nothing below classifies the
alternatives. The word ``canonical'', where it survives from a source, names
this distinguished construction and nothing more. Every intrinsic statement in
this report is relative to this one fixed $\der$.

\emph{Scope amendment, added with member~11.} The last sentence is kept, and
it now holds for every part except \cref{diff:aut:part}. That part works with
the class of \emph{normalized} derivations of \cref{diff:aut:def:normalized} ---
the clauses of \cref{diff:thm:BM} without surjectivity --- and proves that one
sector is the same for all of them: the solution set of every first-order
autonomous equation $F(y,y')=0$ with $F\in\C[Y,Z]$, and the solutions of every
constant meromorphic vector field on a smooth projective curve
(\cref{diff:aut:thm:independence}). That is a statement about a class of
derivations, not a classification of its members: it does not assert that the
class has one member, that two members agree on $\SC$, or that any theorem of
another part holds for another derivation. The only statements of other parts
that it reproves for every normalized derivation are the first-order ones listed
in \cref{diff:aut:cor:earlier}.
\end{warning}

From now on we write $x'=\der x$ interchangeably. The elementary rules follow
at once from \cref{diff:thm:BM}: for $x\neq0$, $u>0$ and $r\in\R$,
\begin{equation}\label{diff:eq:basicder}
 \der(x^{-1})=-x^{-2}\der x,\qquad
 \der\!\left(\frac{x}{y}\right)=\frac{y\der x-x\der y}{y^{2}},\qquad
 \der\log u=\frac{\der u}{u},\qquad
 \der(u^{r})=r\,u^{r-1}\der u,
\end{equation}
where $u^{r}=\exp(r\log u)$; in particular $\der\omega^{r}=r\omega^{r-1}$ and
$\der t^{r}=-r\,t^{r+1}$ for ordinary real $r$ (see
\cref{diff:warn:monomialmap}), $\der t=-t^{2}$, and $\der\log\omega=t$. With
$\ell_{0}=\omega$ and $\ell_{n+1}=\log\ell_{n}$, each $\ell_{n}$ is positive
infinite and induction on \cref{diff:eq:basicder} gives the iterated-logarithm
derivatives
\begin{equation}\label{diff:eq:logtower}
 \der\ell_{n}=\frac{1}{\ell_{0}\ell_{1}\cdots\ell_{n-1}}\qquad(n\geq1),
\end{equation}
with the empty product understood as $1$. A rational power rule
$\der(x^{q})=qx^{q-1}\der x$ for $q\in\Q$ and $x>0$ is available purely
algebraically, without the exponential.

For an algebraic element there is a corresponding implicit rule.

\begin{proposition}[Algebraic implicit differentiation]
\label{diff:prop:algdiff}
Let $F\subseteq\SC$ be a differential subfield and let $\alpha$ be algebraic
over $F$ with minimal polynomial $P=\sum_{k}a_{k}X^{k}\in F[X]$. Then
\[
 \der\alpha=-\frac{P^{\der}(\alpha)}{P_{X}(\alpha)},
 \qquad P^{\der}(X)=\sum_{k}(\der a_{k})X^{k}.
\]
\end{proposition}
\begin{proof}
In characteristic zero an irreducible polynomial is separable, so
$P_{X}(\alpha)\neq0$. Differentiating the finite identity $P(\alpha)=0$ gives
$P^{\der}(\alpha)+P_{X}(\alpha)\der\alpha=0$.
\end{proof}

This is the derivative rule for an exact algebraic root object; it is not a
choice of a numerical root branch.

\subsection{The unique complexification}
\label{diff:sub:complexification}

\begin{proposition}[Unique extension to $\SC$]
\label{diff:prop:complexification}
There is exactly one derivation on $\SC=\No[i]$ extending $\der$, namely
\[
 \der(a+ib)=\der a+i\der b .
\]
It commutes with conjugation, is strongly additive on set-indexed summable
families in $\SC$, is surjective, and its constant field is exactly $\C$.
Moreover, for $z\neq0$ and $r=\abs z$,
\begin{equation}\label{diff:eq:modulusder}
 \frac{\der r}{r}=\Ree\!\left(\frac{\der z}{z}\right).
\end{equation}
\end{proposition}
\begin{proof}
Any extension satisfies $0=\der(i^{2}+1)=2i\,\der i$, hence $\der i=0$, which
forces the displayed formula; direct expansion verifies the Leibniz rule.
Conjugation commutes with it by inspection. For strong additivity, take real
and imaginary parts: their derivative families are summable by
\ref{diff:BM:additive}, a finite union of well-ordered supports is well
ordered, and each exponent retains finite multiplicity. The kernel consists of
the elements both of whose components lie in $\Ker\der=\R$, that is exactly
$\C$. Given $a+ib$, choose real primitives $A,B$ by \ref{diff:BM:surj}; then
$\der(A+iB)=a+ib$. Finally $r^{2}=z\bar z$ gives
$2r\der r=(\der z)\bar z+z\overline{\der z}$; divide by $2r^{2}$ to obtain
\cref{diff:eq:modulusder}.
\end{proof}

All seven manuscripts prove this proposition, by this four-line argument. Two
consequences deserve emphasis. First, the constant field \emph{changes} from
$\R$ to $\C$ on complexification, and several later statements --- the kernel
of the logarithmic derivative, the fact that solution spaces are
$\C$-spans --- depend on that. Second, $\der$ is surjective on $\SC$ and $\SC$
is algebraically closed, and neither fact solves $\der y=ay$ with $y\neq0$:
surjectivity solves $\der y=f$, and algebraic closure supplies roots of
polynomials in existing elements. Algebraic closure is not differential
solution closure, and \cref{diff:sec:phasecrit} measures the gap exactly.

\begin{convention}[What ``a differential equation over $\SC$'' means]
\label{diff:conv:meaning}
In \cref{diff:part:scalar,diff:part:workspaces}, a differential equation over
$\SC$ is an algebraic equation in elements of $\SC$ and their iterated
$\der$-images. It is \emph{not} an equation for a class function of an
independent variable. \Cref{diff:part:coordinate} treats the latter, under a
different operator and with a different constant field, and
\cref{diff:sec:taxonomy} keeps them apart.
\end{convention}

Every real coefficient is a logarithmic derivative in the real field.

\begin{proposition}[Real exponential integration]
\label{diff:prop:realexp}
For every $p\in\No$, the element $s=\exp(\Izero p)$ of
\cref{diff:prop:normprimitive} satisfies $\der s=ps$ and $s>0$. More generally,
for $a,b\in\No$ the equation $\der y=ay+b$ has surreal solutions, and with
$h=\Izero a$ all real solutions are
$y=e^{h}\bigl(c+\Izero(e^{-h}b)\bigr)$, $c\in\R$.
\end{proposition}
\begin{proof}
Apply \ref{diff:BM:exp} to a primitive of $p$, respectively of $a$; the
product rule verifies the displayed formula, and the quotient of the
difference of two solutions by $e^{h}$ has derivative zero.
\end{proof}

Notice the two distinct uses of structure: \ref{diff:BM:surj} supplies the
primitive, and \ref{diff:BM:exp} supplies the exponential. Knowing only that
$\der$ is onto would not justify the second step. The difficulty in the
complex case is therefore not additive integration; it is the absence of a
suitable nonzero integrating factor, and \cref{diff:sec:phasecrit} says exactly
when one exists.

\subsection{Smallness, derived, with strict inclusion}
\label{diff:sub:smallness}

\begin{lemma}[Order reversal on infinitesimals, and smallness]
\label{diff:lem:smallnegative}
If $0<\varepsilon\in\mm$ then $\der\varepsilon<0$. Consequently $\der$ is
strictly order-reversing on $\mm$, and
\[
 \der\Ofin=\der\mm\subseteq\mm .
\]
\end{lemma}
\begin{proof}
Since $\varepsilon$ is a positive infinitesimal, $1/\varepsilon>\R$, so
\ref{diff:BM:positivity} and \cref{diff:eq:basicder} give
\[
 0<\der(1/\varepsilon)=-\frac{\der\varepsilon}{\varepsilon^{2}},
\]
whence $\der\varepsilon<0$. For every positive integer $n$, the element
$\omega+n\varepsilon$ is positive infinite, so \ref{diff:BM:positivity} and
\ref{diff:BM:norm} give $1+n\der\varepsilon>0$. Therefore
$-1/n<\der\varepsilon<0$ for every $n$, that is $\der\varepsilon\in\mm$.
Negative infinitesimals follow by additivity and $0$ is immediate. Strict
order reversal follows by applying the first assertion to the positive
infinitesimal difference of two infinitesimals. Finally every finite element
is $r+\varepsilon$ with $r\in\R$ and $\varepsilon\in\mm$, and $\der r=0$, so
$\der\Ofin=\der\mm$.
\end{proof}

\begin{remark}[An independent route to smallness, and a second route to the containment]
\label{diff:rem:sqrtsmallness}
One source manuscript instead \emph{imports} the finite form
$\der\Ofin\subseteq\Ofin$ --- which is what the argument needs --- and
strengthens it to $\der\Ofin=\der\mm\subseteq\mm$ by a square-root trick: a
positive $\varepsilon\in\mm$ is $q^{2}$ with $q=\sqrt\varepsilon\in\mm$, so
$\der q\in\Ofin$ by the imported smallness and
$\der\varepsilon=2q\,\der q\in\mm$. That argument is correct and is recorded
here as an alternative, but it starts from the containment rather than proving
it, and it does not give the order reversal. A third formulation of the same
one-line proof, used by three of the seven, applies
\ref{diff:BM:positivity} to $\omega\pm nf$ for $f\in\Ofin$, obtaining
$1\pm n\der f>0$ directly; a fourth applies it to $\omega/n\pm f$. All are the
same argument, and none uses a mean-value theorem, a topological limit, or an
estimate about numerical rounding. Smallness is a structural consequence of the displayed positivity axiom
for this derivation.
\end{remark}

\begin{warning}[The containment is strict, and that gap is the whole subject]
\label{diff:warn:strict}
$\der\mm\subseteq\mm$ is a \emph{strict} inclusion. Indeed
$\der\log\omega=1/\omega$, and every primitive of $1/\omega$ differs from
$\log\omega$ by a real constant, so every primitive of $1/\omega$ is infinite
and $1/\omega\notin\der\mm$. On the other hand
$\omega^{-2}=\der(-\omega^{-1})\in\der\mm$. Testing only whether an imaginary
coefficient is infinitesimal therefore loses the decisive information. The
entire obstruction theory of \cref{diff:part:scalar} is the measurement of this
gap, and \cref{diff:warn:collapse} states what goes wrong if it is collapsed.
\end{warning}

A first, coarse consequence is already enough to kill a constant nonreal rate.

\begin{proposition}[A preliminary obstruction]
\label{diff:prop:coarse}
For every $z\in\SC^{\times}$, the imaginary part of $\der z/z$ is
infinitesimal. Hence for $\lambda\in\C\setminus\R$,
\[
 \der z=\lambda z\quad\Longrightarrow\quad z=0 .
\]
\end{proposition}
\begin{proof}
Multiplying $z$ by $1$ or by $i$ does not change $\der z/z$, so we may write
$z=u(1+ih)$ with $u\in\No^{\times}$ and $h\in\Ofin$, taking as $u$ a nonzero
coordinate of maximal absolute value. Then
\[
 \frac{\der z}{z}=\frac{\der u}{u}+\frac{i\der h}{1+ih}
 =\frac{\der u}{u}+\frac{h\der h}{1+h^{2}}+i\frac{\der h}{1+h^{2}} .
\]
By \cref{diff:lem:smallnegative}, $\der h$ is infinitesimal, and division by
$1+h^{2}\geq1$ preserves that. A nonzero ordinary imaginary part is not
infinitesimal.
\end{proof}

This already rules out a surcomplex value behaving as $e^{i\omega}$ under
$\der$. It does \emph{not} yet classify infinitesimal rates: the exact image
is strictly smaller than $\No+i\mm$, and identifying it needs the finite-angle
structure of \cref{diff:sec:phasegeometry}.

\section{One derivation, and seven other operators}
\label{diff:sec:taxonomy}

Almost every apparent paradox in this subject is a type error. The seven
manuscripts between them separate eight operators; the merged report keeps all
eight distinct, with the symbols of \cref{diff:tab:derivatives}. A formula can
be correct for one and false for another with no contradiction, and several
theorems below are true for one and false for another \emph{on the same printed
equation}.

\begin{table}[htbp]
\centering\small
\caption{The eight operators. Only the first is ``the derivation''.}
\label{diff:tab:derivatives}
\begin{tabularx}{\textwidth}{@{}L{0.14\textwidth}YL{0.20\textwidth}L{0.19\textwidth}@{}}
\toprule
Symbol & What it does & What it holds fixed & Sample value \\
\midrule
$\der$ &
The intrinsic Berarducci--Mantova derivation on a \emph{number} of $\No$ or
$\SC$ &
Exactly $\R$, or $\C$ after complexification &
$\der\omega=1$, $\der t=-t^{2}$, $\der(t^{3})=-3t^{4}$ \\[0.35em]
$\der_{s}$ &
The rescaled derivation $\der_{s}f=\der f/\der s$, for $\der s\neq0$ &
Same kernel; but not the same normalization &
$\der_{s}s=1$; for $s=\log\omega$, $\der s=t$ \\[0.35em]
$\derZ$ &
Formal differentiation of an indeterminate $Z$ on $\SC[[Z]]$ &
Every coefficient in $\SC$, however infinite &
$\derZ(\omega Z)=\omega$, $\derZ\omega=0$ \\[0.35em]
$\dertil$ &
Coefficientwise \emph{intrinsic} differentiation on $\SC[[Z]]$,
$\dertil\sum a_{n}Z^{n}=\sum(\der a_{n})Z^{n}$ &
The formal indeterminate $Z$ &
$\dertil(\omega Z)=Z$ \\[0.35em]
$\Dz$ &
Coordinate differentiation of the holomorphic coefficient functions of a
common-domain Hahn family, in an ordinary complex $z$ &
Every Hahn monomial and every scalar parameter &
$\Dz(t^{3})=0$, $\Dz\omega=0$, $\Dz z=1$ \\[0.35em]
$\dercoef$ &
Intrinsic differentiation of the Hahn \emph{monomials} of a common-domain
family &
The ordinary variable $z$ &
$\dercoef(\omega z)=z$, $\dercoef(t^{2}e^{z})=-2t^{3}e^{z}$ \\[0.35em]
$D_{t}$ &
Formal differentiation in a Hahn parameter treated as an indeterminate &
The coefficients &
$D_{t}(t^{3})=3t^{2}$ \\[0.35em]
$F'_{\mathrm{fine}}$ &
The pointwise fine order-field derivative of a class \emph{function}, an
$\epsilon$--$\delta$ difference quotient in the actual field variable &
Determined by the function and its topology; not a field derivation &
$\fin'_{\mathrm{fine}}\equiv1$, $\pp'_{\mathrm{fine}}\equiv0$ \\
\bottomrule
\end{tabularx}
\end{table}

\begin{warning}[Seven separators, each of which has been mistaken for a paradox]
\label{diff:warn:separators}
\begin{enumerate}[label=(\arabic*),leftmargin=2.4em]
\item\emph{$t$ is a number, not a symbol.} With $t=\omega^{-1}$ one has
$\Dz(t^{3})=0$, $D_{t}(t^{3})=3t^{2}$ and $\der(t^{3})=-3t^{4}$, all correct.
Treating $t$ as a formal constant is a coefficient operation; substituting
$t=\omega^{-1}$ and applying the number derivation gives $\der t=-t^{2}$, not
$0$. Replacing $D_{t}$ by $\der$ in a displayed ordinary differential equation
without changing the coefficient functions usually changes the equation. The
notation $D_{t}$ becomes related to $\der$ only after an embedding is specified
and the formal operations are verified to commute with its interpretation.
\item\emph{The constant function is not the number.} $F(Z)=\omega$ has
$\derZ F=0$, while the number $\omega$ has $\der\omega=1$. For $f(Z)=e^{Z}$ at
a real point $r$, $f'(r)=e^{r}$ but $\der(f(r))=\der(e^{r})=0$.
\item\emph{A fine derivative does not determine an intrinsic one.} For every
$x\in\No$ and every finite $h$, exactly,
\[
 \fin(x+h)=\fin(x)+h,\qquad \pp(x+h)=\pp(x),
\]
so the fine derivatives of $x\mapsto\fin(x)$ and $x\mapsto\pp(x)$ are $1$ and
$0$ everywhere, with exact difference quotients for all $0<\abs h<1$ and no
limit subtlety. Yet $\der(\fin(\omega))=\der0=0$ while
$\der(\pp(\omega))=\der\omega=1$. Hence a blanket identity
$\der(F(x))=F'_{\mathrm{fine}}(x)\der x$ is \emph{false} for arbitrary
fine-differentiable $F$; the missing hypothesis is a compatibility structure
for $F$, and fine differentiability is not one. The same example explains why a
zero fine derivative need not force a class function to be constant, whereas
$\der f=0$ for a \emph{scalar} means exactly $f\in\C$.
\item\emph{A second fine-derivative separator, with a phase.} Let
$H(x)=\cis(\fin x)$, using the finite-angle map of
\cref{diff:def:cis} on finite angles only. Then $H(\omega)=1$ and, for
infinitesimal $h$, $H(\omega+h)=\Emm(ih)=1+ih+O(h^{2})$ as an ordered-field
estimate, so $H'_{\mathrm{fine}}(\omega)=i$. Nevertheless
\[
 \der(H(\omega))=\der 1=0\ \neq\ H'_{\mathrm{fine}}(\omega)\,\der\omega=i .
\]
A locally valid Taylor rule grants no universal scalar chain rule for class
functions.
\item\emph{Omitting a variable, not using the wrong rule.} For
$F(z,t)=tz$ at $z=\omega$, $t=\omega^{-1}$, the value is $1$ and
\[
 \der F=t\,\der\omega+\omega\,\der t=t-\omega t^{2}=0,
\]
whereas the coordinate derivative $\Dz F=t$ alone would return $t$. Here the
ordinary multivariable chain rule is valid; what is wrong is dropping a
variable. The same point in one variable: for $F(Z)=tZ$ and $x=\omega$,
$\der(F(\omega))=(-t^{2})\omega+t=0$, as it must be.
\item\emph{A coherent family is not a number.} For $F(z)=tz$ with $z$ a fixed
ordinary complex number, $\Dz F=t$ while $\der(F(z))=-t^{2}z$, because that
fixed ordinary $z$ is a constant of $\der$.
\item\emph{The same printed equation, twice.} $\Dz F=iF$, $F(0)=1$ has the
ordinary entire solution $e^{iz}$, a single-coefficient member of a
common-domain Hahn ring. $\der y=iy$ has no nonzero solution in $\SC$
(\cref{diff:cor:oscillator}). There is no contradiction: one statement is
about a holomorphic function of an ordinary coordinate, the other about a
scalar element of a fixed differential field. The nonexistent bridge would be
to evaluate the entire function at $z=\omega$, call the result a surcomplex
scalar, and assume its scalar derivative comes from the coordinate chain rule.
Neither the fixed-domain Hahn representation nor the ordinary initial-value
theorem supplies that bridge.
\end{enumerate}
\end{warning}

\begin{proposition}[Change of independent scale changes the solvability data]
\label{diff:prop:rescale}
Let $s\in\No$ with $\der s\neq0$. Then $\der_{s}f=\der f/\der s$ is a
derivation with $\der_{s}s=1$, and
\[
 \der_{s}y=qy \quad\Longleftrightarrow\quad \der y=(\der s)q\,y .
\]
Hence the phase test of \cref{diff:thm:phasecriterion} must be applied to
$\Imm\bigl((\der s)q\bigr)$, not to $\Imm q$. With $s=\log\omega$, so
$\der s=t$: the equation $\der_{s}y=iy$ is the forbidden $\der y=ity$, whereas
$\der_{s}y=ity$ becomes $\der y=it^{2}y$ and \emph{is} solvable.
\end{proposition}

The rescaled operator preserves the field but not the normalization at
$\omega$, nor necessarily the same order conventions. A backend cannot change
the independent scale while leaving differential solvability metadata
untouched. At the scale $s=\tau=\log t=-\log\omega$ the rescaled derivation is
the Euler derivation $\derT=-\omega\der$ of \cref{diff:rs:part}, which acts on
real-exponent Hahn series by $t^{a}\mapsto a\,t^{a}$ and reverses the sign rule
\ref{diff:BM:positivity} (\cref{diff:rs:lem:euler,diff:rs:warn:orientation}).

\begin{proposition}[The two formal derivations commute; the valid polynomial chain rule]
\label{diff:prop:totalchain}
On $\SC[[Z]]$ the derivations $\derZ$ and $\dertil$ commute, and
$\dertil+u\derZ$ is a derivation for every coefficient scalar $u$; commutation
with that sum need not hold when $\der u\neq0$. For a polynomial
$F(Z)=\sum_{k=0}^{n}a_{k}Z^{k}\in\SC[Z]$ and $x\in\SC$,
\begin{equation}\label{diff:eq:totalchain}
 \der\bigl(F(x)\bigr)=(\dertil F)(x)+(\derZ F)(x)\,\der x .
\end{equation}
The same identity holds for rational functions at arguments where the
denominator is nonzero.
\end{proposition}
\begin{proof}
Commutation and the derivation property are coefficientwise identities using
only finite convolution sums at each degree. For \cref{diff:eq:totalchain},
apply the finite product rule to each $a_{k}x^{k}$, obtaining
$(\der a_{k})x^{k}+ka_{k}x^{k-1}\der x$, and add. Both the coefficient and the
variable derivative obey the quotient rule, which gives the rational case.
\end{proof}

\begin{example}[Two worked checks of \cref{diff:eq:totalchain}]
\label{diff:ex:totalchainchecks}
For $F(Z)=\omega Z^{2}+\omega^{-1}$ at $x=\omega$: $F(x)=\omega^{3}+\omega^{-1}$
and $\der(F(x))=3\omega^{2}-\omega^{-2}$, whereas the variable derivative
alone gives only $2\omega^{2}$; the missing term is the derivative of the
coefficients after evaluation. For $F(Z)=\omega Z+tZ^{2}$ and $h=\omega^{2}$:
$F(h)=2\omega^{3}$ so $\der(F(h))=6\omega^{2}$, while on the right
$(\dertil F)(h)=h-t^{2}h^{2}=0$ and
$(\derZ F)(h)\der h=(\omega+2th)(2\omega)=6\omega^{2}$. The cancellation in
the coefficient term is a feature of that example, not a licence to omit
coefficient differentiation.
\end{example}

An infinite series needs more, and the extra hypothesis can be made minimal.

\begin{proposition}[Two-derivative chain rule for series, with derived summability]
\label{diff:prop:chain}
Let $F(Z)=\sum_{n\geq0}a_{n}Z^{n}$ with $a_{n}\in\SC$, and let $h\in\SC$.
If the \emph{evaluated} family $(a_{n}h^{n})_{n\geq0}$ is summable, then the
two families defining $(\dertil F)(h)$ and $(\derZ F)(h)$ are summable as
well, and
\[
 \der(F(h))=(\dertil F)(h)+(\derZ F)(h)\,\der h .
\]
\end{proposition}
\begin{proof}
Assume first $h\neq0$. Multiplying each member of a summable family by its
ordinary integer index preserves summability, since it enlarges no support;
multiplying the whole family by the fixed scalar $h^{-1}$ preserves Hahn
summability. Hence
$(na_{n}h^{n-1})_{n\geq1}=(nh^{-1}(a_{n}h^{n}))_{n\geq1}$ is summable, and so
is its multiple by the fixed $\der h$. By \ref{diff:BM:additive} the family
$(\der(a_{n}h^{n}))_{n}$ is summable. Subtracting the already controlled
variable-derivative family shows $((\der a_{n})h^{n})_{n}$ is summable. The
termwise product rule may now be summed. If $h=0$, the evaluated series and
both formal-derivative evaluations have at most one nonzero term each, and the
identity reduces to $\der a_{0}=\der a_{0}+a_{1}\der 0$.
\end{proof}

\begin{remark}[Why this is stronger than three summability certificates]
\label{diff:rem:chainstrength}
An arbitrary sum of two families can be summable through cancellation even
when neither summand family is (\cref{diff:warn:splitting}). The special
power-series form is what lets the variable-derivative family be controlled
\emph{first}, after which the split is justified. So
\cref{diff:prop:chain} assumes one certificate and derives two, which is
strictly stronger than assuming three independent ones. It still does not
assert that an arbitrary formal series is evaluable at an arbitrary scalar; see
\cref{diff:sec:formalexistence} for the relevant separations and for the
counterexample of \cref{diff:rem:coefstream}.
\end{remark}

\begin{remark}[No general composition is assumed]
\label{diff:rem:nocomposition}
No global operation ``replace $\omega$ by $x$ in every surreal'' is ever used.
Berarducci and Mantova construct a composition $f\circ x$ for $f$ in the
$\omega$-series class and positive infinite $x$, identify $(\der f)\circ x$
with the fine derivative of $x\mapsto f\circ x$, give a compatible chain rule
inside that setting, and prove a Taylor theorem there
\cite[Cor.\ 7.6--7.7, Thms 7.14, 8.2]{BMGerms}. They also prove that the
simplest derivation on \emph{all} of $\No$ admits no composition satisfying
their composition and compatibility axioms
\cite[Thm 8.4]{BMGerms}. \Cref{diff:sec:composition} records both results with
their hypotheses intact and explains the exact sense in which the two
derivatives then agree. All differentiation in this report is either the fixed
$\der$, finite algebra, or a summable infinitesimal formal substitution whose
validity is proved.
\end{remark}

\section{The central hazard: three objects called ``phase''}
\label{diff:sec:hazard}

\subsection{Three names, three symbols}
\label{diff:sub:threephases}

Three different objects occur in the source manuscripts under the single word
``phase''. Two of them are angles; the third is not an angle of anything, and
does not even live in $\Ofin$. Merging them --- in a sentence, in a word, or in
a symbol --- produces false statements. This report therefore fixes the
following vocabulary once and uses it without exception.

\begin{convention}[The three phases]
\label{diff:conv:phases}
\begin{enumerate}[label=\textup{(\alph*)},leftmargin=2.4em]
\item The \emph{finite angle} $\theta$. The map $\cis$ of
\cref{diff:def:cis} sends the \emph{finite} surreals $\Ofin$ onto the unit
circle $\Tor(\No)$ with kernel $2\pi\Z$
(\cref{diff:thm:cisgroup}). Its argument $\theta$ is the whole angle, and it is
always finite: $\theta\in\Ofin$. No value of a circular function at an infinite
angle is ever used.
\item The \emph{infinitesimal phase residue} $\eta$. Every unit
$u\in\Tor(\No)$ factors \emph{uniquely} as $u=c\,\Emm(i\eta)$ with $c$ an
ordinary unit complex number and $\eta\in\mm$
(\cref{diff:thm:unitphase}). The constant $c$ carries all of the ordinary
phase, $\der c=0$, and $\eta$ carries all of the variable intrinsic phase:
$u^{\dagger}=i\der\eta$. The two are related to (a) by
$\theta=\st(\theta)+\eta$ and $c=e^{i\st(\theta)}$, so (b) is the
\emph{infinitesimal part} of (a), not a synonym for it.
\item The \emph{accumulated phase} $B$. For a real coefficient $b$, an
accumulated phase is any $B\in\No$ with $\der B=b$; the normalized choice is
$\Izero(b)$. Any two differ by a real constant. $B$ is an arbitrary surreal.
It need not be finite, it is not the angle of any surcomplex number, and when
it is infinite it lies outside $\Ofin$ entirely, with its purely infinite part
$\pp(B)\in\Pcl$ --- the \emph{obstruction} --- nonzero. The whole theory of
\cref{diff:part:scalar} is the statement that $\der y=(a+ib)y$ is solvable
exactly when the accumulated phase of $b$ is finite, that is $B\in\Ofin$.
Only then does $B$ acquire an angular reading at all: a finite $B$ is the
argument of $\cis B$ in sense~(a), with $\eta=B-\st(B)$ its infinitesimal
residue in sense~(b).  An infinite $B$ is the argument of nothing.  So (c) is
not a third kind of angle; it is the object whose \emph{finiteness} is what
sense~(a) requires, and the criterion is exactly the assertion that it can
fail.
\item \textbf{A missing oscillatory generator.}  In
\cref{diff:sub:torus} and in the Picard--Vessiot sections the phrase
\emph{independent phases} counts something else again: the number of
transcendental generators $\mathsf T$ that must be adjoined \emph{outside}
$\SC$ before a diagonal system has a full solution basis.  That count is
$d=m-\operatorname{rank}\Lambda$, a rank of a lattice of obstructions, and it
is not a count of angles, residues or accumulated phases.  Wherever this report
writes ``independent phases'' or ``phase generators'', that fourth sense is
meant, and it is never any of (a)--(c).
\item The \emph{differential phase spectrum} $\Ph(A)$. For a finite matrix
system $\der y=Ay$ over $\SC$, \cref{diff:thm:matnormal} attaches a
\emph{multiset} $\Ph(A)=\{P_{1},\dots,P_{n}\}$ of purely infinite surreals: the
accumulated obstructions of the diagonal entries of its normal form. Its
members are objects of sense~(c) --- more precisely, values of $\Obs$ --- so
$\Ph$ adds no new kind of phase; what is new is that there are $n$ of them at
once and that the multiset, not its sum, is the complete invariant
(\cref{diff:cor:mattrace}). The eighth source manuscript writes $\Psi$ for the
map $b\mapsto\pp(B)$ with $\der B=b$. That map is exactly $\Obs$ of
\cref{diff:def:obs}, so this report keeps the single symbol $\Obs$ for it and
does not adopt $\Psi$, in accordance with \cref{diff:conv:noJ}. The adoption
would in any case have collided: the letter $\Psi$ already occurs once in this
report, in \cref{diff:cor:noglobalPsi}, as a local placeholder for a putative
derivation-compatible exponential $\SC\to\SC^{\times}$ --- a map that is there
proved not to exist, and that is not a phase map of any of the senses
(a)--(e).
\end{enumerate}
\end{convention}

\begin{warning}[``Spectrum'' names three inequivalent things across the collection]
\label{diff:warn:spectrum}
The differential phase spectrum $\Ph(A)$ of sense~(e) is a multiset of
\emph{purely infinite surreal numbers}. It is not a list of eigenvalues: the
companion report \emph{Surcomplex Spectral Theory} \cite{RepoSpectral} uses
``spectrum'' for the algebraic eigenvalue and singular-value lists of a finite
matrix and their valuations, which are ordinary surcomplex numbers of the
field, and \cref{diff:thm:matspectral} is precisely the theorem that relates
the two --- by integrating and taking purely infinite parts, not by identifying
them. \Cref{diff:ex:matnonnormal} shows that for a nonnormal coefficient the
two are unrelated. Nor is $\Ph(A)$ a set of angles: no member of a phase
spectrum is ever routed through the finite-angle map $\cis$ of sense~(a) or
through the finite-angle constructions of \cite{RepoTrig}, since a nonzero
member of $\Pcl$ is infinite and $\cis$ is defined on $\Ofin$ only. And it is
not a prime spectrum of a ring: where the collection's analytic-geometry
report discusses $\operatorname{Spec}$ of a coefficient ring, that is the
ring-theoretic object and is written out as \emph{prime spectrum} there.
\end{warning}

\begin{table}[htbp]
\centering\small
\caption{The phase objects, and where each source's notation lands.}
\label{diff:tab:phases}
\begin{tabularx}{\textwidth}{@{}L{0.16\textwidth}L{0.10\textwidth}L{0.13\textwidth}YL{0.20\textwidth}@{}}
\toprule
Name here & Symbol & Lives in & What it is & Source notations \\
\midrule
Finite angle & $\theta$ & $\Ofin$ &
The argument of $\cis$; the \emph{whole} angle of a unit, finite by theorem &
$\theta$, $y$ \\[0.3em]
Infinitesimal phase residue & $\eta$ & $\mm$ &
What is left of a unit after factoring out the constant ordinary unit complex
number $c$ &
$\eta$, $\theta$, $\varepsilon$, $\delta$ \\[0.3em]
Accumulated phase & $B$ & $\No$ &
A primitive of the imaginary coefficient; \emph{not} an angle; obstructed
exactly when $\pp(B)\neq0$ &
$B$; $\Izero(b)$; and four different letters for the primitive operator \\[0.3em]
Differential phase spectrum & $\Ph(A)$ & multisets in $\Pcl$ &
The $n$ accumulated obstructions of a finite matrix system, as a multiset; the
complete invariant of differential equivalence &
$\operatorname{Ph}$; $\Psi(\lambda_{j})$, where the map $\Psi$ is this
report's $\Obs$; member~12's ``phase labels'' $b$, which are ordinary real
residual frequencies with $\Ph=\{b\log t\}$ (\cref{diff:rs:warn:frequencies}) \\
\bottomrule
\end{tabularx}
\end{table}

\begin{warning}[Where the senses collide in the sources]
\label{diff:warn:phasecollision}
Two of the first seven manuscripts use ``phase'' for sense (a), four use it for
sense (b), and all seven use ``accumulated phase'' for sense (c); the eighth
introduces sense (e) and, with it, a fourth letter for the primitive
obstruction. One of them prints
(a) and (b) twenty lines apart, both called the phase. Sense (c) is the one
that carries the obstruction, and it is the one that is \emph{not} an angle:
$\Phi(i)=\omega$ and $\Phi(i/\omega)=\log\omega$ are purely infinite surreals,
not arguments of anything. A merged sentence such as ``the phase of the
solution is $\log\omega$'' is meaningless; the correct sentence is ``the
accumulated phase of the coefficient $i/\omega$ is $\log\omega$, which is
infinite, so there is no solution''. Similarly, ``every surcomplex number has
a finite phase'' is a theorem about (a); ``the variable phase of a unit is
infinitesimal'' is a theorem about (b); and neither implies anything about (c)
without \cref{diff:thm:phasecriterion}. Two later members add two further uses.
Member~12 calls the ordinary real numbers $b=\Imm\lambda$ ``phase labels''; they
are none of (a)--(e), this report calls them residual frequencies, and the
frequency $b$ corresponds to the phase $b\log t$ of sense (e)
(\cref{diff:rs:warn:frequencies}). Member~11 uses the word only to cite this
report, and its one ``finite angle'' is sense (a); its abelian logarithmic
derivative is a tangent vector, not an angle (\cref{diff:aut:sub:image}).
Member~13 uses ``phase obstruction'' for this report's criterion and writes
$\theta$ for a real infinitesimal primitive of the angular rate of a unit; that
is sense (b), written $\eta$ here (\cref{diff:cp:conv:alphabet}).
\end{warning}

\subsection{Two load-bearing classes that differ by one operator}
\label{diff:sub:twoclasses}

The second half of the hazard is not a word but a symbol. Four of the seven
manuscripts print, within a few pages of each other, the two classes
\[
 \No+i\Acl \qquad\text{and}\qquad \Strip=\No+i\Ofin,
\]
where $\Acl=\der\mm$. The first is the exact image of the logarithmic
derivative (\cref{diff:thm:logimage}); the second is the maximal domain of a
derivation-compatible exponential (\cref{diff:thm:maximalstrip}). They are
different classes, and the difference is one application of $\der$.

\begin{warning}[Do not collapse $\der\Ofin$ to $\Ofin$]
\label{diff:warn:collapse}
$\Acl=\der\mm$ is a \emph{strict} subclass of $\mm$
(\cref{diff:warn:strict}), and $\mm$ is a strict subclass of $\Ofin$. If a
merge collapses $\der\Ofin$ to $\Ofin$, the rank-one criterion becomes ``$b$
finite''. Even the stronger replacement ``$b$ infinitesimal'', obtained by
replacing $\Acl$ with $\mm$, would make $\der y=i\omega^{-1}y$ solvable.
Every one of the seven manuscripts refutes that conclusion, and the counterexample
pair of \cref{diff:ex:logseparator} --- $i/(\omega\log\omega)$ insoluble against
$i/(\omega(\log\omega)^{2})$ soluble --- exists precisely to block it. The
strip $\Strip=\No+i\Ofin$ is $\der$-stable \emph{because}
$\der\Ofin\subseteq\mm$; that is a consequence of the difference between the
two classes, not evidence that they are the same.
\end{warning}

\begin{warning}[$\No+i\der\Ofin$ and $\No+i\der\mm$ are equal only after a lemma]
\label{diff:warn:twoimages}
The sources print the logarithmic-derivative image in both forms. They agree,
but only because $\der\Ofin=\der\mm$, which is a \emph{theorem}
(\cref{diff:lem:smallnegative}, with the independent square-root route of
\cref{diff:rem:sqrtsmallness}). The two forms must never be presented as
notational variants: one is the other after a proved step. This report uses
the single symbol $\Acl$ for the common value and states the lemma before using
it.
\end{warning}

\subsection{The letter \texorpdfstring{$J$}{J}, and the word ``obstruction''}
\label{diff:sub:letterJ}

Across the source batch the letter $J$ denotes three unrelated objects: in one
manuscript it is the normalized intrinsic primitive on $\No$; in another it is
not an operator at all but the representative \emph{class} of purely infinite
surreals; and in a companion report that remains standalone --- the rank-one
Berkovich manuscript, on analytic functions on surcomplex disks and annuli ---
it is the zero-constant \emph{coordinate} primitive of a convergent power
series over a fixed coefficient field, whose ``complete primitive obstruction''
is an annular Laurent residue and has nothing to do with the phase obstruction
here. Two further manuscripts write $J$ for the inverse of $\der|_{\mm}$.

\begin{convention}[No letter $J$; one symbol per object]
\label{diff:conv:noJ}
This report uses no symbol $J$ for any operator or class.  The letter
occurs only as a neutral index set for set-indexed families
(\cref{diff:def:summable} and \cref{diff:thm:BM}), where no confusion with
the objects below is possible.  The normalized intrinsic primitive is
$\Izero$ (\cref{diff:prop:normprimitive}), with values in
$\Pcl\oplus\mm$, the zero-constant-term surreals. The purely infinite
class $\Pcl$ is the target of the obstruction map, not of $\Izero$.
The inverse of $\der|_{\mm}$ is written
$\Izero|_{\Acl}$ and named the \emph{unique infinitesimal primitive}. The
obstruction on $\No$ is $\Obs$ and its surcomplex companion is $\Phi$
(\cref{diff:def:obs}). Cross-references to companion reports name the operator
rather than the letter.
\end{convention}

\begin{remark}[The obstruction does not depend on the normalization]
\label{diff:rem:obsnormalization}
Adding a real constant to a primitive does not change its purely infinite
part, so $\Obs$ is insensitive to the normalization of $\Izero$. That
insensitivity is what makes $\Obs$ well defined, and it is why the
normalization is a convenience rather than a hypothesis. It is also why a
``primitive'' defined only up to $\R$, or selected by a choice function, is a
\emph{different} object: the first would not give a well-defined primitive
operator, and the second would not give a class function without a choice
principle. $\Izero$ avoids both by singling out a unique element; see
\cref{diff:prop:normprimitive} and \cref{diff:warn:Jnotalgorithm}.
\end{remark}
\part{The scalar theory over \texorpdfstring{$\SC$}{No[i]}}
\label{diff:part:scalar}

\section{Infinitesimal exponential, finite angle, polar form}
\label{diff:sec:phasegeometry}

\subsection{The infinitesimal exponential and logarithm}
\label{diff:sub:smallexp}

For $h\in\mm_{\C}$ define the Hahn sums
\begin{equation}\label{diff:eq:explog}
 \Emm(h)=\sum_{n\geq0}\frac{h^{n}}{n!},
 \qquad
 \Lmm(1+h)=\sum_{n\geq1}\frac{(-1)^{n+1}}{n}h^{n}.
\end{equation}
\Cref{diff:lem:neumann} makes both families summable, because the support of a
nonzero infinitesimal is positive and well ordered. These are defined by strong Hahn summation, without a topological limit. Their domain is
\emph{infinitesimal arguments only}: no attempt is made to sum the ordinary
exponential Taylor series at a nonzero ordinary constant, still less at an
infinite argument.

\begin{lemma}[Formal identities survive infinitesimal evaluation]
\label{diff:lem:smallchain}
The maps in \cref{diff:eq:explog} are inverse group isomorphisms
\[
 (\mm_{\C},+)\ \underset{\Lmm}{\overset{\Emm}{\rightleftarrows}}\
 (1+\mm_{\C},\cdot),
\]
they commute with conjugation, and
\begin{equation}\label{diff:eq:diffexplog}
 \der\Emm(h)=\Emm(h)\,\der h,
 \qquad
 \der\Lmm(w)=\frac{\der w}{w}\quad(w\in1+\mm_{\C}).
\end{equation}
Moreover $\Emm$ is injective on $\mm_{\C}$: for $h\in\mm_{\C}$,
\begin{equation}\label{diff:eq:expinjective}
 \Emm(h)=1\quad\Longleftrightarrow\quad h=0 .
\end{equation}
More generally, for any $f(X)=\sum_{n\geq0}c_{n}X^{n}$ with
\emph{ordinary constant} coefficients $c_{n}\in\C$ and any $h\in\mm_{\C}$, the
family $(c_{n}h^{n})_{n}$ is summable, its sum is denoted $f(h)$, and
\begin{equation}\label{diff:eq:formalchain}
 \der\bigl(f(h)\bigr)=f'_{\mathrm{formal}}(h)\,\der h .
\end{equation}
\end{lemma}
\begin{proof}
The formal identities $\exp(U+V)=\exp(U)\exp(V)$, $\log(\exp U)=U$ and
$\exp(\log(1+U))=1+U$ hold over the ordinary characteristic-zero coefficient
field. After substitution at infinitesimals, \cref{diff:lem:neumann} makes
every coefficient calculation a finite one at each fixed exponent, which
justifies the products and regroupings; this is a finite check per exponent,
not an appeal to a limit. Conjugation acts on coefficients and preserves
supports.

For \cref{diff:eq:diffexplog} and \cref{diff:eq:formalchain}, the finite
product rule gives $\der(h^{n})=nh^{n-1}\der h$, and \ref{diff:BM:additive}
permits termwise differentiation, including the assertion that the
differentiated family is summable. The family
$(nc_{n}h^{n-1})_{n\geq1}$ is separately summable by \cref{diff:lem:neumann},
and multiplication by the single fixed scalar $\der h$ preserves summability;
so the factorisation
\[
 \der\Emm(h)=\der h\sum_{m\geq0}\frac{h^{m}}{m!}
\]
and, for the logarithm, $\der h\sum_{n\geq1}(-1)^{n+1}h^{n-1}=\der h/(1+h)$,
are both legitimate. For \cref{diff:eq:expinjective}: if $h\neq0$, every power
$h^{n}$ with $n\geq2$ has valuation strictly larger than that of $h$, so the
leading term of $\Emm(h)-1$ is the leading term of $h$.
\end{proof}

\begin{warning}[Three limits on \cref{diff:lem:smallchain}]
\label{diff:warn:smallexplimits}
(i) The injectivity \cref{diff:eq:expinjective} is \emph{local} to $\mm_{\C}$
and asserts nothing about an injective exponential on the whole class $\SC$.
(ii) The argument does \emph{not} extend to a Taylor series at an arbitrary
infinite argument: the original family need not be summable there, and
\cref{diff:rem:nonsummableexp} exhibits the failure. (iii) The evaluation rule
$f(h)$ requires \emph{constant} coefficients. A formal series with arbitrary
coefficients in $\SC$ need not be summable at any given nonzero infinitesimal:
$\sum_{n\geq0}\omega^{n^{2}}X^{n}$ is not summable at $X=\omega^{-1}$, because
its monomials $\omega^{n^{2}-n}$ are not reverse well ordered. The
infinitesimal chain rule must not be exported to arbitrary
$\SC$-coefficient streams without a joint support certificate
(\cref{diff:rem:coefstream}). No lower bound on an ordinary radius of
convergence is needed, however: factorially growing complex coefficients
remain coefficients at the same Hahn monomials.
\end{warning}

\subsection{The finite angle}
\label{diff:sub:cis}

Let
\[
 \Tor(\No)=\{u\in\SC:u\bar u=1\},
 \qquad
 \Tor(\C)=\{c\in\C:\abs c=1\}.
\]

\begin{theorem}[Unique infinitesimal phase residue]
\label{diff:thm:unitphase}
Every $u\in\Tor(\No)$ has a unique representation
\begin{equation}\label{diff:eq:unitphase}
 u=c\,\Emm(i\eta),\qquad c\in\Tor(\C),\quad\eta\in\mm,
\end{equation}
namely $c=\st(u)$ and $i\eta=\Lmm(u/c)$; and
\begin{equation}\label{diff:eq:unitlogder}
 u^{\dagger}:=\frac{\der u}{u}=i\,\der\eta .
\end{equation}
Consequently $\Tor(\No)\cong\Tor(\C)\times(\mm,+)$ as abelian groups.
\end{theorem}
\begin{proof}
Write $u=a+ib$ with $a^{2}+b^{2}=1$; both coordinates are finite, so
$c=\st(u)$ exists and $c\bar c=1$, whence $c\neq0$. Then $w=u/c$ lies in
$1+\mm_{\C}$ and satisfies $w\bar w=1$, that is $\bar w=w^{-1}$. Put
$q=\Lmm(w)$. By \cref{diff:lem:smallchain},
\[
 \bar q=\Lmm(\bar w)=\Lmm(w^{-1})=-q,
\]
so $q=i\eta$ with $\eta\in\mm$ real, and $w=\Emm(i\eta)$. Conversely every
expression \cref{diff:eq:unitphase} has modulus one, its standard part recovers
$c$, and the injective local logarithm recovers $\eta$. Finally $\der c=0$, so
\cref{diff:eq:diffexplog} gives \cref{diff:eq:unitlogder}. Multiplication
multiplies the $c$'s and adds the $\eta$'s.
\end{proof}

The proof chooses no global value of $\sin\omega$, no branch of a global
logarithm and no global circular character. The ordinary phase is stored as the
\emph{discrete datum} $c$, not as a selected real argument, so there is no
ambiguity by multiples of $2\pi$: the branch ambiguity of an angle is an
ordinary constant $2\pi k$, which $\der$ kills, and the infinitesimal residue
in \cref{diff:eq:unitphase} is branch-free outright.

\begin{definition}[The finite-angle map]
\label{diff:def:cis}
For a \emph{finite} real surreal $\theta=r+\eta$, with $r=\st(\theta)\in\R$
and $\eta\in\mm$, set
\[
 \cis(\theta)=e^{ir}\,\Emm(i\eta),
\]
where $e^{ir}\in\C$ is the ordinary complex exponential, a differential
constant. The real and imaginary coordinates of $\cis$ are the canonical
finite-angle cosine and sine, and they agree with their Taylor construction.
The domain is $\Ofin$ only; no value of a circular function at an infinite
angle is used.
\end{definition}

\begin{theorem}[Finite-angle phase group, with its differential identity]
\label{diff:thm:cisgroup}
The map $\cis:(\Ofin,+)\to(\Tor(\No),\cdot)$ is a surjective group
homomorphism with kernel exactly $2\pi\Z$, and
\begin{equation}\label{diff:eq:cisder}
 \der\cis(\theta)=i(\der\theta)\cis(\theta),
 \qquad\text{equivalently}\qquad
 \frac{\der\cis(\theta)}{\cis(\theta)}=i\,\der\theta .
\end{equation}
Consequently every nonzero $y\in\SC$ is $y=\rho\cis(\theta)$ with $\rho>0$ in
$\No$ and $\theta$ \emph{finite}.
\end{theorem}
\begin{proof}
The group law follows from \cref{diff:lem:smallchain} and the ordinary phase
group law; conjugation inverts $\Emm(i\eta)$, so the image lies in
$\Tor(\No)$. For surjectivity, let $u\in\Tor(\No)$ and take $c,\eta$ from
\cref{diff:thm:unitphase}; choosing an ordinary real $r$ with $e^{ir}=c$ gives
$u=\cis(r+\eta)$. If $\cis(r+\eta)=1$, standard parts give $e^{ir}=1$, so
$r\in2\pi\Z$, and then \cref{diff:eq:expinjective} gives $\eta=0$; conversely
these elements lie in the kernel. Since $r$ and $e^{ir}$ are differential
constants, \cref{diff:eq:diffexplog} gives \cref{diff:eq:cisder}. Applying
surjectivity to $y/\abs y$ gives the last assertion.
\end{proof}

The surjectivity and kernel content of \cref{diff:thm:cisgroup} is the
finite-angle phenomenon already central to the repository's trigonometry report
\cite{RepoTrig}; it is reproved here because the merged development needs it
with the derivative identity attached. The differential clause
\cref{diff:eq:cisder} is the part the present theory adds, and it is where the
obstruction originates: the factor $\der\theta$ is indispensable, and removing
it would replace an infinitesimal angular rate by the forbidden constant rate
$1$. In particular, for an infinitesimal $\varepsilon$ the finite-angle sine
and cosine satisfy
$\der\sin\varepsilon=(\cos\varepsilon)\der\varepsilon$ and
$\der\cos\varepsilon=-(\sin\varepsilon)\der\varepsilon$, so there is no
conflict between the existence of finite-angle trigonometry and the
nonexistence of an intrinsic oscillator (\cref{diff:cor:oscillator}).

\subsection{Polar decomposition and its differential content}
\label{diff:sub:polar}

\begin{definition}[Logarithmic derivative]
\label{diff:def:logder}
For $z\in\SC^{\times}$ write $z^{\dagger}=\der z/z$. The notation does not
presuppose the existence of a complex logarithm. It is a group homomorphism
$(\SC^{\times},\cdot)\to(\SC,+)$: $(zw)^{\dagger}=z^{\dagger}+w^{\dagger}$,
$(z^{-1})^{\dagger}=-z^{\dagger}$, and $\Ker(\dagger)=\C^{\times}$ by
\cref{diff:prop:complexification}.
\end{definition}

\begin{corollary}[Branch-free polar form]
\label{diff:cor:polar}
Every $z\in\SC^{\times}$ has a unique decomposition
\begin{equation}\label{diff:eq:polar}
 z=\rho\,c\,\Emm(i\eta),
 \qquad \rho=\abs z>0,\quad c=\st(z/\abs z)\in\Tor(\C),\quad\eta\in\mm,
\end{equation}
and
\begin{equation}\label{diff:eq:polarlogder}
 z^{\dagger}=\der\log\rho+i\,\der\eta
 =\frac{\der\rho}{\rho}+i\,\der\eta .
\end{equation}
In coordinates $z=u+iv$,
\begin{equation}\label{diff:eq:coordinates}
 \Ree(z^{\dagger})=\frac{uu'+vv'}{u^{2}+v^{2}},
 \qquad
 \Imm(z^{\dagger})=\frac{uv'-vu'}{u^{2}+v^{2}}.
\end{equation}
Equivalently $z=\rho\cis(\theta)$ with $\theta\in\Ofin$, and then
$\Ree(z^{\dagger})=\der\abs z/\abs z$ and
$\Imm(z^{\dagger})=\der\theta$.
\end{corollary}
\begin{proof}
Apply \cref{diff:thm:unitphase} to $z/\abs z$ and use
$\der\log\rho=\der\rho/\rho$ together with \cref{diff:eq:unitlogder}; the
coordinate formulas are direct computation. The $\cis$ form is
\cref{diff:thm:cisgroup}, and $\der\theta=\der\eta$ because
$\theta-\eta=\st(\theta)\in\R$.
\end{proof}

The first summand of \cref{diff:eq:polarlogder} is a radial logarithmic rate
and is unrestricted; the second is an angular rate, and its allowed
primitives are constrained to finite angles. An enormous modulus does not
create an enormous circular phase: the direction of $\omega+i$, for example, is
a small perturbation of $1$. Identity \cref{diff:eq:polarlogder} is precisely
the reason that an arbitrary radial logarithmic derivative is possible while an
arbitrary imaginary one is not.

\begin{warning}[A unit cannot hide an infinite accumulated phase]
\label{diff:warn:nohiding}
A unit $u$ may be \emph{denoted} by an expression involving a large parameter.
As an element of $\SC$ it nevertheless has a finite angle $\theta$, and hence
$u^{\dagger}=i\der\theta\in i\Acl$. A notation such as $e^{i\log\omega}$
cannot, by itself, turn an infinitesimal coefficient into a permissible
logarithmic derivative. This is not a restriction on ordinary complex
exponentiation, where the variable is a coordinate, and it is not a claim that
global phase characters are impossible; it is a restriction on compatibility
with the specified scalar derivation.
\end{warning}

\section{Primitives, accumulated phase, and the bounded-primitive ideal}
\label{diff:sec:primitives}

\subsection{The normalized primitive}
\label{diff:sub:normprimitive}

\begin{proposition}[Normalized primitive]
\label{diff:prop:normprimitive}
For each $b\in\No$ there is a unique $\Izero(b)\in\No$ with
\[
 \der\Izero(b)=b,\qquad \ct(\Izero(b))=0 .
\]
The map $\Izero$ is $\R$-linear and satisfies
$\der\circ\Izero=\mathrm{id}$ and $\Izero\circ\der=\mathrm{id}-\ct$. Its
componentwise extension $\Izero(a+ib)=\Izero(a)+i\Izero(b)$ is the unique
$\C$-linear map $\SC\to\SC$ with $\der\Izero=\mathrm{id}$ and
$\ct\circ\Izero=0$, and it satisfies
$\Izero(\der f)=f-\ct(f)$.
\end{proposition}
\begin{proof}
By \ref{diff:BM:surj} choose any primitive $F$ of $b$ and put
$\Izero(b)=F-\ct(F)$. Any two primitives differ by an element of
$\Ker\der=\R$, and if both have vanishing $\omega^{0}$ coefficient that
difference is zero; this gives existence, uniqueness and independence of the
choice of $F$. Linearity follows from uniqueness, since a sum or real multiple
of normalized primitives has the required derivative and normalization. The
identity $\Izero(\der f)=f-\ct(f)$ holds because $f-\ct(f)$ is a primitive of
$\der f$ with zero constant coefficient. The complex statements follow
componentwise from \cref{diff:prop:complexification}.
\end{proof}

\begin{warning}[$\Izero$ is a normalization, not an algorithm]
\label{diff:warn:Jnotalgorithm}
The definition is by a \emph{uniqueness} property and makes no assertion about
an effective primitive-finding procedure. The normalization refers to the
coefficient of $\omega^{0}$, \emph{not} to evaluation at an initial point;
$\Izero$ is not asserted to be multiplicative, not asserted to preserve
arbitrary Hahn sums, and not asserted to be computable from a supplied symbolic
representation. It does not require a choice function on primitive families
--- the normalization singles out a unique element, so no choice principle is
involved. Its value is semantic precision. Sample values:
$\Izero(1)=\omega$, and $\Izero(\omega^{-1})=\log\omega-\ct(\log\omega)$,
which is infinite regardless of its real constant term.
\end{warning}

The integration constant is not evaluation at an endpoint, and one familiar
algebraic identity accordingly acquires a correction term.

\begin{proposition}[Integration by parts, and a Rota--Baxter correction]
\label{diff:prop:RB}
For $f,g\in\SC$,
\begin{equation}\label{diff:eq:parts}
 \Izero(f\der g+g\der f)=fg-\ct(fg),
\end{equation}
whose final term is not in general $\ct(f)\ct(g)$; and
\begin{equation}\label{diff:eq:RB}
 \Izero(f)\,\Izero(g)=\Izero\bigl(f\,\Izero(g)+\Izero(f)\,g\bigr)
 +\ct\bigl(\Izero(f)\Izero(g)\bigr).
\end{equation}
The correction term in \cref{diff:eq:RB} cannot in general be omitted.
\end{proposition}
\begin{proof}
Both identities follow by differentiating the left-hand sides and applying
$\Izero\circ\der=\mathrm{id}-\ct$. For the last assertion take $f=1$ and
$g=t^{2}$: then $\Izero(f)=\omega$, $\Izero(g)=-t$, so
$\Izero(f)\Izero(g)=-1$, while
$f\Izero(g)+\Izero(f)g=-t+\omega t^{2}=0$. Without the correction the identity
would assert $-1=0$.
\end{proof}

In the terminology sometimes used for integral operators, $\Izero$ is
therefore \emph{not} a weight-zero Rota--Baxter operator on all of $\SC$, and
\cref{diff:eq:RB} is the precise replacement. No such terminology is needed to
use the formula; the point is that the selected constant coefficient is linear
but not multiplicative (\cref{diff:warn:ct}), so a ``zero constant of
integration'' convention is not automatically an integral based at a common
evaluation point.

\subsection{Accumulated phase and its obstruction}
\label{diff:sub:accumulated}

\begin{definition}[Accumulated phase and obstruction]
\label{diff:def:obs}
For $b\in\No$, an \emph{accumulated phase} of $b$ is any $B\in\No$ with
$\der B=b$; the normalized choice is $B=\Izero(b)$. Put
\[
 \Acl:=\der\mm=\der\Ofin,
 \qquad
 \Obs(b):=\pp\bigl(\Izero(b)\bigr)\in\Pcl ,
\]
and for $q\in\SC$ put
\[
 \Phi(q):=\Obs(\Imm q)\in\Pcl .
\]
We call $\Acl$ the \emph{bounded-primitive ideal} (equivalently, the class of
\emph{admissible angular rates}), and $\Obs(b)$, $\Phi(q)$ the \emph{phase
obstruction}.
\end{definition}

By \cref{diff:rem:obsnormalization} the value $\Obs(b)=\pp(B)$ is the same for
every accumulated phase $B$ of $b$, so $\Obs$ is well defined without the
normalization. It has an actual surreal representative, not an ambiguous
antiderivative choice. Sample values:
\begin{equation}\label{diff:eq:obsexamples}
 \Obs(1)=\omega,\qquad
 \Obs(\omega^{-1})=\log\omega,\qquad
 \Obs(\omega^{-2})=0 .
\end{equation}
The obstruction measures accumulated phase, not the instantaneous size of a
coefficient: $\Phi(a)=0$ for every real coefficient $a\in\No$, however
large, whereas $\Phi(ib)=\Obs(b)$ may be nonzero even for infinitesimal $b$.

\begin{theorem}[Structure of the bounded-primitive ideal]
\label{diff:thm:idealstructure}
\begin{enumerate}[label=\textup{(\roman*)},leftmargin=2.4em]
\item $\Acl=\der\mm=\der\Ofin$, and $\der:\mm\to\Acl$ is an $\R$-linear
bijection that \emph{reverses} strict order.
\item For $b\in\No$ the following are equivalent: $b\in\Acl$; $\Izero(b)$ is
infinitesimal; $\Izero(b)$ is finite; some primitive of $b$ is finite; every
primitive of $b$ is finite; $\Obs(b)=0$.
\item $\Acl$ is convex in $\No$, is an additive subgroup, and is an
$\Ofin$-submodule of $(\No,+)$, that is, an ideal of the valuation ring
$\Ofin$.
\item $0\neq\Acl\subsetneq\mm\subsetneq\Ofin$, with witnesses
$\omega^{-2}=\der(-\omega^{-1})\in\Acl$ and $\omega^{-1}\in\mm\setminus\Acl$.
\item There is an exact cut description
\begin{equation}\label{diff:eq:Icut}
 \Acl=\{b\in\No:\abs b<\der f\text{ for every }f>\R\}.
\end{equation}
\item $\Obs:\No\to\Pcl$ is $\R$-linear and surjective with kernel exactly
$\Acl$; indeed $\Izero(\der P)=P$ and $\Obs(\der P)=P$ for $P\in\Pcl$.
\item As additive real vector classes,
\begin{equation}\label{diff:eq:directsum}
 \No=\der\Pcl\oplus\Acl=\der\Pcl\oplus\der\mm,
\end{equation}
with projection $b\mapsto\der\pp\Izero(b)$ onto the first summand; $\der$ is
injective on each of $\Pcl$ and $\mm$ separately, and
$\der:\Pcl\to\der\Pcl$ is an order-\emph{preserving} $\R$-linear bijection.
\item $\der$ induces an order-preserving $\R$-linear isomorphism of ordered
additive quotients
\begin{equation}\label{diff:eq:quotientiso}
 \No/\Ofin\longrightarrow\No/\Acl,
 \qquad x+\Ofin\longmapsto \der x+\Acl,
\end{equation}
and in particular $\der^{-1}(\Acl)=\Ofin$.
\end{enumerate}
\end{theorem}
\begin{proof}
(i) The equality $\der\Ofin=\der\mm$ and the order reversal are
\cref{diff:lem:smallnegative}; the kernel of $\der|_{\mm}$ is
$\mm\cap\R=\{0\}$, giving injectivity, and surjectivity onto $\Acl$ is the
definition.

(ii) A primitive in $\mm$ is finite; conversely a finite primitive becomes
infinitesimal after subtracting its standard part. All primitives differ by a
real constant, so either all are finite or none is. Since $\ct(\Izero(b))=0$,
the splitting \cref{diff:eq:split} makes $\Izero(b)$ finite exactly when it is
infinitesimal, exactly when $\pp\Izero(b)=0$.

(iii) Suppose $0<b\leq d$ with $d\in\Acl$. Write $d=-\der\varepsilon$ with
$0\leq\varepsilon\in\mm$, using (i), and let $B$ be a primitive of $b$. If $B$
were negative infinite then $-B$ would be positive infinite and
\ref{diff:BM:positivity} would give $b=\der B<0$, a contradiction. If $B$ were
positive infinite then $B+\varepsilon$ would be positive infinite, whence
$\der(B+\varepsilon)=b-d>0$, contradicting $b\leq d$. So $B$ is finite and
$b\in\Acl$ by (ii). Sign symmetry and additivity give convexity in general.
For $a\in\Ofin$ choose a positive integer $n>\abs a$; then $\abs{ab}\leq n\abs b$
for $b\in\Acl$, and convexity with real linearity gives $ab\in\Acl$.

(iv) $\Acl\subseteq\mm$ is \cref{diff:lem:smallnegative}. Nonzero and strict:
the displayed witnesses, using \cref{diff:warn:strict}.

(v) If $b=\der\varepsilon$ with $\varepsilon\in\mm$, then for every $f>\R$ both
$f+\varepsilon$ and $f-\varepsilon$ are positive infinite, so their derivatives
are positive and $\der f>\abs b$. Conversely let $b\notin\Acl$ and let $B$ be a
primitive; then $B$ is infinite. Take $f=\abs B/2>\R$; then
$\der f=\abs b/2$ by \ref{diff:BM:positivity}, so the required inequality
$\abs b<\der f$ fails.

(vi) Linearity is inherited from $\Izero$ and $\pp$. If $P\in\Pcl$ then
$\ct(P)=0$, so $P$ is the normalized primitive of $\der P$ and
$\Obs(\der P)=\pp(P)=P$; this gives surjectivity. The kernel is (ii).

(vii) Decompose $\Izero(b)=P+\varepsilon$ with $P=\pp\Izero(b)\in\Pcl$ and
$\varepsilon\in\mm$, and differentiate: $b=\der P+\der\varepsilon$. If
$\der P=\der\varepsilon$ with $P\in\Pcl$, $\varepsilon\in\mm$, then
$P-\varepsilon\in\R$ and uniqueness of \cref{diff:eq:split} forces
$P=\varepsilon=0$; the same argument with one term zero gives the injectivity
assertions. A nonzero positive member of $\Pcl$ is positive infinite, so its
derivative is positive by \ref{diff:BM:positivity}; applying this to
differences gives the order assertion.

(viii) Surjectivity is inherited from \ref{diff:BM:surj}, and
$\der^{-1}(\Acl)=\Ofin$ is (ii) applied to $\der x$. If $x+\Ofin$ is positive
then $x$ is positive infinite, and for every $d=\der\varepsilon\in\Acl$ the
element $x-\varepsilon$ is positive infinite, so $\der x>d$; hence the induced
map preserves positivity, and the converse follows by bijectivity and sign
symmetry.
\end{proof}

\begin{definition}[The unique infinitesimal primitive]
\label{diff:def:uniqueinf}
By \cref{diff:thm:idealstructure}(i) the restriction $\der|_{\mm}$ is a
bijection onto $\Acl$; its inverse is $\Izero|_{\Acl}:\Acl\to\mm$, the
\emph{unique infinitesimal primitive}. Unlike a general primitive, it carries
no residual additive-constant ambiguity. (This is the object that two source
manuscripts call $J$; see \cref{diff:conv:noJ}.)
\end{definition}

\begin{remark}[The represented quotient, concretely]
\label{diff:rem:representedquotient}
\Cref{diff:eq:quotientiso} is the represented form of the familiar
isomorphism ``$\No/\Ofin\cong\No/\Acl$''. Its actual representatives are the
classes $\Pcl$ and $\der\Pcl$, in accordance with
\cref{diff:warn:cosets}; no proper-class coset is made into an object. Inside a
set-sized field (\cref{diff:sec:localization}) the corresponding ordinary set
quotients are available.
\end{remark}
\section{The finite-primitive criterion}
\label{diff:sec:phasecrit}

\subsection{The exact image of the logarithmic derivative}
\label{diff:sub:logimage}

\begin{theorem}[Exact logarithmic-derivative image]
\label{diff:thm:logimage}
For the fixed derivation on $\SC$,
\begin{equation}\label{diff:eq:logimage}
 \boxed{\quad
 \{z^{\dagger}:z\in\SC^{\times}\}
 =\No+i\Acl
 =\No+i\,\der\mm
 =\No+i\,\der\Ofin .
 \quad}
\end{equation}
The kernel of $\dagger$ is $\C^{\times}$, and $\der:\mm\to\Acl$ is injective.
The image is a \emph{proper} subclass of $\SC$: $1\in\No$ is in the image,
being $(e^{\omega})^{\dagger}$, while $i$ is not.
\end{theorem}
\begin{proof}
If $z\neq0$, \cref{diff:cor:polar} gives
$z^{\dagger}=\der\log\abs z+i\der\eta$ with $\eta\in\mm$, so the imaginary
part lies in $\der\mm=\Acl$; this is the inclusion from left to right.

For the reverse inclusion, let $a\in\No$ and $b\in\Acl$. Put $A=\Izero(a)$ and
let $\eta\in\mm$ be the unique infinitesimal primitive of $b$
(\cref{diff:def:uniqueinf}). Then
\begin{equation}\label{diff:eq:generator}
 y_{0}=\exp(A)\,\Emm(i\eta)\neq0
\end{equation}
and, by \ref{diff:BM:exp} and \cref{diff:eq:diffexplog},
$y_{0}^{\dagger}=\der A+i\der\eta=a+ib$. The kernel assertion is
\cref{diff:def:logder}; injectivity of $\der$ on $\mm$ is
\cref{diff:thm:idealstructure}(i). Properness: $1\notin\Acl$ by
\cref{diff:lem:smallnegative}, since $1$ is not infinitesimal.
\end{proof}

\begin{warning}[Two products of two different exponentials]
\label{diff:warn:twoexponentials}
The witness \cref{diff:eq:generator} is a product of \emph{two} exponential
constructions with explicitly different declared domains: $\exp$ is the real
surreal exponential on $\No$, and $\Emm$ is the infinitesimal Hahn exponential
of \cref{diff:eq:explog} on $\mm_{\C}$. The product is not an abbreviation for
an unspecified global surcomplex exponential, and the two are never glued.
\Cref{diff:sec:strip} determines exactly how far a single compatible
exponential can reach, and \cref{diff:cor:noglobalPsi} shows it cannot reach
all of $\SC$.
\end{warning}

\begin{warning}[The image is real-linear, not complex-linear]
\label{diff:warn:notClinear}
$\No+i\Acl$ is closed under addition and under multiplication by real
constants, but not under multiplication by $i$: every real coefficient has
zero $\Phi$-obstruction, while multiplying one by $i$ may not. An implementation must
not model this image as a complex vector subspace.
\end{warning}

\subsection{The criterion, and all solutions}
\label{diff:sub:criterion}

\begin{theorem}[Finite-primitive criterion for rank-one equations]
\label{diff:thm:phasecriterion}
Let $q=a+ib\in\SC$ with $a,b\in\No$, and let $A,B\in\No$ satisfy $\der A=a$
and $\der B=b$. The following are equivalent.
\begin{enumerate}[label=\textup{(\roman*)},leftmargin=2.4em]
\item $\der y=qy$ has a nonzero solution in $\SC$.
\item $b\in\Acl$, that is, $b=\der\eta$ for some $\eta\in\mm$.
\item $b$ has a finite accumulated phase: $B\in\Ofin$.
\item Every accumulated phase of $b$ is finite.
\item $\Izero(b)$ is finite, equivalently infinitesimal.
\item $\pp(B)=0$, equivalently $\Obs(b)=0$, equivalently $\Phi(q)=0$.
\end{enumerate}
When these hold, the infinitesimal $\eta$ with $\der\eta=b$ is unique, namely
$\eta=\Izero(b)=B-\st(B)$, and \emph{all} solutions are
\begin{equation}\label{diff:eq:allsolutions}
 y=c\,\exp(\Izero(a))\,\Emm\bigl(i\Izero(b)\bigr),\qquad c\in\C,
\end{equation}
the nonzero ones being those with $c\in\C^{\times}$. When they fail, the only
solution is $y=0$.
\end{theorem}
\begin{proof}
(i)$\Leftrightarrow$(ii) is \cref{diff:thm:logimage}, since for $y\neq0$ the
equation says $y^{\dagger}=q$. Alternatively and directly: if $y\neq0$ solves
it, write $y=\rho\cis(\theta)$ by \cref{diff:thm:cisgroup}; then
\cref{diff:eq:polarlogder} gives $\der\rho/\rho=a$ and $\der\theta=b$ with
$\theta$ finite, which is (iii). The equivalences among
(ii)--(vi) are \cref{diff:thm:idealstructure}(ii). The displayed family
consists of solutions by \cref{diff:eq:generator}; if $y$ is any solution and
$y_{0}$ is the nonzero one, the quotient rule gives $\der(y/y_{0})=0$, so
$y/y_{0}\in\C$ by \cref{diff:prop:complexification}. The case $c=0$ is the
zero solution.
\end{proof}

This is the theorem that all seven source manuscripts prove, in seven
notations. Three points about it deserve separate statements.

\begin{warning}[Integrated, not pointwise]
\label{diff:warn:integrated}
The criterion is \emph{not} ``$b$ must be infinitesimal''. That condition is
necessary and strictly weaker. The three equations
\[
 \der y=iy,\qquad \der y=i\omega^{-1}y,\qquad \der y=i\omega^{-2}y
\]
have, respectively: only the zero solution (the accumulated phase of $1$ is
$\omega$); only the zero solution, although $\omega^{-1}$ is infinitesimal
(its accumulated phase is $\log\omega$, infinite); and the solutions
$c\,\Emm(-i\omega^{-1})$, $c\in\C$. A sibling statement that ``$b\in\mm$ is
necessary'' is a strictly weaker and different result.
\end{warning}

\begin{remark}[The real coefficient never obstructs]
\label{diff:rem:realnoobstruct}
Adding any real coefficient $a$ introduces no further obstruction in the full
class: multiply the solution by $\exp(\Izero(a))$. Thus
\[
 \der y=\left(\omega^{3}+\frac{i}{\omega^{2}}\right)y
 \quad\text{has solutions}\quad
 y=c\exp(\omega^{4}/4)\,\Emm(-i/\omega),
\]
while replacing $i/\omega^{2}$ by $i/\omega$ leaves only $y=0$; and
$\der y=(\omega+i\omega^{-2})y$ has solutions
$c\exp(\omega^{2}/2)\Emm(-i/\omega)$. The real growth rate may be arbitrarily
large. That assertion can fail in a \emph{restricted workspace}, because $A$
or $e^{A}$ may leave the workspace; \cref{diff:sec:rationalworkspace} computes
that failure exactly.
\end{remark}

\begin{remark}[A one-dimensional equation with a zero-dimensional solution space]
\label{diff:rem:dimension}
A one-dimensional linear equation need not have a one-dimensional nonzero
solution space \emph{inside this field}. The familiar dimension count is an
upper bound until a sufficiently large differential extension is supplied;
\cref{diff:sec:PV} supplies one.
\end{remark}

\subsection{No intrinsic oscillator: three independent proofs}
\label{diff:sub:nooscillator}

\begin{corollary}[The missing oscillator]
\label{diff:cor:oscillator}
For $\lambda\in\C$, the equation $\der y=\lambda y$ has a nonzero solution in
$\SC$ if and only if $\lambda\in\R$, in which case the solutions are exactly
$c\,e^{\lambda\omega}$ with $c\in\C$. In particular $\der y=iy$ and
$\der y=-iy$ have only the zero solution, and for every ordinary real
$\lambda\neq0$ so do $\der y=i\lambda y$ and $\der^{2}y+\lambda^{2}y=0$. In
particular $\der^{2}y+y=0$ has only $y=0$ in $\SC$, while $\der^{2}y-y=0$ has
all solutions $Ae^{\omega}+Be^{-\omega}$ with $A,B\in\C$.
\end{corollary}
\begin{proof}[First proof, from the criterion]
Write $\lambda=\alpha+i\beta$. An accumulated phase of $\beta$ is
$\beta\omega$, which is finite exactly when $\beta=0$; apply
\cref{diff:thm:phasecriterion}. When $\beta=0$, dividing a solution by
$e^{\alpha\omega}$ gives derivative zero. For the second-order equation,
factor $\der^{2}+\lambda^{2}=(\der-i\lambda)(\der+i\lambda)$ over $\SC$; both
factors are injective by the first part, so their composite is.
\end{proof}
\begin{proof}[Second proof, from smallness alone]
Suppose $\der y=i\beta y$ with $\beta\in\R\setminus\{0\}$ and $y\neq0$.
Conjugation gives $\der\bar y=-i\beta\bar y$, so $\der(y\bar y)=0$ and
$y\bar y$ is a positive real constant. Dividing by its positive square root
gives $u+iv$ with
\[
 u^{2}+v^{2}=1,\qquad \der u=-\beta v,\qquad \der v=\beta u .
\]
Both $u$ and $v$ are finite, so both derivatives are infinitesimal by
\cref{diff:lem:smallnegative}; dividing by the nonzero real constant $\beta$
then makes both $u$ and $v$ infinitesimal, so $u^{2}+v^{2}$ is infinitesimal
and cannot equal $1$.
\end{proof}
\begin{proof}[Third proof, an energy argument for the second-order equation]
First let $f\in\No$ be real with $\der^{2}f+f=0$. The energy
$E=f^{2}+(\der f)^{2}$ satisfies
$\der E=2\der f\,(f+\der^{2}f)=0$, so $E$ is a nonnegative real constant; hence
$f$ and $\der f$ are finite. By \cref{diff:lem:smallnegative} both $\der f$ and
$\der^{2}f$ are infinitesimal, and since $f=-\der^{2}f$ the element $f$ is
infinitesimal too. Then $E$ is both real and infinitesimal, so $E=0$, and
positivity gives $f=0$. For complex $f$, apply this to its real and imaginary
parts.
\end{proof}

The three proofs are genuinely different and all three survive. The first uses
the whole criterion. The second uses only the constant field, the $H$-field
positivity clause and the normalization $\der\omega=1$ --- not the finite-angle
parametrization, not surjectivity, and no logarithm --- so the obstruction is
provably not an artefact of a branch choice or of a chosen phase convention.
The third uses only conservation of energy plus smallness, and isolates the
role of smallness in the second-order equation; the factorization proof,
however, scales more efficiently to general operators
(\cref{diff:thm:constant}).

\begin{warning}[What the nonexistence statements are about]
\label{diff:warn:whatnonexistence}
All of these are statements about surcomplex \emph{number} solutions under the
fixed $\der$. They do not deny ordinary oscillatory functions, do not say that
the coordinate functions $\sin z$, $\cos z$ fail their ordinary differential
equations, and do not deny algebraically defined global surcomplex phases. They
do say that $\sin\omega$ and $\cos\omega$ are not Berarducci--Mantova
differential elements of $\No$ in the sense of
\cref{diff:conv:meaning}: the system $c^{2}+s^{2}=1$, $c'=-s$, $s'=c$ has no
solution in $\No$. See \cref{diff:cor:rotation,diff:sub:globalexp}.
\end{warning}

\begin{corollary}[The unit-circle version]
\label{diff:cor:rotation}
For $x\in\No$, the system
\[
 c^{2}+s^{2}=1,\qquad \der c=-(\der x)s,\qquad \der s=(\der x)c
\]
has a solution $c,s\in\No$ if and only if $x$ is finite.
\end{corollary}
\begin{proof}
If $c,s$ solve it then $z=c+is\neq0$ satisfies $\der z=i(\der x)z$, so by
\cref{diff:thm:phasecriterion} $\der x$ has a finite accumulated phase; all
those are $x+r$ with $r\in\R$, so $x$ is finite. Conversely, for finite $x$
take the coordinates of $\cis(x)$ and apply \cref{diff:eq:cisder}.
\end{proof}

This concerns the derivatives of the \emph{values} $c,s$. An externally
differentiable sine function is a different object, and its derivative with
respect to its argument must not be substituted for $\der$ without a proved
chain rule (\cref{diff:warn:separators}).

\begin{corollary}[A Riccati transfer]
\label{diff:cor:riccati}
The only solutions in $\SC$ of $\der u=1+u^{2}$ are $u=i$ and $u=-i$; in
particular there is none in $\No$.
\end{corollary}
\begin{proof}
Both constants are solutions. For any other solution, $v=(u-i)/(u+i)$ is
defined and nonzero, and
\[
 v^{\dagger}=\frac{\der u}{u-i}-\frac{\der u}{u+i}
 =\frac{2i\,\der u}{u^{2}+1}=2i,
\]
contradicting \cref{diff:cor:oscillator}.
\end{proof}

The same Riccati equation for the variable derivative $D_z$, on strongly
entire Hahn power series, has the same two solutions; see
\texttt{hol:nl:prop:riccati} in the holonomic-rigidity report. The
derivations differ, and neither result implies the other.

This is not a general nonlinear solvability theorem. It is a direct way to
transfer the first-order phase obstruction to a familiar nonlinear equation.
(Added with member~11: the sentence stands for this proof, but the class
containing the equation is now classified. \Cref{diff:aut:thm:main} describes all
solutions of every first-order autonomous equation with constant coefficients,
for every normalized derivation, and \cref{diff:aut:ex:riccati} recovers this
corollary from it by a projective route. Beyond that class --- variable
coefficients, higher order, systems --- no nonlinear solvability theorem is
claimed; see items~N21 and~N25.)

\begin{remark}[Algebraic closure is not differential closure]
\label{diff:rem:algnotdiff}
$\SC$ is algebraically closed and $\der$ is additively surjective, yet
$\der y=iy$ has no nonzero solution. These are three different closure
properties: $X^{2}+1=0$ has a solution, $\der y=f$ has a solution, and
$\der y=iy$ with $y\neq0$ does not. Adjoining an algebraic element cannot
repair the last, since algebraic closedness leaves no proper algebraic
extension; and taking an algebraic root cannot cancel a nonzero obstruction,
because every nonzero integer multiple of a nonzero element of $\Pcl$ is
nonzero (\cref{diff:cor:tensor}). A genuine differential extension must
introduce a transcendental element, and \cref{diff:sec:PV} constructs several.
The nonzero requirement can itself be written differentially, by adding an
unknown $v$ and the equation $yv=1$; so this is an explicit failure of
differential closedness, not a complaint about an undefined expression such as
$e^{i\omega}$.
\end{remark}

\section{Gauge normal form and the obstruction sequence}
\label{diff:sec:gauge}

\subsection{Gauge equivalence and the unique normal form}
\label{diff:sub:gaugeform}

A change of unknown $y=gz$ with $g\in\SC^{\times}$ sends
\begin{equation}\label{diff:eq:gauge}
 \der y=qy
 \qquad\text{to}\qquad
 \der z=\bigl(q-g^{\dagger}\bigr)z .
\end{equation}
Call $q,r\in\SC$ \emph{gauge equivalent} when $q-r=g^{\dagger}$ for some
$g\in\SC^{\times}$. A gauge transformation changes the coefficient by a
logarithmic derivative, not by an arbitrary primitive, and it does not change
whether a nonzero solution exists in the base field. No differential-module
terminology is needed to use it.

\begin{theorem}[Purely infinite phase normal form, with uniqueness]
\label{diff:thm:gauge}
The map $\Phi:\SC\to\Pcl$, $\Phi(q)=\Obs(\Imm q)$, is additive and
$\R$-linear, is surjective, and satisfies
\[
 \Phi(q)=0\quad\Longleftrightarrow\quad
 q\text{ is a logarithmic derivative} .
\]
Two coefficients are gauge equivalent if and only if their obstructions agree.
Explicitly, writing $q=a+ib$ and $\Izero(b)=P+\eta$ with $P=\Phi(q)\in\Pcl$
and $\eta\in\mm$, the gauge
\[
 g=\exp(\Izero(a))\,\Emm(i\eta)
\]
transforms $\der y=qy$ into
\begin{equation}\label{diff:eq:normalgauge}
 \der z=i(\der P)z .
\end{equation}
Every equation $\der y=qy$ is therefore gauge equivalent to \emph{exactly one}
equation of the form \cref{diff:eq:normalgauge} with $P\in\Pcl$, and the
canonical parameter is $P=\Phi(q)$.
\end{theorem}
\begin{proof}
Linearity and surjectivity are \cref{diff:thm:idealstructure}(vi) composed
with $\Imm$; the kernel description is \cref{diff:thm:phasecriterion}. The
logarithmic derivative of the displayed $g$ is $a+i\der\eta$, and
$b=\der P+\der\eta$, so \cref{diff:eq:gauge} gives
\cref{diff:eq:normalgauge}. If $q_{1}-q_{2}=g^{\dagger}$ then
$\Imm(q_{1}-q_{2})\in\Acl$ and linearity of $\Obs$ makes the obstructions
equal; conversely equal obstructions put $\Imm(q_{1}-q_{2})$ in $\Acl$, and
\cref{diff:thm:logimage} supplies a $g$ with $g^{\dagger}=q_{1}-q_{2}$.
Uniqueness: $\Phi(i\der P)=\Obs(\der P)=P$ by
\cref{diff:thm:idealstructure}(vi), so distinct $P$ represent distinct classes.
\end{proof}

The normal form says more than ``no solution''. It isolates the precise
accumulated phase that cannot be represented by a surcomplex direction: the
real growth factor and the infinitesimal phase residue are exactly the
removable parts, and the purely infinite part is exactly the irremovable one.
The parameter $P$ is \emph{not} an angular value assigned to an infinite
input; it records the part of an accumulated phase that no existing surcomplex
phase can absorb, the absorbed factor carrying only the infinitesimal residue
$\eta$ (\cref{diff:conv:phases}).

\begin{corollary}[Obstruction exact sequence]
\label{diff:cor:exact}
The following is an exact sequence of class groups, multiplicative before
$\dagger$ and additive after it:
\begin{equation}\label{diff:eq:exact}
 1\longrightarrow\C^{\times}\longrightarrow\SC^{\times}
 \xrightarrow{\ \dagger\ }(\SC,+)
 \xrightarrow{\ \Phi\ }(\Pcl,+)\longrightarrow 0 .
\end{equation}
The corresponding real obstruction sequence is
$0\to\Acl\to\No\xrightarrow{\Obs}\Pcl\to 0$; it uses
$\Phi(ib)=\Obs(b)$, not the restriction $\Phi|_{\No}=0$.
\end{corollary}
\begin{proof}
$\Ker(\dagger)=\C^{\times}$ and $\operatorname{im}(\dagger)=\No+i\Acl$ by
\cref{diff:thm:logimage}; that class is $\Ker\Phi$ by
\cref{diff:thm:phasecriterion}; and $\Phi$ is onto by
\cref{diff:thm:gauge}.
\end{proof}

\begin{warning}[The sequence is elementwise]
\label{diff:warn:elementwise}
The word ``class'' in \cref{diff:eq:exact} is intentional, and the sequence is
shorthand for the displayed kernel and image identities \emph{about elements}.
No set or class of all equivalence classes is formed, and $\Pcl$ appears
precisely so that no coset of $\Ofin$ has to be an object
(\cref{diff:warn:cosets}). Every map in the sequence acts on set-sized
elements, and $\Phi$ together with the representative $i\der P$ carries all the
content.
\end{warning}

\begin{corollary}[Tensor, dual, and no finite-order obstruction]
\label{diff:cor:tensor}
Under multiplication of scalar solution equations --- equivalently, tensor
product of the corresponding rank-one objects, represented by addition of
coefficients --- the obstructions add; under inversion, that is dualization,
they negate. For an integer $n\neq0$,
\[
 \Phi(nq)=n\Phi(q),\qquad \Phi(nq)=0\iff\Phi(q)=0 .
\]
Every nonzero gauge class has infinite order. Hence raising a putative solution
of an obstructed equation to a finite power does not remove the obstruction.
\end{corollary}
\begin{proof}
The coefficient of a product of nonzero solutions is the sum of their
logarithmic derivatives, and of an inverse is the negative. The rest is
$\R$-linearity of $\Phi$ together with the absence of additive torsion in the
real vector class $\Pcl$.
\end{proof}

The integer linearity just recorded --- rather than real linearity --- is what
drives the differential Galois arguments of
\cref{diff:sec:PV,diff:sub:torus}: rational functions in solutions involve
integral monomial exponents, so the relations that matter are integer
relations. This rank-one classification is a differential gauge
classification; it is not the holomorphic line-bundle or divisor
classification of the repository's global analytic reports, and no theory of
higher-rank differential Galois groups is needed to state it.

\section{Sharp thresholds and worked scales}
\label{diff:sec:thresholds}

\subsection{The support threshold over a real-exponent Hahn field}
\label{diff:sub:supportthreshold}

Fix $t=\omega^{-1}$ and, for an additive subgroup $\Gamma\subseteq\R$, embed
the Hahn fields
\[
 H_{\Gamma}=\R((t^{\Gamma})),\qquad
 K_{\Gamma}=\C((t^{\Gamma}))=H_{\Gamma}[i]
\]
into $\No$ and $\SC$ by Conway monomials. On these fields the real-exponent
rule of \cref{diff:eq:basicder} applies, so $\der(t^{\alpha})=-\alpha t^{\alpha+1}$
for ordinary real $\alpha$ and $\der=-t^{2}\,d/dt$
(\cref{diff:warn:monomialmap} applies: this formula is for real exponents
only). The case $\Gamma=\R$ gives the largest real-exponent field; the cases
$\Gamma=\Z$ and $\Gamma=\Q$ recur in \cref{diff:part:workspaces}.

\begin{theorem}[Sharp support threshold]
\label{diff:thm:power}
Let $\Gamma\subseteq\R$ be an additive subgroup and let
$b=\sum_{\alpha\in S}b_{\alpha}t^{\alpha}\in H_{\Gamma}$ with $S$ well
ordered. A primitive of $b$ in $\No$ is
\begin{equation}\label{diff:eq:powerprimitive}
 B=\sum_{\substack{\alpha\in S\\ \alpha\neq1}}
 \frac{b_{\alpha}}{1-\alpha}\,t^{\alpha-1}
 \;+\;b_{1}\log\omega ,
\end{equation}
and for every $p\in\No$ the equation $\der y=(p+ib)y$ has a nonzero solution
in $\SC$ if and only if
\begin{equation}\label{diff:eq:powerthreshold}
 b=0\qquad\text{or}\qquad \min\supp(b)>1 .
\end{equation}
Equivalently, for real $f\in H_{\Gamma}$ one has $f\in\Acl$ if and only if
$f=0$ or $v(f)>1$, where $v$ is the least-exponent valuation. In the soluble
case the sum in \cref{diff:eq:powerprimitive} is the unique infinitesimal
accumulated phase of $b$.
\end{theorem}
\begin{proof}
Shifting a well-ordered support by $-1$ preserves well ordering, and omitting
one exponent while adding one logarithmic term preserves summability, so
\cref{diff:eq:powerprimitive} is a legitimate element of $\No$; the
real-exponent rule and \ref{diff:BM:additive} give $\der B=b$.

If every exponent exceeds $1$, all powers in $B$ have positive exponent, there
is no logarithmic term, and $B$ is infinitesimal. If the least exponent
$\alpha_{0}$ is less than $1$, the corresponding primitive term has negative
exponent and nonzero coefficient; it dominates the remaining power terms and
also $\log\omega$, since $\log\omega\prec\omega^{\delta}$ for every real
$\delta>0$. So $B$ is infinite. If the least exponent is exactly $1$, the
logarithmic term has nonzero coefficient while all remaining power terms are
infinitesimal, so $B$ is again infinite. Adding a real constant changes none of
these conclusions. Apply \cref{diff:thm:phasecriterion}.
\end{proof}

The logarithmic comparison used here is a standard consequence of the surreal
exponential; equivalently it follows from the normal-form description of
$\log\omega$ \cite{BM,Gonshor}. The threshold is a \emph{strict} inequality at
exponent $1$, and the strictness is not an artefact of the statement: it is
the $\log\omega$ term.

\begin{corollary}[Power threshold]
\label{diff:cor:power}
For $p\in\R$, the equation $\der y=i\omega^{-p}y$ has a nonzero surcomplex
solution exactly when $p>1$, and then all solutions are
\[
 y=c\,\Emm\!\left(\frac{i\omega^{1-p}}{1-p}\right),\qquad c\in\C .
\]
Equivalently $\omega^{-p}\in\Acl\iff p>1$.
\end{corollary}
\begin{proof}
For $p\neq1$ a primitive of $\omega^{-p}$ is $\omega^{1-p}/(1-p)$, which is
infinitesimal for $p>1$ and infinite for $p<1$; for $p=1$ a primitive is
$\log\omega$, infinite. Apply
\cref{diff:thm:phasecriterion,diff:thm:power}.
\end{proof}

\begin{corollary}[Iterated-logarithm thresholds at every finite depth]
\label{diff:cor:logscales}
With $\ell_{0}=\omega$, $\ell_{n+1}=\log\ell_{n}$ and, for an ordinary integer
$n\geq0$ and real $p$,
\[
 b_{n,p}=\frac{1}{\ell_{0}\ell_{1}\cdots\ell_{n-1}\ell_{n}^{\,p}}
\]
(the prefix product empty for $n=0$), the equation $\der y=ib_{n,p}y$ has a
nonzero solution in $\SC$ exactly when $p>1$. An accumulated phase is
$\ell_{n}^{1-p}/(1-p)$ for $p\neq1$ and $\ell_{n+1}$ for $p=1$.
\end{corollary}
\begin{proof}
By \cref{diff:eq:logtower}, $\der\ell_{n}=1/(\ell_{0}\cdots\ell_{n-1})$;
differentiate the displayed expressions using the real power and logarithm
rules of \cref{diff:eq:basicder}. For $p>1$ the accumulated phase is
infinitesimal; for $p\leq1$ it is infinite, each $\ell_{n}$ being positive
infinite. Apply \cref{diff:thm:phasecriterion}.
\end{proof}

These are finite iterated constructions: no infinite product and no transfinite
logarithm tower is used. (Added with member~13: \cref{diff:cp:part} imports
that tower, and \cref{diff:cp:cor:nonreal} carries the obstructed boundary case
$p=1$ to every ordinal depth.) The tests resemble familiar convergence thresholds for
improper integrals, but their proof here is algebraic differentiation and
comparison of surreal sizes. No improper real integral, no limiting process
and no integration to a numerical endpoint occurs anywhere.

\subsection{Why the criterion cannot be globalized by valuation}
\label{diff:sub:noglobalvaluation}

\begin{example}[The separating pair]
\label{diff:ex:logseparator}
Put $L=\log\omega$. Then
\[
 \der(\log L)=\frac{1}{\omega L},
 \qquad
 \der\!\left(-L^{-1}\right)=\frac{1}{\omega L^{2}} .
\]
The first accumulated phase, $\log\log\omega$, is infinite; the second,
$-1/\log\omega$, is infinitesimal. Therefore
\begin{equation}\label{diff:eq:logsep}
 \der y=\frac{i}{\omega\log\omega}\,y
 \ \text{has no nonzero solution},
 \qquad
 \der y=\frac{i}{\omega(\log\omega)^{2}}\,y
 \ \text{does}.
\end{equation}
Both coefficients have $v(\Imm a)>1$ in the full Hahn valuation normalized by
$v(t)=1$: indeed $v(L)<0$, so $v(t/L)=1-v(L)>1$.
\end{example}

\begin{warning}[The threshold theorem is restricted, and must stay restricted]
\label{diff:warn:noglobalization}
\Cref{diff:thm:power} is stated over a real-exponent Hahn field and the
restriction is substantive. By \cref{diff:ex:logseparator}, even
$v(\Imm a)>1$ is \emph{not} sufficient over all of $\No$: the missing
information lies in the slower logarithmic scale, which a real-exponent
workspace cannot see. A computer-algebra backend must not apply
\cref{diff:thm:power} outside its declared coefficient domain, and the
support threshold must not be merged into the general-$\Gamma$ resolvent of
\cref{diff:thm:resolvent}, whose hypotheses are different. A primitive-based
representation handles all of these scales uniformly; a single power cutoff
does not.
\end{warning}

\begin{table}[htbp]
\centering\small
\caption{Distinguishable situations at two and three scales. $L=\log\omega$.}
\label{diff:tab:scales}
\begin{tabularx}{\textwidth}{@{}L{0.22\textwidth}L{0.22\textwidth}Y@{}}
\toprule
Imaginary coefficient $b$ & An accumulated phase $B$ &
Nonzero solution of $\der y=iby$? \\
\midrule
$1$ & $\omega$ & No \\
$t=\omega^{-1}$ & $L$ & No, although $b$ is infinitesimal \\
$t^{2}=\omega^{-2}$ & $-t$ & Yes: $\Emm(-it)$ \\
$t^{3/2}$ & $-2t^{1/2}$ & Yes: $\Emm(-2it^{1/2})$ \\
$1/(\omega L)$ & $\log L$ & No: a slower divergence still obstructs \\
$1/(\omega L^{2})$ & $-1/L$ & Yes: $\Emm(-i/L)$ \\
$\der\eta$ with $\eta\in\mm$ & $\eta$ & Yes: $\Emm(i\eta)$ \\
$\omega^{3}+i\omega^{-2}$ (whole coefficient) & --- &
Yes: $c\exp(\omega^{4}/4)\Emm(-i/\omega)$ \\
\bottomrule
\end{tabularx}
\end{table}

\subsection{An exact decision procedure for rational coefficients}
\label{diff:sub:rationalcriterion}

The obstruction map is a semantic classification, not a universal effective
procedure. For rational functions over the ordinary constants the condition is
nevertheless decidable by finite algebra.

\begin{theorem}[Rational phase criterion]
\label{diff:thm:rational}
Let $b(X)\in\R(X)$ be nonzero, in reduced form $b=P/Q$. Then
\begin{equation}\label{diff:eq:rationalcriterion}
 b(\omega)\in\Acl
 \quad\Longleftrightarrow\quad
 \deg Q-\deg P\geq2 .
\end{equation}
Hence for $q(\omega)=a(\omega)+ib(\omega)$ with $a,b\in\R(X)$ this is exactly
the criterion for a nonzero solution of $\der y=q(\omega)y$. The case $b=0$ is
always admissible.
\end{theorem}
\begin{proof}
Let $p=\deg Q-\deg P$. Expanding at infinity gives
$b(\omega)=ct^{p}+(\text{higher integer powers of }t)$ with $c\in\R^{\times}$.
If $p\geq2$, write $b(\omega)=\sum_{n\geq2}b_{n}t^{n}$; then
$\eta=-\sum_{n\geq2}b_{n}t^{n-1}/(n-1)$ is infinitesimal and differentiates
strongly to $b(\omega)$, so $b(\omega)\in\Acl$. If $p\leq1$, the leading term
gives $\abs{b(\omega)}\geq\abs c t^{p}/2$; were $b(\omega)\in\Acl$, real
linearity and the convexity of \cref{diff:thm:idealstructure}(iii) would force
$t^{p}\in\Acl$, contradicting \cref{diff:cor:power}. The assertion for $q$ is
\cref{diff:thm:phasecriterion}.
\end{proof}

\begin{table}[htbp]
\centering\small
\caption{The rational criterion on examples.}
\label{diff:tab:rational}
\begin{tabular}{@{}lll@{}}
\toprule
$b(\omega)$ & Finite accumulated phase? & Reason, or the phase \\
\midrule
$0$ & yes & $0$ \\
$1$ & no & $\omega$ is infinite \\
$1/\omega$ & no & $\log\omega$ is infinite \\
$1/\omega^{2}$ & yes & $-1/\omega$ \\
$\omega/(1+\omega^{2})$ & no & degree gap $1$ \\
$1/(1+\omega^{2})$ & yes & degree gap $2$ \\
$(1-\omega^{2})/(1+\omega^{2})^{2}$ & yes & $\omega/(1+\omega^{2})$ \\
\bottomrule
\end{tabular}
\end{table}

\begin{warning}[The theorem is over $\R(X)$; the implementation is over $\Q(X)$]
\label{diff:warn:rationalimpl}
\Cref{diff:thm:rational} is an algebraic statement over $\R(X)$. The delivered
executable restricts its input to exact rational coefficients so that
polynomial arithmetic and degree comparison are effective, and it
\emph{raises} on inputs outside $\Q(\text{variable})$ --- other functions,
nonrational coefficients, unbound parameters, non-numeric tokens --- rather
than reporting them as mathematical nonexistence. No algorithm for an
arbitrary real coefficient is asserted. See
\cref{diff:sub:certificates,diff:app:verification}.
\end{warning}

\begin{example}[An algebraic solution with a small angular rate]
\label{diff:ex:algebraicphase}
For $b=1/(1+\omega^{2})$ the unit-modulus \emph{algebraic} number
\[
 y=\frac{1+i\omega}{\sqrt{1+\omega^{2}}}
\]
satisfies $\der y=iby$, since
\[
 y^{\dagger}=\frac{i}{1+i\omega}-\frac{\omega}{1+\omega^{2}}
 =\frac{i}{1+\omega^{2}} .
\]
This example needs neither a named inverse trigonometric function nor an
infinite-phase exponential.
\end{example}

\begin{example}[A rational unit-circle (Cayley) coordinate]
\label{diff:ex:cayleyunit}
With $t=\omega^{-1}$ put $u=(1+it)/(1-it)$. Then $u\bar u=1$, $\st(u)=1$, and
\[
 u^{\dagger}=\frac{-2it^{2}}{1+t^{2}} .
\]
Its infinitesimal phase residue is $2\arctan(t)$, understood as the
infinitesimal Taylor evaluation of \cref{diff:lem:smallchain}; indeed
$\der(2\arctan t)=-2t^{2}/(1+t^{2})$. So an admissible solution need not be
stored as an infinite sum: an exact rational expression can already carry the
required phase, and its full expansion is optional. The coefficient domain here
is Gaussian rationals in $t$.
\end{example}

\begin{example}[A solvable coefficient with an infinite Hahn expansion]
\label{diff:ex:hahnphase}
For $q=i/\omega^{2}=it^{2}$ the normalized solution is
\[
 y=\Emm(-it)=1-it-\frac{t^{2}}{2}+\frac{it^{3}}{6}+\frac{t^{4}}{24}-\cdots,
\]
whose support lies in $\N$, so the family is summable. The identity
$\der y=it^{2}y$ is exact, not asymptotic. A finite truncation
$E_{N}(t)=\sum_{n\leq N}(-i)^{n}t^{n}/n!$ has the exact residual
\begin{equation}\label{diff:eq:truncresidual}
 \der E_{N}-it^{2}E_{N}=-i\,\frac{(-i)^{N}}{N!}\,t^{N+2},
\end{equation}
which is what a finite verifier can check; it is not the whole solution.
\end{example}
\section{The maximal derivation-compatible exponential strip}
\label{diff:sec:strip}

\subsection{A surjective partial exponential}
\label{diff:sub:stripexp}

Set
\[
 \Strip:=\No+i\Ofin\subset\SC .
\]
This is an additive subgroup and a real vector subclass, and it is
$\der$-stable precisely because $\der\Ofin\subseteq\mm$
(\cref{diff:lem:smallnegative}). It is \emph{not} a subring:
$\omega\in\Strip$ and $i\in\Strip$, while $i\omega\notin\Strip$. Note
carefully that $\Strip$ involves $\Ofin$, whereas the logarithmic-derivative
image of \cref{diff:thm:logimage} involves $\der\Ofin=\Acl$; see
\cref{diff:warn:collapse}. Define
\begin{equation}\label{diff:eq:stripexp}
 \Estrip(x+i\theta)=\exp(x)\,\cis(\theta)
 =\exp(x)\,e^{i\st(\theta)}\,\Emm\bigl(i(\theta-\st(\theta))\bigr),
 \qquad x\in\No,\ \theta\in\Ofin .
\end{equation}

\begin{theorem}[The strip exponential]
\label{diff:thm:strip}
$\Estrip:(\Strip,+)\to(\SC^{\times},\cdot)$ is a surjective group homomorphism
with kernel exactly $2\pi i\Z$, it extends the ordinary complex exponential,
and it satisfies
\begin{equation}\label{diff:eq:compat}
 \der\Estrip(w)=\Estrip(w)\,\der w\qquad(w\in\Strip).
\end{equation}
Hence $\Strip/(2\pi i\Z)\cong\SC^{\times}$ as abelian groups. Moreover
$\der\Strip=\No+i\Acl$.
\end{theorem}
\begin{proof}
Multiplicativity follows from the real exponential law, the additivity of
standard part on finite elements and the group law of
\cref{diff:thm:cisgroup}; the derivative identity follows from
\ref{diff:BM:exp} and \cref{diff:eq:cisder}, noting that
$\st(\theta)$ and $e^{i\st(\theta)}$ are differential constants. If
$\Estrip(x+i\theta)=1$, taking moduli gives $\exp x=1$, so $x=0$, and then
$\cis\theta=1$ gives $\theta\in2\pi\Z$; conversely these lie in the kernel.
For surjectivity, write $z=\rho c\Emm(i\eta)$ by \cref{diff:cor:polar} and
choose an ordinary real $r$ with $e^{ir}=c$; then
$z=\Estrip(\log\rho+i(r+\eta))$. This uses the ordinary circle only, not a
branch on the whole surcomplex plane. Finally $\der\No=\No$ and
$\der\Ofin=\Acl$ give the image assertion.
\end{proof}

So $\Estrip$ reaches every nonzero surcomplex number while restricting the
allowable \emph{exponents}: surjectivity does not require an exponent for
every element of $\SC$. Its kernel consists of ordinary integer periods, not
of a surreal integer part. Every nonzero $z$ has strip logarithms, all
differing by an ordinary period $2\pi in$, and all with the \emph{same}
intrinsic derivative $z^{\dagger}$ --- a conclusion that is canonical even
when no choice of logarithm branch is.

\subsection{Maximality is pointwise}
\label{diff:sub:maximality}

\begin{theorem}[Exact domain of the exponential differential equation]
\label{diff:thm:maximalstrip}
For $w=x+iB\in\SC$ with $x,B\in\No$, the following are equivalent.
\begin{enumerate}[label=\textup{(\arabic*)},leftmargin=2.4em]
\item There exists $y\in\SC^{\times}$ with $y^{\dagger}=\der w$.
\item $B$ is finite.
\item $w\in\Strip$.
\end{enumerate}
Consequently no map taking nonzero surcomplex values and satisfying
\cref{diff:eq:compat} at every point of its domain can extend $\Estrip$ to
even \emph{one} exponent outside $\Strip$.
\end{theorem}
\begin{theorem}[Maximality restated]
\label{diff:thm:maximalstripalt}
Equivalently, $\der^{-1}(\Acl)=\Ofin$: the condition in
\cref{diff:thm:maximalstrip}(1) is $\der B\in\Acl$, and
\cref{diff:thm:idealstructure}(viii) identifies its solution class as
$\Ofin$.
\end{theorem}
\begin{proof}[Proof of \cref{diff:thm:maximalstrip,diff:thm:maximalstripalt}]
By \cref{diff:thm:logimage}, condition (1) is equivalent to
$\Imm(\der w)=\der B\in\Acl$. Every accumulated phase of $\der B$ is
$B+r$ with $r\in\R$, so by \cref{diff:thm:phasecriterion} such a phase is
finite exactly when $B$ is finite; equivalently
$\der^{-1}(\Acl)=\Ofin$ by \cref{diff:thm:idealstructure}(viii). The
remaining equivalence is the definition of $\Strip$, and
\cref{diff:thm:strip} supplies a value at every point of $\Strip$. The
reasoning is pointwise: it excludes even an isolated proposed value outside
the strip, irrespective of any group law, homomorphism axiom, or regularity
imposed on the map.
\end{proof}

The claimed strength is exactly this pointwise form. It is stronger than the
failure of one candidate formula: at an excluded point there is \emph{no}
admissible nonzero value at all. On the other hand the differential law alone
is not claimed to pin down every normalization on $\Strip$: ratios of
admissible values are differential constants, and \cref{diff:eq:stripexp} is
the specified canonical construction with the ordinary complex normalization.
A symbolic system should therefore attach the domain predicate
$\Imm w\in\Ofin$ to $\Estrip$ rather than advertise an exponential
\emph{operation} on a differential subfield.

\subsection{No globally compatible exponential, in three strengths}
\label{diff:sub:globalexp}

\begin{corollary}[Phase-independent nonexistence]
\label{diff:cor:noglobalPsi}
There is no map $\Psi:\SC\to\SC^{\times}$ with
\begin{equation}\label{diff:eq:globalcompat}
 \der\Psi(z)=\Psi(z)\,\der z\qquad\text{for every }z\in\SC .
\end{equation}
No homomorphism, continuity, analyticity, conjugation or normalization
hypothesis is needed. Likewise there is no map $E:\No\to\SC^{\times}$ with
$\der(E(x))=iE(x)\der x$ for every $x\in\No$, and no extension of the
finite-angle $\cis$ to all real surreals can satisfy
$\der\cis(x)=i(\der x)\cis(x)$ everywhere.
\end{corollary}
\begin{proof}
Evaluate \cref{diff:eq:globalcompat} at $z=i\omega$: the value
$y=\Psi(i\omega)$ would be nonzero with $\der y=iy$, contradicting
\cref{diff:cor:oscillator} (or \cref{diff:prop:coarse}, or
\cref{diff:thm:maximalstrip}). For the second statement evaluate at
$x=\omega$.
\end{proof}

Because no structural hypothesis is used, this cannot be repaired by a
cleverer choice of infinite circular phase. It is stronger in its
phase-independence than an obstruction proved for one chosen canonical
exponential, and it is compatible with the repository's existing
derivation-obstruction theorem \cite{RepoAnalysis}, which it strengthens
rather than contradicts.

The defect can also be quantified, uniformly over all conventions of one
natural shape.

\begin{proposition}[Character-independent invariant defect]
\label{diff:prop:defect}
Let $\chi:(\Pcl,+)\to\Tor(\No)$ be \emph{any} group homomorphism, and for
$y=P+r+\eta$ with $P\in\Pcl$, $r\in\R$, $\eta\in\mm$ define the phase
extension
\[
 E_{\chi}(x+iy)=\exp(x)\,e^{ir}\,\Emm(i\eta)\,\chi(P).
\]
Put $\delta_{\chi}(z)=E_{\chi}(z)^{\dagger}-\der z$ for $z=x+iy$. Then
\begin{equation}\label{diff:eq:defect}
 \Obs\bigl(\Imm\delta_{\chi}(z)\bigr)=-\pp(y).
\end{equation}
\end{proposition}
\begin{proof}
By \cref{diff:thm:unitphase} the individual \emph{value} $\chi(P)$ is
$c_{P}\Emm(i\eta_{P})$ with $c_{P}\in\Tor(\C)$ and $\eta_{P}\in\mm$; we
differentiate that value only, and never differentiate $\chi$ as a function.
Its logarithmic derivative is $i\der\eta_{P}$, so
$E_{\chi}(z)^{\dagger}=\der x+i\der\eta+i\der\eta_{P}$ and
$\delta_{\chi}(z)=i(\der\eta_{P}-\der P)$. Apply $\Obs$, using
$\der\eta_{P}\in\Acl$ and $\Obs(\der P)=P$ from
\cref{diff:thm:idealstructure}(vi).
\end{proof}

So no choice of character repairs the obstruction at a purely infinite
imaginary exponent: the defect may change with $\chi$, but its purely infinite
class cannot. The trivial character is one admissible convention; nothing here
assumes that every conceivable global exponential has this form.

\begin{warning}[What is not being denied]
\label{diff:warn:notdenied}
Ehrlich and Kaplan construct canonical global surcomplex exponential and
circular functions using an initial integer part, with kernel described by
$2\pi i$ times that integer part \cite[\S11]{EK}; the repository carries that
construction alongside its phase-extension discussion \cite{RepoTrig}. These
are algebraic and structural choices, not assertions of compatibility with
\cref{diff:eq:compat}. The no-go theorems above do not invalidate those
functions, do not say that a symbol such as $\operatorname{Exp}(i\omega)$
cannot be assigned a surcomplex value, and do not settle their separate
robustness questions. They say only that \emph{whatever} nonzero value is
assigned, its Berarducci--Mantova derivative cannot equal $i$ times itself.
The trigonometry report's global phase constructions likewise do not assert
the scalar chain rule and are not refuted.
\end{warning}

\begin{example}[A phase convention is not a differential solution]
\label{diff:ex:phaseconvention}
A convention that discards the purely infinite part of an angle may assign
phase $1$ to $\omega$. That value is algebraically meaningful under its own
convention, but its intrinsic derivative is $\der1=0$, whereas differential
exponential compatibility at $i\omega$ would require $\der1=i$. The
constant-value convention makes the distinction especially transparent: a
nonzero ordinary constant has derivative zero, not $i$ times itself. So a
phase convention and a differential solution are different objects even when
both are written $e^{i\omega}$.
\end{example}

\begin{remark}[There is no finite clock]
\label{diff:rem:noclock}
Ordinary sine and cosine lifts on finite halos obey the chain rule with the
factor $\der z$, and for finite $z$ that factor is infinitesimal: there is no
finite surreal $z$ with $\der z=1$. The forbidden oscillator would require a
genuinely infinite intrinsic clock.
\end{remark}

\begin{remark}[The formal exponential series is not summable at $\omega$]
\label{diff:rem:nonsummableexp}
The formal series $F(Z)=\sum_{n\geq0}(iZ)^{n}/n!\in\C[[Z]]$ satisfies
$\derZ F=iF$, an ordinary formal differential equation with no bearing on
$\der y=iy$ until one evaluates. At $Z=\omega$ the terms carry monomials
$\omega^{n}=t^{-n}$, whose additive valuation support
$\{0,-1,-2,\dots\}$ is \emph{not} well ordered; the family therefore has no
Hahn sum and does not define an element of $\SC$. The same holds for
$\sum_{n\geq0}(i\omega)^{n}/n!=\sum_{n}(i^{n}/n!)t^{-n}$, and choosing a
larger exponent group does not repair that support defect. At an
infinitesimal argument the series \emph{is} summable, but then the chain rule
involves the derivative of that infinitesimal argument rather than the value
$1$. This is how formal analytic solutions are reconciled with the
obstruction.
\end{remark}

\section{Inhomogeneous equations: existence is a separate question}
\label{diff:sec:inhom}

\begin{theorem}[Variation of constants]
\label{diff:thm:variation}
Suppose $y_{0}\neq0$ satisfies $\der y_{0}=qy_{0}$ in $\SC$ --- equivalently,
by \cref{diff:thm:phasecriterion}, suppose $\Phi(q)=0$ and take
$y_{0}=\exp(\Izero a)\Emm(i\Izero b)$ for $q=a+ib$. Then for every $f\in\SC$
all solutions of $\der y-qy=f$ are
\begin{equation}\label{diff:eq:variation}
 y=y_{0}\bigl(\Izero(f/y_{0})+c\bigr),\qquad c\in\C .
\end{equation}
In particular $\der-q$ is surjective on $\SC$ in this case. If the homogeneous
equation has no nonzero solution, then for a given $f$ the inhomogeneous
equation has \emph{at most one} solution in $\SC$.
\end{theorem}
\begin{proof}
Put $y=y_{0}u$; since $\der y_{0}=qy_{0}$ the equation becomes
$y_{0}\der u=f$, and \ref{diff:BM:surj} with
\cref{diff:prop:normprimitive} gives $u=\Izero(f/y_{0})+c$ with $c\in\C$, these
being all primitives by \cref{diff:prop:complexification}. Substitution
verifies \cref{diff:eq:variation}. For the last clause, the difference of two
solutions would solve the obstructed homogeneous equation.
\end{proof}

\begin{warning}[``At most one'' must not be replaced by ``none'']
\label{diff:warn:atmostone}
The absence of a nonzero homogeneous solution does \emph{not} imply
inhomogeneous nonexistence. Particular solutions may exist and then be unique:
\begin{align*}
 (\der-i)\,i&=1, &
 (\der-i)(i\omega+1)&=\omega,\\
 \der y-iy=1
 &\ \text{has the unique } y=i, &
 \der y=iy-i
 &\ \text{has the solution } y=1,\\
 \der^{2}y+y=\omega^{2}
 &\ \text{has the unique } y=\omega^{2}-2, &
 \der^{2}y+y=\omega^{3}
 &\ \text{has the unique } y=\omega^{3}-6\omega,
\end{align*}
existence verified by differentiation and uniqueness by
\cref{diff:cor:oscillator}. An integrating-factor expression using an
undefined value $e^{i\omega}$ is unnecessary and would obscure the correct
reasoning; the integrating factor simply does not exist in $\SC$, so variation
of constants yields no existence conclusion there. This argument alone gives
no general inhomogeneous existence theorem outside $\Phi(q)=0$. The additional
import in \cref{diff:thm:matsurjectivity} will supply existence for every
forcing in \cref{diff:part:matrix}. Note also that having a trivial kernel is not the
same as failing to be surjective: \cref{diff:thm:resolvent} exhibits an
operator that is bijective on a workspace although its homogeneous equation
has no nonzero solution anywhere in $\SC$.
\end{warning}

\begin{proposition}[Polynomial forcing, and a terminating inverse]
\label{diff:prop:polyforcing}
Let $f\in\C[\omega]$ of degree $N$ and $\lambda\in\C^{\times}$. Then
\begin{equation}\label{diff:eq:polyforcing}
 y=-\sum_{k=0}^{N}\frac{\der^{k}f}{\lambda^{k+1}}
\end{equation}
solves $(\der-\lambda)y=f$; for $f=0$ read the sum as $0$. If $\lambda$ is
nonreal this is the unique solution in $\SC$; if $\lambda$ is real, all
solutions are obtained by adding $c\exp(\lambda\omega)$, $c\in\C$. In
particular, with $\lambda=i$,
\begin{equation}\label{diff:eq:polyinverse}
 (\der-i)^{-1}f=i\sum_{k=0}^{N}(-i)^{k}\der^{k}f .
\end{equation}
More generally, for real $\mu$ and nonzero $P\in\C[X]$,
$P(\der)\bigl(e^{\mu\omega}q(\omega)\bigr)=e^{\mu\omega}P(\mu+D)q(\omega)$
where $D$ is ordinary polynomial differentiation; writing
$P(\mu+X)=X^{r}Q(X)$ with $Q(0)\neq0$, the operator $Q(D)^{-1}$ is a finite
nilpotent expansion on a finite-dimensional polynomial space and repeated
polynomial integration solves the remaining $D^{r}$ equation. Hence every
forcing term that is a finite sum of real-exponent exponential polynomials has
a particular solution of the same kind.
\end{proposition}
\begin{proof}
Applying $\der-\lambda$ to the finite sum \cref{diff:eq:polyforcing} makes all
intermediate terms cancel, and $\der^{N+1}f=0$ leaves $f$. For
\cref{diff:eq:polyinverse} use $\der-i=-i(1+i\der)$ together with
$\der^{N+1}f=0$, so the finite geometric identity applies: this is a true
terminating formula, with no infinite summation and no asymptotic truncation.
On ordinary polynomials in $\omega$ the operator $\der$ \emph{is} ordinary
polynomial differentiation, since $\der\omega=1$; the displayed conjugation
identity then follows from \ref{diff:BM:exp}, and
\cref{diff:cor:oscillator,diff:thm:constant} give the uniqueness statements.
\end{proof}

\section{All homogeneous constant-coefficient equations}
\label{diff:sec:constant}

Every ordinary complex number is a differential constant, so for
$P(T)\in\C[T]$ the substitution $T\mapsto\der$ gives a well-defined
$\C$-linear operator on $\SC$.

\begin{lemma}[Kernels of powers and real shifts]
\label{diff:lem:kernels}
For every ordinary integer $m\geq1$,
\[
 \Ker(\der^{m})=\Span_{\C}\{1,\omega,\dots,\omega^{m-1}\}
 =\{P(\omega):P\in\C[X],\ \deg P<m\},
\]
with unique coefficients. For $\lambda\in\R$,
$\Ker(\der-\lambda)^{m}=e^{\lambda\omega}\Ker(\der^{m})$. For
$\lambda\in\C\setminus\R$, every power of $\der-\lambda$ is injective.
\end{lemma}
\begin{proof}
The case $m=1$ is the constant-field statement. Inductively, if
$\der^{m}y=0$ then $\der^{m-1}y=c\in\C$, and
$y-c\omega^{m-1}/(m-1)!$ is killed by $\der^{m-1}$, hence is a polynomial of
degree at most $m-2$; the converse is repeated differentiation, using
$\der\omega=1$. Coefficients are unique because a nonzero
ordinary-coefficient polynomial cannot vanish at the infinite $\omega$, its
highest-degree term dominating all lower ones; equivalently, $\omega$ is
transcendental over the ordinary constants. For real $\lambda$,
\ref{diff:BM:exp} gives
\begin{equation}\label{diff:eq:conjugate}
 (\der-\lambda)\bigl(e^{\lambda\omega}f\bigr)=e^{\lambda\omega}\der f,
\end{equation}
and iteration proves the second assertion. For nonreal $\lambda$,
\cref{diff:prop:coarse} or \cref{diff:cor:oscillator} makes $\der-\lambda$
injective --- alternatively multiply by $e^{-\alpha\omega}$ to reduce to
$\der v=i\beta v$ --- and a finite power of an injective operator is
injective.
\end{proof}

\begin{theorem}[Constant-coefficient solution classification]
\label{diff:thm:constant}
Let $P\in\C[T]$ be nonzero and let $m_{\lambda}$ be the multiplicity of a root
$\lambda$. The complete solution space in $\SC$ of $P(\der)y=0$ is
\begin{equation}\label{diff:eq:constantkernel}
 \Ker P(\der)=
 \bigoplus_{\substack{\lambda\in\R\\ P(\lambda)=0}}
 e^{\lambda\omega}\,
 \Span_{\C}\{1,\omega,\dots,\omega^{m_{\lambda}-1}\},
\end{equation}
of complex dimension $\sum_{\lambda\in\R,\,P(\lambda)=0}m_{\lambda}$. Nonreal
characteristic roots contribute nothing; a nonzero constant polynomial gives
the zero space. In particular an operator of positive order may have zero
kernel, and the solution dimension may be strictly less than $\deg P$.
\end{theorem}
\begin{proof}
Factor $P$ over $\C$ into pairwise coprime primary factors
$p_{\lambda}(T)=(T-\lambda)^{m_{\lambda}}$; a nonzero scalar factor does not
change the kernel. For pairwise coprime $f,g\in\C[T]$ choose $A,B$ with
$Af+Bg=1$; then for $v\in\Ker(f(\der)g(\der))$,
\[
 v=A(\der)f(\der)v+B(\der)g(\der)v,
\]
whose first summand is killed by $g(\der)$ and second by $f(\der)$, and whose
two kernels intersect in $0$ by the same identity. This coprime-kernel
decomposition is valid for any linear operator on any vector space; no
finite-dimensionality is needed. Induction over the primary factors, or
equivalently the Chinese-remainder idempotents modulo $P$, decomposes
$\Ker P(\der)$ as the direct sum of the primary kernels. Apply
\cref{diff:lem:kernels}: nonreal-root summands vanish and real-root summands
are those displayed. (Equivalently, factor $P=cQR$ with $Q$ the product of the
real-root factors and $R$ that of the nonreal ones; $R(\der)$ is injective and
commutes with $Q(\der)$, so $\Ker P(\der)=\Ker Q(\der)$.) Each displayed
summand has dimension $m_{\lambda}$, and directness gives the count.
\end{proof}

\begin{example}[Order four, dimension two; order five, dimension three]
\label{diff:ex:constantexamples}
In $\SC$:
\begin{align*}
 \der^{2}y+y=0 &\ \Longrightarrow\ y=0,\\
 \der^{2}y-y=0 &\ \Longleftrightarrow\ y=Ae^{\omega}+Be^{-\omega},\\
 (\der-1)^{2}(\der^{2}+1)y=0 &\ \Longleftrightarrow\
 y=(A+B\omega)e^{\omega},\\
 (\der-2)^{2}(\der^{2}+1)y=0 &\ \Longleftrightarrow\
 y=(A+B\omega)e^{2\omega},\\
 (\der-2)^{3}(\der+1)(\der^{2}+1)y=0 &\ \Longleftrightarrow\
 y=e^{2\omega}(A+B\omega+C\omega^{2})+De^{-\omega},\\
 (\der-2)^{2}(\der+1)(\der^{2}+1)y=0 &\ \Longleftrightarrow\
 y=(A+B\omega)e^{2\omega}+Ce^{-\omega},\\
 \der(\der-1)^{2}(\der^{2}+1)y=0 &\ \Longleftrightarrow\
 y=A+e^{\omega}(B+C\omega),
\end{align*}
with all constants in $\C$. The missing modes are in every case exactly the
oscillatory ones. There is no contradiction with the ordinary theorem that an
order-$n$ linear equation has an $n$-dimensional local solution space: that
theorem concerns functions of a coordinate, or solutions in a sufficiently
large differential extension, not solutions in an arbitrary fixed differential
field.
\end{example}

\begin{remark}[Real solutions, and the coefficient hypothesis]
\label{diff:rem:realsolutions}
If $P\in\R[T]$ and real surreal solutions are wanted, the constants in
\cref{diff:eq:constantkernel} must be taken in $\R$: the imaginary part is a
linear combination of the same independent real basis elements and vanishes
only when all its coefficients do. The hypothesis that the coefficients of $P$
are \emph{ordinary complex constants} is essential: for $a\in\SC\setminus\C$
the operators $\der$ and multiplication by $a$ do not commute, so polynomial
factorization in a commuting symbol cannot be transported to differential
operators. \Cref{diff:thm:constant} says nothing about
variable-coefficient operators.
\end{remark}

\begin{corollary}[Constant matrix systems]
\label{diff:cor:matrix}
Let $A\in M_{n}(\C)$ and consider $\der Y=AY$ with $Y\in\SC^{n}$. The
solution space has complex dimension equal to the total algebraic multiplicity
of the \emph{real} eigenvalues of $A$. On a Jordan block $\lambda I+N$ of size
$m$ with $\lambda\in\R$, all solutions are
\begin{equation}\label{diff:eq:jordan}
 Y=e^{\lambda\omega}
 \left(\sum_{k=0}^{m-1}\frac{\omega^{k}N^{k}}{k!}\right)c,
 \qquad c\in\C^{m},
\end{equation}
a finite polynomial in the nilpotent $N$, not an infinite matrix exponential
requiring convergence; its inverse replaces $\omega$ by $-\omega$. On a block
with nonreal eigenvalue only $Y=0$ occurs. An invertible fundamental matrix in
$\GL_{n}(\SC)$ exists if and only if every eigenvalue of $A$ is real.
\end{corollary}
\begin{proof}
An ordinary complex similarity is constant for $\der$, so pass to Jordan form.
On a block with nonreal $\lambda$ the last coordinate satisfies
$(\der-\lambda)y_{m}=0$ and is zero by \cref{diff:lem:kernels}; upward
substitution through the block makes every coordinate zero. On a real block,
the displayed finite matrix polynomial $B$ satisfies $\der B=NB$ and is
invertible, so multiplication by $e^{\lambda\omega}B$ carries constant vectors
to solutions and its inverse carries solutions to constant vectors. Add the
block dimensions.
\end{proof}

For the real rotation matrix
$\left(\begin{smallmatrix}0&-1\\1&0\end{smallmatrix}\right)$ the solution space
is zero, while a unipotent Jordan block has a fundamental matrix that is a
polynomial in $\omega$. The contrast isolates oscillation, not general
complexity, as the obstruction, and it already rules out a universal
matrix-exponential constructor in the intrinsic field.

\section{Triangular systems, and the limits of a determinant test}
\label{diff:sec:triangular}

Let $A=(a_{jk})$ be an upper-triangular $n\times n$ matrix over $\SC$, with
$n$ an ordinary positive integer. A \emph{fundamental matrix} is
$Y\in\GL_{n}(\SC)$ with $\der Y=AY$, entrywise.

\begin{theorem}[Triangular fundamental-matrix criterion]
\label{diff:thm:triangular}
A fundamental matrix in $\GL_{n}(\SC)$ exists if and only if
\begin{equation}\label{diff:eq:diagcriterion}
 \Phi(a_{jj})=0\qquad(1\leq j\leq n).
\end{equation}
When it does, an upper-triangular one can be constructed by finitely many
first-order homogeneous integrations and additive primitives.
\end{theorem}
\begin{proof}
\emph{Sufficiency.} Choose $s_{j}\neq0$ with $\der s_{j}=a_{jj}s_{j}$ by
\cref{diff:thm:phasecriterion}, put $Y_{jj}=s_{j}$ and $Y_{ij}=0$ for $i>j$.
For each column $j$, construct the entries above the diagonal in
\emph{descending} order of $i$ by
\begin{equation}\label{diff:eq:triangularrecursion}
 Y_{ij}=s_{i}\,\Izero\!\left(
 s_{i}^{-1}\sum_{k=i+1}^{j}a_{ik}Y_{kj}\right),\qquad i<j .
\end{equation}
The entries on the right are already constructed. The product rule gives
$\der Y_{ij}=a_{ii}Y_{ij}+\sum_{k=i+1}^{j}a_{ik}Y_{kj}$, so $\der Y=AY$, and
$\det Y=\prod_{j}s_{j}\neq0$.

\emph{Necessity}, by induction on $n$. The last row of any fundamental matrix
satisfies $\der Y_{nj}=a_{nn}Y_{nj}$, and at least one entry is nonzero, so
$\Phi(a_{nn})=0$. All last-row entries are complex constant multiples of one
nonzero solution, so a change of columns by an invertible constant matrix makes
the last row $(0,\dots,0,s_{n})$; the result is still fundamental and has block
form $\begin{pmatrix}Z&v\\0&s_{n}\end{pmatrix}$ with $\det Z\neq0$, whose
upper-left block satisfies the system for the leading
$(n-1)\times(n-1)$ block of $A$. Induction gives the remaining conditions.
\end{proof}

In particular every strictly upper-triangular system has a unipotent
fundamental matrix over $\SC$: this is a finite iterated-integration fact, and
no infinite exponential series for an arbitrary matrix is needed.

\begin{proposition}[The trace test is necessary and not sufficient]
\label{diff:prop:trace}
For any finite matrix system, triangular or not, Jacobi's formula gives
$\der\det Y=(\tr A)\det Y$ for a fundamental $Y$, so
\begin{equation}\label{diff:eq:tracenecessary}
 \Phi(\tr A)=0
\end{equation}
is necessary. It is not sufficient: for $A=\operatorname{diag}(i,-i)$ the
trace vanishes, yet neither diagonal equation has a nonzero solution and the
vector system has no nonzero solution at all.
\end{proposition}

Opposite phase obstructions can cancel in a determinant without disappearing
from the individual modes, so a determinant-only test discards essential
differential information. For a general matrix system the trace condition is
only the first invariant. The complete one is the \emph{multiset} of phases,
and \cref{diff:part:matrix} supplies it: \cref{diff:thm:matnormal} shows that
every finite system over $\SC$ has a well-defined differential phase spectrum
$\Ph(A)\subset\Pcl$ which is a complete invariant of gauge equivalence, and
\cref{diff:eq:mattrace} recovers \cref{diff:eq:tracenecessary} from it as the
sum
\[
 \Phi(\tr A)=\sum_{P}m_{P}(A)\,P ,
\]
whose failure to determine the multiset is exactly why the trace test is
insufficient. Two of the phase spectrum's consequences belong beside the trace
test. First, by \cref{diff:cor:matinhom} a fundamental matrix exists in
$\GL_{n}(\SC)$ if and only if every member of $\Ph(A)$ is zero, which extends
\cref{diff:thm:triangular} beyond the triangular case. Second, by
\cref{diff:cor:matperturbation}, perturbing every entry of a Hermitian $H$ by
a coefficient with a finite primitive leaves the entire differential unitary
class of $\der y=iHy$ unchanged, with \emph{no} spectral-gap hypothesis --- a
statement the trace invariant cannot see at all, and one that
\cref{diff:prop:sharpperturb} shows is sharp. For non-Hermitian coefficients no
such spectral formula holds, and \cref{diff:ex:matnonnormal} is the
counterexample that marks the line; \cref{diff:sub:nc:scalarscope} separates
this limitation from the complete finite-system classification.
\part{Finite matrix systems over \texorpdfstring{$\SC$}{No[i]}}
\label{diff:part:matrix}

\section{The differential phase spectrum of a finite system}
\label{diff:sec:matrixphase}

\subsection{From one necessary invariant to a complete one}
\label{diff:sub:matbridge}

\Cref{diff:prop:trace} leaves a finite matrix system with exactly one
invariant, $\Phi(\tr A)$, and with a proof that it is not enough:
$\operatorname{diag}(i,-i)$ has vanishing trace obstruction and no nonzero
solution. The missing information is a \emph{multiset}, not a sum. This part
supplies it for every finite system, and computes it from ordinary algebraic
eigenvalues when the coefficient is $iH$ with $H$ Hermitian.

The eighth source manuscript is the origin of this part
(\cref{diff:app:provenance}). Its own audit names the gap in this report's
terms --- general nontriangular systems and higher-rank differential modules
left outside the classification --- and states that the rank-one phase
criterion, the bounded-primitive ideal and the diagonal Galois tori of
\cref{diff:part:scalar,diff:sec:PV} are prior coverage rather than discoveries
of its own. That is how they are used below: nothing already proved in
\cref{diff:part:scalar} is reproved here, and every statement of this part
that coincides with one already in the report is cited, not restated.

Two agreements are worth printing as agreements before anything new is proved.

\begin{remark}[The eighth manuscript refutes the same collapse]
\label{diff:rem:matagreement}
\Cref{diff:warn:collapse} is load-bearing: $\Acl=\der\mm=\der\Ofin$ is a
\emph{strict} subclass of $\mm$, so the criterion is never ``the coefficient
is infinitesimal''. The eighth manuscript opens with exactly that point and
exactly that witness: $i/\omega$ is obstructed because its imaginary part has
the infinite primitive $\log\omega$. Its structure lemma for the ideal ---
order-convexity, the $\Ofin$-module property, containment in $\mm$, and the
separating pair $\omega^{-2}\in\Acl$, $\omega^{-1}\notin\Acl$ --- is
\cref{diff:thm:idealstructure}(iii),(iv), proved there by the same
order-theoretic route from \ref{diff:BM:positivity}. Nothing in this part
weakens that warning, and \cref{diff:prop:sharpperturb} below shows that the
strictness is sharp for matrices as well as for scalars.
\end{remark}

\begin{remark}[The rank-one criterion is imported from \cref{diff:part:scalar}]
\label{diff:rem:matrankone}
The eighth manuscript reproves the exact logarithmic-derivative image
\cref{diff:eq:logimage} by polar decomposition and a formal Hahn logarithm,
and credits it to this collection. It is \cref{diff:thm:logimage} and
\cref{diff:thm:phasecriterion}, with the same proof idea as
\cref{diff:cor:polar}; it is not reproved here. Every uniqueness argument
below rests on the following special case of \cref{diff:thm:phasecriterion},
which is used often enough to deserve a number.
\end{remark}

\begin{lemma}[The rank-one separation test]
\label{diff:lem:matseparation}
Let $P,Q\in\Pcl$ and let $x\in\SC$ satisfy
$\der x=i\,\der(P-Q)\,x$. If $x\neq0$ then $P=Q$ and $x\in\C^{\times}$.
\end{lemma}
\begin{proof}
By \cref{diff:thm:phasecriterion} a nonzero solution forces
$\Phi\bigl(i\der(P-Q)\bigr)=\Obs\bigl(\der(P-Q)\bigr)=0$, which is $P-Q=0$ by
\cref{diff:thm:idealstructure}(vi). Then $\der x=0$, so $x\in\C^{\times}$ by
the constant-field statement of \cref{diff:thm:BM}.
\end{proof}

We also need the order behaviour of $\Obs$ itself, which
\cref{diff:thm:idealstructure} contains but does not display.

\begin{corollary}[The obstruction preserves order]
\label{diff:cor:obsorder}
The map $\Obs:\No\to\Pcl$ is order preserving: $b\leq c$ implies
$\Obs(b)\leq\Obs(c)$. Consequently, for $P,Q\in\Pcl$,
\[
 P<Q\quad\Longleftrightarrow\quad\der P<\der Q ,
\]
and an ordered tuple of real coefficients has an ordered tuple of
obstructions.
\end{corollary}
\begin{proof}
The displayed equivalence is \cref{diff:thm:idealstructure}(vii), which states
that $\der:\Pcl\to\der\Pcl$ is an order-preserving bijection. For the first
assertion, suppose $b\leq c$ and put $P=\Obs(b)$, $Q=\Obs(c)$. If $Q<P$ then
$\der Q<\der P$, and $b-\der P\in\Acl$, $c-\der Q\in\Acl$ by
\cref{diff:thm:idealstructure}(vi),(vii); so
$b-c=(\der P-\der Q)+d$ with $d\in\Acl$ and $\der P-\der Q>0$. Since
$\der P-\der Q=\der(P-Q)$ with $P-Q\in\Pcl\setminus\{0\}$, it does not lie in
$\Acl$, and convexity of $\Acl$
(\cref{diff:thm:idealstructure}(iii)) then gives $b-c>0$, contradicting
$b\leq c$.
\end{proof}

\begin{convention}[The alphabet of this part]
\label{diff:conv:matalphabet}
Throughout \cref{diff:part:matrix}, $n$ is an ordinary positive integer and
all matrices are $n\times n$ over $\SC$ unless said otherwise; $\der$ acts
entrywise, $A^{*}$ is the conjugate transpose for the conjugation of
\cref{diff:prop:complexification}, and $I$ is the identity matrix.
\begin{itemize}[leftmargin=1.6em]
\item $A$ is a coefficient matrix, as in \cref{diff:sec:triangular};
$Y$ is a fundamental matrix.
\item $H$ is a \emph{Hermitian} matrix, $H^{*}=H$. In this report the letter
$H$ otherwise occurs only in the hyphenated adjective ``$H$-field'' and in the
subscripted coefficient fields $H_{\Gamma}=\R((t^{\Gamma}))$; a bare $H$ in
this part is always a matrix and never either of those.
\item $G\in\GL_{n}(\SC)$ is a gauge matrix --- the matrix upgrade of the
scalar gauge $g$ of \cref{diff:eq:gauge} --- and $W$ is a gauge that is
\emph{unitary}, $W^{*}W=I$. The letter $U$ is reserved for an ordinary complex
domain (\cref{diff:warn:concordance}) and is not used for a unitary matrix;
$\Un_{n}(\SC)$ denotes the unitary group.
\item $\mathcal M$ is a horizontal Hermitian metric
(\cref{diff:def:horizontal}), and $\Rot=\left(
\begin{smallmatrix}0&1\\-1&0\end{smallmatrix}\right)$ is the rotation
generator. The sources write $J$ for the latter; this report uses no letter
$J$ (\cref{diff:conv:noJ}).
\item $\Tser$ is the differential field of transseries, which occurs only in
\cref{diff:thm:matsurjectivity}; it is never the unit circle $\Tor(\No)$ and never
a transcendental generator $\mathsf T$.
\item $\Pi$ is the additive \emph{phase lattice} generated inside $\Pcl$ by
the phases of a system, and $\Lambda\subseteq\Z^{n}$ remains the integer
\emph{relation} lattice of \cref{diff:eq:lattice}. The two are different
objects over different groups and are never written with the same letter.
\end{itemize}
\end{convention}

\subsection{Matrix gauge equivalence and the complete phase normal form}
\label{diff:sub:matnormal}

A change of unknown $y=Gz$ with $G\in\GL_{n}(\SC)$ carries $\der y=Ay$ to
$\der z=A^{G}z$, where
\begin{equation}\label{diff:eq:matgauge}
 A^{G}=G^{-1}AG-G^{-1}\der G .
\end{equation}
Call $A,B$ \emph{differentially equivalent} when $B=A^{G}$ for some
$G\in\GL_{n}(\SC)$, and \emph{differentially unitarily equivalent} when $G$
may be taken unitary. The connection term $G^{-1}\der G$ is not optional:
similarity of matrices alone is not differential equivalence, and
\cref{diff:ex:matnonnormal} below turns on exactly that.

Two inputs are needed that \cref{diff:part:scalar} does not contain. Both are
imported, and the import is flagged here and in
\cref{diff:sub:nc:matrixscope}.

\begin{theorem}[Imported: splitting and first-order surjectivity, transferred]
\label{diff:thm:matsurjectivity}
The following two properties hold over $\Tser[i]$, where $\Tser$ is the field of
transseries, and are recorded in the asymptotic-differential-algebra
literature \cite[Lemma 2.4.19]{ADHNorm}:
every monic scalar differential operator of positive order splits into
first-order factors, and every nonzero scalar linear differential operator is
surjective. Both transfer to $\SC$. In particular, for all $a,f\in\SC$ there
is $u\in\SC$ with
\begin{equation}\label{diff:eq:matfirstorder}
 \der u-au=f .
\end{equation}
\end{theorem}
\begin{proof}
For each fixed order the two assertions are first-order sentences of the
ordered differential-field language: expanding a product of first-order
operators uses only the finite rule $\der a=a\der+\der(a)$, so the coefficient
identities are polynomial identities in finitely many coefficients and their
derivatives, and surjectivity of a given order is a first-order sentence with
a nonzero leading coefficient as hypothesis. A complex coefficient is an
ordered pair of real ones, and complex addition, multiplication and
differentiation are definable in that interpretation, so the sentences about
$\Tser[i]$ are sentences about $\Tser$. The elementary embedding
$\Tser\preceq\No$ of \cite{ADH} transfers them to $\No$, hence to $\SC$, for
arbitrary surreal coefficients and not merely for coefficients in the embedded
copy of $\Tser$.
\end{proof}

\begin{warning}[What this import does and does not add]
\label{diff:warn:matimport}
\Cref{diff:thm:matsurjectivity} is an \emph{additional imported input}, on the
same footing as \cref{diff:thm:BM} and not a consequence of anything proved in
\cref{diff:part:scalar}. It must not be read as making $\SC$ differentially
closed, and it must not be read as a homogeneous existence theorem: ``linearly
closed'' there means the splitting property, not that every homogeneous
equation has a full solution set in the base field. Reading it the second way
would make $\der y=iy$ appear solvable and contradict
\cref{diff:cor:oscillator}. It also yields no algorithm: no effective
factorization of an arbitrary surreal-coefficient operator follows from a
first-order transfer (\cref{diff:sub:nc:effectivity}). What it does supply is
\emph{inhomogeneous existence} for first-order scalar equations, which
\cref{diff:thm:variation} deliberately did not assert outside the admissible
locus $\Phi(q)=0$; see \cref{diff:cor:matinhom} and the amended item~N22.
\end{warning}

\begin{lemma}[Cyclic vector, and a rank-one filtration]
\label{diff:lem:matcyclic}
Let $M$ be a differential module over $\SC$ of ordinary finite dimension
$n\geq1$, that is, a free $\SC$-module of rank $n$ with an additive
$\nabla$ satisfying $\nabla(fv)=(\der f)v+f\nabla v$. Then $M$ has a vector
$v$ for which $v,\nabla v,\dots,\nabla^{n-1}v$ is a basis, and $M$ has a
filtration by differential submodules whose successive quotients have
dimension one.
\end{lemma}
\begin{proof}
Fix a basis $e_{0},\dots,e_{n-1}$ and extend $\nabla$ to $M[T]$ by giving the
indeterminate derivative $1$. Define $v_{0}=e_{0}$ and
$v_{k}=e_{k}-\sum_{j<k}\binom{k}{j}\nabla^{k-j}v_{j}$, and put
$v(T)=\sum_{j<n}T^{j}v_{j}/j!$. The Leibniz rule gives
$\nabla^{k}v(T)\big|_{T=0}=e_{k}$ for $k<n$, so the determinant of the columns
$v(T),\dots,\nabla^{n-1}v(T)$ is a polynomial in $T$ with constant term $1$,
hence nonzero. Substituting $T=\omega-c$ for a real $c$ respects
differentiation because $\der(\omega-c)=1$, and a nonzero polynomial has only
finitely many roots while the elements $\omega-c$ are pairwise distinct; choose
$c$ outside that finite exceptional set to get a cyclic vector.

For the filtration, a cyclic vector presents $M$ as $\SC[\der]/\SC[\der]L$
with $L$ monic of order $n$. By \cref{diff:thm:matsurjectivity} write
$L=(\der-b)R$ with $R$ a product of $n-1$ first-order factors. The class of
$R$ is nonzero, spans a one-dimensional differential submodule and satisfies
$\der R\equiv bR$ modulo $L$; the quotient is the module with annihilator $R$,
and induction on $n$ finishes. The dimension counts follow from division by a
monic operator, whose remainder has strictly smaller order.
\end{proof}

\begin{theorem}[Complete phase normal form for finite systems]
\label{diff:thm:matnormal}
For every $A\in\Mat_{n}(\SC)$ there are $G\in\GL_{n}(\SC)$ and
$P_{1},\dots,P_{n}\in\Pcl$ with
\begin{equation}\label{diff:eq:matnormal}
 A^{G}=i\operatorname{diag}(\der P_{1},\dots,\der P_{n}) .
\end{equation}
The multiset
\[
 \Ph(A):=\{P_{1},\dots,P_{n}\}\subset\Pcl
\]
does not depend on $G$, and it is a \emph{complete} invariant of differential
equivalence: two finite systems of the same size are differentially equivalent
if and only if their multisets agree. We call $\Ph(A)$ the \emph{differential
phase spectrum} of $A$, in the sense of \cref{diff:conv:phases}(e).
\end{theorem}
\begin{proof}
Apply \cref{diff:lem:matcyclic} to the connection $\nabla=\der-A$ on $\SC^{n}$;
an adapted basis makes the system upper triangular. Suppose inductively that
the first $n-1$ coordinates have been diagonalized, leaving
\[
 A=\begin{pmatrix}D&b\\0&a\end{pmatrix},
 \qquad D=\operatorname{diag}(d_{1},\dots,d_{n-1}).
\]
For $S=\left(\begin{smallmatrix}I&u\\0&1\end{smallmatrix}\right)$ the
off-diagonal column of $A^{S}$ is $b+Du-ua-\der u$, and each scalar equation
$\der u_{j}-(d_{j}-a)u_{j}=b_{j}$ has a solution by
\cref{diff:eq:matfirstorder}. This kills the column, with no hypothesis on
repetitions among the diagonal entries: when $d_{j}=a$ the equation is
ordinary intrinsic integration, available by \ref{diff:BM:surj}. Induction
gives a diagonal form, and \cref{diff:thm:gauge} applied to each diagonal
entry supplies a further diagonal gauge bringing it to
\cref{diff:eq:matnormal} with $P_{j}=\Phi(a_{j})$.

For independence and completeness, suppose $D_{P}=i\operatorname{diag}(\der
P_{j})$ and $D_{Q}=i\operatorname{diag}(\der Q_{k})$ satisfy
$D_{Q}=(D_{P})^{\Xi}$ for an invertible $\Xi$. Then $\der\Xi=D_{P}\Xi-\Xi D_{Q}$, that
is
\begin{equation}\label{diff:eq:matintertwine}
 \der\Xi_{jk}=i\,\der(P_{j}-Q_{k})\,\Xi_{jk}
\end{equation}
entrywise. By \cref{diff:lem:matseparation} an entry can be nonzero only when
$P_{j}=Q_{k}$, and is then a complex constant. Invertibility of $\Xi$ forces the
two multisets to have the same multiplicities. Conversely equal multisets give
the same normal form after a constant permutation.
\end{proof}

\begin{remark}[This is the matrix upgrade of \cref{diff:thm:gauge}, and of nothing else]
\label{diff:rem:matupgrade}
For $n=1$, \cref{diff:thm:matnormal} is \cref{diff:thm:gauge}. For $n>1$ it is
a genuinely stronger statement, and it is the first classification in this
report of a system that is neither triangular nor constant. It is still
\emph{not} a module-theoretic filtration theorem with computable data: the
normal form is obtained from an existence theorem, the phases are semantic
values of $\Obs$, and \cref{diff:sub:nc:matrixscope} records what remains open.
\end{remark}

\subsection{Solution spaces, projectors, and the refined trace identity}
\label{diff:sub:matprojectors}

For $P\in\Pcl$ write $m_{P}(A)$ for the multiplicity of $P$ in $\Ph(A)$, zero
if $P$ does not occur.

\begin{corollary}[Existence for every forcing, and the exact kernel dimension]
\label{diff:cor:matinhom}
For every $A\in\Mat_{n}(\SC)$ and every $f\in\SC^{n}$ the equation
$\der y-Ay=f$ has a solution in $\SC^{n}$, and its solution set is an affine
space over $\C$ of dimension $m_{0}(A)$. In particular
\[
 \dim_{\C}\Ker\bigl(\der-A\bigr)\big|_{\SC^{n}}=m_{0}(A),
\]
and $\der-A$ is bijective on $\SC^{n}$ exactly when $0\notin\Ph(A)$.
\end{corollary}
\begin{proof}
Transport by the gauge of \cref{diff:thm:matnormal} and solve each scalar
inhomogeneous equation by \cref{diff:eq:matfirstorder}. In the homogeneous
equation a coordinate can be nonzero only at phase $0$, where it is a complex
constant, by \cref{diff:lem:matseparation}. The difference of two particular
solutions is homogeneous.
\end{proof}

\begin{warning}[What \cref{diff:cor:matinhom} changes, and what it does not]
\label{diff:warn:matinhom}
\Cref{diff:thm:variation} proves surjectivity of $\der-q$ only on the
admissible locus $\Phi(q)=0$, and \cref{diff:warn:atmostone} states the
correct weaker conclusion elsewhere: \emph{at most one}, not none.
\Cref{diff:cor:matinhom} upgrades ``at most one'' to ``exactly one'' in the
obstructed case, and it does so \emph{only} by invoking the separately
imported first-order surjectivity of \cref{diff:thm:matsurjectivity}. It does
not follow from the rank-one criterion, it does not follow from
\cref{diff:thm:variation}, and it does not make the homogeneous equation
solvable. The examples of \cref{diff:warn:atmostone} are unaffected and remain
correct; what changes is that their uniqueness is now known to be accompanied
by existence. Item~N22 of \cref{diff:sec:nonclaims} is amended accordingly and
not deleted.
\end{warning}

\begin{proposition}[Canonical phase projectors]
\label{diff:prop:matprojectors}
Let $E_{P}$ be the coordinate projector onto the phase-$P$ block of a normal
form and put $\Pi_{P}=GE_{P}G^{-1}$ for a normalizing gauge $G$. Then
$\Pi_{P}$ is independent of the choice of $G$, and
\[
 \Pi_{P}\Pi_{Q}=\delta_{PQ}\Pi_{P},
 \qquad \sum_{P}\Pi_{P}=I,
 \qquad \der\Pi_{P}=A\Pi_{P}-\Pi_{P}A .
\]
The algebra of horizontal endomorphisms of the system --- matrices $X$ with
$\der X=AX-XA$ --- is isomorphic to $\prod_{P}\Mat_{m_{P}(A)}(\C)$, and for
two systems
\[
 \dim_{\C}\Homd(A,B)=\sum_{P}m_{P}(A)\,m_{P}(B).
\]
Direct sums unite phase multisets, tensor products add phases pairwise, and
duals negate them.
\end{proposition}
\begin{proof}
After aligning the ordered phase blocks, \cref{diff:eq:matintertwine} makes
any change between normalizing gauges a constant block-diagonal matrix, so it commutes with every $E_{P}$; the
displayed identities follow by multiplying and differentiating
$\Pi_{P}=GE_{P}G^{-1}$ and using $\der G=AG-GD_{P}$. The same entrywise
calculation, applied to $\der X=AX-XB$ in normal-form coordinates, identifies
the horizontal homomorphisms with constant matrices supported on matching
phase blocks. The last three statements are the corresponding computations for
$A\oplus B$, $A\otimes I+I\otimes B$ and $-A^{\mathsf T}$, in each case reduced
to the scalar addition rule of \cref{diff:cor:tensor}.
\end{proof}

\begin{corollary}[The determinant sees only the phase sum]
\label{diff:cor:mattrace}
For every finite $A$,
\begin{equation}\label{diff:eq:mattrace}
 \Phi(\tr A)=\sum_{P}m_{P}(A)\,P .
\end{equation}
\end{corollary}
\begin{proof}
Take traces in $A=(\der G)G^{-1}+GD_{P}G^{-1}$ and use Jacobi's identity
$\tr\bigl((\der G)G^{-1}\bigr)=(\det G)^{\dagger}$, whose obstruction is zero
by \cref{diff:thm:phasecriterion}.
\end{proof}

\Cref{diff:eq:mattrace} contains \cref{diff:prop:trace} and explains it: a
fundamental matrix exists only if every phase is zero, so
$\Phi(\tr A)=0$ is necessary, and it fails to be sufficient precisely because
a multiset is not recoverable from its sum. The witness
$\operatorname{diag}(i,-i)$ of \cref{diff:prop:trace} has
$\Ph=\{\omega,-\omega\}$, which sums to zero while
$\{0,0\}$ does too; the two systems are not differentially equivalent, and
\cref{diff:thm:matnormal} says so. The trace condition is the first invariant;
the phase multiset is the complete one.

\subsection{Horizontal metrics and the unitary normal form}
\label{diff:sub:matmetrics}

\begin{definition}[Horizontal metric]
\label{diff:def:horizontal}
A Hermitian $\mathcal M\in\Mat_{n}(\SC)$ is \emph{horizontal} for $A$ when
\begin{equation}\label{diff:eq:horizontal}
 \der\mathcal M+A^{*}\mathcal M+\mathcal M A=0 ,
\end{equation}
and \emph{positive definite} when $v^{*}\mathcal M v>0$ for every nonzero
$v\in\SC^{n}$.
\end{definition}

Equation \cref{diff:eq:horizontal} says formally that $y^{*}\mathcal M y$ is
constant along a solution, in any differential extension carrying the
compatible conjugation. It does not assert that the system has nonzero
solutions in $\SC^{n}$; by \cref{diff:cor:matinhom} it usually does not.

\begin{theorem}[The complete metric cone]
\label{diff:thm:matmetric}
Let $G$ normalize $A$ as in \cref{diff:thm:matnormal}. Every horizontal
Hermitian matrix is uniquely
\begin{equation}\label{diff:eq:matmetriccone}
 \mathcal M=G^{-*}\Bigl(\bigoplus_{P}C_{P}\Bigr)G^{-1},
 \qquad C_{P}\in\Herm_{m_{P}(A)}(\C),
\end{equation}
and is positive definite exactly when every $C_{P}$ is. In particular every
finite system has a positive horizontal metric, and the real dimension of its
space of horizontal Hermitian forms is $\sum_{P}m_{P}(A)^{2}$.
\end{theorem}
\begin{proof}
Put $B=G^{*}\mathcal M G$. Since $\der G=AG-GD_{P}$, differentiating gives
$\der B=D_{P}B-BD_{P}$, that is
$\der B_{jk}=i\,\der(P_{j}-P_{k})B_{jk}$, so
\cref{diff:lem:matseparation} makes every off-phase entry zero and every
within-phase entry a complex constant. Hermitian symmetry leaves exactly the
blocks $C_{P}$. Congruence preserves definiteness, and for a \emph{constant}
Hermitian matrix positivity over $\SC$ is ordinary complex positive
definiteness: diagonalize by an ordinary constant unitary matrix and read off
its real eigenvalues. Taking every $C_{P}=I$ gives existence, and
$\dim_{\R}\Herm_{m}(\C)=m^{2}$.
\end{proof}

\begin{warning}[A positive horizontal metric is not a stability statement]
\label{diff:warn:matlyapunov}
For
\[
 A=\begin{pmatrix}1&1\\0&1\end{pmatrix},
 \qquad
 G=e^{\omega}\begin{pmatrix}1&\omega\\0&1\end{pmatrix},
\]
one has $\der G=AG$, so both phases are zero, and
\[
 \mathcal M=e^{-2\omega}
 \begin{pmatrix}1&-\omega\\-\omega&\omega^{2}+1\end{pmatrix}
\]
is horizontal with $\det\mathcal M=e^{-4\omega}>0$. The ordinary differential
equation with the same printed coefficients has growing solutions. There is no
contradiction, and no Lyapunov or uniform-equivalence claim is made anywhere
in this part: the invariant metric itself has a strongly varying scale, here
$e^{-2\omega}$, and \cref{diff:thm:matmetric} has no hypothesis relating it to
a fixed norm. This is the matrix form of the separation already made in
\cref{diff:warn:separators}(7) between a printed equation read in $\SC$ and the
same equation read in an ordinary coordinate.
\end{warning}

\begin{theorem}[Unitary phase normal form]
\label{diff:thm:matunitary}
If $A^{*}=-A$, the normalizing gauge of \cref{diff:thm:matnormal} can be taken
unitary. For a fixed ordering of the phases, any two unitary normalizing
gauges differ by a constant block unitary matrix, one block per repeated
phase, and the projectors of \cref{diff:prop:matprojectors} are then Hermitian
orthogonal projectors.
\end{theorem}
\begin{proof}
The identity matrix is horizontal for a skew-Hermitian $A$, so
\cref{diff:thm:matmetric} applied to any normalizing $G$ makes
$G^{*}G=C$ a constant positive definite block Hermitian matrix. Let
$R=C^{-1/2}$, computed in the ordinary finite-dimensional complex block
matrices; then $\der R=0$ and $R$ commutes with $D_{P}$, so $W=GR$ is still
normalizing and $W^{*}W=I$. Uniqueness is \cref{diff:eq:matintertwine} again,
with unitarity making the constant diagonal blocks unitary, and
$WE_{P}W^{*}$ is Hermitian and idempotent.
\end{proof}

The square root in that proof is of a \emph{constant complex} matrix. No
analytic functional calculus on a class-sized space is used, and the theorem
is about differential diagonalization: the term $W^{*}\der W$ has not been
removed and reappears, essentially, in \cref{diff:eq:matspectralproof}.

\section{The spectral--integral theorem for Hermitian coefficients}
\label{diff:sec:matspectral}

\subsection{Finite Hermitian algebra over a real closed field}
\label{diff:sub:matweyl}

The two lemmas of this subsection hold over any real closed field and its
complexification, and use no property of $\der$ at all. Related finite
spectral material for $\SC$ is developed in the companion report
\emph{Surcomplex Spectral Theory} \cite{RepoSpectral}, whose ``spectrum'' is
the algebraic eigenvalue list; see \cref{diff:conv:phases}(e) for why that is
not the object named in \cref{diff:thm:matspectral}.

\begin{lemma}[Finite spectral theorem with attained min--max]
\label{diff:lem:matminmax}
Every Hermitian $H\in\Mat_{n}(\SC)$ has a unitary eigenbasis and eigenvalues
in $\No$. Writing $\lambda_{1}(H)\leq\cdots\leq\lambda_{n}(H)$,
\begin{equation}\label{diff:eq:matminmax}
 \lambda_{k}(H)=\min_{\dim S=k}\ \max_{v\in S,\ v^{*}v=1}v^{*}Hv ,
\end{equation}
and both extrema are attained.
\end{lemma}
\begin{proof}
$\SC$ is algebraically closed, so there is an eigenvector $v\neq0$ with
eigenvalue $\lambda$; since $v^{*}v>0$ and $v^{*}Hv$ is fixed by conjugation,
$\lambda=v^{*}Hv/v^{*}v\in\No$. Normalize $v$ by the positive square root of
$v^{*}v$, available because $\No$ is real closed; the orthogonal complement is
$H$-invariant, and finite induction gives a unitary eigenbasis. For a fixed
$k$-dimensional $S$ the compression of $H$ to an orthonormal basis of $S$ is
again finite Hermitian, so its largest eigenvalue attains the inner maximum;
the span of the first $k$ eigenvectors attains the value $\lambda_{k}(H)$, and
every $k$-dimensional $S$ meets the span of the last $n-k+1$ eigenvectors
nontrivially, where a unit vector has Rayleigh quotient at least
$\lambda_{k}(H)$. No compactness and no completeness is used; the argument is
finite algebra over an ordered field.
\end{proof}

\begin{lemma}[An entrywise Weyl bound, and its ideal form]
\label{diff:lem:matweyl}
For Hermitian $H,E\in\Mat_{n}(\SC)$ put $m=\max_{j,k}\abs{E_{jk}}$. Then
\begin{equation}\label{diff:eq:matweyl}
 \abs{\lambda_{j}(H+E)-\lambda_{j}(H)}\leq nm
 \qquad(1\leq j\leq n).
\end{equation}
If every entry of $E$ lies in $\Acl_{\C}:=\Acl+i\Acl$, then every eigenvalue
difference in \cref{diff:eq:matweyl} lies in $\Acl$.
\end{lemma}
\begin{proof}
For $v^{*}v=1$ the triangle and Cauchy--Schwarz inequalities over the real
closed field $\No$ give
$\abs{v^{*}Ev}\leq m\bigl(\sum_{j}\abs{v_{j}}\bigr)^{2}\leq nm$. Using the
span of the first $j$ eigenvectors of $H$ in \cref{diff:eq:matminmax} gives the
upper bound, and interchanging $H$ and $H+E$ the lower one. If
$E_{jk}=a+ib$ with $a,b\in\Acl$ then $\abs{E_{jk}}\leq\abs a+\abs b\in\Acl$,
so convexity of $\Acl$ (\cref{diff:thm:idealstructure}(iii)) puts
$\abs{E_{jk}}$ in $\Acl$; a finite maximum and the integer multiple $nm$ stay
in $\Acl$, and one more application of convexity finishes the proof.
\end{proof}

\subsection{Bounded gauges have integrable connection terms}
\label{diff:sub:matboundedgauge}

\begin{lemma}[Bounded-gauge estimate]
\label{diff:lem:matboundedgauge}
Let $W\in\Un_{n}(\SC)$. Then every entry of $W$ and of $W^{*}$ is finite,
every entry of $\der W$ lies in $\Acl_{\C}$, every entry of $W^{*}\der W$ lies
in $\Acl_{\C}$, and $W^{*}\der W$ is skew-Hermitian.
\end{lemma}
\begin{proof}
A unitary column has squared moduli summing to $1$, so each entry has modulus
at most $1$ and is finite. Differentiating a finite real or imaginary
coordinate lands in $\der\Ofin=\Acl$ by
\cref{diff:thm:idealstructure}(i). Since $\Acl$ is an $\Ofin$-module
(\cref{diff:thm:idealstructure}(iii)), finite sums of products of entries of
$W^{*}$ with entries of $\der W$ stay in $\Acl_{\C}$. Differentiating
$W^{*}W=I$ gives $(\der W)^{*}W+W^{*}\der W=0$.
\end{proof}

\begin{warning}[The conclusion is about primitives, not about size]
\label{diff:warn:matboundedgauge}
It is the conclusion $W^{*}\der W\in\Mat_{n}(\Acl_{\C})$, and not the weaker
statement that its entries are infinitesimal, that makes
\cref{diff:thm:matspectral} true. The two are different by
\cref{diff:warn:collapse}: an infinitesimal entry may accumulate an infinite
phase, as $\omega^{-1}$ does. This is the precise point at which the strictness
$\Acl\subsetneq\mm$ is load-bearing for the matrix theory.
\end{warning}

\subsection{The theorem}
\label{diff:sub:matmain}

\begin{theorem}[Spectral--integral classification of Hermitian systems]
\label{diff:thm:matspectral}
Let $H=H^{*}\in\Mat_{n}(\SC)$ with algebraic eigenvalues ordered
$\lambda_{1}\leq\cdots\leq\lambda_{n}$ in $\No$, and put
\[
 P_{j}:=\Obs(\lambda_{j})\in\Pcl .
\]
Then $P_{1}\leq\cdots\leq P_{n}$ and
\begin{equation}\label{diff:eq:matspectral}
 \Ph(iH)=\{P_{1},\dots,P_{n}\} .
\end{equation}
There is a unitary $W\in\Un_{n}(\SC)$ with
\begin{equation}\label{diff:eq:matspectralgauge}
 W^{*}(iH)W-W^{*}\der W
   =i\operatorname{diag}(\der P_{1},\dots,\der P_{n}),
\end{equation}
and two Hermitian coefficients give differentially unitarily equivalent
systems $\der y=iH_{1}y$ and $\der y=iH_{2}y$ if and only if their ordered
algebraic eigenvalues have identical images under $\Obs$, counted with
multiplicity. Neither commuting coefficients nor a spectral gap is assumed.
\end{theorem}
\begin{proof}
The matrix $iH$ is skew-Hermitian, so \cref{diff:thm:matunitary} supplies a
unitary normalizing gauge $W$ with phases $Q_{1}\leq\cdots\leq Q_{n}$, ordered
by \cref{diff:cor:obsorder}. Rearranging its defining identity gives
\begin{equation}\label{diff:eq:matspectralproof}
 W^{*}HW=\operatorname{diag}(\der Q_{1},\dots,\der Q_{n})
        +E,\qquad E:=-i\,W^{*}\der W .
\end{equation}
By \cref{diff:lem:matboundedgauge} the matrix $E$ is Hermitian with every
entry in $\Acl_{\C}$. Conjugation by a unitary matrix does not change
eigenvalues, and $\der Q_{1}\leq\cdots\leq\der Q_{n}$ by
\cref{diff:cor:obsorder}; so \cref{diff:lem:matweyl}, applied to the diagonal
matrix and the perturbation $E$, gives
$\lambda_{j}(H)-\der Q_{j}\in\Acl$ for every $j$. Apply $\Obs$: its kernel is
exactly $\Acl$ by \cref{diff:thm:idealstructure}(vi), and
$\Obs(\der Q_{j})=Q_{j}$ by the same item, so
$\Obs(\lambda_{j}(H))=Q_{j}$. This is
\cref{diff:eq:matspectral,diff:eq:matspectralgauge}, and the ordering of the
$P_{j}$ is \cref{diff:cor:obsorder}. The equivalence statement is
\cref{diff:thm:matnormal} together with \cref{diff:thm:matunitary}, which makes
the normalizing gauges of two skew-Hermitian systems unitary, so that the
equivalence between them is unitary as well.
\end{proof}

\begin{warning}[The connection term is not removed]
\label{diff:warn:matconnection}
\Cref{diff:thm:matspectral} does \emph{not} say that pointwise diagonalization
commutes with differentiation; it does not. The term $W^{*}\der W$ in
\cref{diff:eq:matspectralproof} is present, and the theorem says only that it
is invisible in the ordered additive quotient by $\Acl$ described in
\cref{diff:thm:idealstructure}(viii). Nor is the classification the assertion
that every phase grade has a horizontal vector in $\SC^{n}$: by
\cref{diff:cor:matinhom} only the grade $P=0$ does. The formula is exact for
arbitrary surreal scales and arbitrary ordinary finite multiplicities; there is
no hypothesis that the $\lambda_{j}$ be distinct, that gaps dominate a
derivative, or that $H$ commute with its derivative or with any of its
primitives.
\end{warning}

\subsection{Gap-free perturbation, the exact kernel count, and sharpness}
\label{diff:sub:matperturbation}

\begin{corollary}[Finite-primitive perturbation invariance, with no gap hypothesis]
\label{diff:cor:matperturbation}
If $H_{1},H_{2}$ are Hermitian and every entry of $H_{1}-H_{2}$ lies in
$\Acl_{\C}$, then $\der y=iH_{1}y$ and $\der y=iH_{2}y$ are differentially
unitarily equivalent.
\end{corollary}
\begin{proof}
\Cref{diff:lem:matweyl} puts corresponding ordered eigenvalue differences in
$\Acl=\Ker\Obs$; apply \cref{diff:thm:matspectral}.
\end{proof}

\begin{corollary}[Exact kernel count]
\label{diff:cor:matkernel}
For Hermitian $H$,
\begin{equation}\label{diff:eq:matkernelcount}
 \dim_{\C}\{y\in\SC^{n}:\der y=iHy\}
   =\#\{j:\lambda_{j}(H)\in\Acl\},
\end{equation}
counted with algebraic multiplicity. The forced equation $\der y-iHy=f$ is
solvable for every $f\in\SC^{n}$, and uniquely so exactly when no eigenvalue
of $H$ lies in $\Acl$.
\end{corollary}
\begin{proof}
Combine \cref{diff:thm:matspectral,diff:cor:matinhom} with
$\Ker\Obs=\Acl$.
\end{proof}

For $H=\operatorname{diag}(\omega^{-2},\omega^{-1},1)$ the kernel has complex
dimension one, because $\omega^{-2}\in\Acl$ while $\omega^{-1}\notin\Acl$ and
$1\notin\Acl$ (\cref{diff:eq:obsexamples}). Replacing $H$ by any Hermitian
matrix differing from it entrywise by $\Acl_{\C}$ preserves that answer, even
when the perturbation is far larger than a gap at some more remote
infinitesimal scale --- which is exactly what a gap hypothesis would forbid.

\begin{proposition}[Sharpness of the perturbation class]
\label{diff:prop:sharpperturb}
For real $h,k\in\No$ the equations $\der y=ihy$ and $\der y=iky$ are
differentially unitarily equivalent exactly when $h-k\in\Acl$. Hence no
strictly larger additive class of scalar perturbations preserves the phase of
every Hermitian system.
\end{proposition}
\begin{proof}
A scalar unitary gauge equivalence is $i(h-k)=u^{\dagger}$ with $u\bar u=1$;
\cref{diff:thm:logimage} together with \cref{diff:thm:unitphase} gives
necessity and constructs a unit solution for sufficiency. If $b\notin\Acl$ the
one-dimensional systems with coefficients $0$ and $b$ have obstructions $0$ and
$\Obs(b)\neq0$.
\end{proof}

``Infinitesimal perturbation'' therefore cannot replace ``finite-primitive
perturbation'': the scalar change $0\mapsto\omega^{-1}$ already moves the phase
from $0$ to $\log\omega$. This is \cref{diff:warn:collapse} again, now in its
sharpest form: $\Acl$ is not merely a convenient subclass of $\mm$, it is the
exact class for which \cref{diff:cor:matperturbation} is true.

\subsection{The two boundaries, exhibited}
\label{diff:sub:matboundaries}

\begin{example}[A nonnormal system whose eigenvalues predict the wrong phases]
\label{diff:ex:matnonnormal}
Let $x=\omega$ and
\[
 G=\begin{pmatrix}1-x^{2}&x\\-x&1\end{pmatrix},
 \qquad
 A=\begin{pmatrix}-x&1+x^{2}\\-1&x\end{pmatrix}.
\]
Then $\det G=1$ and $\der G=AG$, so $A$ has a fundamental matrix in
$\GL_{2}(\SC)$ and, by \cref{diff:cor:matinhom},
$\Ph(A)=\{0,0\}$. Yet
\[
 \det(\zeta I-A)=\zeta^{2}+1 ,
\]
so the instantaneous algebraic eigenvalues of $A$ are $i$ and $-i$, and taking
obstructions of their imaginary parts would predict $\{\omega,-\omega\}$. The
prediction is wrong by \cref{diff:thm:matnormal}.
\end{example}

\begin{warning}[What the counterexample rules out, and what it marks]
\label{diff:warn:matnonnormal}
\Cref{diff:ex:matnonnormal} is not a curiosity; it is the boundary of
\cref{diff:thm:matspectral}. It shows that the Hermitian hypothesis cannot be
dropped and that the main formula does not extend to arbitrary matrices by
extracting imaginary parts of eigenvalues. The mechanism is exactly the one
\cref{diff:lem:matboundedgauge} controls in the Hermitian case: the similarity
gauges needed for a general nonnormal matrix need not be entrywise finite, so
their connection term need not lie in $\Mat_{n}(\Acl_{\C})$ and can carry a
nonzero phase of its own. Here $G$ has the infinite entry $1-\omega^{2}$.

The complete classification of \cref{diff:thm:matnormal} still applies to
this non-Hermitian system: its phase multiset is $\{0,0\}$. What fails is
computing that multiset from the instantaneous eigenvalues by the Hermitian
formula. \Cref{diff:thm:matspectral} applies $\Obs$ to the algebraic
eigenvalues of $H$ when $A=iH$; the general normal form instead uses
differential factorization and gauges. Effective construction of those data
on specified input languages remains a separate question
(\cref{diff:sub:nc:open}).

For the power-Hahn regular-singular systems of \cref{diff:rs:part},
$\Ph(-tA)=\{(\Imm\lambda_j)\log t\}$, with $\lambda_j$ the eigenvalues
of the ordinary residual matrix (\cref{diff:rs:prop:phasespectrum}). This
provides another explicitly computable non-Hermitian class;
\cref{diff:ex:matnonnormal} lies outside it.
\end{warning}

\begin{example}[A genuinely coupled unitary system with two independent phases]
\label{diff:ex:matcoupled}
Let $x=\omega$ and
\[
 a=\frac{1-x^{2}}{1+x^{2}},
 \qquad b=\frac{2x}{1+x^{2}},
 \qquad
 W=\begin{pmatrix}a&-b\\b&a\end{pmatrix},
 \qquad
 D=i\operatorname{diag}(x,x^{-1}),
\]
so that $W$ is a real orthogonal matrix over $\No$ and
$A:=(\der W)W^{\mathsf T}+WDW^{\mathsf T}$ satisfies $A^{*}=-A$. Explicitly
\[
 A=\begin{pmatrix}
 i(a^{2}x+b^{2}x^{-1}) &
 -\dfrac{2}{1+x^{2}}+iab(x-x^{-1})\\[6pt]
 \dfrac{2}{1+x^{2}}+iab(x-x^{-1}) &
 i(b^{2}x+a^{2}x^{-1})
 \end{pmatrix},
\]
which is neither diagonal nor a constant-matrix system. Its exact normal form
is $D$, so
\[
 \Ph(A)=\Bigl\{\tfrac12\omega^{2},\ \log\omega\Bigr\},
\]
using $\Obs(\omega)=\tfrac12\omega^{2}$ and $\Obs(\omega^{-1})=\log\omega$
(\cref{diff:eq:obsexamples}). Writing $H=-iA$, one computes
\[
 \tr H=x+x^{-1},
 \qquad
 \det H=1-\frac{4}{(1+x^{2})^{2}},
\]
so the instantaneous eigenvalues of $H$ are \emph{not} $x$ and $x^{-1}$. Their
discrepancies lie in $\Acl$, so \cref{diff:thm:matspectral} recovers the same
two phases from them. Both phases are nonzero, so by
\cref{diff:cor:matinhom} every forced system $\der y-Ay=f$ has a unique
solution in $\SC^{2}$; and the two phases are $\Q$-linearly independent, which
\cref{diff:sub:mattorus} uses.
\end{example}

\section{Consequences of the matrix normal form}
\label{diff:sec:matconsequences}

\subsection{Normalized non-Abelian integration}
\label{diff:sub:matnonabelian}

Write
\[
 \mathfrak u_{n}(\Acl_{\C})
  =\{A\in\Mat_{n}(\Acl_{\C}):A^{*}=-A\},
\]
and let $\Un_{n}(1+\mm_{\C})$ be the unitary matrices $W$ with
$W-I\in\Mat_{n}(\mm_{\C})$. The latter is a congruence subgroup, not the
requirement that off-diagonal entries be near $1$.

\begin{theorem}[Normalized compact integration]
\label{diff:thm:matnonabelian}
The logarithmic-derivative map
\begin{equation}\label{diff:eq:matnonabelian}
 \Un_{n}(1+\mm_{\C})\longrightarrow\mathfrak u_{n}(\Acl_{\C}),
 \qquad W\longmapsto(\der W)W^{*},
\end{equation}
is a bijection of classes. Equivalently, for skew-Hermitian $A$ the following
are equivalent:
\begin{enumerate}[label=\textup{(\roman*)},leftmargin=2.4em]
\item every entry of $A$ has finite real and imaginary primitives, that is
$A\in\Mat_{n}(\Acl_{\C})$;
\item $\der y=Ay$ has a fundamental matrix in $\GL_{n}(\SC)$;
\item $\der W=AW$ has a unitary solution with $\st(W)=I$ entrywise.
\end{enumerate}
The solution in \textup{(iii)} is unique.
\end{theorem}
\begin{proof}
\Cref{diff:lem:matboundedgauge} makes \cref{diff:eq:matnonabelian} well defined
and gives (iii)$\Rightarrow$(i). Assume (i) and put $H=-iA$, Hermitian with
entries in $\Acl_{\C}$. Comparing $H$ with the zero matrix,
\cref{diff:lem:matweyl} places every eigenvalue of $H$ in $\Acl$, so
\cref{diff:thm:matspectral} gives $\Ph(A)=\{0,\dots,0\}$ and a unitary $V$
with $A^{V}=0$, that is $\der V=AV$. Taking standard parts entrywise in
$V^{*}V=I$ --- legitimate because every entry is finite by
\cref{diff:lem:matboundedgauge} --- shows that $C=\st(V)$ is an ordinary
complex unitary matrix, and $W=VC^{-1}$ has $\der W=AW$ and $\st(W)=I$.

If $W_{1},W_{2}$ are two normalized solutions then $\der(W_{1}^{-1}W_{2})=0$,
so $W_{1}^{-1}W_{2}$ is a constant complex matrix with standard part $I$,
hence $I$. That is uniqueness and injectivity of
\cref{diff:eq:matnonabelian}. Clearly (iii)$\Rightarrow$(ii); conversely a
fundamental matrix makes every phase zero by \cref{diff:cor:matinhom}, and for
a skew-Hermitian system \cref{diff:thm:matunitary} then supplies a unitary
fundamental matrix, which after multiplication by the inverse of its standard
part is normalized.
\end{proof}

\begin{warning}[This is an existence theorem, not an exponential formula]
\label{diff:warn:matnonabelian}
\Cref{diff:thm:matnonabelian} is not the assertion that
$W=\exp\bigl(\int A\bigr)$, and there is no ordering prescription, no chosen
contour, and no claimed convergence of Picard iteration. Its normalization is
residue data, $\st(W)=I$, not an initial value at a real or surreal time ---
the same distinction \cref{diff:warn:Jnotalgorithm} makes for $\Izero$. For the
unique solution $W_{A}$ the ordinary product rule gives
\begin{equation}\label{diff:eq:matcocycle}
 \bigl(\der(W_{A}W_{B})\bigr)(W_{A}W_{B})^{*}
 =A+W_{A}BW_{A}^{*},
\end{equation}
whose right-hand side again lies in $\mathfrak u_{n}(\Acl_{\C})$. So the
bijection transports the group law to $A\star B=A+W_{A}BW_{A}^{*}$, which is
equal to $A+B$ exactly when $W_{A}B=BW_{A}$. Commutation of $A$ and $B$
alone is not sufficient, as the following example also shows.
\end{warning}

\begin{example}[An exact noncommuting Hahn solution]
\label{diff:ex:matnoncommuting}
Let $t=\omega^{-1}$, so $\der t=-t^{2}$, and take the constant skew-Hermitian
matrices
\[
 X=\begin{pmatrix}0&i\\i&0\end{pmatrix},
 \qquad
 Z=\begin{pmatrix}i&0\\0&-i\end{pmatrix} .
\]
For $A=-t^{2}X-2t^{3}Z$ --- every entry of which lies in $\Acl_{\C}$ --- the
normalized solution of \cref{diff:thm:matnonabelian} is the Hahn series
$W=\sum_{n\geq0}W_{n}t^{n}$ determined by
\begin{equation}\label{diff:eq:matnoncommrec}
 W_{0}=I,
 \qquad
 (n+1)W_{n+1}=XW_{n}+2ZW_{n-1},
 \qquad W_{-1}=0 ,
\end{equation}
whose first terms are
\[
 W=I+tX+t^{2}\Bigl(\tfrac12X^{2}+Z\Bigr)
   +t^{3}\Bigl(\tfrac16X^{3}+\tfrac13XZ+\tfrac23ZX\Bigr)+\cdots .
\]
This differs from $\exp(tX+t^{2}Z)$ at order $t^{3}$ by
$(ZX-XZ)/6\neq0$. The support is contained in $\N$, so the series is strongly
summable by \cref{diff:lem:neumann} and the recurrence proves the equation
exactly; $(\der(W^{*}W))=0$ with standard part $I$ makes it unitary, so it is
the unique solution. This is a computation in a restricted Hahn workspace and
is not a claim that such a recurrence is effectively available for a general
coefficient matrix.

For the product law take $B=A$, so $[A,B]=0$ automatically. The displayed
coefficients of $W=W_A$ give
\[
 [t^{4}]\bigl(W_A A W_A^{*}-A\bigr)=-[X,Z]\neq0.
\]
Indeed the two contributions at this order are
$[tX,-2t^{3}Z]=-2t^{4}[X,Z]$ and
$[t^{2}Z,-t^{2}X]=t^{4}[X,Z]$; the terms involving only $X$ cancel.
Thus $A\star A\neq2A$, even though the coefficient matrices $A$ and $B$
commute.
\end{example}

\Cref{diff:ex:matnoncommuting} is the same phenomenon as the two-scale
computation of \cref{diff:sec:twoscale}, which refutes the naive exponential
for the \emph{coordinate} derivative $\Dz$; the two live in different
categories (\cref{diff:tab:categories}) and neither is evidence for the other.
What they share is the mechanism: a noncommuting pair of constant matrices
produces an exact third-order defect that no exponential of an integral can
absorb.

\subsection{Real systems: rotation blocks and metric signatures}
\label{diff:sub:matreal}

For a real system, conjugation sends phase $P$ to phase $-P$, so the two
multiplicities are equal.

\begin{theorem}[Real rotation-block classification]
\label{diff:thm:matreal}
For every $A\in\Mat_{n}(\No)$ there is $G\in\GL_{n}(\No)$ with
\begin{equation}\label{diff:eq:matrealnormal}
 A^{G}=0_{m_{0}}\ \oplus\bigoplus_{P>0}\bigl((\der P)\Rot\bigr)^{\oplus m_{P}},
 \qquad
 n=m_{0}+2\sum_{P>0}m_{P},
\end{equation}
where $\Rot=\left(\begin{smallmatrix}0&1\\-1&0\end{smallmatrix}\right)$. The
positive phases $P\in\Pcl$ and their multiplicities are unique. Every finite
real differential module over $\No$ is a direct sum of trivial
one-dimensional modules and nontrivial two-dimensional rotation modules, and
if $A^{\mathsf T}=-A$ then $G$ can be chosen orthogonal.
\end{theorem}
\begin{proof}
Uniqueness of the complex phase subspaces
(\cref{diff:prop:matprojectors}) makes conjugation interchange the $P$ and
$-P$ subspaces. On the zero-phase subspace the horizontal complex space is
conjugation stable, so a real basis of its fixed part is a complex basis after
complexification. For $P>0$ take normalizing columns
$v_{1},\dots,v_{m_{P}}$ of phase $P$; their conjugates span phase $-P$. Write
$v_{j}=a_{j}+ib_{j}$. From $\der v_{j}=Av_{j}-i(\der P)v_{j}$,
\[
 \der a_{j}=Aa_{j}+(\der P)b_{j},
 \qquad
 \der b_{j}=Ab_{j}-(\der P)a_{j},
\]
which is the block $(\der P)\Rot$. The real columns are independent because
the change between $(v_{j},\bar v_{j})$ and $(a_{j},b_{j})$ is invertible over
$\C$. The only invariant real subspaces inside one nonzero phase pair are $0$
and the whole pair, since a real subspace complexifies to a
conjugation-stable sum of two distinct phase lines. For skew-symmetric $A$ the
identity metric is horizontal, so $G^{\mathsf T}G$ is a positive constant matrix
commuting with the rotation blocks, by \cref{diff:thm:matrealmetric} below;
its positive constant square root commutes with them too, and multiplying $G$
by that inverse square root makes it orthogonal without changing the normal
form.
\end{proof}

\begin{theorem}[The real metric cone and its parity]
\label{diff:thm:matrealmetric}
Every real finite system has a positive symmetric horizontal metric. With the
multiplicities of \cref{diff:eq:matrealnormal}, the cone of positive symmetric
horizontal metrics is isomorphic to
\begin{equation}\label{diff:eq:matrealcone}
 \Symm^{+}_{m_{0}}(\R)\times\prod_{P>0}\Herm^{+}_{m_{P}}(\C),
\end{equation}
a product over a finite set, empty products being a point. The real vector
space of all horizontal symmetric forms has dimension
$\tfrac12 m_{0}(m_{0}+1)+\sum_{P>0}m_{P}^{2}$, and on each nonzero real phase
block both indices of a nondegenerate horizontal symmetric form are even.
\end{theorem}
\begin{proof}
Pull a real horizontal form back by the gauge of \cref{diff:thm:matreal}.
Complex phase separation (\cref{diff:thm:matmetric}) eliminates all terms
between different positive phase pairs and between such a pair and phase zero,
and leaves an arbitrary constant real symmetric matrix on the zero block.
Group a repeated positive phase as $(\der P)\Rot_{m}$ with
$\Rot_{m}=\left(\begin{smallmatrix}0&I_{m}\\-I_{m}&0\end{smallmatrix}\right)$;
the real symmetric horizontal matrices are then the constant matrices
commuting with $\Rot_{m}$, that is
\[
 C=\begin{pmatrix}S&T\\-T&S\end{pmatrix},
 \qquad S^{\mathsf T}=S,\quad T^{\mathsf T}=-T ,
\]
which is the realification of the complex Hermitian matrix $S-iT$. A constant
complex unitary diagonalization of $S-iT$ becomes a real orthogonal
diagonalization of $C$ with every eigenvalue repeated twice; this gives
positivity, the cone formula and the parity statement. Congruence by $G^{-1}$
returns to the original coordinates.
\end{proof}

There is no integration over phases in \cref{diff:eq:matrealcone}: the
products are over a finite multiset of constant-matrix cones. In particular a
real irreducible two-dimensional system has a one-dimensional space of
horizontal symmetric forms, hence a single positive ray. That fact is the
source of the uniqueness in \cref{diff:thm:matpinney}.

\subsection{A positive Ermakov--Pinney dichotomy}
\label{diff:sub:matpinney}

For $q\in\No$ consider the intrinsic equations
\begin{equation}\label{diff:eq:matlinearq}
 \der^{2}y+qy=0
\end{equation}
and
\begin{equation}\label{diff:eq:matpinney}
 \der^{2}\rho+q\rho=\rho^{-3},\qquad \rho>0 .
\end{equation}
Both derivatives are the fixed $\der$. The positive sign on the right of
\cref{diff:eq:matpinney} is part of the statement; the sign convention in the
classical literature varies \cite{Pinney}.

\begin{theorem}[Positive Pinney existence and uniqueness]
\label{diff:thm:matpinney}
For every $q\in\No$, equation \cref{diff:eq:matpinney} has a positive solution
in $\No$, and exactly one of the following holds.
\begin{enumerate}[label=\textup{(\roman*)},leftmargin=2.4em]
\item \Cref{diff:eq:matlinearq} has no nonzero solution in $\No$. Then
\cref{diff:eq:matpinney} has \emph{exactly one} positive solution in $\No$,
and $\der^{2}y+qy=f$ has a unique solution in $\No$ for every $f\in\No$.
\item \Cref{diff:eq:matlinearq} has a real fundamental pair in $\No$. Then the
positive solutions of \cref{diff:eq:matpinney} form the classical positive
quadratic-form family \cref{diff:eq:matpinneyclassical}, with two real
parameters after determinant normalization.
\end{enumerate}
\end{theorem}
\begin{proof}
Let $A_{q}=\left(\begin{smallmatrix}0&1\\-q&0\end{smallmatrix}\right)$. By
\cref{diff:thm:matrealmetric} it has a positive symmetric horizontal metric
$\mathcal M$, and \cref{diff:eq:horizontal} gives
$\der\det\mathcal M=-2(\tr A_{q})\det\mathcal M=0$, so $\det\mathcal M$ is a
positive real constant; scale it to $1$. Writing
$\mathcal M=\left(\begin{smallmatrix}a&b\\b&c\end{smallmatrix}\right)$, the
horizontality equations are
\begin{equation}\label{diff:eq:matmetriccomponents}
 \der a=2qb,\qquad \der b=qc-a,\qquad \der c=-2b .
\end{equation}
Since $c>0$ and $\No$ is real closed, put $\rho=\sqrt c$. The last equation
gives $b=-\rho\,\der\rho$ and $\det\mathcal M=1$ gives
$a=(\der\rho)^{2}+\rho^{-2}$; the middle equation becomes exactly
\cref{diff:eq:matpinney}. Conversely a positive solution of
\cref{diff:eq:matpinney} produces the positive determinant-one horizontal
metric
\begin{equation}\label{diff:eq:matpinneymetric}
 \mathcal M_{\rho}=
 \begin{pmatrix}
 (\der\rho)^{2}+\rho^{-2} & -\rho\,\der\rho\\
 -\rho\,\der\rho & \rho^{2}
 \end{pmatrix},
\end{equation}
positive because its lower diagonal entry is positive and its determinant is
$1$. So positive Pinney solutions correspond bijectively to positive
determinant-one horizontal metrics.

In dimension two the real normal form \cref{diff:eq:matrealnormal} has either
two zero phases or one nonzero conjugate pair. If \cref{diff:eq:matlinearq}
has no nonzero solution in $\No$ it must be the second case, the positive
metric cone is a single ray by \cref{diff:eq:matrealcone}, and the
determinant-one condition selects one point. In the zero-phase case a real
fundamental matrix exists and the cone is the three-dimensional positive
symmetric constant-matrix cone, whose determinant-one slice is the asserted
two-parameter family. In the nonzero-phase case \cref{diff:cor:matinhom} gives
a unique solution in $\SC^{2}$ of every forced system, and conjugating shows
it is real when $f$ is.

There is no intermediate one-dimensional case: if $u\in\No^{\times}$ solves
\cref{diff:eq:matlinearq}, choose $B$ with $\der B=u^{-2}$ by
\ref{diff:BM:surj} and put $v=uB$; then $u\der v-v\der u=1$, so $u,v$ are an
independent real fundamental pair.
\end{proof}

With $u,v$ a real fundamental pair of constant Wronskian
$\mathrm{Wr}=u\der v-v\der u\neq0$, alternative (ii) reads
\begin{equation}\label{diff:eq:matpinneyclassical}
 \rho=\sqrt{\alpha u^{2}+2\beta uv+\gamma v^{2}},
 \qquad
 \alpha,\gamma>0,\quad \alpha\gamma-\beta^{2}=\mathrm{Wr}^{-2},
 \quad \alpha,\beta,\gamma\in\R .
\end{equation}
This is the classical Pinney parametrization \cite{Pinney}, here read in an
ordered differential field; no priority is claimed for it, nor for the
abstract Schwarz-closed-field existence principle that underlies it
\cite{ADHNorm}. What is new to this field is that alternative (i) can occur,
and that it collapses the positive family to a single point. For recent
existence and sharp uniqueness results for amplitude--phase representations
over Hardy fields see \cite{ADHSecond}; no solution of the named Boshernitzan
conjecture discussed there is claimed here.

For $q=1$ alternative (i) holds by \cref{diff:cor:oscillator}, and the unique
positive Pinney solution is $\rho=1$. For $q=-1$ the linear equation has the
fundamental pair $e^{\omega},e^{-\omega}$ and alternative (ii) gives many
positive amplitudes. A unique positive solution of a second-order equation is
not a contradiction to ordinary initial-value theory: intrinsic differentiation
has no two-real-parameter initial-data interpretation
(\cref{diff:warn:separators}). (Added with member~13: for the supercritical
critical potentials of \cref{diff:cp:part} the unique positive solution and the
canonical phase are explicit; see \cref{diff:cp:rem:pinney}.)

\begin{theorem}[Canonical amplitude--phase reduction]
\label{diff:thm:matamplitude}
Call $q\in\No$ \emph{intrinsically oscillatory} when alternative (i) of
\cref{diff:thm:matpinney} holds; this means only that the nonzero homogeneous
oscillations are missing from $\No$, and says nothing about oscillation of a
function of a real variable. Let $\rho$ be the unique positive Pinney
solution, $h=\rho^{-2}$, and $B=\Izero(h)=P+\eta$ with $P\in\Pcl$ and
$\eta\in\mm$ (\cref{diff:eq:split}). Then $P>0$, and with
\[
 G_{\rho}=\begin{pmatrix}\rho&0\\ \der\rho&\rho^{-1}\end{pmatrix},
 \qquad
 R(\eta)=\Ree\bigl(\Emm(i\eta)\bigr)\,I+\Imm\bigl(\Emm(i\eta)\bigr)\,\Rot ,
\]
the matrix $S_{q}=G_{\rho}R(\eta)\in\SL_{2}(\No)$ satisfies
\begin{equation}\label{diff:eq:matcanonicalgauge}
 A_{q}^{S_{q}}=(\der P)\,\Rot .
\end{equation}
The amplitude $\rho$, the normalized primitive $B$, the phase $P$ and this
amplitude-times-infinitesimal-rotation gauge are canonical; no uniqueness
among all symplectic gauges is asserted.
\end{theorem}
\begin{proof}
A direct multiplication using \cref{diff:eq:matpinney} gives
$A_{q}^{G_{\rho}}=\rho^{-2}\Rot=h\Rot$. The rotation $R(\eta)$ is a
well-defined real matrix with entries the infinitesimal cosine and sine of
\cref{diff:thm:unitphase}; it has determinant one, commutes with $\Rot$, and
satisfies $R(\eta)^{-1}\der R(\eta)=(\der\eta)\Rot$. Subtracting gives
\cref{diff:eq:matcanonicalgauge}. If $P=0$ then $B=\eta$ is finite and the same
calculation reduces the system to zero, producing a real fundamental matrix
and contradicting alternative (i). Since $\der B=h>0$, $B$ cannot be negative
infinite by \ref{diff:BM:positivity}, so its nonzero purely infinite part is
positive. The uniqueness assertions are \cref{diff:thm:matpinney} and the
normalization $\ct(B)=0$ of \cref{diff:prop:normprimitive}.
\end{proof}

The conserved quadratic expression attached to $\rho$ is
\begin{equation}\label{diff:eq:matlewis}
 \mathcal E_{\rho}(y,\der y)
  =\bigl(\rho\,\der y-y\,\der\rho\bigr)^{2}+\bigl(y/\rho\bigr)^{2},
\end{equation}
whose derivative vanishes formally along \cref{diff:eq:matlinearq} by
\cref{diff:eq:matpinneymetric}. It is meaningful in an oscillatory extension
of the kind constructed in \cref{diff:sec:PV} even when the only solution in
$\No$ is zero. It is an exact invariant, not an averaged energy.

\subsection{The oscillatory Airy equation as an exact Hahn example}
\label{diff:sub:matairy}

Let $x=\omega$ and consider
\begin{equation}\label{diff:eq:matairy}
 \der^{2}y+xy=0 .
\end{equation}
Put $P=\tfrac23x^{3/2}$, so $\der P=x^{1/2}$ by the real-power rule
(\cref{diff:warn:monomialmap}, applied at the ordinary real exponent
$\tfrac32$). Let $\mathsf T_{P}$ denote a formal oscillatory symbol with
$\der\mathsf T_{P}=i(\der P)\mathsf T_{P}$ in the sense of
\cref{diff:thm:PVone}; it is a new differential element and never a surcomplex
number. Seek a solution $a\mathsf T_{P}$ with
\begin{equation}\label{diff:eq:matairyamp}
 a=x^{-1/4}\sum_{n\geq0}c_{n}x^{-3n/2},\qquad c_{0}=1 .
\end{equation}
Substituting into \cref{diff:eq:matairy} gives
\begin{equation}\label{diff:eq:matairyampeq}
 \der^{2}a+2ix^{1/2}\der a+\frac{i}{2x^{1/2}}\,a=0 ,
\end{equation}
and with $\alpha_{n}=-\tfrac14-\tfrac32n$ coefficient comparison yields
\begin{equation}\label{diff:eq:matairyrec}
 c_{n+1}=\frac{\alpha_{n}(\alpha_{n}-1)}{3i(n+1)}c_{n}
        =-i\,\frac{(6n+1)(6n+5)}{48(n+1)}\,c_{n},
 \qquad
 c_{n}=(-i)^{n}\Bigl(\tfrac34\Bigr)^{n}
       \frac{(1/6)_{n}(5/6)_{n}}{n!} .
\end{equation}
The exponents in \cref{diff:eq:matairyamp} decrease strictly to $-\infty$, so
the displayed expression is a legitimate surcomplex Hahn sum however fast the
ordinary coefficients grow, and its termwise differentiation is justified by
\ref{diff:BM:additive}. This is the same phenomenon as the factorially
divergent solutions of \cref{diff:thm:factorial}, and it carries the same
disclaimer (\cref{diff:warn:antitopological}): a strictly descending support is
a summability certificate, not convergence in the fine topology on $\SC$.

Put $b=\der a+ix^{1/2}a$ and $G=\left(\begin{smallmatrix}a&\bar a\\
b&\bar b\end{smallmatrix}\right)$. The two columns satisfy
$\der v=A_{x}v\mp i(\der P)v$ with
$A_{x}=\left(\begin{smallmatrix}0&1\\-x&0\end{smallmatrix}\right)$, so
$\der G=A_{x}G-G\operatorname{diag}(i\der P,-i\der P)$. The determinant
$\mathrm{Wr}=a\bar b-\bar ab$ has zero derivative, hence lies in $\C$, and its
leading term is $-2i$ because $a\sim x^{-1/4}$ and $b\sim ix^{1/4}$; so
$\mathrm{Wr}=-2i$ exactly, $G$ is invertible, and
\begin{equation}\label{diff:eq:matairyphases}
 \Ph(A_{x})=\Bigl\{\tfrac23\omega^{3/2},\ -\tfrac23\omega^{3/2}\Bigr\}.
\end{equation}
By \cref{diff:cor:matinhom} equation \cref{diff:eq:matairy} has no nonzero
solution in $\SC$, and $\der^{2}y+\omega y=f$ has a unique solution in $\SC$
for every $f$. As always in this report that is a statement about surcomplex
\emph{number} solutions under $\der$; it does not deny the ordinary Airy
functions of a complex coordinate (\cref{diff:warn:whatnonexistence}).

The unique positive solution of $\der^{2}\rho+\omega\rho=\rho^{-3}$ is
\begin{equation}\label{diff:eq:matairyrho}
 \rho=\sqrt{a\bar a}.
\end{equation}
Indeed, with $v=\binom ab$ the real matrix $C=\Ree(vv^{*})$ satisfies
$\der C=A_{x}C+CA_{x}^{\mathsf T}$ --- the scalar imaginary phase cancels ---
and $\det C=\bigl(\Imm(a\bar b)\bigr)^{2}=1$ since $\mathrm{Wr}=-2i$; so
$\mathcal M=C^{-1}$ is a positive determinant-one horizontal metric and the
bijection of \cref{diff:thm:matpinney} gives
\cref{diff:eq:matairyrho}, with uniqueness because the phases
\cref{diff:eq:matairyphases} are nonzero. Its expansion begins
\begin{equation}\label{diff:eq:matairyjet}
 \rho=\omega^{-1/4}\left(
 1-\frac{5}{64}\omega^{-3}
 +\frac{2285}{8192}\omega^{-6}
 -\frac{1690275}{524288}\omega^{-9}
 +\frac{10371973875}{134217728}\omega^{-12}+\cdots\right),
\end{equation}
these being coefficients of an exact Hahn series and not estimates for a
convergent ordinary power series. The normalized primitive of $\rho^{-2}$ is
$B=\tfrac23\omega^{3/2}-\tfrac{5}{48}\omega^{-3/2}+\cdots$, whose purely
infinite part is exactly the positive phase in
\cref{diff:eq:matairyphases}; the remaining correction is infinitesimal and is
absorbed by the rotation $R(\eta)$ of \cref{diff:thm:matamplitude}.

For a truncation $a_{N}=\sum_{n\leq N}c_{n}x^{\alpha_{n}}$ the recurrence
cancels everything but the last second derivative, giving the exact residual
\begin{equation}\label{diff:eq:matairyres}
 \der^{2}a_{N}+2ix^{1/2}\der a_{N}+\frac{i}{2x^{1/2}}a_{N}
  =\alpha_{N}(\alpha_{N}-1)c_{N}x^{\alpha_{N}-2},
\end{equation}
which the delivered suite checks at $N=16$
(\cref{diff:app:verification}). As \cref{diff:warn:truncationtype} requires,
the residual records the \emph{type} of the omitted information: it is a formal
remainder at a stated order, and the infinite-support statement is proved by
the recurrence together with strict descent of exponents, not by the finite
check.

\subsection{First integrals from phase resonance}
\label{diff:sub:matresonance}

Let $y_{1},\dots,y_{n}$ be algebraically independent formal variables and
extend $\der$ to $\SC[y_{1},\dots,y_{n}]$ by
\[
 \mathcal D_{A}f=\der_{\mathrm{coeff}}f+\sum_{j}(Ay)_{j}\frac{\partial f}{\partial y_{j}},
\]
where $\der_{\mathrm{coeff}}$ differentiates coefficients only. A polynomial or
rational \emph{first integral} is an element killed by $\mathcal D_{A}$. This
definition presupposes no nonzero solution in $\SC^{n}$. In phase coordinates
$z=G^{-1}y$ one has $\mathcal D_{A}z_{j}=i(\der P_{j})z_{j}$.

Put
\[
 \Pi=\sum_{j=1}^{n}\Z P_{j}\subseteq\Pcl,
 \qquad
 \Lambda=\Bigl\{m\in\Z^{n}:\textstyle\sum_{j}m_{j}P_{j}=0\Bigr\},
 \qquad
 \Lambda^{+}=\Lambda\cap\N^{n},
\]
and $r=\dim_{\Q}\Span_{\Q}\{P_{1},\dots,P_{n}\}$. This $\Lambda$ is the
relation lattice of \cref{diff:eq:lattice} for the diagonal system
$\Ph(A)$; it is saturated by \cref{diff:lem:saturated}, so
$\operatorname{rank}\Lambda=n-r$ and $\Pi$ is free abelian of rank $r$.

\begin{theorem}[The resonance algebra]
\label{diff:thm:matresonance}
In phase coordinates,
\begin{equation}\label{diff:eq:matinvariantring}
 \Ker\mathcal D_{A}\big|_{\SC[z]}=\C[z^{m}:m\in\Lambda^{+}],
 \qquad
 \Ker\mathcal D_{A}\big|_{\SC(z)}=\C(z^{m}:m\in\Lambda).
\end{equation}
The polynomial algebra is finitely generated, and the rational field is purely
transcendental over $\C$ of transcendence degree $n-r$.
\end{theorem}
\begin{proof}
For $f=\sum_{\alpha}c_{\alpha}z^{\alpha}$ the derivative vanishes exactly when
$\der c_{\alpha}+i\bigl(\der\sum_{j}\alpha_{j}P_{j}\bigr)c_{\alpha}=0$ for each
$\alpha$, and \cref{diff:lem:matseparation} makes a nonzero coefficient
require $\alpha\in\Lambda^{+}$ and $c_{\alpha}\in\C$.

For the rational statement work in the Laurent ring and write $f/g$ with
coprime nonzero Laurent polynomials. A zero derivative gives
$f\mid\mathcal D_{A}f$ and $g\mid\mathcal D_{A}g$; since
$\mathcal D_{A}$ enlarges no Laurent support and the highest and lowest degrees
in each variable are additive under multiplication, one gets
$\mathcal D_{A}f=hf$ and $\mathcal D_{A}g=hg$ for a common $h\in\SC$, the case
$h=0$ included. Comparing two coefficients of $f$ shows their ratio has
logarithmic derivative $i\der(\beta\cdot P-\alpha\cdot P)$, so by
\cref{diff:lem:matseparation} their phases agree and the ratio is constant;
thus all monomials of $f$ share one phase, and after factoring out one
coefficient and one monomial the rest has constant coefficients and exponents
in $\Lambda$. The same holds for $g$, and equality of the two values of $h$
forces the factored phases to agree and the factored coefficient ratio to be
constant. Hence $f/g$ is a rational function of monomials with exponents in
$\Lambda$; the reverse inclusion is immediate.

For finite generation, call $0\neq u\in\Lambda^{+}$ indecomposable when it is
not a sum of two nonzero members of $\Lambda^{+}$. Two distinct
indecomposables are incomparable in the coordinatewise order on $\N^{n}$,
since $u<v$ would give $v-u\in\Lambda^{+}\setminus\{0\}$. To see that
such an antichain is finite, take any infinite sequence of distinct members
of $\N^{n}$. Its first coordinates admit an infinite nondecreasing
subsequence: either one value occurs infinitely often, or every tail is
unbounded and an increasing subsequence can be chosen recursively. Repeat
this selection for each of the remaining finitely many coordinates. The
result contains two coordinatewise comparable members, a contradiction. This
is Dickson's finite-antichain lemma. Hence there are finitely many
indecomposables, and induction on total degree writes every member of
$\Lambda^{+}$ as a sum of them. Finally a
$\Z$-basis of $\Lambda$ gives $n-r$ Laurent monomials generating the rational
field, algebraically independent because distinct integer combinations of
basis exponents give distinct monomials. No analytic finiteness principle is
used.
\end{proof}

\begin{example}[Complete reducibility does not make the invariant algebra a polynomial ring]
\label{diff:ex:matthreephase}
Let $x=\omega$ and
\[
 G=\begin{pmatrix}1&x^{-1}&0\\0&1&x^{-2}\\0&0&1\end{pmatrix},
 \qquad
 D=\operatorname{diag}(i,i,-2i),
 \qquad
 A=(\der G)G^{-1}+GDG^{-1},
\]
so that
\begin{equation}\label{diff:eq:matthreematrix}
 A=\begin{pmatrix}
 i & -x^{-2} & x^{-4}\\
 0 & i & (-3ix-2)x^{-3}\\
 0 & 0 & -2i
 \end{pmatrix}
\end{equation}
and $\Ph(A)=\{\omega,\omega,-2\omega\}$, so by \cref{diff:cor:matinhom} the
system has no nonzero solution in $\SC^{3}$. Its formal first integrals are
nevertheless nontrivial: with $z=G^{-1}y$ the resonance condition is
$\alpha_{1}+\alpha_{2}=2\alpha_{3}$ and
\begin{equation}\label{diff:eq:mattoric}
 \Ker\mathcal D_{A}\big|_{\SC[z]}\cong\C[\mathcal U,\mathcal V,\mathcal W]
   /(\mathcal U\mathcal W-\mathcal V^{2}),
 \qquad
 \mathcal U=z_{1}^{2}z_{3},\ \mathcal V=z_{1}z_{2}z_{3},\
 \mathcal W=z_{2}^{2}z_{3},
\end{equation}
the rational invariant field being $\C(z_{1}/z_{2},z_{2}^{2}z_{3})$ of
transcendence degree two. Here $r=1$, matching
\cref{diff:ex:torusexamples}(i) in kind: $\omega$ and $-2\omega$ are
$\Q$-dependent. The three first-integral identities are checked in the
\emph{original} coupled coordinates by the delivered suite.
\end{example}

\subsection{The phase lattice as a differential Galois torus}
\label{diff:sub:mattorus}

\Cref{diff:thm:torus} classifies the Picard--Vessiot extension of a
\emph{diagonal} system by its integer relation lattice.
\Cref{diff:thm:matnormal} reduces every finite system to a diagonal one, so
that theorem now applies to all of them. The point of this subsection is only
to record what the reduction gives, and to add the two ingredients the
diagonal treatment did not need: an explicit group-algebra presentation in
rank $r$, and the real form. Nothing here is new Galois theory, and the eighth
manuscript says so of itself.

\begin{corollary}[Every finite system has a phase torus]
\label{diff:cor:mattorus}
Let $A\in\Mat_{n}(\SC)$ with $\Ph(A)=\{P_{1},\dots,P_{n}\}$, phase lattice
$\Pi=\sum_{j}\Z P_{j}$ of rank $r$, and relation lattice $\Lambda$ as in
\cref{diff:eq:lattice}. Then the number $d=n-\operatorname{rank}\Lambda$ of
independent phase generators of \cref{diff:conv:phases}(d) equals $r$, and
\cref{diff:thm:torus} applies to the normal form of $A$: a Picard--Vessiot
extension for $\der y=Ay$ is obtained by adjoining $r$ algebraically
independent generators, and its differential Galois group is
$\gm(\C)^{r}\cong\Hom(\Pi,\C^{\times})$.
\end{corollary}
\begin{proof}
$\Lambda$ is the kernel of the surjection $\Z^{n}\to\Pi$,
$m\mapsto\sum m_{j}P_{j}$, so $\Z^{n}/\Lambda\cong\Pi$ is free of rank $r$ by
\cref{diff:lem:saturated}, giving $d=r$; and
$\Hom(\Pi,\C^{\times})\cong(\C^{\times})^{r}$ is the description of
\cref{diff:eq:torus} in coordinates dual to a $\Z$-basis of $\Pi$. The
Picard--Vessiot statement is \cref{diff:thm:torus} applied to
$i\operatorname{diag}(\der P_{1},\dots,\der P_{n})$, transported back by the
gauge of \cref{diff:thm:matnormal}, which lies in $\GL_{n}(\SC)$ and therefore
generates no new constants.
\end{proof}

\begin{proposition}[The group algebra of the phase lattice, and its real form]
\label{diff:prop:matPValgebra}
Choose a $\Z$-basis $Q_{1},\dots,Q_{r}$ of $\Pi$ and form the Laurent
differential algebra
\begin{equation}\label{diff:eq:matPVring}
 \mathcal L_{\Pi}=\SC[\mathsf T_{1}^{\pm1},\dots,\mathsf T_{r}^{\pm1}],
 \qquad \der\mathsf T_{k}=i(\der Q_{k})\mathsf T_{k},
\end{equation}
writing $\mathsf T_{Q}=\prod_{k}\mathsf T_{k}^{m_{k}}$ for
$Q=\sum_{k}m_{k}Q_{k}\in\Pi$, so that $\mathsf T_{Q}\mathsf T_{Q'}=\mathsf
T_{Q+Q'}$ and $\der\mathsf T_{Q}=i(\der Q)\mathsf T_{Q}$. Then:
\begin{enumerate}[label=\textup{(\roman*)},leftmargin=2.4em]
\item $\mathcal L_{\Pi}$ has no nonzero proper differential ideal;
\item if $G$ normalizes $A$ then
$Y=G\operatorname{diag}(\mathsf T_{P_{1}},\dots,\mathsf T_{P_{n}})$ is a
fundamental matrix for $\der y=Ay$, and its entries together with
$(\det Y)^{-1}$ generate $\mathcal L_{\Pi}$ over $\SC$;
\item for a real system $\Pi$ is stable under negation, the involution
$\overline{c\,\mathsf T_{Q}}=\bar c\,\mathsf T_{-Q}$ commutes with $\der$,
the fixed combinations $(\mathsf T_{Q}+\mathsf T_{-Q})/2$ and
$(\mathsf T_{Q}-\mathsf T_{-Q})/2i$ supply the entries of a real fundamental
matrix in the normal form \cref{diff:eq:matrealnormal}, and the real form of
the torus is $\bigl(\operatorname{Res}^{1}_{\C/\R}\gm\bigr)^{r}$, whose real
points are $(S^{1})^{r}$.
\end{enumerate}
Every element of $\mathcal L_{\Pi}$ is a \emph{finite} sum; no Hahn
summability question arises, and no exponential $e^{iP}$ is defined in $\SC$.
\end{proposition}
\begin{proof}
(i) Let $\mathfrak a\neq0$ be a differential ideal and choose $f\in\mathfrak a$
with the fewest nonzero monomials; multiplying by a Laurent monomial and a
nonzero scalar, arrange $f=1+\sum_{Q\neq0}c_{Q}\mathsf T_{Q}$. Differentiation
kills the constant term and introduces no new monomials, so $\der f$ has fewer
terms and must vanish; then $c_{Q}^{\dagger}=-i\der Q$ for each remaining $Q$,
impossible for $Q\neq0$ by \cref{diff:thm:phasecriterion}. Hence $f=1$.
This is the rank-$r$ form of the rank-one statement proved in
\cref{diff:thm:extension}.

(ii) $\der G=AG-i G\operatorname{diag}(\der P_{j})$ gives $\der Y=AY$ at once.
Conversely the entries of $G^{-1}Y$ supply every $\mathsf T_{P_{j}}$, and
$(\det Y)^{-1}$ supplies the inverse of their product, so multiplying by all
but one supplies each $\mathsf T_{P_{j}}^{-1}$; these characters generate the
group algebra of $\Pi$.

(iii) Stability under negation is the conjugation symmetry used in
\cref{diff:thm:matreal}, and
$\overline{\der(c\mathsf T_{Q})}=\overline{(\der c+ic\der Q)\mathsf T_{Q}}
=(\der\bar c-i\bar c\der Q)\mathsf T_{-Q}=\der(\bar c\mathsf T_{-Q})$. The
fixed combinations satisfy $\der$-relations of the rotation block
$(\der Q)\Rot$, by the computation of
\cref{diff:prop:oscillatorExtension} applied with $a=i\der Q$. An automorphism
commuting with the involution has $\chi(Q)\overline{\chi(Q)}=1$, which is the
asserted real form.
\end{proof}

\begin{warning}[What the torus statement is about]
\label{diff:warn:mattorus}
\Cref{diff:cor:mattorus,diff:prop:matPValgebra} concern the very large base
field $\SC$, with the factorization and surjectivity properties of
\cref{diff:thm:matsurjectivity}. They say nothing comparable for the same
printed equation over a small rational-function base: in particular
\cref{diff:eq:matairyphases} is \emph{not} a claim that the ordinary Airy
differential Galois group over $\C(x)$ is a torus. Enlarging the base to $\SC$
has already supplied the nonoscillatory gauge entries. Nor is the real form of
\cref{diff:prop:matPValgebra}(iii) a construction of trigonometric functions on
surreals; it is an algebraic-group statement, and
\cref{diff:thm:orderobstruction} records what such an extension costs. Finally,
when $\Pi\neq0$ no differential embedding $\mathcal L_{\Pi}\to\SC$ over $\SC$
exists, by \cref{diff:cor:nonembedding}: the extension represents the
obstruction, it does not remove it.
\end{warning}

\subsection{What this part does not settle}
\label{diff:sub:matlimits}

\begin{warning}[The boundaries of the matrix theory]
\label{diff:warn:matlimits}
\Cref{diff:part:matrix} classifies systems $\der y=Ay$ for $A$ an $n\times n$
matrix over $\SC$ with $n$ an ordinary finite integer, and computes the
invariant from algebraic eigenvalues only when $A=iH$ with $H$ Hermitian. It
does not classify non-Hermitian coefficients by their spectra ---
\cref{diff:ex:matnonnormal} refutes that --- does not treat infinite matrices,
does not supply an algorithm for the phase map or the normalizing gauge, and
does not settle which nontriangular differential modules admit a filtration
whose rank-one factors carry computable obstructions, which needs genuine
module-theoretic information rather than the eigenvalues of an instantaneous
coefficient matrix. Every ordinary finite rank over $\SC$ is covered;
classifications over smaller prescribed bases and in infinite dimension are
outside this theorem. All statements are for the one fixed $\der$ of
\cref{diff:thm:BM} and for no other derivation, no class-function derivative
and none of the coordinate operators of \cref{diff:tab:derivatives}. The
consolidated list is \cref{diff:sub:nc:matrixscope}. (Added with member~12: for
one class of coefficients the filtration question and the algorithm question are
answered in part. \Cref{diff:rs:part} splits every power-Hahn regular-singular
system over $\KR[\log t]$ into rank-one factors whose obstructions are read off
from the residual matrix, and for finite rational supports with algebraic
coefficients computes its Hahn invariants by a finite procedure
(\cref{diff:rs:rem:effective}). Everything in
this warning stands outside that class.)
\end{warning}

\begin{remark}[Proper classes, and a set-sized reading]
\label{diff:rem:matclasses}
$\No$ is a proper class, and \cref{diff:sub:classes} fixes the framing this
report uses. Nothing in \cref{diff:part:matrix} makes a proper class an element
of a set or needs a class-sized basis of $\Pcl$: every phase lattice $\Pi$ is
finitely generated and every element of $\mathcal L_{\Pi}$ is a finite sum.
There is also a set-sized reading of each finite-data argument, in the spirit
of \cref{diff:thm:localize}: starting from a set of given coefficients together
with $\R\cup\{\omega\}$, form a countable chain of set-sized subfields of
$\No$, adjoining at each stage the real algebraic roots, derivatives, real
exponentials, purely infinite parts, primitives and finite angles needed by the
previous stage, together with witnesses for all finite factorizations and
first-order inhomogeneous equations of \cref{diff:thm:matsurjectivity}
with coefficients at that stage. These are a \emph{set} of requests, each
specified by a finite tuple; there need not be only finitely many requests.
Witnesses can be chosen set-theoretically: for each request take the least
ordinal birthday bound at which a witness tuple exists, then choose from the
nonempty set of witness tuples at that bound. Replacement and set choice
therefore produce a set of witnesses. The union is a
real closed differential subfield $\No_{0}$ inheriting the $H$-field rule and
the designated witnesses, whose complexification has constant field exactly
$\C$; every finite tuple appears at some stage and its witnesses at a later
one, so the elementary lemmas of \cref{diff:sec:primitives} and all the finite
matrix arguments hold internally. \Cref{diff:prop:matPValgebra} may therefore
be read as an ordinary set-sized Picard--Vessiot statement over such an
$\No_{0}[i]$.

This is a size justification for the proofs and nothing more. It does
\emph{not} say that any fixed Hahn workspace already contains the required
primitives or solutions --- \cref{diff:ex:logmissing} shows one that does not
--- and it gives no effective presentation, no birthday bound and no minimal
such subfield.
\end{remark}
\part{Workspaces: Laurent, rational, Hahn, and set-sized}
\label{diff:part:workspaces}

\section{The Laurent parent and its resolvents}
\label{diff:sec:laurent}

\subsection{A differential subfield with a simple valuation}
\label{diff:sub:laurentparent}

Throughout this section $t=\omega^{-1}$ and, unless stated otherwise,
$\Gamma\subseteq\R$ is an additive subgroup with $1\in\Gamma$, so that
$\Gamma+1=\Gamma$. Put
\[
 K_{\Gamma}=\C((t^{\Gamma}))\subseteq\SC,
 \qquad
 H_{\Gamma}=\R((t^{\Gamma}))\subseteq\No,
\]
embedded by Hahn summation of Conway monomials, and write
$K_{0}=\C((t))$ for the case $\Gamma=\Z$, the ordinary formal Laurent series
field with ordinary complex coefficients. For nonzero $f$ let $v(f)$ be the
least exponent with nonzero coefficient, with $v(0)=+\infty$.

\begin{proposition}[The derivation on a real-exponent workspace]
\label{diff:prop:hahnderiv}
$K_{\Gamma}$ is $\der$-stable, its constants are exactly $\C$, and
\begin{equation}\label{diff:eq:hahnderiv}
 \der\!\left(\sum_{r\in S}a_{r}t^{r}\right)=-\sum_{r\in S}r\,a_{r}t^{r+1},
 \qquad\text{so}\qquad
 v(\der f)\geq v(f)+1 ,
\end{equation}
and hence $v(\der^{n}f)\geq v(f)+n$ for every $n\geq0$. On $K_{0}$ this is
exactly the operator
\begin{equation}\label{diff:eq:laurentD}
 D:=\der|_{K_{0}}=-t^{2}\frac{d}{dt}.
\end{equation}
Furthermore
\begin{equation}\label{diff:eq:hahnimage}
 \der K_{\Gamma}=\{f\in K_{\Gamma}:[t^{1}]f=0\},
\end{equation}
and for $f=\sum a_{r}t^{r}$ with $a_{1}=0$ a primitive inside $K_{\Gamma}$ is
$\sum_{r\neq1}\frac{a_{r}}{1-r}t^{r-1}$; for general $f\in K_{\Gamma}$ the
normalized primitive in the ambient $\SC$ is
\begin{equation}\label{diff:eq:logprimitive}
 \Izero(f)=a_{1}\log\omega+\sum_{r\neq1}\frac{a_{r}}{1-r}\,t^{r-1},
\end{equation}
which is normalized because $\log\omega$ is purely infinite and so has zero
constant coefficient.
\end{proposition}
\begin{proof}
The real-exponent rule $\der t^{r}=-rt^{r+1}$ of \cref{diff:eq:basicder} and
strong additivity \ref{diff:BM:additive} give \cref{diff:eq:hahnderiv};
translation of a well-ordered support by $1$ preserves well ordering and
distinct exponents stay distinct, and $\Gamma+1=\Gamma$ keeps the result in
the field. A nonconstant term has a nonzero derivative at a unique shifted
exponent, so the constants are the exponent-$0$ terms, that is $\C$. The only
term that could contribute at exponent $1$ has exponent $0$, whose multiplier
is zero, which proves necessity in \cref{diff:eq:hahnimage}; conversely the
displayed primitive has a support that is a translate of a subset of
$\supp f$, hence well ordered and inside $\Gamma$, and termwise
differentiation returns $f$. For \cref{diff:eq:logprimitive} use
$\der\log\omega=t$.
\end{proof}

\begin{example}[A missing logarithmic primitive]
\label{diff:ex:logmissing}
$K_{0}=\C((t))$ is a differential field, yet $t$ has \emph{no} primitive in
it: by \cref{diff:eq:hahnderiv} the coefficient of $t^{1}$ in $\der f$ is
always zero, since it could only arise from exponent $0$, whose derivative
coefficient vanishes. In the full surreal field $\der\log\omega=t$. The same
holds over $H_{\Q}=\R((t^{\Q}))$, which is $\der$-stable while $t$ has no
primitive in it, its ambient primitives being $\log\omega+c$ with $c\in\R$.
\end{example}

\begin{warning}[Differentiation closure is not integration closure]
\label{diff:warn:diffnotint}
So $\der$ is \emph{not} surjective on $K_{\Gamma}$, even though $\der$ is
surjective on $\SC$. This is a failure of closure of a chosen representation
parent, not a failure of ambient surjectivity, and the whole-field
surjectivity clause \ref{diff:BM:surj} must not be silently transported to a
localized workspace. Algebraic localization, algebraic closure, derivative
stability, and closure under primitives are four separate properties; see
\cref{diff:sec:localization} for the counterexamples separating them. There is
a formal differential-form reading of the obstruction: since
$d\omega=-t^{-2}\,dt$, the coefficient of $t^{-1}dt$ in $f\,d\omega$ is
$-a_{1}$. That is a formal Laurent/Hahn residue statement only; it is not an
identification with a cluster residue or an actual-point contour residue from
the repository's analytic reports.
\end{warning}

Over the real form, the ideal $\Acl$ intersects the Laurent field exactly.

\begin{corollary}[A valuation test, on real exponents only]
\label{diff:cor:valphase}
For real $f\in H_{\Gamma}$ with $\Gamma\subseteq\R$,
\[
 f\in\Acl
 \quad\Longleftrightarrow\quad
 f=0\ \text{ or }\ v(f)>1,
\]
where $\Acl$ is computed in the \emph{full} surreal field. In particular
$t^{2}\R[[t]]=\Acl\cap\R((t))$ on the real Laurent field.
\end{corollary}
\begin{proof}
This is \cref{diff:thm:power} restated through \cref{diff:eq:logprimitive}. If
$v(f)>1$ every exponent of the primitive is positive, so it is infinitesimal.
If the least exponent $r$ is below $1$, the primitive has the nonzero leading
infinite term $a_{r}t^{r-1}/(1-r)$, which dominates any logarithmic term and
cannot be removed by a real constant. If the least exponent is exactly $1$ the
logarithmic coefficient is nonzero and all other primitive terms are
infinitesimal, so the primitive is again infinite. A logarithmic term cannot
cancel a Laurent principal part, and primitives differ only by real constants.
\end{proof}

\begin{warning}[Not a valuation test for arbitrary coefficients]
\label{diff:warn:valonly}
\Cref{diff:cor:valphase} uses \emph{real exponents relative to
$t=\omega^{-1}$}. It is not a valuation-only classification of arbitrary
surreal coefficients: \cref{diff:ex:logseparator} is exactly the information
that this real-exponent workspace cannot see.
\end{warning}

\subsection{Resolvents as exact Hahn sums}
\label{diff:sub:resolvents}

\begin{theorem}[First-order resolvent on a real-exponent workspace]
\label{diff:thm:resolvent}
Let $\Gamma\subseteq\R$ be an additive subgroup with $1\in\Gamma$, and let
$\lambda\in\C^{\times}$. Then $\der-\lambda$ is a bijection of $K_{\Gamma}$,
with inverse the strongly summable operator series
\begin{equation}\label{diff:eq:resolvent}
 (\der-\lambda)^{-1}f=-\sum_{n\geq0}\lambda^{-n-1}\der^{n}f ,
\end{equation}
and for the truncation
$y_{N}=-\sum_{n=0}^{N}\lambda^{-n-1}\der^{n}f$ the residual is exactly
\begin{equation}\label{diff:eq:residual}
 (\der-\lambda)y_{N}-f=-\lambda^{-(N+1)}\der^{N+1}f .
\end{equation}
\end{theorem}
\begin{proof}
For $f\neq0$ put $S=\supp f$ and $s_{0}=\min S$. By
\cref{diff:eq:hahnderiv}, $\supp\der^{n}f\subseteq S+n$, and $S+\N$ is well
ordered: a decreasing sequence in it would have integer parts bounded above by
its first term minus $s_{0}$, and an infinite subsequence with the same
integer part would be a decreasing sequence in $S$. At a fixed exponent
$\gamma$, a contribution from $\der^{n}f$ requires $\gamma-n\in S$, hence
$n\leq\gamma-s_{0}$, and $\gamma-s_{0}$ is an \emph{ordinary real number}, so
only finitely many $n$ occur. This proves strong summability. Strong
additivity \ref{diff:BM:additive} allows $\der$ to be applied termwise, and
the terms cancel pairwise after reindexing, leaving $f$; the finite identity
\cref{diff:eq:residual} is the same cancellation truncated. For injectivity,
if $y\neq0$ then either $\der y=0$ or $v(\der y)\geq v(y)+1>v(y)$ by
\cref{diff:eq:hahnderiv}, so the leading term of $\lambda y$ cannot be
cancelled; hence $(\der-\lambda)y\neq0$.
\end{proof}

\begin{warning}[The real-exponent hypothesis is substantive]
\label{diff:warn:resolventreal}
The step ``$\gamma-s_{0}$ is an ordinary real number, so only finitely many
$n$ occur'' is where $\Gamma\subseteq\R$ is used. In a higher-rank value group
a fixed increment $1$ need not dominate all relevant scales, and no claim is
made here that successive valuation improvements always converge. A support
argument has not been replaced by a convergence claim, and
\cref{diff:thm:resolvent} is not asserted for an arbitrary set-sized ordered
$\Gamma$.
\end{warning}

\begin{proposition}[Polynomial Laurent inverse]
\label{diff:prop:polyresolvent}
Let $P\in\C[T]$ with $P(0)\neq0$ and formal reciprocal
$1/P(T)=\sum_{n\geq0}q_{n}T^{n}\in\C[[T]]$. For every $f\in K_{0}$ the family
$(q_{n}\der^{n}f)_{n\geq0}$ is strongly summable, $P(\der):K_{0}\to K_{0}$ is
bijective, and
\begin{equation}\label{diff:eq:polyresolvent}
 P(\der)^{-1}f=\sum_{n\geq0}q_{n}\der^{n}f .
\end{equation}
If $Q_{N}(T)=\sum_{n=0}^{N}q_{n}T^{n}$ and $y_{N}=Q_{N}(\der)f$, then
$P(T)Q_{N}(T)-1$ has order at least $N+1$, so
\begin{equation}\label{diff:eq:polyresidual}
 v\bigl(P(\der)y_{N}-f\bigr)\geq v(f)+N+1 .
\end{equation}
\end{proposition}
\begin{proof}
Summability is the valuation bound of \cref{diff:eq:hahnderiv} exactly as in
\cref{diff:thm:resolvent}. Strong additivity permits applying each of the
finitely many derivatives in $P(\der)$ termwise; at every exponent the
calculation is finite, so the formal identity $P(T)\sum q_{n}T^{n}=1$ gives
$P(\der)y=f$. For injectivity, the leading term of $P(\der)y$ is $P(0)$ times
the leading term of $y$, every positive-order derivative having strictly
larger valuation. Alternatively one may factor $P$ over $\C$ and compose the
first-order resolvents of \cref{diff:thm:resolvent} with multiplicities, which
gives the same conclusion; an effective implementation then needs exact
representations of the roots, or can derive a coefficient recurrence from $P$
directly.
\end{proof}

The hypothesis $P(0)\neq0$ is essential, and its failure is exactly
\cref{diff:ex:logmissing}: when $P(0)=0$ the factor $\der$ carries the
obstruction \cref{diff:eq:hahnimage}, which must not be handled by a division
by zero inside a resolvent routine. \Cref{diff:prop:polyresolvent} covers both
$P(T)=T-i$ and $P(T)=T^{2}+1$. So the absence of a nonzero homogeneous
oscillator coexists with unique solvability of \emph{every} inhomogeneous
oscillator equation whose forcing lies in $K_{0}$ --- an explicit instance of
\cref{diff:warn:atmostone}.

\begin{remark}[Parent selection changes homogeneous solution spaces]
\label{diff:rem:parentselection}
For nonreal $\lambda$, $\der-\lambda$ has trivial kernel in all of $\SC$ and
\cref{diff:thm:resolvent}'s uniqueness is uniqueness in $\SC$. For real
nonzero $\lambda$ the Laurent solution is unique \emph{inside} $K_{\Gamma}$,
while the full field also contains $c\,e^{\lambda\omega}$, which lies outside
$K_{\Gamma}$. The same inverse formula holds in both cases. A theorem whose
uniqueness statement changes when the ambient field is enlarged must carry its
parent as part of its statement, and an algorithm must return that parent.
\end{remark}

\subsection{Factorially divergent exact solutions}
\label{diff:sub:factorial}

Since $\der^{n}t=(-1)^{n}n!\,t^{n+1}$, the resolvents produce exact Hahn
series whose coefficients grow factorially. Three instances occur in the
source manuscripts and all three are recorded here.

\begin{example}[$(\der-i)y=t$]
\label{diff:ex:factorialfirst}
Applying \cref{diff:thm:resolvent} with $\lambda=i$ and $f=t$, or equivalently
$1/(s-i)=i\sum_{n\geq0}(-i)^{n}s^{n}$ in
\cref{diff:eq:polyresolvent}, gives the unique solution in $K_{0}$
\begin{equation}\label{diff:eq:factorialfirst}
 y=\sum_{n\geq0}i^{n+1}n!\,t^{n+1}
 =it-t^{2}-2it^{3}+6t^{4}+24it^{5}-120t^{6}-\cdots,
\end{equation}
which by \cref{diff:cor:oscillator} is the unique solution in all of $\SC$.
Its truncation residual is \cref{diff:eq:residual}; explicitly, writing
$y_{N}$ for the sum through $n=N$,
$(\der-i)y_{N}-t=-i^{-(N+1)}(-1)^{N+1}(N+1)!\,t^{N+2}$.
\end{example}

\begin{example}[$(\der+1)Y=t$, and a beyond-all-orders ambiguity]
\label{diff:ex:factorialsecond}
With $\lambda=-1$ and $f=t$,
\begin{equation}\label{diff:eq:factorialsecond}
 Y_{H}=\sum_{n\geq0}n!\,t^{n+1},\qquad (\der+1)Y_{H}=t,
\end{equation}
with support $\{1,2,3,\dots\}$, each exponent occurring once; and for
$Y_{N}=\sum_{n=0}^{N-1}n!\,t^{n+1}$ the exact residual is
\begin{equation}\label{diff:eq:factorialresidual}
 (\der+1)Y_{N}-t=-N!\,t^{N+1}.
\end{equation}
\end{example}

\begin{proposition}[Uniqueness depends on the ambient field]
\label{diff:prop:factorialambiguity}
Equation \cref{diff:eq:factorialsecond} has exactly one solution in
$K_{\Gamma}$, but all its solutions in $\SC$ are
\begin{equation}\label{diff:eq:beyondallorders}
 Y_{H}+C\,e^{-\omega},\qquad C\in\C .
\end{equation}
Every nonzero correction $Ce^{-\omega}$ is smaller in modulus than $t^{r}$ for
every ordinary real $r$, and lies outside $K_{\Gamma}$.
\end{proposition}
\begin{proof}
Workspace uniqueness is \cref{diff:thm:resolvent}. In the full field the
difference of two solutions satisfies $\der Z=-Z$, so
\cref{diff:cor:oscillator} gives $Z=Ce^{-\omega}$. From
$\log\omega\prec\omega$ one gets $-\omega<-r\log\omega$ for every real $r$
after allowing any fixed real multiplicative coefficient; exponentiating gives
$\abs Ce^{-\omega}\prec t^{r}$. A nonzero real-exponent Hahn series has a
least real exponent with a nonzero ordinary leading coefficient, so it cannot
be smaller than every real power of $t$.
\end{proof}

The whole sequence of algebraic-order asymptotic coefficients therefore fails
to distinguish the full-field solutions in
\cref{diff:eq:beyondallorders}. The exact Hahn workspace selects one solution
by its allowed support, not by any limiting process. No Borel summation, no
Stokes multiplier and no analytic continuation rule is asserted by this
example; those would require additional data and separate theorems.

\begin{theorem}[An exact factorially divergent forced oscillator]
\label{diff:thm:factorial}
The equation
\begin{equation}\label{diff:eq:forcedosc}
 \der^{2}y+y=\omega^{-1}
\end{equation}
has the unique solution in $\SC$
\begin{equation}\label{diff:eq:factorial}
 y=\sum_{n\geq0}(-1)^{n}(2n)!\,\omega^{-(2n+1)},
\end{equation}
and for $N\geq1$ its $N$-term truncation $y_{N}$ satisfies exactly
\begin{equation}\label{diff:eq:factorialres}
 y_{N}''+y_{N}-\omega^{-1}=(-1)^{N-1}(2N)!\,\omega^{-(2N+1)} .
\end{equation}
\end{theorem}
\begin{proof}
The formal inverse of $1+T^{2}$ is $\sum_{n\geq0}(-1)^{n}T^{2n}$; apply
\cref{diff:prop:polyresolvent} to $f=t$, using
$\der^{k}t=(-1)^{k}k!\,t^{k+1}$. Existence can also be checked directly:
$\der^{2}t^{2n+1}=(2n+1)(2n+2)t^{2n+3}$, so every positive odd exponent beyond
the first cancels between $\der^{2}y$ and $y$; the same cancellation for a
finite truncation leaves exactly \cref{diff:eq:factorialres}. The difference of
two surcomplex solutions satisfies $h''+h=0$, hence is zero by
\cref{diff:thm:constant}, so uniqueness holds in the whole class field and not
merely in $K_{0}$.
\end{proof}

\begin{warning}[An exact value that is not a limit of its truncations]
\label{diff:warn:antitopological}
The ordinary scalar series $\sum(-1)^{n}(2n)!\,t^{2n+1}$ has radius of
convergence zero: for nonzero ordinary complex $t$ the ratio of successive
term magnitudes is $(2n+2)(2n+1)\abs t^{2}$, which grows without bound. That
causes no difficulty for \cref{diff:eq:factorial}: its coefficients are
ordinary complex numbers and its support is the well-ordered set
$\{1,3,5,\dots\}$, which is all Hahn summation requires. Conversely the
existence of the surreal value does not make the scalar series an ordinary
holomorphic germ, nor is the value a fine-topology limit of its truncations:
the tail $y-y_{N}$ has leading term
$(-1)^{N}(2N)!\,\omega^{-(2N+1)}$, so
\[
 \abs{y-y_{N}}>\omega^{-\omega}\qquad\text{for every }N\geq1 ,
\]
and the errors never enter that fixed infinitesimal neighbourhood of zero. The
residual identities are exact algebraic identities and scale certificates;
they certify a formal remainder at a stated order, not convergence of the
sequence $(y_{N})$ in the fine surreal topology. A numerical substitution such as
$t=10^{-6}$ would not verify a Hahn identity and is not used anywhere in this
report. A closed formula for a coefficient stream is not a finite list of every
coefficient.
\end{warning}

\begin{remark}[The coefficient recursion, and its exact guarantee]
\label{diff:rem:coefrecursion}
Writing $(\der-\lambda)y=f$ coefficientwise in $K_{0}$ gives
\begin{equation}\label{diff:eq:coefrecursion}
 -\lambda y_{n}-(n-1)y_{n-1}=f_{n},
 \qquad\text{that is}\qquad
 y_{n}=-\lambda^{-1}\bigl(f_{n}+(n-1)y_{n-1}\bigr),
\end{equation}
and below a known lower exponent bound all coefficients vanish. Each requested
finite prefix costs one inversion of $\lambda$ and a linear number of
coefficient operations, and the recursion works for negative lower bounds. It
does \emph{not} make equality of arbitrary program-generated coefficient
streams decidable.
\end{remark}

\section{The rational-exponent workspace}
\label{diff:sec:rationalworkspace}

A particularly informative set-sized workspace is
\begin{equation}\label{diff:eq:HahnQ}
 \Hring:=\C((t^{\Q})),\qquad t=\omega^{-1},
 \qquad \Hring_{\R}:=\R((t^{\Q})),
\end{equation}
embedded in $\SC$ through normal forms. Every element has a well-ordered
\emph{rational} support: this is a \emph{full} Hahn field, not merely the union
of the finite-ramification Laurent-series fields, and in particular
well-ordered supports need not have a common denominator. By
\cref{diff:prop:hahnderiv} it is $\der$-stable with
$\der(\sum c_{q}t^{q})=\sum(-q)c_{q}t^{q+1}$, so the restriction of the
abstract derivation is completely explicit here. That explicitness does not
mean the full Berarducci--Mantova construction has been implemented: the
special rational monomial subgroup is doing essential work
(\cref{diff:warn:monomialmap}).

\begin{remark}[A companion report that uses this workspace]
\label{diff:rem:tailspans}
The companion report \emph{Tail-Spans and Differential Transcendence in
Surreal and Surcomplex Hahn Fields}
(\path{docs/surreal/tail-spans-and-differential-transcendence/},
\cite{RepoTailSpans}) takes \cref{diff:eq:HahnQ} and
\cref{diff:prop:hahnderiv} as its starting point rather than reproving
them, and works inside the real-algebraic subfield
$\mathbb A_{\R}((t^{\Q}))\subseteq\Hring_{\R}$, where the displayed
formula still applies and keeps that subfield $\der$-stable, since it
multiplies each coefficient by a rational number and shifts its exponent
by one. Over the base $\Q((t^{\Q}))$ it produces a family of actual
positive infinitesimal surreal numbers, of cardinality $2^{\aleph_{0}}$,
whose $\der$-jets are all algebraically independent; the numbers are
$\sum_{n\in A}\sqrt{p_{n}}\,\omega^{-n}$ for an explicit almost disjoint
family of infinite sets $A\subseteq\N_{>0}$. That is a statement about
differential transcendence over a \emph{set-sized} Hahn base and does not
bear on the image and primitive questions settled here, nor does it supply
any derivative formula for a monomial outside the rational exponents
(\cref{diff:warn:monomialmap} still applies).
\end{remark}

\begin{theorem}[Exact derivative image: one resonant coefficient]
\label{diff:thm:Hahnintegral}
For $f=\sum_{q}c_{q}t^{q}\in\Hring$,
\begin{equation}\label{diff:eq:Hahnintegrability}
 f\in\der(\Hring)\quad\Longleftrightarrow\quad [t^{1}]f=0 .
\end{equation}
If $c_{1}=0$ the normalized primitive inside $\Hring$ is
$-\sum_{q\neq1}\frac{c_{q}}{q-1}t^{q-1}$; for arbitrary $f\in\Hring$ the
normalized primitive in the ambient $\SC$ is
$c_{1}\log\omega-\sum_{q\neq1}\frac{c_{q}}{q-1}t^{q-1}$. Hence $\Hring$ is
closed under differentiation but not under antidifferentiation, and
$\log\omega\notin\Hring$.
\end{theorem}
\begin{proof}
This is \cref{diff:prop:hahnderiv} with $\Gamma=\Q$. The formula remains valid
for infinite rational supports: translation of an already well-ordered support
is the precise support certificate, and no denominator bound or discrete
increasing sequence of exponents is needed. No appeal to topology,
cardinality or computational limitation is involved; the obstruction is the
single resonant coefficient $[t]f$.
\end{proof}

\begin{theorem}[Exact logarithmic-derivative image in $\Hring$]
\label{diff:thm:Hahnlogder}
\begin{equation}\label{diff:eq:Hahnlogderimage}
 \{y^{\dagger}:y\in\Hring^{\times}\}
 =\{-rt+g:\ r\in\Q,\ g=0\text{ or }v(g)>1\}.
\end{equation}
For $q=-rt+g$ in this class, put
\begin{equation}\label{diff:eq:workspacefactor}
 k=-\sum_{\alpha>1}\frac{[t^{\alpha}]g}{\alpha-1}\,t^{\alpha-1},
\end{equation}
which has positive valuation (or is $0$) and satisfies $\der k=g$. Then all
solutions in $\Hring$ of $\der y=qy$ are
\begin{equation}\label{diff:eq:Hahnall}
 y=c\,t^{r}\Emm(k),\qquad c\in\C .
\end{equation}
\end{theorem}
\begin{proof}
A nonzero $y\in\Hring$ has a unique leading-term factorization
$y=c\,t^{r}(1+h)$ with $c\in\C^{\times}$, $r\in\Q$ and $v(h)>0$ (possibly
$h=0$), and then
\[
 y^{\dagger}=-rt+\frac{\der h}{1+h}=-rt+\der\Lmm(1+h)
\]
by \cref{diff:lem:smallchain}. Since $\Lmm(1+h)$ has positive valuation or is
zero, every nonzero term of its derivative has exponent strictly greater than
$1$ by \cref{diff:eq:hahnderiv}. This gives containment in the right-hand
class. Conversely, for $v(g)>1$ the series \cref{diff:eq:workspacefactor} has
positive valuation and $\der k=g$, so $\Emm(k)\in\Hring$ is defined by a
positive-support Hahn sum and
$\bigl(t^{r}\Emm(k)\bigr)^{\dagger}=-rt+g$. The ratio of two nonzero solutions
lies in the constant field $\C$.
\end{proof}

\begin{proposition}[Workspace versus ambient]
\label{diff:prop:Hahnambient}
Let $q=a+ib\in\Hring$ with $a,b\in\Hring_{\R}$. A nonzero solution of
$\der y=qy$ exists in the \emph{ambient} $\SC$ if and only if
\begin{equation}\label{diff:eq:Hahnambient}
 b=0\qquad\text{or}\qquad v(b)>1 .
\end{equation}
\end{proposition}
\begin{proof}
This is \cref{diff:cor:valphase} with $\Gamma=\Q$, equivalently
\cref{diff:thm:power}; the growth comparison $\log\omega\ll\omega^{\delta}$
for positive rational $\delta$ is what makes a nonzero $\log\omega$ term
infinite.
\end{proof}

So there are \emph{three} distinct outcomes, not a binary solvable/unsolvable
flag: a solution in the declared workspace, a solution only in the ambient
field, and no nonzero solution at all.

\begin{table}[htbp]
\centering\small
\caption{Where a solution lives. $\Hring=\C((t^{\Q}))$, $t=\omega^{-1}$.}
\label{diff:tab:workspaceambient}
\begin{tabularx}{\textwidth}{@{}L{0.25\textwidth}L{0.16\textwidth}Y@{}}
\toprule
Coefficient $q$ in $\der y=qy$ & Location & Explanation, or one solution \\
\midrule
$1$ & $\SC$, not $\Hring$ &
$e^{\omega}$; the exponent-$0$ coefficient violates
\cref{diff:eq:Hahnlogderimage} \\
$\sqrt2\,t$ & $\SC$, not $\Hring$ &
$\exp(\sqrt2\log\omega)$; the exponent-$1$ coefficient is irrational \\
$it$ & Neither & The accumulated phase is $\log\omega$ \\
$it^{2}$ & $\Hring$ & $\Emm(-it)$ \\
$-\tfrac32 t+(1+2i)t^{2}$ & $\Hring$ & $t^{3/2}\Emm(-(1+2i)t)$ \\
\bottomrule
\end{tabularx}
\end{table}

The first two rows mean that no \emph{nonzero} solution belongs to $\Hring$;
multiplication by an ordinary complex constant cannot remove either
obstruction, since it does not change a logarithmic derivative. Note that the
second row is the workspace failure announced in
\cref{diff:rem:realnoobstruct}: the real coefficient never obstructs
\emph{ambiently}, but $A$ or $e^{A}$ may leave a declared workspace.

\begin{remark}[A resonant normal form]
\label{diff:rem:resonant}
\Cref{diff:thm:Hahnlogder} has a normal-form reading. A scalar equation over
$\Hring$ has a fundamental solution in the same field only when its
logarithmic coefficient splits into a rational monomial contribution $-rt$ and
a higher-valuation perturbation: the monomial supplies the exponent-$1$
residue, and the infinitesimal exponential absorbs the perturbation. Lower
exponents require larger-scale real exponentials, and lower \emph{imaginary}
exponents face the stronger ambient phase obstruction. The condition at
exponent $1$ is much more restrictive for logarithmic than for additive
integration: a logarithmic derivative in $\Hring$ has no exponents below $1$,
its exponent-$1$ coefficient is a real \emph{rational} number, and its higher
exponents may carry arbitrary ordinary complex coefficients. This resembles a
one-dimensional normal-form calculation, but every assertion is algebraic and
exact: there is no asymptotic error term in \cref{diff:eq:Hahnall}, and the
rational exponent $r$ is not an approximation to a more general exponent.
\end{remark}
\section{Differentially stable workspaces for set-sized data}
\label{diff:sec:localization}

The repository's localization principle is algebraic: a set of surcomplex
values fits inside a set-sized Hahn workspace \cite{RepoFound}. That is a
\emph{size} statement. It does not say that the workspace contains the
derivatives, primitives, exponentials or finite phases of its members. Two
genuinely different refinements are available, and both survive the merge: a
\emph{full Hahn field} indexed by a closed exponent group, and a set-sized
\emph{subfield} closed under a specified list of finitary operations. They
prove different things and neither implies the other.

\subsection{Full Hahn fields: the monomial criterion and the hull}
\label{diff:sub:hull}

\begin{definition}
\label{diff:def:monomialclosed}
A set-sized additive subgroup $\Gamma\subseteq\No$ is \emph{differentially
closed on monomials} if $\supp\der(t^{\gamma})\subseteq\Gamma$ for every
$\gamma\in\Gamma$, the support being taken in the $t$-exponent convention of
\cref{diff:eq:tmonomial}.
\end{definition}

\begin{lemma}[Monomial criterion]
\label{diff:lem:monomialcriterion}
For a set-sized additive group $\Gamma\subseteq\No$,
\[
 \der K_{\Gamma}\subseteq K_{\Gamma}
 \quad\Longleftrightarrow\quad
 \Gamma\text{ is differentially closed on monomials},
\]
and the same equivalence holds for $H_{\Gamma}=\R((t^{\Gamma}))$.
\end{lemma}
\begin{proof}
Necessity: apply the closure condition to each $t^{\gamma}$. Sufficiency: for
$x=\sum_{\gamma\in S}c_{\gamma}t^{\gamma}$ with $S\subseteq\Gamma$, strong
additivity \ref{diff:BM:additive} gives the summable family
$(c_{\gamma}\der t^{\gamma})_{\gamma\in S}$ and
$\der x=\sum_{\gamma\in S}c_{\gamma}\der t^{\gamma}$. Every term has support in
$\Gamma$, so the sum does too; ordinary real and complex coefficients are
constants of $\der$.
\end{proof}

\begin{theorem}[Differential Hahn hull]
\label{diff:thm:hull}
Let $\Gamma_{0}\subseteq\No$ be a set-sized additive subgroup. Define
\begin{equation}\label{diff:eq:hullstage}
 \Gamma_{n+1}=\Span_{\Q}\!\left(\Gamma_{n}\cup
 \bigcup_{\gamma\in\Gamma_{n}}\supp\der(t^{\gamma})\right),
 \qquad
 \Gamma_{*}=\bigcup_{n<\omega}\Gamma_{n}.
\end{equation}
Then $\Gamma_{*}$ is a set-sized \emph{divisible} additive group,
$K_{\Gamma_{*}}$ is algebraically closed, and
$\der K_{\Gamma_{*}}\subseteq K_{\Gamma_{*}}$. It is the least such divisible
monomial-closed exponent group containing $\Gamma_{0}$. Consequently every set
$A\subseteq\SC$ lies in a set-sized algebraically closed Hahn field stable
under $\der$: take $\Gamma_{0}$ to be the group generated by the union of the
normal-form supports of the members of $A$.
\end{theorem}
\begin{proof}
Each $\der(t^{\gamma})$ is an individual surreal and has a set support, so by
replacement and union the displayed union is a set, as is its rational span.
The stages increase, so their countable union is a set-sized additive group;
every member lies in a stage whose successor contains all its rational
multiples, so the union is divisible. If $\gamma\in\Gamma_{*}$ it lies in some
$\Gamma_{n}$, and then $\supp\der(t^{\gamma})\subseteq\Gamma_{n+1}$. Now take
any $f=\sum_{\gamma\in S}c_{\gamma}t^{\gamma}\in K_{\Gamma_{*}}$: strong
additivity gives the summable family $(c_{\gamma}\der t^{\gamma})_{\gamma\in S}$,
each member has support in $\Gamma_{*}$, so their sum has support in
$\Gamma_{*}$ and equals $\der f$. Divisibility makes $H_{\Gamma_{*}}$ real
closed and $K_{\Gamma_{*}}$ algebraically closed; this is the standard Hahn
field fact, an \emph{algebraic} statement, not a differential one. Minimality:
any divisible monomial-closed group containing $\Gamma_{0}$ contains each stage
by induction, hence their union.
\end{proof}

\begin{warning}[The proof must not use a union of fields]
\label{diff:warn:antilemma}
The last step of the proof appeals to strong additivity applied to every
series over $\Gamma_{*}$ \emph{simultaneously}. It does not use
\[
 K_{\bigcup_{n}\Gamma_{n}}=\bigcup_{n}K_{\Gamma_{n}},
\]
which is generally \emph{false}: one Hahn series can use exponents appearing at
arbitrarily late stages. The correct route is the one displayed, through the
monomial supports.
\end{warning}

\begin{warning}[Three scope limits on \cref{diff:thm:hull}]
\label{diff:warn:hulllimits}
(i) Minimality is only among \emph{full Hahn fields indexed by divisible
exponent groups}, not among all differential subfields. (ii) The hull group
$\Gamma_{*}$ need not be countable: it is the number of closure \emph{stages}
that is countable, and the countably many stages do not constitute a
terminating algorithm for every input representation. (iii) The theorem is a
semantic existence construction, not an effective algorithm for computing
$\der(t^{\gamma})$ at arbitrary surreal exponents, and it supplies no optimal
birthday or cardinal bound. It is not primitive closure, and it asserts no
closure under real exponentiation, arbitrary logarithms, or solutions of all
first-order linear equations; \cref{diff:ex:logmissing} is the witness.
\end{warning}

\begin{proposition}[The real-exponent case, exactly]
\label{diff:prop:realgroup}
Let $0\neq\Gamma\subseteq\R$ be an additive subgroup. Then $K_{\Gamma}$ is
$\der$-stable if and only if $1\in\Gamma$; and if $\Gamma$ is divisible its
differential Hahn hull is $\Gamma+\Q$. The group $\{0\}$ gives the constant
field and is already stable.
\end{proposition}
\begin{proof}
For nonzero $\gamma\in\Gamma$, \cref{diff:eq:basicder} gives
$\der t^{\gamma}=-\gamma t^{\gamma+1}$, so closure requires
$\gamma+1\in\Gamma$, hence $1\in\Gamma$; conversely containing $1$ provides
every exponent $\gamma+1$. The hull assertion is \cref{diff:thm:hull}.
\end{proof}

\begin{example}[Divisible, algebraically closed, and not derivative-stable]
\label{diff:ex:sqrt2}
$\Gamma=\Q\sqrt2$ is divisible, so $K_{\Gamma}$ is algebraically closed and
contains $t^{\sqrt2}$; but
\[
 \der t^{\sqrt2}=-\sqrt2\,t^{\sqrt2+1}\notin K_{\Q\sqrt2},
 \qquad\text{since}\quad \sqrt2+1\notin\Q\sqrt2 .
\]
Its hull is $\Q+\Q\sqrt2$. So a divisible Hahn value group need not define a
derivative-stable workspace, and an arbitrary algebraically closed Hahn
workspace does not automatically support the ambient derivation. Together with
\cref{diff:ex:logmissing} this separates \emph{four} requirements: input
containment, algebraic closure, derivative stability, and closure under
primitives. This example does not dispute the repository's theorem that any set
of surcomplex inputs fits inside \emph{some} set-sized Hahn workspace
\cite{RepoFound}; it shows that the workspace must be chosen with the
derivation in view. Closure properties of surreal subfields under exponential,
logarithmic, derivative and antiderivative operations are studied much more
finely by Bournez and Guilmant \cite{BG}; that literature addresses constraints
not imposed by the elementary constructions here, and no theorem of it is used
to justify \cref{diff:thm:hull}.
\end{example}

\subsection{Set-sized subfields closed under a named operation list}
\label{diff:sub:subfieldclosure}

\begin{theorem}[Closure localization]
\label{diff:thm:localize}
Let $S\subseteq\SC$ be a set. There is a set-sized \emph{real closed} subfield
$F\subseteq\No$ with $\R\cup\{\omega\}\subseteq F$ and $S\subseteq F[i]$ such
that $F$ is closed under
\[
 \der,\quad \Izero,\quad \pp,\ \ct,\ \fin,\quad \exp,\quad
 \log\text{ on positive inputs},
\]
and under the real and imaginary components of $\Emm$ and $\Lmm$ on their
stated domains --- equivalently, under finite-angle sine and cosine on
$F\cap\Ofin$. One may arrange
\[
 \abs F\leq\kappa:=\max\{\abs S,\abs\R,\aleph_{0}\}.
\]
Then $L=F[i]$ is algebraically closed, is a differential subfield of $\SC$
with $\der L=L$ and constants exactly $\C$, and is closed under componentwise
normalized primitives; $\Ker(\der|_{F})=\R$ and $\der F=F$. Moreover the
criterion of \cref{diff:thm:phasecriterion} holds \emph{internally}: for
$a,b\in F$ a nonzero solution of $\der y=(a+ib)y$ exists in $L$ if and only if
$b$ has a finite accumulated phase in $\No$, and then $\Obs$ takes values in
$F\cap\Pcl$ and is onto it, so the gauge classification and the exact sequence
of \cref{diff:sec:gauge} read as ordinary set-level results for this
workspace.

One may in addition require that $f(h)\in L$ for every $f\in\C[[X]]$ with
ordinary complex coefficients and every $h\in L\cap\mm_{\C}$, at the same
cardinality bound, using $\abs\C^{\aleph_{0}}=\abs\R$ to count the formal
series.
\end{theorem}
\begin{proof}
Let $F_{0}$ be the real closure inside $\No$ of the field generated by $\R$,
$\omega$ and the real and imaginary parts of the members of $S$. Given
$F_{n}$, adjoin the images of all eligible finite tuples under the listed class
functions --- with their declared domains: positive arguments for $\log$,
finite arguments for the phase components, infinitesimal arguments for
$\Emm,\Lmm$ and for the optional formal-series values --- take the field
generated by those images, and let $F_{n+1}$ be its real closure inside $\No$,
that is, adjoin all real roots in $\No$ of nonzero polynomials over the
current field.

Each stage is a set: images under fixed definable class functions are sets by
replacement, field generation uses finite expressions, and a set of nonzero
polynomials of finite degree has a set of roots, finitely many each. Each stage
has size at most $\kappa$, since the operations have finite arity; for the
optional clause there are $\abs\C^{\aleph_{0}}=\abs\R$ formal series, so at
most $\kappa$ elements are adjoined per stage and a real algebraic closure of a
field of that cardinality has the same cardinality.

Put $F=\bigcup_{n<\omega}F_{n}$, a set of size at most $\kappa$. Any input to
one of the operations lies in a stage and its image in the next; any finite
tuple of polynomial coefficients lies in a single stage and all its real
algebraic roots are adjoined at the next. So $F$ has the advertised closures
and is real closed, and $L=F[i]$ is algebraically closed. Differential closure
holds componentwise, $\der F=F$ because $\Izero F\subseteq F$, and the constant
fields are $F\cap\R=\R$ and $L\cap\C=\C$ because $\R\subseteq F$.

For the internal criterion, necessity is inherited from the ambient statement.
For sufficiency, if $b\in F$ has a finite accumulated phase then
$\Izero(b)\in F\cap\mm$ by \cref{diff:thm:idealstructure}(ii); for $a\in F$ we
have $\exp(\Izero a)\in F$, and $\Emm(i\Izero b)\in L$ because its real and
imaginary coordinates are finite-angle cosine and sine of an element of
$F\cap\Ofin$. So the witness \cref{diff:eq:generator} lies in $L$. Whenever a
nonzero solution exists, \emph{every} solution lies in $L$ as well, since the
multipliers are in $\C\subseteq L$; the same applies to the solution spaces of
\cref{diff:thm:constant,diff:cor:matrix}. Finally $\pp$ and $\Izero$ preserve
$F$, so $\Obs$ maps $F$ into $F\cap\Pcl$, and
$\Obs(\der P)=P$ for $P\in F\cap\Pcl$ gives surjectivity.
\end{proof}

\begin{warning}[What \cref{diff:thm:localize} does not give]
\label{diff:warn:localizelimits}
It adjoins values of \emph{specified finitary operations}, some of which are
defined in the ambient field by Hahn sums. It does \emph{not} assert closure
under all set-indexed summable families: an arbitrary-support closure is a
different and strictly stronger demand, and $F$ need not contain every Hahn
series over its own monomials. The optional clause adjoins formal-series
values with \emph{constant} coefficients at a \emph{single} infinitesimal
argument; that is a specified set of operations, not Hahn completeness. A
particular coherent family and its requested evaluations can be included as
set-sized data, but arbitrary future families are not automatically covered.
The theorem is not proved to produce a full Hahn field of the form
$\C((t^{\Gamma}))$ --- \cref{diff:thm:hull} is the construction that does that,
and the two are not interchangeable --- and it provides no finite data
representation, no algorithm for the adjoined maps, no birthday bound, no
computable enumeration, and no termination claim for symbolic integration.
Normalization of primitives need not be computable from a supplied symbolic
representation. Structured fields with simultaneous exponential, logarithmic,
differential and integration closure are studied much more finely in
\cite{BG}.
\end{warning}

\begin{remark}[A weaker but sometimes sufficient variant]
\label{diff:rem:finitaryclosure}
If a full Hahn field is not needed, a lighter construction suffices: start
from the field generated by the data and the ordinary constants; at each finite
stage adjoin the values of a specified finitary operation set --- for instance
$\der$, $\Izero$, conjugation, and real $\exp$ and $\log$ where defined --- and
take the generated field; the union over $\N$ is a set-sized field stable
under each listed operation, because any finite tuple from the union occurs at
a common stage and its operation values at the next. This variant is
\emph{not} proved to be a full Hahn field and need not contain every summable
family of its own members. The distinction is useful for formalization: a small
operational workspace can support a particular proof even when it is not closed
under all possible Hahn sums.
\end{remark}

\begin{remark}[Localization cannot supply an oscillator]
\label{diff:rem:nolocaloscillator}
Adjoining a solution of $T'=iT$ cannot be accomplished by passing to a larger
workspace of either kind \emph{inside} $\SC$: the proof of
\cref{diff:cor:oscillator} already applies to the full class. The next section
changes the differential field itself rather than its set-sized container.
\end{remark}

\section{Where an oscillatory solution can live}
\label{diff:sec:PV}

We use Picard--Vessiot terminology in the standard sense of
\cite{PS,Singer}: a solution-generated differential extension with \emph{no
new constants}. Five explicit constructions occur across the source
manuscripts, over four base fields (with two constructions over $\SC$).
They are genuinely different theorems and all are recorded; none of them
appeals to an unqualified existence theorem for differential extensions of a
proper-class field. Other parts compute further differential Galois groups:
the phase-lattice torus of an arbitrary finite system
(\cref{diff:cor:mattorus}), the additive group of one logarithm
(\cref{diff:rs:thm:minlog}), and, over a set-sized Hahn base at every limit
stage of the log-atomic tower, the Borel subgroup of $\SL_{2}(\C)$
(\cref{diff:cp:thm:pv}), the only non-commutative one in this report.

\subsection{The missing-solution extension over \texorpdfstring{$\SC$}{No[i]}}
\label{diff:sub:PVone}

\begin{theorem}[Adjoining one missing rank-one solution]
\label{diff:thm:PVone}
Let $a\in\SC$ with $\Phi(a)\neq0$. Introduce a transcendental $\mathsf T$ and extend
$\der$ to the rational function field $\SC(\mathsf T)$ by
\begin{equation}\label{diff:eq:DU}
 \der\mathsf T=a\mathsf T ,
\end{equation}
using the product rule on polynomials and then the quotient rule. Then the
constant field of $\SC(\mathsf T)$ is exactly $\C$; $\SC(\mathsf T)$ is a Picard--Vessiot
extension for $\der y=ay$; its differential automorphisms over $\SC$ are
exactly $\mathsf T\mapsto c\mathsf T$ with $c\in\C^{\times}$; and its differential Galois
group is $\gm(\C)$. If $\Phi(a)=0$ the corresponding extension is trivial.
\end{theorem}
\begin{proof}
The product and quotient rules extend the derivation uniquely from $\SC[\mathsf T]$ to
$\SC(\mathsf T)$. Let $R=P/Q\neq0$ have derivative zero, with nonzero coprime
$P,Q\in\SC[\mathsf T]$. Then $Q\der P=P\der Q$, and coprimality gives $P\mid\der P$
and $Q\mid\der Q$; since $\der\mathsf T=a\mathsf T$, differentiation does not increase the
degree in $\mathsf T$, so $\der P=hP$ and $\der Q=hQ$ for the same $h\in\SC$. If
$p_{j}\mathsf T^{j}$ and $p_{k}\mathsf T^{k}$ are two nonzero terms of $P$ with $j\neq k$,
comparing coefficients gives $\der p_{j}/p_{j}+ja=h=\der p_{k}/p_{k}+ka$, that
is
\[
 (p_{j}/p_{k})^{\dagger}=(k-j)a .
\]
Applying $\Phi$ and using \cref{diff:cor:tensor} yields
$(k-j)\Phi(a)=0$, impossible since $\Pcl$ is torsion-free. So $P$ is a single
term, and likewise $Q$, whence $R=c\mathsf T^{m}$ with $c\in\SC^{\times}$ and
$m\in\Z$. The constancy equation is $c^{\dagger}+ma=0$; applying $\Phi$ again
forces $m=0$, and then $\der c=0$ gives $c\in\C^{\times}$.

An automorphism over $\SC$ sends $\mathsf T$ to another nonzero solution of
\cref{diff:eq:DU}; the quotient by $\mathsf T$ has derivative zero, hence lies in
$\C^{\times}$. Conversely every such scaling is a differential automorphism.
When $\Phi(a)=0$ a nonzero solution already lies in $\SC$ by
\cref{diff:thm:phasecriterion}.
\end{proof}

Transcendence is necessary: algebraic closedness of $\SC$ leaves no nontrivial
algebraic extension in which a missing solution could appear, and an algebraic
root cannot cancel a nonzero obstruction (\cref{diff:cor:tensor}). $\mathsf T$ is a
genuinely new differential element, not an alternative notation for a
pre-existing surcomplex number.

\begin{theorem}[The same extension, with differential simplicity]
\label{diff:thm:extension}
For $a=i$ the conclusions of \cref{diff:thm:PVone} hold, and in addition the
Laurent polynomial ring $\SC[\mathsf T,\mathsf T^{-1}]$ has no nonzero proper differential
ideal; so $\SC(\mathsf T)$ has the explicit Picard--Vessiot properties for
$\der y=iy$ with differential Galois group $\C^{\times}$. A second, independent
proof of the constant-field statement runs as follows.
\end{theorem}
\begin{proof}
Let $R\in\SC(\mathsf T)^{\times}$ be constant. Since $\SC$ is algebraically closed,
factor numerator and denominator:
\[
 R=c\,\mathsf T^{m}\prod_{j=1}^{k}(\mathsf T-a_{j})^{n_{j}},
\]
with $c\in\SC^{\times}$, $m\in\Z$, the $a_{j}\in\SC^{\times}$ distinct and the
$n_{j}$ nonzero integers. Logarithmic differentiation gives
\begin{equation}\label{diff:eq:rationalconstant}
 0=c^{\dagger}+i\left(m+\sum_{j}n_{j}\right)
 +\sum_{j}\frac{n_{j}(ia_{j}-\der a_{j})}{\mathsf T-a_{j}} .
\end{equation}
Each $ia_{j}-\der a_{j}$ is nonzero by \cref{diff:cor:oscillator}, and
distinct simple poles cannot cancel, so $k=0$. What remains is
$c^{\dagger}=-im$, which \cref{diff:thm:phasecriterion} excludes unless
$m=0$; then $c\in\C^{\times}$.

For differential simplicity, suppose $\mathfrak{q}$ is a nonzero proper
differential ideal of $\SC[\mathsf T,\mathsf T^{-1}]$. Choose a nonzero element of
$\mathfrak{q}$ with the fewest nonzero monomials and multiply by a unit to make
its constant coefficient $1$.
If it has another term $a_{j}\mathsf T^{j}$ then $j\neq0$ and
$\der(a_{j}\mathsf T^{j})=(\der a_{j}+ija_{j})\mathsf T^{j}\neq0$, again by
\cref{diff:thm:phasecriterion}. Its derivative is therefore a nonzero element
of $\mathfrak{q}$ with one fewer nonzero monomial, the constant term
differentiating to zero, contradicting minimality. So the chosen element is $1$
and $\mathfrak{q}$ is the whole ring.
\end{proof}

\subsection{Set-sized bases, and a non-embedding}
\label{diff:sub:PVsetsized}

The proper-class base can be avoided entirely. Each of the following is a
complete construction over a set-sized base, with its own no-new-constants
proof.

\begin{theorem}[Over a localized workspace]
\label{diff:thm:PVlocal}
Let $F$ be as in \cref{diff:thm:localize}, put $L=F[i]$, and let $\mathsf T$ be
transcendental over $L$ with $\der\mathsf T=i\mathsf T$, the derivation extended by the
product and quotient rules. Then the constant field of $L(\mathsf T)$ is exactly $\C$,
$L(\mathsf T)$ is a scalar Picard--Vessiot extension for $\der y=iy$, and its
differential automorphisms over $L$ are exactly $\mathsf T\mapsto c\mathsf T$ with
$c\in\C^{\times}$, so the group is $\gm(\C)$.
\end{theorem}
\begin{proof}
Every rational $R(T)\in L(\mathsf T)$ has an injective formal Laurent expansion at
$\mathsf T=0$, $R=\sum_{n\geq n_{0}}a_{n}T^{n}\in L((T))$, with supports bounded below
in $\Z$; the induced derivation is
\[
 \der R=\sum_{n\geq n_{0}}\bigl(\der a_{n}+ni\,a_{n}\bigr)T^{n},
\]
which agrees with the rational-function derivation by the product and inverse
rules. If $\der R=0$ then $\der a_{n}=-ni\,a_{n}$ for every $n$. For $n\neq0$
the constant $-ni$ is nonreal, so $a_{n}=0$ by \cref{diff:prop:coarse} applied
in the ambient field. For $n=0$, $\der a_{0}=0$ gives $a_{0}\in\C$. Hence
$R=a_{0}$ and there are no new constants. The automorphism statement is as in
\cref{diff:thm:PVone}. No black-box existence theorem is used: the field and
its constants have been constructed and checked.
\end{proof}

\begin{theorem}[Over the rational Hahn field]
\label{diff:thm:PVhahn}
Let $K_{0}^{\Q}=\C((t^{\Q}))$ with its explicit coefficientwise derivation
$\der(\sum_{q}c_{q}t^{q})=\sum_{q}(-q)c_{q}t^{q+1}$. This field is set-sized
and algebraically closed, its differential constants are $\C$, it contains
$\omega=t^{-1}$ with $\der\omega=1$, and --- being a differential subfield of
$\SC$ --- it contains no nonzero solution of $\der y=niy$ for
$0\neq n\in\Z$. Adjoin a transcendental $\mathsf T$ with $\der\mathsf T=i\mathsf T$. Then the
differential constants of $K_{0}^{\Q}(\mathsf T)$ are exactly $\C$, it is a
Picard--Vessiot field for $\der y=iy$, and its differential Galois group is
$\gm$ over $\C$.
\end{theorem}
\begin{proof}
Let $f\in K_{0}^{\Q}(\mathsf T)^{\times}$ have $\der f=0$. Write $f=cP/Q$ with
$c\in(K_{0}^{\Q})^{\times}$ and monic coprime $P,Q\in K_{0}^{\Q}[\mathsf T]$ of
degrees $m,n$. Clearing denominators gives
\[
 \frac{\der c}{c}PQ+(\der P)Q-P(\der Q)=0,
\]
and reduction modulo $P$ and modulo $Q$ gives $P\mid\der P$ and
$Q\mid\der Q$. The derivation does not raise the degree in $\mathsf T$, and on a monic
polynomial of degree $m$ its leading coefficient is $im$. Therefore
\[
 \der P=imP,\qquad \der Q=inQ,\qquad \frac{\der c}{c}=i(n-m),
\]
and the last equation is impossible in $K_{0}^{\Q}$ unless $m=n$, whence
$c\in\C^{\times}$ by the constant-field statement. Comparing the coefficient
of $\mathsf T^{j}$ in $P$ for $j<m$ gives $\der a_{j}=i(m-j)a_{j}$, again impossible
for a nonzero $a_{j}$; so $P=\mathsf T^{m}$, similarly $Q=\mathsf T^{n}$, and coprimality with
$m=n$ forces $m=n=0$. Hence $f\in\C^{\times}$. The automorphism computation is
as before.
\end{proof}

\begin{remark}[Field-transcendental is not differentially transcendental]
\label{diff:rem:transcendence}
Since $K_{0}^{\Q}$ is already algebraically closed, $\mathsf T$ is necessarily
transcendental as a \emph{field} element. It is nevertheless
\emph{differentially algebraic}, satisfying the nonzero differential
polynomial $\der\mathsf T-i\mathsf T=0$. The two notions must not be confused. One might try
to identify $\mathsf T$ with $\sum_{n\geq0}(i\omega)^{n}/n!$; that fails for the
support reason of \cref{diff:rem:nonsummableexp}, and choosing a larger
exponent group does not repair the defect. The differential extension is
therefore not obtained by appending a missing value to a list of permitted
exponents.
\end{remark}

\begin{theorem}[Over the rational function field $\C(X)$]
\label{diff:thm:PVrational}
Let $F=\C(X)$ with $\der X=1$, let $\mathsf T$ be transcendental over $F$, and define
a derivation on $E=F(\mathsf T)$ by $\der X=1$, $\der\mathsf T=i\mathsf T$, extended by the
polynomial and quotient rules; this is well defined because $X$ and $\mathsf T$ are
algebraically independent. Then the differential constant field of $E$ is
exactly $\C$, $E$ is a rank-one Picard--Vessiot field for $\der Y=iY$ over
$F$, its differential automorphisms over $F$ are exactly $\mathsf T\mapsto c\mathsf T$ with
$c\in\C^{\times}$, and its differential Galois group is $\gm(\C)$.
\end{theorem}
\begin{lemma}
\label{diff:lem:rationallogder}
For $r\in\C(X)^{\times}$, the identity $r'/r=c$ with $c\in\C^{\times}$ is
impossible.
\end{lemma}
\begin{proof}
Write $r=p/q$ with nonzero complex polynomials of degrees $d,e$. Expansion at
infinity gives
\[
 \frac{r'}{r}=\frac{p'}{p}-\frac{q'}{q}=\frac{d-e}{X}+O(X^{-2}),
\]
an identity of formal Laurent expansions in $1/X$ whose constant term is zero,
so it cannot equal a nonzero constant. Equivalently, compare degrees after
clearing denominators in $p'q-pq'=cpq$.
\end{proof}
\begin{proof}[Proof of \cref{diff:thm:PVrational}]
Expand $R\in E$ at $\mathsf T=0$ as an injective formal Laurent series
$R=\sum_{n\geq n_{0}}r_{n}(X)T^{n}\in\C(X)((\mathsf T))$ --- a formal expansion, with
no analytic convergence asserted --- on which the derivation acts by
\[
 \der\!\left(\sum_{n}r_{n}T^{n}\right)
 =\sum_{n}\bigl(r_{n}'+in\,r_{n}\bigr)T^{n},
\]
compatible with the rational subfield. If $\der R=0$ then $r_{n}'+inr_{n}=0$
for all $n$; for $n\neq0$ \cref{diff:lem:rationallogder} forces $r_{n}=0$, and
for $n=0$ the equation $r_{0}'=0$ gives $r_{0}\in\C$. Injectivity of the
expansion gives $R\in\C$. The automorphism statement is as before.
\end{proof}

\begin{corollary}[No differential embedding]
\label{diff:cor:nonembedding}
$F=\C(X)$ embeds differentially into $\SC$ by $X\mapsto\omega$, since
$\der\omega=1$. That embedding does \emph{not} extend to $E=\C(X)(\mathsf T)$: the
image of $\mathsf T$ would be a nonzero solution of $\der Y=iY$, contradicting
\cref{diff:cor:oscillator}. The same argument applies to the extension
$\C(\omega)(\mathcal E)$ with $\der\omega=1$ and $\der\mathcal E=i\mathcal E$,
defined by the polynomial product rule and then the quotient rule: there is no
differential embedding of it into $\SC$ fixing $\C(\omega)$, and this
conclusion requires no assertion whatever about that extension's full constant
field. Likewise no differential embedding of $\SC(\mathsf T)$ into $\SC$ over $\SC$
exists, and $\mathsf T$ cannot be algebraic over $\SC$.
\end{corollary}

So a field can be algebraically closed, contain the required base
coefficients, and possess additive primitives, and still fail to contain this
simplest oscillatory differential extension. The generator $\mathsf T$ is
transcendental as a field element, but \emph{differentially algebraic}, since
it satisfies $\der\mathsf T=i\mathsf T$. It is not an algebraic root
missing from $\No[i]$, and no change of normal-form encoding inside the same
field can make it exist. A symbolic constructor for such an extension is mathematically
reasonable, but it must announce a change of ambient differential field:
calling its new element a surcomplex number would be a mathematical error, not
a harmless choice of representation.

\begin{table}[htbp]
\centering\small
\caption{Five Picard--Vessiot constructions for $\der y=iy$ (or $\der y=ay$).}
\label{diff:tab:PV}
\begin{tabularx}{\textwidth}{@{}L{0.21\textwidth}L{0.15\textwidth}YL{0.13\textwidth}@{}}
\toprule
Base field & Size & Route to ``no new constants'' & Reference \\
\midrule
$\SC=\No[i]$ &
proper class &
Coprimality, then $\Phi$ on coefficient ratios, using torsion-freeness of
$\Pcl$ &
\cref{diff:thm:PVone} \\
$\SC=\No[i]$ &
proper class &
Partial fractions: distinct simple poles cannot cancel; plus differential
simplicity of $\SC[\mathsf T,\mathsf T^{-1}]$ &
\cref{diff:thm:extension} \\
$L=F[i]$, $F$ localized &
set &
Formal Laurent expansion at $\mathsf T=0$, coefficientwise, plus
\cref{diff:prop:coarse} &
\cref{diff:thm:PVlocal} \\
$\C((t^{\Q}))$ &
set &
Degree and divisibility in $K_{0}^{\Q}[\mathsf T]$, leading coefficient $im$ &
\cref{diff:thm:PVhahn} \\
$\C(X)$, $\der X=1$ &
set &
Formal Laurent expansion at $\mathsf T=0$ plus
\cref{diff:lem:rationallogder} &
\cref{diff:thm:PVrational} \\
\bottomrule
\end{tabularx}
\end{table}

\subsection{The recovered oscillator, its real form, and what it costs}
\label{diff:sub:realform}

Inside any of the extensions above, with $\der\mathsf T=i\mathsf T$, put
\begin{equation}\label{diff:eq:csU}
 \mathsf C=\frac{\mathsf T+\mathsf T^{-1}}{2},\qquad
 \mathsf S=\frac{\mathsf T-\mathsf T^{-1}}{2i}.
\end{equation}

\begin{proposition}[A genuine oscillator]
\label{diff:prop:oscillatorExtension}
The elements \cref{diff:eq:csU} satisfy
\[
 \der\mathsf C=-\mathsf S,\qquad
 \der\mathsf S=\mathsf C,\qquad
 \mathsf C^{2}+\mathsf S^{2}=1 .
\]
They are linearly independent over $\C$, and every solution of
$\der^{2}y+y=0$ in the extension is a $\C$-linear combination of them.
\end{proposition}
\begin{proof}
$\der(\mathsf T^{-1})=-i\mathsf T^{-1}$ gives the differential identities, and direct
expansion the quadratic one. $\mathsf T$ and $\mathsf T^{-1}$ are $\C$-independent, since a
nontrivial relation would make the transcendental $\mathsf T$ algebraic, and the linear
change \cref{diff:eq:csU} is invertible. For the spanning statement,
coprimeness of $\der-i$ and $\der+i$ decomposes the oscillator kernel as in
\cref{diff:thm:constant}; any solution of $\der y=iy$ is a constant multiple of
$\mathsf T$ and any solution of $\der y=-iy$ is a constant multiple of $\mathsf T^{-1}$,
because the quotient has derivative zero and the constants are exactly $\C$.
\end{proof}

These are genuine oscillator elements in a field that actually contains them.
They are \emph{not} values of surcomplex sine and cosine at $\omega$:
identifying $\mathsf T$ with an element of $\No[i]$ would contradict
\cref{diff:prop:coarse}. Calling $\mathsf T$ an exponential is harmless only when its
status as a new differential element is retained.

For the real form, extend conjugation by
$\overline{\mathsf T}=\mathsf T^{-1}$. This is an involution
commuting with the derivation, since
$\overline{\der\mathsf T}=\overline{i\mathsf T}
=-i\mathsf T^{-1}=\der(\mathsf T^{-1})$, and it fixes
$\mathsf C$ and $\mathsf S$. Two equivalent real presentations occur in the
sources. First, over a real base $F$ one may take a transcendental $\mathsf T_{\R}$
with
\begin{equation}\label{diff:eq:realT}
 \der \mathsf T_{\R}=\frac{1+\mathsf T_{\R}^{2}}{2},
 \qquad
 \mathsf C=\frac{1-\mathsf T_{\R}^{2}}{1+\mathsf T_{\R}^{2}},
 \qquad
 \mathsf S=\frac{2\mathsf T_{\R}}{1+\mathsf T_{\R}^{2}},
\end{equation}
the three identities following by the quotient rule, with
$\mathsf T=(1+i\mathsf T_{\R})/(1-i\mathsf T_{\R})$; the real rational function field then has
constant field $\R$. Second, and equivalently, the conjugation-fixed Cayley
coordinate
\begin{equation}\label{diff:eq:cayley}
 \mathsf v=\frac{\mathsf T-1}{i(\mathsf T+1)},
 \qquad
 \mathsf T=\frac{1+i\mathsf v}{1-i\mathsf v},
 \qquad
 \mathsf C=\frac{1-\mathsf v^{2}}{1+\mathsf v^{2}},
 \qquad
 \mathsf S=\frac{2\mathsf v}{1+\mathsf v^{2}},
 \qquad
 \der\mathsf v=\frac{1+\mathsf v^{2}}{2},
\end{equation}
is transcendental over $F$, so $L(\mathsf T)=F(\mathsf v)[i]$ and the fixed field of the
involution is $F(\mathsf v)$: an ordinary rational function field over an ordered
field, which admits orderings extending the order of $F$.

\begin{theorem}[Oscillation cannot retain the $H$-field ordering rule]
\label{diff:thm:orderobstruction}
Let $F\subseteq\No$ be an ordered differential subfield containing
$\R\cup\{\omega\}$, as for the real bases above. No ordering of
$F(\mathsf v)$ extending the given order of $F$ can satisfy
\[
 h>\R\quad\Longrightarrow\quad \der h>0
\]
with the inherited derivation and $\der\omega=1$.
\end{theorem}
\begin{proof}
In any such ordering the identity $\mathsf C^{2}+\mathsf S^{2}=1$ bounds
$\mathsf C$ and $\mathsf S$ in absolute value by $1$. The proof of
\cref{diff:lem:smallnegative} uses the displayed positivity rule,
$\der\omega=1$, the inverse rule \eqref{diff:eq:basicder} and
$\ker\der\supseteq\R$, and nothing else; all four hold in this ordered
extension, so the lemma applies there and forces
$\der\mathsf C$ and $\der\mathsf S$ to be infinitesimal relative to $\R$.
Since $\der\mathsf C=-\mathsf S$ and $\der\mathsf S=\mathsf C$, both
$\mathsf C$ and $\mathsf S$ are then infinitesimal, so
$\mathsf C^{2}+\mathsf S^{2}$ is infinitesimal and cannot equal $1$.
\end{proof}

Equivalently, in the presentation \cref{diff:eq:realT}: $F(\mathsf T_{\R})$ admits
orders extending that of $F$, but no such realization can preserve
$\der\Ofin\subseteq\mm$ with constant field $\R$, both oscillator coordinates
being bounded by $1$ and so having infinitesimal derivatives. So it is not the
bare possibility of ordering a real field that blocks oscillation. The
obstruction is the \emph{compatibility} of that ordering with the Hardy-type
differential positivity rule. One may enlarge the differential field to obtain
the oscillator, but one cannot keep all the structure that made the base a
surreal $H$-field. This says precisely what must be surrendered.

\subsection{Several phases: the differential Galois torus}
\label{diff:sub:torus}

Let $a_{1},\dots,a_{m}\in\SC$ and consider the diagonal system
$\der y_{j}=a_{j}y_{j}$. Write $\Phi_{j}=\Phi(a_{j})\in\Pcl$ and define the
\emph{phase-relation lattice}
\begin{equation}\label{diff:eq:lattice}
 \Lambda=\Bigl\{n=(n_{1},\dots,n_{m})\in\Z^{m}:
 \textstyle\sum_{j}n_{j}\Phi_{j}=0\Bigr\}.
\end{equation}
If $n\in\Lambda$ then $\sum_{j}n_{j}a_{j}$ is a logarithmic derivative in
$\SC$, so a product of integer powers of solutions can be normalized to an
element of the base field. These are \emph{integer} relations because rational
functions in solutions involve integral monomial exponents; real linear
dependence of the obstructions is not enough.

\begin{lemma}[The lattice is saturated]
\label{diff:lem:saturated}
$\Lambda$ is saturated in $\Z^{m}$: if $kn\in\Lambda$ for a nonzero integer
$k$ then $n\in\Lambda$. Hence $\Z^{m}/\Lambda$ is free abelian of rank
$d=m-\operatorname{rank}\Lambda$.
\end{lemma}
\begin{proof}
If $k\sum_{j}n_{j}\Phi_{j}=0$, torsion-freeness of the real vector class $\Pcl$
gives $\sum_{j}n_{j}\Phi_{j}=0$. So the quotient is finitely generated and
torsion-free, hence free; equivalently the elementary integer normal-form
theorem extends a basis of the saturated subgroup to a basis of $\Z^{m}$.
\end{proof}

\begin{theorem}[The diagonal phase torus]
\label{diff:thm:torus}
The diagonal system has a Picard--Vessiot extension obtained by adjoining $d$
algebraically independent solution symbols after base-field gauges and an
invertible integer monomial change of coordinates. Its differential Galois
group is isomorphic to $\gm(\C)^{d}$, and in the original coordinates its
action on the solutions is described by
\begin{equation}\label{diff:eq:torus}
 G=\Bigl\{(c_{1},\dots,c_{m})\in(\C^{\times})^{m}:
 \textstyle\prod_{j}c_{j}^{\,n_{j}}=1
 \ \text{for every }n\in\Lambda\Bigr\}.
\end{equation}
\end{theorem}
\begin{proof}
Choose an integer basis of $\Z^{m}$ whose first
$\ell=\operatorname{rank}\Lambda$ vectors form a basis of $\Lambda$
(\cref{diff:lem:saturated}). The associated unimodular matrix changes the
diagonal coefficients to integer linear combinations $a_{1}',\dots,a_{m}'$.
The first $\ell$ have zero obstruction, hence nonzero solutions
$g_{j}\in\SC$ by \cref{diff:thm:phasecriterion}. The remaining obstructions
$\Phi'_{\ell+1},\dots,\Phi'_{m}$ are $\Z$-independent, since any relation among them
would be an element of $\Lambda$ with nonzero complementary coordinates.

For those $d$ coefficients introduce algebraically independent symbols
$\mathsf T_{1},\dots,\mathsf T_{d}$ with $\der \mathsf T_{k}=a'_{\ell+k}\mathsf T_{k}$. If a nonzero constant
of $\SC(\mathsf T_{1},\dots,\mathsf T_{d})$ is written $P/Q$ with coprime multivariate
polynomials, the argument of \cref{diff:thm:PVone} gives $\der P=hP$ and
$\der Q=hQ$ with $h\in\SC$, total degree not increasing under this derivation.
For distinct nonzero terms $p_{\alpha}\mathsf T^{\alpha}$ and $p_{\beta}\mathsf T^{\beta}$ of
$P$, coefficient comparison gives
\[
 (p_{\alpha}/p_{\beta})^{\dagger}
 =\sum_{k=1}^{d}(\beta_{k}-\alpha_{k})\,a'_{\ell+k},
\]
and applying $\Phi$ contradicts integer independence unless
$\alpha=\beta$. So $P$ and $Q$ are monomials, and the same argument on their
quotient shows that a rational constant has zero exponent vector and a
coefficient in $\C$.

Combining the $g_{j}$ with the $\mathsf T_{k}$ and inverting the unimodular monomial
change supplies nonzero solutions of all the original diagonal equations,
generating the same differential field. Each independent symbol may be scaled
by an arbitrary element of $\C^{\times}$, giving $\gm(\C)^{d}$; in the
original coordinates the trivialized integer combinations must stay fixed,
which is exactly \cref{diff:eq:torus}, and those equations leave precisely $d$
independent constant multipliers.
\end{proof}

\begin{example}[Three separations]
\label{diff:ex:torusexamples}
(i) For $a_{1}=i$, $a_{2}=2i$ the obstructions are $\omega$ and $2\omega$, the
lattice is generated by $(2,-1)$, one can choose solutions with
$y_{2}=y_{1}^{2}$, and $G=\{(c,c^{2}):c\in\C^{\times}\}$ is one-dimensional.
(ii) For $a_{1}=i$, $a_{2}=\sqrt2\,i$ the obstructions $\omega$ and
$\sqrt2\,\omega$ are dependent over $\R$ but have no nonzero integer relation,
so there are two algebraically independent generators and the group is
$\gm(\C)^{2}$; the field exponential on real arguments does not supply a
canonical irrational power of a new oscillatory symbol inside a rational
Picard--Vessiot extension. (iii) For $a_{1}=i/\omega$,
$a_{2}=i/(\omega\log\omega)$ the obstructions are $\log\omega$ and
$\log\log\omega$; the first dominates the second by an infinite factor, so no
nonzero integer relation exists and the torus again has dimension two ---
smallness of both coefficients does not reduce the number of independent
missing phases.
\end{example}

The theorem treats diagonal systems completely and supplies useful information
for triangular extensions. Taken by itself it is \emph{not} a classification of
Picard--Vessiot groups for arbitrary nontriangular matrices over $\SC$. It
becomes one in combination with the matrix normal form of \cref{diff:part:matrix}:
\cref{diff:thm:matnormal} reduces an arbitrary finite system to a diagonal one
over $\SC$, and \cref{diff:cor:mattorus} then reads off its torus, identifying
the count $d$ of \cref{diff:lem:saturated} with the rational rank $r$ of the
phase lattice $\Pi\subseteq\Pcl$ generated by $\Ph(A)$.
\Cref{diff:prop:matPValgebra} adds the group-algebra presentation in rank $r$
and the real form; the differential simplicity proved there in rank $r$ is the
general form of the rank-one statement of \cref{diff:thm:extension}. For a
computational backend the lattice is additional required data: treating
$i\omega$ and $2i\omega$ as independent phases would produce the wrong Galois
group, and treating $i\omega$ and $i\sqrt2\,\omega$ as algebraically dependent
merely because their obstructions are real multiples of one another would also
be wrong.

\section{The one genuine bridge: \texorpdfstring{$\omega$}{omega}-series and composition}
\label{diff:sec:composition}

There is an established setting in which intrinsic differentiation really is
recoverable as differentiation of a represented function. It is not all of
$\No$ with an arbitrary choice of functions.

Let $\mathscr W=\R\langle\!\langle\omega\rangle\!\rangle$ denote the
$\omega$-series: the smallest subfield of $\No$ containing $\R$ and $\omega$
and closed under exponential, positive logarithm, and set-indexed summable
families. It is a \emph{proper class}, explicitly distinct from the set-sized
workspaces of \cref{diff:sec:rationalworkspace,diff:sec:localization}.
Berarducci and Mantova construct a composition $f\circ x$ for
$f\in\mathscr W$ and positive infinite $x$; their results identify
$(\der f)\circ x$ with the \emph{fine} derivative of $x\mapsto f\circ x$ and
give a compatible chain rule inside that setting, and their Taylor theorem
states that for sufficiently small $h$
\begin{equation}\label{diff:eq:omegaTaylor}
 f\circ(x+h)=\sum_{n\geq0}\frac{(\der^{n}f)\circ x}{n!}\,h^{n},
\end{equation}
the sum being support-controlled and the permitted size of $h$ depending on
$f$ and $x$ \cite[Cor.\ 7.6--7.7, Thms 7.14, 8.2]{BMGerms}.

\begin{remark}[How the two derivatives agree, and where]
\label{diff:rem:agreement}
For $f=\log\omega$ the represented function is $x\mapsto\log x$ and
$\der f=1/\omega$; composition sends the latter to $1/x$, exactly its fine
function derivative. For $f=\omega^{3}$ both descriptions give $3x^{2}$.
Equation \cref{diff:eq:omegaTaylor} explains this agreement in a much larger
class than finite elementary expressions. This is a different construction
from a formal polynomial: in $P(Z)=\omega Z$ the symbol $\omega$ is a
coefficient held fixed by $\derZ$, whereas in the $\omega$-series
$f=\omega^{2}$ both occurrences belong to the distinguished compositional
scale and $f\circ x=x^{2}$. A representation must record which interpretation
is intended before any chain rule can be chosen. The local agreement of
derivative operations does not say that every surreal scalar is an
$\omega$-series or that every primitive belongs to that class: the full-field
surjectivity clause \ref{diff:BM:surj} and the restricted composition theorem
have different domains and must be applied separately.
\end{remark}

\begin{warning}[Imported, with its hypotheses intact]
\label{diff:warn:composition}
The composition, Taylor and incompatibility results are \emph{separate
imported theorems}; they do not follow from \cref{diff:thm:BM}, and they are
not hypotheses hidden inside the phase criterion. In particular
$f\circ x$ is an additional operation, not the result of substituting into an
arbitrary normal form. Berarducci and Mantova also prove that the simplest
derivation on \emph{all} of $\No$ admits no composition compatible with the
composition and compatibility axioms of that paper
\cite[Thm 8.4]{BMGerms}. Their hypotheses matter --- the compatibility includes
a constant rule, an order condition for positive derivative, and a chain rule
at positive infinite arguments --- and the theorem is quoted here for this
particular derivation. It is \emph{not} strengthened into an assertion about
all conceivable notions of composition, or about every other choice of
derivation, and nothing here settles alternative global composition
structures. The present finite-phase obstruction is a direct
differential-algebraic argument, not a restatement or a strengthening of that
composition theorem.

The practical lesson is to retain at least three input kinds: a formal
expression with named coefficient constants, an $\omega$-series with a
distinguished compositional scale, and a general scalar with an intrinsic
derivation. They can overlap as \emph{values} while supporting different
operations, and identifying their carriers in memory does not justify
identifying their semantic interfaces.
\end{warning}

\begin{remark}[Four tempting inferences, each refuted by a witness]
\label{diff:rem:fourinferences}
The elementary embedding of the ordered differential field of
logarithmic-exponential transseries into $(\No,\der)$ \cite{ADH} places this
derivation in a powerful nonoscillatory asymptotic setting; it is a theorem
about a specified ordered differential language and must not be silently
upgraded to a statement about every simultaneously added analytic,
compositional or complex exponential operation, and it is not a theorem that
every algebraic differential equation has a solution. Four specific inferences
fail, each with an explicit witness:
\begin{enumerate}[label=(\arabic*),leftmargin=2.4em]
\item algebraic closedness does not imply surjectivity of the logarithmic
derivative --- $i$ is missing from the image (\cref{diff:thm:logimage});
\item surjectivity of $\der$ does not imply closure under exponential
integrals --- $i\omega$ is a primitive of $i$, and no compatible exponential
of it exists in $\SC$ (\cref{diff:cor:noglobalPsi});
\item differential stability of a Hahn workspace does not imply primitive
closure --- $\log\omega\notin H_{\Q}$ (\cref{diff:ex:logmissing});
\item an infinitesimal imaginary part is necessary but not sufficient for a
nonzero rank-one solution --- $i/\omega$ is excluded while $i/\omega^{2}$ is
included (\cref{diff:warn:integrated}).
\end{enumerate}
\end{remark}
\part{Regular-singular systems over real-exponent Hahn workspaces}
\label{diff:rs:part}

This part is the contribution of member~12 (\cref{diff:app:provenance}). It
treats the finite matrix equations
\[
 \der y=-t\,Ay,\qquad A=A_{0}+A_{+},\quad A_{0}\in\Mat_{n}(\C),\quad
 A_{+}\in\Mat_{n}\bigl(\KR^{>0}\bigr),
\]
whose coefficient is a real-power Hahn series of order at least one in $t$ with
an ordinary complex leading matrix, and it classifies them completely, both
over the Hahn field $\KR=\C((t^{\R}))$ and over the whole of $\SC$. Three
relations to the rest of the report fix its place.
\begin{itemize}[leftmargin=1.6em]
\item It is the \emph{matrix version of the rank-one resonant normal form}
of \cref{diff:rem:resonant,diff:thm:Hahnlogder}: in rank one it returns that
theorem (\cref{diff:rs:rem:rankone}), and in rank $n$ it replaces the single
exponent-$1$ coefficient by a residual matrix and its Levelt resonances.
\item Over $\SC$ it computes one explicit \emph{non-Hermitian} sector of the
differential phase spectrum of \cref{diff:part:matrix}:
$\Ph(-tA)=\{(\Imm\lambda_{j})\,\tau\}$ with $\tau=\log t$ and $\lambda_{j}$ the
eigenvalues of $A_{0}$ (\cref{diff:rs:prop:phasespectrum}). It does so by an
elementary route that uses neither the import of
\cref{diff:thm:matsurjectivity} nor the surjectivity clause
\ref{diff:BM:surj}. It therefore answers item~N24 and the corresponding open
continuation \emph{for this class only}: an explicit phase formula is available
here without the matrix import. The abstract classification already covers all
finite systems over $\SC$; computable formulas for broader non-Hermitian
classes remain open (\cref{diff:sub:nc:open}).
\item It supplies, for this class, an explicit splitting whose rank-one factors
carry obstructions read off from $A_{0}$ alone, together with a finite exact
procedure for finite rational supports with algebraic coefficients
(\cref{diff:rs:sub:procedure}). That is
a partial answer to item~N57 and to the ``computable filtration'' question of
\cref{diff:warn:matlimits}; it is \emph{not} an answer to item~N41, which
concerns irregular singularities of \emph{coordinate} equations.
\end{itemize}

\section{The Euler derivation, the coefficient class, and three hazards}
\label{diff:rs:sec:setting}

\subsection{The Euler derivation is a rescaled derivation}
\label{diff:rs:sub:euler}

Put
\begin{equation}\label{diff:rs:eq:tau}
 \tau:=\log t=-\log\omega,
 \qquad
 \der\tau=-t\neq0 .
\end{equation}
The rescaled derivation of \cref{diff:prop:rescale} with $s=\tau$ is
\begin{equation}\label{diff:rs:eq:euler}
 \derT f=\frac{\der f}{\der\tau}=-\omega\,\der f .
\end{equation}
It is a derivation of $\SC$ with the same constant field $\C$, it satisfies
$\derT\tau=1$, and it is \emph{not} normalized at $\omega$: $\derT\omega=-\omega$.

\begin{lemma}[The Euler derivation on real-exponent Hahn fields]
\label{diff:rs:lem:euler}
For every additive subgroup $\Gamma\subseteq\R$, the field
$K_{\Gamma}=\C((t^{\Gamma}))\subseteq\SC$ is $\derT$-stable and
\begin{equation}\label{diff:rs:eq:eulerhahn}
 \derT\Bigl(\sum_{\gamma}c_{\gamma}t^{\gamma}\Bigr)
 =\sum_{\gamma}\gamma\,c_{\gamma}t^{\gamma},
 \qquad
 \derT\tau=1 .
\end{equation}
Its constants in $K_{\Gamma}$ are exactly $\C$, and $\derT$ maps
$K_{\Gamma}^{\geq0}=\{f:f=0\text{ or }v(f)\geq0\}$ into
$K_{\Gamma}^{>0}=\{f:f=0\text{ or }v(f)>0\}$. In particular
$\der=-t\,\derT$ preserves $K_{\Gamma}$ exactly when
$\Gamma=\{0\}$ or $1\in\Gamma$, whereas $\derT$ preserves every
$K_{\Gamma}$. The nonzero-group case is \cref{diff:prop:realgroup}.
\end{lemma}
\begin{proof}
For ordinary real $a$ the real-power rule of \cref{diff:eq:basicder}, which
\cref{diff:warn:monomialmap} allows at ordinary real exponents, gives
$\der t^{a}=-a\,t^{a+1}$, that is $\derT t^{a}=a\,t^{a}$. Strong additivity
\ref{diff:BM:additive} extends this to Hahn sums; the support can only shrink,
so every $K_{\Gamma}$ is stable. Finally $\derT\tau=-\omega\der(-\log\omega)
=\omega t=1$. The coefficient $\gamma c_{\gamma}$ vanishes only at $\gamma=0$,
which gives the constants and the valuation statement.
If $\Gamma=\{0\}$, then $K_\Gamma=\C$ and $\der$ vanishes on it.
If $1\in\Gamma$, the formula $\der=-t\,\derT$ preserves $K_\Gamma$.
Conversely, if $\Gamma\ne\{0\}$ and $\der$ preserves $K_\Gamma$, choose
$0\ne\gamma\in\Gamma$. The nonzero derivative
$\der t^\gamma=-\gamma t^{\gamma+1}$ belongs to $K_\Gamma$, so
$\gamma+1\in\Gamma$ and hence $1\in\Gamma$.
\end{proof}

\begin{remark}[One operator, several names]
\label{diff:rs:rem:names}
$\derT$ is the row $\der_{s}$ of \cref{diff:tab:derivatives} at the scale
$s=\tau$. Member~12 writes it $\delta=-\omega\der$ and writes $\ell$ for
$\tau$; the companion report on tail-spans calls the same operator the Euler
derivation and records $\der|_{H_{0}}=-t$ times it \cite{RepoTailSpans}. This
report writes $\derT$ and $\tau$ only: the letter $\delta$ already occurs in
\cref{diff:tab:phases} as a source notation for the infinitesimal phase
residue, and $\ell_{1}=\log\omega=-\tau$ is the first iterated logarithm of
\cref{diff:eq:logtower}. \Cref{diff:rs:eq:eulerhahn} lets one calculate
as though $\derT=t\,d/dt$, and it is exactly the formula of
\cref{diff:prop:hahnderiv} divided by $-t$.
\end{remark}

\begin{remark}[Not a coordinate derivative, and the cross-category pair]
\label{diff:rs:rem:notcoordinate}
The symbol $t$ names the actual surreal number $\omega^{-1}$, not a moving
ordinary complex variable (\cref{diff:warn:separators}(1)).
\Cref{diff:rs:eq:eulerhahn} is an identity for the intrinsic derivation; no
differentiability in the fine order topology is assumed, and where the
verification program of member~12 uses a formal variable it checks the
algebraic identities that this embedding then realizes. The contrast with the
coordinate theory is sharp. For $\varepsilon\in\KR^{>0}$ the \emph{intrinsic}
scalar equation $\derT y=(\lambda+\varepsilon)y$ has a nonzero solution in
$\SC$ exactly when $\Imm\lambda=0$ (\cref{diff:rs:thm:selection}), while the
\emph{coordinate} equation $\Dz Y=(\lambda+\varepsilon)Y/z$ of
\cref{diff:ex:infmonodromy} has the coherent solution
$z^{\lambda}\Emm(\varepsilon\log z)$ for every $\lambda\in\C$. This is the
cross-category pair of \cref{diff:warn:crosscategory} once more. Likewise the
coordinate Frobenius theory of the companion non-Abelian-support report
concerns $d/dz$ with positive-support residue, a different operator, and is
not merged here.
\end{remark}

\begin{warning}[Orientation]
\label{diff:rs:warn:orientation}
$\derT$ reverses the sign rule \ref{diff:BM:positivity}: if $x>\R$ then
$\der x>0$ and hence $\derT x<0$. So $\derT y=ay$ with real $a>0$ has the
\emph{infinitesimal} solution $t^{a}$, whereas $\der y=\lambda y$ has the
infinitesimal solutions $c\,e^{\lambda\omega}$ only for $\lambda<0$, and the
same holds for every normalized derivation in \cref{diff:aut:part}
(\cref{diff:aut:lem:eigen}). The positivity $\lambda-\mu>0$ in the Levelt
resonances below is oriented to $\derT$. No sign statement is moved between the
two frames without the translation $\derT=-\omega\der$.
\end{warning}

\subsection{The coefficient class and the alphabet of this part}
\label{diff:rs:sub:class}

For $\Gamma\subseteq\R$ we keep the least-exponent valuation $v$ of
\cref{diff:sec:laurent}, the valuation ring $K_{\Gamma}^{\geq0}$ and its
maximal ideal $K_{\Gamma}^{>0}$ from \cref{diff:rs:lem:euler}. In the
embedding into $\SC$, their elements are finite, respectively infinitesimal.

\begin{definition}[Power-Hahn regular-singular system]
\label{diff:rs:def:class}
A \emph{power-Hahn regular-singular system} is an equation
\begin{equation}\label{diff:rs:eq:system}
 \derT y=Ay,
 \qquad A=A_{0}+A_{+},
 \quad A_{0}\in\Mat_{n}(\C),
 \quad A_{+}\in\Mat_{n}\bigl(\KR^{>0}\bigr),
\end{equation}
equivalently $\der y=-t\,Ay$. Its \emph{residual matrix} is $A_{0}$. The unknown
$y$ ranges over all of $\SC^{n}$, with no support restriction.
\end{definition}

The adjective is a definition of the coefficient class, not an assertion that
arbitrary finite surreal matrices have such an expansion;
\cref{diff:rs:prop:boundary} makes the distinction indispensable.

\begin{convention}[The alphabet of this part]
\label{diff:rs:conv:alphabet}
Throughout \cref{diff:rs:part}:
\begin{itemize}[leftmargin=1.6em]
\item $\KR=\C((t^{\R}))$, and $K_{\Gamma}=\C((t^{\Gamma}))$ for an
\emph{arbitrary} additive subgroup $\Gamma\subseteq\R$. The standing
hypothesis $1\in\Gamma$ of \cref{diff:sec:laurent} is dropped in this part,
because $\derT$ preserves every $K_{\Gamma}$ (\cref{diff:rs:lem:euler});
contrast \cref{diff:ex:sqrt2}.
\item $\Eig(A_{0})$ is the set of ordinary complex eigenvalues of $A_{0}$,
written $\lambda=a_{\lambda}+ib_{\lambda}$. The word ``spectrum'' is not used
for it (\cref{diff:warn:spectrum}).
\item $G_{0}\in\GL_{n}(\C)$ is a constant gauge bringing $A_{0}$ to a Jordan
form $A_{0}^{\circ}=G_{0}^{-1}A_{0}G_{0}$, with generalized eigenspaces
$V=\C^{n}=\bigoplus_{\lambda}V_{\lambda}$, $d_{\lambda}=\dim V_{\lambda}$, and
$A_{0}^{\circ}|_{V_{\lambda}}=\lambda I+N^{\circ}_{\lambda}$ with
$N^{\circ}_{\lambda}$ nilpotent; a single $N^{\circ}_{\lambda}$ may contain
several Jordan blocks. The source writes $J$ for $A_{0}^{\circ}$ and
$J_{\lambda}$ for $N^{\circ}_{\lambda}$; this report uses no letter $J$
(\cref{diff:conv:noJ}).
\item $S\subset\R_{>0}$ is the positive support of $G_{0}^{-1}AG_{0}-A_{0}^{\circ}$
(a constant change of basis redistributes coefficients among entries but does
not change the union of supports), $S^{*}$ is its additive monoid as in
\cref{diff:lem:neumann}, and $S^{*}_{>0}=S^{*}\setminus\{0\}$. The source
writes $M=\langle S\rangle_{\N}$.
\item $\Lev(A_{0})=\{\lambda-\mu\in\R_{>0}:\lambda,\mu\in\Eig(A_{0})\}$ is the
finite set of \emph{Levelt resonances}. The source writes $\mathcal R$, which is
the rotation generator $\Rot$ here.
\item $G_{\natural}$ is the normalizing Hahn gauge and $A^{\natural}$ the
normalized coefficient of \cref{diff:rs:thm:normal} (the source's $H$ and $B$:
a bare $H$ is a Hermitian matrix in \cref{diff:part:matrix}, and $B$ is an
accumulated phase throughout); $\pr_{\gamma}$ is a block projection (the
source's $\Pi_{\gamma}$; $\Pi$ is the phase lattice); $\mathcal H_{\gamma}$ and
$\mathcal N_{\lambda\mu}$ are the homological operators (the source's
$L_{\gamma}$ and $K_{\lambda\mu}$).
\item $D=\operatorname{diag}(b_{\lambda}I_{d_{\lambda}})$ and
$D_{\mathrm{re}}=\operatorname{diag}(a_{\lambda}I_{d_{\lambda}})$ are real
diagonal matrices, $t^{D_{\mathrm{re}}}=\operatorname{diag}(t^{a_{\lambda}}I_{d_{\lambda}})$
is the real-monomial shear, $A^{\flat}=iD+N$ is the constant normal form,
$F=G_{0}G_{\natural}t^{D_{\mathrm{re}}}$ and $G_{\tau}=F\exp(\tau N)$ (the
source's $T$, $R_{0}$, $C$, $PHT$ and $Q$).
\item $\hob_{A}$ is the \emph{Hahn obstruction} of \cref{diff:rs:thm:forced}
(the source's $\operatorname{ob}_{A}$); it is not the phase obstruction $\Obs$
or $\Phi$. The operator $y\mapsto\derT y-Ay$ is written out; the source's name
$\mathcal D_{A}$ for it is the first-integral derivation of
\cref{diff:sub:matresonance} here.
\end{itemize}
For a $\derT$-stable subfield $\C\subseteq E\subseteq\SC$ with
$A\in\Mat_{n}(E)$ and $G\in\GL_{n}(E)$, the substitution $y=Gz$ carries \cref{diff:rs:eq:system} to
\begin{equation}\label{diff:rs:eq:gauge}
 \derT z=A^{G}z,\qquad A^{G}=G^{-1}AG-G^{-1}\derT G .
\end{equation}
Multiplying by $-t$ gives \cref{diff:eq:matgauge} for the
$\der$-coefficient $-tA$ as an identity in $\SC$. If $t\in E$, then $E$
is stable under both derivations, and $E$-gauge equivalence of the two
systems is the same in either frame. Without this hypothesis the
$\derT$-system can be defined over $E$ while its $\der$-coefficient
lies outside $E$: for $E=K_{\Z\sqrt2}$, the equation $\derT y=y$
has coefficient $1\in E$ but its translated coefficient $-t$ is not in $E$.
A morphism from a system $B'$ to a system $A$, of possibly different sizes,
is a rectangular $X$ over the stated field with $\derT X=AX-XB'$.
Because the field contains $\C$, these morphism spaces are $\C$-vector spaces.
\end{convention}

\begin{warning}[Residual frequencies are not phases]
\label{diff:rs:warn:frequencies}
The source calls the real numbers $b=\Imm\lambda$ ``phase labels'', and speaks
of ``phase multiplicities'' and ``mixed-phase'' systems. Such a $b$ is an
ordinary real number. It is none of the senses (a)--(e) of
\cref{diff:conv:phases}, and this report calls it a \emph{residual frequency}.
The relation to sense~(e) is a translation, not an identification:
$\derT z=ibz$ is $\der z=i(-bt)z$, whose accumulated phase is
$-b\log\omega=b\tau$. So the frequency $b$ corresponds to the purely infinite
phase $b\tau\in\Pcl$, and \cref{diff:rs:prop:phasespectrum} shows that the
multiset of frequencies determines $\Ph(-tA)=\{b_{j}\tau\}$.
\end{warning}

\begin{warning}[Three senses of ``resonance'']
\label{diff:rs:warn:resonance}
With this part the word acquires a third meaning, and every occurrence is
qualified.
(1) A \emph{Levelt resonance} is a positive real difference $\lambda-\mu$ of
two residual eigenvalues; the set $\Lev(A_{0})$ is used only in this part.
(2) \emph{Phase resonance} is an integer relation among the phases of a finite
system, and the resonance algebra of \cref{diff:thm:matresonance} is its
invariant ring.
(3) The \emph{resonant coefficient} is the single coefficient $[t^{1}]f$ of
\cref{diff:thm:Hahnintegral}, and the resonant normal form of
\cref{diff:rem:resonant} is its logarithmic version. In $\der$-form the residual
matrix is the $t^{1}$-coefficient $-A_{0}$ of $-tA$; so sense (3) is the
rank-one residual coefficient, sense (1) concerns differences of residual
eigenvalues and is empty in rank one, and sense (2) is unrelated to both.
\end{warning}

\begin{remark}[What this part imports]
\label{diff:rs:rem:imports}
The hypotheses used in this part are the clauses \ref{diff:BM:leibniz},
\ref{diff:BM:const}, \ref{diff:BM:additive}, \ref{diff:BM:exp},
\ref{diff:BM:norm} and \ref{diff:BM:positivity} of \cref{diff:thm:BM}, the
complexification of \cref{diff:prop:complexification} (whose extension formula
and constant field do not use surjectivity), and the ordered-word count of
\cref{diff:lem:neumann}, proved in \cref{diff:app:support}. Neither
\ref{diff:BM:surj} nor the transferred splitting and surjectivity of
\cref{diff:thm:matsurjectivity} is used, and no linear-operator factorization
theorem is needed. Where results of earlier parts are cited below, the cited
steps use only these clauses; the source proves each such step directly and
the report records the agreement instead of printing it twice. Classical
regular-singular theory --- Levelt's Jordan decomposition for singular
differential operators \cite{Levelt}, its exposition for formal connections
by Kamgarpour and Weatherhog \cite{KW}, and van der Hoeven's operator methods
on generalized power series \cite{vdH} --- is the direct antecedent of the
homological reduction; the ordinary formal theory over $\C((t))$ does not by
itself give the support statements below for an arbitrary well-ordered subset
of $\R_{>0}$.
\end{remark}

\section{The Hahn--Levelt normal form and finite resonance data}
\label{diff:rs:sec:normal}

\subsection{Supports, words, and the resonance-reaching letters}
\label{diff:rs:sub:words}

\begin{lemma}[Finite words at a fixed weight, with a length bound]
\label{diff:rs:lem:words}
Let $S\subset\R_{>0}$ be well ordered. Then $S^{*}$ is well ordered, and for
each $r\in\R$ only finitely many ordered words $(s_{1},\dots,s_{k})$ in $S$
satisfy $s_{1}+\dots+s_{k}=r$. If $S\neq\varnothing$ and $\varepsilon=\min S$,
all such words have length at most $\lfloor r/\varepsilon\rfloor$.
\end{lemma}
\begin{proof}
The first two assertions are \cref{diff:lem:neumann}, proved in
\cref{diff:app:support}; the length bound follows by adding
$s_{j}\geq\varepsilon$.
\end{proof}

Here is member~12's second, Higman-free proof for real exponents. If
$S=\varnothing$, only the empty word occurs and $S^{*}=\{0\}$. Otherwise
put $\varepsilon=\min S>0$. Every sequence in $S$ has an infinite
nondecreasing subsequence, as proved in \cref{diff:sub:app:twofacts}.
Applying this extraction successively to the $k$ coordinates gives a
componentwise nondecreasing subsequence of any sequence in $S^{k}$.
If infinitely many distinct words of length $k$ had one sum, this
subsequence would be constant: increasing any coordinate strictly would
increase the sum. This contradicts distinctness. The length bound leaves
only finitely many possible $k$, proving finiteness at each weight.

For well-ordering, suppose $r_{1}>r_{2}>\cdots$ were word sums. They are
nonnegative, and all chosen representing words have length at most
$\lfloor r_{1}/\varepsilon\rfloor$. Some length occurs infinitely often;
coordinatewise extraction at that length makes the corresponding sums
nondecreasing, a contradiction. The argument uses the Archimedean
property of $\R$. At higher rank, $k\varepsilon\leq r$ need not bound
$k$: in $\Z\oplus_{\mathrm{lex}}\Z$, take $\varepsilon=(0,1)$ and
$r=(1,0)$. The general finite-word lemma still gives finitely many
representations of each weight; this numerical length estimate does not
replace it (\cref{diff:warn:summability}).

\begin{corollary}[Products, geometric inverses, and well-founded recursion]
\label{diff:rs:cor:support}
If $X$ is a finite matrix whose entries are supported in $S$, then the family
$(X^{k})_{k\geq0}$ is strongly summable and $(I-X)^{-1}=\sum_{k\geq0}X^{k}$ has
support in $S^{*}$. For each $\gamma\in S^{*}$ only finitely many pairs
$(\alpha,\beta)\in(S^{*}_{>0})^{2}$ satisfy $\alpha+\beta=\gamma$, and each
member of such a pair is strictly smaller than $\gamma$. Consequently a
coefficient recursion using only those pairs is a well-defined transfinite
recursion on $S^{*}$.
\end{corollary}
\begin{proof}
A contribution to a fixed coefficient of a matrix power is labelled by an
ordered word in $S$ together with a finite sequence of matrix indices, so
\cref{diff:rs:lem:words} makes the number of contributions finite; this is
where \emph{ordered} words are needed (\cref{diff:rem:orderedwords}). The
geometric identity follows coefficientwise by finite cancellation. The pair
assertion is \cref{diff:lem:twosupports} applied to the well-ordered set
$S^{*}_{>0}$, and positivity of the other member gives the strict inequality.
\end{proof}

\begin{definition}[Resonance-reaching support]
\label{diff:rs:def:critical}
For a finite set $R\subset\R_{>0}$ let $S_{\crit}(R)$ be the set of letters
$s\in S$ that occur in some finite ordered word in $S$ whose sum lies in $R$.
\end{definition}

\begin{lemma}[Finite dependence at finitely many weights]
\label{diff:rs:lem:critical}
$S_{\crit}(R)$ is finite. Any coefficient expression formed from input
matrices $A_{s}$ by addition, fixed linear operators and ordered multiplication
that adds weights, and having total weight in $R$, depends only on the $A_{s}$
with $s\in S_{\crit}(R)$, provided every term has positive input weight and the
expression is defined by the well-founded recursion of
\cref{diff:rs:cor:support}.
\end{lemma}
\begin{proof}
If $S=\varnothing$, there are no positive-weight inputs and
$S_{\crit}(R)=\varnothing$. Otherwise the union of the words with sums
in the finite set $R$ is finite by \cref{diff:rs:lem:words}, so their
letters form a finite set.

To expand a coefficient of weight $r$, replace each recursively defined
factor by its defining expression. Every product splits its weight into
two positive members of $S^{*}$, each at least $\varepsilon=\min S$.
Thus following either factor decreases the remaining weight by at least
$\varepsilon$. A dependency branch has at most
$\lfloor r/\varepsilon\rfloor$ such steps, and at each step there are
only finitely many predecessor pairs. The entire dependency tree is
therefore finite. Addition and fixed linear operators introduce no new
input letters or weights. At its leaves, every term is labelled by an
ordered word of input letters summing to $r$, so if $r\in R$ those
letters all belong to $S_{\crit}(R)$.
\end{proof}

\subsection{The homological operator and the normalized resonant form}
\label{diff:rs:sub:normal}

Replace $A$ by $G_{0}^{-1}AG_{0}$ and write
$A=A_{0}^{\circ}+\sum_{s\in S}A_{s}t^{s}$, with $A_{\gamma}=0$ for
$\gamma\notin S$. For $\gamma>0$ define
\begin{equation}\label{diff:rs:eq:homological}
 \mathcal H_{\gamma}(X)=\gamma X-A_{0}^{\circ}X+XA_{0}^{\circ},
 \qquad
 \mathcal N_{\lambda\mu}(X)=N^{\circ}_{\lambda}X-XN^{\circ}_{\mu}
 \quad\text{on }\Hom(V_{\mu},V_{\lambda}).
\end{equation}
On $\Hom(V_{\mu},V_{\lambda})$ one has
$\mathcal H_{\gamma}=(\gamma-\lambda+\mu)I-\mathcal N_{\lambda\mu}$ and
$\mathcal N_{\lambda\mu}^{\,d_{\lambda}+d_{\mu}-1}=0$, by expanding a power of
the two commuting operators $X\mapsto N^{\circ}_{\lambda}X$ and
$X\mapsto-XN^{\circ}_{\mu}$. If $c=\gamma-\lambda+\mu\neq0$ the inverse is the
finite operator
\begin{equation}\label{diff:rs:eq:hinverse}
 \mathcal H_{\gamma}^{-1}
 =\sum_{k=0}^{d_{\lambda}+d_{\mu}-2}c^{-k-1}\mathcal N_{\lambda\mu}^{\,k}.
\end{equation}
No bound on the ordinary complex numbers $c^{-1}$ is imposed: a Hahn series
admits arbitrary complex coefficients, and a numerical small-divisor estimate
is not a condition in this category. Let $\pr_{\gamma}$ be the projection of
$\Mat_{n}(\C)$ onto the direct sum of the \emph{whole} blocks
$\Hom(V_{\mu},V_{\lambda})$ with $\gamma=\lambda-\mu\in\R_{>0}$. The whole
block is retained, not only the kernel of $\mathcal H_{\gamma}$ on it; on the
other blocks $\mathcal H_{\gamma}$ is invertible by \cref{diff:rs:eq:hinverse}.

The distinction matters for a nonsemisimple residual matrix. Take
\[
 A_{0}^{\circ}=\begin{pmatrix}1&1&0\\0&1&0\\0&0&0\end{pmatrix},
 \qquad A=A_{0}^{\circ}+tE_{23},
\]
where $E_{23}$ has its only nonzero entry, equal to $1$, in row $2$,
column $3$. On the resonant block
$\Hom(V_{0},V_{1})\cong\C^{2}$ at weight $1$, write a column as
$u=(u_{1},u_{2})^{\mathsf T}$. Then
$\mathcal H_{1}u=(-u_{2},0)^{\mathsf T}$ and
$\Ker\mathcal H_{1}=\C(1,0)^{\mathsf T}$.
If the retained coefficient $b$ were required to lie only in this kernel,
the equation
$\mathcal H_{1}u=(0,1)^{\mathsf T}-b$ would be impossible in its
second coordinate. Retaining the whole block and setting $u=0$ gives
$b=(0,1)^{\mathsf T}$. This normalization does not claim the smallest
possible resonant representative; it specifies a unique one by fixing the
whole resonant gauge block to zero.

\begin{theorem}[Normalized resonant form]
\label{diff:rs:thm:normal}
There is a unique pair
\begin{equation}\label{diff:rs:eq:GA}
 G_{\natural}=I+\sum_{\gamma\in S^{*}_{>0}}G_{\gamma}t^{\gamma},
 \qquad
 A^{\natural}=A_{0}^{\circ}+\sum_{\gamma\in S^{*}_{>0}}A^{\natural}_{\gamma}t^{\gamma}
\end{equation}
satisfying
\begin{equation}\label{diff:rs:eq:normalidentity}
 \derT G_{\natural}=AG_{\natural}-G_{\natural}A^{\natural},
 \qquad
 \pr_{\gamma}G_{\gamma}=0,
 \qquad
 A^{\natural}_{\gamma}=\pr_{\gamma}A^{\natural}_{\gamma}.
\end{equation}
The inverse of $G_{\natural}$ has support in $S^{*}$, and
$A^{\natural}-A_{0}^{\circ}$ has finite support contained in
$\Lev(A_{0})\cap S^{*}$.
\end{theorem}
\begin{proof}
Suppose the coefficients of weight below $\gamma$ have been constructed. The
coefficient of $t^{\gamma}$ in \cref{diff:rs:eq:normalidentity} is
\begin{equation}\label{diff:rs:eq:recursion}
 \mathcal H_{\gamma}(G_{\gamma})=R_{\gamma}-A^{\natural}_{\gamma},
 \qquad
 R_{\gamma}=A_{\gamma}+\sum_{\substack{\alpha+\beta=\gamma\\ \alpha,\beta>0}}
 \bigl(A_{\alpha}G_{\beta}-G_{\alpha}A^{\natural}_{\beta}\bigr),
\end{equation}
where $\alpha,\beta$ range over $S^{*}_{>0}$. The sum is finite,
and every recursively constructed coefficient on the right has smaller
weight (\cref{diff:rs:cor:support}); $A_{\gamma}$ is given input. On resonant blocks impose $G_{\gamma}=0$, which
forces the corresponding part of $A^{\natural}_{\gamma}$ to equal $R_{\gamma}$;
on the other blocks impose $A^{\natural}_{\gamma}=0$ and invert
$\mathcal H_{\gamma}$:
\begin{equation}\label{diff:rs:eq:explicit}
 \boxed{\quad
 A^{\natural}_{\gamma}=\pr_{\gamma}R_{\gamma},
 \qquad
 G_{\gamma}=\mathcal H_{\gamma}^{-1}(I-\pr_{\gamma})R_{\gamma},
 \quad}
\end{equation}
the inverse acting on the complementary block space only. Transfinite
recursion on $S^{*}$ defines all coefficients; their supports lie in the
well-ordered $S^{*}$, so they are Hahn matrices, and the same equations prove
the identity coefficient by coefficient and uniqueness under the stated
normalization. Since $G_{\natural}-I$ has strictly positive support in $S^{*}$,
the geometric series for its inverse is strongly summable and stays in $S^{*}$.
Finally $\pr_{\gamma}=0$ unless $\gamma\in\Lev(A_{0})$, a finite set.
If $S=\varnothing$, the construction is simply
$G_{\natural}=I$ and $A^{\natural}=A_{0}^{\circ}$.
\end{proof}

Already at this stage, if $S\subseteq\Gamma$ for an additive subgroup
$\Gamma\subseteq\R$, then $S^{*}\subseteq\Gamma$ and the gauge, its
inverse and the normalized coefficient all lie over $K_{\Gamma}$.
There is no need for $1\in\Gamma$ in the Euler frame. The later
real-monomial shear can require additional exponents, as recorded in
\cref{diff:rs:thm:workspace,diff:rs:rem:notbackwards}.

\begin{remark}[Where the difficulty lies]
\label{diff:rs:rem:difficulty}
Solving a formal Sylvester equation at each exponent would not prove
\cref{diff:rs:thm:normal}. One must also show that the recursion has only
finite coefficient contributions, produces a well-ordered support and gives an
invertible Hahn matrix; those assertions are what permit a support that
accumulates below a finite exponent. Conversely there is no analytic norm
estimate to prove: arbitrarily large coefficients do not invalidate a Hahn
series.
\end{remark}

\subsection{Finite resonance dependence, not merely high-tail invariance}
\label{diff:rs:sub:finite}

\begin{theorem}[Finite resonance theorem]
\label{diff:rs:thm:finite}
Fix $A_{0}^{\circ}$, its generalized-eigenspace decomposition, and a
well-ordered positive support set $S$. The matrices $A^{\natural}_{\gamma}$ of
\cref{diff:rs:thm:normal} depend only on the input coefficient matrices at the
finite set $S_{\crit}=S_{\crit}(\Lev(A_{0}))$. Every entry of every
$A^{\natural}_{\gamma}$ is a polynomial in those input entries, with
coefficients obtained from the nilpotent blocks $N^{\circ}_{\lambda}$ and the
inverses of the nonzero numbers $\gamma'-\lambda+\mu$ occurring in the finite
dependency calculation. If $S\neq\varnothing$ and $\varepsilon=\min S$, the
polynomial for $A^{\natural}_{r}$, $r\in\Lev(A_{0})$, has degree at most
$\lfloor r/\varepsilon\rfloor$. In particular, replacing $A$ by
\begin{equation}\label{diff:rs:eq:skeleton}
 A^{\crit}=A_{0}^{\circ}+\sum_{s\in S_{\crit}}A_{s}t^{s}
\end{equation}
does not change $A^{\natural}$.
\end{theorem}
\begin{proof}
When $S=\varnothing$, both systems in the assertion are
$A_{0}^{\circ}$ and there is no positive coefficient to compute.
Otherwise put $\varepsilon=\min S$ and give $A_{s}$ the input weight $s$.
\Cref{diff:rs:eq:recursion,diff:rs:eq:explicit} build $G_{\gamma}$ and
$A^{\natural}_{\gamma}$ by addition, weight-preserving fixed linear operators
and weight-adding products, so each positive-weight coefficient is a
polynomial in input entries all of whose terms have total input weight
$\gamma$; finiteness and termination are
\cref{diff:rs:lem:words,diff:rs:lem:critical}, and each term has at most
$\lfloor\gamma/\varepsilon\rfloor$ input factors. The only possibly nonzero
positive coefficients of $A^{\natural}$ sit at $\gamma\in\Lev(A_{0})$, where all
letters belong to $S_{\crit}$ by \cref{diff:rs:lem:critical}; setting every
other input coefficient to zero leaves these polynomials unchanged, and
uniqueness in \cref{diff:rs:thm:normal} gives the same $A^{\natural}$ for
\cref{diff:rs:eq:skeleton}.
\end{proof}

\begin{corollary}[Gauge-equivalent finite coefficient replacement]
\label{diff:rs:cor:replacement}
The system $A$ is $\KR$-gauge equivalent to the finite-support system
$A^{\crit}$. More generally, two systems with the same residual matrix, a
common well-ordered positive support container $S$, and the same coefficients
on $S_{\crit}(\Lev(A_{0}))$ are $\KR$-gauge equivalent. If
$\Lev(A_{0})\cap S^{*}=\varnothing$ then $A^{\natural}=A_{0}^{\circ}$.
\end{corollary}
\begin{proof}
If $G_{\natural}$ and $G_{\natural}^{\crit}$ normalize the two systems to the
same $A^{\natural}$, put $X=G_{\natural}(G_{\natural}^{\crit})^{-1}$.
Differentiating the inverse gives
\[
 \derT X=AX-XA^{\crit},
\]
since the two terms containing $A^{\natural}$ cancel. Thus $y=Xz$
carries the $A^{\crit}$-system to the $A$-system. These matrices are in
the common residual Jordan basis fixed above; conjugation by $G_{0}$
returns the statement to the original basis. The general assertion has
the same proof. The gauge also lies in $K_{\Gamma}$ whenever the
common support container lies in $\Gamma$. If no Levelt resonance lies
in $S^{*}$, every projection used in the recursion vanishes.
\end{proof}

\begin{corollary}[A useful but weaker cutoff statement]
\label{diff:rs:cor:cutoff}
If $\Lev(A_{0})\neq\varnothing$, changing or adding coefficients at exponents
strictly greater than $\max\Lev(A_{0})$ cannot change $A^{\natural}$, provided
the union support stays well ordered and positive. If
$\Lev(A_{0})=\varnothing$, every positive perturbation normalizes to
$A_{0}^{\circ}$.
\end{corollary}
\begin{proof}
A letter larger than $\max\Lev(A_{0})$ cannot occur in a positive word whose
sum lies in $\Lev(A_{0})$; apply \cref{diff:rs:thm:finite} to the union
support.
\end{proof}

\Cref{diff:rs:thm:finite} is stronger than \cref{diff:rs:cor:cutoff}. A
support such as $\{1-1/j:j\geq2\}$ has infinitely many members below $1$,
yet only finitely many of them can occur in exact word sums reaching a fixed
finite list of resonances. The critical set is sufficient, not necessarily
minimal: matrix products can vanish and word contributions can cancel, and the
definition deliberately does not predict those coefficient-dependent
cancellations from the support alone (\cref{diff:rs:ex:accumulation}).

\section{Spectral selection and the two classifications}
\label{diff:rs:sec:selection}

\subsection{The real-monomial shear}
\label{diff:rs:sub:shear}

Every real power belongs to $\KR$, so $t^{D_{\mathrm{re}}}\in\GL_{n}(\KR)$ and
$\derT t^{D_{\mathrm{re}}}=D_{\mathrm{re}}\,t^{D_{\mathrm{re}}}$ by
\cref{diff:rs:lem:euler}.

\begin{theorem}[Constant frequency-and-nilpotent form]
\label{diff:rs:thm:constantform}
For the original system, $F=G_{0}G_{\natural}t^{D_{\mathrm{re}}}\in\GL_{n}(\KR)$
satisfies
\begin{equation}\label{diff:rs:eq:Fidentity}
 \derT F=AF-FA^{\flat},
 \qquad A^{\flat}=iD+N,
\end{equation}
where $N\in\Mat_{n}(\C)$ is nilpotent, $N^{n}=0$, and $[D,N]=0$. The diagonal
generalized-eigenspace blocks of $N$ are the $N^{\circ}_{\lambda}$; its other
possibly nonzero blocks are exactly the constant matrices obtained from the
resonant $A^{\natural}_{\gamma}$ by the shear.
\end{theorem}
\begin{proof}
After $G_{0}$ and $G_{\natural}$ the coefficient is $A^{\natural}$, and the
shear gives
$A^{\flat}=t^{-D_{\mathrm{re}}}A^{\natural}t^{D_{\mathrm{re}}}-D_{\mathrm{re}}$.
The diagonal block $\lambda I+N^{\circ}_{\lambda}$ becomes
$ib_{\lambda}I+N^{\circ}_{\lambda}$. A positive coefficient block from
$V_{\mu}$ to $V_{\lambda}$ is nonzero only if
$\gamma=\lambda-\mu\in\R_{>0}$; then $b_{\lambda}=b_{\mu}$ and
$\gamma=a_{\lambda}-a_{\mu}$, so conjugation by the shear multiplies its
monomial by $t^{-a_{\lambda}+a_{\mu}}$ and makes it constant. Every remaining
coefficient is therefore constant and connects only spaces with the same
frequency, which proves $[D,N]=0$. Within one frequency space, order the
distinct values $a_{\lambda}$ increasingly; every off-diagonal block of $N$
maps a smaller value to a larger one. Thus $N$ is block triangular
with nilpotent diagonal blocks. Its characteristic polynomial is the
product of their characteristic polynomials, hence is $X^{n}$;
Cayley--Hamilton gives $N^{n}=0$. The gauge identity follows from the
product rule.
\end{proof}

For a fixed residual Jordan basis, $N$ is determined polynomially by the
finite resonance data of \cref{diff:rs:thm:finite}. Its block similarity
invariants do not depend on the choices made, by \cref{diff:rs:thm:Kclass}
below; the displayed \emph{entries} of $N$ are not asserted to be canonical
across choices of Jordan basis.

\subsection{Euler modes in the ambient field}
\label{diff:rs:sub:eulermodes}

This is the point at which the ambient field, not an abstract Hahn field,
becomes essential: it excludes solutions with completely arbitrary surreal
supports.

\begin{lemma}[Euler modes]
\label{diff:rs:lem:eulermode}
\begin{enumerate}[label=\textup{(\roman*)},leftmargin=2.4em]
\item If $x\in\No$ is finite, then $\derT x$ is infinitesimal.
\item For $b\in\R\setminus\{0\}$ the equation $\derT z=ibz$ has only $z=0$ in
$\SC$; for $b=0$ its solutions are exactly $\C$.
\item For $\lambda=a+ib\in\C$, the solutions in $\SC$ of $\derT z=\lambda z$
are $\C\,t^{a}$ if $b=0$, and $0$ if $b\neq0$.
\end{enumerate}
\end{lemma}
\begin{proof}
(i) For every positive real $r$ both $\log\omega\pm rx$ are positive infinite,
so \ref{diff:BM:positivity} gives $t\pm r\der x>0$; hence
$\abs{\omega\der x}<1/r$ for every $r$. This strengthens
$\der\Ofin\subseteq\mm$ of \cref{diff:lem:smallnegative} to
$\omega\der\Ofin\subseteq\mm$, and it is the forward half of the cut
description \cref{diff:eq:Icut} with $f=r^{-1}\log\omega$.
(ii) This is the second proof of \cref{diff:cor:oscillator} with $\derT$ in
place of $\der$ and (i) in place of \cref{diff:lem:smallnegative}: since
$\derT$ commutes with conjugation, $\derT(z\bar z)=0$, so $z\bar z$ is a
nonnegative ordinary real constant. If $z\neq0$, this constant is
positive, so the real and imaginary parts $u,v$ are finite. By (i),
$\derT u=-bv$ and $\derT v=bu$ are infinitesimal. Since $b$ is a
nonzero ordinary real, both $u$ and $v$ are infinitesimal, making
$u^{2}+v^{2}$ infinitesimal, contrary to its positive constant value.
For $b=0$, the assertion is the constant-field statement for $\derT$. Equivalently $\der z=i(-bt)z$, whose accumulated phase
$b\tau$ is infinite for $b\neq0$; this is the case $p=1$ of
\cref{diff:cor:power} scaled by $-b$, but the route above needs no primitive
and no \ref{diff:BM:surj}.
(iii) Write $z=t^{a}w$; then $\derT w=ibw$, and apply (ii).
\end{proof}

Member~12 states that this scalar obstruction is not new: it is the
Euler-scale specialization of the finite-primitive criterion of
\cref{diff:part:scalar}, proved again only so that the ambient exclusion below
rests on the published derivation clauses alone.

\subsection{Spectral selection over the whole surcomplex field}
\label{diff:rs:sub:selection}

\begin{theorem}[Spectral selection and explicit solution envelope]
\label{diff:rs:thm:selection}
For a power-Hahn regular-singular system \cref{diff:rs:eq:system} put
\[
 m_{\R}(A_{0})=\sum_{\lambda\in\Eig(A_{0})\cap\R}\operatorname{algmult}_{A_{0}}(\lambda).
\]
The explicitly constructed $G_{\tau}=F\exp(\tau N)\in\GL_{n}(\KR[\tau])$ and
the real diagonal matrix $D$, whose entries are the imaginary parts of the
eigenvalues of $A_{0}$ with algebraic multiplicity, satisfy
\begin{equation}\label{diff:rs:eq:Gtau}
 \derT G_{\tau}=AG_{\tau}-G_{\tau}\,iD .
\end{equation}
Every surcomplex solution, without any support restriction on the unknown, is
uniquely of the form
\begin{equation}\label{diff:rs:eq:allsol}
 y=G_{\tau}c,\qquad c\in\Ker D\subseteq\C^{n}.
\end{equation}
Hence $\dim_{\C}\Sol_{\SC}(A)=m_{\R}(A_{0})$, a fundamental matrix exists in
$\GL_{n}(\SC)$ if and only if every eigenvalue of $A_{0}$ is real, and every
solution already lies in $\KR[\tau]^{n}$.
\end{theorem}
\begin{proof}
Take $F,D,N$ from \cref{diff:rs:thm:constantform}. Since $N^{n}=0$,
$\exp(\tau N)=\sum_{k=0}^{n-1}\tau^{k}N^{k}/k!$ is a finite polynomial matrix
in the actual surreal $\tau$, with inverse $\exp(-\tau N)$, with
$\derT\exp(\tau N)=N\exp(\tau N)$, and commuting with $D$; no complex exponential
at an infinite imaginary argument is used. The product rule gives
$\derT G_{\tau}=AF\exp(\tau N)-F(iD+N)\exp(\tau N)+FN\exp(\tau N)
=AG_{\tau}-G_{\tau}iD$. For $y\in\SC^{n}$ the substitution $y=G_{\tau}w$ turns
the system into $\derT w=iDw$, and \cref{diff:rs:lem:eulermode}(ii) kills the
coordinates with nonzero frequency and leaves the others arbitrary constants.
The columns of $G_{\tau}$ indexed by zero frequencies are independent over
$\SC$, hence over $\C$, and there are $m_{\R}(A_{0})$ of them. If that number is
$n$, $G_{\tau}$ is a fundamental matrix; otherwise all solution columns lie in a
proper $\SC$-subspace and no invertible solution matrix exists.
\end{proof}

\begin{corollary}[Perturbation invariance in the specified sector]
\label{diff:rs:cor:invariance}
The dimension of the full surcomplex solution space is unchanged by arbitrary
positive-order power-Hahn perturbations of the coefficient, even when the
perturbation does not commute with $A_{0}$ or changes the resonant nilpotent
matrix. Two systems with the same residual matrix are $\SC$-gauge equivalent.
\end{corollary}
\begin{proof}
$D$ depends only on the residual eigenvalues, and the gauges $G_{\tau}$ of the
two systems reduce both to the same $iD$.
\end{proof}

\begin{corollary}[Real systems]
\label{diff:rs:cor:real}
If $A$ has entries in $\R((t^{\R}))$, then $\dim_{\R}\Sol_{\No}(A)=m_{\R}(A_{0})$,
a real fundamental matrix exists if and only if every residual eigenvalue is
real, and all real solutions lie in $\R((t^{\R}))[\tau]^{n}$.
\end{corollary}
\begin{proof}
If $y=u+iv$ solves the real system, so do $u$ and $v$; the complex solution
space is the complexification of the real one. The envelope follows by taking
real and imaginary parts of the $\KR[\tau]$ representation. When every
residual eigenvalue is real, choose $G_{0}$ real and run every recursion over
the real coefficient field.
\end{proof}

\begin{remark}[The solution collection is a set]
\label{diff:rs:rem:setsized}
Before \cref{diff:rs:thm:selection} one may regard $\Sol_{\SC}(A)$ as a class of
tuples. \Cref{diff:rs:eq:allsol} identifies it with the image of the ordinary
set $\C^{m_{\R}(A_{0})}$, so it is a set and its dimension is an ordinary finite
vector-space dimension; no dimension theory for a proper-class vector space is
invoked. The fields $\KR$ and $\KR(\tau)$ are sets.
\end{remark}

\subsection{Classification over the Hahn field and over the whole field}
\label{diff:rs:sub:classification}

Decompose the constant form by residual frequency,
$V=\bigoplus_{b}V_{b}$ with $A^{\flat}|_{V_{b}}=ibI+N_{b}$, $N_{b}$ nilpotent,
$b$ ranging over the finite set $\Imm\Eig(A_{0})$, and $m_{b}=\dim V_{b}$.

\begin{theorem}[Classification over the power-Hahn field]
\label{diff:rs:thm:Kclass}
Let $A^{\flat}=iD+N$ and $A'^{\flat}=iD'+N'$ be two constant forms, possibly of
different sizes. Every $\KR$-valued solution $X$ of
$\derT X=A^{\flat}X-XA'^{\flat}$ is a constant matrix, vanishes between distinct
residual frequencies, and satisfies
\begin{equation}\label{diff:rs:eq:intertwine}
 N_{b}X_{b}=X_{b}N'_{b}
\end{equation}
on equal-frequency blocks. Consequently two power-Hahn regular-singular systems
are $\KR$-gauge equivalent exactly when, for every $b$, their nilpotent matrices
$N_{b}$ have the same Jordan block sizes: the residual frequencies with these
partitions form a complete invariant.
\end{theorem}
\begin{proof}
On a block from $V'_{c}$ to $V_{b}$ expand $X=\sum_{\gamma}X_{\gamma}t^{\gamma}$.
Coefficient comparison gives
\begin{equation}\label{diff:rs:eq:morphismcoeff}
 \bigl(\gamma-i(b-c)\bigr)X_{\gamma}=N_{b}X_{\gamma}-X_{\gamma}N'_{c}.
\end{equation}
The operator $X\mapsto N_bX-XN'_c$ has
$(m_b+m'_c-1)$-st power zero: its two summands commute, and every
term in that power contains either $N_b^{m_b}$ or $(N'_c)^{m'_c}$.
If the scalar $\gamma-i(b-c)$ is nonzero, a finite geometric inverse
therefore forces $X_{\gamma}=0$. This scalar vanishes exactly when
$\gamma=0$ and $b=c$, because $\gamma,b,c$ are real. The surviving
coefficient satisfies \cref{diff:rs:eq:intertwine}, and conversely
every such constant matrix is a morphism.

An invertible morphism is invertible on each frequency space, where
\cref{diff:rs:eq:intertwine} makes the nilpotent matrices similar;
similarity of nilpotent complex matrices is determined by Jordan block
sizes. For the original systems, transport a constant intertwiner $C$
to $FC(F')^{-1}$. Thus constancy is a statement in the reduced frames;
the original morphism need not have constant entries.
\end{proof}

\begin{corollary}[Exact morphism dimensions over $\KR$]
\label{diff:rs:cor:HomK}
If the nilpotent Jordan block sizes at frequency $b$ are
$\nu_{b,1},\dots,\nu_{b,r_{b}}$ for $A$ and $\nu'_{b,1},\dots,\nu'_{b,r'_{b}}$
for $A'$, then
\[
 \dim_{\C}\Hom_{\KR}(A',A)=\sum_{b}\sum_{j=1}^{r_{b}}\sum_{k=1}^{r'_{b}}
 \min(\nu_{b,j},\nu'_{b,k}).
\]
\end{corollary}
\begin{proof}
A map between nilpotent Jordan blocks of sizes $p,q$ is a $\C[X]$-linear map
$\C[X]/(X^{q})\to\C[X]/(X^{p})$, determined by the image of $1$.
That image must be annihilated by $X^{q}$, and the permitted subspace
has basis
\[
 X^{\max(p-q,0)},\,X^{\max(p-q,0)+1},\,\ldots,\,X^{p-1}.
\]
There are $\min(p,q)$ basis elements. Add over block pairs with the
same frequency; distinct frequencies contribute zero.
\end{proof}

\begin{theorem}[Classification over the whole surcomplex field]
\label{diff:rs:thm:SCclass}
Two power-Hahn regular-singular systems are $\SC$-gauge equivalent if and only
if the multisets
\[
 \{\Imm\lambda:\lambda\in\Eig(A_{0}),\ \text{counted with algebraic multiplicity}\}
\]
agree. More generally $\dim_{\C}\Hom_{\SC}(A',A)=\sum_{b}m_{b}m'_{b}$, and every
ambient morphism already has entries in $\KR[\tau]$.
\end{theorem}
\begin{proof}
With $G_{\tau},G'_{\tau}$ from \cref{diff:rs:thm:selection}, a morphism is
$X=G_{\tau}Z(G'_{\tau})^{-1}$ with $\derT Z=iDZ-ZiD'$, whose scalar entries solve
$\derT z=i(b-c)z$. By \cref{diff:rs:lem:eulermode}(ii) unequal frequencies give
zero and equal ones give arbitrary complex constants. Hence the
morphism space is the direct sum of the matrix spaces
$\Hom_{\C}(V'_b,V_b)$, of dimensions $m_bm'_b$.
An invertible morphism exists exactly when $m_b=m'_b$ for every $b$,
which is equality of the frequency multisets.
Both $G_{\tau}=F\exp(\tau N)$ and its inverse
$\exp(-\tau N)F^{-1}$ have entries in $\KR[\tau]$, so the same is
true of every transported morphism.
\end{proof}

\begin{remark}[What the two classifications distinguish]
\label{diff:rs:rem:twolevels}
The real parts of the residual eigenvalues are already removable over $\KR$,
because every real monomial $t^{a}$ is available. Nilpotent Levelt data are not
in general removable over $\KR$, but become removable after adjoining $\tau$.
The residual frequencies remain invariants even over $\SC$. Passing from $\KR$ to
$\SC$ therefore erases exactly the nilpotent data in this class, and not the
frequency multiplicities.
\end{remark}

For a two-dimensional example, take the constant Euler coefficients
\[
 N=\begin{pmatrix}0&1\\0&0\end{pmatrix}
 \quad\text{and}\quad 0_{2}.
\]
Both frequency multisets are $\{0,0\}$. Over $\KR$, their Jordan
partitions are $(2)$ and $(1,1)$, so they are not gauge equivalent.
Over $\SC$, the matrix $I+\tau N$, with inverse $I-\tau N$, satisfies
$\derT(I+\tau N)=N(I+\tau N)$ and gives an equivalence.
The endomorphisms of the first system over $\KR$ are precisely
$\left(\begin{smallmatrix}a&b\\0&a\end{smallmatrix}\right)$ with
$a,b\in\C$, a two-dimensional space. Over $\SC$ they are
$(I+\tau N)C(I-\tau N)$ for arbitrary $C\in\Mat_{2}(\C)$, a
four-dimensional space.

The normal forms also make tensor constructions explicit.
For systems of sizes $n,n'$, the tensor coefficient is
$A\otimes I_{n'}+I_n\otimes A'$, as the product rule on $y\otimes y'$
shows. The notation $A\otimes A'$ for a system in the next corollary
means this differential tensor product, not the matrix $A\otimes A'$
alone. If $A$ and $A'$ have
frequency blocks $(b,N_{b})$ and $(c,N'_{c})$, their tensor product has blocks
$(b+c,\ N_{b}\otimes I+I\otimes N'_{c})$, the nilpotent sum being nilpotent
because its two terms commute; a dual has frequencies $-b$ and nilpotents
$-N_{b}^{\mathsf T}$. These follow by differentiating a tensor product and a
pairing.

\begin{corollary}[Solution dimensions of tensor products]
\label{diff:rs:cor:tensor}
Over $\SC$,
\[
 \dim_{\C}\Sol_{\SC}(A\otimes A')=\sum_{b}m_{b}\,m'_{-b}.
\]
In particular two rank-one systems with no nonzero ambient solution, with
frequencies $b\neq0$ and $-b$, have a tensor product with a fundamental
solution.
\end{corollary}
\begin{proof}
Use the tensor gauge $G_{\tau}\otimes G'_{\tau}$. It reduces the
coefficient to $i(D\otimes I+I\otimes D')$, whose scalar frequencies
are the pairwise sums $b+c$. The Euler-mode lemma leaves exactly the
constant coordinates with $b+c=0$, of total dimension
$\sum_bm_bm'_{-b}$. This proves the formula, including the stated
rank-one case.
\end{proof}

This does not contradict the obstruction: a nonzero horizontal vector of a
tensor product need not be a tensor product of horizontal vectors of its
factors, and it is the frequency sum, not the separate existence of
solutions, that controls the answer. It agrees with the rank-one statement of
\cref{diff:cor:tensor} (obstructions add) and with the tensor rule of
\cref{diff:prop:matprojectors} (phases add pairwise), and it is proved here
without the import of \cref{diff:thm:matsurjectivity}.

\subsection{The phase spectrum of a regular-singular system}
\label{diff:rs:sub:phasespectrum}

\begin{proposition}[An explicit non-Hermitian sector of the phase spectrum]
\label{diff:rs:prop:phasespectrum}
For every power-Hahn regular-singular system, with $\lambda_{1},\dots,\lambda_{n}$
the eigenvalues of $A_{0}$ counted with algebraic multiplicity,
\begin{equation}\label{diff:rs:eq:phasespectrum}
 \Ph(-tA)=\bigl\{(\Imm\lambda_{1})\,\tau,\ \dots,\ (\Imm\lambda_{n})\,\tau\bigr\}
 \subset\Pcl,
\end{equation}
and the gauge $G_{\tau}\in\GL_{n}(\KR[\tau])$ of \cref{diff:rs:thm:selection}
realizes the normal form of \cref{diff:thm:matnormal} explicitly:
$(-tA)^{G_{\tau}}=i\operatorname{diag}\bigl(\der((\Imm\lambda_{j})\tau)\bigr)$.
Consequently $m_{0}(-tA)=m_{\R}(A_{0})$, and \cref{diff:rs:thm:SCclass} is the
completeness statement of \cref{diff:thm:matnormal} for this class.
\end{proposition}
\begin{proof}
Multiply \cref{diff:rs:eq:Gtau} by $-t$: $\der G_{\tau}=(-tA)G_{\tau}-G_{\tau}(-itD)$,
so $(-tA)^{G_{\tau}}=-itD=i\operatorname{diag}(-b_{j}t)$, and
$-b_{j}t=\der(b_{j}\tau)$ by \cref{diff:rs:eq:tau}. Each $b_{j}\tau=-b_{j}\log\omega$
lies in $\Pcl$. The multiset is independent of the gauge by the uniqueness half
of \cref{diff:thm:matnormal}, which is \cref{diff:lem:matseparation} and uses no
import.
\end{proof}

Only the \emph{existence} half of \cref{diff:thm:matnormal} rests on the
imported \cref{diff:thm:matsurjectivity}; for this class it is replaced by the
explicit $G_{\tau}$. The coefficient $-tA$ is in general neither Hermitian nor
normal, so \cref{diff:rs:prop:phasespectrum} is a genuinely non-Hermitian sector
in which the phase spectrum is computed from finite data: the residual
eigenvalues. It answers item~N24 and the open continuation ``compute the phase
spectrum for broader non-Hermitian classes'' for this class and for no other.
Member~12 states the comparison in the same terms: \cref{diff:part:matrix}
gives an abstract phase-multiset normal form and a constructive Hermitian
sector, and this part computes one concrete non-Hermitian sector by a
different, elementary route.

\begin{remark}[The eigenvalue recipe in this class --- an observation of the merge]
\label{diff:rs:rem:eigenrecipe}
The following observation is not stated by either source; it is recorded, with
its proof, because it locates \cref{diff:rs:prop:phasespectrum} relative to
\cref{diff:thm:matspectral,diff:ex:matnonnormal}. In this class the naive
recipe --- integrate the imaginary parts of the instantaneous eigenvalues of the
$\der$-coefficient and keep purely infinite parts --- returns the correct
spectrum. The characteristic polynomial of $A$ is monic with coefficients in
the valuation ring $\KR^{\geq0}$ of the algebraically closed field $\KR$, and it
reduces under $\st$ to the characteristic polynomial of $A_{0}$.
To see that every root $\mu$ is finite, suppose $v(\mu)<0$.
In a monic equation of degree $n$, the term $\mu^n$ then has
strictly smaller valuation than every lower-degree term, whose
coefficient has nonnegative valuation; a unique least term cannot
cancel. Thus every eigenvalue $\mu_j$ belongs to $\KR^{\geq0}$.
Reducing the factorization $\prod_j(X-\mu_j)$ coefficientwise gives
$\prod_j(X-\st(\mu_j))$, the characteristic polynomial of $A_0$.
Its root multiset is therefore the residual eigenvalue multiset,
including multiplicities. After renumbering,
$\mu_j-\lambda_j\in\KR^{>0}$. The eigenvalues of $-tA$ are $-t\mu_j$, and
$\Imm(-t\mu_{j})=-t\,b_{j}-t\,\Imm(\mu_{j}-\lambda_{j})$, whose second term is a
real Hahn series of valuation greater than $1$, hence in $\Acl$ by
\cref{diff:cor:valphase}. So $\Obs(\Imm(-t\mu_{j}))=\Obs(-b_{j}t)=b_{j}\tau$.
\Cref{diff:ex:matnonnormal} lies outside this class: the $\derT$-coefficient
$-\omega A$ of its system has infinite entries. The observation does not
extend the recipe beyond the class. Finding other classes with computable phase
formulas remains open; existence and completeness of the abstract normal form
are already supplied by \cref{diff:thm:matnormal}.
\end{remark}

\section{Hahn solutions, the logarithm, and forced equations}
\label{diff:rs:sec:logs}

Let $V_{0}$ be the zero-frequency space and $N_{0}$ the restriction of $N$ to
it; it may have dimension zero.

\begin{theorem}[The Hahn solution space]
\label{diff:rs:thm:Hahnsol}
$\Sol_{\KR}(A)=F\Ker N_{0}$ and $\dim_{\C}\Sol_{\KR}(A)=\dim_{\C}\Ker N_{0}$.
A fundamental matrix exists in $\GL_{n}(\KR)$ if and only if $D=0$ and $N=0$,
and
\[
 \dim_{\C}\Sol_{\SC}(A)-\dim_{\C}\Sol_{\KR}(A)=\dim_{\C}V_{0}-\dim_{\C}\Ker N_{0}
\]
counts exactly the additional solutions supplied by the logarithm.
\end{theorem}
\begin{proof}
Apply \cref{diff:rs:eq:morphismcoeff} to a column whose source is the trivial
rank-one system: only its exponent-zero, zero-frequency coefficient survives,
and $N_{0}$ must annihilate it. A fundamental matrix needs $n$ independent
columns, which happens exactly when every frequency vanishes and $N=0$; then
$F$ itself is one.
\end{proof}

\begin{lemma}[The logarithm is transcendental over $\KR$]
\label{diff:rs:lem:tautrans}
No element of $\KR$ has $\derT$-derivative $1$, and $\tau$ is transcendental
over $\KR$.
\end{lemma}
\begin{proof}
$\derT f=1$ means $\der f=-t$, and $-t\notin\der\KR$ by
\cref{diff:eq:hahnimage} (\cref{diff:ex:logmissing}); directly, the
$t^{0}$-coefficient of $\derT f$ vanishes by \cref{diff:rs:eq:eulerhahn}.

If $\tau$ were algebraic, let
$P(X)=X^d+a_{d-1}X^{d-1}+\cdots+a_0$ be its monic minimal
polynomial over $\KR$, with $d\geq1$. Differentiating $P(\tau)=0$
and using $\derT\tau=1$ gives
\[
 \bigl(P^{\derT}+P'\bigr)(\tau)=0,
\]
where $P^{\derT}$ differentiates the coefficients and $P'$ is the
formal polynomial derivative. The polynomial in parentheses has
degree at most $d-1$, so minimality forces it to vanish identically.
Its $X^{d-1}$-coefficient is $\derT a_{d-1}+d$. Hence
$\derT(-a_{d-1}/d)=1$ with $-a_{d-1}/d\in\KR$, a contradiction.
This argument uses characteristic zero and the coefficient formula,
without a trace or algebraic-closedness assertion.
\end{proof}

\begin{theorem}[Minimal logarithmic repair]
\label{diff:rs:thm:minlog}
Assume every residual eigenvalue is real, so $D=0$. If $N=0$ a fundamental
matrix already lies in $\GL_{n}(\KR)$. If $N\neq0$, let $\nu$ be its nilpotence
index. Then:
\begin{enumerate}[label=\textup{(\roman*)},leftmargin=2.4em]
\item every ambient fundamental matrix is $F\exp(\tau N)C_{0}$ with
$C_{0}\in\GL_{n}(\C)$, of polynomial degree exactly $\nu-1$ in $\tau$ over
$\KR$;
\item the field over $\KR$ generated by the entries of a fundamental matrix and
its inverse determinant is exactly $\KR(\tau)$;
\item this extension has no new constants, and its group of differential
$\KR$-automorphisms is $(\C,+)$, acting by $\tau\mapsto\tau+c$.
\end{enumerate}
So one actual logarithm repairs every finite nilpotent Levelt resonance; no
tower of logarithms is needed in this class.
\end{theorem}
\begin{proof}
(i) follows from \cref{diff:rs:eq:allsol} with $D=0$. The coefficient of
$\tau^{\nu-1}$ is $FN^{\nu-1}C_{0}/(\nu-1)!\neq0$, and transcendence
(\cref{diff:rs:lem:tautrans}) makes polynomial degree unambiguous. (ii) The
field generated contains $\exp(\tau N)=F^{-1}(\cdot)C_{0}^{-1}$; the finite
logarithm identity $\sum_{j=1}^{n-1}(-1)^{j+1}(E-I)^{j}/j=\tau N$ for
$E=\exp(\tau N)$ recovers $\tau$ from a nonzero entry of $N$.
The identity is the formal identity $\log(\exp X)=X$ modulo $X^n$,
evaluated at the nilpotent matrix $\tau N$; every displayed sum is
finite. The reverse inclusion is the fundamental-matrix formula. (iii) As a differential subfield of
$\SC$, $\KR(\tau)$ has constants $\C$. A differential automorphism sends $\tau$
to another primitive of $1$, that is to $\tau+c$ with $c\in\C$; conversely
transcendence makes each translation a field automorphism, and
$\derT(\tau+c)=1$ makes it differential.
\end{proof}

\begin{remark}[An additive extension inside $\SC$]
\label{diff:rs:rem:additivePV}
\Cref{diff:rs:thm:minlog} is the usual additive Picard--Vessiot extension,
proved directly. It contrasts with the oscillatory extensions of
\cref{diff:sec:PV}: their groups are tori $\gm(\C)^{r}$ and they admit no
differential embedding into $\SC$ (\cref{diff:cor:nonembedding}), whereas the
group here is $(\C,+)$ and the extension $\KR(\tau)$ already lies inside $\SC$.
The logarithm is missing from the small base $\KR$, not from the ambient field.
If only some frequencies vanish, the logarithmic statement applies to the
zero-frequency summand; adjoining $\tau$ does \emph{not} manufacture a
surcomplex fundamental matrix on a nonzero-frequency summand, because the
obstruction of \cref{diff:rs:lem:eulermode}(ii) is untouched.
\end{remark}

We now solve $\derT y=Ay+f$ for arbitrary $f\in\KR^{n}$, which separates an
obstruction to a Hahn solution from an obstruction to any surcomplex solution.
For $f\in\KR^{n}$ put $g=F^{-1}f$, write $g_{b}=\sum_{\gamma}g_{\gamma,b}t^{\gamma}$
on the frequency space $V_{b}$, and define the \emph{Hahn obstruction}
\begin{equation}\label{diff:rs:eq:hob}
 \hob_{A}(f)=[g_{0,0}]\in V_{0}/N_{0}V_{0}.
\end{equation}
The map depends on the chosen normalizing identification; its vanishing and the
resulting cokernel do not.

\begin{theorem}[Exact Hahn obstruction sequence]
\label{diff:rs:thm:forced}
There is an exact sequence of $\C$-vector spaces
\begin{equation}\label{diff:rs:eq:exact}
 0\longrightarrow\Ker N_{0}
 \xrightarrow{\ v\mapsto Fv\ }\KR^{n}
 \xrightarrow{\ \derT-A\ }\KR^{n}
 \xrightarrow{\ \hob_{A}\ }V_{0}/N_{0}V_{0}
 \longrightarrow0 .
\end{equation}
So $\derT-A$ has finite-dimensional kernel and cokernel of equal dimension, and
a Hahn solution of $\derT y-Ay=f$ exists if and only if $g_{0,0}\in N_{0}V_{0}$.
Every forcing $f\in\KR^{n}$ nevertheless has a solution in $\KR[\tau]^{n}$, of
$\tau$-degree at most the nilpotence index $\nu_{0}$ of $N_{0}$ ($\nu_{0}=1$ when
$N_{0}=0$). If $y_p$ is the constructed particular solution, all ambient
solutions are
\[
 y=y_p+G_\tau c,\qquad c\in\Ker D\subseteq\C^n.
\]
They all belong to $\KR[\tau]^n$ and have the same degree bound.
If $V_{0}=0$ there is a unique solution in $\KR^{n}$, and it is also
the unique solution in $\SC^{n}$.
\end{theorem}
\begin{proof}
The substitution $y=Fz$ gives $\derT(Fz)-AFz=F(\derT z-A^{\flat}z)$. For
$(\gamma,b)\neq(0,0)$ the coefficient equation has the unique solution
\begin{equation}\label{diff:rs:eq:coeffresolvent}
 z_{\gamma,b}=\bigl((\gamma-ib)I-N_{b}\bigr)^{-1}g_{\gamma,b}
 =\sum_{k=0}^{m_{b}-1}\frac{N_{b}^{k}g_{\gamma,b}}{(\gamma-ib)^{k+1}},
\end{equation}
with support inside $\supp g_b$: at each exponent it applies a
finite ordinary complex linear map to the coefficient vector and
introduces no new exponent. Taking the finite union over the frequency
blocks therefore gives a Hahn vector, even if the inverse coefficients
are unbounded as ordinary complex numbers. At
$(0,0)$ the equation is $-N_{0}z_{0,0}=g_{0,0}$, solvable exactly when
$g_{0,0}\in N_{0}V_{0}$, with homogeneous freedom $\Ker N_{0}$. This proves
exactness at the two copies of $\KR^{n}$ and on the left; surjectivity of
$\hob_{A}$ follows from $f=Fv$ with constant $v\in V_{0}$, and rank--nullity for
$N_{0}$ gives equal dimensions. For an arbitrary resonant forcing, replace the
missing coefficient solution by
\begin{equation}\label{diff:rs:eq:particular}
 p(\tau)=\sum_{k=0}^{\nu_{0}-1}\frac{\tau^{k+1}}{(k+1)!}N_{0}^{k}g_{0,0}.
\end{equation}
Indeed,
\[
 \derT p=\sum_{k=0}^{\nu_0-1}\frac{\tau^k}{k!}N_0^kg_{0,0},
 \qquad
 N_0p=\sum_{k=1}^{\nu_0}\frac{\tau^k}{k!}N_0^kg_{0,0}.
\]
The last term of the second sum vanishes and all positive-degree
terms cancel, leaving $(\derT-N_0)p=g_{0,0}$. Add $p$ to the
nonresonant Hahn solution and multiply by $F$ to obtain a particular
solution in $\KR[\tau]^n$. The difference of any two ambient
solutions is homogeneous, so \cref{diff:rs:thm:selection} gives exactly
the stated affine family. Its homogeneous term uses only $N_0$ and
has degree at most $\nu_0-1$, proving the common degree bound.
If $V_0=0$, nothing is missing and the homogeneous solution space is
zero, giving uniqueness even in $\SC^n$.
\end{proof}

The degree bound for forcing is sharp. Let $A=N$ be the nilpotent
Jordan block of size $\nu$, with $Ne_j=e_{j-1}$ for $j>1$ and
$Ne_1=0$, and force by the constant vector $e_\nu$. Formula
\eqref{diff:rs:eq:particular} gives
\[
 p(\tau)=\tau e_\nu+\frac{\tau^2}{2!}e_{\nu-1}
             +\cdots+\frac{\tau^\nu}{\nu!}e_1.
\]
Every ambient solution differs from $p$ by $\exp(\tau N)c$,
$c\in\C^\nu$, which has degree at most $\nu-1$. The leading
coefficient $e_1/\nu!$ cannot cancel, so every solution has degree
exactly $\nu$ over $\KR$. This includes $\nu=1$, where the equation
is simply $\derT y=1$.

\begin{remark}[Algebraic, not topological, Fredholm terminology; relation to item~N22]
\label{diff:rs:rem:fredholm}
The equal-dimension statement may be called an algebraic index-zero statement.
No Banach space, closed-range argument, compact operator or topology on $\SC$ is
involved, and the existence after adjoining $\tau$ is proved only for forcing
in $\KR^{n}$; nothing is inferred for arbitrary surcomplex forcing. For this
class and this forcing, \cref{diff:rs:thm:forced} gives inhomogeneous existence
\emph{without} the import of \cref{diff:thm:matsurjectivity}: the existence
statement of \cref{diff:cor:matinhom} for $\der y-(-tA)y=-tf$ is recovered
import-free when $f\in\KR^{n}$. Item~N22 is amended accordingly and not deleted.
\end{remark}

\section{Worked examples and exact resonance certificates}
\label{diff:rs:sec:examples}

Write $E_{jk}$ for a matrix unit. All examples use $\derT t^{a}=at^{a}$ and
$\derT\tau=1$.

\begin{example}[A logarithm created by two nonresonant coefficients]
\label{diff:rs:ex:interaction}
Let
\[
 A=\operatorname{diag}(2,1,0)+t^{1/2}E_{12}+t^{3/2}E_{23}.
\]
Here $\Lev(A_{0})=\{1,2\}$ and $S=\{\tfrac12,\tfrac32\}$. The first coefficient
connects eigenvalues with difference $1$ at weight $\tfrac12$, the second at
weight $\tfrac32$; neither lies in its resonant block. Their ordered product maps
the $0$-space to the $2$-space at weight $2$. The exact normalized form is
\[
 G_{\natural}=I-2t^{1/2}E_{12}+2t^{3/2}E_{23},
 \qquad
 A^{\natural}=\operatorname{diag}(2,1,0)+2t^{2}E_{13},
\]
the resonant coefficient arising from $A_{1/2}G_{3/2}=2E_{13}$; no resonant
input coefficient was present. With $t^{D_{\mathrm{re}}}=\operatorname{diag}(t^{2},t,1)$
the constant nilpotent is $N=2E_{13}$, and
\[
 Y=G_{\natural}t^{D_{\mathrm{re}}}(I+2\tau E_{13})
 =\begin{pmatrix}t^{2}&-2t^{3/2}&2\tau t^{2}\\0&t&2t^{3/2}\\0&0&1\end{pmatrix}
\]
is a fundamental matrix with $\det Y=t^{3}$; its third column differentiates to
$(2t^{2}+4\tau t^{2},\,3t^{3/2},\,0)^{\mathsf T}$, which is $A$ times the column.
The surcomplex solution space has dimension $3$, the Hahn solution space
dimension $\dim\Ker(2E_{13})=2$, and since $N^{2}=0\neq N$ one logarithm of
degree one is unavoidable. The whole input is critical; adding an arbitrary
coefficient at weight $\tfrac73$, or any positive well-ordered family entirely
above $2$, leaves $N$ unchanged, although $G_{\natural}$ changes.
\end{example}

\begin{example}[Accumulating small divisors without a logarithm]
\label{diff:rs:ex:accumulation}
Let $A=\operatorname{diag}(1,0)+fE_{12}$ with $f=\sum_{j\geq2}t^{1-1/j}$, whose
support increases to $1$ and is well ordered. Put
$h=-\sum_{j\geq2}j\,t^{1-1/j}$. Coefficientwise
$\bigl((1-\tfrac1j)-1\bigr)(-j)=1$ gives $(\derT-1)h=f$, and since $E_{12}^{2}=0$,
$G_{\natural}=I+hE_{12}$ satisfies $\derT G_{\natural}=AG_{\natural}-G_{\natural}A_{0}$.
So $A^{\natural}=A_{0}$, $N=0$, and $G_{\natural}\operatorname{diag}(t,1)$ is
already a Hahn fundamental matrix. The divisors $(1-\frac1j)-1$ tend to zero and
the coefficients of $h$ grow without bound; neither is an obstruction. Every letter is at least $\tfrac12$, so a word summing to $1$ has
length at most two. There is no letter $1$, and a length-two word must
have both letters equal to $\tfrac12$. Its matrix product $E_{12}^{2}$
vanishes, so $S_{\crit}=\{\tfrac12\}$ is a valid finite dependence set
while the minimal one here is empty.
\end{example}

\begin{example}[A nonsemisimple residual matrix]
\label{diff:rs:ex:nonsemisimple}
Take $A=E_{12}+tE_{21}$. Both residual eigenvalues are $0$ and $\Lev(A_{0})$ is
empty, so $A^{\natural}=A_{0}^{\circ}=E_{12}$ and $N=E_{12}$. In the first step
$\mathcal N(X)=E_{12}X-XE_{12}$ has $\mathcal N^{3}=0$ and
\[
 G_{1}=(I-\mathcal N)^{-1}E_{21}=E_{21}+\operatorname{diag}(1,-1)-2E_{12};
\]
the recursion supplies all later coefficients. The ambient dimension is $2$,
the Hahn dimension $1$, and the fundamental matrix needs $\tau$-degree one.
Dividing entrywise by eigenvalue differences would miss the nilpotent terms of
\cref{diff:rs:eq:hinverse}.
\end{example}

\begin{example}[A noncommuting mixed-frequency system]
\label{diff:rs:ex:mixed}
Let $A=\operatorname{diag}(0,i)+t\left(\begin{smallmatrix}1&2\\3&4\end{smallmatrix}\right)$.
The two displayed matrices do not commute, the residual matrix is semisimple,
and no difference of distinct eigenvalues is a positive real, so
$A^{\natural}=\operatorname{diag}(0,i)$ and $N=0$, with first normalizing
coefficient
\[
 G_{1}=\begin{pmatrix}1&1-i\\ \tfrac32(1+i)&4\end{pmatrix};
\]
the remaining $G_{j}$ exist in $\C[[t]]$ by the same recursion. Exactly one
column of $G_{\natural}$ gives a nonzero family of homogeneous surcomplex
solutions; the other residual mode acquires none, even from an arbitrary
surreal support. The dimension is $1$, not $2$, and
$\Ph(-tA)=\{0,\tau\}$ by \cref{diff:rs:prop:phasespectrum}.
\end{example}

\begin{example}[A forced equation with no Hahn solution]
\label{diff:rs:ex:forced}
For $N=E_{12}$ and $g=(0,1)^{\mathsf T}$, the equation $\derT z=Nz+g$ has no
solution in $\KR^{2}$, since its second row would be $\derT z_{2}=1$
(\cref{diff:rs:lem:tautrans}); equivalently $g\notin NV$. It has the particular
solution $z=(\tau^{2}/2,\ \tau)^{\mathsf T}$ in $\KR[\tau]^{2}$, since
$\derT z=(\tau,1)^{\mathsf T}=Nz+g$. The $\tau$-degree two for this forcing exceeds
the degree one needed for a homogeneous fundamental matrix, as
\cref{diff:rs:eq:particular} allows.
\end{example}

\begin{example}[A residue change can change the answer discontinuously]
\label{diff:rs:ex:discontinuity}
The scalar equation $\derT y=(a+ib)y$ has a one-dimensional solution space for
$b=0$ and a zero one for every nonzero real $b$, however small in the ordinary
metric. The invariance of \cref{diff:rs:cor:invariance} is invariance under
positive \emph{Hahn-order} perturbations, not under numerically small changes of
the residual matrix; confusing the two notions of smallness would give a false
stability claim.
\end{example}

\section{Smaller workspaces, the power-Hahn boundary, and an exact procedure}
\label{diff:rs:sec:workspaces}

\subsection{Finite monomial adjunctions}
\label{diff:rs:sub:workspace}

\begin{theorem}[Finite-monomial and one-logarithm envelope]
\label{diff:rs:thm:workspace}
Let $\Gamma\subseteq\R$ be any additive subgroup, let $A\in\Mat_{n}(K_{\Gamma}^{\geq0})$,
and let its residual eigenvalues be $\lambda=a_{\lambda}+ib_{\lambda}$. The
normalizing gauge $G_{\natural}$ of \cref{diff:rs:thm:normal} and its inverse
lie in $\Mat_{n}(K_{\Gamma})$, and the full constant normal form is obtained over
the finitely generated extension
$K_{\Gamma}(t^{a_{\lambda}}:\lambda\in\Eig(A_{0}))$. Every ambient homogeneous
solution already lies in
\begin{equation}\label{diff:rs:eq:solworkspace}
 K_{\Gamma}\bigl(t^{a}:a\in\Eig(A_{0})\cap\R\bigr)[\tau]^{n}.
\end{equation}
So finitely many actual real monomials and one logarithm suffice; no general
differential closure is needed.
\end{theorem}
\begin{proof}
The recursion only adds positive input exponents, so all coefficients of
$G_{\natural}$ and $G_{\natural}^{-1}$ have exponents in $\Gamma$, and so does the
finite $A^{\natural}$. The shear uses only the listed monomials. For an ambient
solution only the zero-frequency columns of $G_{0}G_{\natural}t^{D_{\mathrm{re}}}\exp(\tau N)$
occur; $N$ preserves frequency spaces, so these columns use only generalized
eigenspaces with real residual eigenvalues, whose shear entries are the
monomials in \cref{diff:rs:eq:solworkspace}.
\end{proof}

These adjunctions are ordinary field generation inside $\SC$: an element
of $K_\Gamma(t^{a_1},\ldots,t^{a_k})$ is a rational expression in
finitely many added monomials with coefficients in $K_\Gamma$.
The assertion does not replace this field by the entire Hahn field on
$\Gamma+\Z a_1+\cdots+\Z a_k$, or claim closure under arbitrary Hahn
sums there. The generated field is $\derT$-stable, since
$\derT t^{a_j}=a_jt^{a_j}$ and the product and quotient rules preserve
rational expressions; adjoining $\tau$, with $\derT\tau=1$, also
preserves this stability.

The monomial adjunction can be necessary. For $\Gamma=\Q$ the equation
$\derT y=\sqrt2\,y$, that is $\der y=-\sqrt2\,t\,y$, has the nonzero ambient
solutions $c\,t^{\sqrt2}$, none of which lies in $K_{\Q}$. This is the second
row of \cref{diff:tab:workspaceambient} with $\sqrt2$ replaced by $-\sqrt2$ ---
the same obstruction, an irrational exponent-$1$ coefficient in $\der$-form.
Together with \cref{diff:ex:sqrt2}, where $\der$
fails to preserve $K_{\Q\sqrt2}$, it shows why this part lets $\Gamma$ be
arbitrary: $\derT$ preserves every $K_{\Gamma}$, and the obstruction moves from
stability of the workspace to membership of the monomial $t^{a}$.

\begin{remark}[Rank one is the resonant normal form]
\label{diff:rs:rem:rankone}
For $n=1$ and $\Gamma=\Q$ there are no Levelt resonances, and
\cref{diff:rs:thm:workspace,diff:rs:thm:selection} say: $\derT y=(\lambda+g)y$
with $\lambda\in\C$ and $g\in K_{\Q}^{>0}$ has a nonzero solution in $K_{\Q}$
exactly when $\lambda\in\Q$, the solutions being $c\,G_{\natural}t^{\lambda}$ with
$G_{\natural}\in1+K_{\Q}^{>0}$. In $\der$-form the coefficient is
$-\lambda t-tg$ with $v(tg)>1$, and this is exactly \cref{diff:thm:Hahnlogder}:
the logarithmic-derivative image of $\Hring=\C((t^{\Q}))$ is
$\{-rt+g':r\in\Q,\ g'=0\text{ or }v(g')>1\}$, every member of which has
valuation at least $1$. So \cref{diff:rs:part} is the matrix version of
\cref{diff:rem:resonant}: the rank-one exponent-$1$ residue becomes the residual
matrix $-A_{0}$, its rational real part becomes the real parts $a_{\lambda}$ to be
sheared away, its imaginary part becomes the residual frequencies that no
ambient gauge removes, and the new phenomenon in rank $n$ is the finite
nilpotent Levelt data $N$.
\end{remark}

\begin{remark}[Do not transfer the simplified classification backwards]
\label{diff:rs:rem:notbackwards}
Over $K_{\Gamma}$ the real parts need not all be removable; even rank one retains
the real exponent modulo $\Gamma$. The complete invariants of
\cref{diff:rs:thm:Kclass} are for $\KR=K_{\R}$ only. Applied unchanged to $K_{\Q}$
they would identify $\derT y=\sqrt2\,y$ with the trivial equation, which is false
over $K_{\Q}$.

More precisely, two constant rank-one Euler systems with coefficients
$\lambda,\mu\in\C$ are $K_\Gamma$-gauge equivalent exactly when
$\lambda-\mu\in\Gamma$. A gauge $x\neq0$ must solve
$\derT x=(\lambda-\mu)x$. Writing
$x=\sum_{\gamma\in\Gamma}c_\gamma t^\gamma$ gives
$(\gamma-\lambda+\mu)c_\gamma=0$ at every exponent. Thus a nonzero
gauge exists only when $\lambda-\mu=\gamma\in\Gamma$, and all such
gauges are $c\,t^\gamma$ with $c\in\C^\times$. Conversely these
monomials are gauges. The invariant is therefore the imaginary part
together with the real part modulo $\Gamma$. Rank-one positive
perturbations first normalize over $K_\Gamma$, so the same criterion
applies to their residual coefficients.
\end{remark}

\subsection{The power-Hahn hypothesis cannot be dropped}
\label{diff:rs:sub:boundary}

The theorems above are not a principle that ``the standard part of every finite
surreal coefficient controls all solutions''. Logarithmically slow scales refute
that.

\begin{proposition}[Infinitesimal coefficient, zero residue, no solution]
\label{diff:rs:prop:boundary}
The scalar equation
\begin{equation}\label{diff:rs:eq:boundary}
 \derT y=\frac{i}{\log\omega}\,y
\end{equation}
has no nonzero solution in $\SC$, although its coefficient is infinitesimal with
standard part zero. The coefficient does not belong to $\KR$.
\end{proposition}
\begin{proof}
In $\der$-form \cref{diff:rs:eq:boundary} is $\der y=-i\,y/(\omega\log\omega)$,
the first equation of \cref{diff:ex:logseparator} with $b$ replaced by $-b$; its
accumulated phase $-\log\log\omega$ is infinite, and a nonzero solution would
make it finite by the polar decomposition (\cref{diff:thm:phasecriterion}, whose
direction (i)$\Rightarrow$(iii) uses no primitive). Member~12's own route: the
derivation $\der_{1}=(\log\omega)\derT=-\omega\log\omega\,\der$ satisfies
$\der_{1}x\in\mm$ for finite real $x$, by \ref{diff:BM:positivity} applied to the
positive infinite $\log\log\omega\pm rx$, and the norm-square argument of
\cref{diff:rs:lem:eulermode}(ii) applied to $\der_{1}y=iy$ excludes every nonzero
solution. Finally $\log\omega<\omega^{a}$ for every real $a>0$, a property of the
surreal exponential \cite{BM,ADH}, so $1/\log\omega$ exceeds every $t^{a}$ while
being infinitesimal. If $1/\log\omega$ were in $\KR$, its least
exponent would be some $a_0>0$. Dividing by $t^{a_0}$ would give a
positive finite element with nonzero ordinary residue, hence
$1/\log\omega<Ct^{a_0}$ for some ordinary $C>0$. But
$Ct^{a_0}<t^{a_0/2}<1/\log\omega$, a contradiction. The use of
$a_0/2$ makes the comparison independent of the leading coefficient.
\end{proof}

\begin{remark}[What the counterexample shows, and what it does not]
\label{diff:rs:rem:boundary}
\Cref{diff:rs:prop:boundary} refutes the extension of this part from positive
power-Hahn order to arbitrary infinitesimal surreal coefficients. It does
\emph{not} refute a classification based on richer logarithmic or differential
phase data, and the scalar criterion of \cref{diff:part:scalar} already predicts
the obstruction: in $\der$-form the imaginary coefficient
$-1/(\omega\log\omega)$ has the infinite primitive $-\log\log\omega$
(\cref{diff:cor:logscales}). A layered analogue for coefficient fields
containing $\log\omega$, $\log\log\omega$ and their real powers is recorded as
an open continuation in \cref{diff:sub:nc:open}. (Added with member~13: the
equation is the case $s=\ell_{1}$, $\nu=-i$ of \cref{diff:cp:cor:nonreal}, the
first-order obstruction at stage one of the log-atomic tower.)
\end{remark}

\subsection{An exact finite procedure, and what it does not compute}
\label{diff:rs:sub:procedure}

A concrete effective input class consists of finite positive rational
supports and complex algebraic matrix entries, represented exactly with
their algebraic roots specified. In this class one can compute a residual
Jordan basis and eigenvalues, decide equality and whether an eigenvalue
difference is positive real, and perform exact field arithmetic and
matrix rank calculations. These coefficient operations are part of the
input contract: finite rational support alone does not supply them for
arbitrary named complex constants. The same finite procedure applies to
another coefficient presentation only when it supplies these operations.

Fix the resulting $A_0^\circ$ and finite rational $S$. Write
$R=\Lev(A_0)$. If $R\neq\varnothing$, let $r_{\max}=\max R$ and
choose a common denominator $d$ with $S\subseteq d^{-1}\Z_{>0}$.
Every monoid exponent at most $r_{\max}$ belongs to the finite grid
\[
 \{j/d:0\leq j\leq\lfloor dr_{\max}\rfloor\}.
\]
Starting at $0$ and adding letters of $S$ while staying in this grid
enumerates the required exponents. Apply
\cref{diff:rs:eq:recursion,diff:rs:eq:explicit} in increasing order,
record the nonzero $A^\natural_r$, perform the shear and compute the
Jordan forms of the $N_b$. If $R=\varnothing$, the normalized
coefficient is already $A_0^\circ$; the nilpotent blocks of that
residual matrix must still be retained. Empty support requires no
recursion either. This determines all $\KR$-gauge invariants; the
ambient invariants need only the residual eigenvalues.

A smaller certificate keeps only the weights
\[
 W=\{\gamma\in S^*:\text{ there is }r\in R
                         \text{ with }r-\gamma\in S^*\}.
\]
These are precisely the partial sums of all words reaching a resonance,
including their empty and full prefixes. The finite-word lemma makes
$W$ finite, even for an accumulating support when those words have
been determined. If $\gamma\in W$ and
$\alpha+\beta=\gamma$ with $\alpha,\beta\in S^*_{>0}$, then both
$\alpha,\beta\in W$: for a witnessing $r$,
$r-\alpha=(r-\gamma)+\beta$ and
$r-\beta=(r-\gamma)+\alpha$ lie in $S^*$.
Thus the recursion on these weights is closed under every predecessor
pair. Moreover $S\cap W=S_{\crit}(R)$, since a letter is in $W$
exactly when it can be completed to a word reaching $R$.
The certificate consists of these critical input coefficients, the
homological inverses at positive weights in $W$, the resulting
$A^\natural_r$ and the nilpotent Jordan data, each verified by finite
identities. At a resonant weight the inverse is restricted to the
complementary block space, as in \cref{diff:rs:eq:explicit}.
For example, if $S=\{\tfrac23,1,\tfrac32\}$ and $R=\{1,2\}$, then
$W=\{0,\tfrac23,1,\tfrac43,2\}$ and
$S_{\crit}(R)=\{\tfrac23,1\}$. The coefficient at $\tfrac32$ is below
the largest resonance but cannot affect the normalized resonant data.

Three limits travel with this procedure. The full gauge $G_{\natural}$ may still
have infinite support; a finite certificate for the invariants is not a finite
listing of that gauge, whose exact specification is the support theorem and the
recursion. Already the scalar coefficient $A=t$ has no Levelt
resonances and normalizes to $0$, but its normalized gauge is
$\sum_{k\geq0}t^k/k!$, with infinite support.
For an arbitrary well-ordered support supplied only indirectly the
critical set is finite, but finding all exact word sums that hit a resonance
need not be an effective task under an unspecified representation; finitely
many black-box coefficient queries cannot certify that an unseen coefficient at
a resonant weight is absent, so \cref{diff:rs:thm:finite} is a finite-dependence
theorem and not an oracle algorithm. And an accumulating support means that
``iterate through all exponents below $\max\Lev(A_{0})$'' need not terminate,
which is exactly why the finite-word form of \cref{diff:rs:thm:finite} is
stronger than a bounded-exponent prescription; for a symbolic family such as
\cref{diff:rs:ex:accumulation} the support description must be used to find the
words.

\begin{remark}[Relation to the open effectivity questions]
\label{diff:rs:rem:effective}
Item~N57 poses the identification of restricted, support-certified input
languages on which eigenvalue primitives, their purely infinite parts and a
normalizing \emph{unitary} gauge become effectively computable, and
\cref{diff:warn:matlimits} leaves open which nontriangular differential modules
admit a filtration whose rank-one factors carry computable obstructions. For
power-Hahn regular-singular systems this part answers both \emph{in part}: the
gauge $G_{\tau}$ splits the module over $\KR[\tau]$ into rank-one factors whose
obstructions $(\Imm\lambda_{j})\tau$ are read off from the residual matrix, and
for finite rational supports with the effective coefficient operations
specified above (in particular, algebraic coefficients), the procedure
computes the $\KR$-invariants by a finite calculation. The gauge $G_{\tau}$ is not
unitary, $G_{\natural}$ is not finitely listed, and nothing is said about other
input languages; the unitary-gauge question of item~N57 and the general
filtration question stay open.
\end{remark}

\begin{table}[htbp]
\centering\small
\caption{Proof dependencies of \cref{diff:rs:part}, after member~12's audit
register.}
\label{diff:rs:tab:dependencies}
\begin{tabularx}{\textwidth}{@{}L{0.25\textwidth}YL{0.28\textwidth}@{}}
\toprule
Result & Essential inputs & What is not used \\
\midrule
Hahn normal form (\cref{diff:rs:thm:normal}) & Ordered-word lemma; Sylvester
block inverse; transfinite recursion & Norm convergence; semisimple residue \\
Finite resonance data (\cref{diff:rs:thm:finite}) & Finite words at a fixed
total weight; polynomial recursion & Finite support below a cutoff \\
Constant reduction (\cref{diff:rs:thm:constantform}) & Real powers; finite
directed resonance blocks & Analytic continuation; global complex
exponential \\
Ambient mode exclusion (\cref{diff:rs:lem:eulermode}) & \ref{diff:BM:positivity};
constants; conjugation & Prescribed support for the unknown \\
Spectral selection (\cref{diff:rs:thm:selection}) & Constant reduction; finite
nilpotent exponential; mode exclusion & Operator factorization or transferred
surjectivity \\
Hahn classification (\cref{diff:rs:thm:Kclass}) & Coefficient comparison;
scalar-plus-nilpotent inverse & Completeness of the ambient field \\
Ambient classification (\cref{diff:rs:thm:SCclass}) & Explicit logarithmic
gauge; scalar mode exclusion & The abstract normal form
\cref{diff:thm:matnormal} \\
Forced equation (\cref{diff:rs:thm:forced}) & Finite coefficient resolvents;
polynomial resonant primitive & An ambient forcing theorem \\
Minimal logarithm (\cref{diff:rs:thm:minlog}) & Minimal-polynomial
differentiation; finite matrix logarithm; constants & Algebraic closedness of the Hahn field \\
\bottomrule
\end{tabularx}
\end{table}

Several details are easy to miss and are recorded as the source records them.
The sign in $\derT=-\omega\der$ makes $\derT t=t$ and $\derT\tau=1$. In
\cref{diff:rs:eq:recursion} the second product is $G_{\alpha}A^{\natural}_{\beta}$,
in that order and with a minus sign. The retained resonant space is a whole
generalized-eigenspace $\Hom$ block. A Levelt resonance requires a
\emph{positive real} eigenvalue difference; an imaginary or negative difference
cannot meet a positive real exponent. The nilpotent exponential is a finite
polynomial. The solution count uses algebraic multiplicity, not the number of
distinct real eigenvalues. And $\Sol_{\KR}(A)$ can be smaller than
$\Sol_{\SC}(A)$ although both fields have the constant field $\C$.
\part{Autonomous first-order equations: curves, derivation independence, and abelian logarithmic derivatives}
\label{diff:aut:part}

This part is the contribution of member~11 (\cref{diff:app:provenance}). It
answers, for one precisely delimited nonlinear class, a question this report
had explicitly declined: which solutions does a fixed autonomous algebraic
first-order equation with ordinary complex coefficients have in $\SC$? The
answer is complete for equations $F(y,\Dn y)=0$ with $F\in\C[Y,Z]$, and more
generally for constant meromorphic vector fields on smooth projective complex
curves (\cref{diff:aut:thm:main}). Three features set the part apart.
\begin{itemize}[leftmargin=1.6em]
\item It is proved for \emph{every normalized derivation} $\Dn$
(\cref{diff:aut:def:normalized}), not only for the fixed $\der$ of
\cref{diff:thm:BM}, and it shows that the solution sets in this sector do not
depend on the choice (\cref{diff:aut:thm:independence}). That is a scoped
amendment to the one-derivation rule of
\cref{diff:warn:nouniqueness,diff:warn:BMclauses} and items~N2--N3, which are
amended and not deleted (\cref{diff:aut:conv:scope}).
\item It lifts, for this class, the non-claims N21 and N25 and the remark after
\cref{diff:cor:riccati} that the Riccati corollary ``is not a general nonlinear
solvability theorem''. Those statements are kept as provenance and re-scoped;
beyond the class they stand (\cref{diff:sub:nc:scalarscope}).
\item It adds a projective obstruction for $(\Dn y)^{2}=Q(y)$ and the exact
logarithmic-derivative image of a constant abelian variety, the proper-group
analogue of \cref{diff:thm:logimage}.
\end{itemize}
The part depends only on \cref{diff:part:foundations,diff:part:scalar} and on
classical curve and formal-group theory.

\section{Normalized derivations}
\label{diff:aut:sec:derivations}

\begin{definition}[Normalized surreal derivation]
\label{diff:aut:def:normalized}
A \emph{normalized surreal derivation} is an additive map $\Dn:\No\to\No$ with
\[
 \Dn(xy)=x\Dn y+y\Dn x,\qquad \Dn(e^{x})=e^{x}\Dn x,\qquad \Ker\Dn=\R,
 \qquad x>\R\ \Rightarrow\ \Dn x>0,\qquad \Dn\omega=1,
\]
and \emph{strong additivity}: the derivative family of every Hahn-summable
family is summable and $\Dn\bigl(\sum_{j}x_{j}\bigr)=\sum_{j}\Dn x_{j}$. It is
extended to $\SC$ by $\Dn(x+iy)=\Dn x+i\Dn y$; the extension has constant
field $\C$ and commutes with conjugation. In the numbering of
\cref{diff:thm:BM} these are the clauses \ref{diff:BM:leibniz},
\ref{diff:BM:const}, \ref{diff:BM:additive}, \ref{diff:BM:exp},
\ref{diff:BM:norm} and \ref{diff:BM:positivity}, without the surjectivity
clause \ref{diff:BM:surj}.
\end{definition}

The Berarducci--Mantova derivation $\der$ is one normalized derivation, and it
is also surjective; its existence is the input for this part \cite{BM}.

\begin{convention}[The symbol $\Dn$, and the scope of the one-derivation rule]
\label{diff:aut:conv:scope}
Three rules fix the use of the new symbol.
\begin{enumerate}[label=(\arabic*),leftmargin=2.4em]
\item Throughout \cref{diff:aut:part}, $\Dn$ denotes an \emph{arbitrary}
normalized derivation, and every statement holds for every such $\Dn$ unless it
is flagged. The abelian finite-primitive criterion
(\cref{diff:aut:cor:abeliancriterion}) assumes that $\Dn$ is surjective;
the explicit speed threshold of \cref{diff:aut:cor:speeds} needs no
such assumption and holds for every normalized derivation. The
source writes $D$. In this report $D$ never denotes a derivation other than the
fixed one (\cref{diff:warn:concordance}), and the new symbol keeps the
$\der$-specific statements of the other parts visibly distinct.
\item Outside this part every intrinsic statement remains relative to the one
fixed $\der$ (\cref{diff:warn:nouniqueness}, items~N2--N3). No theorem of
another part is upgraded to all normalized derivations. Where a theorem of this
part reproves a $\der$-statement for every $\Dn$ it says so, and it does so only
for the first-order autonomous sector
(\cref{diff:aut:cor:earlier}).
\item Nothing here classifies normalized derivations, asserts that $\No$ has
only one, or asserts that two of them agree on $\SC$: they are shown to agree
on the solutions of every nonzero $F(y,\Dn y)=0$ and on the fields
$K_{\xi}$ of \cref{diff:aut:thm:scales}, and nowhere else.
\end{enumerate}
\end{convention}

\begin{convention}[The alphabet of this part]
\label{diff:aut:conv:alphabet}
$\mathcal X$ is a smooth connected projective curve over $\C$ and $\C(\mathcal X)$
its function field (not the rational function field $\C(X)$ of
\cref{diff:thm:PVrational}); $\xi$ is a nonzero meromorphic vector field, that
is a derivation of $\C(\mathcal X)$, and $\varphi_{\xi}$ the meromorphic
differential with $\varphi_{\xi}(\xi)=1$; $\sigma$ is a local parameter at a point
of $\mathcal X(\C)$; points of $\mathcal X(\SC)$ or of an abelian variety are
bold, $\mathbf y$, and $\st(\mathbf y)$ is their reduction (the source writes
$\bar P$, but a bar is complex conjugation here); $\iota$ is a differential
embedding; $\varphi,\varphi_{j}$ are $1$-forms (the source's $\eta,\theta_{j}$,
letters reserved here for the phase residue and the finite angle); $\mathbf A$
is a constant abelian variety of dimension $d$ (the source's $A$), with formal
group law $\mathcal F$, formal logarithm $\log_{\mathcal F}$ and formal exponential
$\exp_{\mathcal F}$; $\Log$ is a strip logarithm (\cref{diff:aut:lem:log}); $\mathfrak t$
is a time primitive; $\dlog_{\Dn}$ is the abelian logarithmic derivative;
$\Obs^{\C}_{\Dn}$ is a complex primitive obstruction (\cref{diff:aut:sub:primitive});
$\Pcl_{\C}=\Pcl+i\Pcl$. For $v\neq0$ we write $u\prec v$ when
$u/v\in\mm_{\C}$ and $u\sim v$ when $u/v\in1+\mm_{\C}$.
\end{convention}

\begin{lemma}[Finite coordinates have infinitesimal derivatives]
\label{diff:aut:lem:small}
For every normalized derivation,
$\Dn\Ofin_{\C}\subseteq\mm_{\C}$, $\Dn\mm_{\C}\subseteq\mm_{\C}$ and
$\Dn\Ofin_{\C}=\Dn\mm_{\C}$.
\end{lemma}
\begin{proof}
For finite real $x$ and real $r$ the number $\omega-rx$ is positive infinite,
so $1-r\Dn x>0$; as $r$ is arbitrary, $\Dn x\in\mm$. Apply this to both
coordinates, and use $\Dn(c+\varepsilon)=\Dn\varepsilon$ for $c\in\C$.
\end{proof}

This is the ``third formulation'' of \cref{diff:rem:sqrtsmallness}. As the proof
of \cref{diff:thm:orderobstruction} already notes, the argument of
\cref{diff:lem:smallnegative} uses only the positivity rule, $\der\omega=1$ and
the constant field, so it holds verbatim for every normalized derivation; no
surjectivity is used.

\begin{lemma}[Formal evaluation and the chain rule]
\label{diff:aut:lem:evaluation}
For a finite tuple $\varepsilon=(\varepsilon_{1},\dots,\varepsilon_{r})\in\mm_{\C}^{r}$
there is a ring homomorphism
$\C[[T_{1},\dots,T_{r}]]\to\Ofin_{\C}$, $f\mapsto f(\varepsilon)$. It preserves the
constant coefficient modulo $\mm_{\C}$, respects formal composition when the
inner series have zero constant term, and satisfies
\begin{equation}\label{diff:aut:eq:chain}
 \Dn f(\varepsilon)=\sum_{j=1}^{r}(\partial_{T_{j}}f)(\varepsilon)\,\Dn\varepsilon_{j}.
\end{equation}
Formal maps with zero constant term and invertible linear part have
formal inverses, and the two maps evaluate to mutually inverse maps on
infinitesimal tuples.
\end{lemma}
\begin{proof}
Zero coordinates can first be set to zero in the formal series; if
all coordinates vanish, evaluation is the constant coefficient and the
claims are immediate. Otherwise the union $S$ of their supports is a
well-ordered set of positive exponents in the convention of
\cref{diff:eq:tmonomial}. For a fixed output exponent, the ordered-word
lemma \cref{diff:lem:neumann} gives only finitely many contributing
words in $S$. Each occurrence of a letter has at most $r$ labels,
one for each input coordinate. Thus the fully expanded family of
monomials $\varepsilon^\nu$ has finitely many contributions at each
exponent and support in the well-ordered monoid $S^*$. Ordinary
complex coefficients add no exponents. Every nonconstant monomial has
positive exponent, proving the residue assertion.

This finiteness licenses the Cauchy product and proves
multiplicativity. For composition, let each inner series have zero
constant term. In a fully expanded term, each inner factor then
contributes at least one input letter. A fixed output exponent permits
only finitely many input words, hence only finitely many partitions
into nonempty inner factors and finitely many formal monomials at
each resulting degree. The composed evaluation is therefore strongly
summable before regrouping, and agrees with the formal composition.

Strong additivity and Leibniz give
\[
 \Dn(\varepsilon^\nu)=
 \sum_{\substack{1\leq j\leq r\\\nu_j>0}}
       \nu_j\varepsilon^{\nu-e_j}\Dn\varepsilon_j.
\]
For fixed $j$, the resulting sum is the evaluation of
$\partial_{T_j}f$ multiplied by the fixed element
$\Dn\varepsilon_j$. There are only finitely many $j$, giving
\cref{diff:aut:eq:chain}. Finally, formal inversion is constructed
degree by degree: at each positive degree the unknown vector of
coefficients is multiplied by the same invertible ordinary linear
matrix. The two inverse identities then evaluate by the composition
clause.
\end{proof}

For $r=1$ and $\Dn=\der$ this is \cref{diff:lem:smallchain}; the multivariable
form and the composition clause are what the curve and abelian arguments need,
and \cref{diff:warn:smallexplimits}(iii) applies unchanged: the coefficients are
ordinary constants.

\begin{lemma}[Formal implicit uniqueness in a monad]
\label{diff:aut:lem:implicit}
Let $\Theta(X_{1},\dots,X_{r},W)\in\C[[X_{1},\dots,X_{r},W]]$ with
$\Theta(0,0)=0$ and $\partial_{W}\Theta(0,0)=a\neq0$. There is a unique
$W=w(X)\in\C[[X]]$ with zero constant term and $\Theta(X,w(X))=0$. For every
$\varepsilon\in\mm_{\C}^{r}$, $w(\varepsilon)$ is the unique $u\in\mm_{\C}$ with
$\Theta(\varepsilon,u)=0$.
\end{lemma}
\begin{proof}
Solve by homogeneous degree: the unknown degree-$n$ term occurs with coefficient
$a$, and all other terms of that degree are determined by lower degrees.
Evaluation (\cref{diff:aut:lem:evaluation}) gives a solution. For
$u,v\in\mm_{\C}$, use the formal divided difference
\[
 \Theta(X,W)-\Theta(X,W')=(W-W')Q(X,W,W'),
 \qquad Q(0,0,0)=a.
\]
Indeed, each difference $W^k-(W')^k$ is divisible by $W-W'$;
the quotient of the linear term contributes $a$, and every other
term has positive total degree. Evaluation therefore gives
$\Theta(\varepsilon,u)-\Theta(\varepsilon,v)
 =(u-v)(a+\varepsilon')$ with $\varepsilon'\in\mm_{\C}$.
The factor $a+\varepsilon'$ is a unit, so two roots must coincide. The argument also
applies after translating the unknown to any ordinary constant value.
\end{proof}

This is a formal coefficient argument plus a unit argument; it is not a claim
that a countable Newton iteration converges in the fine topology.

\begin{lemma}[Strip logarithms, for every normalized derivation]
\label{diff:aut:lem:log}
For $z\in\SC^{\times}$ write its leading-monomial decomposition
$z=a\mu(1+\varepsilon)$ with $a\in\C^{\times}$, $\mu>0$ a monomial and
$\varepsilon\in\mm_{\C}$, choose an ordinary logarithm $\log a$ of $a$, and put
\begin{equation}\label{diff:aut:eq:Log}
 \Log z=\log a+\log\mu+\Lmm(1+\varepsilon).
\end{equation}
Then $\Log z\in\Strip$ and $\Estrip(\Log z)=z$, so $\Log z$ is a strip logarithm
in the sense of \cref{diff:thm:strip}, and every choice satisfies
$\Dn\Log z=\Dn z/z$. For $z\in\mm_{\C}\setminus\{0\}$ and real $q>0$,
$\Log z\prec\abs z^{-q}$. If $z=\beta\omega^{-1/m}(1+\varepsilon)$ with
$\beta\in\C^{\times}$ and $m\geq1$, then
$\Log z=\log\beta-\tfrac1m\log\omega+\Lmm(1+\varepsilon)$.
\end{lemma}
\begin{proof}
The real part of $\Log z$ is $\log\abs a+\log\mu+\Ree\Lmm(1+\varepsilon)=\log\abs z$
and its imaginary part $\arg a+\Imm\Lmm(1+\varepsilon)$ is finite; since the
real surreal exponential agrees with $\Emm$ on real infinitesimals
\cite{Gonshor}, \cref{diff:eq:stripexp} gives
$\Estrip(\Log z)=a\mu\,\Emm\bigl(\Lmm(1+\varepsilon)\bigr)=z$. Differentiate the three
summands: the first is a constant, the second uses exponential compatibility,
and the third uses \cref{diff:aut:lem:evaluation}; the logarithmic derivatives
add. For the estimate, the last summand is infinitesimal and $\log a$ ordinary,
and $\abs{\log\rho}\prec\rho^{-q}$ for positive infinitesimal $\rho$ follows from
the domination of every real power by the surreal exponential \cite{BM}.
The last formula is the normalized monomial decomposition of
$\beta\omega^{-1/m}(1+\varepsilon)$.
\end{proof}

For $\Dn=\der$ the derivative statement is the remark after
\cref{diff:thm:strip}: all strip logarithms of $z$ have the same intrinsic
derivative $z^{\dagger}$. $\Log$ is used pointwise or on a specified
leading-term class; it is not a global additive logarithm and it does not
define a global surcomplex exponential. With the leading decomposition fixed,
changing the ordinary choice of $\log a$ adds $2\pi i k$ for an ordinary
integer $k$; it changes no derivative or growth estimate.

\section{The realization classification on curves}
\label{diff:aut:sec:curves}

\subsection{Solutions, and the statement}
\label{diff:aut:sub:statement}

A \emph{nonconstant solution} of $(\mathcal X,\xi)$ is a differential-field
embedding
\begin{equation}\label{diff:aut:eq:embedding}
 \iota:(\C(\mathcal X),\xi)\hookrightarrow(\SC,\Dn),
 \qquad \Dn\iota(g)=\iota(\xi g)\quad(g\in\C(\mathcal X)),
\end{equation}
equivalently a nonconstant point $\mathbf y\in\mathcal X(\SC)$ with
$\Dn\mathbf y=\xi(\mathbf y)$. Here a nonconstant point means one not in $\mathcal X(\C)$.
Its Zariski closure over $\C$ is the whole curve: a proper closed
subset of a smooth connected curve consists of finitely many closed
points, all ordinary because $\C$ is algebraically closed.
Consequently no nonzero rational-function denominator vanishes there,
and evaluation gives an embedding of $\C(\mathcal X)$.
Constant points are treated separately. This
curve-and-derivation formalism, with classical formal local normal forms, is
that of Noordman, van der Put and Top \cite{NPT}, whose local discussion
concerns a formal time parameter at zero; here the time is the fixed infinite
surreal $\omega$ with $\Dn\omega=1$, and the realization and exclusion conditions
are different.

Every point of $\mathcal X(\SC)$ has an ordinary reduction $p\in\mathcal X(\C)$
(\cref{diff:aut:lem:projective}). Choose a local parameter $\sigma$ at $p$ and
write
\begin{equation}\label{diff:aut:eq:localfield}
 \xi(\sigma)=f(\sigma)=c_{k}\sigma^{k}+c_{k+1}\sigma^{k+1}+\cdots,
 \qquad c_{k}\neq0,\quad k\in\Z .
\end{equation}
At a simple zero ($k=1$) the coefficient $\lambda=c_1$ does not
depend on the local parameter. Indeed, another parameter has expansion
$\widetilde\sigma=a\sigma+O(\sigma^2)$ with $a\in\C^\times$, and
$\xi(\widetilde\sigma)=(a+O(\sigma))(\lambda\sigma+O(\sigma^2))
=\lambda\widetilde\sigma+O(\widetilde\sigma^2)$.

\begin{theorem}[Autonomous realization classification]
\label{diff:aut:thm:main}
Let $\Dn$ be a normalized derivation. The nonconstant solutions of
$(\mathcal X,\xi)$ reducing to $p$ are as follows.
\begin{enumerate}[label=\textup{(\roman*)},leftmargin=2.4em]
\item If $k\leq0$ there are none.
\item If $k=1$ there are solutions exactly when $\lambda=c_{1}\in\R_{<0}$. With
the normalized formal linearizer $h_{p}$ of \cref{diff:aut:prop:simple} they are
in bijection with $c\in\C^{\times}$, with local parameter
$\sigma=h_{p}^{-1}\bigl(c\,e^{\lambda\omega}\bigr)$.
\item If $k=m+1\geq2$, put $c_{\ast}=-1/(mc_{m+1})$. There are exactly $m$
branches, indexed by $\beta^{m}=c_{\ast}$, each in bijection with $C\in\C$, and
every local parameter has the form
\begin{equation}\label{diff:aut:eq:master}
 \sigma=\beta\,\omega^{-1/m}\,
 w_{p,\beta,C}\Bigl(\omega^{-1/m},\ \frac{\log\omega}{\omega}\Bigr),
 \qquad w_{p,\beta,C}\in1+(X,Y)\C[[X,Y]],
\end{equation}
the series being specified uniquely by the formal implicit equation of
\cref{diff:aut:prop:multiple}.
\end{enumerate}
Distinct listed parameters give distinct solutions, and all nonconstant
solutions occur in the list. The constant solutions of $\Dn\mathbf y=\xi(\mathbf y)$
are exactly the ordinary zeros of $\xi$.
\end{theorem}

Three features make this more than local existence. Projective reduction puts
\emph{every} solution, including those with infinite affine coordinates, into
the local classification. Nonreal simple eigenvalues are excluded, not merely
those with nonnegative real part. And the formulas bound the scales of every
solution independently of its parameters and of the normalized derivation
(\cref{diff:aut:sec:scales}).

\subsection{Projective reduction}
\label{diff:aut:sub:reduction}

\begin{lemma}[Reduction of projective points]
\label{diff:aut:lem:projective}
Every point $\mathbf y$ of a projective variety $\mathcal Y/\C$ over $\SC$ has a
well-defined reduction $\st(\mathbf y)\in\mathcal Y(\C)$, and regular morphisms
over $\C$ commute with reduction. For a smooth curve $\mathcal X$ and
$p\in\mathcal X(\C)$, any local parameter $\sigma$ at $p$ gives a bijection
$\{\mathbf y\in\mathcal X(\SC):\st(\mathbf y)=p\}\leftrightarrow\mm_{\C}$, whose
inverse is given by the formal local coordinate expansions at $p$.
\end{lemma}
\begin{proof}
Embed $\mathcal Y$ in $\mathbb P^{N}_{\C}$ and write $\mathbf y=[z_{0}:\cdots:z_{N}]$.
Choose a coordinate of maximal modulus, say $z_{j}$: the ratios $z_{i}/z_{j}$ are
finite and their standard parts are not all zero, the $j$th being $1$; reducing
the homogeneous equations gives a point of $\mathcal Y(\C)$. Another coordinate
of maximal modulus differs by a finite unit, whose nonzero standard part
rescales all reduced coordinates by one scalar, so the reduction is well
defined. A regular map near the reduction point can be written in
affine coordinates using denominators nonzero at that ordinary point.
Their evaluations are finite units, so taking standard parts commutes
with these expressions. This proves compatibility with regular
morphisms, and hence independence of the projective embedding.

For the valuative interpretation, take the set-sized field $E$
generated over $\C$ by the finitely many homogeneous coordinates and
use its valuation ring $E\cap\Ofin_{\C}$, with residue field $\C$.
The properness criterion \cite[Tag~0BX5]{Stacks} applies in that
set-sized setting; no scheme with a proper-class coefficient field is
needed for this interpretation.

At a smooth point of a curve, choose an affine presentation, with
denominators nonzero at $p$, in which the equations for the coordinates
other than $\sigma$ have an invertible ordinary Jacobian minor.
Translate their ordinary values to zero and write them as a tuple
$Z$. The formal map consisting of $\sigma$ and these equations has
invertible linear part. Its formal inverse gives the other coordinates
as ordinary-coefficient series in $\sigma$. Evaluating at any
$\varepsilon\in\mm_{\C}$ gives a point reducing to $p$; all defining
identities are preserved by \cref{diff:aut:lem:evaluation}, and every
chart denominator remains a unit.

For uniqueness, suppose two tuples $Z,Z'$ have the same value of
$\sigma$ and the same reduction. Subtract the polynomial equations,
or first clear their unit denominators. Their difference factors as
\[
 0=M(\varepsilon,Z,Z')(Z-Z'),
 \qquad M(\varepsilon,Z,Z')=M_0+\Delta,
\]
where $M_0$ is the invertible ordinary Jacobian minor and every entry
of $\Delta$ is infinitesimal. Thus
$\st(\det M)=\det M_0\neq0$, so $M$ is invertible and $Z=Z'$.
This proves the claimed bijection without completeness of an ambient
valuation topology.
\end{proof}

\begin{lemma}[A nonconstant trajectory reduces to a zero]
\label{diff:aut:lem:zero}
If $\mathbf y$ is a nonconstant solution and $p=\st(\mathbf y)$, then $k\geq1$ in
\cref{diff:aut:eq:localfield}.
\end{lemma}
\begin{proof}
Put $u=\sigma(\mathbf y)$. Reduction gives $u\in\mm_{\C}$.
The parameter value $0$ corresponds to the ordinary point $p$ under
\cref{diff:aut:lem:projective}; since $\mathbf y$ is nonconstant,
$u\neq0$. Hence $\Dn u\in\mm_{\C}$ by \cref{diff:aut:lem:small}. The local
expression gives $\Dn u=f(u)=c_{k}u^{k}(1+\varepsilon)$ with $\varepsilon\in\mm_{\C}$.
For $k=0$ the right side has nonzero standard part, and for $k<0$ it is
infinite; both contradict the infinitesimal left side.
\end{proof}

The lemma includes the poles of $\xi$, and because projective reduction also
controls the infinite affine solutions it is stronger than the statement that a
bounded affine solution tends to an equilibrium. Nothing here treats an ordinary
path on a Riemann surface as a fine-continuous path in $\SC$.

\subsection{Simple zeros}
\label{diff:aut:sub:simple}

\begin{lemma}[Constant-coefficient eigenvectors]
\label{diff:aut:lem:eigen}
If $\lambda\in\C$ and $u\in\SC^{\times}$ satisfy $\Dn u=\lambda u$, then
$\lambda\in\R$ and $u=c\,e^{\lambda\omega}$ with $c\in\C^{\times}$; conversely these
solve the equation for real $\lambda$. Such a solution is infinitesimal exactly
when $\lambda<0$.
\end{lemma}
\begin{proof}
Write $\lambda=\alpha+i\beta$, $\rho=\abs u$ and $\zeta=u/\rho$. Conjugation gives
$\Dn(u\bar u)=2\alpha u\bar u$, so $\Dn\rho/\rho=\alpha$, $\Dn\zeta=i\beta\zeta$ and
$\zeta\bar\zeta=1$. The coordinates of $\zeta$ are finite, so $\Dn\zeta$ is
infinitesimal (\cref{diff:aut:lem:small}), while $\st\zeta$ has modulus one;
reducing the equation forces $\beta=0$. Then $\Dn(ue^{-\alpha\omega})=0$ and the
constant field is $\C$. The converse uses $\Dn\omega=1$, and $e^{\alpha\omega}$ is
infinitesimal, ordinary or infinite according as $\alpha$ is negative, zero or
positive.
\end{proof}

For $\Dn=\der$ this is the first-order part of \cref{diff:cor:oscillator}, and the
proof is its second proof; it uses only a finite unit and its standard part, and
assumes no infinite imaginary exponential. The new use is inside an exhaustive
nonlinear classification.

\begin{proposition}[Linearization and all simple-zero solutions]
\label{diff:aut:prop:simple}
Suppose $f(\sigma)=\lambda\sigma+O(\sigma^{2})$ with $\lambda\neq0$. There is a
unique $h_{p}(\sigma)=\sigma+O(\sigma^{2})\in\C[[\sigma]]$ with
$h_{p}'(\sigma)f(\sigma)=\lambda h_{p}(\sigma)$. The nonzero infinitesimal
solutions of $\Dn u=f(u)$ are exactly $u=h_{p}^{-1}(c\,e^{\lambda\omega})$ with
$c\in\C^{\times}$ when $\lambda\in\R_{<0}$, and there are none otherwise.
\end{proposition}
\begin{proof}
Write $f(\sigma)=\lambda\sigma(1+\sigma q(\sigma))$. Then
$\lambda/f(\sigma)-1/\sigma$ is in $\C[[\sigma]]$, because
$(1+\sigma q)^{-1}-1$ is divisible by $\sigma$. Let $H(\sigma)$
be its formal primitive with constant term zero and put
\[
 h_p(\sigma)=\sigma\exp H(\sigma),\qquad
 H(\sigma)=\int_0^\sigma\bigl(\lambda/f(s)-1/s\bigr)\,ds .
\]
The integral divides the coefficient of $s^j$ by the nonzero integer
$j+1$, and the exponential is a formal series in the zero-constant
$H$. Thus $h_p=\sigma+O(\sigma^2)$. Its logarithmic derivative is
$\lambda/f$, proving the required identity. Conversely, if
$h_p=\sigma U$ with $U(0)=1$, the identity forces
$U'/U=\lambda/f-1/\sigma=H'$. The derivative of $U\exp(-H)$
vanishes, so this series equals its constant term $1$. This proves
uniqueness; the linear coefficient $1$ gives the formal inverse. For a
nonzero infinitesimal solution the chain rule gives $\Dn h_{p}(u)=\lambda h_{p}(u)$
with $h_{p}(u)\sim u\neq0$, and \cref{diff:aut:lem:eigen} gives necessity and the
form. Conversely each displayed input is infinitesimal when $\lambda<0$, all
evaluations exist, and differentiating the inverse identity gives $\Dn u=f(u)$;
$c$ is unique by injectivity of the inverse formal map.
\end{proof}

For the quadratic local field $f(\sigma)=\lambda\sigma+a\sigma^2$,
the normalized maps are explicitly
\[
 h_p(\sigma)=\frac{\sigma}{1+(a/\lambda)\sigma},
 \qquad h_p^{-1}(v)=\frac{v}{1-(a/\lambda)v}.
\]
Their denominators are units on infinitesimal inputs, and
$h'_pf=\lambda h_p$ follows directly. This includes the local
linearization used in the logistic example
(\cref{diff:aut:ex:logistic}).

The formal linearization is classical \cite[Lemma~6.1(i), arXiv v1]{NPT}. The exclusion of
every nonreal $\lambda$, including those with negative real part, is the
specifically surcomplex restriction: a classical spiralling decay at
$\lambda=-1+i$ does not become a nonzero element of $\SC$ satisfying that
eigenvalue equation.

\subsection{Multiple zeros: complete power--logarithm families}
\label{diff:aut:sub:multiple}

Suppose $f(\sigma)=c_{m+1}\sigma^{m+1}+O(\sigma^{m+2})$ with $m\geq1$. Integrating
the Laurent series $1/f$ termwise gives the \emph{time primitive}
\begin{equation}\label{diff:aut:eq:time}
 \mathfrak t(\sigma)=c_{\ast}\sigma^{-m}+\sum_{r=1}^{m-1}e_{r}\sigma^{-r}
 +\varrho\Log\sigma+g(\sigma),
 \qquad c_{\ast}=-\frac{1}{mc_{m+1}},\quad g\in\sigma\C[[\sigma]],
\end{equation}
where $\varrho$ is the coefficient of $\sigma^{-1}$ in $1/f$, the sum is empty
for $m=1$, and all other integration constants are fixed to zero (changing them
only relabels the parameter $C$ below). For every nonzero infinitesimal $u$,
\cref{diff:aut:lem:evaluation,diff:aut:lem:log} give
$\Dn\mathfrak t(u)=\Dn u/f(u)$, the denominator being
$c_{m+1}u^{m+1}(1+\mm_{\C})\neq0$. So a solution of $\Dn u=f(u)$ satisfies
\begin{equation}\label{diff:aut:eq:timeconstant}
 \mathfrak t(u)=\omega+C\qquad\text{for exactly one }C\in\C,
\end{equation}
because the constant field is $\C$. Here uniqueness of $C$ is for
the chosen strip logarithm in $\mathfrak t$; changing that choice
can relabel $C$. The scale argument below is independent of the choice.

\begin{lemma}[The leading scale is forced]
\label{diff:aut:lem:leading}
Every nonzero infinitesimal solution satisfies $u\sim\beta\,\omega^{-1/m}$ for
exactly one $\beta\in\C^{\times}$ with $\beta^{m}=c_{\ast}$.
\end{lemma}
\begin{proof}
Divide the remaining terms of \cref{diff:aut:eq:time} by
$u^{-m}$. The negative powers give multiples of $u^{m-r}$ with
$m-r>0$; the regular part gives $u^mg(u)$; all are infinitesimal.
The logarithmic term is infinitesimal after this division by
\cref{diff:aut:lem:log} with $q=m$. Thus
$\mathfrak t(u)=c_*u^{-m}(1+\eta)$ for some $\eta\in\mm_\C$.
Since $(\omega+C)/\omega=1+C/\omega$ also has residue $1$,
\cref{diff:aut:eq:timeconstant} implies
$c_*u^{-m}/\omega\in1+\mm_\C$. Inverting this unit gives
\[
 (u/x_m)^m\in c_*(1+\mm_\C),\qquad x_m=\omega^{-1/m}.
\]
Writing $v=u/x_m$, the positive number $|v|^m$ has ordinary
residue $|c_*|>0$. Hence $|v|$ is bounded above and below by
positive ordinary real numbers: an infinite or infinitesimal $|v|$
would make $|v|^m$ infinite or infinitesimal. Therefore $v$ is a
finite unit. Its unique standard part $\beta\in\C^\times$ satisfies
$\beta^m=c_*$, and $v/\beta\in1+\mm_\C$, proving the claim.
\end{proof}

Fix $\beta^{m}=c_{\ast}$ and an ordinary logarithm $\log\beta$, and put
$x_{m}=\omega^{-1/m}$ and $y_{\omega}=(\log\omega)/\omega$. Both are infinitesimal,
though not treated as algebraically independent after evaluation. Writing
$u=\beta x_{m}W$ with $W\in1+\mm_{\C}$, dividing \cref{diff:aut:eq:timeconstant} by
$\omega=x_{m}^{-m}$ and using the last formula of \cref{diff:aut:lem:log} gives
\begin{equation}\label{diff:aut:eq:implicit}
 0=W^{-m}-1+\sum_{r=1}^{m-1}e_{r}\beta^{-r}x_{m}^{m-r}W^{-r}
 -\frac{\varrho}{m}\,y_{\omega}
 +x_{m}^{m}\bigl(\varrho\log\beta+\varrho\Lmm(W)+g(\beta x_{m}W)-C\bigr),
\end{equation}
all powers and logarithms of $W$ being expanded at $W=1$.
In particular, the logarithmic contribution is
\[
 x_m^m\varrho\Log(\beta x_mW)
 =-\frac{\varrho}{m}y_\omega
   +x_m^m\varrho\bigl(\log\beta+\Lmm(W)\bigr),
\]
which fixes the sign and the factor $1/m$ in
\cref{diff:aut:eq:implicit}.

\begin{proposition}[Existence, uniqueness and completeness at a multiple zero]
\label{diff:aut:prop:multiple}
For every $\beta^{m}=c_{\ast}$ and $C\in\C$ there is a unique
$w_{\beta,C}(X,Y)\in1+(X,Y)\C[[X,Y]]$ satisfying \cref{diff:aut:eq:implicit}
formally with $(x_{m},y_{\omega})$ replaced by $(X,Y)$. Then
$u_{\beta,C}=\beta x_{m}w_{\beta,C}(x_{m},y_{\omega})$ is a nonzero infinitesimal
solution of $\Dn u=f(u)$, and once the logarithm of each $\beta$ is fixed every
such solution equals exactly one $u_{\beta,C}$.
\end{proposition}
\begin{proof}
Put $W=1+U$. The negative powers of $1+U$ and
$\Lmm(1+U)$ are ordinary formal series in $U$, and the argument
$\beta X(1+U)$ of $g$ has zero constant term. Thus the right side
of \cref{diff:aut:eq:implicit} defines a series in $(X,Y,U)$.
Every term except $(1+U)^{-m}-1$ contains $X$ or $Y$, since
$m-r\geq1$ and $m\geq1$. Its constant term is zero and its
$U$-derivative at the origin is exactly $-m\neq0$.
The formal implicit lemma gives a unique zero-constant
$U_{\beta,C}(X,Y)$, hence $w_{\beta,C}=1+U_{\beta,C}$.

Both $x_m$ and $y_\omega$ are infinitesimal, so the evaluation
lemma gives a root in $1+\mm_\C$. Reversing the algebra yields
$\mathfrak t(u_{\beta,C})=\omega+C$. Differentiation gives
$\Dn u_{\beta,C}/f(u_{\beta,C})=1$; the denominator is nonzero
and the leading term of $u_{\beta,C}$ is $\beta x_m$, so this is
a nonzero solution of $\Dn u=f(u)$.

Conversely, for any nonzero infinitesimal solution the leading-scale
lemma fixes $\beta$, and the time identity fixes $C$ after the
logarithm of that $\beta$ has been fixed. Its quotient
$u/(\beta x_m)$ solves the same evaluated implicit equation.
Uniqueness in the entire monad of $1$ identifies it with the
constructed value, regardless of its support. This step uses the
unit divided difference from \cref{diff:aut:lem:implicit}; it does
not assume that formal evaluation is injective. Distinct $\beta$
have different leading coefficients, and distinct $C$ have different
values of $\mathfrak t$. Finally $c_*\neq0$ and characteristic zero
give exactly $m$ distinct roots of $\beta^m=c_*$ in $\C$.
\end{proof}

The time parameter also has an exact first-order separation. On one
fixed branch, let $C\neq C'$ and set
$W_C=u_{\beta,C}/(\beta x_m)$ and
$W_{C'}=u_{\beta,C'}/(\beta x_m)$. Subtracting the two implicit
equations gives
\[
 (-m+\eta)(W_C-W_{C'})=(C-C')x_m^m,
 \qquad \eta\in\mm_\C.
\]
Indeed, only the term $-Cx_m^m$ depends explicitly on $C$, and the
divided difference in $W$ has residue $-m$. Consequently
\[
 u_{\beta,C}-u_{\beta,C'}
 \sim-\frac{\beta}{m}(C-C')\,\omega^{-1-1/m}.
\]
Thus the constant distinguishes solutions at a specified next scale.

When $f(\sigma)=c_{m+1}\sigma^{m+1}$ exactly, all $e_r$, $\varrho$
and $g$ vanish, and the implicit solution has the explicit form
\[
 w_{\beta,C}(X,Y)=(1+CX^m)^{-1/m},\qquad
 u_{\beta,C}=\beta\omega^{-1/m}(1+C/\omega)^{-1/m}.
\]
The fractional power is the formal binomial series with constant
term $1$, evaluated at the infinitesimal $C/\omega$. Its inverse
$m$th power is $1+C/\omega$, so it gives the time identity
$c_*u_{\beta,C}^{-m}=\omega+C$ directly, for all $m\geq1$.

\begin{remark}[No hidden solutions at another scale]
\label{diff:aut:rem:nohidden}
A formal family alone would not prove the theorem. Completeness has two steps:
every solution satisfies $\mathfrak t(u)=\omega+C$ because the constant field is
exactly $\C$, and that identity forces the scale $\omega^{-1/m}$. Uniqueness in
the actual monad then excludes corrections with supports outside the displayed
two-variable series; in particular there is no additional beyond-all-orders
parameter.
\end{remark}

\begin{remark}[Coordinate and branch choices]
\label{diff:aut:rem:choices}
A change of local parameter transports the solutions by an invertible
ordinary formal series. If
$\widetilde\sigma=a\sigma+O(\sigma^2)$, $a\in\C^\times$, then
$\widetilde c_{m+1}=a^{-m}c_{m+1}$, so
$\widetilde c_*=a^mc_*$ and the corresponding leading coefficient is
$\widetilde\beta=a\beta$. The number $m$ of branches is unchanged.

On a fixed coordinate branch, replacing $\log\beta$ by
$\log\beta+2\pi i k$ adds $2\pi i k\varrho$ to $\mathfrak t$ and
relabels the same solution by $C'=C+2\pi i k\varrho$.
The combination $\varrho\log\beta-C$ in the implicit equation is
unchanged. Thus the parametrizations depend on the choices, while
the family of points and the existence criterion do not; in
particular the leading coefficient $\beta$ is not a coordinate
invariant.
\end{remark}

\begin{proof}[Proof of \cref{diff:aut:thm:main}]
Every nonconstant solution reduces to an ordinary point
(\cref{diff:aut:lem:projective}), which is a zero of $\xi$ by
\cref{diff:aut:lem:zero}; \cref{diff:aut:prop:simple,diff:aut:prop:multiple} give
exactly the local lists. Conversely, evaluate the completed local coordinate map
of $\mathcal X$ at each listed nonzero infinitesimal: the point is nonconstant,
hence generic, and $\Dn\sigma=f(\sigma)$ with the chain rule and the rule for
inverses gives $\Dn g(\mathbf y)=(\xi g)(\mathbf y)$ for every rational function
$g$. To check injectivity directly, a nonzero local Laurent
expansion factors as $a\sigma^r(1+\sigma q(\sigma))$ with
$a\in\C^\times$, $r\in\Z$ and $q\in\C[[\sigma]]$.
At a nonzero infinitesimal $u$ its value is
$au^r(1+u q(u))\neq0$, since the last factor has residue $1$.
Thus no nonzero rational function is killed, and the construction
gives a field embedding, not only an assignment of coordinates. A constant point has zero derivative and solves the equation
exactly when $\xi$ is regular and vanishes there.
\end{proof}

\section{One finite-scale Hahn field, and derivation independence}
\label{diff:aut:sec:scales}

\subsection{All solutions in one set-sized field}
\label{diff:aut:sub:scales}

Let $p_{1},\dots,p_{r}$ be the multiple zeros of $\xi$, of orders $m_{j}+1$, and
$\lambda_{1},\dots,\lambda_{s}$ the negative real coefficients at simple zeros.
Put $M=\operatorname{lcm}(m_{1},\dots,m_{r})$ (the empty least common multiple being
$1$) and let
\begin{equation}\label{diff:aut:eq:group}
 \mathfrak G_{\xi}=\bigl\langle\omega^{1/M},\ \log\omega,\
 e^{\lambda_{1}\omega},\dots,e^{\lambda_{s}\omega}\bigr\rangle
\end{equation}
be the multiplicative group they generate. These are positive Conway
monomials: $\omega$ is log-atomic, so $\log\omega$ is a monomial;
positive rational powers of $\omega$ are monomials; and
$\lambda_j\omega$ is purely infinite, whose exponential is a monomial
\cite[\S3.2, Definition~5.1 and Remark~5.18]{BM}. Thus $\mathfrak G_{\xi}$ is a finitely
generated ordered monomial group, relations among the generators being allowed.
Let $K_{\xi}=\C((\mathfrak G_{\xi}))\subseteq\SC$ be the field of series with
complex coefficients and reverse well-ordered support in $\mathfrak G_{\xi}$. In
the additive notation of \cref{diff:sec:localization} it is $K_{\Gamma_{\xi}}$ for
the finitely generated group $\Gamma_{\xi}$ of exponents of its monomials, a group
not contained in $\R$. Indeed, $\log\omega$ is positive infinite but
$(\log\omega)^n\prec\omega$ for every positive integer $n$; the
ordered group already has distinct logarithmic and power scales.

\begin{theorem}[Uniform finite-scale localization]
\label{diff:aut:thm:scales}
For every normalized derivation $\Dn$:
\begin{enumerate}[label=\textup{(\roman*)},leftmargin=2.4em]
\item $K_{\xi}$ is a set-sized differential subfield of $\SC$ of cardinality
$2^{\aleph_{0}}$, and $\Dn|_{K_{\xi}}$ does not depend on $\Dn$;
\item every constant or nonconstant solution of $(\mathcal X,\xi)$ lies in
$\mathcal X(K_{\xi})$;
\item the same $K_{\xi}$ serves simultaneously for every ordinary parameter of
the families of \cref{diff:aut:thm:main}.
\end{enumerate}
\end{theorem}
\begin{proof}
A Hahn field over an ordered group is a field: after extracting a leading term,
invert a unit $1+\varepsilon$ by its geometric series
(\cref{diff:lem:neumann}). Every monomial in that inverse remains in the
same group. The finitely generated group $\mathfrak G_\xi$ is countable,
so its series form a set of cardinality at most
$|\C|^{\aleph_0}=2^{\aleph_0}$; this bound includes every permitted
support and coefficient assignment. The constant series give the reverse
inequality, and ordinary coefficients are constants of every $\Dn$.

For a monomial
$\mathfrak g=\omega^{a/M}(\log\omega)^b\prod_j e^{n_j\lambda_j\omega}$
with $a,b,n_j\in\Z$, exponential compatibility, the product rule and
$\Dn\omega=1$ give the \emph{three-shift formula}
\begin{equation}\label{diff:aut:eq:threeshift}
 \frac{\Dn\mathfrak g}{\mathfrak g}=\frac{a}{M\omega}+\frac{b}{\omega\log\omega}
 +\sum_{j=1}^{s}n_{j}\lambda_{j}.
\end{equation}
In particular $\Dn\log\omega=1/\omega$, and the same rules apply to
negative integer powers by differentiating inverses. Relations among the
generators cause no ambiguity: two product representations are equal in
$\SC$, hence their differentiated values are equal for each $\Dn$.
The displayed formula is independent of that derivation.

For a reverse well-ordered support $S\subseteq\mathfrak G_\xi$, all
resulting derivative supports lie in
\[
 S\ \cup\ \omega^{-1}S\ \cup\ (\omega\log\omega)^{-1}S.
\]
Each translate is reverse well-ordered, and a finite union has the same
property. For a fixed output monomial, each translate has at most one
preimage in $S$, so at most three input monomials contribute. Thus the
termwise derivative family is Hahn-summable, including possible collisions
or cancellations between translates. Strong additivity now gives
$\Dn K_\xi\subseteq K_\xi$, and the coefficient formula proves that
all normalized derivations have the same restriction to the whole Hahn
field, not just to its finite sums.

At a multiple zero, $m_j\mid M$, so
$x_{m_j}=(\omega^{1/M})^{-M/m_j}$ and
$y_\omega=(\log\omega)\omega^{-1}$ belong to $\mathfrak G_\xi$.
Every evaluated term of \cref{diff:aut:eq:master} has its monomial
in that group, and the evaluation lemma supplies summability.
At a stable simple zero, the inverse linearizer is an ordinary
zero-constant series in $c e^{\lambda_j\omega}$, so its value is also
in $K_\xi$. Varying $c$, $C$ or $\beta$ changes only ordinary
coefficients and does not enlarge the group.

Finally choose an affine chart about the ordinary reduction point.
Its coordinate functions have ordinary formal expansions in the local
parameter. Evaluation uses only monomials from $\mathfrak G_\xi$,
by the same support lemma, and remains in $K_\xi$; rational functions
may also have a finite Laurent principal part, handled by inversion
in the field. Hence the entire point lies in $\mathcal X(K_\xi)$.
Constant solutions already have coordinates in $\C$. This proves all
three assertions, including cases where either list of zeros is empty.
\end{proof}

\begin{remark}[An instance of the monomial criterion, and a multi-scale workspace]
\label{diff:aut:rem:monomial}
\Cref{diff:aut:eq:threeshift} says that $\mathfrak G_{\xi}$ is differentially
closed on monomials in the sense of \cref{diff:def:monomialclosed}, so for
$\Dn=\der$ the stability of $K_{\xi}$ is an instance of
\cref{diff:lem:monomialcriterion}. The field $K_{\xi}$ is an explicit
\emph{multi-scale} $\der$-stable Hahn workspace, with power, logarithmic and
exponential scales. The open continuation ``classify logarithmic derivatives in
multi-scale $\der$-stable Hahn workspaces'' (\cref{diff:sub:nc:open}) is adjacent:
member~11 supplies such workspaces but does not classify logarithmic derivatives
on them, so the continuation stands.
\end{remark}

\begin{corollary}[A continuum of independent solutions inside two scales]
\label{diff:aut:cor:independent}
The Hahn field $\C\bigl((\langle\omega,\log\omega\rangle)\bigr)$ contains $2^{\aleph_{0}}$
distinct solutions of $\Dn u=u^{3}-u^{2}$ that are algebraically independent
over $\C$ as a family; the field they generate has transcendence degree
$2^{\aleph_{0}}$ over $\C$.
\end{corollary}
\begin{proof}
At $u=0$ the vector field has a double zero with leading coefficient
$-1$. Thus $m=1$, $c_*=1$, and \cref{diff:aut:prop:multiple} gives
a nonzero infinitesimal solution $u_C\sim\omega^{-1}$ for every
$C\in\C$, in the two-scale field
$K=\C((\langle\omega,\log\omega\rangle))$.
Fixing the logarithm of $\beta=1$, distinct $C$ give distinct time
values and therefore distinct solutions. They are nonconstant:
a nonzero infinitesimal is not an ordinary complex number.

The three-shift proof makes $K$ a set-sized differential field, and
its constant field is exactly $\C$, since
$\C\subseteq K\subseteq\SC$ and $\Ker\Dn=\C$ on $\SC$.
Rosenlicht's algebraic-independence theorem, in the form of
\cite[Lemma~7.2 and Proposition~7.3, arXiv v1]{NPT}, applies to
any finite set of distinct nonconstant solutions in this differential
field. Every polynomial relation uses finitely many variables, so the
whole family $(u_C)_{C\in\C}$ is algebraically independent.
It has cardinality $|\C|=2^{\aleph_0}$ and generates a field of
exactly that transcendence degree. This uses the classical theorem
as an import, entirely within a set-sized field.
\end{proof}

A finitely generated \emph{monomial group} allows arbitrary well-based infinite
sums, which is why \cref{diff:aut:cor:independent} is consistent with
\cref{diff:aut:thm:scales}: finite monomial generation is not finite
transcendence degree. The notion is algebraic independence of differentially
\emph{algebraic} solutions; it is a different notion from the differential
transcendence proved in the companion tail-spans report \cite{RepoTailSpans},
and there is no conflict.

\subsection{Independence of the normalized derivation}
\label{diff:aut:sub:independence}

\begin{theorem}[Derivation-independent autonomous sector]
\label{diff:aut:thm:independence}
Let $\Dn_{1},\Dn_{2}$ be normalized derivations. For every smooth connected
projective curve $\mathcal X/\C$ and nonzero meromorphic $\xi$, the differential
embeddings $(\C(\mathcal X),\xi)\hookrightarrow(\SC,\Dn_{1})$ and
$(\C(\mathcal X),\xi)\hookrightarrow(\SC,\Dn_{2})$ are the same embeddings of
underlying fields. For every $F\in\C[Y,Z]$,
\begin{equation}\label{diff:aut:eq:independence}
 \{y\in\SC:F(y,\Dn_{1}y)=0\}=\{y\in\SC:F(y,\Dn_{2}y)=0\},
\end{equation}
and if $F\neq0$ and $y$ lies in either set then $\Dn_{1}y=\Dn_{2}y$. No equality
of the two derivations on all of $\SC$ is asserted.
\end{theorem}
\begin{proof}
For curves, \cref{diff:aut:thm:scales} places every solution for
either derivation in the same $K_\xi$, where the derivations agree.
The corresponding field embedding is therefore differential for one
exactly when it is differential for the other.

For the scalar assertion, $F=0$ gives the whole class $\SC$ for
both solution collections; no derivative equality is claimed in
that case. Now let $F\neq0$. A nonzero constant has no solutions.
Otherwise factor $F$ over $\C$ and consider a solution for $\Dn_1$;
some irreducible factor $R$ vanishes at $(y,\Dn_1y)$.
If $y\in\C$, both derivatives vanish and there is nothing to prove.
A factor involving only $Y$ forces this case. A factor involving
only $Z$ is $Z-c$ up to a nonzero scalar, and
$\Dn_1y=c$ implies $\Dn_1(y-c\omega)=0$, hence
$y=c\omega+d$ with $d\in\C$. Both derivations agree on that value.

It remains to treat $y\notin\C$ with $R$ involving both variables.
The evaluation map from the one-dimensional integral algebra
$A=\C[Y,Z]/(R)$ to $\SC$ is injective. Indeed, a nonzero kernel
would be a nonzero prime of this affine curve, hence maximal;
its residue field would be $\C$ by the weak Nullstellensatz,
forcing $y\in\C$. Thus evaluation extends to an embedding of
$L=\operatorname{Frac}(A)$. Here $Y$ is transcendental and
$L/\C(Y)$ is finite separable. Irreducibility in characteristic
zero gives $R_Z\neq0$ in $L$: otherwise $R$ would divide its
nonzero partial derivative of smaller $Z$-degree.

On the smooth projective curve with function field $L$, define
\begin{equation}\label{diff:aut:eq:inducedfield}
 \xi Y=Z,\qquad \xi Z=-\frac{R_Y}{R_Z}\,Z.
\end{equation}
These values annihilate $R$ and define the unique derivation on
$L$ with $\xi Y=Z$, by separability. It is nonzero, since $Z=0$
in $A$ would force the irreducible $R$ to be associated to $Z$,
contrary to its involving both variables.

Write $z=\Dn_1y$. Differentiating the evaluated relation gives
\[
 R_Y(y,z)z+R_Z(y,z)\Dn_1z=0.
\]
The denominator $R_Z(y,z)$ is nonzero by injectivity, so this
identity proves that evaluation intertwines the derivations on
both generators, hence on all of $L$. The curve assertion then
makes the same embedding differential for $\Dn_2$, giving
$\Dn_2y=z=\Dn_1y$ and $F(y,\Dn_2y)=0$.
Interchanging the derivations proves the equality of the two
solution collections.
\end{proof}

For example, take $F(Y,Z)=Z^2-Y$. On its rational curve the induced
vector field has $\xi Y=Z$ and $\xi Z=1/2$. If $y$ is a nonconstant
solution, put $z=\Dn y\neq0$; differentiating $z^2=y$ gives
$2z\Dn z=z$, hence $z=(\omega+C)/2$ and
\[
 y=\frac{(\omega+C)^2}{4},\qquad C\in\C.
\]
Every displayed value solves the equation. There is also the scalar
constant solution $y=0$, but the point $(Y,Z)=(0,0)$ is not a
constant solution of the induced vector field, since $\xi Z=1/2$
there. This is why constant scalar solutions are separated before
passing to an embedding of a curve's function field.


\begin{corollary}[Uniform scalar bound]
\label{diff:aut:cor:scalarbound}
For every nonzero $F\in\C[Y,Z]$ there is one Hahn subfield
$K_{F}=\C((\mathfrak G_{F}))\subseteq\SC$ with a finitely generated monomial group
and cardinality $2^{\aleph_{0}}$ containing all solutions of $F(y,\Dn y)=0$ for
every normalized derivation simultaneously. In particular that solution
collection is a set.
\end{corollary}
\begin{proof}
For each irreducible factor $R$ involving both variables, form the
curve and vector field of \cref{diff:aut:eq:inducedfield} and take
its finitely generated group $\mathfrak G_R$ from
\cref{diff:aut:thm:scales}. Let $\mathfrak G_F$ be generated by the
union of these finitely many groups together with $\omega$ and
$\log\omega$. It is finitely generated even if there are no such
factors, and the construction depends only on $F$.

The inclusion of each monomial group includes its entire Hahn field,
so all nonconstant solutions belonging to a mixed factor are in
$K_F$. Factors in $Y$ alone contribute constants, and factors
$Z-c$ contribute $c\omega+d$; these too lie in $K_F$.
This covers every factor and every normalized derivation at once.
The three-shift and countability arguments apply to the union of
generators, proving differential stability, a common derivative
restriction and cardinality $2^{\aleph_0}$. A nonzero constant
$F$ has no solutions and is covered by the same choice of the
two base generators.
\end{proof}

The hypothesis $F\neq0$ is necessary for a set-sized bound: if
$F=0$, its solution collection is all of $\SC$.


\begin{corollary}[Real autonomous equations]
\label{diff:aut:cor:realindep}
For $F\in\R[Y,Z]$ the set of solutions in $\No$ does not depend on the
normalized derivation, and on every solution of a nonzero such equation all
normalized derivations take the same value.
\end{corollary}
\begin{proof}
Intersect \cref{diff:aut:eq:independence} with $\No$ and use the pointwise
equality of derivatives.
\end{proof}

\begin{corollary}[Two earlier corollaries hold for every normalized derivation]
\label{diff:aut:cor:earlier}
For every normalized derivation $\Dn$: the equation $\Dn y=\lambda y$ with
$\lambda\in\C$ has a nonzero solution in $\SC$ exactly when $\lambda\in\R$, the
solutions then being $c\,e^{\lambda\omega}$, $c\in\C$; and the only solutions in
$\SC$ of $\Dn u=1+u^{2}$ are $u=\pm i$. These are the first-order statements of
\cref{diff:cor:oscillator} and \cref{diff:cor:riccati}. The second-order
statements of \cref{diff:cor:oscillator}, such as the absence of a nonzero
solution of $\der^{2}y+y=0$, are not first-order autonomous equations in one
unknown, are not covered by \cref{diff:aut:thm:independence}, and are not
upgraded.
\end{corollary}
\begin{proof}
The first statement is \cref{diff:aut:lem:eigen}. The second is
\cref{diff:aut:ex:riccati}, or \cref{diff:aut:thm:independence} with
$F=Z-1-Y^{2}$ together with \cref{diff:cor:riccati} for $\der$.
\end{proof}

\begin{remark}[What the independence theorem does not say]
\label{diff:aut:rem:notsay}
It does not say that $\No$ has only one normalized derivation, nor that every
surreal is first-order autonomous differential-algebraic over $\C$; it concerns a
sharply specified common sector. Coefficients in $\SC\setminus\C$, higher
differential order and simultaneous systems are outside its statement, and
their coefficients or solutions may probe scales on which normalized
derivations differ. Which higher-order sectors distinguish normalized
derivations is left open, and nothing here implies a uniform finite-scale bound
for arbitrary higher-order equations.
\end{remark}

\section{Real branches and worked equations}
\label{diff:aut:sec:examples}

Suppose $\mathcal X,\xi,p,\sigma$ are defined over $\R$ and
$f(\sigma)=c\sigma^{m+1}+O(\sigma^{m+2})$ with $c\in\R^{\times}$. The real version
of the time primitive uses $\log\abs\sigma$ in place of $\Log\sigma$, and on a
branch $\sigma\sim\beta\omega^{-1/m}$ with real $\beta$ one uses $\log\abs\beta$ in
\cref{diff:aut:eq:implicit}; all coefficients are then real for real $C$.

\begin{corollary}[Exact real branch count]
\label{diff:aut:cor:realbranches}
At such a multiple zero the real nonconstant solutions are exactly the families
of \cref{diff:aut:prop:multiple} with real $\beta$, real $C$ and the real logarithm
convention. There is one real branch when $m$ is odd; when $m$ is even there are
two if $c<0$ and none if $c>0$; each branch is in bijection with $\R$. At a simple
zero, real nonconstant solutions exist exactly when $\lambda<0$, with
real nonzero amplitude in the normalized linearizer.
\end{corollary}
\begin{proof}
For nonzero real $u$, the sign of $u$ is fixed and
$\Dn\log|u|=\Dn u/u$, whether $u>0$ or $u<0$.
The real time primitive therefore satisfies
$\mathfrak t_{\R}(u)=\omega+C$ with $C\in\R$.
The leading-scale argument gives a real
$\beta=\st(u/\omega^{-1/m})\neq0$, and
$W=u/(\beta\omega^{-1/m})\in1+\mm$ is positive.
Thus the real implicit equation uses $\log|\beta|$ and $\log W$.
Conversely, with real $\beta$ and $C$ that formal equation has real
coefficients. Its unique solution has real coefficients, either by
the coefficient recursion or by uniqueness after complex conjugation,
and evaluates to a real solution. Distinct time constants give
distinct solutions as before.

If $m$ is odd, $\beta^m=-1/(mc)$ has exactly one real root.
If $m$ is even, it has two real roots precisely when $c<0$,
and none when $c>0$. These roots are nonzero and index the stated
families. At a real simple zero the normalized linearizer and its
inverse have real coefficients. Therefore
$h_p^{-1}(a e^{\lambda\omega})$ is real exactly when its amplitude
$a$ is real; nonconstant infinitesimal solutions require
$\lambda<0$ and $a\neq0$.
\end{proof}

The sign restrictions reflect the orientation $\Dn\omega=1$, not a choice of
complex branch cut (compare \cref{diff:rs:warn:orientation}). For example,
if $\beta<0$ and the complex convention chooses
$\log\beta=\log|\beta|+i\pi$, then
$\mathfrak t_{\C}=\mathfrak t_{\R}+i\pi\varrho$ and
$C_{\C}=C_{\R}+i\pi\varrho$ on that branch. Thus a real solution
need not have a real time constant in an arbitrary complex-logarithm
convention. The real primitive removes this fixed imaginary shift.

\begin{example}[Growth can be a stable branch at infinity]
\label{diff:aut:ex:growth}
$\Dn u=u$ has $u=0$ and $u=c\,e^{\omega}$, $c\in\C^{\times}$. At the affine zero the
eigenvalue is $+1$, so no nonconstant solution reduces to $0$; at infinity the
parameter $\sigma=1/u$ has $\Dn\sigma=-\sigma$, eigenvalue $-1$, and every
nonconstant solution reduces to infinity. ``Negative simple eigenvalue'' is a
projective statement, not a restriction to bounded affine solutions.
Directly, $\Dn(u e^{-\omega})=0$ forces $u=c e^\omega$,
so the family is exhaustive. For $c\neq0$ its local coordinate
is $\sigma=c^{-1}e^{-\omega}\in\mm_\C$, displaying the stable
branch in the chart at infinity.
\end{example}

\begin{example}[The logistic equation]
\label{diff:aut:ex:logistic}
For $\Dn u=-u+u^{2}$ the equilibria $0$ and $1$ have eigenvalues $-1$ and $+1$,
and at infinity $\sigma=1/u$ satisfies $\Dn\sigma=\sigma-1$, nonzero at
$\sigma=0$. Every nonconstant solution is therefore
\[
 u=\frac{c e^{-\omega}}{1+c e^{-\omega}},\qquad c\in\C^\times.
\]
Indeed, for $u\neq1$ put $v=u/(1-u)$. The quotient rule gives
$\Dn v=-v$, hence $v=c e^{-\omega}$ by
\cref{diff:aut:lem:eigen}. Conversely $1+c e^{-\omega}$ is a
finite unit of residue $1$, so the formula is always defined and
solves the logistic equation. The value $c=0$ gives the equilibrium
$u=0$; the other equilibrium $u=1$ is separate and is never obtained
by an ordinary $c$. Since $\st(u)=0$ for every member of the
nonconstant family, none has been lost at the other affine zero
or at infinity. Real solutions correspond exactly to real $c$,
together with the separate equilibrium $1$.
\end{example}

\begin{example}[Riccati without a tangent trajectory]
\label{diff:aut:ex:riccati}
For $\Dn u=1+u^{2}$ the zeros $\pm i$ are simple with eigenvalues $\pm2i$, and at
infinity $\Dn(1/u)=-(1+(1/u)^{2})$ is nonzero at $1/u=0$. By
\cref{diff:aut:thm:main} the only solutions in $\SC$ are $u=i$ and $u=-i$. This
reproves \cref{diff:cor:riccati} for every normalized derivation by a second,
projective route, and it rules out a surcomplex element playing the differential
role of $\tan\omega$, without any assertion about a separately defined symbolic
trigonometric function.
\end{example}

\begin{example}[An explicit logarithmic correction]
\label{diff:aut:ex:logcorrection}
Consider the branch at $u=0$ of $\Dn u=-u^{2}(1+u)$. A time primitive is
$\mathfrak t(u)=1/u+\Log u-\Lmm(1+u)$, with $\mathfrak t'(u)=-1/(u^{2}(1+u))$;
here $m=1$, $\beta=1$, $\varrho=1$, and we choose $\log\beta=0$.
Use formal variables $X,Y$ before evaluating them at
$x_1=1/\omega$ and $y_\omega=(\log\omega)/\omega$.
The formal implicit equation for the unit $W$ is
\[
 W^{-1}+X\Lmm(W)-Y-X\Lmm(1+XW)-CX-1=0,
\]
and with $b=Y+CX$ its first homogeneous terms are
$W=1-b+(b^{2}-Xb-X^{2})+O((X,Y)^{3})$. Thus, with $L=\log\omega$,
\[
 u_{C}=\frac1\omega-\frac{L+C}{\omega^{2}}+\frac{(L+C)^{2}-(L+C)-1}{\omega^{3}}
 +\cdots,
\]
the remainder being the evaluation of a series in
$X(X,Y)^3\C[[X,Y]]$. This is a formal total-degree statement;
it does not assert convergence of truncations in the fine topology.
The formula describes all \emph{nonconstant} solutions reducing to
$0$; the constant solution $u=0$ must also be included.

There is a second family at $u=-1$. In the parameter $v=u+1$ the
equation is $\Dn v=-v(1-v)^2$, with eigenvalue $-1$.
Its normalized linearizer is explicitly
\[
 h(v)=\frac{v}{1-v}\exp\!\left(\frac{v}{1-v}\right),
 \qquad h(v)=v+O(v^2).
\]
Here the exponential is the ordinary formal series, evaluated only
at an infinitesimal argument. Its logarithmic derivative is
$h'(v)/h(v)=1/(v(1-v)^2)$, so $h'(v)(-v(1-v)^2)=-h(v)$.
Thus the nonconstant solutions reducing to $-1$ are
$u=-1+h^{-1}(a e^{-\omega})$, $a\in\C^\times$.
At infinity, $\sigma=1/u$ satisfies
$\Dn\sigma=1+1/\sigma$, which has a pole at zero, so there are
no nonconstant solutions reducing to infinity. Together the two
families and the equilibria $0,-1$ give the complete list.
Both families are real exactly for real time constant $C$ or
real amplitude $a$, respectively.
\end{example}

\section{Holomorphic differentials and abelian varieties}
\label{diff:aut:sec:abelian}

\subsection{Holomorphic differentials forbid constant-speed compact flows}
\label{diff:aut:sub:holomorphic}

\begin{proposition}[Infinitesimal contraction]
\label{diff:aut:prop:contraction}
Let $\mathcal Y/\C$ be smooth and projective, of any finite dimension, and let
$\varphi$ be a global regular $1$-form on $\mathcal Y$. For every
$\mathbf y\in\mathcal Y(\SC)$, $\varphi_{\mathbf y}(\Dn\mathbf y)\in\mm_{\C}$.
\end{proposition}
\begin{proof}
Reduce $\mathbf y$ to $p\in\mathcal Y(\C)$ by
\cref{diff:aut:lem:projective}. On an affine neighbourhood of $p$,
choose finitely many regular coordinate functions $t_1,\ldots,t_N$.
After shrinking around $p$ if necessary, write
$\varphi=\sum_j a_j\,dt_j$ with each $a_j$ regular there.
The coordinates $t_j(\mathbf y)$ and coefficients $a_j(\mathbf y)$
are finite: regular functions are ratios with denominators nonzero
at $p$, whose evaluations are units with ordinary nonzero residue.
By \cref{diff:aut:lem:small}, each
$\Dn t_j(\mathbf y)$ is infinitesimal. Consequently
\[
 \varphi_{\mathbf y}(\Dn\mathbf y)
 =\sum_j a_j(\mathbf y)\,\Dn t_j(\mathbf y)\in\mm_\C,
\]
since $\mm_\C$ is an ideal of the finite-element ring.
The algebraic chain rule makes this contraction invariant under
changing the chosen coordinates. No analytic integration or
choice of a global coordinate system is used.
\end{proof}

\begin{corollary}[Holomorphic dual obstruction]
\label{diff:aut:cor:dual}
Let $\mathcal X/\C$ be a smooth projective curve, $0\neq\varphi\in H^{0}(\mathcal X,\Omega^{1})$,
and $\xi$ the meromorphic vector field with $\varphi(\xi)=1$. There is no
nonconstant differential embedding $(\C(\mathcal X),\xi)\hookrightarrow(\SC,\Dn)$.
\end{corollary}
\begin{proof}
A nonconstant solution would give
$1=\varphi_{\mathbf y}(\xi(\mathbf y))=\varphi_{\mathbf y}(\Dn\mathbf y)\in\mm_{\C}$.
\end{proof}

The same obstruction is visible in the divisor of $\xi$.
Locally write $\varphi=\sigma^r a(\sigma)\,d\sigma$ with
$r\geq0$ and $a(0)\neq0$. Then
$\xi=\sigma^{-r}a(\sigma)^{-1}\,\partial/\partial\sigma$ has
order $-r\leq0$: it is nonzero where $\varphi$ is nonzero and
has a pole at a zero of $\varphi$. Thus it has no zeros anywhere,
in agreement with \cref{diff:aut:thm:main}.

\begin{theorem}[Squarefree hyperelliptic obstruction]
\label{diff:aut:thm:hyperelliptic}
Let $Q\in\C[Y]$ be squarefree with $\deg Q\geq3$. Then the solutions in $\SC$ of
\begin{equation}\label{diff:aut:eq:hyperelliptic}
 (\Dn y)^{2}=Q(y)
\end{equation}
are exactly the constants $y=c\in\C$ with $Q(c)=0$.
\end{theorem}
\begin{proof}
Every constant root of $Q$ solves the equation. Suppose instead that
$y$ is nonconstant and put $z=\Dn y\neq0$.
Squarefreeness makes $Q$ a nonsquare in $\C(Y)$, since each root
has odd order one. Hence $Z^2-Q(Y)$ is irreducible, and the
scalar-to-curve reduction in \cref{diff:aut:thm:independence}
identifies $(y,z)$ with a nonconstant differential point on its
smooth projective model.

We check directly that $\varphi=dY/Z$ is a global regular
$1$-form on this model. At finite points with $Z\neq0$ this is
immediate. At a root $a$ of $Q$, squarefreeness gives $Q'(a)\neq0$,
so $Z$ is a local parameter and the identity
$2Z\,dZ=Q'(Y)\,dY$ gives
\[
 \frac{dY}{Z}=\frac{2\,dZ}{Q'(Y)}.
\]
The right side is regular at $(a,0)$, with unit denominator.
These cases cover the affine curve, which is smooth because
$Q$ and $Q'$ have no common root.

At infinity put $w=1/Y$, $v=Z w^{\mathrm g+1}$ and let
$q(w)=w^{\deg Q}Q(1/w)$, so $q(0)\neq0$ is the leading coefficient
of $Q$. If $\deg Q=2\mathrm g+2$, $\mathrm g\geq1$, the local
equation is $v^2=q(w)$. There are two points over infinity,
with $v(0)=\pm\sqrt{q(0)}\neq0$; at each, $w$ is a local
parameter and
\[
 \frac{dY}{Z}=-\frac{w^{\mathrm g-1}}{v}\,dw
\]
is regular. If $\deg Q=2\mathrm g+1$, the equation is
$v^2=wq(w)$. Its unique point at infinity has parameter $v$,
since the derivative in $w$ at $(0,0)$ is $-q(0)\neq0$.
Solving formally gives $w=q(0)^{-1}v^2+O(v^4)$. The same
formula $dY/Z=-w^{\mathrm g-1}dw/v$ is then a unit multiple
of $v^{2\mathrm g-2}\,dv$, again regular. The exponents are
nonnegative precisely in the degrees under consideration.

Thus \cref{diff:aut:prop:contraction} applies to $\varphi$.
But its contraction at the nonconstant point is
$\varphi_{(y,z)}(\Dn(y,z))=\Dn y/z=1$, whereas that proposition
makes it infinitesimal. This contradiction excludes every
nonconstant solution.
\end{proof}

\begin{remark}[A necessary distinction at branch points]
\label{diff:aut:rem:branchpoints}
Differentiating $z^{2}=Q(y)$ with $\Dn y=z\neq0$ gives $\Dn z=Q'(y)/2$; the
cancellation is valid for a nonconstant solution, not for a constant solution
at $z=0$. So \cref{diff:aut:eq:hyperelliptic} and the vector system
$\Dn y=z$, $\Dn z=Q'(y)/2$ have different constant-point behaviour: a simple root
$c$ gives the constant solution $y=c$ of the scalar equation, but $(c,0)$ is not a
stationary point of the vector system. In particular an elliptic translation
system has no points in $\SC$ solving that vector equation, although the squared
scalar equation keeps its branch-point constants.
\end{remark}

\begin{example}[Squarefreeness and degree are real hypotheses]
\label{diff:aut:ex:hypotheses}
For $Q(Y)=Y^{3}$, which is not squarefree, $y=4/(\omega+C)^{2}$ with $C\in\C$ is
nonconstant and satisfies $(\Dn y)^{2}=y^{3}$. For degree two,
$y=(e^{\omega}+e^{-\omega})/2$ satisfies $(\Dn y)^{2}=y^{2}-1$.
Indeed, the first derivative is $-8/(\omega+C)^3$, whose square
is $64/(\omega+C)^6=y^3$; in the second case
$\Dn y=(e^\omega-e^{-\omega})/2$, and subtracting the two
squares gives $1$. Both solutions are nonconstant for every
normalized derivation. Neither hypothesis of
\cref{diff:aut:thm:hyperelliptic} can simply be deleted.
\end{example}

\subsection{The formal logarithm and the identity monad}
\label{diff:aut:sub:formallog}

This subsection concerns a \emph{constant} abelian variety $\mathbf A/\C$ of
dimension $d$, not one whose coefficients vary in the differential field. Let
$\mathbf A_{1}(\SC)$ be the kernel of reduction $\mathbf A(\SC)\to\mathbf A(\C)$,
a group homomorphism by \cref{diff:aut:lem:projective}; the constant inclusion is
a section. The formal completion of $\mathbf A$ at its identity carries a
commutative formal group law
$\mathcal F(T,S)\in\C[[T_{1},\dots,T_{d},S_{1},\dots,S_{d}]]^{d}$,
$\mathcal F(T,S)=T+S+{}$higher terms. The geometric setting of abelian varieties
and tangent spaces is standard \cite[Chapter~I, \S1 and Theorem~6.4]{Milne}; the characteristic-zero logarithm
fact is proved here.

\begin{lemma}[Normalized formal-group logarithm]
\label{diff:aut:lem:formallog}
For a commutative $d$-dimensional formal group over $\C$ there is a unique formal
map $\log_{\mathcal F}(T)=T+O(T^{2})$ with
$\log_{\mathcal F}(\mathcal F(T,S))=\log_{\mathcal F}(T)+\log_{\mathcal F}(S)$. It has
a formal inverse $\exp_{\mathcal F}(T)=T+O(T^{2})$, and its differential is the
normalized invariant differential matrix of the formal group.
\end{lemma}
\begin{proof}
Put $J(T)=\partial_S\mathcal F(T,0)$. Its constant matrix is the
identity, so it is invertible over $\C[[T]]$. Translating the
cotangent basis at zero gives the column of invariant formal forms
\[
 (\varphi_1,\ldots,\varphi_d)^{\mathsf t}=J(T)^{-1}\,dT.
\]
The dual invariant tangent fields commute: their infinitesimal
translations commute because $\mathcal F$ is commutative.
For such fields $v,w$, the identity
$d\varphi_j(v,w)=v(\varphi_j(w))-w(\varphi_j(v))-\varphi_j([v,w])$
has all three terms zero. These fields form a basis over
$\C[[T]]$, so each $\varphi_j$ is closed.

Here is the coefficientwise integration step. If
$\alpha_n=\sum_i a_i(T)\,dT_i$ is closed with all $a_i$
homogeneous of degree $n$, define
\[
 H_{n+1}(T)=\frac{1}{n+1}\sum_i T_i a_i(T).
\]
Closedness gives $\partial_j a_i=\partial_i a_j$; Euler's
homogeneous identity then gives
$\partial_j H_{n+1}=(a_j+\sum_iT_i\partial_i a_j)/(n+1)=a_j$.
Integrating each homogeneous part in this way produces the unique
zero-constant primitive of each $\varphi_j$. The resulting vector
$\ell(T)$ has identity linear part.

Associativity, and commutativity for the first variable, give
\[
 J(\mathcal F(T,S))^{-1}\,d\mathcal F(T,S)
 =J(T)^{-1}\,dT+J(S)^{-1}\,dS.
\]
Thus the total differential in $(T,S)$ of
$\ell(\mathcal F(T,S))-\ell(T)-\ell(S)$ vanishes.
In characteristic zero a formal series with all partial derivatives
zero is constant; here its constant term is zero. This proves the
homomorphism identity and defines $\log_{\mathcal F}=\ell$.
The identity linear part gives the formal inverse
$\exp_{\mathcal F}$ by degree recursion.

For uniqueness, compose another normalized logarithm with
$\exp_{\mathcal F}$. The result is an additive formal map with
identity linear part. For any homogeneous term $H_n$ of degree
$n>1$, additivity gives $2^nH_n(T)=H_n(2T)=2H_n(T)$,
so $H_n=0$. The map is the identity, as required.
\end{proof}

By the smooth-coordinate argument of \cref{diff:aut:lem:projective},
now with $d$ free parameters and an invertible complementary Jacobian,
the formal coordinates give a bijection between
$\mathbf A_1(\SC)$ and $\mm_\C^d$. They describe every point of
the monad, not only points already presented by a chosen series.
All coordinates of $\mathcal F$, $\log_{\mathcal F}$ and
$\exp_{\mathcal F}$ have zero constant term. Their evaluations
therefore remain infinitesimal, and the composition clause of
\cref{diff:aut:lem:evaluation} evaluates both inverse identities
and the group identity on every infinitesimal tuple.

For a local check, the one-dimensional multiplicative formal group
$\mathcal F(T,S)=T+S+TS$ has $J(T)=1+T$ and
\[
 \log_{\mathcal F}(T)=\log(1+T),\qquad
 \exp_{\mathcal F}(U)=e^U-1.
\]
These are ordinary formal series at zero, evaluated only on the
infinitesimal domain. This example illustrates the local lemma;
the later global splitting uses properness of an abelian variety.

A formal coordinate change $\widetilde T=H(T)$ with $H(0)=0$
and ordinary invertible linear part $B$ satisfies
\[
 \log_{\widetilde{\mathcal F}}(H(T))=B\log_{\mathcal F}(T).
\]
Indeed, composition with $\exp_{\mathcal F}$ is an additive
formal map with linear part $B$, and the same homogeneous argument
kills every higher term. Thus the logarithm is intrinsically
Lie-algebra valued; its coordinate vectors change by the ordinary
linear tangent map, not by a nonlinear reparametrization.

\begin{proposition}[Global splitting through the identity monad]
\label{diff:aut:prop:abeliansplit}
In a fixed ordinary basis of $\Lie(\mathbf A)$ there is a group isomorphism
\begin{equation}\label{diff:aut:eq:abeliansplit}
 \mathbf A(\SC)\longrightarrow\mathbf A(\C)\oplus\mm_{\C}^{d},
 \qquad
 \mathbf y\longmapsto\bigl(\st(\mathbf y),\ \log_{\mathcal F}(\mathbf y-\st(\mathbf y))\bigr),
\end{equation}
with inverse $(c,u)\mapsto c+\exp_{\mathcal F}(u)$. Intrinsically the second
summand is $\Lie(\mathbf A)(\mm_{\C})$, independent of the basis.
\end{proposition}
\begin{proof}
Reduction is a homomorphism because the ordinary group law commutes
with reduction, and the inclusion of ordinary points is a section.
Every point has the unique decomposition
$\mathbf y=c+\mathbf q$ with
$c=\st(\mathbf y)\in\mathbf A(\C)$ and
$\mathbf q=\mathbf y-c\in\mathbf A_1(\SC)$.
Uniqueness also follows from
$\mathbf A(\C)\cap\mathbf A_1(\SC)=\{0\}$.
Commutativity makes this a direct-product decomposition of groups.

The evaluated formal logarithm identifies the second factor with
the additive group $\mm_\C^d$. Its inverse sends $u$ to
$\exp_{\mathcal F}(u)$, which has reduction zero. Consequently
$(c,u)\mapsto c+\exp_{\mathcal F}(u)$ is inverse to the displayed
map in both directions. The coordinate-change identity above
identifies its second factor canonically with
$\Lie(\mathbf A)(\mm_\C)$.
\end{proof}

No complex analytic uniformization, no period lattice and no exponential defined
on every surcomplex Lie-algebra element is used; $\exp_{\mathcal F}$ is evaluated
only on $\mm_{\C}^{d}$. In the terms of the companion Tate report this is the
formal-group treatment of a \emph{constant} good-reduction variety over $\SC$,
with reduction by standard part (\cref{diff:aut:rem:tate}).

\subsection{The exact logarithmic-derivative image}
\label{diff:aut:sub:image}

For $\mathbf y\in\mathbf A(\SC)$ its \emph{algebraic logarithmic derivative} is
\begin{equation}\label{diff:aut:eq:dlog}
 \dlog_{\Dn}\mathbf y=d\bigl(\,\cdot\,-\mathbf y\bigr)_{\mathbf y}(\Dn\mathbf y)
 \in\Lie(\mathbf A)(\SC),
\end{equation}
the differential of translation by $-\mathbf y$ applied to $\Dn\mathbf y$;
equivalently its coordinates are the contractions of the normalized invariant
differentials with $\Dn\mathbf y$. It is a tangent vector, not the argument of a
complex number. For $\mathbf A$ replaced by the multiplicative group it would be
$z\mapsto z^{\dagger}$. The class $\Dn\mm_{\C}=\Dn\Ofin_{\C}$ is a $\C$-linear
subspace of $\mm_{\C}$ (\cref{diff:aut:lem:small}), and $\Dn:\mm_{\C}\to\Dn\mm_{\C}$
is a bijection because $\mm_{\C}\cap\C=\{0\}$.

\begin{theorem}[Abelian logarithmic-derivative image]
\label{diff:aut:thm:abelianimage}
For every normalized derivation and every constant abelian variety $\mathbf A/\C$
of dimension $d$,
\begin{equation}\label{diff:aut:eq:abelianimage}
 \dlog_{\Dn}\mathbf y=\Dn\log_{\mathcal F}\bigl(\mathbf y-\st(\mathbf y)\bigr),
 \qquad
 \dlog_{\Dn}\bigl(\mathbf A(\SC)\bigr)=\Lie(\mathbf A)(\Dn\mm_{\C}).
\end{equation}
The kernel is exactly $\mathbf A(\C)$, and in any ordinary Lie basis
\begin{equation}\label{diff:aut:eq:abelianexact}
 0\longrightarrow\mathbf A(\C)\longrightarrow\mathbf A(\SC)
 \xrightarrow{\ \dlog_{\Dn}\ }(\Dn\mm_{\C})^{d}\longrightarrow0
\end{equation}
is split exact, with splitting $a\mapsto\exp_{\mathcal F}\bigl((\Dn|_{\mm_{\C}})^{-1}a\bigr)$.
\end{theorem}
\begin{proof}
For each fixed point $\mathbf y$, the point
$c=\st(\mathbf y)$ is ordinary and has zero derivative.
Translation by $-c$ therefore preserves the logarithmic derivative
and gives $\mathbf q=\mathbf y-c$ in the identity monad.
Writing its formal coordinate tuple as $T(\mathbf q)$, the evaluated
chain rule and the differential formula from
\cref{diff:aut:lem:formallog} give
\[
 \Dn\log_{\mathcal F}(\mathbf q)
 =J(T(\mathbf q))^{-1}\Dn T(\mathbf q)
 =\dlog_{\Dn}\mathbf y.
\]
This differentiates the coordinate tuple at the fixed point; it does
not require differentiating a standard-part map.

Under \cref{diff:aut:prop:abeliansplit}, the logarithmic derivative
is exactly the homomorphism $(c,u)\mapsto\Dn u$.
The restriction $\Dn:\mm_\C\to\Dn\mm_\C$ is a
$\C$-linear bijection: surjectivity onto that image is its
definition, and its kernel is $\mm_\C\cap\C=\{0\}$.
Applied coordinatewise, it proves the asserted image and shows
that the kernel consists precisely of the points with $u=0$,
namely $\mathbf A(\C)$.

Let $P_{\Dn}$ denote this restriction's inverse. Then
$a\mapsto\exp_{\mathcal F}(P_{\Dn}a)$ is a homomorphism,
has logarithmic derivative $a$, and reduces to the identity point.
It is the unique section with image in the identity monad.
In particular, all points over one value $a$ form its translate
by the kernel $\mathbf A(\C)$. No surjectivity of $\Dn$ on
all of $\SC$ was used.
\end{proof}

\begin{remark}[Agreement with, and contrast to, the multiplicative group]
\label{diff:aut:rem:contrast}
For $\Dn=\der$, $\der\mm_{\C}=\Acl+i\Acl=\Acl_{\C}$ in the notation of
\cref{diff:lem:matweyl}, so the image is $\Lie(\mathbf A)(\Acl_{\C})$. The
multiplicative group is not proper, and for it \cref{diff:thm:logimage} constrains
only the imaginary part: $\{z^{\dagger}\}=\No+i\Acl$. For an abelian variety the
\emph{whole} Lie coordinate is constrained, real and imaginary parts alike. The
theorem is not a consequence of the algebraic closedness of $\SC$: it specifies a
generally proper differential image.
\end{remark}

\subsection{A primitive obstruction when primitives exist}
\label{diff:aut:sub:primitive}

For the obstruction construction and finite-primitive criterion, assume
$\Dn$ is surjective, as $\der$ is. By
\cref{diff:eq:split}, applied coordinatewise, every $z\in\SC$ is
$z=\pp(z)+\ct(z)+\varepsilon$ with $\pp(z)\in\Pcl_{\C}$, $\ct(z)\in\C$ and
$\varepsilon\in\mm_{\C}$. For $a\in\SC$ choose $\mathfrak b$ with $\Dn\mathfrak b=a$ and put
\begin{equation}\label{diff:aut:eq:Obsc}
 \Obs^{\C}_{\Dn}(a)=\pp(\mathfrak b)\in\Pcl_{\C}.
\end{equation}
Two primitives differ by an ordinary complex constant, so this is well defined;
it is $\C$-linear and surjective, with $\Obs^{\C}_{\Dn}(\Dn P)=P$ for
$P\in\Pcl_{\C}$. For $\Dn=\der$ it is the componentwise obstruction
$\Obs^{\C}_{\der}(a)=\Obs(\Ree a)+i\Obs(\Imm a)$. It is \emph{not} the surcomplex
phase obstruction $\Phi(a)=\Obs(\Imm a)$ of \cref{diff:def:obs}, which ignores the
real part, and the two must not be conflated: \cref{diff:aut:rem:contrast} is the
reason.

\begin{corollary}[Finite-primitive criterion for abelian motion; surjective $\Dn$]
\label{diff:aut:cor:abeliancriterion}
Let $\Dn$ be surjective and $a\in\Lie(\mathbf A)(\SC)$. The equation
$\dlog_{\Dn}\mathbf y=a$ is solvable in $\mathbf A(\SC)$ if and only if
$\Obs^{\C}_{\Dn}(a)=0$ componentwise in one, equivalently every, ordinary Lie
basis. Then, for a finite primitive vector $\mathfrak b$ with $\Dn\mathfrak b=a$, all solutions are
\begin{equation}\label{diff:aut:eq:abeliansolutions}
 \mathbf y=c+\exp_{\mathcal F}\bigl(\mathfrak b-\st(\mathfrak b)\bigr),\qquad c\in\mathbf A(\C).
\end{equation}
Moreover $\Ker\Obs^{\C}_{\Dn}=\Dn\mm_{\C}$.
\end{corollary}
\begin{proof}
Choose any primitive vector $\mathfrak b$. Vanishing of its purely
infinite part means precisely that it is finite; the remaining
parts are an ordinary constant vector and an infinitesimal vector.
For a finite primitive, put
$u=\mathfrak b-\st(\mathfrak b)\in\mm_\C^d$.
It still satisfies $\Dn u=a$. If $\mathfrak b'$ is another
primitive, then $\mathfrak b'-\mathfrak b\in\C^d$; if one is
finite, so is the other, and subtracting standard parts gives
the same $u$. This is the unique infinitesimal primitive.
Conversely an infinitesimal primitive has zero purely infinite
part. This proves $\Ker\Obs^\C_{\Dn}=\Dn\mm_\C$
coordinatewise, and \cref{diff:aut:thm:abelianimage} gives
solvability and the full translate of its kernel displayed above.

A change of ordinary Lie basis is an invertible matrix over $\C$.
It commutes with $\Dn$, with the $\C$-linear obstruction and,
on finite vectors, with standard part. Thus vanishing in one
basis is equivalent to vanishing in every basis, and the solution
formula represents the same points.
\end{proof}

The standard part in \cref{diff:aut:eq:abeliansolutions} is taken only after the
finiteness of $\mathfrak b$ has been established; ``the standard part'' of an arbitrary
infinite primitive would not be a valid substitute. Surjectivity is used
to define the obstruction on every Lie-algebra input. For an input
with an explicit primitive, the image theorem tests finiteness
directly, without a surjectivity assumption.

\begin{corollary}[Sharp one-power speed threshold; every normalized derivation]
\label{diff:aut:cor:speeds}
Fix $0\neq v\in\Lie(\mathbf A)(\C)$ and $p\in\R$.
For every normalized derivation $\Dn$, the equation
$\dlog_{\Dn}\mathbf y=\omega^{-p}v$ is solvable exactly when
$p>1$. In that case all solutions are
\[
 \mathbf y=c+\exp_{\mathcal F}\!\left(
 \frac{\omega^{1-p}}{1-p}v\right),\qquad c\in\mathbf A(\C).
\]
In particular no nonzero constant speed and no speed
$\omega^{-1}v$ is realizable.
\end{corollary}
\begin{proof}
Exponential compatibility and $\Dn\omega=1$ give
$\Dn\omega^r=r\omega^{r-1}$ for every real $r$, and
$\Dn\log\omega=1/\omega$. Thus an explicit primitive vector is
$\omega^{1-p}v/(1-p)$ for $p\neq1$, and $(\log\omega)v$
for $p=1$, for every normalized derivation.

When $p>1$, the first vector is infinitesimal, so the image theorem
and its splitting give exactly the displayed solutions.
When $p<1$, any coordinate with $v_j\neq0$ is infinite in
modulus; the same is true of $(\log\omega)v_j$ when $p=1$.
Every other primitive differs by an ordinary constant vector,
which cannot make that coordinate finite. There is therefore
no infinitesimal primitive, and \cref{diff:aut:thm:abelianimage}
excludes a solution. This argument uses neither surjectivity nor
a globally defined primitive obstruction.
\end{proof}

This is the abelian analogue of \cref{diff:cor:power}. The integration is
elementary; its force is that it is an exact obstruction for every constant
abelian variety in every finite dimension, and that it constrains a real speed
$\omega^{-p}v$ as well as an imaginary one, which the multiplicative group does not.

\section{What the autonomous theory settles, and what it does not}
\label{diff:aut:sec:limits}

\begin{remark}[A finite geometric decision criterion]
\label{diff:aut:rem:decision}
For a presented smooth constant projective curve and a nonzero
meromorphic vector field, the divisor of that field has finite support.
Thus \cref{diff:aut:thm:main} reduces existence of a nonconstant
surcomplex trajectory to a finite test: find the zeros, compute their
orders, and at each simple zero compute the linear coefficient.
A multiple zero contributes solutions; a simple zero contributes
solutions exactly when its coefficient is real and negative.
The classification then specifies all solutions as Hahn evaluations.

To call this an effective procedure, the presentation must supply
algorithms for that zero divisor, the local orders and leading
coefficients, and exact tests for vanishing imaginary part and real
sign. Exact coefficient-field arithmetic alone does not supply
these geometric operations. Algebraic complex presentations equipped
with polynomial factorization and local-expansion algorithms provide
a sufficient setting. These requirements are a contract for a
procedure, not an implementation in the repository.

The output is a finite list of local families with their ordinary
parameters and formal recursions. An infinite series is represented
by that defining data, not by listing every coefficient. A chosen
parameter must itself have a representation if coefficients are to
be computed effectively. No unconditional computability claim is
made for arbitrary unencoded complex constants. In this qualified
sense the criterion lifts item~N25 for the stated class; beyond it
N25 stands.
\end{remark}

\begin{remark}[Why ordinary formal initial-value results are not enough]
\label{diff:aut:rem:formalivp}
At an ordinary finite time a vector field that is nonzero at a point has a
nonconstant local solution through it. In $\SC$ every projective point has an
ordinary reduction and its finite coordinates have infinitesimal derivatives, so
a nonzero constant velocity at the reduction is impossible: admissible
trajectories are centred at zeros, and their rate of approach is constrained by
the differential structure. This explains the apparent reversal between the local
formal picture and \cref{diff:aut:thm:main} without conflating their time
parameters, and it is the nonlinear form of \cref{diff:rem:noclock}.
For example, the ordinary formal equation $dy/ds=1$ has the
solution $y=s$ centred at zero. The intrinsic equation $\Dn y=1$
in $\SC$ instead has exactly $y=\omega+C$, $C\in\C$.
These points reduce to infinity in $\mathbb P^1$. There the local
parameter $\sigma=1/y$ satisfies $\Dn\sigma=-\sigma^2$,
a double zero of the projective vector field, just as the
classification requires. Substituting an infinitesimal for the
ordinary time $s$ cannot preserve $ds/ds=1$ under $\Dn$, since
its intrinsic derivative is infinitesimal.
\end{remark}

\begin{remark}[Established inputs and proposed contributions]
\label{diff:aut:rem:inputs}
The published construction of a surreal derivation is an input. The theory of
smooth curves, formal implicit equations, invariant differentials and
characteristic-zero formal-group logarithms is classical; the local formal normal
forms are \cite[Lemma~6.1]{NPT}, and Rosenlicht's independence theorem is imported
from \cite[Lemma~7.2, Proposition~7.3]{NPT}, whose numbered statements are cited
from the consulted arXiv version~1 and are not assumed to match the journal
numbering. The scalar imaginary obstruction overlaps \cref{diff:part:scalar}.
Member~11 proposes as new only the exhaustive ambient realization theorem with its
finite-scale support bound, the derivation independence of the constant
first-order solution sets, the projective differential-form obstruction and the
exact abelian logarithmic-derivative image; historical novelty is a separate
question (item~N93).
\end{remark}

\begin{remark}[Boundaries]
\label{diff:aut:rem:boundaries}
The curve theorem does not classify higher-dimensional autonomous systems:
\cref{diff:aut:prop:contraction} survives in every dimension, but near a singular
point several eigendirections interact, and integer relations among the
eigenvalues of the linear part together with noncommuting normal-form terms
prevent the one-coordinate time primitive from proving a
general classification. The abelian theorem assumes a constant variety;
nonconstant abelian varieties can acquire a connection term after coordinate
changes, which the constant-section splitting does not address, and
noncommutative algebraic groups need another argument because their invariant
forms need not be closed, which is where the proof of
\cref{diff:aut:lem:formallog} deliberately stops.
For instance, the affine group with law
$(a,b)(a',b')=(aa',b+ab')$ has left-invariant forms
$da/a$ and $db/a$. Near its identity put $a=1+x$, $b=y$.
The second form is $dy/(1+x)$, and
\[
 d\!\left(\frac{dy}{1+x}\right)
 =-\frac{dx\wedge dy}{(1+x)^2}\neq0.
\]
Thus its invariant forms cannot all be differentials of formal
logarithm coordinates with additive group law. This pinpoints the
commutativity step needed above; it is not a classification of
noncommutative surcomplex differential equations.
\end{remark}

\begin{remark}[Foundational scope and a formalization route]
\label{diff:aut:rem:formalization}
$\mathcal X(\SC)$, $\mathbf A(\SC)$ and the exact sequences over them are
class-level statements, read in a conservative class theory such as NBG with the
usual pointwise algebraic definitions (\cref{diff:sub:classes}); no collection of
all surreal supports is used as a set. Each formal evaluation uses a set-sized
support, \cref{diff:aut:thm:scales} supplies one set-sized field for the whole
solution set of a fixed equation, and the imported independence theorem is
applied in that field; for the abelian theorem each formal computation uses
finitely many input series and the global sequence is stated pointwise.
A dependency-ordered formalization uses four interfaces:
\begin{enumerate}[label=\textup{(\roman*)},leftmargin=2.4em]
\item a valued differential field with ordinary residue section,
constant field $\C$, small derivative on finite elements and a
specified element of derivative $1$, together with the exponential,
order and strong-additivity laws of \cref{diff:aut:def:normalized};
\item finite-variable strong formal evaluation, zero-constant
composition and the chain rule for that intrinsic derivation;
\item smooth algebraic local coordinates identifying the entire
residue monad with infinitesimal tuples, with projective reduction;
\item the one-variable time primitive, forced scale and
constant-field uniqueness argument, followed by commutative formal
groups and invariant differentials for the abelian extension.
\end{enumerate}
A Hahn identity alone supplies neither the geometric coordinates
nor the summability bridge needed to evaluate them.

The current Lean library already proves actual real and complex
finite-variable evaluation, strong summability, the ring laws and
zero-constant composition. The live ledger
\path{docs/FORMALIZATION.md} records this reusable algebraic portion
of \cref{diff:aut:lem:evaluation}. It does not yet prove that lemma's
chain rule for $\Dn$, the full projective-monad statement
\cref{diff:aut:lem:projective}, or the autonomous classification.
Checked fine-differentiation rules concern a different derivative
and cannot fill the intrinsic-derivation interface by a change of
notation. The source package supplied no Lean file; that historical
fact is distinct from current partial library coverage
(\cref{diff:warn:noformalization}).
\end{remark}

\begin{remark}[Relation to other reports]
\label{diff:aut:rem:tate}
The companion report on Hahn--Tate uniformization notes, among its limitations,
that curves with nonnegative valuation of $j$ need a different treatment, such as
formal groups around a good-reduction model and a residue elliptic curve
\cite{RepoTate}. \Cref{diff:aut:prop:abeliansplit} does this for \emph{constant}
abelian varieties over $\SC$, with reduction by standard part; it does not do it
over that report's Hahn fields with a coarse valuation, so this is a cross-link
and not an answer. The Tate construction and the present one use
different residue maps; identifying their formal-group vocabulary
does not identify their hypotheses or fields. The Tate result
itself imposes no Berarducci--Mantova differential equation, whereas
this part determines an intrinsic logarithmic-derivative image.
Likewise, ``autonomous flow'' here means a solution of an intrinsic
differential equation. The dynamics report's embedding of a
near-identity map into a formal flow is a statement about formal
iteration and its infinitesimal generator. These are different
constructions, even when they use some of the same formal calculus.
\end{remark}

\begin{table}[htbp]
\centering\small
\caption{Proof dependencies of \cref{diff:aut:part}, after member~11's audit.}
\label{diff:aut:tab:dependencies}
\begin{tabularx}{\textwidth}{@{}L{0.25\textwidth}Y@{}}
\toprule
Conclusion & Necessary inputs, and the check against overreach \\
\midrule
Finite derivative (\cref{diff:aut:lem:small}) & $\Dn\omega=1$, the sign rule and
finite coordinates. No surjectivity. \\
Formal evaluation (\cref{diff:aut:lem:evaluation}) & Strong Hahn summability and
the ordered-word lemma. No fine-topological limit. \\
Projective reduction (\cref{diff:aut:lem:projective}) & Finitely many homogeneous
coordinates, standard part, smooth local implicit equations; poles checked in
local coordinates. \\
Simple zeros (\cref{diff:aut:prop:simple}) & Formal linearization, constant field
$\C$, \cref{diff:aut:lem:eigen}. A negative real part alone is insufficient. \\
Multiple zeros (\cref{diff:aut:prop:multiple}) & A time primitive, the constant-field
identity $\mathfrak t(u)-\omega\in\C$, a forced leading scale, implicit unit
uniqueness. Formal construction alone is insufficient. \\
Finite-scale field (\cref{diff:aut:thm:scales}) & All local families and the
three-shift formula. Finite monomial generation is not finite transcendence
degree. \\
Derivation independence (\cref{diff:aut:thm:independence}) & Constant coefficients
and the normalization. Variable coefficients and higher order are not covered. \\
Hyperelliptic exclusion (\cref{diff:aut:thm:hyperelliptic}) & Squarefreeness,
degree at least three, holomorphy of $dY/Z$. Branch-point constants are
retained. \\
Abelian image (\cref{diff:aut:thm:abelianimage}) & A constant abelian variety, its
residue section, the commutative formal logarithm on the identity monad. No
global uniformization. \\
Primitive obstruction (\cref{diff:aut:cor:abeliancriterion}) & In addition,
surjectivity of $\Dn$. Standard part only of finite primitive vectors. \\
Power-speed threshold (\cref{diff:aut:cor:speeds}) & Explicit power or
logarithmic primitives and the abelian image theorem; no surjectivity
assumption. \\
\bottomrule
\end{tabularx}
\end{table}

\part{Transfinite critical potentials: second-order equations at every logarithmic depth}
\label{diff:cp:part}

This part is the contribution of member~13 (\cref{diff:app:provenance}). It
treats the scalar second-order equations
\[
 \der^{2}y+\bigl(Q_{\alpha}+c\,(\ell_{\alpha}^{\dagger})^{2}\bigr)y=0,
 \qquad
 Q_{\alpha}=\frac14\sum_{\beta<\alpha}(\ell_{\beta}^{\dagger})^{2},
\]
where $\alpha$ is an arbitrary ordinal, $c$ is an ordinary real number,
$\ell_{\alpha}=\omega^{\omega^{-\alpha}}$ is the log-atomic tower that continues
the iterated logarithms of \cref{diff:eq:logtower} through every ordinal, and
$\ell^{\dagger}=\der\ell/\ell$. The unknown ranges over all of $\No$ or $\SC$.
For every ordinal the solution space is two-dimensional when $c\leq\frac14$ and
zero when $c>\frac14$ (\cref{diff:cp:thm:main}); no set of earlier logarithmic
scales decides which case occurs (\cref{diff:cp:thm:invisible}); and at every
nonzero limit ordinal three nested, $\der$-stable, real closed Hahn fields that
contain the same potential have solution dimensions $0$, $1$ and $2$, the
solution field over the first having the upper triangular Borel subgroup of
$\SL_{2}(\C)$ as differential Galois group
(\cref{diff:cp:thm:fieldjump,diff:cp:thm:pv}). Four relations to the rest of the
report fix its place.
\begin{itemize}[leftmargin=1.6em]
\item \emph{The missing oscillator.} The supercritical half is the
second-order statement of \cref{diff:cor:oscillator} read at another scale.
\Cref{diff:cp:rem:oscillator} makes this precise: one gauge identity carries
$\der^{2}y+b^{2}y=0$ to the Euler operator at the scale $e^{\omega}$, and the
family to the same Euler operator at the scale
$\ell_{\alpha}$; in both cases the obstruction is an infinite accumulated
phase, $\abs b\,\omega$ in the first and $\sqrt{c-\frac14}\,\ell_{\alpha+1}$ in
the second. Neither statement contains the other.
\item \emph{First-order thresholds and \cref{diff:rs:part}.} The coefficient
$\ell_{n}^{\dagger}$ is the boundary case $p=1$ of \cref{diff:cor:logscales},
and \cref{diff:cp:cor:nonreal} carries that obstructed boundary to every
ordinal depth. Stage $0$ is a power-Hahn regular-singular system, and its
classification in \cref{diff:rs:part} agrees with the one here.
For $\alpha\geq1$ the potential lies outside $\KR$ except at
$(\alpha,c)=(1,0)$, which is the previous stage's critical endpoint
$Q_{1,0}=Q_{0,1/4}$. The Conway exponents outside $\R$ in the other
potentials put them beyond that coefficient field
(\cref{diff:cp:rem:stagezero}).
\item \emph{\Cref{diff:part:matrix}.} For a real coefficient,
\cref{diff:thm:matpinney} excludes a one-dimensional solution space in $\No$.
The family confirms this without that part's import, exhibits the
one-dimensional case inside a set-sized subfield exactly where a primitive is
missing (\cref{diff:cp:thm:fieldjump,diff:cp:prop:gaps}), and has an explicit
canonical amplitude and phase (\cref{diff:cp:rem:pinney}).
\item \emph{Workspaces.} The fields of \cref{diff:cp:sec:fields} are full Hahn
fields $K_{\Gamma}$ in the sense of \cref{diff:sec:localization}, of higher
rank and with logarithmic monomials. Like \cref{diff:ex:sqrt2,diff:ex:logmissing}
they separate closure properties, here with algebraic closure, derivative
stability and closure under logarithms all present while an exponential and
then a primitive are missing (\cref{diff:cp:rem:workspaces}).
\end{itemize}
Every statement of this part is relative to the one fixed $\der$
(\cref{diff:warn:nouniqueness}). The main construction uses the clauses of
\cref{diff:thm:BM} other than \ref{diff:BM:surj}, because the primitives it
needs are supplied by the tower itself; surjectivity enters only the real Riccati
remark (\cref{diff:cp:rem:riccati}) and the comparison with
\cref{diff:part:matrix} (\cref{diff:cp:rem:pinney}), which also uses that
part's import. The part adds one input that \cref{diff:thm:BM} does not list:
the ordinal log-atomic formulas of \cite{ADH}, imported with their limit cases
in \cref{diff:cp:prop:tower}.

\section{The log-atomic tower, and the alphabet of this part}
\label{diff:cp:sec:tower}

\begin{convention}[The alphabet of this part]
\label{diff:cp:conv:alphabet}
Throughout \cref{diff:cp:part}:
\begin{itemize}[leftmargin=1.6em]
\item Ordinals are written $\alpha,\beta,\zeta$, and $\lambda$ always denotes a
\emph{nonzero limit ordinal}; in this part $\lambda$ is never an eigenvalue.
An ordinal is regarded as a surreal number, and $\omega^{x}$ is the Conway
monomial map of \cref{diff:eq:normal}, with \cref{diff:warn:monomialmap} in
force. Exponent-group elements are written $\gamma$. Member~13's ordinal index
$\eta$ is $\zeta$ here.
\item $\ell_{\alpha}=\omega^{\omega^{-\alpha}}$,
$\varpi_{\alpha}=\sum_{\beta<\alpha}\omega^{-\beta}$,
$\ell_{<\alpha}=\omega^{\varpi_{\alpha}}$ and
$u_{\alpha}=\omega^{\varpi_{\alpha}/2}$. Member~13 writes $p_{\alpha}$,
$s_{\alpha}$, $P_{\alpha}$, $U_{\alpha}$ for $\omega^{-\alpha}$,
$\varpi_{\alpha}$, $\ell_{<\alpha}$, $u_{\alpha}$: here $p$ is the real
threshold exponent of \cref{diff:cor:power,diff:cor:logscales}, $s$ is a scale
as in $\der_{s}$, the $P_{j}$ are phases, and $U$ is an ordinary complex domain.
\item Logarithmic derivatives are written with $\dagger$. Member~13's
$a_{\alpha}$, $A_{\alpha}$ and $h_{\alpha}$ are $\ell_{\alpha}^{\dagger}$,
$\ell_{<\alpha}^{\dagger}$ and $u_{\alpha}^{\dagger}=\frac12\ell_{<\alpha}^{\dagger}$
(\cref{diff:cp:lem:triangular}); $a$ is a real coefficient and $A$ a matrix here.
\item Potentials are written $Q$: $Q_{\alpha}$ as above (\cref{diff:cp:lem:triangular}),
$Q_{\alpha,c}=Q_{\alpha}+c(\ell_{\alpha}^{\dagger})^{2}$, and a general
potential $Q\in\SC$. Member~13 writes $q_{\alpha,c}$ and $q$; in this report
$q$ is the rank-one coefficient of \cref{diff:thm:phasecriterion}.
\item A scale is written $s$, with the rescaled derivation $\der_{s}$ of
\cref{diff:tab:derivatives}; the Euler operator is $s\der_{s}$, which is
$\der_{\log s}$ when $s>0$. Member~13 writes $T$, $D_{T}$ and $E_{T}$; the
sans-serif $\mathsf T$ is an adjoined generator here, never a number. Its
partner $U$ with $U^{2}T'=1$ is $u$ here.
\item $\mathrm{Wr}(y_{1},y_{2})=y_{1}\der y_{2}-y_{2}\der y_{1}$, as in
\cref{diff:eq:matpinneyclassical}, and $\mathrm{Wr}_{s}$ is the same expression
for $\der_{s}$; member~13 writes $W$, which is a unitary matrix here. For a
differential subfield $F\subseteq\SC$ and $Q\in F$,
$\Sol_{F}(Q)=\{y\in F:\der^{2}y+Qy=0\}$.
\item At a limit ordinal the exponent groups are $\Delta_{0}\subset
\Delta_{1}\subset\Delta_{2}$ and the fields are
$H_{\Delta_{j}}=\R((t^{\Delta_{j}}))$ and $K_{\Delta_{j}}=\C((t^{\Delta_{j}}))$,
as in \cref{diff:tab:notation}. Member~13 writes $\Gamma_{j}$, $K_{j}$ and
$k_{j}$; the report's $K_{0}$ is $\C((t))$, and $\Gamma_{n}$ are the stages of
\cref{diff:thm:hull}.
\item $Y$ is a fundamental matrix, $\psi_{a,b}$ a differential automorphism,
$\Bor_{2}(\C)$ the upper triangular Borel subgroup of $\SL_{2}(\C)$, and
$\Sch(s)$ the Schwarzian. Member~13 writes $\Phi$, $\sigma_{a,b}$, $B(\C)$ and
$\mathcal S_{\der}$; here $\Phi$ is the phase map, $\sigma$ the local parameter
of \cref{diff:aut:part}, $B$ an accumulated phase and $\Strip$ the strip. The
word ``Borel'' in this part names that algebraic subgroup; it is unrelated to
the Borel summation disclaimed in item~N37.
\item Member~13's real infinitesimal primitive $\theta$ of an angular rate is
the infinitesimal phase residue $\eta$ of \cref{diff:conv:phases}(b), and its
Riccati solution $b$ is $w$ here.
\item For $g\neq0$, $f\prec g$ means $f/g\in\mm_{\C}$, and $f\sim g$ means
$f/g-1\in\mm_{\C}$, as in \cref{diff:aut:part}.
\end{itemize}
\end{convention}

\begin{warning}[Member~13's $\mathfrak m$ is not this report's $\mm$]
\label{diff:cp:warn:frakm}
Member~13 writes $\mathfrak m(x)=\omega^{x}$ for the Conway monomial map, and
$\mathfrak i$ and $\mathcal O$ for the infinitesimals and the finite surreals.
In this report $\mm$ is the ideal of infinitesimals and never takes an argument,
$\Ofin$ is the class of finite surreals, and the monomial map is written
$\omega^{x}$ or, in increasing orientation, $t^{-x}$ (\cref{diff:eq:tmonomial}).
Every formula of the member has been transcribed accordingly. A reader moving
between the two texts should translate $\mathfrak m(\,\cdot\,)$ before reading
anything else.
\end{warning}

\begin{proposition}[The ordinal log-atomic tower: imported]
\label{diff:cp:prop:tower}
For every ordinal $\alpha$ the families $(\omega^{-\beta})_{\beta<\alpha}$ and
$(\ell_{\beta+1})_{\beta<\alpha}$ are summable, $\ell_{\alpha}$ is positive
infinite and log-atomic, and
\begin{equation}\label{diff:cp:eq:tower}
 \log\ell_{\alpha}=\ell_{\alpha+1},\qquad
 \der\ell_{\alpha}=\ell_{<\alpha}^{-1},\qquad
 \ell_{<\alpha}=\omega^{\varpi_{\alpha}}
 =\exp\Bigl(\sum_{\beta<\alpha}\ell_{\beta+1}\Bigr).
\end{equation}
For every real $r$,
$\ell_{\alpha}^{r}=\exp(r\ell_{\alpha+1})=\omega^{r\omega^{-\alpha}}$ and
$\ell_{<\alpha}^{r}=\omega^{r\varpi_{\alpha}}$.
\end{proposition}

These are imported, not proved here: \cite[Lemma~2.3, Proposition~2.5,
Lemmas~2.6--2.8 and the pre-derivation formula on p.~10]{ADH}, where the notation is
$\lambda_{-\alpha}=\log_{\alpha}\omega$, and for the real powers the
real-linear exponent formula in the proof of \cite[Lemma~2.3]{ADH}; the
numbering is that of arXiv version~3, checked for this review. For finite
$n$, $\ell_{n}$ is the $n$-fold logarithm of $\omega$ and
$\ell_{<n}=\ell_{0}\cdots\ell_{n-1}$, so \cref{diff:cp:eq:tower} returns
\cref{diff:eq:logtower}; $\ell_{0}=\omega$ and $\ell_{1}=\log\omega$. At an
infinite limit $\lambda$, $\ell_{\lambda}$ is the simplest positive infinite
surreal below every $\ell_{\beta}$ with $\beta<\lambda$. The monomial
$\ell_{<\alpha}$ is the monomial product $\prod_{\beta<\alpha}\ell_{\beta}$, and
its two readings --- sum the exponents $\omega^{-\beta}$ in normal form, or
exponentiate the sum of the logarithms $\ell_{\beta+1}$ --- agree by the
monomial-product logarithm identity of \cite[Lemma~2.3]{ADH}. At a limit
ordinal that identity is an imported fact; replacing it by ordinary
convergence of finite products would be invalid (\cref{diff:warn:summability}).
Likewise the limit-stage derivative formula is a fact about this derivation,
not the result of formally differentiating a finite iterate, and the
normalization $\der\omega=1$ alone does not determine it
(\cref{diff:cp:rem:derivation}). The displayed real-power identities do
not define $\omega^{z}$ for non-real $z$, or a complex power of a
positive infinite scale; \cref{diff:cp:cor:nonreal} will give the
relevant differential obstruction. The real powers rest on
two facts from \cite[Section~2]{ADH}: the logarithm of a
monomial is purely infinite, and the correspondence between the exponent of a
monomial and its logarithm is real-linear. Since the monomial map is
multiplicative, $u_{\alpha}^{2}=\ell_{<\alpha}$.
The empty-index conventions are
$\varpi_{0}=0$, $\ell_{<0}=u_{0}=1$ and $\der\ell_{0}=1$.
Every displayed family is indexed by the set of ordinals below one
fixed $\alpha$; none is summed over the class of all ordinals.

\begin{remark}[What the notation does not assert]
\label{diff:cp:rem:notation}
The symbols $\ell_{\alpha}$ designate particular log-atomic numbers. They do
not assert a globally defined transfinite iteration or composition of $\log$
on $\No$ (\cref{diff:rem:nocomposition}), and \cref{diff:cp:prop:noend} shows
that the tower has no final stage inside $\No$. The remark after
\cref{diff:cor:logscales} --- ``no transfinite logarithm tower is used'' ---
remains true of \cref{diff:part:scalar}; this part imports that tower.
\end{remark}

\begin{lemma}[Separation of scales]
\label{diff:cp:lem:scales}
The families $(\omega^{-\beta})_{\beta<\alpha}$ and
$(\ell_{\beta+1})_{\beta<\alpha}$ are strictly decreasing, and
$\varpi_{\beta}<\varpi_{\zeta}$ for $\beta<\zeta$. Consequently
$\ell_{\beta}^{\dagger}=\omega^{-\varpi_{\beta+1}}=\ell_{<\beta+1}^{-1}=\der\ell_{\beta+1}$,
\begin{equation}\label{diff:cp:eq:ratio}
 \frac{\ell_{\beta+1}^{\dagger}}{\ell_{\beta}^{\dagger}}=\frac{1}{\ell_{\beta+1}},
 \qquad
 \ell_{\zeta}^{\dagger}\prec\ell_{\beta}^{\dagger}\quad(\beta<\zeta),
\end{equation}
and at a limit $\lambda$, $\ell_{\lambda}\prec\ell_{\beta}$ for every
$\beta<\lambda$.
\end{lemma}
\begin{proof}
The monomial map is strictly increasing and $-\zeta<-\beta$, which gives the
first assertion, including the decrease of
$\ell_{\beta+1}=\omega^{\omega^{-(\beta+1)}}$.
For $\beta<\zeta$ the difference
$\varpi_{\zeta}-\varpi_{\beta}=\sum_{\beta\leq\xi<\zeta}\omega^{-\xi}$
has leading term $\omega^{-\beta}$ with coefficient $1$; hence it
is positive and $\varpi$ increases. By
\cref{diff:cp:eq:tower}, $\ell_{\beta}^{\dagger}=\ell_{<\beta}^{-1}\ell_{\beta}^{-1}
=\omega^{-\varpi_{\beta}-\omega^{-\beta}}=\omega^{-\varpi_{\beta+1}}$, and
$\der\ell_{\beta+1}=\der\log\ell_{\beta}=\ell_{\beta}^{\dagger}$. Distinct
monomials are Archimedean separated, and
$\varpi_{\beta+2}-\varpi_{\beta+1}=\omega^{-(\beta+1)}$ gives the ratio.
Finally
$\ell_{\lambda}/\ell_{\beta}
=\omega^{\omega^{-\lambda}-\omega^{-\beta}}$ has strictly negative
Conway exponent, so it is infinitesimal. The same argument works
for any $\beta<\zeta$, not only at a limit.
\end{proof}

For finite $n$, $\ell_{n}^{\dagger}=1/(\ell_{0}\ell_{1}\cdots\ell_{n})$ is the
coefficient $b_{n,1}$ of \cref{diff:cor:logscales}: the obstructed boundary
case $p=1$ at depth $n$.

\section{The critical identity at every ordinal}
\label{diff:cp:sec:support}

\begin{lemma}[Summation and triangular differentiation]
\label{diff:cp:lem:triangular}
For every ordinal $\alpha$ the families $(\ell_{\beta}^{\dagger})_{\beta<\alpha}$,
$((\ell_{\beta}^{\dagger})^{2})_{\beta<\alpha}$ and
$(\ell_{\beta}^{\dagger}\ell_{\zeta}^{\dagger})_{\beta,\zeta<\alpha}$ are summable,
and
\begin{align}
 \ell_{<\alpha}^{\dagger}&=\sum_{\beta<\alpha}\ell_{\beta}^{\dagger},
 \label{diff:cp:eq:Adef}\\
 \der\bigl(\ell_{\beta}^{\dagger}\bigr)&=-\ell_{\beta}^{\dagger}\,\ell_{<\beta+1}^{\dagger}
 \qquad\text{for every ordinal }\beta,\label{diff:cp:eq:ader}\\
 \der\bigl(\ell_{<\alpha}^{\dagger}\bigr)&=-\frac12\Bigl((\ell_{<\alpha}^{\dagger})^{2}
 +\sum_{\beta<\alpha}(\ell_{\beta}^{\dagger})^{2}\Bigr).\label{diff:cp:eq:Ader}
\end{align}
In particular $Q_{\alpha}$ is a well-defined surreal, $Q_{0}=0$, and
$Q_{\alpha}>0$ for $\alpha\geq1$.
\end{lemma}
\begin{proof}
In the orientation of \cref{diff:eq:tmonomial},
$\ell_{\beta}^{\dagger}=t^{\varpi_{\beta+1}}$ with exponents strictly increasing
in $\beta$ (\cref{diff:cp:lem:scales}); an ordinal-indexed increasing family of
exponents is well ordered and each exponent occurs once, so the first two
families are summable (\cref{diff:def:summable}); doubling the
exponents preserves their strict increase. Summability of
$(\ell_{\beta+1})_{\beta<\alpha}$ was supplied by
\cref{diff:cp:prop:tower}.

Here the double family's support and finite incidence can also be
seen directly. For $\beta\leq\zeta$ its exponent is
\[
 \varpi_{\beta+1}+\varpi_{\zeta+1}
 =2\sum_{\xi\leq\beta}\omega^{-\xi}
   +\sum_{\beta<\xi\leq\zeta}\omega^{-\xi}.
\]
The first place where these coefficient strings differ shows that
these exponents increase in the lexicographic order on
$(\beta,\zeta)$ with $\beta\leq\zeta<\alpha$. This is a well-order:
choose the least first coordinate and then the least second
coordinate in any nonempty subset. The coefficient string also
recovers the two indices uniquely. Thus each exponent receives
one ordered pair on the diagonal and two off it. This proves both
conditions for summability and licenses restrictions to either
triangle, multiplication of the two sums and regrouping. It is
the Hahn product theorem in this particular support
\cite[Section~2.4]{BM}.

Since $\ell_{<\alpha}=\exp(\sum_{\beta<\alpha}\ell_{\beta+1})$,
clauses \ref{diff:BM:exp} and \ref{diff:BM:additive} give
$\ell_{<\alpha}^{\dagger}=\sum_{\beta<\alpha}\der\ell_{\beta+1}
=\sum_{\beta<\alpha}\ell_{\beta}^{\dagger}$, which is \cref{diff:cp:eq:Adef}.
As $\ell_{\beta}^{\dagger}=\ell_{<\beta+1}^{-1}$, this gives
\cref{diff:cp:eq:ader}. Differentiating \cref{diff:cp:eq:Adef} by
\ref{diff:BM:additive},
\[
 \der\bigl(\ell_{<\alpha}^{\dagger}\bigr)
 =-\sum_{\beta<\alpha}\sum_{\zeta\leq\beta}\ell_{\beta}^{\dagger}\ell_{\zeta}^{\dagger},
\]
a subfamily of the summable double family. The square
$(\ell_{<\alpha}^{\dagger})^{2}$ counts each off-diagonal unordered pair twice and
each diagonal pair once, so twice the triangular sum is
$(\ell_{<\alpha}^{\dagger})^{2}+\sum_{\beta<\alpha}(\ell_{\beta}^{\dagger})^{2}$;
all regroupings are licensed by the joint summability just proved
(\cref{diff:warn:summability}). For $\alpha\geq1$ the leading term of
$Q_{\alpha}$ is $\frac14(\ell_{0}^{\dagger})^{2}=\frac14\omega^{-2}>0$.
\end{proof}

\begin{theorem}[Exact critical solutions and Wronskian]
\label{diff:cp:thm:critical}
For every ordinal $\alpha$,
\begin{equation}\label{diff:cp:eq:critical}
 u_{\alpha}^{\dagger}=\frac12\,\ell_{<\alpha}^{\dagger},
 \qquad
 \der^{2}u_{\alpha}+Q_{\alpha}u_{\alpha}=0,
 \qquad
 u_{\alpha}^{2}\,\der\ell_{\alpha}=1 .
\end{equation}
Both $u_{\alpha}$ and $u_{\alpha}\ell_{\alpha}$ solve $\der^{2}y+Q_{\alpha}y=0$,
and $\mathrm{Wr}(u_{\alpha},u_{\alpha}\ell_{\alpha})=1$.
\end{theorem}
\begin{proof}
The first identity is the derivative of
$\log u_{\alpha}=\frac12\sum_{\beta<\alpha}\ell_{\beta+1}$. Hence, by
\cref{diff:cp:eq:Ader},
\[
 \frac{\der^{2}u_{\alpha}}{u_{\alpha}}
 =\der\bigl(u_{\alpha}^{\dagger}\bigr)+\bigl(u_{\alpha}^{\dagger}\bigr)^{2}
 =\frac12\der\bigl(\ell_{<\alpha}^{\dagger}\bigr)+\frac14(\ell_{<\alpha}^{\dagger})^{2}
 =-\frac14\sum_{\beta<\alpha}(\ell_{\beta}^{\dagger})^{2}=-Q_{\alpha}.
\]
The third identity is $\ell_{<\alpha}\,\der\ell_{\alpha}=1$. Write $u=u_{\alpha}$
and $s=\ell_{\alpha}$. Differentiating $u^{2}\der s=1$ gives
$2\der u\,\der s+u\,\der^{2}s=0$, so
$\der^{2}(us)+Q_{\alpha}us=s(\der^{2}u+Q_{\alpha}u)+2\der u\,\der s+u\,\der^{2}s=0$,
and $\mathrm{Wr}(u,us)=u^{2}\der s=1$.
\end{proof}

\begin{remark}[A limit stage is not a missing induction step]
\label{diff:cp:rem:limitstage}
At a successor, $\ell_{<\alpha+1}=\ell_{<\alpha}\ell_{\alpha}$ and
$Q_{\alpha+1}=Q_{\alpha}+\frac14(\ell_{\alpha}^{\dagger})^{2}$. At a limit, the
same identity for $\der^{2}u_{\alpha}$ comes from the single summable
triangular calculation of \cref{diff:cp:lem:triangular}. No derivative of a
nonconvergent net and no interchange of an uncontrolled double series occurs,
and $Q_{\lambda}$ is defined by the Hahn sum of its increments,
without assuming topological convergence of the partial sums
$Q_{\beta}$. Summing those partial sums would be a different,
invalid operation: at every nonzero limit $\lambda$, the family
$(Q_{\beta})_{0<\beta<\lambda}$ has infinitely many contributions
$\frac14\omega^{-2}$ at the same exponent and is not summable.
\end{remark}

\begin{lemma}[A fundamental pair spans every solution]
\label{diff:cp:lem:span}
Let $F$ be a differential field of characteristic zero with constants $C_{F}$,
let $Q\in F$, and let $y_{1},y_{2}\in F$ solve $\der^{2}y+Qy=0$ with
$\mathrm{Wr}(y_{1},y_{2})\neq0$. Then $\mathrm{Wr}(y_{1},y_{2})\in C_{F}$, and
every solution in $F$ has a unique expression in
$C_{F}y_{1}+C_{F}y_{2}$. Thus this solution space has dimension two
over $C_{F}$.
\end{lemma}
\begin{proof}
For any two solutions $v,w$,
$\der\mathrm{Wr}(v,w)=v\der^{2}w-w\der^{2}v=-Qvw+Qwv=0$.
In particular $\mathcal W=\mathrm{Wr}(y_{1},y_{2})$ is a nonzero constant.
For a third solution $y$, solve the two linear equations
\[
 \begin{pmatrix}y_{1}&y_{2}\\ \der y_{1}&\der y_{2}\end{pmatrix}
 \begin{pmatrix}a\\b\end{pmatrix}
 =\begin{pmatrix}y\\\der y\end{pmatrix}.
\]
Their determinant is $\mathcal W$, so their unique coefficients are
$a=\mathrm{Wr}(y,y_{2})/\mathcal W$ and $b=\mathrm{Wr}(y_{1},y)/\mathcal W$.
The numerator and denominator in each quotient have derivative
zero, hence $a,b\in C_{F}$. Conversely, constant linear
combinations remain solutions. If such a combination is zero,
its derivative is zero too, and the same invertible matrix makes
both coefficients zero. This proves spanning and independence.
\end{proof}

Applied to \cref{diff:cp:thm:critical}, these formulas give exactly
$y=u_{\alpha}(a+b\ell_{\alpha})$, with unique ordinary constants
$a,b$ (real in $\No$, complex in $\SC$). The first two stages make
the normalization concrete: for $\alpha=0$ the equation is
$\der^{2}y=0$ and the pair is $(1,\omega)$; for $\alpha=1$ the
potential is $\frac14\omega^{-2}$ and the pair is
$(\sqrt\omega,\sqrt\omega\log\omega)$. Both Wronskians are $1$.

With $F=\SC$ the solution collection is first a class; the spanning formulas
of \cref{diff:cp:thm:main} show that it is a set, isomorphic to $\C^{2}$ or to
$0$, so no dimension theory of a proper-class vector space is used, as in
\cref{diff:rs:rem:setsized}. All Hahn and Picard--Vessiot fields of
\cref{diff:cp:sec:fields,diff:cp:sec:pv} are set-sized from the outset.

\section{The Euler threshold at an infinite scale, and the missing oscillator}
\label{diff:cp:sec:euler}

\begin{corollary}[No nonreal powers of an infinite scale]
\label{diff:cp:cor:nonreal}
Let $s\in\No$ with $s>\R$, so that $\der s>0$ by \ref{diff:BM:positivity},
and let $\nu\in\C\setminus\R$. Then
\begin{equation}\label{diff:cp:eq:forbidden}
 \der z=\nu\,s^{\dagger}z,\qquad\text{equivalently}\qquad s\der_{s}z=\nu z,
\end{equation}
has no nonzero solution in $\SC$. For real $\nu$ its solutions are exactly
$\C\,s^{\nu}$, with $s^{\nu}=\exp(\nu\log s)$.
\end{corollary}
\begin{proof}
Write $\nu=a+ib$ with $a,b\in\R$ and $b\neq0$.
The imaginary coefficient $b\,s^{\dagger}$ has the explicit real
primitive $b\log s$. This primitive is infinite in modulus:
$\log s$ exceeds every real number, since an ordinary upper bound
would exponentiate to an ordinary upper bound on $s$.
If a nonzero solution existed, the direction
(i)$\Rightarrow$(ii) of \cref{diff:thm:phasecriterion}, or its direct
polar proof, would supply an infinitesimal real $\eta$ with
$\der\eta=b\,s^{\dagger}$. Then
$\der(\eta-b\log s)=0$, so $\eta-b\log s\in\R$ by the
constant-field property. An ordinary constant cannot turn an
infinite primitive into a finite one, a contradiction. This uses
the displayed primitive and no surjectivity of $\der$.

For real $\nu$, $s^{\nu}=\exp(\nu\log s)$ is nonzero and has
logarithmic derivative $\nu s^{\dagger}$. If $z$ is any solution,
$\der(z/s^{\nu})=0$, so $z/s^{\nu}\in\C$; conversely each such
constant multiple solves the equation, including zero.
\end{proof}

Member~13 proves this from an infinitesimal real primitive of the angular rate
of a unit, obtained from the strong Hahn logarithm; that primitive is the
infinitesimal phase residue $\eta$ of \cref{diff:thm:unitphase}, and its
argument is the direction (i)$\Rightarrow$(iii) of
\cref{diff:thm:phasecriterion}. The route is recorded as agreement and not
printed twice. The forbidden coefficient $\nu s^{\dagger}$ may be very small ---
at $s=\ell_{\alpha}$ it lies below $\ell_{\beta}^{\dagger}$ for every
$\beta<\alpha$ --- and the obstruction is the infinite primitive $\log s$, not
the size of the coefficient (\cref{diff:warn:integrated}). Three earlier
statements are instances: $s=e^{\omega}$ gives the first-order statement of
\cref{diff:cor:oscillator}; $s=\ell_{n}$ gives the case $p=1$ of
\cref{diff:cor:logscales}, multiplied by $\Imm\nu$; and $s=\ell_{1}$,
$\nu=-i$ gives \cref{diff:rs:prop:boundary}, since
$\derT y=iy/\log\omega$ is $\der y=-i\ell_{1}^{\dagger}y$.

\begin{lemma}[Liouville transformation as an operator identity]
\label{diff:cp:lem:gauge}
Let $F$ be a differential field of characteristic zero, and let
$s,u,Q\in F$ with $s\neq0$, $\der s\neq0$, $u\neq0$,
$\der^{2}u+Qu=0$ and $u^{2}\der s=1$. Define
$\der_{s}=(\der s)^{-1}\der$ and $s^{\dagger}=\der s/s$.
Then for every $v,c\in F$,
\begin{equation}\label{diff:cp:eq:gauge}
 \Bigl(\der^{2}+Q+c\,(s^{\dagger})^{2}\Bigr)(uv)
 =u\,(s^{\dagger})^{2}\Bigl((s\der_{s})^{2}-s\der_{s}+c\Bigr)v .
\end{equation}
\end{lemma}
\begin{proof}
Differentiating $u^{2}\der s=1$ gives $2\der u\,\der s+u\,\der^{2}s=0$. Since
$\der v=(\der s)\der_{s}v$ and $\der^{2}v=(\der^{2}s)\der_{s}v+(\der s)^{2}\der_{s}^{2}v$,
the finite product rule gives
\[
 \der^{2}(uv)+Quv=(\der^{2}u+Qu)v+(2\der u\,\der s+u\,\der^{2}s)\der_{s}v
 +u(\der s)^{2}\der_{s}^{2}v=u(\der s)^{2}\der_{s}^{2}v .
\]
Since $\der_{s}s=1$, $(s\der_{s})^{2}-s\der_{s}=s^{2}\der_{s}^{2}$; adding the
term $c(s^{\dagger})^{2}uv$ gives \cref{diff:cp:eq:gauge}.
\end{proof}

The rescaled operator $\der_{s}$ is a derivation and has the same
constant field as $\der$, since its multiplying scalar is nonzero.
The identity does not differentiate $c$, so $c$ need not be constant
in this lemma; the Euler classification below does require an
ordinary constant. It gives a bijection between the two kernels
by $v\mapsto uv$, because both $u$ and $u(s^{\dagger})^{2}$ are
nonzero. It also transports Wronskians exactly:
\[
 \mathrm{Wr}(uv_{1},uv_{2})
 =u^{2}(\der s)\mathrm{Wr}_{s}(v_{1},v_{2})
 =\mathrm{Wr}_{s}(v_{1},v_{2}).
\]
These are pointwise field identities, so they also apply in $\SC$.
They treat $v$ as an arbitrary element of the field and do not
assert that every surreal is a composite $v(s)$ and invoke no substitution
operator; such a distinction is essential in view of the global
composition obstruction of \cite[Theorem~8.4]{BMGerms}
(\cref{diff:rem:nocomposition}).

\begin{proposition}[Euler classification at a positive infinite scale]
\label{diff:cp:prop:euler}
Let $s\in\No$ with $s>\R$ (so $\der s>0$), let $c\in\R$, and put
$r_{\pm}=\frac12\bigl(1\pm\sqrt{1-4c}\bigr)$ when $c\leq\frac14$. Then
\[
 \Ker_{\SC}\Bigl((s\der_{s})^{2}-s\der_{s}+c\Bigr)=
 \begin{cases}
 \C\,s^{r_{+}}+\C\,s^{r_{-}}, & c<\frac14,\\[0.2em]
 \C\,s^{1/2}+\C\,s^{1/2}\log s, & c=\frac14,\\[0.2em]
 0, & c>\frac14,
 \end{cases}
\]
and the real kernel is the same with $\R$ in place of $\C$.
\end{proposition}
\begin{proof}
The operator is $s^{2}\bigl(\der_{s}^{2}+cs^{-2}\bigr)$, and $\der_{s}$ is a
derivation of $\SC$ with constants $\C$. For real $r$, $s\der_{s}(s^{r})=rs^{r}$.
If $c<\frac14$ the two distinct real roots of $r^{2}-r+c$ give solutions with
$\mathrm{Wr}_{s}(s^{r_{+}},s^{r_{-}})=(r_{-}-r_{+})s^{r_{+}+r_{-}-1}=r_{-}-r_{+}\neq0$,
and \cref{diff:cp:lem:span} for $\der_{s}$ gives completeness. If $c=\frac14$,
put $v_{0}=s^{1/2}$ and $v_{1}=s^{1/2}\log s$. Then
\[
 (s\der_{s}-\tfrac12)v_{0}=0,\qquad
 (s\der_{s}-\tfrac12)v_{1}=v_{0}.
\]
Both therefore lie in the squared factor's kernel, and
$\mathrm{Wr}_{s}(v_{0},v_{1})=v_{0}^{2}\der_{s}(\log s)=1$.
The same spanning lemma proves completeness at the repeated root.
If $c>\frac14$ the roots are
$\nu_{\pm}=\frac12\pm i\sqrt{c-\frac14}$, the operator factors as
$(s\der_{s}-\nu_{+})(s\der_{s}-\nu_{-})$ with commuting constant-coefficient
factors. If their product kills $v$, first apply
\cref{diff:cp:cor:nonreal} to
$(s\der_{s}-\nu_{-})v$: it must be zero. Apply that corollary
again to obtain $v=0$. Thus the conclusion does not assume
that a candidate solution has a power or exponential form.
For the real statement, the displayed bases are real and
independent over $\C$. Taking imaginary parts of a real solution's
unique complex expansion forces both coefficients to be real.
\end{proof}

The number $\frac14$ and the repeated-root logarithm are the classical Euler
phenomenon. The nonreal-root argument excludes \emph{every} nonzero element of
$\SC$, not merely a particular trial formula.

\begin{remark}[The Euler operator is a constant-coefficient operator at a rescaled scale]
\label{diff:cp:rem:constant}
This remark is an observation of this merge; member~13 does not state it. For
$s>0$ with $\der s\neq0$, one has
$s\der_{s}=\der_{\log s}$ (\cref{diff:prop:rescale}), and
$\log s$ is positive infinite when $s>\R$. So
\cref{diff:cp:prop:euler} is the order-two case of the constant-coefficient
classification of \cref{diff:thm:constant}, with $P(X)=X^{2}-X+c$, read for
$\der_{\log s}$ and $\log s$ in place of $\der$ and $\omega$:
the roots are real exactly when $c\leq\frac14$, and
$\exp(r\log s)=s^{r}$ for each real root $r$.
\Cref{diff:thm:constant} is stated for $\der$ only; its proof
goes through for $\der_{\log s}$ with \cref{diff:cp:cor:nonreal} in place of
\cref{diff:cor:oscillator}, but nothing below uses that, because member~13's
direct proof is complete.
\end{remark}

\begin{remark}[The missing oscillator is the Euler problem at the scale $e^{\omega}$]
\label{diff:cp:rem:oscillator}
This remark is an observation of this merge, with its proof; member~13 does
not state it. Take $s=e^{\omega}$. Then $s^{\dagger}=1$, $s\der_{s}=\der$, and
$u=e^{-\omega/2}$ satisfies $u^{2}\der s=1$ and $\der^{2}u-\frac14u=0$, so
\cref{diff:cp:lem:gauge} with $Q=-\frac14$ reads
\[
 \Bigl(\der^{2}+c-\tfrac14\Bigr)\bigl(e^{-\omega/2}v\bigr)
 =e^{-\omega/2}\bigl(\der^{2}-\der+c\bigr)v .
\]
For $c=\frac14+b^{2}$ with $b\in\R\setminus\{0\}$ the left operator is
$\der^{2}+b^{2}$, and the case $c>\frac14$ of \cref{diff:cp:prop:euler} is the
second-order statement of \cref{diff:cor:oscillator}; for $c=\frac14-\kappa^{2}$
with $\kappa>0$ real it returns exactly
$Ae^{\kappa\omega}+Be^{-\kappa\omega}$, $A,B\in\C$, solving
$\der^{2}y-\kappa^{2}y=0$. At the repeated endpoint
$c=\frac14$ it instead gives $A+B\omega$, as
\cref{diff:thm:constant} does. At the scale $s=\omega$ the same lemma, with
$u=1$ and $Q=0$, gives stage~$0$ of the family, and at the scale
$s=\ell_{\alpha}$, with $u=u_{\alpha}$ and $Q=Q_{\alpha}$, it gives
\cref{diff:cp:thm:main}. So the constant-potential oscillator, stage~$0$ and
stage~$\alpha$ are one Euler threshold read at the scales
$\log s=\omega$, $\ell_{1}$ and $\ell_{\alpha+1}$, and in every
supercritical case the obstruction is the infinite accumulated phase
$\sqrt{c-\frac14}\,\log s$ of \cref{diff:cp:cor:nonreal}; the constant-potential
threshold $b^{2}>0$ is the Euler threshold $c>\frac14$ shifted by the gauge
potential $-\frac14$.

This is the precise sense, and the only sense, in which the supercritical half
of \cref{diff:cp:thm:main} generalizes the missing oscillator: it moves the
obstructed accumulated phase from $\omega$ down through every stage of the
log-atomic tower, where the forbidden potential
$(c-\frac14)(\ell_{\alpha}^{\dagger})^{2}$ lies below every earlier scale
(\cref{diff:cp:thm:invisible}). The oscillator is not a member of the family:
its potential is an ordinary constant, its gauge potential $-\frac14$ is
negative, and its scale $e^{\omega}$ lies above the whole tower, since
$\log e^{\omega}=\ell_{0}$. Conversely the family is not a corollary of
\cref{diff:cor:oscillator}: it needs the imported tower, the critical identity
\cref{diff:cp:thm:critical} and the criterion at the scale $\ell_{\alpha+1}$.
Neither statement says anything about a general potential
(\cref{diff:cp:rem:notarbitrary}).
\end{remark}

\section{The transfinite classification}
\label{diff:cp:sec:classification}

\begin{theorem}[Transfinite critical classification]
\label{diff:cp:thm:main}
For every ordinal $\alpha$ and every $c\in\R$, the potential
$Q_{\alpha,c}=Q_{\alpha}+c(\ell_{\alpha}^{\dagger})^{2}$ is a well-defined
surreal, and the solution space of
\begin{equation}\label{diff:cp:eq:main}
 \der^{2}y+Q_{\alpha,c}\,y=0
\end{equation}
in $\No$ has dimension two over $\R$ if $c\leq\frac14$ and dimension zero if
$c>\frac14$. The identical statement holds in $\SC$ over $\C$. For $c<\frac14$
a basis is
\begin{equation}\label{diff:cp:eq:basissub}
 u_{\alpha}\ell_{\alpha}^{r_{+}},\qquad u_{\alpha}\ell_{\alpha}^{r_{-}},
 \qquad r_{\pm}=\tfrac12\bigl(1\pm\sqrt{1-4c}\bigr),
\end{equation}
and for $c=\frac14$ a basis is
\begin{equation}\label{diff:cp:eq:basiscrit}
 u_{\alpha+1},\qquad u_{\alpha+1}\ell_{\alpha+1}.
\end{equation}
Every displayed basis element is a single Conway monomial, and
$Q_{\alpha,1/4}=Q_{\alpha+1}$.
\end{theorem}
\begin{proof}
The coefficients exist by \cref{diff:cp:lem:triangular}. Apply
\cref{diff:cp:lem:gauge} in $F=\SC$ with $s=\ell_{\alpha}$, $u=u_{\alpha}$ and
$Q=Q_{\alpha}$; its hypotheses are \cref{diff:cp:thm:critical}, and
$s^{\dagger}=\ell_{\alpha}^{\dagger}$. Multiplication by $u_{\alpha}$ is a
bijection of $\SC$ and of $\No$, and the factor $u(s^{\dagger})^{2}$ is nonzero,
so \cref{diff:cp:eq:main} is equivalent to
$\bigl((s\der_{s})^{2}-s\der_{s}+c\bigr)v=0$, and
\cref{diff:cp:prop:euler} gives the dimensions and bases. At $c=\frac14$ use
$u_{\alpha}\ell_{\alpha}^{1/2}=\omega^{\varpi_{\alpha}/2+\omega^{-\alpha}/2}=u_{\alpha+1}$
and $\log\ell_{\alpha}=\ell_{\alpha+1}$. The monomial assertion is
$u_{\alpha}\ell_{\alpha}^{r}=\omega^{\varpi_{\alpha}/2+r\omega^{-\alpha}}$, by
\cref{diff:cp:prop:tower}.
\end{proof}

The theorem is proved at limit ordinals by the same argument as at successors;
it is not an induction that silently skips limits. The critical endpoint is
exactly the next stage, including after a limit. In Conway normal form,
\begin{align}
 Q_{\alpha,c}&=\frac14\sum_{\beta<\alpha}\omega^{-2\varpi_{\beta+1}}
 +c\,\omega^{-2\varpi_{\alpha+1}},\label{diff:cp:eq:qnormal}\\
 y_{\pm}&=\omega^{\varpi_{\alpha}/2+r_{\pm}\omega^{-\alpha}}
 &&(c<\tfrac14),\label{diff:cp:eq:ynormal}\\
 y_{0}&=\omega^{\varpi_{\alpha+1}/2},\qquad
 y_{1}=\omega^{\varpi_{\alpha+1}/2+\omega^{-(\alpha+1)}}
 &&(c=\tfrac14),\label{diff:cp:eq:ycrit}
\end{align}
with no unspecified primitive, inverse operator or integration constant; for
$c=0$ the last term of \cref{diff:cp:eq:qnormal} is omitted. The Wronskian transport following \cref{diff:cp:lem:gauge} gives,
in the printed order,
$\mathrm{Wr}(y_{+},y_{-})=-\sqrt{1-4c}$ and
$\mathrm{Wr}(y_{0},y_{1})=1$. Consequently every solution is a
unique constant linear combination of its displayed pair; the
constants are the Wronskian ratios of \cref{diff:cp:lem:span}.

\begin{remark}[A Riccati form]
\label{diff:cp:rem:riccati}
For a real root $r$ of $r^{2}-r+c$, the logarithmic derivative of
$u_{\alpha}\ell_{\alpha}^{r}$ is
\begin{equation}\label{diff:cp:eq:riccati}
 w=\tfrac12\ell_{<\alpha}^{\dagger}+r\,\ell_{\alpha}^{\dagger},
 \qquad \der w+w^{2}+Q_{\alpha,c}=0 .
\end{equation}
Indeed $\der\ell_{\alpha}^{\dagger}=-\ell_{\alpha}^{\dagger}(\ell_{<\alpha}^{\dagger}
+\ell_{\alpha}^{\dagger})$ by \cref{diff:cp:eq:ader}, and the coefficient of
$(\ell_{\alpha}^{\dagger})^{2}$ in the residual is $r^{2}-r+c$. When
$c>\frac14$ there is no \emph{real} Riccati solution: if such a
$w$ existed, surjectivity would supply $H\in\No$ with $\der H=w$.
Then $y=\exp H\neq0$ would satisfy
$\der^{2}y+Q_{\alpha,c}y=(\der w+w^{2}+Q_{\alpha,c})y=0$,
contradicting \cref{diff:cp:thm:main}. This implication uses
\ref{diff:BM:surj}, together with the real exponential law;
the second-order nonexistence itself does not.
\end{remark}

\begin{remark}[Not a classification of arbitrary potentials]
\label{diff:cp:rem:notarbitrary}
The theorem is exact for the family \cref{diff:cp:eq:qnormal}. It does not
say that every surreal potential has this form, that a comparison of leading
coefficients settles arbitrary perturbations, or that every
$\varepsilon\prec(\ell_{\alpha}^{\dagger})^{2}$ preserves subcritical or
supercritical behaviour of $Q_{\alpha,c}+\varepsilon$; that would need further
differential asymptotic estimates, which the gauge identity alone does not
give. At the critical endpoint a smaller perturbation can change
the answer: for any real $d$,
\[
 Q_{\alpha,1/4}+d(\ell_{\alpha+1}^{\dagger})^{2}
 =Q_{\alpha+1,d}.
\]
The added term is $\prec(\ell_{\alpha}^{\dagger})^{2}$ when
$d\neq0$, yet $d=0$ gives dimension two and $d=1$ gives dimension
zero. No comparison theorem for general surreal Riccati equations
is used.
\end{remark}

\begin{remark}[Stage zero is a regular-singular system; the finite stages are first-order thresholds]
\label{diff:cp:rem:stagezero}
At $\alpha=0$: $\ell_{0}=\omega$, $u_{0}=1$, $Q_{0}=0$ and
$\ell_{0}^{\dagger}=t$, so \cref{diff:cp:eq:main} is $\der^{2}y+ct^{2}y=0$. With
$\derT=-\omega\der$ one has $\omega^{2}\der^{2}=\derT^{2}+\derT$, so the
equation is $\derT^{2}y+\derT y+cy=0$, the power-Hahn regular-singular system
$\derT Y=A_{0}Y$ of \cref{diff:rs:def:class} for $Y=(y,\derT y)$, with
$A_{0}=\left(\begin{smallmatrix}0&1\\-c&-1\end{smallmatrix}\right)$ and
$A_{+}=0$. Its eigenvalues $\mu$ solve $\mu^{2}+\mu+c=0$ and are real exactly
when $c\leq\frac14$, so \cref{diff:rs:thm:selection} gives dimension $2$ or
$0$, in agreement; the bases correspond through $t^{\mu}=\omega^{r}$ with
$\mu=-r$, and at $c=\frac14$ the Jordan block gives $\omega^{1/2}$ and
$\omega^{1/2}\log\omega$. This comparison is made by this merge and is
recorded as agreement.

There is one overlap at the next stage:
$Q_{1,0}=Q_{1}=\frac14\omega^{-2}=Q_{0,1/4}$ belongs to $\KR$
and is the repeated-root case just described. For $\alpha=1$
with $c\neq0$, the term
$c\omega^{-2\varpi_{2}}=c\omega^{-2-2\omega^{-1}}$ has Conway
exponent $-2-2\omega^{-1}\notin\R$, so $Q_{1,c}\notin\KR$. For every $\alpha\geq2$
the same monomial occurs with coefficient $\frac14$ in
$Q_{\alpha}$; its exponent differs from all later terms, including
the perturbation, so it cannot cancel. Thus these potentials
also lie outside $\KR$. This exact support test, rather than
merely the appearance of a logarithm in a formula, establishes
the boundary of \cref{diff:rs:part}. For finite $n$ the
supercritical obstruction is the boundary case $p=1$ of
\cref{diff:cor:logscales} at depth $n$, and \cref{diff:cp:cor:nonreal} extends
it to every ordinal.
\end{remark}

\begin{remark}[Amplitude and phase of a supercritical member]
\label{diff:cp:rem:pinney}
This remark is an observation of this merge, with its proof; member~13 does
not state it. Let $c>\frac14$ and put
$\rho=(c-\frac14)^{-1/4}u_{\alpha+1}>0$. Then
\[
 \der^{2}\rho+Q_{\alpha,c}\,\rho=\rho^{-3},
\]
because $\der^{2}u_{\alpha+1}=-Q_{\alpha+1}u_{\alpha+1}$ by
\cref{diff:cp:thm:critical}, $Q_{\alpha,c}-Q_{\alpha+1}=(c-\frac14)(\ell_{\alpha}^{\dagger})^{2}$,
and $(\ell_{\alpha}^{\dagger})^{2}u_{\alpha+1}^{4}=1$ since
$u_{\alpha+1}^{2}=\ell_{<\alpha+1}=(\ell_{\alpha}^{\dagger})^{-1}$. Moreover
$\rho^{-2}=\sqrt{c-\frac14}\;\der\ell_{\alpha+1}$, whose normalized primitive is
$\sqrt{c-\frac14}\;\ell_{\alpha+1}$, a purely infinite monomial with zero
infinitesimal part. By \cref{diff:cp:thm:main} the equation is intrinsically
oscillatory in the sense of \cref{diff:thm:matamplitude}, so
\cref{diff:thm:matpinney}(i) makes $\rho$ the \emph{unique} positive Pinney
solution, \cref{diff:thm:matamplitude} has canonical phase
$P=\sqrt{c-\frac14}\;\ell_{\alpha+1}$ with $\eta=0$, and the companion system has
differential phase spectrum $\{P,-P\}$. The Pinney identity is direct; the
uniqueness, the canonical phase and the spectrum use \cref{diff:part:matrix},
hence the import \cref{diff:thm:matsurjectivity} (item~N53) and
\ref{diff:BM:surj}. At the scale $e^{\omega}$ of \cref{diff:cp:rem:oscillator}
the same formula gives $\rho=\abs b^{-1/2}$ and $P=\abs b\,\omega$ for
$\der^{2}y+b^{2}y=0$, which for $b=1$ is the value $\rho=1$ recorded after
\cref{diff:thm:matpinney}.
\end{remark}

\section{Invisible perturbations, and no final stage}
\label{diff:cp:sec:invisible}

\begin{proposition}[The exact earlier-stage asymptotics]
\label{diff:cp:prop:earlier}
For $\beta<\alpha$ and $c,d\in\R$,
\begin{equation}\label{diff:cp:eq:earlier}
 Q_{\alpha,c}-Q_{\beta}\sim\tfrac14(\ell_{\beta}^{\dagger})^{2},
 \qquad
 Q_{\alpha,c}-Q_{\alpha,d}=(c-d)(\ell_{\alpha}^{\dagger})^{2}
 \prec(\ell_{\beta}^{\dagger})^{2}\quad(c\neq d).
\end{equation}
\end{proposition}
\begin{proof}
After subtracting $Q_{\beta}$, the largest monomial left in
\cref{diff:cp:eq:qnormal} is $(\ell_{\beta}^{\dagger})^{2}$, with coefficient
$\frac14$; every other term is strictly smaller by \cref{diff:cp:lem:scales}.
The second assertion is the explicit difference with the same separation.
\end{proof}

\begin{theorem}[No set of earlier scales decides solvability]
\label{diff:cp:thm:invisible}
For every set $\mathcal B$ of ordinals there are positive surreals $Q_{-},Q_{+}$
with
\[
 Q_{+}-Q_{-}\prec(\ell_{\beta}^{\dagger})^{2}
 \quad\text{and}\quad
 Q_{\pm}-Q_{\beta}\sim\tfrac14(\ell_{\beta}^{\dagger})^{2}
 \qquad(\beta\in\mathcal B),
\]
such that $\der^{2}y+Q_{-}y=0$ has two independent surcomplex solutions and
$\der^{2}y+Q_{+}y=0$ has only $y=0$ in $\SC$.
\end{theorem}
\begin{proof}
Since $\mathcal B$ is a set, there is an ordinal $\alpha\geq1$ with
$\alpha>\beta$ for every $\beta\in\mathcal B$. Take $Q_{-}=Q_{\alpha,0}=Q_{\alpha}$
and $Q_{+}=Q_{\alpha,1}$. Both are positive, with leading term
$\frac14\omega^{-2}$; \cref{diff:cp:thm:main} gives dimensions two and zero, and
\cref{diff:cp:prop:earlier} gives the two precision assertions.
\end{proof}

\begin{example}[The first infinite stage]
\label{diff:cp:ex:omegaexpansion}
At $\alpha=\omega$, with $\ell_{n}=\log_{n}\omega$,
\begin{equation}\label{diff:cp:eq:Qomega}
 Q_{\omega}=\frac{1}{4\omega^{2}}
 +\frac{1}{4\omega^{2}(\log\omega)^{2}}
 +\frac{1}{4\omega^{2}(\log\omega)^{2}(\log_{2}\omega)^{2}}+\cdots .
\end{equation}
The coefficient $(\ell_{\omega}^{\dagger})^{2}$ is smaller than every displayed
monomial, yet adding it changes the dimension from two to zero, and
$Q_{\omega}+c(\ell_{\omega}^{\dagger})^{2}$ has the exact threshold
$c=\frac14$. Specifying the entire countable logarithmic expansion
\cref{diff:cp:eq:Qomega} does not specify solvability.
\end{example}

\begin{warning}[What the invisibility statement is about]
\label{diff:cp:warn:quotient}
\Cref{diff:cp:thm:invisible} concerns the explicitly named precision hierarchy
$(\ell_{\beta}^{\dagger})^{2}$, not every conceivable observation or algorithm.
It is a statement about a quotient by an additive class of errors smaller than
all the designated monomials. It is not a discontinuity claim in the fine
surreal topology (\cref{diff:warn:summability}), and that additive class is not
called an ideal of the surreal field.
\end{warning}

\begin{proposition}[There is no all-ordinal critical potential]
\label{diff:cp:prop:noend}
There is no surreal $Q$ with
$Q-Q_{\beta}\sim\frac14(\ell_{\beta}^{\dagger})^{2}$ for every ordinal $\beta$.
In particular $\frac14\sum_{\beta\in\mathbf{On}}(\ell_{\beta}^{\dagger})^{2}$ is not a
Hahn sum in $\No$.
\end{proposition}
\begin{proof}
The relation at $\beta$ forces the coefficient of $(\ell_{\beta}^{\dagger})^{2}$
in the normal form of $Q$ to be $\frac14$, since the support of $Q_{\beta}$
consists of earlier monomials. The support of $Q$ would then contain the
proper class of distinct monomials $(\ell_{\beta}^{\dagger})^{2}$,
contradicting the set-sized support of a surreal.
\end{proof}

This does not conflict with the existence of $Q_{\alpha}$ for every set
ordinal, uncountable ones included: every initial segment is a set, their
union is not. The absence of a final stage is a boundary of the construction,
not a convergence problem to be repaired by a summation convention
(\cref{diff:sub:classes}).

\begin{remark}[Relation to the valuation-insufficiency pair]
\label{diff:cp:rem:separator}
\Cref{diff:ex:logseparator,diff:warn:noglobalization} show that for first-order
equations a real-exponent valuation cannot see the slower logarithmic scale
that decides solvability. \Cref{diff:cp:thm:invisible} is the second-order,
transfinite counterpart: within this family no set of logarithmic scales
below the relevant stage decides, while the single coefficient $c$ at the
stage $\alpha$ does. It concerns this family only
(\cref{diff:cp:rem:notarbitrary}).
\end{remark}

\section{Three differential Hahn fields at a limit ordinal}
\label{diff:cp:sec:fields}

Fix a nonzero limit ordinal $\lambda$ and put
\begin{equation}\label{diff:cp:eq:groups}
 \Delta_{0}=\Span_{\R}\{\varpi_{\beta}:\beta<\lambda\},\qquad
 \Delta_{1}=\Delta_{0}+\R\,\varpi_{\lambda},\qquad
 \Delta_{2}=\Delta_{1}+\R\,\omega^{-\lambda},
\end{equation}
set-sized ordered real vector subgroups of the additive surreals. Let
$H_{\Delta_{j}}=\R((t^{\Delta_{j}}))$ and $K_{\Delta_{j}}=H_{\Delta_{j}}[i]
=\C((t^{\Delta_{j}}))$ be the full Hahn fields, viewed in $\No$ and $\SC$;
since each $\Delta_{j}$ is a group, $t^{\Delta_{j}}=\omega^{\Delta_{j}}$, and
membership of an exponent does not depend on the orientation.

\begin{warning}[Two levels of support]
\label{diff:cp:warn:supports}
$\Span_{\R}$ means \emph{finite} real linear combinations. It is not closure
under Hahn sums inside the exponent group: an element of $\Delta_{0}$ has
bounded \emph{inner} support, while an element of $H_{\Delta_{0}}$ may have
arbitrary well-ordered \emph{outer} support made from such exponents. The
outer sum defining $Q_{\lambda}$ is allowed; the inner sum $\varpi_{\lambda}$ is
not an allowed exponent of $H_{\Delta_{0}}$. Everything below depends on this
distinction, and confusing the two levels would destroy the example.
\end{warning}

\begin{lemma}[Bounded support in the exponent group]
\label{diff:cp:lem:exponents}
Every $\gamma\in\Delta_{0}$ has Conway support in
$\{\omega^{-\zeta}:\zeta<\beta\}$ for some $\beta<\lambda$. Consequently
$\varpi_{\lambda}\notin\Delta_{0}$, $\omega^{-\lambda}\notin\Delta_{1}$, and
$\varpi_{\lambda},\omega^{-\lambda}$ are linearly independent over $\R$
modulo $\Delta_{0}$.
\end{lemma}
\begin{proof}
A finite combination of the $\varpi_{\beta}$ has support bounded by the largest
of its finitely many indices, which is below $\lambda$. The support of
$\varpi_{\lambda}$ is all of $\{\omega^{-\zeta}:\zeta<\lambda\}$ and is not so
bounded. Every element of $\Delta_{1}$ has support inside that family, hence
zero coefficient at $\omega^{-\lambda}$. If
$a\varpi_{\lambda}+b\omega^{-\lambda}\in\Delta_{0}$, the coefficient at
$\omega^{-\lambda}$ gives $b=0$ and the unbounded support then gives $a=0$.
\end{proof}

All three groups are divisible, so the classical Hahn real-closedness theorem
makes every $H_{\Delta_{j}}$ real closed and every $K_{\Delta_{j}}$
algebraically closed \cite[Section~2.3]{BM}, the standard algebraic fact used in
\cref{diff:thm:hull}. \Cref{diff:cp:lem:exponents} shows that both inclusions
$H_{\Delta_{0}}\subsetneq H_{\Delta_{1}}\subsetneq H_{\Delta_{2}}$ are strict.

\begin{proposition}[Derivative stability is a support assertion]
\label{diff:cp:prop:fieldclosure}
Each $H_{\Delta_{j}}$ and $K_{\Delta_{j}}$ is $\der$-stable and contains
$Q_{\lambda}$; their constants are $\R$ and $\C$. More precisely, the
logarithmic derivative of every monomial $\omega^{\gamma}$ with
$\gamma\in\Delta_{1}$ lies in $H_{\Delta_{0}}$, and with $\gamma\in\Delta_{2}$ in
$H_{\Delta_{2}}$.
\end{proposition}
\begin{proof}
If $\beta<\lambda$ then $\beta+1<\lambda$, so $\varpi_{\beta+1}\in\Delta_{0}$
and $\ell_{\beta}^{\dagger}=\omega^{-\varpi_{\beta+1}}\in H_{\Delta_{0}}$; hence
$\ell_{<\beta}^{\dagger}$ for $\beta\leq\lambda$, and $Q_{\lambda}$, have support
in $\Delta_{0}$ and lie in the full Hahn field $H_{\Delta_{0}}$. For
$\gamma=\sum_{j=1}^{m}r_{j}\varpi_{\beta_{j}}\in\Delta_{0}$ with $\beta_{j}<\lambda$,
$\omega^{\gamma}=\prod_{j}\ell_{<\beta_{j}}^{r_{j}}$ by \cref{diff:cp:prop:tower}
and the multiplicativity of the monomial map; so
$(\omega^{\gamma})^{\dagger}=\sum_{j}r_{j}\ell_{<\beta_{j}}^{\dagger}\in H_{\Delta_{0}}$.
A further term $r\varpi_{\lambda}$ adds $r\ell_{<\lambda}^{\dagger}\in
H_{\Delta_{0}}$, and a further $r\omega^{-\lambda}$ adds
$r\ell_{\lambda}^{\dagger}=r\omega^{-\varpi_{\lambda}-\omega^{-\lambda}}\in
H_{\Delta_{2}}$. So each $\Delta_{j}$ is differentially closed on monomials
(\cref{diff:def:monomialclosed}), and \cref{diff:lem:monomialcriterion} gives
$\der$-stability of $H_{\Delta_{j}}$ and $K_{\Delta_{j}}$; that step uses
strong additivity \ref{diff:BM:additive}, not continuity. A constant of a
subfield is an ambient constant, and the fields contain $\R$, respectively $\C$.
\end{proof}

\begin{theorem}[Zero, one and two solutions over closed Hahn fields]
\label{diff:cp:thm:fieldjump}
For $j=0,1,2$,
\begin{equation}\label{diff:cp:eq:jump}
 \dim_{\R}\Sol_{H_{\Delta_{j}}}(Q_{\lambda})
 =\dim_{\C}\Sol_{K_{\Delta_{j}}}(Q_{\lambda})=j,
\end{equation}
and the complex solution spaces are $0$, $\C\,u_{\lambda}$ and
$\C\,u_{\lambda}+\C\,u_{\lambda}\ell_{\lambda}$.
\end{theorem}
\begin{proof}
By \cref{diff:cp:thm:critical,diff:cp:lem:span} every ambient solution is
$au_{\lambda}+bu_{\lambda}\ell_{\lambda}$ with $a,b\in\C$, a combination of the
two distinct monomials with exponents $\varpi_{\lambda}/2$ and
$\varpi_{\lambda}/2+\omega^{-\lambda}$. By \cref{diff:cp:lem:exponents} neither
lies in $\Delta_{0}$, only the first lies in $\Delta_{1}$, and both lie in
$\Delta_{2}$. A nonzero constant combination of distinct monomials cannot cancel
either support position, so Hahn membership gives exactly the stated spaces;
the real case is the same with real constants.
\end{proof}

\begin{proposition}[An exponential gap followed by a primitive gap]
\label{diff:cp:prop:gaps}
$H_{\Delta_{0}}$ and $H_{\Delta_{1}}$ are closed under $\log$ on positive
elements. Nevertheless
\[
 \log u_{\lambda}\in H_{\Delta_{0}},\quad u_{\lambda}\notin H_{\Delta_{0}};
 \qquad
 u_{\lambda}^{-2}\in H_{\Delta_{1}},\quad
 \text{no primitive of }u_{\lambda}^{-2}\text{ lies in }H_{\Delta_{1}} .
\]
\end{proposition}
\begin{proof}
For $\beta\leq\lambda$, $\log\ell_{<\beta}=\sum_{\zeta<\beta}\ell_{\zeta+1}$, and
the exponent of $\ell_{\zeta+1}$ is
$\omega^{-(\zeta+1)}=\varpi_{\zeta+2}-\varpi_{\zeta+1}\in\Delta_{0}$ because
$\zeta+2<\lambda$ for $\zeta<\lambda$; so $\log\ell_{<\beta}\in H_{\Delta_{0}}$,
and the logarithm of any monomial with exponent in $\Delta_{1}$ is a finite
real combination of these, hence in $H_{\Delta_{0}}$. A positive element of
$H_{\Delta_{j}}$ is $r\omega^{\gamma}(1+\varepsilon)$ with $r>0$ real and
$\varepsilon\prec1$ in $H_{\Delta_{j}}$, and $\log(1+\varepsilon)$ is a summable
Taylor series in $\varepsilon$ with support in $\Delta_{j}$ (\cref{diff:lem:neumann}).
This gives the closure. Now $\log u_{\lambda}=\frac12\log\ell_{<\lambda}\in
H_{\Delta_{0}}$ while $\varpi_{\lambda}/2\notin\Delta_{0}$. Also
$u_{\lambda}^{-2}=\omega^{-\varpi_{\lambda}}\in H_{\Delta_{1}}$ and
$\der\ell_{\lambda}=u_{\lambda}^{-2}$; every real primitive is $\ell_{\lambda}+r$
with $r\in\R$, and $\omega^{-\lambda}\notin\Delta_{1}$.
\end{proof}

\begin{remark}[Which closure is missing, and the relation to \cref{diff:part:workspaces,diff:part:matrix}]
\label{diff:cp:rem:workspaces}
Full Hahn summation, real closedness, derivative stability and closure under
logarithms are all present in $H_{\Delta_{0}}$, and still
$\Sol_{H_{\Delta_{0}}}(Q_{\lambda})=0$: the first missing operation is an
exponential, the second a primitive, which is the phenomenon of
\cref{diff:ex:logmissing} at a limit stage. No closure assertion is made about
$H_{\Delta_{2}}$, logarithms included. $\Delta_{0}$ is divisible and
differentially closed on monomials, so it is its own differential Hahn hull
(\cref{diff:thm:hull}); the example is a further witness that the hull is not
solution closure (\cref{diff:warn:hulllimits}). Over the whole of $\No$,
\cref{diff:thm:matpinney} excludes a one-dimensional solution space for a real
potential, because a primitive of $u^{-2}$ produces a second solution from one
(\ref{diff:BM:surj}); $H_{\Delta_{1}}$ lacks exactly that primitive, which is
why the one-dimensional case occurs there and nowhere in $\No$. The three
complex fields are algebraically closed, so no step of the jump is an
algebraic root; compare \cref{diff:rem:algnotdiff}.
\end{remark}

\section{A Borel Picard--Vessiot group at every limit stage}
\label{diff:cp:sec:pv}

Keep $\lambda$ and take the set-sized differential field $k=K_{\Delta_{0}}$, with
constants $\C$. By \cref{diff:cp:thm:critical},
\begin{equation}\label{diff:cp:eq:pvder}
 \der u_{\lambda}=h\,u_{\lambda},\quad h=\tfrac12\ell_{<\lambda}^{\dagger}\in k,
 \qquad
 \der\ell_{\lambda}=u_{\lambda}^{-2}.
\end{equation}
The calculation below needs neither an abstract existence theorem for
Picard--Vessiot extensions nor a black-box Galois computation, in keeping with
\cref{diff:sec:PV}.

\begin{lemma}[Algebraic independence of the two generators]
\label{diff:cp:lem:algind}
$u_{\lambda}$ and $\ell_{\lambda}$ are algebraically independent over $k$.
\end{lemma}
\begin{proof}
Their exponents $\varpi_{\lambda}/2$ and $\omega^{-\lambda}$ are linearly
independent over $\Q$ modulo $\Delta_{0}$ (\cref{diff:cp:lem:exponents}). In a
finite relation $\sum_{(n,m)}f_{nm}u_{\lambda}^{n}\ell_{\lambda}^{m}=0$ with
$f_{nm}\in k$, a nonzero term has support in the coset
$\Delta_{0}+n\varpi_{\lambda}/2+m\omega^{-\lambda}$, distinct pairs give
disjoint cosets, and terms in disjoint cosets cannot cancel.
\end{proof}

\begin{theorem}[Borel differential Galois group]
\label{diff:cp:thm:pv}
The field $L=k(u_{\lambda},\ell_{\lambda})\subseteq K_{\Delta_{2}}$ is a
Picard--Vessiot field for $\der^{2}y+Q_{\lambda}y=0$ over $k$. Its constants are
$\C$, its transcendence degree over $k$ is two, and its differential
$k$-automorphisms are exactly
\begin{equation}\label{diff:cp:eq:pvaction}
 \psi_{a,b}(u_{\lambda})=a\,u_{\lambda},\qquad
 \psi_{a,b}(\ell_{\lambda})=a^{-2}\ell_{\lambda}+b,
 \qquad a\in\C^{\times},\ b\in\C .
\end{equation}
On the ordered solution basis $(u_{\lambda},u_{\lambda}\ell_{\lambda})$ the group
is
\begin{equation}\label{diff:cp:eq:borel}
 \Bor_{2}(\C)=\left\{\begin{pmatrix}a&d\\0&a^{-1}\end{pmatrix}:
 a\in\C^{\times},\ d\in\C\right\}\subseteq\SL_{2}(\C).
\end{equation}
\end{theorem}
\begin{proof}
By \cref{diff:cp:eq:pvder}, $L$ is a differential field. It contains the
fundamental matrix
\[
 Y=\begin{pmatrix}u_{\lambda}&u_{\lambda}\ell_{\lambda}\\
 hu_{\lambda}&hu_{\lambda}\ell_{\lambda}+u_{\lambda}^{-1}\end{pmatrix},
 \qquad \det Y=1,
\]
whose entries generate $u_{\lambda}$ and $\ell_{\lambda}=(u_{\lambda}\ell_{\lambda})/u_{\lambda}$
over $k$. No new constant occurs, because $L\subseteq\SC$, whose constants are
$\C\subseteq k$; these are the Picard--Vessiot requirements. The transcendence
degree is \cref{diff:cp:lem:algind}. For a differential $k$-automorphism
$\psi$, $\psi(u_{\lambda})^{\dagger}=h=u_{\lambda}^{\dagger}$, so
$\psi(u_{\lambda})=au_{\lambda}$ with $a\in\C^{\times}$; then
$\der\psi(\ell_{\lambda})=a^{-2}u_{\lambda}^{-2}$, so
$\psi(\ell_{\lambda})-a^{-2}\ell_{\lambda}=b\in\C$. Conversely,
\cref{diff:cp:lem:algind} makes \cref{diff:cp:eq:pvaction} a well-defined
$k$-automorphism of the rational function field $k(u_{\lambda},\ell_{\lambda})$;
\cref{diff:cp:eq:pvder} shows that it commutes with $\der$ on the generators,
hence on $L$. It sends $u_{\lambda}\ell_{\lambda}$ to
$ab\,u_{\lambda}+a^{-1}u_{\lambda}\ell_{\lambda}$, so its matrix is
$\left(\begin{smallmatrix}a&ab\\0&a^{-1}\end{smallmatrix}\right)$, and $d=ab$ is
arbitrary. The parameters and the action are algebraic, so the algebraic group
is the Borel subgroup itself, not only an abstract list of automorphisms.
\end{proof}

\begin{corollary}[After the first Hahn enlargement]
\label{diff:cp:cor:additive}
Over $K_{\Delta_{1}}$ the Picard--Vessiot field is $K_{\Delta_{1}}(\ell_{\lambda})$,
with differential Galois group $\ell_{\lambda}\mapsto\ell_{\lambda}+b$, $b\in\C$.
\end{corollary}
\begin{proof}
$K_{\Delta_{1}}$ contains $u_{\lambda}$ and not $\ell_{\lambda}$, which is therefore
transcendental over it, $K_{\Delta_{1}}$ being algebraically closed. Its
derivative lies in $K_{\Delta_{1}}$, an automorphism fixing $u_{\lambda}$ has
$a=1$, and the same fundamental matrix and no-new-constants argument apply.
\end{proof}

\begin{remark}[Scope of the Galois computation]
\label{diff:cp:rem:pvscope}
``Minimal'' refers to the solution-generated field $L$, not to the much larger
$K_{\Delta_{2}}$: every fundamental system is a constant invertible
combination of $(u_{\lambda},u_{\lambda}\ell_{\lambda})$ and generates $L$. The
intermediate Hahn fields are explicit workspaces, and no minimality is claimed
for them among arbitrary fields. The other differential Galois groups computed
in this report are commutative --- $\gm(\C)$ in \cref{diff:thm:PVone}, the tori
of \cref{diff:thm:torus,diff:cor:mattorus}, and $(\C,+)$ in
\cref{diff:rs:thm:minlog} and \cref{diff:cp:cor:additive} --- so
$\Bor_{2}(\C)$ is the only non-commutative one; it reflects the two successive
missing operations of \cref{diff:cp:prop:gaps}, an exponential and a primitive.
\end{remark}

\section{Schwarzian form, the first limit stage, and what this part does not claim}
\label{diff:cp:sec:schwarzian}

\subsection{Schwarzian coordinates}
\label{diff:cp:sub:schwarzian}

For $s\in\SC$ with $\der s\neq0$ put
\begin{equation}\label{diff:cp:eq:schwarzian}
 \Sch(s)=\frac{\der^{3}s}{\der s}-\frac32\Bigl(\frac{\der^{2}s}{\der s}\Bigr)^{2},
\end{equation}
a finite differential rational expression; no topology and no composition is
involved.

\begin{lemma}[Schwarzian and a normalized fundamental pair]
\label{diff:cp:lem:schwarzian}
For $Q\in\SC$ the following are equivalent: some $s\in\SC$ has $\der s\neq0$ and
$\Sch(s)=2Q$; the equation $\der^{2}y+Qy=0$ has a fundamental pair in $\SC$. If
$y_{1},y_{2}$ is such a pair, every such $s$ is a constant fractional linear
transform of $y_{2}/y_{1}$.
\end{lemma}
\begin{proof}
Given $s$, algebraic closedness supplies $v=(\der s)^{-1/2}\neq0$, for either
root. Then $v^{\dagger}=-\frac12\der^{2}s/\der s$ and
$\der^{2}v/v=-\frac12\der^{3}s/\der s+\frac34(\der^{2}s/\der s)^{2}=-\frac12\Sch(s)$,
so $v$ and $vs$ are solutions with $\mathrm{Wr}(v,vs)=v^{2}\der s=1$.
Conversely, for a fundamental pair with constant $\mathrm{Wr}\neq0$, the
quotient $s=y_{2}/y_{1}$ has $\der s=\mathrm{Wr}/y_{1}^{2}$, and the same
differentiation gives $\Sch(s)=-2\der^{2}y_{1}/y_{1}=2Q$. If $s$ and
$s_{0}=y_{2}/y_{1}$ both have Schwarzian $2Q$, their reconstructed fundamental
pairs differ by an invertible constant matrix $(m_{jk})$
(\cref{diff:cp:lem:span}), so $s=(m_{11}s_{0}+m_{12})/(m_{21}s_{0}+m_{22})$, the
denominator being nonzero because
$\der s_{0}\neq0$ rules out $s_{0}\in\C$; conversely every such quotient comes
from another fundamental pair.
\end{proof}

\begin{corollary}[Transfinite Schwarzian potentials]
\label{diff:cp:cor:schwarzian}
For every ordinal $\alpha$, $\Sch(\ell_{\alpha})=2Q_{\alpha}$. The $s$ with
$\der s\neq0$ and $\Sch(s)=2Q_{\alpha,c}$ are the constant fractional linear
transforms of $\ell_{\alpha}^{\sqrt{1-4c}}$ if $c<\frac14$ and of
$\ell_{\alpha+1}$ if $c=\frac14$; if $c>\frac14$ there is no such $s$ in $\SC$.
\end{corollary}
\begin{proof}
Use the pair $u_{\alpha},u_{\alpha}\ell_{\alpha}$ for $Q_{\alpha}$, the quotient
of the pair \cref{diff:cp:eq:basissub} for $c<\frac14$, and of
\cref{diff:cp:eq:basiscrit} for $c=\frac14$; then apply
\cref{diff:cp:lem:schwarzian,diff:cp:thm:main}.
\end{proof}

So the invisible perturbation of \cref{diff:cp:thm:invisible} also crosses the
image of the scalar Schwarzian map. This is an exact projective reformulation
of \cref{diff:cp:thm:main}, not an independent result, and the classical
Schwarzian correspondence is reproved only to make its application to
surcomplex scalars unambiguous.

\subsection{The first limit stage in normal form}
\label{diff:cp:sub:omega}

\begin{example}[$\lambda=\omega$]
\label{diff:cp:ex:omega}
The inner exponent sum is the strong geometric series
$\varpi_{\omega}=\sum_{n\geq0}\omega^{-n}=\omega/(\omega-1)$; the identity does
not turn an outer monomial into a variable-base exponential. Then
\[
 u_{\omega}=\omega^{\omega/(2(\omega-1))},\qquad
 \ell_{\omega}=\omega^{\omega^{-\omega}},\qquad
 \ell_{\omega}^{\dagger}=\omega^{-\omega/(\omega-1)-\omega^{-\omega}},\qquad
 \der\ell_{\omega}=\omega^{-\omega/(\omega-1)}=u_{\omega}^{-2}.
\]
$\ell_{\omega}$ is a positive infinite number below every $\log_{n}\omega$; its
role as a primitive of the infinite inverse logarithmic product is recorded in
the proof of \cite[Theorem~8.4]{BMGerms}, and its existence is not a new
assertion. The whole surcomplex solution space of $\der^{2}y+Q_{\omega}y=0$ is
\[
 \C\,\omega^{\omega/(2(\omega-1))}+\C\,\omega^{\omega/(2(\omega-1))+\omega^{-\omega}},
\]
$Q_{\omega}+(\ell_{\omega}^{\dagger})^{2}$ has only the zero solution, and at
the boundary $Q_{\omega}+\frac14(\ell_{\omega}^{\dagger})^{2}=Q_{\omega+1}$ the
basis is $\omega^{\varpi_{\omega+1}/2}$,
$\omega^{\varpi_{\omega+1}/2+\omega^{-(\omega+1)}}$. All of this is exact. The
groups of \cref{diff:cp:eq:groups} are
\[
 \Delta_{0}=\bigoplus_{n\geq0}\R\,\omega^{-n},\qquad
 \Delta_{1}=\Delta_{0}+\R\,\frac{\omega}{\omega-1},\qquad
 \Delta_{2}=\Delta_{1}+\R\,\omega^{-\omega};
\]
the direct sum is inside the additive surreals, its members have finitely
many terms, and it is not the field of Laurent series in $\omega^{-1}$. Every
term of $Q_{\omega}$ has exponent $-2(1+\omega^{-1}+\dots+\omega^{-n})\in\Delta_{0}$,
so $H_{\Delta_{0}}$ contains the potential; adjoining the direction
$\omega/(\omega-1)$ admits the first solution monomial, and adjoining
$\omega^{-\omega}$ the second. At higher limits ``finite inner support'' is
replaced by ``support bounded below $\lambda$'' (\cref{diff:cp:lem:exponents});
no countability assumption is used there or in \cref{diff:cp:thm:pv}.
\end{example}

\subsection{Inputs, boundaries and continuations}
\label{diff:cp:sub:limits}

\begin{remark}[Which assumptions matter, and the derivation]
\label{diff:cp:rem:derivation}
The normal-form and strong-additivity inputs justify every limit stage. The
ordinal derivative formula of \cref{diff:cp:prop:tower} fixes the coefficients
and the bases; the normalization $\der\omega=1$ alone does not determine it.
The exact constant field is essential both to the angular obstruction and to
the spanning argument, and the complex argument needs only infinitesimal
logarithms, not an exponential at arbitrary imaginary arguments. For another
surreal derivation the finite identities may still hold, but the limit-stage
derivatives need not be the Berarducci--Mantova ones; no derivation-independent
version is claimed, and in particular nothing here is asserted for the class
of normalized derivations of \cref{diff:aut:part}. These are not statements
about coordinate differentiation in $\Hol(U)((t^{\Gamma}))$
(\cref{diff:part:coordinate}) or about fine-differentiable class functions
(\cref{diff:warn:separators}).
\end{remark}

\begin{remark}[Established inputs and proposed contributions]
\label{diff:cp:rem:inputs}
The finite critical hierarchy is classical Kneser--Weber--Hartman--Hille
theory, with a modern treatment and the positive-solution reduction in
\cite[Section~2]{KT}, whose Schr\"odinger potential has the opposite sign to
$Q$ here; neither the finite threshold nor reduction of order is new. The
derivation is Berarducci and Mantova's \cite{BM}, the ordinal log-atomic
formulas are \cite[Section~2]{ADH}, Hahn real closedness and Neumann's lemma are
standard, and the finite-primitive obstruction is \cref{diff:part:scalar}. The
general logical theory of such solvability is not new either:
\cite[p.~17]{ADH} discusses the predicate $\Omega(a)$ for the existence of a
nonzero solution of $4y''+ay=0$, and the asymptotic differential algebra of
\cite{ADHBook} is neither replaced nor claimed. Mantova's monotonicity and
Taylor results for transseries \cite{Mantova2026} concern functional
composition and are not used or superseded. The main proofs use no model
completeness, no differential closure and no finite-system splitting theorem.
Member~13 proposes as its contributions only the transfinite classification
(\cref{diff:cp:thm:main}), the invisibility theorem
(\cref{diff:cp:thm:invisible}), and especially the limit-stage filtration and
its Galois group (\cref{diff:cp:thm:fieldjump,diff:cp:thm:pv}); the gauge and Schwarzian mechanisms
are standard, and no isolated elementary identity is promoted. Priority is
provisional (item~N103).
\end{remark}

\begin{remark}[A formalization interface]
\label{diff:cp:rem:formalization}
The proofs split into an imported surreal interface --- Conway normal forms,
Hahn summability, the log-atomic derivative identities and the strongly
additive derivation --- and explicit finite differential algebra: the operator
identity, the Euler factorization, Wronskian spanning, exponent-coset
independence and the automorphism calculation. Member~13 recommends isolating
that interface rather than trying to infer the whole derivation from
$\der\omega=1$. No Lean file is supplied (\cref{diff:warn:noformalization}).
\end{remark}

\begin{table}[htbp]
\centering\small
\caption{Proof dependencies of \cref{diff:cp:part}.}
\label{diff:cp:tab:dependencies}
\begin{tabularx}{\textwidth}{@{}L{0.27\textwidth}Y@{}}
\toprule
Conclusion & Necessary inputs, and the check against overreach \\
\midrule
Tower (\cref{diff:cp:prop:tower}) & Imported from \cite{ADH}, limit cases
included. Not derived from $\der\omega=1$. \\
Critical identity (\cref{diff:cp:thm:critical}) & Strong additivity, the
exponential clause, joint summability of the double family. No limit of partial
sums. \\
Nonreal powers (\cref{diff:cp:cor:nonreal}) & \Cref{diff:thm:phasecriterion},
direction (i)$\Rightarrow$(iii); constant field $\C$. No primitive. \\
Euler classification (\cref{diff:cp:prop:euler}) & Wronskian spanning for
$\der_{s}$, real powers, \cref{diff:cp:cor:nonreal}. \\
Classification (\cref{diff:cp:thm:main}) & The gauge identity with
$s=\ell_{\alpha}$. Not a classification of arbitrary potentials. \\
Invisibility (\cref{diff:cp:thm:invisible}) & Scale separation; a set of
ordinals has a strict upper bound. The named hierarchy only. \\
Field jump (\cref{diff:cp:thm:fieldjump}) & Bounded inner support, the monomial
criterion, Hahn real closedness. Finite spans, not Hahn closure, of exponents. \\
Borel group (\cref{diff:cp:thm:pv}) & Coset independence, constants of $\SC$.
No abstract Picard--Vessiot existence theorem. \\
\bottomrule
\end{tabularx}
\end{table}

\part{Coordinate differential equations in common-domain Hahn rings}
\label{diff:part:coordinate}

We now change operators. Throughout \cref{diff:part:coordinate} the
derivative is $\Dz$ of \cref{diff:tab:derivatives}: it differentiates
holomorphic coefficient \emph{functions} in an ordinary complex coordinate
$z$ and holds every Hahn monomial and every scalar parameter constant. No
scalar derivation on the coefficient field is involved in
\cref{diff:sec:coherentring,diff:sec:coherentlinear,diff:sec:coherentnonlinear,diff:sec:monodromy,diff:sec:twoscale};
\cref{diff:sec:formalexistence} then puts the two operators side by side and
proves the compatibility statements that do hold.

The two halves of this report have different constant fields. The general
solution of a coherent linear equation is parametrized by \emph{Hahn}
constants; the general solution of a scalar equation for $\der$ is
parametrized, when it exists at all, by \emph{ordinary complex} constants. A
number can be a coefficient constant without being a differential-field
constant --- $\Dz\omega=0$ while $\der\omega=1$ --- and that difference is the
whole reason the two theories coexist.

\section{The fixed-common-domain ring}
\label{diff:sec:coherentring}

Let $\Gamma$ be a set-sized ordered abelian subgroup of the additive surreals,
interpret $t^{\gamma}=\omega^{-\gamma}$ as in \cref{diff:eq:tmonomial}, and put
\[
 K_{\Gamma}=\C((t^{\Gamma}))
 =\Bigl\{\textstyle\sum_{\gamma\in S}c_{\gamma}t^{\gamma}:
 S\subseteq\Gamma\text{ well ordered},\ c_{\gamma}\in\C\Bigr\},
\]
with valuation $v$ the least exponent and $v(0)=+\infty$; so $v(t^{\gamma})>0$
means the monomial is infinitesimal. Here $\Gamma$ may have higher rank than
$\R$, and \emph{no divisibility assumption on $\Gamma$ is needed} for any
construction in this part. Let $U\subseteq\C$ be an ordinary \emph{nonempty connected}
open set, let $\Hol(U)$ denote the ordinary holomorphic functions on $U$, and
define
\begin{equation}\label{diff:eq:coherentring}
 \Hring_{\Gamma}(U)=\Hol(U)((t^{\Gamma}))
 =\Bigl\{F(z)=\textstyle\sum_{\gamma\in S}f_{\gamma}(z)t^{\gamma}:
 S\text{ well ordered},\ f_{\gamma}\in\Hol(U)\Bigr\}.
\end{equation}
Every coefficient is holomorphic on the \emph{same} $U$. Differentiation is
coefficientwise,
\[
 \Dz F=\sum_{\gamma\in S}f'_{\gamma}(z)\,t^{\gamma},
\]
a derivation of the ring that creates no new exponents; its support is a
subset of the original support. For vectors and matrices, support means the
union of the supports of the finitely many entries. A ``nonzero coefficient
function'' means one not identically zero, so the global support is not
redefined each time a coefficient happens to vanish at a particular $z$.

\begin{warning}[Two neighbours this ring is not, and one notation clash]
\label{diff:warn:ringneighbours}
$\Hring_{\Gamma}(U)$ is \emph{not} the larger ring with an unrelated
convergent germ of shrinking radius for each exponent and no common domain,
and it is \emph{not} the formal power-series ring $\C[[z]]((t^{\Gamma}))$ or
$K_{\Gamma}[[z]]$. These distinctions are already emphasized in the
repository's analytic and foundational reports
\cite{RepoAnalysis,RepoFound,RepoCatalogue}, and
\cref{diff:ex:singular,diff:sub:threepromises} exhibit the separating
witnesses. There is no requirement that the family $(f_{\gamma})$ be uniformly
bounded in $\gamma$, and no numerical series over exponents need converge:
holomorphicity and summability impose different conditions and are used at
different steps. Notation: in \cref{diff:eq:coherentring} the symbol
$\Hol(U)$ denotes ordinary holomorphic functions; it is \emph{not} the
finite-element valuation class $\Ofin$, which never carries an argument.
\end{warning}

Evaluation at an ordinary $z_{0}\in U$ is a ring homomorphism to $K_{\Gamma}$,
coefficientwise, and
\begin{equation}\label{diff:eq:coherentconstants}
 \Ker\bigl(\Dz:\Hring_{\Gamma}(U)\to\Hring_{\Gamma}(U)\bigr)=K_{\Gamma},
\end{equation}
because an ordinary holomorphic function with zero derivative on a connected
$U$ is constant. If $U$ is simply connected, define the normalized primitive
coefficientwise,
\[
 I_{z_{0}}F=\sum_{\gamma}\left(\int_{z_{0}}^{z}f_{\gamma}(\zeta)\,d\zeta\right)
 t^{\gamma},
\]
the ordinary integrals being path independent and holomorphic; they do not
enlarge the support.

A section also has a value at a finite-halo point, after enlarging the
exponent workspace if necessary: for $z_{0}\in U$ and infinitesimal $h$,
\begin{equation}\label{diff:eq:halo}
 F^{\sharp}(z_{0}+h)=\sum_{\gamma\in S}\sum_{n\geq0}
 \frac{f_{\gamma}^{(n)}(z_{0})}{n!}\,t^{\gamma}h^{n}.
\end{equation}
The union of supports lies in $S+(\supp h)^{*}$, and the two-support lemma
together with the ordered-word count of
\cref{diff:lem:neumann,diff:app:support} gives strong summability of the double
family. No uniform bound on the ordinary Taylor coefficients is required. This
is \emph{infinitesimal Taylor evaluation only}, not evaluation at an arbitrary
infinite coordinate.

\section{Linear coherent systems with positive perturbation}
\label{diff:sec:coherentlinear}

\begin{theorem}[Positive-perturbation fundamental matrix]
\label{diff:thm:coherentlinear}
Let $U$ be simply connected, $z_{0}\in U$, and
\[
 A=A_{0}+E\in\Mat_{d}\Hring_{\Gamma}(U),
 \qquad A_{0}\in\Mat_{d}\Hol(U),
 \qquad S=\supp E\subseteq\Gamma_{>0}.
\]
Then there is a \emph{unique} Hahn-coherent matrix $X$ with
\begin{equation}\label{diff:eq:fundamental}
 \Dz X=AX,\qquad X(z_{0})=I_{d}.
\end{equation}
It is invertible in $\Mat_{d}\Hring_{\Gamma}(U)$, its support is contained in
$S^{*}$, and its coefficient at exponent $0$ is the ordinary normalized
fundamental matrix $\Upsilon_{0}$ of $\Dz\Upsilon_{0}=A_{0}\Upsilon_{0}$,
$\Upsilon_{0}(z_{0})=I_{d}$. Uniqueness holds among \emph{all} Hahn-coherent
matrices, not only those whose support was prescribed in advance.
\end{theorem}
\begin{proof}
Ordinary holomorphic linear theory gives $\Upsilon_{0}$ on $U$, normalized at
$z_{0}$, with everywhere nonzero determinant: local existence and uniqueness,
analytic continuation along ordinary paths, and simple connectedness
\cite[Ch.\ 4]{Teschl}. Put $G=\Upsilon_{0}^{-1}E\Upsilon_{0}$, whose support lies in
$S$, and $T H=I_{z_{0}}(GH)$. For each $n$ the support of $T^{n}I_{d}$ lies in
the set of sums of $n$ members of $S$, and at a fixed exponent only finitely
many ordered words in $S$ contribute, by \cref{diff:lem:neumann}. Hence
\begin{equation}\label{diff:eq:peano}
 V=\sum_{n\geq0}T^{n}I_{d}
\end{equation}
is strongly summable with support in $S^{*}$ and satisfies $V=I_{d}+TV$, so
$\Dz V=GV$ and $V(z_{0})=I_{d}$; then $X=\Upsilon_{0}V$ solves
\cref{diff:eq:fundamental}. Writing $V=I_{d}+W$ with $W$ of positive support,
the geometric matrix series $\sum_{n\geq0}(-W)^{n}$ is strongly summable by the
same lemma and is a two-sided inverse, so $X$ is invertible with the claimed
exponent-$0$ coefficient.

For uniqueness let $Z$ be the difference of two coherent solutions and put
$H=\Upsilon_{0}^{-1}Z$. If $H\neq0$, take the least exponent $\gamma$ of its global
support. The equation $\Dz H=GH$ forces $H'_{\gamma}=0$: any contribution on
the right would need a coefficient of $H$ at $\gamma-s<\gamma$ for some
$s\in S$, $s>0$. The initial condition gives $H_{\gamma}(z_{0})=0$, so
$H_{\gamma}=0$ --- a contradiction. Note that this argument never used a
prescribed support monoid.
\end{proof}

The proof yields an explicit recursion, which is a \emph{support certificate}
as well as an existence proof: with $V_{0}=I_{d}$ and all unlisted
coefficients zero,
\begin{equation}\label{diff:eq:linearrecursion}
 V'_{\gamma}=\sum_{s+\delta=\gamma}G_{s}V_{\delta},
 \qquad V_{\gamma}(z_{0})=0 \qquad(\gamma>0),
\end{equation}
where every dependency satisfies $\delta<\gamma$ and every displayed sum is
finite.

\begin{warning}[Not an imported convergence proof]
\label{diff:warn:notpicard}
The familiar iterated-integral form \cref{diff:eq:peano} does \emph{not}
import an ordinary convergence proof into the surreal fine topology. The
primitive is an ordinary coefficientwise primitive and the infinite sum is a
strongly summable Hahn family. At higher rank the successive support levels
need not become cofinal in the value group, and no such cofinality is claimed
or used.
\end{warning}

\begin{corollary}[Inhomogeneous coherent systems]
\label{diff:cor:coherentinhom}
Under the hypotheses of \cref{diff:thm:coherentlinear}, for
$f\in\Hring_{\Gamma}(U)^{d}$ and $c\in K_{\Gamma}^{d}$ the unique solution of
$\Dz y=Ay+f$, $y(z_{0})=c$, is
\[
 y=X\bigl(c+I_{z_{0}}(X^{-1}f)\bigr),
\]
with support contained in
$\bigl(\supp c+S^{*}\bigr)\cup\bigl(\supp f+S^{*}\bigr)$.
\end{corollary}
\begin{proof}
The ordinary product rule holds coefficientwise because each Hahn coefficient
is a finite sum, which verifies the formula; the support follows from the
supports of $X$ and $X^{-1}$ and from $S^{*}+S^{*}=S^{*}$; uniqueness follows
from \cref{diff:thm:coherentlinear} applied to differences.
\end{proof}

\begin{corollary}[Liouville determinant identity]
\label{diff:cor:determinant}
With the same normalization,
\[
 \det X=\det\Upsilon_{0}\cdot\Emm\bigl(I_{z_{0}}(\tr E)\bigr).
\]
Here $\Emm$ is the formal exponential of a \emph{positive-support coherent
section}, not an unrestricted scalar complex exponential.
\end{corollary}
\begin{proof}
Finite matrix algebra gives Jacobi's formula $\Dz\det X=(\tr A)\det X$. The
displayed expression satisfies the same scalar equation and equals $1$ at
$z_{0}$, so the $d=1$ case of \cref{diff:thm:coherentlinear} identifies them.
The exponent has positive support, so \cref{diff:lem:neumann} admits its
exponential.
\end{proof}

\begin{proposition}[A negative leading exponent forbids a coherent scalar exponential]
\label{diff:prop:negativecoherent}
If $a\in\Hring_{\Gamma}(U)$ has $v(a)<0$, then $\Dz Y=aY$ has no nonzero
solution in $\Hring_{\Gamma}(U)$.
\end{proposition}
\begin{proof}
Differentiation can delete coefficient functions but cannot create smaller
exponents, so $v(\Dz Y)\geq v(Y)$. Since $U$ is connected, $\Hol(U)$ is an
integral domain and $v(aY)=v(a)+v(Y)<v(Y)$ for $Y\neq0$. The leading exponents
cannot agree.
\end{proof}

In particular $\Dz Y=i\omega Y$ has no nonzero common-domain Hahn-coherent
solution on \emph{any} connected ordinary domain. The argument excludes not
merely one series representation but every nonzero coherent solution in the
stated ring.

\begin{example}[Positive support is not necessary for matrix systems]
\label{diff:ex:nilpotent}
If $N\in\Mat_d(\C)$ satisfies $N^{2}=0$ and $1\in\Gamma$, then
\[
 \Dz X=t^{-1}NX,\qquad X(0)=I_{d},
 \qquad X=I_{d}+z\,t^{-1}N
\]
is a coherent \emph{polynomial} matrix despite its negative exponent.
Nilpotence permits a finite algebraic cancellation that the scalar argument of
\cref{diff:prop:negativecoherent} excludes. So the positive-support hypothesis
is sufficient, not necessary, and negative-exponent matrix systems need finer
analysis; this example bounds \cref{diff:prop:negativecoherent} deliberately.
\end{example}

\begin{example}[A scale obstruction, and how rescaling resolves it]
\label{diff:ex:singular}
Consider
\begin{equation}\label{diff:eq:singular}
 \Dz Y=t^{-1}Y,\qquad Y(0)=1,
\end{equation}
with $1\in\Gamma$. It has no nonzero solution in $\Hring_{\Gamma}(U)$ on any
ordinary nonempty domain, by \cref{diff:prop:negativecoherent}:
$v(\Dz Y)\geq v(Y)$ while $v(t^{-1}Y)=v(Y)-1$. The formal $z$-series
$\sum_{n\geq0}t^{-n}z^{n}/n!$ \emph{does} exist as an element of
$K_{\Gamma}[[z]]$, but its coefficient-family exponents are unbounded below, so
it is not an element of the fixed-domain ring \cref{diff:eq:coherentring}:
formal existence and coherent existence have separated. Now rescale $z=tw$;
since $t$ is $\Dz$-constant the equation becomes $D_{w}Y=Y$, $Y(0)=1$, whose
solution on an ordinary $w$-domain is the ordinary holomorphic $e^{w}$. That
describes a \emph{microscopic} $z$-domain, not a common ordinary one: at
$v(z)>1$ the original Taylor family is also strongly summable, because $z/t$ is
infinitesimal, whereas at $z=tw$ with nonzero ordinary $w$ one uses the
ordinary coefficient value $e^{w}$ and \emph{not} a Hahn sum of infinitely
many exponent-zero terms $w^{n}/n!$. This is a domain-and-scale obstruction,
not a failure of the formal differential equation, and it is why the
positive-support hypothesis of \cref{diff:thm:ivp} cannot simply be dropped.
\end{example}

\section{A polynomially nonlinear initial-value theorem}
\label{diff:sec:coherentnonlinear}

A nonlinear theorem needs both a domain on which an ordinary background
solution exists and a reason why every new Hahn coefficient is determined by
earlier ones. Positivity of the support supplies the latter. The source
manuscripts prove this theorem twice, in two notations, with two different
admissibility hypotheses on the nonlinearity: polynomial dependence on the
dependent variables, or holomorphic dependence on a neighbourhood of the graph
of the background solution with Taylor evaluation in the vector variable. Both
versions are recorded; the second needs the following admissibility lemma.

\begin{lemma}[Admissibility of nonlinear evaluation]
\label{diff:lem:nonlinear}
Let $d\geq1$, let $y_{0}:U\to\C^{d}$ be holomorphic, let
$W\subseteq\C\times\C^{d}$ be an ordinary open neighbourhood of its graph, let
$S\subseteq\Gamma_{>0}$ be well ordered, and let $F_{0}$ and each $F_{s}$,
$s\in S$, be holomorphic maps $W\to\C^{d}$. Write formally
\begin{equation}\label{diff:eq:vectorfield}
 F(z,Y)=F_{0}(z,Y)+\sum_{s\in S}t^{s}F_{s}(z,Y).
\end{equation}
If $u\in\Hring_{\Gamma}(U)^{d}$ has positive support, the Taylor evaluation of
\cref{diff:eq:vectorfield} at $y_{0}+u$ is well defined and every coefficient
is holomorphic on $U$. If $\supp u\subseteq M$ and $M=E^{*}$ for a positive
well-ordered $E\supseteq S$, then the result has support in $M$.
\end{lemma}
\begin{proof}
Let $T$ be the finite union of the supports of the components of $u$; then
$S\cup T$ is well ordered and positive. A term in a Taylor coefficient has
exponent $s+\beta_{1}+\dots+\beta_{k}$ with $s\in S\cup\{0\}$ and each
$\beta_{j}\in T$. \Cref{diff:lem:neumann} gives a well-ordered union and only
finitely many ordered decompositions at each exponent, and there are finitely
many vector-component choices at each fixed length $k$. So each coefficient is
a finite sum of products of holomorphic derivatives of the $F_{s}$ along
$y_{0}$ with coefficients of $u$, hence holomorphic on $U$. The support
containment follows from closure of $M$ under addition.
\end{proof}

No ordinary analytic convergence theorem is applied to a sum indexed by
surreal exponents: the ordinary Taylor coefficients supply holomorphic
functions, and \cref{diff:lem:neumann} supplies the summation permission.

\begin{theorem}[Support-certified positive-support initial-value problem]
\label{diff:thm:ivp}
In the setting of \cref{diff:lem:nonlinear}, assume $U$ is simply connected,
$z_{0}\in U$, and $y_{0}'(z)=F_{0}(z,y_{0}(z))$. Let $\mathbf{v}\in K_{\Gamma}^{d}$
have positive support, and put
\[
 E=S\cup\supp\mathbf{v},\qquad M=E^{*}.
\]
Then there exists a solution of the form
\begin{equation}\label{diff:eq:coherentivp}
 Y=y_{0}+u,\qquad u\in\Hring_{\Gamma}(U)^{d}\ \text{of positive support},
\end{equation}
to
\begin{equation}\label{diff:eq:ivp}
 \Dz Y=F(z,Y),\qquad Y(z_{0})=y_{0}(z_{0})+\mathbf{v},
\end{equation}
and it is unique. Its support satisfies the explicit certificate
\begin{equation}\label{diff:eq:ivpsupport}
 \supp u\subseteq M\setminus\{0\}.
\end{equation}
\end{theorem}

\begin{convention}[The uniqueness class, stated exactly]
\label{diff:conv:ivpuniqueness}
Uniqueness in \cref{diff:thm:ivp} is uniqueness among coherent solutions of
the form $y_{0}+(\text{positive-support vector})$ on $U$ --- that is, among all
positive-support coherent solutions with the specified ordinary reduction
$y_{0}$. The support containment \cref{diff:eq:ivpsupport} is a
\emph{conclusion}, not a hypothesis: no support was prescribed, and uniqueness
therefore holds among all positive-support coherent solutions and not merely
among those whose support was assumed to lie in $M$. Nothing wider is claimed.
In the \emph{linear} case the least-exponent argument delivers more, and
\cref{diff:thm:coherentlinear} states that stronger conclusion separately:
uniqueness there is among all Hahn-coherent matrices.
\end{convention}

\begin{proof}
Put $A(z)=D_{Y}F_{0}(z,y_{0}(z))$, the Jacobian along the background solution.
Ordinary complex linear theory gives an invertible holomorphic fundamental
matrix $\Upsilon$ on $U$ with $\Upsilon'=A\Upsilon$, $\Upsilon(z_{0})=I_{d}$; simple
connectedness removes monodromy \cite[Ch.\ 3--4]{Teschl}.

$M$ is well ordered by \cref{diff:lem:neumann}. Perform recursion along
$M_{>0}$. For a proposed positive-support $u$, expansion gives
$F_{0}(z,y_{0}+u)=F_{0}(z,y_{0})+Au+R_{0}(z,u)$ where every monomial of
$R_{0}$ has degree at least two in $u$, and each perturbation term
$t^{s}F_{s}(z,y_{0}+u)$ carries the additional positive exponent $s$. Hence at
a fixed exponent $\gamma$ the equation reads
\begin{equation}\label{diff:eq:coefode}
 u'_{\gamma}=A(z)u_{\gamma}+R_{\gamma}(z),
 \qquad u_{\gamma}(z_{0})=\mathbf{v}_{\gamma},
\end{equation}
where $R_{\gamma}$ depends only on coefficients $u_{\beta}$ with
$0<\beta<\gamma$: a contribution from $R_{0}$ has at least two positive
exponents summing to $\gamma$, so each is smaller, and a contribution from a
perturbation has $s>0$ before any factor of $u$, so every remaining exponent is
smaller. The only term containing $u_{\gamma}$ is the isolated linear term
$Au_{\gamma}$.

At a fixed $\gamma$ only finitely many contributions occur: expand each
previously admitted exponent as a finite word in $E$; all contributions then
determine ordered words in $E$ with sum $\gamma$, of which
\cref{diff:lem:neumann} permits only finitely many, and there are finitely many
vector indices at each fixed length. Each $F_{s}$ has finitely many polynomial
monomials, or finitely many relevant Taylor terms at each order by
\cref{diff:lem:nonlinear}, and only finitely many $s$ can participate. So
unbounded degrees across different $s$ do not introduce an infinite
coefficient sum, and \emph{no common bound on the polynomial degrees of the
$F_{s}$ is required}.

By induction $R_{\gamma}$ is an ordinary holomorphic vector on $U$, and the
ordinary linear inhomogeneous problem \cref{diff:eq:coefode} has the unique
holomorphic solution
\begin{equation}\label{diff:eq:coefsolution}
 u_{\gamma}(z)=\Upsilon(z)
 \left(\mathbf{v}_{\gamma}+\int_{z_{0}}^{z}\Upsilon(\zeta)^{-1}R_{\gamma}(\zeta)\,
 d\zeta\right)
\end{equation}
on $U$: the integral is an \emph{ordinary} complex integral of a holomorphic
integrand on a simply connected set, so it is path independent and holomorphic
there. This completes the recursion; the recursion is set-sized and no limit
value at an accumulation exponent is needed, each coefficient being assigned
by its own equation. The resulting support lies in $M_{>0}$, hence is well
ordered, and \cref{diff:lem:nonlinear} makes the substitution summable, so
\cref{diff:eq:ivp} holds coefficientwise.

For uniqueness in the sense of \cref{diff:conv:ivpuniqueness}, let $u$ and
$\widetilde u$ be two positive-support coherent solutions with nonzero
difference. The finite union of their supports is well ordered, so the
difference has a least exponent $\gamma$, below which the coefficients agree.
At $\gamma$ all nonlinear and perturbation difference contributions therefore
cancel, leaving
$(u_{\gamma}-\widetilde u_{\gamma})'=A(z)(u_{\gamma}-\widetilde u_{\gamma})$
with zero initial value; ordinary linear uniqueness makes that coefficient
identically zero, a contradiction.
\end{proof}

\begin{warning}[What the proof does not assume, and what the theorem does not promise]
\label{diff:warn:antibanach}
There is no Banach-space norm on the collection of coefficients, no
contraction in the surreal fine topology, no convergence of Picard iterates,
and no shrinking domain per exponent: every $R_{\gamma}$ and every
$u_{\gamma}$ in \cref{diff:eq:coefsolution} is holomorphic on the
\emph{original} $U$. The recursion uses well-founded support dependence and
ordinary linear uniqueness, and nothing else.

The ordinary background solution is a \emph{hypothesis}. The theorem does not
promise to continue it past its ordinary singularities; its existence on an
arbitrary prescribed domain is not asserted, and it may encounter poles or
other obstructions beyond $U$. Simple connectedness is required. Negative-valued
perturbations and arbitrary radius-free coefficient germs are not covered
(\cref{diff:ex:nilpotent,diff:ex:singular}), and unrestricted substitution into
an arbitrary holomorphic function of $Y$ without specifying its domain and
support conditions is out of scope. This is a differential-equation
specialization of the support-recursion viewpoint already present in the
repository, not a claim to have invented positive-support recursion; its
additional content is the explicit initial-value contract, the common-domain
argument through a single fundamental matrix, and the uniqueness clause of
\cref{diff:conv:ivpuniqueness}.
\end{warning}

\begin{remark}[The linear case, and one practical rule]
\label{diff:rem:linearrule}
For $F_{0}(z,Y)=A_{0}(z)Y+b_{0}(z)$ with perturbations affine in $Y$, the
nonlinear terms vanish and every higher coefficient needs one linear
inhomogeneous solve with the \emph{same} fundamental matrix $\Upsilon$: the
expensive ordinary homogeneous problem is solved once, not once per Hahn
exponent. An effective implementation still needs a way to enumerate the
finitely many decompositions of a requested exponent. A well-ordered support is
a \emph{semantic} certificate, not automatically an effective presentation;
for a finite positive grid, finite truncations often reduce the requirement to
integer-composition enumeration, but for an arbitrary program-described
high-rank support no general terminating enumeration algorithm is claimed.
\end{remark}

\subsection{Two Riccati witnesses}
\label{diff:sub:riccati}

\begin{example}[A coherent entire solution with a pole at an infinite coordinate]
\label{diff:ex:riccatipole}
Take $d=1$, $F_{0}(Y)=Y^{2}$, $y_{0}=0$, no coefficient perturbation, and
initial value $\mathbf{v}=t$. Then
\begin{equation}\label{diff:eq:riccatipole}
 y=\frac{t}{1-tz}=\sum_{n\geq0}t^{n+1}z^{n}
\end{equation}
satisfies $\Dz y=y^{2}$, $y(0)=t$. Every coefficient is a \emph{polynomial} in
$z$, so this is a coherent section on the whole ordinary complex plane. Yet
the rational expression has a pole at the infinite point $z=t^{-1}=\omega$,
which is not in the finite halo of $\C$: evaluating the Hahn series there
would produce infinitely many copies of $t$ and is not a strongly summable
evaluation. So being coherent on the whole ordinary plane does not mean being
regular at every surcomplex coordinate. The location of the pole explains,
rather than contradicts, the exact geometric-series identity.
\end{example}

\begin{example}[An entire coefficient family with an exact recursion]
\label{diff:ex:riccatitan}
Take $F_{0}(z,Y)=Y^{2}$, $y_{0}=0$, $U=\C$, $S=\{1\}$, so
\begin{equation}\label{diff:eq:riccati}
 \Dz Y=Y^{2}+t,\qquad Y(0)=0 .
\end{equation}
\Cref{diff:thm:ivp} produces a unique positive-support coherent solution in
$\Hol(\C)((t^{\Z}))$; writing $Y=\sum_{n\geq1}a_{n}z^{2n-1}t^{n}$, the
coefficient equations give
\begin{equation}\label{diff:eq:riccatirecur}
 a_{1}=1,\qquad (2n-1)a_{n}=\sum_{j=1}^{n-1}a_{j}a_{n-j}\quad(n\geq2),
\end{equation}
so
\begin{equation}\label{diff:eq:riccatiexpand}
 Y=tz+\frac{t^{2}z^{3}}{3}+\frac{2t^{3}z^{5}}{15}
 +\frac{17t^{4}z^{7}}{315}+\frac{62t^{5}z^{9}}{2835}+\cdots,
\end{equation}
every coefficient an entire polynomial. Formally this is
$\sqrt t\,\tan(\sqrt t\,z)$; that notation denotes the displayed Taylor
identity, and the final series has integer exponents even when the chosen
workspace does not contain $\sqrt t$. For finite surcomplex $z$ the argument
$\sqrt t\,z$ is infinitesimal and a small-argument tangent interpretation is
legitimate in an enlarged workspace; but substituting $z=\omega=t^{-1}$ gives
exponents $1-n$, unbounded below, and the family is not summable. Entire
coefficient functions on the ordinary plane do not justify evaluation at every
infinite surcomplex input. This is an exact coherent solution on a common
ordinary domain, not an all-scale surcomplex entire function, and it therefore
does not contradict an all-scale rigidity theorem stated in a different
function class.
\end{example}

\section{Monodromy stays an ordinary-domain construction}
\label{diff:sec:monodromy}

On a non-simply-connected ordinary domain $U$, lift the coefficient system to
its ordinary universal covering surface and choose a lift of the base point.
The proof of \cref{diff:thm:coherentlinear} applies there: ordinary
holomorphic primitives of coefficient one-forms are path independent on the
cover, and the support calculation is unchanged.

\begin{theorem}[Monodromy valued in $\GL_{d}(K_{\Gamma})$]
\label{diff:thm:monodromy}
Let $U$ be a nonempty connected ordinary domain with base point $z_0$, and
let $A=A_0+E$ satisfy the holomorphic and positive-support hypotheses of
\cref{diff:thm:coherentlinear}, with perturbation support $S$, without
assuming that $U$ is simply connected. Analytic continuation of the normalized
fundamental matrix around an ordinary based loop $\ell$ yields another fundamental matrix $X_{\ell}$, and their ratio
has zero coordinate derivative, so
\begin{equation}\label{diff:eq:monodromy}
 X_{\ell}=X\,M_{\ell},\qquad M_{\ell}\in\GL_{d}(K_{\Gamma}).
\end{equation}
The coefficient of $M_{\ell}$ at exponent $0$ is the ordinary monodromy matrix
of $A_{0}$, and all its other exponents lie in $S^{*}\setminus\{0\}$. Under a
fixed convention for composition of based loops these matrices form the usual
monodromy representation; reversing the loop-composition convention reverses
the multiplication order.
\end{theorem}
\begin{proof}
Immediate from \cref{diff:eq:coherentconstants} applied on the cover, together
with $X(z_{0})=I_{d}$, which lets $M_{\ell}$ be read off at the lifted
endpoint; the support statement is \cref{diff:thm:coherentlinear}.
\end{proof}

No fine-continuous surcomplex path is asserted: the loop is an ordinary
complex path, continuation is coefficientwise, and the result is a matrix of
Hahn constants. This is the same domain separation that makes coefficientwise
contour theory meaningful.

\begin{example}[A residue invisible to ordinary monodromy]
\label{diff:ex:infmonodromy}
On a branch domain for the ordinary logarithm containing $1$, consider
\begin{equation}\label{diff:eq:scalarmonodromy}
 \Dz Y=\frac{\lambda+\varepsilon}{z}\,Y,\qquad Y(1)=1,
 \qquad\lambda\in\C,\quad v(\varepsilon)>0 .
\end{equation}
The coherent solution is $Y=z^{\lambda}\Emm(\varepsilon\log z)$, with
$z^{\lambda}=\exp(\lambda\log z)$ in ordinary complex analysis. A positive
loop about $0$ changes $\log z$ by $2\pi i$ and gives
\begin{equation}\label{diff:eq:infmonodromy}
 M=e^{2\pi i\lambda}\,\Emm(2\pi i\varepsilon).
\end{equation}
If $\lambda\in\Z$ the \emph{ordinary} monodromy is $1$, while the full
monodromy is not $1$ when $\varepsilon\neq0$; by
\cref{diff:eq:expinjective} its leading change is exactly
$2\pi i\varepsilon$. So reduction to the ordinary coefficient can erase a
genuine global obstruction to a single-valued coherent solution.
\end{example}

\begin{warning}[The cross-category pair]
\label{diff:warn:crosscategory}
In \cref{diff:eq:scalarmonodromy} the element $\varepsilon$ is constant for
$\Dz$ but need not be constant for the scalar $\der$. It is \emph{not} an
infinitesimal phase residue in the sense of \cref{diff:conv:phases}(b); it is
a Hahn coefficient of a coordinate problem. The finite-primitive criterion of
\cref{diff:thm:phasecriterion} is therefore not the criterion for solving
\cref{diff:eq:scalarmonodromy}. Concretely:
\[
 \varepsilon=\omega^{-1}
 \quad\text{is \emph{permitted}}
 \text{ in the coordinate problem }
 \Dz Y=\tfrac{\lambda+\varepsilon}{z}Y,
\]
\[
 \varepsilon=\omega^{-1}
 \quad\text{is \emph{prohibited}}
 \text{ as the imaginary coefficient of }
 \der y=i\omega^{-1}y .
\]
This is the sharpest available illustration that the two halves of this report
are not competing answers to one problem. They are answers to precisely typed
problems that happen to share differential notation.
\end{warning}

\section{A genuinely two-scale, noncommuting example}
\label{diff:sec:twoscale}

Let $\Gamma=\Z^{2}$ with the lexicographic order and use monomials with
\[
 v(t)=(0,1),\qquad v(u)=(1,0),
 \qquad\text{so}\quad u\prec t^{n}\ \text{for every ordinary }n\geq1 .
\]
This is a Hahn workspace for \emph{coordinate} differentiation: no formula for
$\der u$ is assumed, and none is needed. Set
\[
 N=\begin{pmatrix}0&1\\0&0\end{pmatrix},\qquad
 M=\begin{pmatrix}0&0\\1&0\end{pmatrix},\qquad
 A(z)=tN+uzM .
\]
Both matrices are nilpotent and $NM\neq MN$. The normalized fundamental matrix
has the expansion $X(z)=\sum_{p,q\geq0}t^{p}u^{q}X_{pq}(z)$ with
$X_{00}=I_{2}$, and the recursion
\begin{equation}\label{diff:eq:twoscale}
 X'_{pq}=N X_{p-1,q}+zM X_{p,q-1},
 \qquad X_{pq}(0)=0\quad(p+q>0),
\end{equation}
a negative index meaning zero. Ordinary polynomial integration gives
\[
 X_{10}=zN,\qquad X_{01}=\frac{z^{2}}{2}M,
 \qquad
 X_{11}=z^{3}\left(\frac{NM}{6}+\frac{MN}{3}\right)
 =\begin{pmatrix}z^{3}/6&0\\0&z^{3}/3\end{pmatrix},
\]
and since $\tr A=0$, \cref{diff:cor:determinant} gives $\det X=1$ exactly.

\begin{warning}[The naive exponential is wrong, and total degree is not a valuation cutoff]
\label{diff:warn:twoscale}
It would be incorrect to replace $X$ by $\Emm(tzN+uz^{2}M/2)$: that
expression's mixed coefficient is $z^{3}(NM+MN)/4$, not $X_{11}$, and the
difference at $tu$ is
\[
 \frac{z^{3}}{12}(MN-NM)\neq0 .
\]
The order of multiplication inside iterated integrals matters even though
every individual coefficient calculation is finite. This is exactly why
\cref{diff:app:support} counts \emph{ordered} words: an unordered or multiset
version of the support count would not carry
\cref{diff:thm:coherentlinear,diff:thm:ivp}.

Separately: the delivered verification program computes all $X_{pq}$ with
$p+q\leq6$, checks the coefficient equations and initial conditions, and
verifies $\det X=1$ to the same total degree. \emph{Total degree is not a
valuation cutoff here}: $t^{100}$ has higher total degree but lower valuation
than $u$. The support theorem justifies the complete series; the finite jet is
a different, explicitly declared representation of part of it.
\end{warning}

\section{Formal, microscopic, and coherent: three different promises}
\label{diff:sec:formalexistence}

\subsection{Formal solvability is the weakest of the three}
\label{diff:sub:threepromises}

\begin{proposition}[Formal existence over any characteristic-zero field]
\label{diff:prop:formalrecursion}
Let $K$ be any field of characteristic zero and
$A(Z)=\sum_{n\geq0}A_{n}Z^{n}\in\Mat_{d}(K[[Z]])$. The formal problem
$\derZ X=A(Z)X$, $X(0)=I_{d}$ has a unique solution
$X=\sum_{n\geq0}X_{n}Z^{n}$, given by
\begin{equation}\label{diff:eq:formalrecursion}
 X_{0}=I_{d},\qquad (n+1)X_{n+1}=\sum_{k=0}^{n}A_{k}X_{n-k}.
\end{equation}
Division by $n+1$ is permitted and the right side uses only known
coefficients; no topology and no analytic estimate is involved. In particular
the formal problem is always solvable for $K=\SC$.
\end{proposition}

\begin{warning}[Three promises, worked through $A=i\omega$]
\label{diff:warn:threepromises}
\Cref{diff:prop:formalrecursion} does not assert that the solution can be
\emph{evaluated} at a given surcomplex number. For $A=i\omega$ it gives
\[
 X(Z)=\sum_{n\geq0}\frac{(i\omega Z)^{n}}{n!} .
\]
At any nonzero \emph{ordinary} $z$ the valuations of the terms are unbounded
below, so this family is not strongly summable
(\cref{diff:rem:nonsummableexp}). At an infinitesimal $h$ with $\omega h$
infinitesimal it \emph{is} summable, and equals $\Emm(i\omega h)$. And
common-domain coherence is excluded outright by
\cref{diff:prop:negativecoherent}, which rules out not just this series
representation but every nonzero coherent solution of $\Dz Y=i\omega Y$ in the
stated ring. So
\[
 \text{formal solvability}
 \ \neq\
 \text{microscopic evaluability}
 \ \neq\
 \text{common-domain coherence},
\]
and the three must never be substituted for one another.
\end{warning}

\subsection{The coefficient derivation on a common-domain family}
\label{diff:sub:coefderivation}

The one place where the two operators of this report legitimately meet is the
following. Fix a common-domain family
$F(z)=\sum_{\gamma\in S}f_{\gamma}(z)t^{\gamma}$ as in
\cref{diff:eq:coherentring} and write, for each monomial,
\begin{equation}\label{diff:eq:monomialmatrix}
 \der t^{\gamma}=\sum_{\delta}d_{\delta\gamma}\,t^{\delta},
 \qquad d_{\delta\gamma}\in\R .
\end{equation}
These scalars form a support-controlled matrix; only set-sized columns are ever
used at a time, and it is not a matrix indexed by all surreals as a set.

\begin{lemma}[Finite incoming contributions]
\label{diff:lem:finiteincoming}
For every well-ordered exponent set $S$, the union
$T_{S}=\bigcup_{\gamma\in S}\supp\der t^{\gamma}$ is well ordered, and for
each $\delta$ only finitely many $\gamma\in S$ have
$d_{\delta\gamma}\neq0$.
\end{lemma}
\begin{proof}
The family $(t^{\gamma})_{\gamma\in S}$ is summable. Apply
\ref{diff:BM:additive}, including its clause that the \emph{derived} family is
summable; the two conclusions are exactly the two clauses of
\cref{diff:def:summable}.
\end{proof}

\begin{theorem}[Differentiation preserves common-domain coherence]
\label{diff:thm:coefderivation}
The rule
\begin{equation}\label{diff:eq:coefderivation}
 \dercoef F
 =\sum_{\delta\in T_{S}}
 \left(\sum_{\gamma\in S}d_{\delta\gamma}f_{\gamma}\right)t^{\delta}
 \qquad\Bigl(\text{equivalently }
 \dercoef F=\sum_{\gamma\in S}f_{\gamma}(z)\,\der t^{\gamma}\Bigr)
\end{equation}
defines a $\C$-linear derivation with values in the common-domain Hahn
algebra over the differential hull of $\Gamma$. Every
inner sum is finite; each new coefficient is holomorphic on the \emph{same}
ordinary domain $U$, so differentiating the scalar coefficients neither
shrinks nor destroys the common domain; and
\begin{equation}\label{diff:eq:commutederivations}
 \dercoef\Dz F=\Dz\dercoef F .
\end{equation}
If the exponent group $\Gamma$ is differentially stable
(\cref{diff:def:monomialclosed}) the values stay in
$\Hring_{\Gamma}(U)$; for an arbitrary initial $\Gamma$ the same assertions
hold after passing to its differential Hahn hull (\cref{diff:thm:hull}).
\end{theorem}
\begin{proof}
\Cref{diff:lem:finiteincoming} gives the well-ordered output support and the
finiteness of each inner sum, so each new coefficient is a finite real-linear
combination of functions holomorphic on the same $U$, hence holomorphic there.
Linearity is immediate. For the product rule, expand the Hahn product of two
families and use $\der(t^{\gamma+\eta})=t^{\gamma}\der t^{\eta}+t^{\eta}\der t^{\gamma}$;
standard product summability and \ref{diff:BM:additive} make the expanded
families summable, and coefficient comparison then uses only finite sums.
Finally ordinary differentiation is $\C$-linear and commutes with each finite
inner sum in \cref{diff:eq:coefderivation}, which gives
\cref{diff:eq:commutederivations}. Stability of $\Gamma$, or passage to the
hull, keeps the supports inside the workspace.
\end{proof}

\begin{warning}[The derivative may need a larger monomial workspace]
\label{diff:warn:largerworkspace}
$\dercoef$ need \emph{not} preserve a previously chosen smaller monomial
group: the support of $\der t^{\gamma}$ can contain monomials outside a
preselected subgroup (\cref{diff:ex:sqrt2}). Its representing matrix
$(d_{\delta\gamma})$ is only row-finite on each specified summable input
support; no claim is made that an individual derivative has finite support or
that the matrix has a finite global representation. A derivative routine
should therefore return both a value \emph{and} the parent in which that value
is represented.
\end{warning}

\begin{theorem}[Total intrinsic derivative of an evaluation]
\label{diff:thm:totalchainhalo}
Let $F$ be a common-domain family, let $c\in U$ and $h\in\mm_{\C}$, and let
$F^{\sharp}(c+h)$ be the finite-halo evaluation \cref{diff:eq:halo}. Then
\begin{equation}\label{diff:eq:totalchainhalo}
 \der\bigl(F^{\sharp}(c+h)\bigr)
 =(\dercoef F)^{\sharp}(c+h)
 +\der h\cdot(\Dz F)^{\sharp}(c+h).
\end{equation}
Equivalently, for $z\in\Ofin_{\C}$ with $c=\st(z)\in U$ and $h=z-c$,
$\der(\mathrm{ev}_{z}F)=\mathrm{ev}_{z}(\dercoef F)
+\mathrm{ev}_{z}(\Dz F)\,\der z$. All terms may be placed in one set-sized
differentially stable Hahn workspace after adjoining $\supp h$ and taking the
hull (\cref{diff:thm:hull}).
\end{theorem}
\begin{proof}
Apply \ref{diff:BM:additive} to the summable double family
\cref{diff:eq:halo}. Since $c$ is an ordinary complex point, all the numbers
$f_{\gamma}^{(n)}(c)/n!$ are constants of $\der$, and $\der h=\der z$. The
product rule splits the derivative of each term into
\[
 (\der t^{\gamma})\frac{f_{\gamma}^{(n)}(c)}{n!}h^{n}
 \qquad\text{and}\qquad
 t^{\gamma}\frac{f_{\gamma}^{(n)}(c)}{(n-1)!}h^{n-1}\der h
 \quad(n\geq1),
\]
the second omitted for $n=0$. The first family is summable by
\cref{diff:lem:finiteincoming} followed by the same positive-support Taylor
argument; the second is the corresponding Taylor family for $\Dz F$ multiplied
by the fixed scalar $\der h$. So the split and the rearrangement are
legitimate (\cref{diff:warn:splitting} is satisfied), and reindexing gives the
two terms of \cref{diff:eq:totalchainhalo}.
\end{proof}

An ordinary holomorphic lift $f^{\sharp}$ has $\der$-constant coefficients, so
\cref{diff:eq:totalchainhalo} reduces to
$\der(f^{\sharp}(z))=(f')^{\sharp}(z)\der z$. A general coherent family carries
the additional coefficient term, and this distinction is independent of
whether the evaluated value has a short expression.

\begin{example}[Two worked total derivatives]
\label{diff:ex:totalhalo}
(i) $F(z)=t^{2}e^{z}$ at $z=t$: here $\dercoef F=-2t^{3}e^{z}$,
$\Dz F=t^{2}e^{z}$, $\der t=-t^{2}$, and
\[
 \der\bigl(t^{2}\Emm(t)\bigr)=(-2t^{3}-t^{4})\,\Emm(t),
\]
the extra $-t^{4}\Emm(t)$ being the moving-argument term. Differentiating only
the scalar prefactor, or only the analytic argument, would miss part of the
answer.
(ii) $F(Z)=\omega Z^{2}+tZ$ at $z=1+t$: here $\dercoef F=Z^{2}-t^{2}Z$,
$\Dz F=2\omega Z+t$, $\der z=-t^{2}$, and
\[
 F(z)=\omega+2+2t+t^{2},
 \qquad
 \der(F(z))=1-2t^{2}-2t^{3},
\]
which \cref{diff:eq:totalchainhalo} reproduces exactly; keeping only the
argument derivative would lose the leading constant $1$. The family
$F(Z)=\omega$ already disproves that omission.
\end{example}

\subsection{Fixed contours, fixed centres, and moving poles}
\label{diff:sub:contours}

\begin{corollary}[Coefficientwise linear naturality]
\label{diff:cor:naturality}
Any fixed $\C$-linear operation on ordinary coefficient functions commutes
with $\dercoef$ wherever that operation is defined coefficientwise. In
particular, for a \emph{fixed} ordinary piecewise smooth contour
$\gamma_{c}\subseteq U$,
\begin{equation}\label{diff:eq:contourcommute}
 \der\!\left(\int_{\gamma_{c}}F(z)\,dz\right)
 =\int_{\gamma_{c}}\dercoef F(z)\,dz,
 \qquad
 \int_{\gamma_{c}}F\,dz:=\sum_{\gamma}
 \left(\int_{\gamma_{c}}f_{\gamma}\,dz\right)t^{\gamma},
\end{equation}
and the analogue holds for a coefficientwise Laurent residue at a
\emph{fixed} ordinary centre, whenever all coefficient functions belong to the
stated punctured-domain class.
\end{corollary}
\begin{proof}
Apply the fixed linear operation to the finite inner sums of
\cref{diff:eq:coefderivation}: both orders of application give the coefficient
$\sum_{\gamma}d_{\delta\gamma}L(f_{\gamma})$ at $\delta$.
\end{proof}

No surcomplex-valued Riemann integral is introduced: each side of
\cref{diff:eq:contourcommute} is a coefficientwise functional. For example, on
a punctured ordinary disk about $a\in\C$ the family $F(z)=t/(z-a)$ has contour
integral $2\pi it$ around a positive simple loop, and both sides of
\cref{diff:eq:contourcommute} equal $-2\pi it^{2}$; this is an exact Hahn
identity.

\begin{warning}[A moving contour or a moving pole is a different statement]
\label{diff:warn:movingpole}
The contour in \cref{diff:cor:naturality} is fixed and ordinary. A moving
contour, moving endpoints, or a residue with a different centre convention
requires additional terms or a different statement; none is hidden in the
corollary. At an \emph{actual moving pole} the coefficient term cannot be
dropped. Use the formal rational operators $\dertil$ and $\derZ$ of
\cref{diff:prop:totalchain}: for $a\in\SC$ and $F(Z)=1/(Z-a)$,
\begin{equation}\label{diff:eq:movingpole}
 \dertil F=\frac{\der a}{(Z-a)^{2}},\qquad
 \derZ F=-\frac{1}{(Z-a)^{2}},\qquad
 \der\frac{1}{z-a}=-\frac{\der z-\der a}{(z-a)^{2}},
\end{equation}
these rational identities holding whenever the denominators are nonzero and
needing no coherent expansion. A fixed-centre residue must not be silently
replaced by an actual residue at $a$. The repository already distinguishes
those residues; intrinsic differentiation supplies another reason to keep them
apart. Likewise, the formal Hahn residue reading of
\cref{diff:warn:diffnotint} is not an identification with a cluster residue or
an actual-point contour residue.
\end{warning}

\begin{remark}[A pointwise-only evaluation identity]
\label{diff:rem:pointwiseonly}
For a \emph{fixed} ordinary $z\in U$, strong additivity does give
\[
 \der\bigl(F(z)\bigr)=\sum_{\gamma}f_{\gamma}(z)\,\der(t^{\gamma}),
\]
a pointwise scalar identity. Turning it into a statement about a common-domain
coherent family, differentiating again in $z$, or integrating it along a path
requires separate control of the coefficient-family supports and of the domain
--- which is precisely what
\cref{diff:thm:coefderivation,diff:thm:totalchainhalo,diff:cor:naturality}
supply, under their hypotheses. The distinct coefficient rings documented in
the repository remain distinct after $\der$ is added.
\end{remark}

\subsection{Formal power series need their own admissibility condition}
\label{diff:sub:formaladmissibility}

Let $K$ be a differential subfield of $\SC$ (or take $K=\SC$).
On $K[[Z]]$ the formal coefficient and variable derivatives commute:
\[
 \dertil F=\sum_{n}(\der c_{n})Z^{n},
 \qquad
 \derZ F=\sum_{n\geq1}n c_{n}Z^{n-1},
 \qquad
 \dertil\derZ F=\derZ\dertil F .
\]
Evaluation at a scalar $a$ requires summability. \emph{If} the families
defining $F(a)$, $(\dertil F)(a)$ and $(\derZ F)(a)$ are all summable, then the
same computation gives
\begin{equation}\label{diff:eq:formaladmissible}
 \der\bigl(F(a)\bigr)=(\dertil F)(a)+\der a\cdot(\derZ F)(a).
\end{equation}
\Cref{diff:prop:chain} is the sharper version, which assumes only the
summability of the evaluated family and \emph{derives} the other two.

\begin{warning}[Do not apply this to an unevaluable germ]
\label{diff:rem:coefstream}
\Cref{diff:eq:formaladmissible} must not be applied to an arbitrary
unevaluable formal germ. The finite-halo theorem
\cref{diff:thm:totalchainhalo} supplies its own summability proof; an
unrestricted formal coefficient sequence does not. The limiting case is
\cref{diff:warn:smallexplimits}(iii): the evaluation rule $f(h)$ of
\cref{diff:lem:smallchain} requires \emph{ordinary constant} coefficients, and
$\sum_{n\geq0}\omega^{n^{2}}X^{n}$ is not summable at $X=\omega^{-1}$ because
its monomials $\omega^{n^{2}-n}$ are not reverse well ordered. The
infinitesimal chain rule must not be exported to arbitrary
$\SC$-coefficient streams without a joint support certificate. On the specific
ring $\Hol(U)((t))$ one may define $\dercoef=-t^{2}D_{t}$ and $\Dz$
coefficientwise, and they commute because both sides have identical
coefficient formulas; that statement need not hold as a meaningful endomorphism
assertion on an arbitrarily chosen $K_{\Gamma}$, since the scalar derivative
can require monomials outside that workspace
(\cref{diff:warn:largerworkspace}).
\end{warning}

\subsection{The four categories, side by side}
\label{diff:sub:fourcategories}

\begin{table}[htbp]
\centering\small
\caption{Four typed categories that share differential notation. They are not
competing answers to one problem.}
\label{diff:tab:categories}
\begin{tabularx}{\textwidth}{@{}L{0.18\textwidth}L{0.11\textwidth}YY@{}}
\toprule
Category & Constants & Constructive positive result & Boundary \\
\midrule
Scalar $(\SC,\der)$ & $\C$ &
Finite-primitive criterion, finite-phase solution formula, gauge normal form,
constant-coefficient classification &
$\der y=iy$ has no nonzero solution; no global compatible exponential \\[0.4em]
Real-exponent Hahn scalar $K_{\Gamma}$, $\Gamma\subseteq\R$, $1\in\Gamma$ &
$\C$ &
$\der-\lambda$ bijective for $\lambda\in\C^{\times}$; $P(\der)$ bijective on
$\C((t))$ when $P(0)\neq0$; exact residuals &
$t$ has no primitive there; the real-exponent hypothesis is substantive \\[0.4em]
Formal $K[[Z]]$ & $K$ &
Linear coefficient recursion, always solvable &
Evaluation may not be summable \\[0.4em]
Common-domain $\Hring_{\Gamma}(U)$ & $K_{\Gamma}$ &
Positive-perturbation linear and polynomially nonlinear initial-value
theorems, support certificates, monodromy in $\GL_{d}(K_{\Gamma})$ &
A negative scalar leading exponent prevents a coherent exponential; nilpotence
is an exception \\
\bottomrule
\end{tabularx}
\end{table}
\part{Computation, formalization, and the register of non-claims}
\label{diff:part:limits}

\section{Certificate-oriented computational contracts}
\label{diff:sec:computation}

\subsection{Theorem-driven outputs, not a universal solver}
\label{diff:sub:certificates}

The exact criterion suggests a usable symbolic interface. Given a represented
coefficient $q=a+ib$, a backend may try to produce primitives $A,B$ with
$\der A=a$ and $\der B=b$ together with a certified decomposition
$B=P+r+\eta$ with $P\in\Pcl$, $r\in\R$, $\eta\in\mm$. It need not use the
normalization $\Izero$ to do so, since $\pp(B)=\Obs(b)$ is unchanged by the
real integration constant (\cref{diff:rem:obsnormalization}). The logical
workflow is
\begin{equation}\label{diff:eq:workflow}
\begin{gathered}
 q=a+ib\ \longmapsto\ B\ \longmapsto\ \pp(B)\ \longmapsto\\[0.25em]
 \begin{cases}
 \text{a base-field solution},&\pp(B)=0,\\
 \text{a certified obstruction and an extension symbol},&\pp(B)\neq0 .
 \end{cases}
\end{gathered}
\end{equation}

\begin{warning}[This is a specification, not a terminating algorithm]
\label{diff:warn:notalgorithm}
\Cref{diff:eq:workflow} is a mathematical specification. For arbitrary
program-described surreal normal forms, integration, exact support access and
zero testing may all be unavailable. A failed search for a primitive is not a
proof that none exists --- in the full field every element has one --- and a
failed representation of an infinitesimal exponential is not a proof that no
mathematical solution exists. A backend should return a well-defined symbolic
primitive and state that the purely infinite part has not been decided; it
must not report ``no solution'' merely because its representation cannot
compute that part. The repository's capability-tier distinctions are therefore
especially important at this differential boundary \cite{RepoCAS}.
\end{warning}

Three logically distinct outcomes arise for the ambient question, and a fourth
once a workspace is declared.

\begin{table}[htbp]
\centering\small
\caption{Result contracts. The last row is not a nonexistence theorem.}
\label{diff:tab:outcomes}
\begin{tabularx}{\textwidth}{@{}L{0.22\textwidth}YL{0.26\textwidth}@{}}
\toprule
Result & Certificate required & Mathematical conclusion \\
\midrule
Workspace solution &
Rational exponent $r$, an element $k\in\Hring$ with positive support and
$-rt+\der k=q$ &
A certified nonzero solution $t^{r}\Emm(k)$ lies in the declared Hahn field
(\cref{diff:thm:Hahnlogder}) \\[0.4em]
Ambient-only solution &
A verified primitive $A$ of $a$, a \emph{finite} accumulated phase $B$ of $b$
with verified $\pp(B)=0$, an admissible infinitesimal-residue expression,
and certified failure of the declared workspace solvability test &
A nonzero solution $\exp(A)\Emm(i\eta)$ exists in $\SC$, but the workspace
obstruction is nonzero (\cref{diff:thm:phasecriterion}) \\[0.4em]
Ambient phase obstruction &
A verified accumulated phase $B$ of $b$ with certified nonzero $\pp(B)$ &
No nonzero surcomplex solution exists; the gauge normal form returns $P$ and
the removable factor $g$ (\cref{diff:thm:gauge}) \\[0.4em]
Undetermined representation &
A required primitive, decomposition, equality or support certificate is
unavailable &
\emph{No} existence or nonexistence conclusion \\
\bottomrule
\end{tabularx}
\end{table}

For the finite exact core --- finite rational Hahn polynomials with
Gaussian-rational coefficients, represented as a finite map
$q\in\Q\mapsto c_{q}\in\Q(i)$ with zero coefficients deleted --- the first
three outcomes are decidable. \Cref{diff:prop:hahnderiv} gives a terminating
derivative; additive integration inside $\Hring$ succeeds exactly when the
exponent-$1$ coefficient vanishes, and otherwise
\cref{diff:eq:logprimitive} gives an ambient primitive using one explicit
$\log\omega$ constructor. Two tests must be offered \emph{separately}: the
\emph{workspace} test rejects every nonzero term below exponent $1$ and
rejects a nonreal exponent-$1$ coefficient (with Gaussian-rational
coefficients a real exponent-$1$ coefficient is automatically rational), and
the \emph{ambient} test examines the imaginary part $b$ and succeeds precisely
when $b=0$ or every exponent of $\supp b$ exceeds $1$. For arbitrary
computably described Hahn supports none of this is advertised as a total
decision procedure, and the semantic formula $\Phi(q)=\pp(\Izero\Imm q)$
remains exact without promising an algorithm for $\Izero$ or for equality of
arbitrary represented scalars.

\begin{warning}[What a certificate must actually establish]
\label{diff:warn:certificatecontent}
A solution certificate consists of exact derivative identities plus support
and size evidence. For a finite expression the identities may reduce to field
algebra plus the real exponential and logarithm rules. For an infinitesimal
Taylor expression the certificate must identify positive support and justify
the summations used in the derivative rule. For a \emph{nonsolvability}
certificate it is not enough to report a large leading term numerically: one
must prove that an accumulated phase has nonzero purely infinite part. The
output $t^{r}\Emm(k)$ is a finite exact expression denoting an infinite Hahn
series, not a finite list of its coefficients.

A finite-phase certificate need not compute the whole normalized primitive: it
suffices to exhibit a finite $B$ and prove $\der B=\Imm q$. This is useful when
semantic existence is known but a general symbolic integration procedure is
unavailable.

The two inputs $q_{1}=i\omega^{-2}$ and $q_{2}=i\omega^{-1}$ are
distinguished by the exact accumulated phases $-\omega^{-1}$ and $\log\omega$.
A numerical specialization $\omega:=M$ for a large real $M$ would instead give
ordinary complex equations with nonzero oscillatory solutions in \emph{both}
cases. Such a specialization is not an embedding of differential structures and
cannot certify the surreal answer.
\end{warning}

\subsection{A derivation-aware representation interface}
\label{diff:sub:interface}

Every exact parent should record both its ambient embedding and its closure
capabilities. A Laurent parent may advertise: closed under $\der$; closed
under $(\der-\lambda)^{-1}$ for $\lambda\neq0$; \emph{not} closed under
primitives. A log-exponential extension can support more primitives without
thereby supporting arbitrary global phases, and adding a constructor named
\texttt{Exp} is not enough: its domain and its derivative contract must be
stated.

\begin{table}[htbp]
\centering\small
\caption{Operations, required certificates, and result boundaries.}
\label{diff:tab:interface}
\begin{tabularx}{\textwidth}{@{}L{0.24\textwidth}YL{0.27\textwidth}@{}}
\toprule
Object or request & Required certificate & Result boundary \\
\midrule
Intrinsic derivative &
A selected derivation and a parent field, with a supported image in a known
parent &
May require a larger parent: a monomial derivative need not lie in the old one
(\cref{diff:warn:largerworkspace}) \\[0.35em]
Coordinate derivative &
One common coefficient domain &
Every Hahn monomial and scalar parameter stays constant \\[0.35em]
Hahn evaluation &
A joint reverse-well-ordered support with finite multiplicity &
A formal expression may not denote a sum at the input
(\cref{diff:rem:nonsummableexp}) \\[0.35em]
Coherent total derivative &
One common coefficient domain and \emph{both} derivatives &
Dropping $\dercoef$ changes the answer
(\cref{diff:ex:totalhalo}) \\[0.35em]
Finite-phase integrating factor &
A real accumulated phase $B$ with verified $\pp(B)=0$ &
A nonzero solution in $\SC$ \\[0.35em]
Phase obstruction &
A verified nonzero purely infinite part of an accumulated phase &
Nonexistence of a nonzero homogeneous solution \\[0.35em]
Hahn resolvent &
A real-exponent group with $1\in\Gamma$, and a support-shift certificate; a
lower exponent bound for a prefix &
An exact Hahn solution \emph{in that workspace}; prefix computation is not
stream equality \\[0.35em]
Coherent initial-value solver &
Base solution, common domain, positive support, finite decomposition access &
A certified coefficient recursion \\[0.35em]
Global phase &
An explicit phase convention with its supported domain &
A group law does not imply compatibility with $\der$
(\cref{diff:ex:phaseconvention}) \\[0.35em]
Oscillatory extension &
A distinct parent and the rule $\der\mathsf T=i\mathsf T$ &
Not a surcomplex scalar (\cref{diff:cor:nonembedding}) \\
\bottomrule
\end{tabularx}
\end{table}

\begin{warning}[A truncation must record the \emph{type} of omitted information]
\label{diff:warn:truncationtype}
A truncation of $Y_{H}$ in \cref{diff:eq:factorialsecond} can carry the exact
residual \cref{diff:eq:factorialresidual} and a real-power support tail. That
certificate does \emph{not} rule out an exponentially small homogeneous
correction in the full field (\cref{diff:prop:factorialambiguity}). A
computation that returns ``unique solution'' must say whether uniqueness is in
$K_{\Gamma}$, in the full $\SC$, in an oscillatory extension, or among
coherent sections of a named ring. Likewise, returning
$\exp\bigl(\int A\bigr)$ is not a certificate for a complex integrating
factor: the software must either verify the finite-phase criterion, select a
differential extension, or report that the required operation is not
implemented. An unevaluated expression and a proven element of the requested
field are different outputs. A matrix jet specified by total degree, a Laurent
jet specified by valuation, and an exact recursively defined series must remain
three different output types
(\cref{diff:warn:twoscale}). Branch continuation and monodromy should use
ordinary-domain path data rather than an implied fine-continuous path in the
surreal class (\cref{diff:thm:monodromy}). For several extension symbols, the
integer relation lattice of \cref{diff:eq:lattice} is additional required data.
\end{warning}

A realistic first implementation can stay inside an explicit differential
algebra: rational expressions in a finite collection of real-power and
logarithmic scales; finite constant-coefficient operators; formal Laurent
series represented by a recurrence or by coefficient access. At each extension
a derivative routine should return both a value and its parent, and an
integration routine should distinguish a proved primitive, a proved
obstruction \emph{in the selected parent}, and a failure of the available
algorithm to decide. The missing $t$-coefficient of
\cref{diff:eq:hahnimage} is a proved parent obstruction; failure to recognize
a complicated primitive in the full surreal field is not.

\subsection{Formalization boundaries}
\label{diff:sub:formalization}

A proof-assistant development can begin from a set-sized ordered field with a
derivation, a chosen infinite scale $w$ with $\der w=1$, real constants, and
specified local exponential and logarithm operations with their support
theorems. The bounded-derivative argument
(\cref{diff:lem:smallnegative}), the complexification
(\cref{diff:prop:complexification}), the phase decomposition
(\cref{diff:thm:unitphase}) and the first-order criterion
(\cref{diff:thm:phasecriterion}) then become lemmas over explicit hypotheses.
Surjectivity of the derivation is a \emph{genuine additional hypothesis}, not
something inferred from real closedness.

Four concrete proof obligations are worth naming, in increasing order of
difficulty.
\begin{enumerate}[label=(\arabic*),leftmargin=2.4em]
\item The explicit rational monomial derivative formula
\cref{diff:eq:hahnderiv} on $\Hring=\C((t^{\Q}))$. The primitive-image and
logarithmic-image theorems
(\cref{diff:thm:Hahnintegral,diff:thm:Hahnlogder}) then reduce to support
translation, leading-term factorization and the infinitesimal exponential
theorem. This is a tractable independent target.
\item The implication ``monomial support closure $\Rightarrow$ closure of the
full small-support Hahn type under $\der$'', which is the content of
\cref{diff:lem:monomialcriterion,diff:thm:hull}. It complements, rather than
duplicates, a generic Hahn convolution development.
\item \Cref{diff:lem:finiteincoming}: for each admissible input support, each
output coefficient receives only finitely many monomial-derivative
contributions. Once available, the common-domain coefficient interface becomes
finite linear algebra at every output exponent.
\item The represented quotient $\Pcl$ of \cref{diff:warn:cosets}, so that the
rank-one classification does not accidentally require a type whose elements
are proper classes.
\end{enumerate}
The interfaces should keep $\der$, $\derZ$, $\dertil$, $\dercoef$, $\Dz$ and
fine differentiability separate. Evaluation is a map with a support
certificate, and \cref{diff:prop:chain} is a compatibility theorem for that
map; it is not a coercion identifying operators. An exponential operation
should carry the domain predicate $\Imm z\in\Ofin$ whenever it promises the
intrinsic differential law. The difficult Berarducci--Mantova existence
construction can remain an explicit imported interface until separately
formalized; it is a much larger dependency and is not needed to check the
logical consequences of a stated interface.

\begin{warning}[Formalization plan and present scope]
\label{diff:warn:noformalization}
The report-specific differential theorems remain pending in the repository
ledger \path{docs/FORMALIZATION.md}. Existing checked prerequisites there
must be distinguished from a formalization of the Berarducci--Mantova
derivation, its phase obstruction, the finite-matrix classification or the
coherent differential-equation package. The preceding paragraphs specify
further proof obligations and a possible order of development; they do not
assert that the concrete surreal carrier has already been connected to every
required Hahn, real-closed-field, exponential or derivation interface.
\end{warning}

\section{Relation to the literature and to the companion reports}
\label{diff:sec:literature}

The proof has three layers. At the foundation are the established surreal
field and exponential constructions and the chosen strongly additive,
surjective derivation \cite{BM,Gonshor}. The normal-form splitting and the
infinitesimal formal calculus connect that structure to the repository's
finite-angle geometry \cite{RepoTrig}. Above them are the explicitly proved
differential consequences: the finite-primitive criterion, the gauge normal
form, the maximal strip, the scalar and constant-coefficient classifications,
the workspace images, and the coordinate initial-value theorems.

The universal $H$-field results \cite{ADH} explain why the differential
viewpoint is a major part of surreal theory rather than an ad hoc extra
operation; through \cref{diff:part:workspaces} they are cited for context and
are not a substitute for an unstated linear-solvability theorem. In
\cref{diff:part:matrix} that changes: the elementary embedding of \cite{ADH}
is used as a transfer principle, and the splitting and surjectivity facts of
\cite[Lemma 2.4.19]{ADHNorm} are used as premises
(\cref{diff:thm:matsurjectivity}, item~N53). The amplitude--phase results of
\cite{ADHSecond} over Hardy fields, and the classical Pinney parametrization
\cite{Pinney}, are the antecedents of
\cref{diff:thm:matpinney,diff:thm:matamplitude} and are kept visible there; no
solution of the Boshernitzan conjecture named in \cite{ADHSecond} is claimed.
The $\omega$-series composition results
\cite{BMGerms} show that an external function interpretation is valid on a
substantial specified class while warning against unrestricted composition
(\cref{diff:warn:composition}). The global exponential construction \cite{EK}
supplies the crucial contrasting structure: algebraic global exponentiation
and number-derivation compatibility are different requirements
(\cref{diff:warn:notdenied}). Bournez and Guilmant \cite{BG} study closure
properties of selected surreal fields far beyond the elementary constructions
of \cref{diff:sec:localization}, and no theorem of theirs is used to justify
them. Ordinary complex initial-value theory \cite{Teschl} supplies the
analytic ingredient of \cref{diff:part:coordinate}, and
\cite{PS,Singer} supply the Picard--Vessiot and differential-module
terminology in their standard sense. Neumann \cite{Neumann} and Higman
\cite{Higman} are the published sources for \cref{diff:app:support}.

The repository documents are treated as unrefereed project manuscripts that
motivate the gap, not as replacements for the published differential
foundations; none of their analytic-geometry or contour theorems is needed for
the main proofs here. The following cross-links matter.
\begin{itemize}[leftmargin=1.6em]
\item \textbf{Trigonometry.} That report supplies the finite-angle phase group
and explicitly declines to choose a derivation \cite{RepoTrig}.
\Cref{diff:thm:cisgroup} reproves the part needed here and adds the
derivative identity; \cref{diff:warn:notdenied} states precisely what the
present no-go theorems do and do not say about its global phase extensions.
Its existing constructions need not be removed or renamed as false; their
contracts need only be distinguished from the intrinsic differential one, and
the derivation should be named wherever it is used.
\item \textbf{Analysis.} That report's derivation-obstruction theorem, for one
chosen canonical global surcomplex exponential, is \emph{prior coverage}
\cite{RepoAnalysis}. \Cref{diff:cor:noglobalPsi} strengthens it in
phase-independence and does not contradict it; the two are compatible.
\Cref{diff:part:coordinate} reuses its common-domain coherent-section
language and refines the germ-versus-common-domain-versus-formal distinction
into the explicit separations of \cref{diff:tab:categories}.
\item \textbf{Computer algebra and foundations.} The capability taxonomy and
the subsection ``Differential equations must name their category''
\cite{RepoCAS}, and the four-interface discipline of \cite{RepoFound}, are
answered directly by \cref{diff:tab:derivatives,diff:tab:outcomes,diff:tab:interface}
and by the tiering contract of \cref{diff:sub:certificates}.
\item \textbf{Spectral theory.} The companion report \emph{Surcomplex Spectral
Theory} \cite{RepoSpectral} supplies the finite algebraic spectral vocabulary
that \cref{diff:sub:matweyl} uses, and its ``spectrum'' is the algebraic
eigenvalue and singular-value list. That is \emph{not} the differential phase
spectrum of \cref{diff:conv:phases}(e), and
\cref{diff:warn:spectrum} states the distinction in the terms both reports
need. \Cref{diff:thm:matspectral} is the bridge between them for Hermitian
coefficients; \cref{diff:ex:matnonnormal} shows there is no bridge in general.
\textbf{Cross-link both ways.}
\item \textbf{Hahn--Tate uniformization.} That report records that curves with
nonnegative valuation of $j$ need formal groups around a good-reduction model
\cite{RepoTate}, and it makes no assertion about the Berarducci--Mantova
derivation. \Cref{diff:aut:prop:abeliansplit} treats constant good-reduction
abelian varieties over $\SC$ by their formal groups, and
\cref{diff:aut:thm:abelianimage} is the collection's first contact between
surreal derivations and abelian varieties; it is a cross-link, not an answer to
that report's question over Hahn fields with a coarse valuation
(\cref{diff:aut:rem:tate}).
\item \textbf{Tail-spans.} The Euler derivation $\derT$ of \cref{diff:rs:part} is
the operator that report uses at the scale $\log t$ \cite{RepoTailSpans}.
\Cref{diff:aut:cor:independent} concerns algebraic independence of
differentially \emph{algebraic} solutions and is a different notion from that
report's differential transcendence; there is no conflict.
\item \textbf{Classical regular-singular and autonomous theory.} Levelt's Jordan
decomposition \cite{Levelt}, its exposition for formal connections \cite{KW}, and
generalized-series operator methods \cite{vdH} are the antecedents of
\cref{diff:rs:part}; the geometric formalism of autonomous first-order equations,
their local normal forms and Rosenlicht's theorem \cite{NPT} are the antecedents
of \cref{diff:aut:part}, with \cite{Milne,Stacks} for the algebraic geometry.
\item \textbf{Classical oscillation theory and the log-atomic tower.} The
finite critical hierarchy of \cref{diff:cp:part} is classical
Kneser--Weber--Hartman--Hille theory, treated in \cite{KT}; the ordinal
log-atomic formulas are \cite[Section~2]{ADH}; the solvability predicate
$\Omega$ of \cite[p.~17]{ADH} and the asymptotic differential algebra of
\cite{ADHBook} are the general theory the family instantiates; and Mantova's
composition results \cite{Mantova2026} are neither used nor superseded
(\cref{diff:cp:rem:inputs}).
\end{itemize}

\section{The register of non-claims}
\label{diff:sec:nonclaims}

This section is part of the result. Each item below was stated explicitly by
at least one of the eleven source manuscripts, or is an artifact-level finding of
the audit of a source package and is marked as such; all of them survive the
merge. Several are also stated in place above, and the cross-references say
where. Items N1--N52 come from the first seven manuscripts, N53--N64 from the
eighth, N65--N79 from member~12, N80--N94 from member~11 and N95--N106 from
member~13; where a later
member changed the status of an earlier item, the item is amended in place and
its original wording is kept.

\subsection{Recovered on audit}
\label{diff:sub:nc:recovered}
An audit of this merge against its seven sources found two scope disclaimers
stated in a source and missing here.  They are restored, since each limits a
claim the report otherwise appears to make.

\begin{enumerate}[label=\textup{R\arabic*.},leftmargin=3.0em,itemsep=0.35em]
\item \textbf{The interface table proposes contracts.} One source explicitly
distinguishes its certificate and interface specifications from implemented
functions. Accordingly, \cref{diff:tab:interface} is a proposed contract, not
evidence that a complete differential solver exists. The companion
computer-algebra report supplies the capability taxonomy; current checked
prerequisites and pending obligations are tracked separately in
\path{docs/FORMALIZATION.md}. This historical source disclaimer is not an
audit asserting that no related implementation exists anywhere in the repository.
\item \textbf{A recorded run is not a reproduction promise.}  One source notes
that a local environment's Python version and elapsed time need not match its
bundled run.  \Cref{diff:app:verification} reproduces the delivered records
and states the versions they were recorded under; it does not assert that a
rerun elsewhere will reproduce those versions or timings, and a mismatch in
either is not evidence of a failed check.
\end{enumerate}

\subsection{Provenance of the mathematics}
\label{diff:sub:nc:provenance}

\begin{enumerate}[label=\textup{N\arabic*.},leftmargin=3.0em,itemsep=0.35em]
\item The Berarducci--Mantova derivation is \emph{imported}, not constructed
and not reproved (\cref{diff:thm:BM}). Its construction via log-atomic
numbers, a pre-derivation there, path derivatives and nested truncation ranks
is not reproduced. The summability and surjectivity proofs are substantial
imported inputs, not consequences of the product rule alone. Assigning a
derivative to a few familiar monomials is not a construction of a strongly
additive derivation on every surreal, and this report uses the published
conclusions rather than a pretended implementation.
\item No uniqueness characterization of $\der$ is offered
(\cref{diff:warn:nouniqueness}). Other surreal derivations exist; none of them
is classified here; the listed structural properties do not characterize every
possible surreal derivation; and the word ``canonical'', where it survives, is
a naming convention for this distinguished construction. \emph{Amended with
member~11:} \cref{diff:aut:part} does not classify them either. It proves, for
the class of normalized derivations of \cref{diff:aut:def:normalized}, that one
sector --- first-order autonomous equations with constant coefficients --- is
the same for all of them; that is a statement about a class of derivations, not
a classification of its members, and no uniqueness of a normalized derivation
is claimed (item~N80).
\item Every intrinsic statement is relative to this one fixed $\der$ with the
normalization $\der\omega=1$, which is indispensable for discussing a named
equation such as $\der y=iy$. The results are \emph{not} derivation-agnostic.
\emph{Amended with member~11:} this holds for every part except
\cref{diff:aut:part}, whose statements are proved for every normalized
derivation and flag the two that need surjectivity
(\cref{diff:aut:conv:scope}). No theorem of another part is upgraded; the
first-order statements of \cref{diff:cor:oscillator,diff:cor:riccati} are
reproved for every normalized derivation there (\cref{diff:aut:cor:earlier}),
and the second-order statements are not. Member~12's hypotheses
(\cref{diff:rs:rem:imports}) are clauses of the normalized-derivation
definition, but \cref{diff:rs:part} is stated, and here kept, for the fixed
$\der$ only; see \cref{diff:sub:nc:open}.
\item Smallness is \emph{derived} here from the $H$-field positivity clause
(\cref{diff:lem:smallnegative}), and is therefore not advertised as part of
the imported theorem. Two of the source manuscripts instead import it, one of
them flagging it as specific to the chosen derivation and not a general fact
about derivations of non-Archimedean fields; that route is recorded in
\cref{diff:rem:sqrtsmallness}. Either way, smallness is a structural
hypothesis about this derivation.
\item The positive-support summability lemma is imported in its standard
additive-exponent form (\cref{diff:lem:neumann}) and is additionally proved in
\cref{diff:app:support}. One source audit records that the Neumann paper was
\emph{not retrieved in full} and that the publisher page did not provide
accessible full text in that session; the article therefore imports the
standard lemma rather than claiming to reprove or re-audit the original paper,
and \cref{diff:app:support} exists so that nothing depends on an uninspected
passage. Higman's publication metadata was checked against the publisher.
Gonshor's book is used as a standard reference for normal forms and the
surreal exponential; only its publisher catalogue was identified, and the
whole book was not downloaded or audited.
\item The restricted composition, Taylor and global-incompatibility results
for $\omega$-series are \emph{separately} imported
(\cref{diff:warn:composition}); they do not follow from \cref{diff:thm:BM} and
are not hypotheses hidden inside the phase criterion. The composition
obstruction is quoted with its hypotheses intact, for this particular
derivation, and is \emph{not} strengthened into an assertion about every
conceivable notion of composition or about every other choice of derivation.
Nothing here settles alternative global composition structures. The
finite-phase obstruction of this report is a direct differential-algebraic
argument, not a restatement or strengthening of that composition theorem.
\item The universal $H$-field results are cited for context, not as a
substitute for an unstated linear-solvability theorem, and the elementary
embedding of transseries is a theorem about a specified ordered differential
language which must not be upgraded to a statement about every simultaneously
added analytic, compositional or complex exponential operation
(\cref{diff:rem:fourinferences}). Bournez--Guilmant is cited as contextual
work only; no theorem of it justifies
\cref{diff:thm:hull,diff:thm:localize}.
\item The class-theoretic framing is an \emph{assumption}
(\cref{diff:sub:classes}): definable classes, or a class-theoretic
presentation with enough replacement for the stated constructions. The exact
sequences are elementwise and are not the formation of a set of proper-class
cosets (\cref{diff:warn:elementwise}).
\end{enumerate}

\subsection{Novelty, refereeing, verification}
\label{diff:sub:nc:novelty}

\begin{enumerate}[label=\textup{N\arabic*.},leftmargin=3.0em,itemsep=0.35em,start=9]
\item No priority or novelty claim is made, and no claim to resolve a named
published open problem. These are proved deductions from explicitly stated
hypotheses; some may admit formulations already familiar in the general theory
of $H$-fields, and the deductions are \emph{not} asserted to be absent from
every published treatment of $H$-fields, transseries or differential algebra.
Ordinary constant-coefficient differential algebra and formal resolvents have
a long history; their specialization here is exposition.
\item No independent refereeing and no peer review is claimed. The source
manuscripts were research expositions prepared with AI assistance; this merge
is of the same character.
\item No proof-assistant verification, no Lean source, no machine-checked
proof object (\cref{diff:warn:noformalization}). The general proofs remain
human-readable mathematical arguments.
\item The literature consultation was targeted, not an exhaustive novelty or
priority search, and the bibliography is correspondingly targeted.
\item The seven repository audits were \emph{targeted coverage audits}, pinned
to one revision, and read-only. None was a line-by-line verification of the
fifteen reports or of their archived sources; none asserts that the word
``derivation'', or Berarducci--Mantova, or transseries, is absent from every
source manuscript; none assesses the correctness of the repository's other
claims. The reading scopes recorded in the sources include the documentation
catalogue and manifest; the analysis, trigonometry, foundations and
computer-algebra READMEs; the opening foundational section of the analysis
source; selected computer-algebra source ranges; and lines 1--220 of the
trigonometry article. Several of the audits record that code-search or
connector calls returned no usable matches or were truncated, and that those
empty or partial responses were \emph{not} used as evidence of absence and not
treated as a complete inventory. The assessments are revision-specific: no
claim is made that the identified gap remains unfilled on a later commit or
branch. The repository's existing phase decomposition and chosen-exponential
obstruction are acknowledged as prior coverage, not presented as missing
results. No repository file was modified by any of the seven source packages,
and the placement suggestions they made were suggestions. The later, equally
scoped audits of members~11 and~12, pinned to \CommitC, are items N93 and N77,
and that of member~13, pinned to \CommitD, is item~N104.
\end{enumerate}

\subsection{Effectivity: what is exact but not computable}
\label{diff:sub:nc:effectivity}

\begin{enumerate}[label=\textup{N\arabic*.},leftmargin=3.0em,itemsep=0.35em,start=14]
\item $\Izero$ is a normalization defined by a uniqueness property, not an
algorithm (\cref{diff:warn:Jnotalgorithm}). It is not asserted to be
multiplicative, not asserted to preserve arbitrary Hahn sums, and not asserted
to be effective for arbitrary surreal inputs or computable from a supplied
symbolic representation. No implementation of it on arbitrary
program-described surreals, no universal zero test, and no full
differential-algebraic solver is supplied.
\item $\Obs$ and $\Phi$ are exact but not presented as computable functions on
arbitrary finitely described surreal inputs. The semantic formula
$\Phi(q)=\pp(\Izero\Imm q)$ is exact without promising an algorithm for
$\Izero$ or for equality of arbitrary represented scalars. An unevaluated
primitive, or an undecided finiteness test, is \emph{not} a negative
solvability proof (\cref{diff:warn:notalgorithm}).
\item \Cref{diff:eq:workflow} is a specification, not a terminating algorithm.
Only the first three outcomes of \cref{diff:tab:outcomes} are claimed
decidable, and only for the finite exact core; for arbitrary computably
described Hahn supports they are not advertised as total decision procedures.
\item The rational decision procedure is restricted to $\Q(X)$ although
\cref{diff:thm:rational} is stated over $\R(X)$, and it \emph{raises} on
unsupported input rather than reporting mathematical nonexistence
(\cref{diff:warn:rationalimpl}). No algorithm for an arbitrary real
coefficient is asserted.
\item A well-ordered support is a \emph{semantic} certificate, not
automatically an effective presentation
(\cref{diff:rem:linearrule}). No general terminating enumeration algorithm for
arbitrary high-rank program-described supports is claimed. A closed formula
for a coefficient stream is not a finite list of every coefficient, and the
coefficient recursion of \cref{diff:eq:coefrecursion} does not make equality
of arbitrary program-generated streams decidable.
\item $\dercoef$ is represented by a matrix that is only \emph{row-finite on
each specified summable input support}; no claim is made that an individual
monomial derivative has finite support, or that the matrix has a finite global
representation, or that $\dercoef$ preserves a previously chosen smaller
monomial group (\cref{diff:warn:largerworkspace}).
\item \Cref{diff:thm:hull} is a semantic existence construction, not an
effective algorithm for computing $\der(t^{\gamma})$ at arbitrary surreal
exponents, and its countably many stages do not constitute a terminating
algorithm for every input representation
(\cref{diff:warn:hulllimits}).
\end{enumerate}

\subsection{Scope of the scalar theory}
\label{diff:sub:nc:scalarscope}

\begin{enumerate}[label=\textup{N\arabic*.},leftmargin=3.0em,itemsep=0.35em,start=21]
\item The criterion classifies \emph{scalar homogeneous first-order}
multiplicative equations, and nothing more. It is not a claim that all
nonlinear equations have solutions, that $\SC$ is differentially closed --- it
is not (\cref{diff:rem:algnotdiff}) --- that arbitrary nonlinear scalar
equations or all higher-order operators on the full class are classified, or
that primitives are computable from finite input descriptions. \emph{Amended
with member~11:} one nonlinear class is now classified.
\Cref{diff:aut:thm:main,diff:aut:thm:independence} describe every solution in
$\SC$ of every first-order autonomous equation $F(y,y')=0$ with $F\in\C[Y,Z]$,
and more generally of every constant meromorphic vector field on a smooth
projective curve, for every normalized derivation. The item stands for
everything beyond that class: variable coefficients in $\SC\setminus\C$, higher
order, simultaneous systems and higher-dimensional autonomous systems
(items~N81--N82), and $\SC$ remains not differentially closed.
\emph{Amended with member~13:} one explicit family of linear second-order
equations with variable coefficients is classified at every ordinal depth,
$\der^{2}y+(Q_{\alpha}+c(\ell_{\alpha}^{\dagger})^{2})y=0$
(\cref{diff:cp:thm:main}). The item stands for every other higher-order
operator, and in particular for arbitrary perturbations of that family
(item~N99).
\item No general inhomogeneous existence theorem is claimed \emph{by
\cref{diff:part:scalar}}. Outside the admissible locus $\Phi(q)=0$ its only
assertion is \emph{uniqueness}
(\cref{diff:thm:variation,diff:warn:atmostone}); the absence of a nonzero
homogeneous solution does not imply inhomogeneous nonexistence, and the
additional differential-field structure that a general inhomogeneous existence
theorem would need is not invoked there. Nor is it claimed that such broader
solvability questions are open in the literature. \Cref{diff:thm:variation}
proves surjectivity of $\der-q$ only in the admissible case, and no claim is
made \emph{by that scalar argument} about surjectivity of all linear
differential operators on $\SC$: a
trivial kernel is explicitly distinguished from a failure of surjectivity, and
$\der-i$ on $\C((t))$ is the witness (\cref{diff:rem:parentselection}).
\Cref{diff:part:matrix} does prove first-order inhomogeneous existence
(\cref{diff:eq:matfirstorder}) and, from it, existence for every forcing of
every finite system (\cref{diff:cor:matinhom}); that conclusion rests
\emph{entirely} on the separately imported
\cref{diff:thm:matsurjectivity} and not on the rank-one criterion, as
\cref{diff:warn:matinhom} states. It leaves the present item standing for the
scalar theory as such, it does not make any homogeneous equation solvable, and
it changes none of the witnesses of \cref{diff:warn:atmostone}: what it adds is
that their uniqueness is now known to be accompanied by existence. No
effectivity follows (item~N57). \emph{Amended with member~12:} for power-Hahn
regular-singular systems and forcing in $\KR^{n}$, inhomogeneous existence holds
\emph{without} that import, with a solution in $\KR[\log t]^{n}$
(\cref{diff:rs:thm:forced,diff:rs:rem:fredholm}); for arbitrary surcomplex
forcing nothing is inferred there (item~N68).
\item The constant-coefficient theorem and its matrix corollary concern
\emph{ordinary complex constant} coefficients and matrices only, not arbitrary
surcomplex-coefficient systems (\cref{diff:rem:realsolutions}). The kernel
theorems classify \emph{homogeneous} solutions in $\SC$ only; a full kernel
does not settle an inhomogeneous equation over a smaller subfield.
\emph{Amended with member~12:} the item stands as written. The variable-coefficient
class of \cref{diff:rs:part} is a different statement, for the Euler derivation
$\derT$ and power-Hahn coefficients $A_{0}+A_{+}$, and it does not extend the
constant-coefficient theorem to arbitrary surcomplex coefficients; its own
limits are items N65--N79. \emph{Amended with member~13:} the item still
stands. \Cref{diff:cp:prop:euler} is the order-two constant-coefficient
classification read for a rescaled derivation $\der_{\log s}$
(\cref{diff:cp:rem:constant}); it does not extend \cref{diff:thm:constant} to
arbitrary coefficients, and its limits are items N95--N106.
\item \Cref{diff:thm:matnormal} gives a complete phase normal form for every
ordinary finite system, including non-Hermitian ones, and \cref{diff:thm:matspectral} computes
its invariant from ordinary algebraic eigenvalues \emph{when the coefficient is
$iH$ with $H$ Hermitian}. This formula does not extend to arbitrary
coefficients: \cref{diff:ex:matnonnormal} gives a
nonnormal coefficient with instantaneous eigenvalues $\pm i$ whose system
nevertheless has a fundamental matrix in $\GL_{2}(\SC)$, so that integrating
imaginary parts of eigenvalues would predict $\{\omega,-\omega\}$ instead of
the correct $\{0,0\}$. Within the scalar theory of \cref{diff:part:scalar}
taken alone, none of this is available: there the trace condition is only the
first invariant (\cref{diff:prop:trace}), the torus theorem treats
\emph{diagonal} systems only (\cref{diff:thm:torus}), constant matrices are
settled while arbitrary matrices with surreal entries are not, and no general
matrix normal form follows from the scalar result --- the normal form of
\cref{diff:part:matrix} needs the separate import of
\cref{diff:thm:matsurjectivity}. The normal form covers every finite rank over
$\SC$, but does not compute the normalizing gauges or phase obstructions
on arbitrary input descriptions. Effective procedures on specified input
languages require additional information beyond instantaneous eigenvalues
(\cref{diff:warn:matnonnormal,diff:warn:matlimits}).

For power-Hahn regular-singular systems, $\Ph(-tA)=\{(\Imm\lambda_j)\log t\}$,
with the $\lambda_j$ the eigenvalues of the ordinary residual matrix. The
normalizing gauge is explicit and lies in $\GL_n(\KR[\log t])$
(\cref{diff:rs:prop:phasespectrum}); within this class the module splits into
rank-one factors with obstructions read off from that matrix
(\cref{diff:rs:rem:effective}). \Cref{diff:ex:matnonnormal} lies outside the
class. Explicit computation for broader input languages remains open.
\item No differential closedness, no universal Picard--Vessiot closure, no
classification of all differential extensions, and no decision procedure for
arbitrary nonlinear differential equations is claimed. The Picard--Vessiot
constructions of \cref{diff:sec:PV} are explicit; the set-sized ones
deliberately avoid any blanket application of Picard--Vessiot existence theory
to a proper-class base, and the proper-class ones are proved by hand for
exactly that reason. The extension $\C(\omega)(\mathcal E)$ of
\cref{diff:cor:nonembedding} is used only to exclude a differential embedding,
and no assertion whatever is made about its full constant field.
\emph{Amended with member~11:} for presented constant curves and meromorphic
vector fields with exact algebraic coefficients, \cref{diff:aut:rem:decision}
gives a finite geometric decision criterion for the existence of nonconstant
solutions, conditional on exact arithmetic and decidable sign tests
(item~N90). No decision procedure for arbitrary nonlinear differential equations
is claimed.
\item \Cref{diff:thm:maximalstrip}'s maximality does not claim that the
differential law alone uniquely specifies every normalization on $\Strip$;
ratios of admissible values are differential constants.
\item All nonexistence results concern surcomplex \emph{number} solutions
under $\der$ (\cref{diff:warn:whatnonexistence}). They do not deny ordinary
oscillatory functions, do not say that $\sin z$ and $\cos z$ fail their
ordinary differential equations, and do not deny algebraically defined global
surcomplex phase extensions or the Ehrlich--Kaplan canonical global surcomplex
exponentials, which satisfy different structural contracts and should not be
described as nonexistent (\cref{diff:warn:notdenied}). The obstruction is only
to \emph{also} satisfying pointwise differential compatibility for this
$\der$.
\item No general composition law on $\No$ is constructed and unrestricted
substitution ``replace $\omega$ by $x$ in every surreal'' is never used
(\cref{diff:rem:nocomposition}).
\item Splitting a differentiated infinite sum into a coefficient-derivative
part and a variable-derivative part requires \emph{both} separated families to
be summable independently (\cref{diff:warn:splitting}); summability of the sums
does not imply it.
\item $\ct$ is not a multiplicative standard part
(\cref{diff:warn:ct}), and the real-power rule $\der\omega^{r}=r\omega^{r-1}$
is asserted for ordinary real $r$ only and must not be extended by silently
replacing $r$ with an arbitrary surreal exponent
(\cref{diff:warn:monomialmap}).
\item The infinitesimal exponential's injectivity is \emph{local} to
$\mm_{\C}$ and is not a global injectivity claim; the infinitesimal evaluation
rule requires \emph{ordinary constant} coefficients and must not be exported
to arbitrary $\SC$-coefficient streams without a joint support certificate
(\cref{diff:warn:smallexplimits,diff:rem:coefstream}).
\item Strong summability is not valuation convergence and is not
fine-topological convergence (\cref{diff:warn:summability}). No cofinal growth
of $n\sigma$ in a higher-rank value group is established, and no argument uses
one. The truncation residual identities certify a formal remainder at a stated
order, not convergence of the truncation sequence in the fine surreal topology, and
the factorial Hahn series is not an ordinary analytic function germ and not a
fine-topology limit of its truncations
(\cref{diff:warn:antitopological}). Numerical substitution is declined as a
surrogate for a Hahn identity.
\end{enumerate}

\subsection{Scope of the workspaces}
\label{diff:sub:nc:workspacescope}

\begin{enumerate}[label=\textup{N\arabic*.},leftmargin=3.0em,itemsep=0.35em,start=33]
\item Neither localization theorem gives closure under \emph{all} set-indexed
summable families (\cref{diff:warn:localizelimits}): an arbitrary-support
closure is a different and stronger demand, a localized field need not contain
every Hahn series over its own monomials, and there is no Hahn completeness, no
finite data representation, no algorithm for the adjoined maps, no birthday or
cardinal bound, no computable enumeration and no termination claim for
symbolic integration. Arbitrary future coherent families are not automatically
covered. \Cref{diff:thm:localize} is not proved to produce a full Hahn field
with a fixed exponent group, and the lighter variant of
\cref{diff:rem:finitaryclosure} is proved to be neither.
\item \Cref{diff:thm:hull}'s minimality holds only among full Hahn fields
indexed by divisible exponent groups, not among all differential subfields;
the hull group itself need not be countable; and the theorem is not primitive
closure, nor closure under real exponentiation, arbitrary logarithms, or
solutions of all first-order linear equations
(\cref{diff:warn:hulllimits,diff:ex:logmissing}).
\item The sharp support threshold is deliberately restricted to real-exponent
Hahn coefficients, and the criterion provably does \emph{not} globalize to
$\No$ by valuation alone (\cref{diff:warn:noglobalization}). The valuation
test of \cref{diff:cor:valphase} is not a valuation-only classification of
arbitrary surreal coefficients. The threshold must not be merged into the
general resolvent.
\item \Cref{diff:thm:resolvent}'s real-exponent hypothesis is substantive: in
a higher-rank group a fixed increment $1$ need not dominate all relevant
scales, and no claim replaces the support argument by an assertion that
successive valuation improvements converge
(\cref{diff:warn:resolventreal}). \Cref{diff:prop:polyresolvent} is confined to
the rank-one Laurent field with $P(0)\neq0$, and its residual bound is a
certificate for a finite truncation: it does not assert that the truncation
equals the full solution.
\item The factorial examples assert no Borel summation, no Stokes multiplier
and no analytic continuation rule
(\cref{diff:warn:antitopological,diff:prop:factorialambiguity}); the formal
Hahn residue reading of \cref{diff:warn:diffnotint} is not an identification
with a cluster residue or an actual-point contour residue from the
repository's analytic reports.
\item The finite-rank classification over $\SC$ does not automatically
descend to a prescribed smaller workspace: a phase-normalizing gauge or a
required inhomogeneous solution can leave that workspace. A classification
over such restricted bases is not supplied here.
\end{enumerate}

\subsection{Scope of the coordinate theory}
\label{diff:sub:nc:coordinatescope}

\begin{enumerate}[label=\textup{N\arabic*.},leftmargin=3.0em,itemsep=0.35em,start=39]
\item The coherent theory requires \emph{positive} support of the
perturbation. That hypothesis is sufficient, not necessary: the nilpotent
example \cref{diff:ex:nilpotent} shows a coherent solution with a negative
exponent, and negative-exponent matrix systems need finer analysis
(\cref{diff:warn:antibanach}).
\item \Cref{diff:thm:ivp} requires either polynomial dependence on the
dependent variables or holomorphic dependence with the Taylor admissibility of
\cref{diff:lem:nonlinear}, together with a \emph{supplied} ordinary background
solution on $U$ and $U$ simply connected. The background solution's existence
on an arbitrary prescribed domain is not asserted, the theorem does not
continue it past its ordinary singularities, negative-valued perturbations and
arbitrary radius-free coefficient germs are not covered, and unrestricted
substitution into an arbitrary holomorphic function of $Y$ without specifying
its domain and support conditions is out of scope. Uniqueness is exactly the
class of \cref{diff:conv:ivpuniqueness} and nothing wider.
\item No general irregular-singularity classification, no Stokes
classification and no Riemann--Hilbert correspondence over arbitrary Hahn
fields is offered.
\item Monodromy (\cref{diff:thm:monodromy}) is an \emph{ordinary}-domain
construction on the ordinary universal cover with coefficientwise
continuation; no fine-continuous surcomplex path is asserted. Monodromy for
the \emph{nonlinear} coherent problem on non-simply-connected domains remains
future work.
\item The finite-halo evaluation \cref{diff:eq:halo} is infinitesimal Taylor
evaluation only, not evaluation at an arbitrary infinite coordinate
(\cref{diff:ex:riccatipole,diff:ex:riccatitan}).
\item The analytic interface is stated for a \emph{fixed} ordinary open
domain and a \emph{fixed} ordinary contour or residue centre. Moving contours,
moving endpoints, moving poles and other residue conventions are explicitly
outside it (\cref{diff:warn:movingpole}), and statements are deliberately not
transported between arbitrary formal germs and fixed-domain holomorphic
families.
\item Formal existence, microscopic evaluability and common-domain coherence
are three different promises and are never substituted for one another
(\cref{diff:warn:threepromises}). The formal recursion asserts nothing about
evaluability at a given surcomplex number.
\item The Riccati example of \cref{diff:ex:riccatitan} gives an exact coherent
solution on a common ordinary domain, not an all-scale surcomplex entire
function; it does not contradict an all-scale rigidity theorem stated in a
different function class.
\item \Cref{diff:thm:coherentlinear} assumes no norm on coefficients, no
contraction, no convergence of Picard iterates and no shrinking domain per
exponent (\cref{diff:warn:antibanach}); nothing here claims to have invented
positive-support recursion.
\end{enumerate}

\subsection{Scope of the executable checks}
\label{diff:sub:nc:checks}

\begin{enumerate}[label=\textup{N\arabic*.},leftmargin=3.0em,itemsep=0.35em,start=48]
\item The first eight delivered verification suites check \emph{finite exact
symbolic identities and finite jets}, and nothing else. In total they record
$1{,}395+115+149+258+291+150+259+62$ passing checks
(\cref{diff:app:verification}); a matrix identity or a whole checked
coefficient prefix counts as one check in two of the suites. The suites of
members~11 and~12 record $10$ and $8$ passing checks and carry their own scope
items, N94 and N78--N79, and the suite of member~13 records $70$, with scope
items N105--N106; every limitation stated in the rest of this item applies
to them as well. They do not
construct the surreal numbers, do not implement Conway normal forms or the
Berarducci--Mantova path construction, do not verify arbitrary Hahn
summability or strong additivity for arbitrary supports, do not verify
class-size or proper-class statements, do not verify the Hahn--Neumann support
lemma or Higman's lemma, do not verify transfinite recursion, and do not
verify any of the general nonexistence or infinite-support theorems. They are
regression suites for examples, not proof assistants and not implementations
of a general surreal differential solver.
\item The suites model $\der$ by ordinary symbolic differentiation: as
$-t^{2}\,d/dt$ on $t$-expressions, as $d/dx$ on expressions in an ordinary
variable standing in for $\omega$, and as $f\mapsto\mathrm{cancel}(i\mathsf T\,
\partial f/\partial\mathsf T)$ on the oscillatory extension. The primitive
tests replace $\omega$ by an ordinary positive variable and check finite
algebraic differential identities; they do not infer that an ordinary
exponential is a scalar-compatible surcomplex exponential. A finite Taylor
check is a formal coefficient check up to its stated order.
\item Declared cutoffs are part of the claim, not incidental: one suite uses a
formal-series cutoff of $16$ with a deterministic seed; another uses matrix
\emph{total degree} $6$ --- explicitly not a valuation cutoff, since
$t^{100}$ has higher total degree and lower valuation than $u$ --- a Laurent
reciprocal-series order of $12$, and nonlinear checks through $t^{12}$;
another checks resolvent residuals for $\lambda$ in a listed finite set and
$N$ in a listed finite range; another checks factorial residuals to
$N=16$. A successful check of a residual at finitely many orders is not a proof
of its general formula, and symbolic agreement of finite Taylor coefficients is
not a proof of Hahn summability.
\item The restricted rational classifier accepts exact Gaussian-rational or
rational coefficients and rejects floats, strings, non-Gaussian-rational
coefficients, unbound parameters, and non-rational functions by \emph{raising};
those rejections are counted as passing checks in one suite, and the script is
explicit that it is not an expression parser
(\cref{diff:warn:rationalimpl}).
\item Two artifact-level discrepancies in the source packages are recorded
rather than silently repaired: one package's PowerShell build script was
supplied but not executed on Windows in the preparation environment, and one
package's README lists a checksum file that is not present among the delivered
files. Neither affects any mathematical statement.
\end{enumerate}

\subsection{Scope of the matrix theory}
\label{diff:sub:nc:matrixscope}

The eighth manuscript states twelve explicit non-claims of its own. All survive
here, together with the audit findings recorded against its package.

\begin{enumerate}[label=\textup{N\arabic*.},leftmargin=3.0em,itemsep=0.35em,start=53]
\item \Cref{diff:thm:matsurjectivity} is an \emph{import}, not a result of this
report. The splitting of monic scalar operators into first-order factors and
the surjectivity of nonzero scalar operators are taken from
\cite[Lemma 2.4.19]{ADHNorm} and transferred along the elementary embedding of
\cite{ADH}; neither is reproved. The Berarducci--Mantova derivation and that
elementary embedding are likewise imported wholesale
(items~N1 and~N7). ``Linearly closed'' in the cited source means the splitting
property and \emph{not} that every homogeneous equation has a full solution set
in the base field; reading it the second way would contradict
\cref{diff:cor:oscillator} (\cref{diff:warn:matimport}).
\item The rank-one phase obstruction, the bounded-primitive ideal, the
diagonal Galois tori and the universal exponential algebra used in
\cref{diff:part:matrix} are existing repository or literature theory, reproved
or specialized where a proof was wanted, and are not discoveries of the eighth
manuscript; its own README and its comparison table say so
(\cref{diff:rem:matagreement,diff:rem:matrankone,diff:sub:mattorus}). The
classical Pinney parametrization \cref{diff:eq:matpinneyclassical} and the
abstract Schwarz-closed-field existence principle behind it are likewise
disclaimed as antecedents, not contributions
\cite{Pinney,ADHNorm,ADHSecond}.
\item No named published open problem is claimed solved by
\cref{diff:part:matrix}, and in particular no new solution of the Boshernitzan
conjecture discussed in \cite{ADHSecond} is claimed. Its own novelty
assessment is provisional and is narrowed to three statements --- the
arbitrary-dimensional Hermitian spectral--integral formula
\cref{diff:eq:matspectral}, the sharp entrywise perturbation criterion
\cref{diff:cor:matperturbation} with \cref{diff:prop:sharpperturb}, and the
normalized non-Abelian integration theorem
\cref{diff:thm:matnonabelian} --- on the basis of a \emph{targeted} literature
search that did not locate an exact prior statement. That is evidence for
treating them as candidate new results, not a priority determination, and this
report makes no novelty claim of its own (item~N9).
\item No proof-assistant verification, no Lean formalization and no
independent refereeing of \cref{diff:part:matrix} is claimed; its strongest
claims require independent mathematical review. This repeats items~N11
and~N12 for the eighth member, whose own text states it in those words.
\item Non-effectivity. The obstruction map $\Obs$ and the normalizing gauges
$G$ and $W$ are exact \emph{semantic} operations, as $\Izero$ is
(\cref{diff:warn:Jnotalgorithm}). Even for finite matrices, arbitrary surreal
input carries no assumed finite computable encoding, and no factorization
algorithm follows from the first-order transfer in
\cref{diff:thm:matsurjectivity}. Identifying restricted, support-certified
input languages for which eigenvalue primitives, their purely infinite parts
and a normalizing unitary gauge admit certified effective computation is posed
as an open problem, not solved; the exact recurrences of
\cref{diff:ex:matnoncommuting,diff:eq:matairyrec} are instances, not a
solution. \emph{Amended with member~12:} one such language is now supplied for
one class. For power-Hahn regular-singular systems with finite rational supports
and algebraic coefficients (or the effective coefficient operations of
\cref{diff:rs:sub:procedure}), that procedure computes the Hahn invariants
by a finite calculation, and the phase spectrum is read off from the residual
eigenvalues (\cref{diff:rs:prop:phasespectrum}). Its gauge is not unitary and is
not finitely listed, and for indirectly supplied supports it is not an oracle
algorithm (item~N69); the question as posed, for unitary gauges and general
inputs, stays open.
\item Every statement of \cref{diff:part:matrix} is relative to the one fixed
intrinsic $\der$ of \cref{diff:thm:BM} with $\der\omega=1$. None is a claim
about every surreal derivation, about derivatives of a class function of a
surreal variable, or about any of the coordinate operators of
\cref{diff:tab:derivatives}. Infinite-dimensional matrices, topological
initial-value problems on $\No$, Stokes data over smaller analytic fields and a
global surcomplex exponential are all outside its assertions
(\cref{diff:warn:matlimits}).
\item All dimensions in \cref{diff:part:matrix} are ordinary finite integers,
and every displayed Hahn sum there --- \cref{diff:eq:matnoncommrec},
\cref{diff:eq:matairyamp}, \cref{diff:eq:matairyjet} --- is a
support-certified formal sum, not a limit in the fine topology on $\SC$. This is
item~N32 applied to the matrix part.
\item The metric theorems make no Lyapunov and no uniform-equivalence claim.
\Cref{diff:warn:matlyapunov} exhibits a positive horizontal metric of scale
$e^{-2\omega}$ for printed coefficients whose ordinary solutions grow, and
\cref{diff:thm:matmetric} has no hypothesis relating a horizontal metric to any
fixed norm.
\item \Cref{diff:thm:matspectral} does not remove the connection term
$W^{*}\der W$ of \cref{diff:eq:matspectralproof}; that term is present and is
merely invisible in the ordered additive quotient by $\Acl$. It is not the
assertion that diagonalization commutes with differentiation, and the
classification is not the assertion that every phase grade has a horizontal
vector in $\SC^{n}$ --- by \cref{diff:cor:matinhom} only the zero grade does
(\cref{diff:warn:matconnection}).
\item The proper-class handling of \cref{diff:rem:matclasses} is a size
justification for the proofs, not an algorithm. It does not say that any fixed
Hahn workspace already contains the needed primitives or solutions
(\cref{diff:ex:logmissing}), and it gives no effective presentation, no
birthday or cardinal bound, and no minimal witness-closed subfield.
\item The eighth suite's own scope statement is narrow by its authors'
wording: its checks do not verify the arbitrary-support Hahn arguments, the
published model-theoretic inputs, the existence of a gauge for an arbitrary
surreal matrix, or the general theorems. Its deliberate counterexamples are
asserted \emph{as} counterexamples --- the nonnormal coefficient is required to
have eigenvalues $\pm i$ together with a base-field fundamental matrix, and the
noncommuting exponential defect is required to be nonzero --- so a suite that
passed while those assertions failed would be reporting the opposite of what it
appears to report.
\item The eighth package's verification program \emph{rewrites its own
delivered evidence}. Its closing lines compute the package root as the parent
of the script's directory, create \path{data/} there, and write
\path{data/verification.json}; its README documents the behaviour as
regeneration. \Cref{diff:app:verification} therefore reproduces the record
\emph{as delivered} and does not re-run it, and \cref{diff:app:build} repeats
the instruction to run any source suite on a copy. A byte-identity check made
after such a run compares two equally modified copies and establishes nothing.
This is an artifact-level hazard, like the two in item~N52; it affects no
mathematical statement. The delivered records remain historical; the later
review rerun on isolated copies is recorded separately at the end of
\cref{diff:app:verification}.
\end{enumerate}

\subsection{Scope of the regular-singular theory}
\label{diff:rs:sub:nc}

Member~12 states fifteen explicit limitations of its own, by this merge's
count. All survive here, together with the audit findings recorded against its
package (item~N79).

\begin{enumerate}[label=\textup{N\arabic*.},leftmargin=3.0em,itemsep=0.35em,start=65]
\item \emph{The class.} The residual matrix must have ordinary complex entries,
and the positive perturbation must be power-Hahn with well-ordered real
exponents (\cref{diff:rs:def:class}); only the unknown is unrestricted. The
theorems are not a principle that the standard part of every finite surreal
coefficient controls all solutions: \cref{diff:rs:prop:boundary} is an
infinitesimal coefficient with zero residue and no solution. That counterexample
does not refute a classification based on richer logarithmic or differential
phase data (\cref{diff:rs:rem:boundary}).
\item \emph{No higher rank.} The proofs do not extend unchanged to a general
higher-rank exponent group: the positive real lower bound of $S$ bounds word
lengths, and the shear uses a diagonal Euler derivation with real weights. General
Hahn support principles may survive at higher rank, but an intrinsic
Berarducci--Mantova derivation on arbitrary monomials need not act by
multiplication by their Hahn exponents, and merely replacing $\R$ by an ordered
subgroup of $\No$ in the formulas is not legitimate.
\item \emph{No global objects.} No global surcomplex exponential, no
trigonometric function at an infinite argument, no fine-topological
initial-value theorem and no extension of all ordinary complex local solutions
is constructed; the nonexistence of an imaginary Euler mode is part of the
conclusion. Adjoining $\log t$ does not manufacture a surcomplex fundamental
matrix on a nonzero-frequency summand (\cref{diff:rs:rem:additivePV}).
\item \emph{Forcing.} Existence after adjoining $\log t$ is proved only for
forcing in $\KR^{n}$; no theorem for arbitrary surcomplex forcing is inferred
from the coefficient calculation. ``Fredholm'' and ``index zero'' are algebraic:
no Banach space, closed-range argument, compact operator or topology on $\SC$ is
involved (\cref{diff:rs:rem:fredholm}).
\item \emph{Finite dependence is not an algorithm.} The critical set $S_{\crit}$
is sufficient, not minimal: matrix products can vanish and word contributions
can cancel (\cref{diff:rs:ex:accumulation}). A finite certificate for the
invariants is not a finite listing of the gauge $G_{\natural}$, which may have
infinite support. For a support supplied only indirectly, finding the exact word
sums that hit a resonance need not be effective, finitely many black-box
coefficient queries cannot certify that an unseen resonant coefficient is absent,
and iterating through all exponents below $\max\Lev(A_{0})$ need not terminate
(\cref{diff:rs:sub:procedure}).
\item \emph{The Hahn classification is for $\Gamma=\R$ only.} The complete
invariants of \cref{diff:rs:thm:Kclass} are for $\KR$; over $K_{\Gamma}$ real parts
need not be removable, even rank one retains the real exponent modulo $\Gamma$,
and transferring the classification to $K_{\Q}$ would identify $\derT y=\sqrt2\,y$
with the trivial equation (\cref{diff:rs:rem:notbackwards}).
\item \emph{Which smallness.} The invariance of \cref{diff:rs:cor:invariance} is
under perturbations of positive Hahn order, not under numerically small changes
of the residual matrix (\cref{diff:rs:ex:discontinuity}).
\item \emph{Non-canonical entries.} Only the block similarity invariants of the
nilpotent matrix $N$ are independent of choices; its displayed entries are not
asserted to be canonical across Jordan bases.
\item \emph{Sets, not classes.} The ambient solution collection is identified
with a set, and no dimension theory for a proper-class vector space is used
(\cref{diff:rs:rem:setsized}).
\item \emph{Not a coordinate derivative.} $\derT$ is the intrinsic derivation
rescaled; $t$ names the surreal $\omega^{-1}$,
\cref{diff:rs:eq:eulerhahn} is an identity for the intrinsic derivation, and no
differentiability in the fine topology is assumed; symbolic checks with a formal
variable verify algebraic identities realized by the embedding
(\cref{diff:rs:rem:notcoordinate}).
\item \emph{Antecedents.} Regular-singular normal forms, nilpotent logarithmic
factors and the classification of unipotent formal differential modules are
classical \cite{Levelt,KW}; generalized-series operator methods are van der
Hoeven's \cite{vdH}; and the scalar nonoscillation phenomenon already exists in
this report (\cref{diff:rs:lem:eulermode}). None is relabelled as new. Member~12
proposes as its contribution only the finite-resonance support theorem, the
constructive normal form with support and finite-monomial certificates, and the
exact comparison of Hahn and full-surcomplex solutions and morphisms; general
theorems about generalized-series differential modules may imply parts of the
construction in other terminology.
\item \emph{Novelty and verification.} No independent refereeing and no Lean
verification is claimed. No named published open conjecture is claimed resolved.
Priority is provisional: the searches identified the direct antecedents
\cite{BM,ADH,Levelt,KW,vdH} and this report, and did not locate in them the
specific finite-resonance statement together with the full ambient spectral
selection; that is a bounded assessment, not evidence that no equivalent result
exists, and a specialist should check priority before any publication claim.
\item \emph{The repository comparison.} Member~12 inspected the repository on
22 September 2026 at \CommitC: the root README, the documentation catalogue and
this report's README, together with neighbouring analytic-geometry coverage. The
comparison is to the material returned in that review and is not an assertion of
permanent absence from later revisions; the whole repository was not searched
and its other drafts were not audited; and its README distinguishes the pinned
revision from every README read later in the review. The repository was read,
not modified.
\item \emph{The checks.} The program handles finite positive rational
supports, a specified rational cutoff and an explicitly supplied
generalized-eigenvalue block decomposition. It is not a parser or decision
procedure for arbitrary surreal numbers or indirectly described infinite
supports. Its eight checks detect sign errors, missing matrix-order terms and
incorrect nilpotent inverse bounds; they do not replace the support lemma, the
proof excluding arbitrary ambient solutions, or an independent review, and the
infinite-support theorems are proved in the text, not by the tests.
\item \emph{Audit findings on package~12} (artifact level; no mathematical
statement is affected). The program writes its \texttt{--output} file, by
default \path{verification.json} in the current working directory, and prints
the same JSON, and the package's \texttt{verify} target overwrites both delivered
records; placed here as \path{code/12-regular-singular-Makefile}, that file still
names the package's original file names, not the placed ones. Run it only on a
copy. The field \texttt{ambient\_dimension: 1} in the mixed-frequency record is a
label written by the script, not a computed quantity. The tail-invariance check
compares resonant coefficients only through weight $4$. The two delivered records
carry the same JSON content, and there is no checksum manifest. The record states
SymPy 1.14.0; the package README says the run used Python~3 without a minor
version, and the program's docstring requires Python~3.10 or later.
\end{enumerate}

\subsection{Scope of the autonomous theory}
\label{diff:aut:sub:nc}

Member~11 states fifteen explicit limitations of its own, by this merge's count.
All survive here, together with the audit findings recorded against its package
(item~N94).

\begin{enumerate}[label=\textup{N\arabic*.},leftmargin=3.0em,itemsep=0.35em,start=80]
\item \emph{Derivations.} Every theorem of \cref{diff:aut:part} holds for every
normalized derivation, but no equality of two normalized derivations on $\SC$ is
asserted: they agree on the solutions of every nonzero $F(y,y')=0$ and on the
fields $K_{\xi}$, and nowhere else is claimed. No claim is made that $\No$ has a
unique normalized derivation, or that every surreal is first-order autonomous
differential-algebraic over $\C$, and no normalized derivation is classified
(\cref{diff:aut:conv:scope,diff:aut:rem:notsay}).
\item \emph{Coverage.} The classification is complete for first-order equations
on curves only. Coefficients in $\SC\setminus\C$, higher differential order and
simultaneous systems are outside it and may probe scales on which normalized
derivations differ; which higher-order sectors distinguish normalized
derivations is left open, and no uniform finite-scale bound for arbitrary
higher-order equations is implied.
\item \emph{No higher-dimensional classification.} \Cref{diff:aut:prop:contraction}
survives in every dimension, but higher-dimensional autonomous systems are not
classified: near a singular point several eigendirections interact, and integer
relations among eigenvalues and noncommuting normal-form terms defeat the
one-coordinate time primitive (\cref{diff:aut:rem:boundaries}).
\item \emph{Differential-algebraic only.} Nothing is asserted about derivatives
of arbitrary fine-topological class functions, global complex exponentiation or
topological convergence of surreal series; the formal implicit uniqueness is not
a convergence statement about Newton iteration; and no ordinary path on a
Riemann surface is treated as a fine-continuous path in $\SC$.
\item \emph{Monomial generation is not transcendence degree.} The finite
monomial generation of $K_{\xi}$ is not a finite transcendence-degree bound:
\cref{diff:aut:cor:independent} places $2^{\aleph_{0}}$ algebraically independent
solutions inside a two-generator monomial group.
\item \emph{The strip logarithm is pointwise.} $\Log$ is used pointwise or on a
specified leading-term class; it is not a global additive logarithm and defines
no global surcomplex exponential (\cref{diff:aut:lem:log}).
\item \emph{Classical and imported inputs.} The local formal normal forms are
classical \cite[Lemma~6.1]{NPT}, Rosenlicht's algebraic-independence theorem is
imported \cite[Lemma~7.2, Proposition~7.3]{NPT}, and the curve-and-derivation
formalism is that of \cite{NPT}, whose numbered statements are cited from the
consulted arXiv version~1 and are not assumed to match the journal numbering. The
construction of surreal derivations, formal implicit equations, invariant
differentials and characteristic-zero formal-group logarithms are inputs, not
contributions; \cite{Milne} supplies the standard setting only, and the
formal-logarithm argument is proved in full (\cref{diff:aut:lem:formallog}).
\item \emph{Non-canonical parametrizations.} Parametrizations depend on the local
parameter and on the choice of $\log\beta$, and the leading coefficient $\beta$ is
coordinate-dependent; only the number of branches, the family of points and the
existence criterion are canonical (\cref{diff:aut:rem:choices}).
\item \emph{Constant commutative groups only.} The abelian results hold for
constant abelian varieties over $\C$ only. Nonconstant abelian varieties, which can
acquire a connection term, and noncommutative algebraic groups, whose invariant
forms need not be closed, are excluded. No complex analytic uniformization, no
period lattice and no exponential on all surcomplex Lie-algebra elements is
used; $\exp_{\mathcal F}$ is evaluated only on $\mm_{\C}^{d}$.
\item \emph{Surjectivity and standard part.} Surjectivity of the derivation is
used to define $\Obs^{\C}_{\Dn}$ on every input and in
\cref{diff:aut:cor:abeliancriterion}. The explicit primitives in
\cref{diff:aut:cor:speeds} establish that threshold for every
normalized derivation without surjectivity. A standard part is
taken only after finiteness of the primitive has been established.
\item \emph{Effectivity.} The finite geometric decision criterion is an effective
procedure only with effective divisor and local-expansion data, as well
as exact coefficient operations, equality and sign tests. Algebraic
complex presentations with the stated geometric algorithms suffice.
It is not an
unconditional computability claim about arbitrary unencoded complex constants,
and producing an entire formal solution as a finite string requires its
defining implicit equation and coefficient recursion
(\cref{diff:aut:rem:decision}).
\item \emph{Classes.} $\mathcal X(\SC)$, $\mathbf A(\SC)$ and the exact sequences
over them are class-level statements read NBG-style and pointwise; no collection
of all surreal supports is used as a set (\cref{diff:aut:rem:formalization}).
\item \emph{No trigonometry.} The Riccati example asserts nothing about
separately defined symbolic trigonometric functions
(\cref{diff:aut:ex:riccati}).
\item \emph{Novelty, verification and the repository comparison.} No peer
review or Lean verification of the autonomous theorems is claimed.
The source supplied no Lean file. Current library results do cover the
algebraic finite-variable evaluation and composition clauses; the
intrinsic chain rule and geometric reduction bridge remain pending,
as distinguished in \cref{diff:aut:rem:formalization}. No named published open problem is claimed solved and no
exhaustive priority certificate exists: the results are candidate new theorems
with proofs, and a specialist may identify earlier equivalents. The repository
comparison was scoped to the root README, the documentation inventory and this
report's README at \CommitC, with scoped keyword searches; irrelevant keyword
results were not treated as evidence of absence, not every repository manuscript
or publication was read, and the repository was read, not modified. The source's
review log is kept in this directory as
\path{11-autonomous-dynamics-sources.md}.
\item \emph{The checks, and audit findings on package~11.} The program's ten
exact finite checks --- the logistic solution, the two counterexamples of
\cref{diff:aut:ex:hypotheses}, the time primitive and a total-degree truncation
of \cref{diff:aut:ex:logcorrection} --- catch sign, normalization and coefficient
errors in concrete instances; they do not establish the support lemma, the
algebraic-geometric reduction, the completeness of the solution list, or
novelty. At the artifact level: the program itself writes only to standard
output and accepts \texttt{--degree} from $2$ to $7$ (default $4$), but the
package's build script, placed as \path{code/11-autonomous-dynamics-build.sh},
and its README's rebuild instructions redirect that output onto the delivered
\path{verification.txt}; run them only on a copy. There is no checksum manifest.
\end{enumerate}

\subsection{Scope of the transfinite critical potentials}
\label{diff:cp:sub:nc}

Member~13 states twenty-one explicit limitations of its own, by this merge's
count. All survive here, together with the audit findings recorded against its
package (item~N106); its two open directions, which it does not claim to have
solved, are among the continuations of \cref{diff:sub:nc:open}.

\begin{enumerate}[label=\textup{N\arabic*.},leftmargin=3.0em,itemsep=0.35em,start=95]
\item \emph{Scalar unknowns.} The unknown of \cref{diff:cp:part} is a scalar of a
specified differential field, not a function on an ordinary interval
(\cref{diff:conv:meaning}); ``no solution in $\SC$'' is not a claim that an
ordinary complex initial-value problem has no analytic solution, and nothing
is asserted about coordinate differentiation in $\Hol(U)((t^{\Gamma}))$ or about
fine-differentiable class functions (\cref{diff:cp:rem:derivation}).
\item \emph{One derivation, and an imported tower.} Every statement is relative
to the fixed $\der$ (item~N3). The ordinal log-atomic formulas, limit cases
included, are imported from \cite[Section~2]{ADH}; they are not obtained by
differentiating finite iterates, and the normalization $\der\omega=1$ alone
does not determine them. No derivation-independent version is claimed: for
another derivation the finite identities may still hold, but the limit-stage
derivatives need not be the Berarducci--Mantova ones, and nothing is asserted
for the class of normalized derivations of \cref{diff:aut:part}
(\cref{diff:cp:prop:tower,diff:cp:rem:derivation}).
\item \emph{No transfinite iteration or composition.} The $\ell_{\alpha}$ are
particular numbers; no globally defined transfinite iteration or composition of
$\log$, no substitution $v(s)$ and no global composition is used
(\cref{diff:cp:rem:notation,diff:cp:lem:gauge}, \cite[Theorem~8.4]{BMGerms}).
The Conway monomial map is not $x\mapsto\exp(x\log\omega)$, and no
variable-exponent power rule is used (\cref{diff:warn:monomialmap}).
\item \emph{Sums and classes.} Every sum is a Hahn sum of a set-indexed summable
family; none is a limit of partial sums in the fine topology, $Q_{\lambda}$ is
the sum of its increments and not of the $Q_{\beta}$, the imported product
identity is not replaced by convergence of finite products, and no proper-class
family is summed. There is no all-ordinal potential
(\cref{diff:cp:prop:noend}), and that is a boundary of the construction, not a
convergence problem. Solution collections are sets; no dimension theory of a
proper-class vector space is used.
\item \emph{This family only.} No classification of arbitrary potentials, and
no comparison of leading coefficients, is asserted; no comparison theorem for
surreal Riccati equations is used; that every
$\varepsilon\prec(\ell_{\alpha}^{\dagger})^{2}$ preserves subcritical or
supercritical behaviour is not proved and would need further differential
asymptotic estimates, and at $c=\frac14$ a smaller perturbation can matter
(\cref{diff:cp:rem:notarbitrary}).
\item \emph{What invisibility means.} \Cref{diff:cp:thm:invisible} concerns the
named hierarchy $(\ell_{\beta}^{\dagger})^{2}$, not every conceivable
observation or algorithm. It is a statement about a quotient by an additive
class of errors, not a discontinuity in the fine topology, and that class is
not called an ideal of the field (\cref{diff:cp:warn:quotient}).
\item \emph{Field-theoretic scope.} $\Span_{\R}$ means finite combinations, not
Hahn closure inside the exponent group, and at $\lambda=\omega$ the group
$\Delta_{0}$ is not a Laurent-series field (\cref{diff:cp:warn:supports,diff:cp:ex:omega}).
No closure property of $H_{\Delta_{2}}$, logarithms included, is asserted;
``minimal'' refers to the solution-generated field only, and no minimality of
the intermediate Hahn fields among arbitrary fields is claimed
(\cref{diff:cp:rem:workspaces,diff:cp:rem:pvscope}).
\item \emph{Inputs and antecedents.} The finite critical hierarchy and
reduction of order are classical \cite{KT}, whose potential has the opposite
sign; the derivation \cite{BM}, the log-atomic formulas \cite{ADH}, Hahn real
closedness and the positive-support lemma are imported; the obstruction for
nonreal powers is this report's criterion, reproved there in a special case;
the gauge and Schwarzian mechanisms are standard, and the Schwarzian form is
not an independent result. The general theory of such solvability, including
the predicate $\Omega$ of \cite[p.~17]{ADH}, is not new; \cite{ADHBook} is
neither replaced nor claimed; \cite{Mantova2026} is neither used nor
superseded. The main proofs use no model completeness, no differential closure
and no finite-system splitting theorem, and surjectivity only in the real
Riccati remark (\cref{diff:cp:rem:inputs,diff:cp:rem:riccati}).
\item \emph{Novelty and verification.} Member~13 proposes as new only the
transfinite classification, the invisibility theorem, and the limit-stage
filtration with its Borel group; no isolated elementary identity is promoted.
Priority is provisional: the literature search was targeted and did not
certify exhaustive coverage, an expert may recognize a conclusion as a known
consequence of a broader theorem, and the arguments do not depend on the
novelty assessment. No named published conjecture is claimed resolved. No
independent refereeing, no proof-assistant verification and no Lean
implementation is claimed. This report itself makes no novelty claim
(item~N9).
\item \emph{The repository comparison.} Member~13 read the repository at
\CommitD: its tree, root README, the catalogue \path{docs/README.md}, the listing
of \path{docs/new/}, the introductions of the entire-functions and
Gamma-functions reports, and source ranges of this article. Connector searches
for ``Kneser'' and ``iterated-logarithm'' returned no hits, including none for
wording visible in an inspected source, and were not treated as evidence of
absence. The comparison is bounded and is not a proof of absence from the
repository, its history or later commits; the repository was read, not
modified. \emph{Recorded by this merge:} the article it read already
contained the first-order counterpart \cref{diff:cor:logscales}, titled
``Iterated-logarithm thresholds at every finite depth'', which the member does
not cite; \cref{diff:cp:lem:scales,diff:cp:cor:nonreal} now connect the two.
\item \emph{The checks.} The $70$ exact symbolic checks --- the critical and
Riccati identities at lengths $0$ to $12$ ($39$), Euler powers, bases and
Wronskians ($8$), the gauge identity ($4$), Schwarzian identities ($4$) and the
fundamental matrix and Borel action ($15$) --- test finite identities, with the
Euler checks modelling the scale by an ordinary positive variable. They do not
construct the surreal numbers or the derivation, do not prove the imported
derivative formula, do not certify transfinite Hahn summability or
exponent-field membership, and do not establish novelty; a count is not a
measure of proof strength, and no finite list of checks certifies a theorem for
every ordinal.
\item \emph{Audit findings on package~13} (artifact level; no mathematical
statement is affected). The program writes to standard output and, with
\texttt{--output}, to the named file; the package README's
\texttt{python verify.py --output verification.txt} would overwrite a record
of that name, so run it only on a copy. The placed manifest
\path{data/13-critical-potentials-manifest.json} lists five files with SHA-256
digests; all five match the delivered package, but only the program and its
record are shipped here, as \path{code/13-critical-potentials-verify.py} and
\path{data/13-critical-potentials-verification.txt}, and the package's
\path{article.tex}, $21$-page \path{article.pdf} and \path{README.md} are not.
The record states Python 3.13.5 and SymPy 1.14.0.
\end{enumerate}

\subsection{Open continuations, hedged}
\label{diff:sub:nc:open}

The following are directions for further development. They are \emph{not}
assertions that the corresponding questions are unresolved in the literature.

Two items that earlier editions of this list posed have since been answered
inside the report, by \cref{diff:part:matrix}, and are recorded here as
answered rather than deleted; two more are answered in part by
\cref{diff:rs:part}, and a third is adjacent to \cref{diff:aut:part}; all three are
re-scoped below.
\begin{itemize}[leftmargin=1.6em]
\item \emph{Which variable-coefficient matrix equations keep their fundamental
matrices in $\No[i]$.} \Cref{diff:cor:matinhom} answers this for every finite
system: a fundamental matrix exists in $\GL_{n}(\SC)$ exactly when every member
of $\Ph(A)$ is zero, and $\dim_{\C}\Ker(\der-A)=m_{0}(A)$ in general. For
Hermitian coefficients \cref{diff:cor:matkernel} makes the test explicit in
terms of ordinary algebraic eigenvalues: the dimension is the number of
eigenvalues lying in $\Acl$. The distinction the item asked for --- genuine
phase obstruction against workspace obstruction --- is exactly the distinction
between $\Ph(A)\neq\{0,\dots,0\}$ and the workspace phenomena of
\cref{diff:part:workspaces}. For $A=iH$ with $H$ Hermitian,
\cref{diff:cor:matperturbation} shows that the entire phase multiset is
unchanged by Hermitian entrywise perturbations of $H$ in $\Acl_{\C}$.

What the item also asked --- which nontriangular differential modules admit a
filtration whose rank-one factors have \emph{computable} obstructions ---
remains open, and for the reason the item gave: that needs genuine
module-theoretic information, not the eigenvalues of an instantaneous
coefficient matrix. \Cref{diff:ex:matnonnormal} confirms the reason rather than
removing it. The remaining question here is effective computation on a
specified input language; finite-rank classification over $\SC$ itself is
already supplied by \cref{diff:thm:matnormal}.
\item \emph{Connecting the phase invariant with the strongest available
linear-solvability theorems for this $H$-field.} This is done in
\cref{diff:thm:matsurjectivity,diff:cor:matinhom,diff:cor:matkernel}, and it
separates existence from dimension exactly as the item asked: existence for
\emph{every} forcing holds unconditionally, while the dimension of the
homogeneous solution space is the multiplicity $m_{0}(A)$ of the phase zero.
The separation is not free --- existence comes from the imported
\cref{diff:thm:matsurjectivity} and not from the rank-one criterion
(\cref{diff:warn:matinhom}, item~N53) --- and no effective version of either
half is claimed (item~N57).
\end{itemize}

Three of the remaining continuations are touched by \cref{diff:rs:part,diff:aut:part}: two
are answered \emph{in part}, and one only has an adjacent result. They are
re-scoped here rather than deleted.
\begin{itemize}[leftmargin=1.6em]
\item \emph{Find explicit formulas for the already defined phase spectrum of further
classes of non-Hermitian coefficients,} the obstruction being that a nonnormal
similarity gauge need not be entrywise finite (\cref{diff:warn:matnonnormal}).
Answered for one class: for power-Hahn regular-singular systems
$\Ph(-tA)=\{(\Imm\lambda_{j})\log t\}$, with an explicit gauge and no import
(\cref{diff:rs:prop:phasespectrum}, member~12). A general explicit formula outside the currently computed classes remains
open; \cref{diff:ex:matnonnormal} rules out the instantaneous-eigenvalue formula.
\item \emph{Identify restricted, support-certified input languages on which the
phase map, the purely infinite parts of eigenvalue primitives and a normalizing
unitary gauge become effectively computable.} Answered in part for one class:
finite rational supports with algebraic coefficients for power-Hahn
regular-singular systems
(\cref{diff:rs:sub:procedure,diff:rs:rem:effective}). The gauge there is not
unitary and not finitely listed, and the continuation stands for unitary gauges
and for every other class.
\item \emph{Classify logarithmic derivatives in multi-scale $\der$-stable Hahn
workspaces where support shifts lack the uniform rational form of
\cref{diff:eq:hahnderiv}.} Not answered. Member~11 supplies explicit multi-scale
$\der$-stable Hahn fields $K_{\xi}$, with power, logarithmic and exponential
scales (\cref{diff:aut:rem:monomial}), but does not classify logarithmic
derivatives on them; member~13 supplies the higher-rank logarithmic fields
$H_{\Delta_{j}}$ of \cref{diff:cp:sec:fields}, with an explicit
logarithmic-derivative formula for their monomials
(\cref{diff:cp:prop:fieldclosure}), and does not classify them either.
\end{itemize}

Members~11 and~12 add continuations of their own. From member~12: a
\emph{layered} analogue of the regular-singular theory for coefficient fields
containing $\log\omega$, $\log\log\omega$ and their real powers, where
\cref{diff:rs:prop:boundary} shows residual data at the first power scale to be
insufficient and a finite sequence of scale-dependent residual matrices might
decide existence for a restricted logarithmic class --- any answer must separate
derivative effects from monomial valuation and justify every support interchange
anew; which representations of semialgebraically or recursively specified
well-ordered supports make exact resonance-word extraction decidable; and a
modular formalization whose Hahn part needs only finite-dimensional linear
algebra, strong summation and the positive-support lemma, and whose ambient part
needs only the derivation clauses, complexification and the finite-element
estimate. From member~11: autonomous systems in higher dimension; nonconstant
abelian varieties and noncommutative algebraic groups; and which higher-order
differential-algebraic sectors distinguish normalized derivations. One question
is recorded by this merge rather than by either source: every hypothesis
member~12 lists (\cref{diff:rs:rem:imports}) is a clause of the definition of a
normalized derivation, and on $\KR$ the operator $-\omega\Dn$ acts by
$t^{a}\mapsto at^{a}$ for every normalized $\Dn$, so \cref{diff:rs:part} appears to
hold for every normalized derivation; that has not been audited line by line, and
no statement of \cref{diff:rs:part} is upgraded.

Member~13 adds two directions and does not claim either solved: characterize
the ordered exponent groups $\Gamma$ for which $Q_{\lambda}$ lies in
$\R((t^{\Gamma}))$ and $\der$ restricts, and classify the solution dimension there
by the two basis exponents; and study controlled nonconstant perturbations of
$c$ after the exact reduction to $\der_{s}^{2}+cs^{-2}$. Its part is adjacent to
member~12's layered continuation above without answering it: within its one
family the coefficient at the last scale decides existence and no set of
earlier scales does (\cref{diff:cp:thm:invisible,diff:cp:rem:separator}), but no
restricted logarithmic class of regular-singular type is classified. It is
likewise adjacent to the extension of the workspace algorithms to logarithmic
and higher-rank parents below: the fields $H_{\Delta_{j}}$ are explicit such
parents, with no algorithm.

The remaining continuations stand.
Build explicitly typed analytic and coefficient interfaces that prove exchange
rules with contour functionals under their correct common-domain and
summability hypotheses, always fixing whether centres and contours move.
Extend the workspace algorithms from real-exponent Hahn fields to logarithmic
and higher-rank parents with explicit support-change certificates, preserving
the distinction between a representation of a primitive and a decision about
its infinite part. Separate negative-exponent coherent systems into finite
nilpotent cases, cases reducible by a gauge transformation, and genuinely
non-coherent cases, each proposed reduction carrying its own support proof.
Develop monodromy for coherent coordinate equations on non-simply-connected
domains in the nonlinear case.

\section{Conclusion}
\label{diff:sec:conclusion}

The internal derivative does not turn a surcomplex number into an ordinary
complex-valued function at infinity. Real exponential integration is
extremely rich; unit-circle phases are severely constrained, because only a
constant ordinary phase and an infinitesimal residue already live in the base
field. The exact consequence is
\[
 \boxed{\quad
 \{y^{\dagger}:y\in\No[i]^{\times}\}=\No+i\,\der\mm
 =\No+i\,\der\Ofin,
 \qquad \der\mm\subsetneq\mm .\quad}
\]
The purely infinite part of an accumulated phase of the imaginary coefficient
is therefore the complete rank-one obstruction. It supplies a canonical normal
form, distinguishes algebraic closure from differential solution closure,
explains the missing oscillator, determines the maximal domain of any
derivation-compatible exponential, removes every nonreal characteristic root
from a constant-coefficient solution space, and fixes how many formal phase
generators a diagonal Picard--Vessiot extension needs. At the workspace level
two sharper identities are available,
\[
 \der(\Hring)=\{f\in\Hring:[t]f=0\},
 \qquad
 \{y^{\dagger}:y\in\Hring^{\times}\}
 =\Q\,t+\{g:g=0\text{ or }v(g)>1\},
\]
giving exact auditable tests rather than heuristic simplification rules, and
showing why derivative closure, primitive closure and exponential closure must
be separate capabilities.

The same obstruction survives the passage to finite matrices, and there it
becomes a multiset. Every finite system over $\SC$ has a differential phase
spectrum $\Ph(A)\subset\Pcl$, complete for gauge equivalence, of which the
trace test \cref{diff:eq:tracenecessary} sees only the sum; and when the
coefficient is $iH$ with $H$ Hermitian the spectrum is computed by integrating
the ordinary algebraic eigenvalues of $H$ and keeping the purely infinite parts
of their primitives. The connection term left behind by pointwise
diagonalization is neither ignored nor approximated: by
\cref{diff:lem:matboundedgauge} a unitary gauge is entrywise finite, so that
term lands in $\Mat_{n}(\Acl_{\C})$, which is exactly the kernel the invariant
annihilates. That is why the perturbation theorem needs no spectral gap. What
the argument needs instead is the strictness $\Acl\subsetneq\mm$ that
\cref{diff:warn:collapse} insists on, and an additional imported input.
The eigenvalue computation needs Hermitian symmetry, which
\cref{diff:ex:matnonnormal} shows cannot be dropped from that formula.

Two further strata show the same obstruction from other sides. For
regular-singular systems $\der y=-t\,Ay$ with an ordinary residual matrix and a
positive real-power perturbation, finitely many resonance-reaching coefficients
determine the Hahn invariants, real residual parts are sheared away by real
monomials, nilpotent Levelt data are removed by one actual logarithm, and what
survives in $\SC$ is exactly the imaginary parts of the residual eigenvalues:
$\Ph(-tA)=\{(\Imm\lambda_{j})\log t\}$, a non-Hermitian sector computed without
the matrix part's imported input. And for first-order autonomous equations with
constant coefficients --- constant vector fields on curves --- every surcomplex
trajectory is centred at a zero of the field, a simple zero contributes exactly
when its eigenvalue is negative real, a multiple zero contributes finitely many
power--logarithm families, and the whole answer is the same for every normalized
surreal derivation. Constant abelian varieties obey the same rule with their
whole Lie coordinate: $\dlog$ lands exactly in $\Lie(\mathbf A)(\der\mm_{\C})$.

The same obstruction also runs down the log-atomic tower. For the second-order
potentials $Q_{\alpha}+c(\ell_{\alpha}^{\dagger})^{2}$ one gauge identity
reduces the equation to the Euler operator at the scale $\ell_{\alpha}$, and
the threshold $c=\frac14$ is exact at every ordinal depth: the missing
oscillator is the same Euler problem at the scale $e^{\omega}$, and in every
supercritical case the obstruction is an infinite accumulated phase, now
$\sqrt{c-\frac14}\,\ell_{\alpha+1}$. No set of earlier logarithmic scales
decides solvability, and at every limit stage the solutions appear one at a
time in differential Hahn fields that are already real closed and closed under
logarithms --- first after an exponential, then after a primitive --- with the
Borel subgroup of $\SL_{2}(\C)$ as the Galois group of the two steps together.

A distinct theory of coordinate initial-value problems remains available and
useful. On a common ordinary complex domain, positive-support Hahn
perturbations of a holomorphic problem are solved uniquely by support
recursion through ordinary linear equations, with explicit support
certificates, a Liouville determinant identity, and monodromy valued in
$\GL_{d}(K_{\Gamma})$. The two theories coexist without conflict once the
derivative, the parent field, the coefficient domain and the summability
certificate are part of every statement.

The resulting picture is coherent rather than paradoxical: $\SC$ is
algebraically closed; $\der$ is surjective; real exponential integrals exist;
finite angular motion is exact; and a sustained nonzero ordinary angular
velocity is nevertheless absent. These are different closure properties, and
the obstruction sequence \cref{diff:eq:exact} records their precise boundary.
The boundary worth remembering is this:
\begin{quote}
Every surcomplex direction has a finite angle, but a differential evolution
prescribing a purely infinite \emph{accumulated} phase need not have a
surcomplex value. A phase convention may assign a value; it cannot force that
value to satisfy the incompatible number-derivation law.
\end{quote}
And the distinctions are not a notation policy. Their compatibility and
incompatibility are theorems: an intrinsic derivative of a scalar, an external
derivative of a function, and the substitution that might connect them are
three pieces of structure, not one.
\appendix
\part{Appendices}
\label{diff:part:appendices}

\section{The positive-support calculus, proved}
\label{diff:app:support}

This appendix supplies the combinatorial facts used throughout. It works in a
set-sized ordered abelian group and uses no completeness assumption on the
group. The finite-word principle is the well-ordered-alphabet case of Higman's
lemma \cite{Higman}; the support consequence is the classical Hahn--Neumann
fact \cite{Neumann}. Both are cited as published sources, and a proof is
included as well, so that nothing in this report depends on an uninspected
passage. The published citations are \emph{not} replaced by this appendix, and
this appendix is not replaced by them.

\subsection{Two elementary facts about well-ordered supports}
\label{diff:sub:app:twofacts}

First, every sequence in a well-ordered set has an infinite nondecreasing
subsequence: choose an occurrence of the least value in the whole tail, then an
occurrence of the least value in the tail strictly after that index, and
continue; the selected indices increase and the selected values do not
decrease.

\begin{lemma}[Sum of two supports]
\label{diff:lem:twosupports}
If $A,B$ are well-ordered subsets of an ordered abelian group, then $A+B$ is
well ordered, and for each $\gamma$ only finitely many pairs
$(a,b)\in A\times B$ satisfy $a+b=\gamma$. A finite union of well-ordered
subsets is well ordered.
\end{lemma}
\begin{proof}
If $a_{n}+b_{n}$ were strictly decreasing, select a nondecreasing subsequence
of the $a_{n}$; the corresponding $b_{n}$ would then be strictly decreasing,
contradicting well ordering of $B$. In a linear order, the absence of an
infinite strictly decreasing sequence implies well ordering in the classical
set-theoretic setting used here: a nonempty subset with no least element would
produce such a sequence by recursive choice.

If infinitely many distinct pairs had a fixed sum, select a sequence of
distinct pairs. Their $a$ entries are distinct, so a nondecreasing subsequence
of them is strictly increasing, and the corresponding $b$ entries strictly
decrease --- impossible. Finally, an infinite decreasing sequence in a finite
union would have an infinite subsequence inside one member of the union.
\end{proof}

\subsection{Finite words over a well-ordered alphabet}
\label{diff:sub:app:higman}

A word $a_{1}\cdots a_{m}$ \emph{embeds} in a word $b_{1}\cdots b_{n}$ if
there are indices $1\leq j_{1}<\cdots<j_{m}\leq n$ with
$a_{k}\leq b_{j_{k}}$ for every $k$. The empty word embeds in every word.

\begin{lemma}[Finite-word lemma]
\label{diff:lem:higman}
Every infinite sequence of finite words over a well-ordered alphabet contains
two words, in their original order, such that the earlier embeds in the later.
\end{lemma}
\begin{proof}
Call a sequence \emph{bad} if it contains no such pair, and suppose a bad
sequence exists. Choose one minimally in the following sense: after a fixed
prefix $w_{0},\dots,w_{n-1}$, choose $w_{n}$ of smallest possible length among
words that permit continuation to an infinite bad sequence. Natural-number
lengths make each choice possible, and no $w_{n}$ is empty, since the empty
word embeds in everything.

Write $w_{n}=v_{n}a_{n}$ with $a_{n}$ the last letter, and select indices
$n_{0}<n_{1}<\cdots$ for which the letters $a_{n_{j}}$ are nondecreasing.
Consider the sequence
\[
 w_{0},\dots,w_{n_{0}-1},\,v_{n_{0}},v_{n_{1}},v_{n_{2}},\dots
\]
It cannot be bad, because $v_{n_{0}}$ is shorter than the minimal choice
$w_{n_{0}}$ after the same prefix. So it has a good pair. That pair cannot lie
wholly inside the prefix, since the original sequence was bad. If it runs from
a prefix word to some $v_{n_{j}}$, then the prefix word also embeds in
$w_{n_{j}}$, again contradicting badness. And if it runs from $v_{n_{i}}$ to
$v_{n_{j}}$ with $i<j$, append the nondecreasing last letters to the
embedding: $w_{n_{i}}=v_{n_{i}}a_{n_{i}}$ embeds in
$v_{n_{j}}a_{n_{j}}=w_{n_{j}}$, once more a contradiction. Hence no bad
sequence exists.
\end{proof}

\subsection{From words to Hahn sums}
\label{diff:sub:app:towords}

Let $S\subseteq\Gamma_{>0}$ be well ordered. If a word in $S$ embeds in
another, the sum of its letters is at most the sum of the other's; and
equality forces the words to be identical, because any unmatched letter is
strictly positive and any strict comparison of matched letters contributes a
strictly positive difference, so a finite sum of such contributions cannot
vanish.

An infinite strictly decreasing sequence of word sums would therefore be a bad
word sequence, contradicting \cref{diff:lem:higman}; hence $S^{*}$ is well
ordered. If infinitely many distinct words had the same sum, an embedded pair
among them --- which \cref{diff:lem:higman} supplies --- would force two of
them to be identical. So each exponent has only finitely many word
representations. This proves \cref{diff:lem:neumann}.

\begin{remark}[Ordered words, and what this appendix does \emph{not} prove]
\label{diff:rem:orderedwords}
The count is over \emph{ordered} words, not over unordered multisets. That is
what makes the lemma suffice for noncommuting matrix products: each word
retains the order in which its coefficient matrices were multiplied, and
finite matrix dimensions introduce only finitely many intermediate indices for
a word of fixed length. \Cref{diff:warn:twoscale} is the concrete witness that
the ordering matters. Combined with
\cref{diff:lem:twosupports}, this justifies multiplying an arbitrary Hahn
series by a positive-support power series, the Taylor-evaluation families of
\cref{diff:eq:halo,diff:eq:explog}, and every coefficient recursion in this
report.

It does \emph{not} establish cofinal growth of $n\sigma$ in a higher-rank
value group, and no assertion of that kind is used anywhere
(\cref{diff:warn:summability,diff:warn:notpicard}).
\end{remark}

\section{Notation, and a concordance with the source manuscripts}
\label{diff:app:notation}

\begin{longtable}{@{}L{0.17\textwidth}L{0.76\textwidth}@{}}
\caption{Notation used in this report.}
\label{diff:tab:notation}\\
\toprule
Symbol & Meaning \\
\midrule\endfirsthead
\toprule Symbol & Meaning \\\midrule\endhead
$\No$, $\SC$ & The real surreal class field and its complexification
$\No[i]$; class-sized carriers, with all supports and summed families sets. \\
$\der$, $x'$ & The fixed simplest Berarducci--Mantova derivation, normalized
by $\der\omega=1$, extended to $\SC$ by $\der i=0$. \\
$\der_{s}$ & The rescaled derivation $\der f/\der s$, with
$\der_{s}s=1$. \\
$\derZ$, $\dertil$ & Formal differentiation of the indeterminate on
$\SC[[Z]]$, and coefficientwise intrinsic differentiation there. \\
$\Dz$ & Coordinate differentiation of holomorphic coefficient functions in an
ordinary complex $z$; every monomial and scalar is held constant. \\
$\dercoef$ & Intrinsic differentiation of the Hahn monomials of a
common-domain family, holding $z$ fixed. \\
$D_{t}$ & Formal differentiation in a Hahn parameter treated as an
indeterminate. \\
$F'_{\mathrm{fine}}$ & The pointwise fine order-field derivative of a class
function; not a field derivation. \\
$\Ofin$, $\mm$ & The finite real surreals and the real infinitesimals, with
$\Ofin=\R\oplus\mm$. \\
$\Ofin_{\C}$, $\mm_{\C}$ & $\Ofin+i\Ofin$ and $\mm+i\mm$. \\
$\Pcl$ & The additive class of purely infinite normal forms, including $0$;
\emph{not} the class of all infinite surreals. \\
$\pp$, $\ct$, $\fin$, $\st$ & The purely infinite projection; the coefficient
of $\omega^{0}$, defined on all of $\No$; $\ct+{}$infinitesimal part; and the
standard part, defined on $\Ofin$ only. $\ct$ and $\st$ agree only on
$\Ofin$. \\
$t^{\gamma}$, $t$ & The Conway monomial $\omega^{-\gamma}$, and
$t=\omega^{-1}$; not a generic exponential power. \\
$\Emm$, $\Lmm$ & The infinitesimal Hahn exponential on $\mm_{\C}$ and
logarithm on $1+\mm_{\C}$. \\
$\cis$ & The finite-angle phase map on $\Ofin$, with kernel $2\pi\Z$. \\
$\theta$ & A \textbf{finite angle}: the argument of $\cis$, always in
$\Ofin$. \\
$\eta$ & An \textbf{infinitesimal phase residue}: the unique element of $\mm$
left after factoring out the constant ordinary unit $c$. \\
$B$ & An \textbf{accumulated phase}: any primitive of the imaginary
coefficient. Not an angle; obstructed exactly when $\pp(B)\neq0$. \\
$z^{\dagger}$ & The logarithmic derivative $\der z/z$ for $z\neq0$; presumes
no logarithm. \\
$\Izero$ & The unique primitive with $\ct\Izero(b)=0$; a normalization, not
an algorithm. $\Izero|_{\Acl}$ is the unique infinitesimal primitive. \\
$\Acl$ & The bounded-primitive ideal $\der\mm=\der\Ofin$; a convex ideal of
$\Ofin$, strictly inside $\mm$. \\
$\Obs$, $\Phi$ & $\Obs(b)=\pp\Izero(b)\in\Pcl$, and
$\Phi(q)=\Obs(\Imm q)$. \\
$\Strip$, $\Estrip$ & The additive strip $\No+i\Ofin$, not a subring, and its
derivation-compatible exponential. \\
$\Tor(\No)$, $\Tor(\C)$ & $\{u\in\SC:u\bar u=1\}$ and the ordinary complex
unit circle. \\
$H_{\Gamma}$, $K_{\Gamma}$ & $\R((t^{\Gamma}))$ and $\C((t^{\Gamma}))$;
$K_{0}=\C((t))$, on which $\der=-t^{2}d/dt$. \\
$\Hring$ & The rational-exponent workspace $\C((t^{\Q}))$, a full Hahn field.
\\
$\Hring_{\Gamma}(U)$ & The fixed-common-domain ring
$\Hol(U)((t^{\Gamma}))$, all coefficients holomorphic on one ordinary
$U$. \\
$\Hol(U)$ & Ordinary holomorphic functions on $U$; \emph{not} $\Ofin$. \\
$\mathscr W$ & The proper-class $\omega$-series field with its restricted
composition theory. \\
$\mathsf T$ & A transcendental differential generator with
$\der\mathsf T=i\mathsf T$ (or $=a\mathsf T$); never a surcomplex number. \\
$\mathsf C$, $\mathsf S$, $\mathsf v$ & The recovered oscillator elements and
the conjugation-fixed Cayley coordinate in the extension. \\
$\Lambda$, $\gm(\C)$ & The integer phase-relation lattice, and the
multiplicative algebraic group. \\
$\tau$, $\derT$ & In \cref{diff:rs:part}: $\tau=\log t=-\log\omega$, and the
Euler derivation $\derT=-\omega\der$, the rescaled $\der_{s}$ at $s=\tau$. \\
$\KR$, $K_{\Gamma}^{\geq0}$, $K_{\Gamma}^{>0}$ & $\C((t^{\R}))$, and the elements
of $K_{\Gamma}$ of nonnegative, respectively positive, valuation; in
\cref{diff:rs:part} $\Gamma\subseteq\R$ need not contain $1$. \\
$\Eig(A_{0})$, $\Lev(A_{0})$ & The ordinary eigenvalues of a residual matrix, and
its Levelt resonances $\lambda-\mu\in\R_{>0}$. \\
$A_{0}^{\circ}$, $A^{\natural}$, $A^{\flat}$ & A Jordan form of the residual
matrix; the normalized Hahn coefficient; the constant form $iD+N$. \\
$G_{\natural}$, $F$, $G_{\tau}$ & The normalizing Hahn gauge, the Hahn gauge to the
constant form, and the gauge in $\GL_{n}(\KR[\tau])$ to $iD$. \\
$\hob_{A}$ & The Hahn obstruction of a regular-singular system; not $\Obs$. \\
$\Dn$ & In \cref{diff:aut:part}: an arbitrary normalized derivation; $\der$ is
one of them. \\
$\mathcal X$, $\xi$, $\sigma$ & A smooth projective curve, a meromorphic vector
field on it, and a local parameter. \\
$\mathbf A$, $\log_{\mathcal F}$, $\exp_{\mathcal F}$ & A constant abelian variety
and its formal logarithm and exponential. \\
$\dlog_{\Dn}$, $\Obs^{\C}_{\Dn}$ & The abelian logarithmic derivative, and the
complex primitive obstruction $\pp(\mathfrak b)$ with $\Dn\mathfrak b=a$; not $\Phi$. \\
$\Log$, $\mathfrak t$ & A strip logarithm, used pointwise, and a time
primitive. \\
$\ell_{\alpha}$, $\varpi_{\alpha}$ & In \cref{diff:cp:part}: the log-atomic number
$\omega^{\omega^{-\alpha}}$ for an ordinal $\alpha$, extending the iterated
logarithms $\ell_{n}$; and the exponent sum $\sum_{\beta<\alpha}\omega^{-\beta}$. \\
$\ell_{<\alpha}$, $u_{\alpha}$ & The monomials $\omega^{\varpi_{\alpha}}$ (the
monomial product of the $\ell_{\beta}$, $\beta<\alpha$) and
$\omega^{\varpi_{\alpha}/2}$. \\
$Q_{\alpha}$, $Q_{\alpha,c}$ & The critical potential
$\frac14\sum_{\beta<\alpha}(\ell_{\beta}^{\dagger})^{2}$ and
$Q_{\alpha}+c(\ell_{\alpha}^{\dagger})^{2}$; $Q_{\alpha,1/4}=Q_{\alpha+1}$. \\
$\lambda$, $\Delta_{j}$ & In \cref{diff:cp:sec:fields,diff:cp:sec:pv}: a nonzero limit
ordinal, never an eigenvalue; and the exponent groups of
\cref{diff:cp:eq:groups}. \\
$\mathrm{Wr}$, $\Sch$ & The Wronskian $y_{1}\der y_{2}-y_{2}\der y_{1}$, and the
Schwarzian of \cref{diff:cp:eq:schwarzian}. \\
$\Bor_{2}(\C)$, $\psi_{a,b}$ & The upper triangular Borel subgroup of
$\SL_{2}(\C)$, and the automorphisms of \cref{diff:cp:eq:pvaction}; unrelated to
Borel summation. \\
\bottomrule
\end{longtable}

\begin{warning}[The concordance, and three symbols that must not be reused]
\label{diff:warn:concordance}
Across the source manuscripts the derivation is written $\der$, $\partial$,
$D$, $\delta$, $\mathrm{der}$, $\partial_{\mathrm{BM}}$ and with a prime; two
of them write $\partial$ for the intrinsic derivation and $D_{z}$ for the
coordinate one, and one writes $D$ for a Picard--Vessiot derivation on a
wholly different field. This report fixes $\der$ for the intrinsic derivation
and keeps \cref{diff:tab:derivatives} as the reference. The letter $D$ never
names a second derivation of $\No$ or $\SC$. As an operator it occurs only where
it is named: as the abbreviation $D:=\der|_{K_{0}}=-t^{2}\,d/dt$ in
\cref{diff:prop:hahnderiv}, which is the fixed derivation restricted to
$\C((t))$ and not another one; as ordinary polynomial differentiation on
$\C[x]$ in \cref{diff:prop:polyforcing}, where nothing surreal is
differentiated; and, with a subscript, as the coordinate operators $D_{z}$ and
$D_{t}$ of \cref{diff:tab:derivatives} and the Jacobian $D_{Y}F_{0}$ of
\cref{diff:part:coordinate}. Otherwise $D$, with or without a subscript $P$ or
$Q$, is a diagonal matrix, in \cref{diff:part:matrix,diff:rs:part}, and the
calligraphic $\mathcal D_{A}$ of \cref{diff:sub:matresonance} is the
first-integral derivation of a polynomial ring. (An earlier edition of this
warning said that $D$ appears only as polynomial differentiation; that was not
accurate. The correction changes no label and no statement.) Member~11 writes
$D$ for an arbitrary normalized derivation, which is $\Dn$ here
(\cref{diff:aut:conv:scope}); member~12 writes $\delta=-\omega\der$, which is the
rescaled derivation $\derT=\der_{s}$ at $s=\tau=\log t$ here
(\cref{diff:rs:rem:names}), because $\delta$ is already a source notation for the
infinitesimal phase residue in \cref{diff:tab:phases}.

The letter $J$ is not used in this report
(\cref{diff:conv:noJ}). In the sources it denotes the normalized intrinsic
primitive (here $\Izero$), the representative class of purely infinite
surreals (here $\Pcl$), the inverse of $\der|_{\mm}$ (here
$\Izero|_{\Acl}$), and --- in the companion \emph{rank-one Berkovich} report,
which remains standalone --- the zero-constant \emph{coordinate} primitive of
a convergent power series, whose ``complete primitive obstruction'' is an
annular Laurent residue and is unrelated to the phase obstruction. Any
cross-link between the two reports must name the operator, not the letter.

The letter $U$ is used in this report for an ordinary complex domain
(\cref{diff:part:coordinate}), and once more, with $V$, as a pair of formal
placeholders in the statement of the exponential and logarithm identities
(\cref{diff:lem:smallchain}), where they stand for nothing at all. One source manuscript flagged its own collision
between $U$ as a transcendental differential generator and $U$ as an ordinary
domain; the generator is $\mathsf T$ here, and neither use ever denotes a
surreal number or a matrix over one. The eighth manuscript writes $U$ for a
unitary matrix; this report writes $W$, and
\cref{diff:conv:matalphabet} fixes the rest of that part's alphabet:
$H$ for a Hermitian matrix --- never the hyphenated adjective in ``$H$-field''
and never the coefficient field $H_{\Gamma}$ --- $G$ for a gauge matrix, the
upper-case upgrade of the scalar gauge $g$ of \cref{diff:eq:gauge};
$\mathcal M$ for a horizontal metric, where the eighth manuscript writes $G$;
$\Rot$ for the rotation generator, where it writes the forbidden $J$; and
$\Pi$ for the additive phase lattice inside $\Pcl$, because $\Lambda$ is
already the integer relation lattice of \cref{diff:eq:lattice}, which is the
object that manuscript calls $L$. Its $\Psi$ is this report's $\Obs$
(\cref{diff:conv:phases}(e)) and is \emph{not} the $\Psi$ of
\cref{diff:cor:noglobalPsi}, which is a placeholder for a putative
derivation-compatible exponential and is the letter's only occurrence here; its
$\mathcal I$ is $\Acl$. In
\cref{diff:part:matrix} the letter $W$ denotes a unitary matrix and $M$ a
finite differential module; \cref{diff:part:coordinate} uses the same two
letters for a positive-support matrix, an ordinary neighbourhood and a
well-ordered exponent set, and the two parts never share a statement. The
unitary group itself is $\Un_{n}(\SC)$, printed upright so that it is not read
as a domain.

Member~12's letters are renamed in \cref{diff:rs:conv:alphabet}: its Jordan form
$J$ and nilpotent parts $J_{\lambda}$ are $A_{0}^{\circ}$ and $N^{\circ}_{\lambda}$
(the letter $J$ stays retired); $\ell=\log t$ is $\tau$ and $\delta$ is $\derT$;
$K$ is $\KR$, and $\mathcal O_{K}$, $\mathfrak m_{K}$ are $\KR^{\geq0}$, $\KR^{>0}$;
its $P$, $H$, $B$ are $G_{0}$, $G_{\natural}$, $A^{\natural}$; its $M$ and
$\mathcal R$ are $S^{*}$ and $\Lev(A_{0})$; its $\Pi_{\gamma}$, $L_{\gamma}$,
$K_{\lambda\mu}$ are $\pr_{\gamma}$, $\mathcal H_{\gamma}$, $\mathcal N_{\lambda\mu}$;
its $T$, $R_{0}$, $C$, $Q$ are $t^{D_{\mathrm{re}}}$, $D_{\mathrm{re}}$, $A^{\flat}$,
$G_{\tau}$; its set of labels $B$ is $\Imm\Eig(A_{0})$; its obstruction $\operatorname{ob}_{A}$ is
the Hahn obstruction $\hob_{A}$; its operator $\mathcal D_{A}$ is written
$\derT-A$; a matrix it calls $U$ is $X$; a real weight it calls $\eta$ is
$\gamma'$; and its ``phase labels'' are residual frequencies
(\cref{diff:rs:warn:frequencies}). It keeps $D$ for a real diagonal matrix, as
\cref{diff:part:matrix} does, and $V,V_{\lambda},V_{b},V_{0}$ for $\C^{n}$ and its
generalized-eigenspace and frequency subspaces.

Member~11's letters are renamed in \cref{diff:aut:conv:alphabet}: its derivation
$D$ is $\Dn$ and its field $K=\No(i)$ is $\SC$; its curve parameter $t$ is
$\sigma$, so that $t$ remains the number $\omega^{-1}$
(\cref{diff:warn:separators}(1)); its curve $X$, vector field $V$, points $P$ and
reductions $\bar P$ are $\mathcal X$, $\xi$, $\mathbf y$ and $\st(\mathbf y)$; its
$1$-forms $\eta,\theta_{j}$ are $\varphi,\varphi_{j}$, since $\eta$ and $\theta$
are the phase residue and the finite angle here; its abelian variety $A$ is
$\mathbf A$ and its constant $A=-1/(mc_{m+1})$ is $c_{\ast}$; its branch logarithm
$L(z)$ is the strip logarithm $\Log z$, and its formal-group $L,E$ are
$\log_{\mathcal F},\exp_{\mathcal F}$; its time primitive $T$ is $\mathfrak t$, with
$B_{r},\rho,H$ renamed $e_{r},\varrho,g$; its $\ell_{D}$ is $\dlog_{\Dn}$, its
$\mathfrak B_{D}=D\mathfrak m_{K}$ is $\Dn\mm_{\C}$ (which is $\Acl_{\C}$ for
$\Dn=\der$), and its $\operatorname{Obs}_{D}$ is $\Obs^{\C}_{\Dn}$, which is
\emph{not} $\Phi$, and its complex primitive vector $B$ is $\mathfrak b$, since $B$
is an accumulated phase here; its infinitesimal error $\delta$ is an
$\varepsilon$, its formal series $F_{p,\beta,C}$ is $w_{p,\beta,C}$, and its
evaluation points $x,y$ are $x_{m},y_{\omega}$.

Member~13's letters are renamed in \cref{diff:cp:conv:alphabet}, and its most
dangerous letter first (\cref{diff:cp:warn:frakm}): its $\mathfrak m(x)$ is the
Conway monomial map $\omega^{x}$, and its $\mathfrak i$ and $\mathcal O$ are
$\mm$ and $\Ofin$. Its $p_{\alpha}$, $s_{\alpha}$, $P_{\alpha}$, $U_{\alpha}$ are
$\omega^{-\alpha}$, $\varpi_{\alpha}$, $\ell_{<\alpha}$, $u_{\alpha}$; its $a_{\alpha}$,
$A_{\alpha}$, $h_{\alpha}$ are $\ell_{\alpha}^{\dagger}$, $\ell_{<\alpha}^{\dagger}$,
$u_{\alpha}^{\dagger}$; its potentials $q_{\alpha,c}$ and $q$ are $Q_{\alpha,c}$ and
$Q$, since $q$ is the rank-one coefficient here; its coordinate $T$ with $D_{T}$,
$E_{T}$ and partner $U$ is a scale $s$ with $\der_{s}$, $s\der_{s}$ and $u$; its
Wronskian $W$ is $\mathrm{Wr}$; its $\Gamma_{j}$, $K_{j}$, $k_{j}$ are $\Delta_{j}$,
$H_{\Delta_{j}}$, $K_{\Delta_{j}}$; its fundamental matrix $\Phi$, automorphisms
$\sigma_{a,b}$, Borel group $B(\C)$ and Schwarzian $\mathcal S_{\der}$ are $Y$,
$\psi_{a,b}$, $\Bor_{2}(\C)$ and $\Sch$; its infinitesimal primitive $\theta$ of
an angular rate is $\eta$; its Riccati solution $b$ is $w$; and its ordinal
index $\eta$ is $\zeta$.

Finally, $\Ofin$ never takes an argument. Where a source wrote
$\mathcal O(V)$ for holomorphic functions, this report writes $\Hol(V)$.
\end{warning}

\subsection*{Dependency chain}
\label{diff:sub:app:dependency}

\[
 \text{\cref{diff:thm:BM} (imported)}
 \ \Longrightarrow\
 \text{\cref{diff:prop:complexification}, \cref{diff:lem:smallchain}}
 \ \Longrightarrow\
 \text{\cref{diff:thm:unitphase}, \cref{diff:cor:polar}}
\]
\[
 \Longrightarrow\
 \text{\cref{diff:thm:logimage}, \cref{diff:thm:phasecriterion}}
 \ \Longrightarrow\
 \text{\cref{diff:thm:gauge}, \cref{diff:thm:maximalstrip},
 \cref{diff:thm:constant}}.
\]
The ideal-structure theorem \cref{diff:thm:idealstructure} needs only the
constant field, surjectivity, the $H$-field positivity clause and
$\der\omega=1$; the polar decomposition additionally needs real closedness,
standard part and the infinitesimal exponential; the image theorem
additionally needs real exponential compatibility. The bounded-derivative
lemma alone gives a quick independent proof that $\der y=iy$ forces $y=0$
(\cref{diff:cor:oscillator}, second proof). The constant-coefficient theorem
needs only the nonreal-eigenvalue obstruction, the constant field and the real
exponential. The Laurent resolvents need strong additivity and $\der t=-t^{2}$.
The extension theorems need the constant-coefficient obstruction and a formal
Laurent expansion in a new indeterminate. Strong additivity is essential for
the infinitesimal chain rule and for
\cref{diff:thm:coefderivation}; the Laurent identities can be verified inside
formal series algebra independently of any construction of all surreals, the
embedding being what identifies their operator with $\der$. Neither the
composition theorem nor the Picard--Vessiot background is a dependency of the
phase criterion. So the most striking negative conclusions do not depend on
delicate analytic-geometry results or on a full model-theoretic
classification.

\section{Provenance: what each archive contributed}
\label{diff:app:provenance}

This report merges eleven source manuscripts. The manuscripts themselves are
not distributed with this report; the disposition of each source result is
recorded in the provenance material below and in the package README. The
identifiers below are the merge identifiers; among the first seven the original archive names collided, three pairs of them
differing only by a trailing marker, so the mapping is recorded here and in the
package README because it is not recoverable from the filenames. Members 01--09
are pinned to \Commit; member 10 joined later and is pinned to \CommitB;
members 11 and 12 joined together after it, from the archives
\texttt{surreal\_autonomous\_dynamics} and
\texttt{surcomplex\_regular\_singular}, and are pinned to \CommitC. The first of
those archive names is misleading on two counts: its field is $\No[i]$ with real
corollaries, not $\No$ alone, and its ``dynamics'' are flows of autonomous
differential equations, not iteration of maps. Member~13 joined after them,
from the archive \texttt{surreal\_transfinite\_critical\_potentials}, and is pinned to
\CommitD; its code and data are placed under the prefix
\texttt{13-critical-potentials}.

\begin{longtable}{@{}L{0.13\textwidth}L{0.80\textwidth}@{}}
\caption{Provenance by member.}
\label{diff:tab:provenance}\\
\toprule
Member & Distinctive contribution retained in this report \\
\midrule\endfirsthead
\toprule Member & Distinctive contribution retained in this report \\
\midrule\endhead
01 \newline\emph{thresholds and Galois tori} &
The exact image as the boxed organizing identity
(\cref{diff:thm:logimage}); the obstruction as an $\R$-linear surjection with
the gauge \emph{uniqueness} theorem and the rank-one tensor corollary
(\cref{diff:thm:gauge,diff:cor:tensor}); the sharp support threshold over a
real-exponent Hahn field with strictness at exponent $1$ and the
valuation-insufficiency counterexample
(\cref{diff:thm:power,diff:ex:logseparator}); the upper-triangular criterion
with its descending recursion, the Jacobi trace test and the
$\operatorname{diag}(i,-i)$ counterexample
(\cref{diff:thm:triangular,diff:prop:trace}); the rational and diagonal
Picard--Vessiot extensions with from-scratch no-new-constants proofs, the
saturated integer lattice, the torus, and the $(i,2i)$ versus
$(i,\sqrt2 i)$ examples
(\cref{diff:thm:PVone,diff:thm:torus,diff:ex:torusexamples}); the formal
oscillator in complex and real rational form with the order-incompatibility
argument (\cref{diff:eq:realT,diff:thm:orderobstruction}); the
non-summability of the exponential series at $\omega$
(\cref{diff:rem:nonsummableexp}); the phase-independent no-global-exponential
corollary (\cref{diff:cor:noglobalPsi}); the ``at most one, not none''
warning with its two witnesses (\cref{diff:warn:atmostone}); the
class-hygiene discipline (\cref{diff:warn:cosets}); the coefficient-domain
tiering contract (\cref{diff:sub:certificates}). \\[0.3em]
02 \newline\emph{word lemma and nonlinear IVP} &
The entire coordinate half: the common-domain ring with the explicit
disowning of both the shrinking-radius germ ring and the full formal ring
(\cref{diff:eq:coherentring,diff:warn:ringneighbours}); the linear
fundamental-matrix theorem with uniqueness among \emph{all} Hahn-coherent
matrices, the support certificate and the explicit recursion
(\cref{diff:thm:coherentlinear}); the inhomogeneous corollary, the Liouville
determinant identity, the negative-leading-exponent proposition and the
nilpotent counterexample bounding it
(\cref{diff:cor:coherentinhom,diff:cor:determinant,diff:prop:negativecoherent,diff:ex:nilpotent});
the nonlinear deformation theorem with no common degree bound
(\cref{diff:thm:ivp}); monodromy valued in $\GL_{d}(K_{\Gamma})$ with the
example where ordinary monodromy is trivial and full monodromy is not
(\cref{diff:thm:monodromy,diff:ex:infmonodromy}); the two-scale noncommuting
computation and the refutation of the naive exponential
(\cref{diff:sec:twoscale}); the coherent entire solution poled at an infinite
coordinate (\cref{diff:ex:riccatipole}); the cross-category pair
(\cref{diff:warn:crosscategory}); the three-promise separation worked through
$A=i\omega$ (\cref{diff:warn:threepromises}); the strong Laurent resolvent
with its residual certificate and the exact factorial solution
(\cref{diff:prop:polyresolvent,diff:ex:factorialfirst}); the four-category
separation table (\cref{diff:tab:categories}); and the from-scratch proof of
the ordered-word lemma (\cref{diff:app:support}). \\[0.3em]
05 \newline\emph{logarithmic-derivative image} &
The exact decision criterion for rational coefficients and its raising-rather-than-reporting
contract (\cref{diff:thm:rational,diff:warn:rationalimpl}); the
order-theoretic structure and convexity of the bounded-primitive ideal
(\cref{diff:thm:idealstructure}(iii)); the countable-stage differential Hahn
hull with its monomial criterion, its minimality caveat, the real-group
criterion, the $\Q\sqrt2$ counterexample and the
``differentiation closure is not integration closure'' example
(\cref{diff:lem:monomialcriterion,diff:thm:hull,diff:prop:realgroup,diff:ex:sqrt2,diff:ex:logmissing});
the analytic coefficient interface with coefficientwise linear naturality,
the total-derivative chain rule and the separate admissibility condition for
formal power series
(\cref{diff:thm:coefderivation,diff:cor:naturality,diff:thm:totalchainhalo,diff:sub:formaladmissibility});
and the algebraic unit-modulus example
(\cref{diff:ex:algebraicphase}). \\[0.3em]
06 \newline\emph{finite primitive and oscillation} &
The base skeleton, and in particular: the strip theorem with
\emph{pointwise} maximality
(\cref{diff:thm:strip,diff:thm:maximalstrip}); the character-independent
differential defect (\cref{diff:prop:defect}); the products, duals and
no-finite-order corollary (\cref{diff:cor:tensor}); the direct-sum
description (\cref{diff:eq:directsum}); the summary table of distinguishable
situations (\cref{diff:tab:scales}); the factorially divergent forced
oscillator with its closed residual and anti-topological disclaimer
(\cref{diff:thm:factorial,diff:warn:antitopological}); and the theorem that
oscillation cannot retain the $H$-field rule
(\cref{diff:thm:orderobstruction}). \\[0.3em]
07 \newline\emph{exact sequence and rational exponents} &
The explicit $\C((t^{\Q}))$ workspace chapter, with the exact derivative
image, the exact logarithmic-derivative image, the single resonant
coefficient, the resonant normal form and the workspace-versus-ambient
discussion
(\cref{diff:thm:Hahnintegral,diff:thm:Hahnlogder,diff:prop:Hahnambient,diff:rem:resonant,diff:tab:workspaceambient});
the derivative taxonomy including the fine order-field derivative and the
proof that a fine derivative does not determine an intrinsic one
(\cref{diff:tab:derivatives,diff:warn:separators}); the change-of-scale
observation (\cref{diff:prop:rescale}); the two-derivative chain rule with
\emph{derived} summability (\cref{diff:prop:chain,diff:rem:chainstrength});
the integration-by-parts correction for normalized primitives
(\cref{diff:prop:RB}); the restricted $\omega$-series composition bridge and
the precise sense in which two derivatives agree
(\cref{diff:sec:composition,diff:rem:agreement}); the terminating finite
inverse of $\der-i$ on $\C[\omega]$ (\cref{diff:eq:polyinverse}); the
union-of-fields anti-lemma (\cref{diff:warn:antilemma}); and the
four-outcome contract that returns a reason rather than a failure flag
(\cref{diff:tab:outcomes}). \\[0.3em]
08 \newline\emph{coherent IVP and general resolvent} &
The resolvent over a general real-exponent $K_{\Gamma}$ with its substantive
hypothesis (\cref{diff:thm:resolvent,diff:warn:resolventreal}); the
ambient-versus-workspace uniqueness of the factorial solution, with its
exponentially small full-field ambiguity
(\cref{diff:prop:factorialambiguity}); the valuation test for finite phase
(\cref{diff:cor:valphase}); the negative-valued coefficient that destroys
common-domain coherence, with the rescaling that resolves it
(\cref{diff:ex:singular}); ``the same printed equation under different
derivatives'' (\cref{diff:warn:separators}(7)); the Riccati example with an
entire coefficient family (\cref{diff:ex:riccatitan}); the oscillatory
extension proved Picard--Vessiot via the absence of proper differential ideals
in the Laurent ring (\cref{diff:thm:extension}); the anti-Banach disclaimer
(\cref{diff:warn:antibanach}); the certificate table
(\cref{diff:tab:interface}); and the rule that a truncation must record the
\emph{type} of omitted information
(\cref{diff:warn:truncationtype}). \\[0.3em]
09 \newline\emph{BM bridge and derived smallness} &
The derivation of smallness and of strict order reversal on the
infinitesimals from the $H$-field positivity axiom --- the provenance this
report uses (\cref{diff:lem:smallnegative}); the theorem that intrinsic
differentiation preserves common-domain coherence, with the total intrinsic
derivative and the fixed-contour and moving-pole discussion
(\cref{diff:thm:coefderivation,diff:thm:totalchainhalo,diff:cor:naturality,diff:warn:movingpole});
the bounded-primitive ideal as a named object with its full structure theory,
including the cut description and the induced order-preserving quotient
isomorphism (\cref{diff:thm:idealstructure});
the increasingly-slow-scale examples
(\cref{diff:cor:logscales,diff:tab:scales}); the resolvent with its support
certificate and the whole-prefix coefficient recurrence with its exact
guarantee (\cref{diff:thm:resolvent,diff:rem:coefrecursion}); the
constant-coefficient coefficient stream and the
$\sum\omega^{n^{2}}X^{n}$ support counterexample
(\cref{diff:rem:coefstream}); and the derivation-aware representation
interface (\cref{diff:tab:interface}). \\[0.3em]
10 \newline\emph{Hermitian phase spectral integral} &
The whole of \cref{diff:part:matrix}: matrix gauge equivalence and the complete
phase normal form for every finite system, with the differential phase spectrum
as its complete invariant
(\cref{diff:thm:matnormal,diff:conv:phases}(e),\cref{diff:warn:spectrum});
the imported splitting and first-order surjectivity input with its transfer
(\cref{diff:thm:matsurjectivity,diff:warn:matimport}); existence for every
forcing and the exact kernel dimension
(\cref{diff:cor:matinhom,diff:warn:matinhom}); the canonical phase projectors
and the refined trace identity
(\cref{diff:prop:matprojectors,diff:cor:mattrace}); the complete horizontal
metric cone with its anti-Lyapunov example and the unitary normal form
(\cref{diff:thm:matmetric,diff:warn:matlyapunov,diff:thm:matunitary}); the
attained min--max, the entrywise Weyl bound and the bounded-gauge estimate
(\cref{diff:lem:matminmax,diff:lem:matweyl,diff:lem:matboundedgauge,diff:warn:matboundedgauge});
the spectral--integral theorem itself with gap-free perturbation invariance,
the exact kernel count and its sharpness
(\cref{diff:thm:matspectral,diff:cor:matperturbation,diff:cor:matkernel,diff:prop:sharpperturb});
the two deliberate boundary counterexamples and the coupled unitary example
(\cref{diff:ex:matnonnormal,diff:warn:matnonnormal,diff:ex:matcoupled});
normalized non-Abelian integration with its cocycle law and the exact
noncommuting Hahn recurrence
(\cref{diff:thm:matnonabelian,diff:warn:matnonabelian,diff:ex:matnoncommuting});
the real rotation-block descent and the real metric cone with its parity
(\cref{diff:thm:matreal,diff:thm:matrealmetric}); the positive
Ermakov--Pinney dichotomy and the canonical amplitude--phase reduction
(\cref{diff:thm:matpinney,diff:thm:matamplitude}); the exact Airy Hahn
expansion with its Wronskian, amplitude jet and truncation residual
(\cref{diff:sub:matairy}); the resonance algebra with its Dickson finiteness
and the singular three-phase invariant algebra
(\cref{diff:thm:matresonance,diff:ex:matthreephase}); the phase lattice as a
Galois torus, its group algebra and real form
(\cref{diff:cor:mattorus,diff:prop:matPValgebra,diff:warn:mattorus}); and the
set-sized witness-closed reading of the proper-class arguments
(\cref{diff:rem:matclasses}). \\[0.3em]
11 \newline\emph{autonomous dynamics} &
The whole of \cref{diff:aut:part}: the class of normalized derivations and the
scope amendment it forces (\cref{diff:aut:def:normalized,diff:aut:conv:scope});
formal evaluation in several variables, formal implicit uniqueness in a monad,
and strip logarithms for every normalized derivation
(\cref{diff:aut:lem:evaluation,diff:aut:lem:implicit,diff:aut:lem:log});
projective reduction and the reduction-to-a-zero lemma
(\cref{diff:aut:lem:projective,diff:aut:lem:zero}); the realization
classification with its simple-zero and multiple-zero families and the
no-hidden-solution argument
(\cref{diff:aut:thm:main,diff:aut:prop:simple,diff:aut:prop:multiple,diff:aut:rem:nohidden});
uniform finite-scale localization with the three-shift formula, and the
continuum of algebraically independent solutions inside two scales
(\cref{diff:aut:thm:scales,diff:aut:cor:independent}); derivation independence
with its scalar and real corollaries
(\cref{diff:aut:thm:independence,diff:aut:cor:scalarbound,diff:aut:cor:realindep});
the exact real branch count and the worked equations
(\cref{diff:aut:cor:realbranches,diff:aut:ex:logcorrection}); infinitesimal
contraction, the holomorphic dual obstruction and the squarefree hyperelliptic
exclusion with the two counterexamples bounding it
(\cref{diff:aut:prop:contraction,diff:aut:cor:dual,diff:aut:thm:hyperelliptic,diff:aut:ex:hypotheses});
the normalized formal-group logarithm, the identity-monad splitting, the abelian
logarithmic-derivative image, the finite-primitive abelian criterion and the
speed threshold
(\cref{diff:aut:lem:formallog,diff:aut:prop:abeliansplit,diff:aut:thm:abelianimage,diff:aut:cor:abeliancriterion,diff:aut:cor:speeds});
the finite geometric decision criterion (\cref{diff:aut:rem:decision}); and the
review log \path{11-autonomous-dynamics-sources.md}. \\[0.3em]
12 \newline\emph{regular-singular spectral selection} &
The whole of \cref{diff:rs:part}: the Euler derivation on arbitrary real-exponent
Hahn fields (\cref{diff:rs:lem:euler}); the power-Hahn regular-singular class;
the resonance-reaching support and the finite-dependence lemma
(\cref{diff:rs:def:critical,diff:rs:lem:critical}); the normalized Hahn--Levelt
form with support control, and the finite resonance theorem with its degree bound
and finite coefficient replacement
(\cref{diff:rs:thm:normal,diff:rs:thm:finite,diff:rs:cor:replacement,diff:rs:cor:cutoff});
the real-monomial shear (\cref{diff:rs:thm:constantform}); spectral selection with
the ambient dimension $m_{\R}(A_{0})$, perturbation invariance and real descent
(\cref{diff:rs:thm:selection,diff:rs:cor:invariance,diff:rs:cor:real}); the
classifications over $\KR$ and over $\SC$ with exact morphism dimensions and
tensor dimensions
(\cref{diff:rs:thm:Kclass,diff:rs:cor:HomK,diff:rs:thm:SCclass,diff:rs:cor:tensor});
the Hahn solution space, the transcendence of $\log t$, minimal logarithmic
repair with Galois group $(\C,+)$, and the exact Hahn obstruction sequence
(\cref{diff:rs:thm:Hahnsol,diff:rs:lem:tautrans,diff:rs:thm:minlog,diff:rs:thm:forced});
the six worked examples (\cref{diff:rs:sec:examples}); the finite-monomial,
one-logarithm envelope (\cref{diff:rs:thm:workspace}); the power-Hahn boundary
(\cref{diff:rs:prop:boundary}); and the exact finite procedure with its limits
(\cref{diff:rs:sub:procedure}). \\[0.3em]
13 \newline\emph{transfinite critical potentials} &
The whole of \cref{diff:cp:part}: the imported ordinal log-atomic tower with its
limit-stage caveats and scale separation
(\cref{diff:cp:prop:tower,diff:cp:lem:scales}); the summable triangular
differentiation and the exact critical solutions with Wronskian one at every
ordinal (\cref{diff:cp:lem:triangular,diff:cp:thm:critical}); the Wronskian
spanning lemma, the Liouville gauge identity and the Euler classification at a
positive infinite scale (\cref{diff:cp:lem:span,diff:cp:lem:gauge,diff:cp:prop:euler});
the transfinite classification with its Conway normal forms, Wronskians and
Riccati form (\cref{diff:cp:thm:main,diff:cp:rem:riccati}); the earlier-stage
asymptotics, the invisibility theorem, the first infinite stage and the absence
of an all-ordinal potential
(\cref{diff:cp:prop:earlier,diff:cp:thm:invisible,diff:cp:ex:omegaexpansion,diff:cp:prop:noend});
the three Hahn fields at a limit ordinal with bounded inner support, derivative
stability, the $0,1,2$ jump and the exponential-then-primitive gap
(\cref{diff:cp:lem:exponents,diff:cp:prop:fieldclosure,diff:cp:thm:fieldjump,diff:cp:prop:gaps});
the Borel Picard--Vessiot group and its additive remainder
(\cref{diff:cp:lem:algind,diff:cp:thm:pv,diff:cp:cor:additive}); the Schwarzian
form (\cref{diff:cp:lem:schwarzian,diff:cp:cor:schwarzian}); the first limit
stage in normal form (\cref{diff:cp:ex:omega}); and its boundaries, inputs and
formalization interface (\cref{diff:cp:sub:limits}). \\
\bottomrule
\end{longtable}

\begin{remark}[How the merge was performed, and the count of non-claims]
\label{diff:rem:mergemethod}
A result that all seven prove is printed once. Where two members proved the
% (the paragraph below continues to describe the original seven-member merge)
same statement by genuinely different routes, both routes are kept and marked:
\cref{diff:lem:smallnegative} against \cref{diff:rem:sqrtsmallness} for
smallness; the three proofs of \cref{diff:cor:oscillator}; the two
constant-field proofs in
\cref{diff:thm:PVone,diff:thm:extension}; \cref{diff:app:support} against the
published citation; and the two localization constructions of
\cref{diff:sec:localization}, which prove different things. Where two members
stated the same theorem with \emph{different} uniqueness qualifiers, the
qualifier printed is the one its proof supports and no wider
(\cref{diff:conv:ivpuniqueness}), with the stronger linear statement kept
separately (\cref{diff:thm:coherentlinear}). Where one name covered different
coefficient rings, the rings are named
(\cref{diff:thm:resolvent,diff:thm:power,diff:prop:polyresolvent}). Nothing is
strengthened beyond what a source proves.

The first seven manuscripts state, between them,
$13+16+18+13+24+21+13=118$ explicit non-claims, with heavy overlap, and the
eighth states a further $12$, giving $130$. Members~11 and~12 state fifteen
each, by this merge's count, giving $160$, and member~13 states twenty-one,
giving $181$.
\Cref{diff:sec:nonclaims} consolidates them into $106$ numbered items,
organized by theme; every one of the $181$ has a home there, and many items
carry several sources' phrasings at once. Two of the $118$ are artifact-level
discrepancies in the source packages rather than mathematical limitations, and
they are recorded as such in item~N52; a third, the eighth package's
self-rewriting verification script, is item~N64; and the audit findings on the
packages of members~12, 11 and~13 are items N79, N94 and N106.

Integrating member~13 followed the same rules. Its reproof of the obstruction
for nonreal powers of an infinite scale, through an infinitesimal primitive of
the angular rate, is \cref{diff:thm:unitphase} with the direction
(i)$\Rightarrow$(iii) of \cref{diff:thm:phasecriterion}, and is recorded as
agreement (\cref{diff:cp:cor:nonreal}); its derivative-stability argument is an
instance of \cref{diff:lem:monomialcriterion} and is cited as such
(\cref{diff:cp:prop:fieldclosure}); and its stage~$0$ agrees with
\cref{diff:rs:thm:selection} (\cref{diff:cp:rem:stagezero}). Its own proofs are
kept: they use neither surjectivity nor the import of
\cref{diff:thm:matsurjectivity}. Three observations that the member does not
state are recorded as such, with proofs: the Euler operator as a
constant-coefficient operator at a rescaled scale (\cref{diff:cp:rem:constant}),
the missing oscillator as the Euler problem at the scale $e^{\omega}$
(\cref{diff:cp:rem:oscillator}), and the amplitude and phase of a supercritical
member, which use \cref{diff:part:matrix} and flag its import
(\cref{diff:cp:rem:pinney}); the comparison of stage~$0$ with
\cref{diff:rs:part} is likewise made by this merge. Hedges were amended in place
and not removed: items N13, N21, N23 and N48, the remark after
\cref{diff:cor:logscales}, \cref{diff:rs:rem:boundary}, the opening of
\cref{diff:sec:PV}, the remark after \cref{diff:thm:matpinney},
\cref{diff:warn:phasecollision}, and the continuation on multi-scale Hahn
workspaces. One omission in the member's repository comparison is recorded,
with its pin kept (item~N104).

Integrating members~11 and~12 followed the same rules. Their reproofs of results
already in the report are recorded as agreements and cited: smallness
(\cref{diff:aut:lem:small} and \cref{diff:rs:lem:eulermode}(i), against
\cref{diff:lem:smallnegative}), the missing oscillator
(\cref{diff:aut:lem:eigen} and \cref{diff:rs:lem:eulermode}(ii), against the
second proof of \cref{diff:cor:oscillator}), the ordered-word lemma (\cref{diff:rs:lem:words}
against \cref{diff:lem:neumann}, with member~12's Higman-free real-exponent route
kept as a marked second proof), the strip logarithm (\cref{diff:aut:lem:log}
against \cref{diff:thm:strip}), the logarithmic separator
(\cref{diff:rs:prop:boundary} against \cref{diff:ex:logseparator}), the missing
logarithmic primitive (\cref{diff:rs:lem:tautrans} against
\cref{diff:ex:logmissing}), and the monomial adjunction example, printed only as
the second row of \cref{diff:tab:workspaceambient}. Where a reproof holds under
weaker hypotheses --- every normalized derivation, or no surjectivity and no
import --- the statement is kept in its new part with that scope, and no statement
of an older part is widened. Genuinely different routes are kept and marked:
\cref{diff:aut:ex:riccati} beside \cref{diff:cor:riccati}, and
\cref{diff:rs:prop:phasespectrum} beside the existence half of
\cref{diff:thm:matnormal}. Hedges that the new parts affect were amended in place
and not removed: items N2, N3, N21, N22, N23, N24, N25, N48 and N57, the remark
after \cref{diff:cor:riccati}, \cref{diff:warn:BMclauses,diff:warn:nouniqueness,diff:warn:matnonnormal,diff:warn:matlimits},
and three open continuations. One observation that neither source states is
recorded as such (\cref{diff:rs:rem:eigenrecipe}), and one question is recorded
rather than claimed among the open continuations (\cref{diff:sub:nc:open}).

Integrating the eighth member followed the same rules. Its reproofs of results
already in the report --- the convexity and strictness of $\Acl$, the exact
logarithmic-derivative image, the saturation of the relation lattice --- are
recorded as agreements and cited, not printed again
(\cref{diff:rem:matagreement,diff:rem:matrankone,diff:sub:mattorus}). Its two
genuinely different routes are kept and marked: the rank-$r$ differential
simplicity of \cref{diff:prop:matPValgebra}(i) beside the rank-one argument of
\cref{diff:thm:extension}, and the metric proof of the Pinney dichotomy beside
the classical parametrization \cref{diff:eq:matpinneyclassical}. Where it
strengthened a statement this report had deliberately hedged, the hedge was
amended rather than removed and the new hypothesis was named: item~N22 now
records that first-order inhomogeneous existence holds but only by the import
of \cref{diff:thm:matsurjectivity}. Item~N24 distinguishes the complete
finite-system phase classification from the Hermitian eigenvalue formula; the
non-Hermitian counterexample limits the latter. The maintained review corrects
earlier summaries that had conflated those two scopes.
\end{remark}

\section{Verification record, exactly as delivered}
\label{diff:app:verification}

Each source package shipped a deterministic exact-symbolic program together
with its \emph{recorded} run. This appendix reproduces those records. It does
not re-run them and does not merge their counts into a single suite: eleven
suites with eleven scopes remain eleven suites.

\begin{longtable}{@{}L{0.09\textwidth}L{0.13\textwidth}L{0.70\textwidth}@{}}
\caption{The eleven recorded verification runs.}
\label{diff:tab:verification}\\
\toprule
Member & Recorded & Scope and declared cutoffs \\
\midrule\endfirsthead
\toprule Member & Recorded & Scope and declared cutoffs \\\midrule\endhead
01 & $1{,}395$ checks passed &
Standard-library Python only, exact Gaussian rationals, deterministic seed
$20260921$, formal-series cutoff $16$. Categories: Leibniz rule $250$;
additivity $250$; primitive identity including the logarithmic term $250$;
conjugation compatibility $250$; power-threshold regressions $49$;
positive-phase primitives $100$; small-phase differential equation as a jet
$100$; unit-circle exponential identity as a jet $100$; Cayley logarithmic
derivative as a jet $1$; exact polynomial-forcing inverses $36$; repeated
real-root amplitudes $9$. Not a surreal-number implementation and not a proof
checker, by its own docstring. \\[0.3em]
02 & $115$ checks passed &
Python with SymPy, exact symbolic only. Groups: power primitives $7$;
finite-phase formula identities $3$; iterated-logarithm primitives $15$;
Laurent resolvent residuals $10$; noncommuting matrix coefficient equations
and initial values $56$; matrix ODE and determinant jets $5$; nonlinear
rational solution and jets $15$; infinitesimal monodromy jets $2$; total
chain-rule identities $2$. Matrix cutoff is \emph{total degree} $6$, not a
valuation cutoff; Laurent reciprocal order $12$; nonlinear checks through
$t^{12}$. \\[0.3em]
05 & $149$ checks passed &
Python with SymPy. Exposes a restricted rational-phase classifier as an
importable API that type-checks its arguments, rejects nan, complex infinity
and infinities, rejects free symbols other than the declared variable,
coerces to polynomials over $\Q$, and \emph{raises} outside
$\Q(\text{variable})$. Check families include the real power rule over $22$
rationals; power and logarithmic primitives to depth $3$; twelve hand-listed
rational-phase examples; a $6\times6$ degree sweep; eight rejection tests;
the algebraic unit $(1+i\omega)/\sqrt{1+\omega^{2}}$; truncated
$\Emm(-it)$ residuals for $N=0..12$ against
\cref{diff:eq:truncresidual}; the total chain rule with its moving-argument
term; commutation of the two coefficient operations; a fixed-contour sample;
and a constant-coefficient example. \\[0.3em]
06 & $258$ checks passed &
Python with SymPy, models $\der$ as $-t^{2}\,d/dt$ on $t$-expressions and the
extension derivation by $f\mapsto\mathrm{cancel}(i\mathsf T\,\partial
f/\partial\mathsf T)$; uses an ordinary variable as a stand-in for $\omega$.
Categories: total chain rule $7$; scale primitives $23$; formal
exponential/logarithm $45$; phase and gauge algebra $5$; oscillatory
extension $8$; constant-coefficient modes $18$; Laurent resolvent residuals
$60$; factorial oscillator $35$; polynomial resolvent cutoffs $21$; Laurent
valuation bounds $36$. Covers real-power primitives at
$p\in\{-2,0,\tfrac12,1,\tfrac32,2,3\}$, the log-scale ladder at depths
$0..3$ and $p\in\{0,1,2,3\}$, and truncated infinitesimal identities at
orders $1..9$. Its own trailing scope note reads: exact finite symbolic
identities only; these checks do not prove the infinite-support, class-size,
or general differential-field theorems. \\[0.3em]
07 & $291$ checks passed &
Python with SymPy plus a finite Hahn core implemented as a validating map
from rational exponents to Gaussian-rational coefficients, with a terminating
derivative, a primitive that rejects a nonzero resonant $[t^{1}]$, an ambient
primitive pair returning $(p,c)$ for $p+c\log\omega$, a classifier returning
\emph{workspace} / \emph{ambient-only} / \emph{phase-obstructed}, and a
workspace certificate $(r,k)$ asserting that $k$ has positive support. The
normalizer \emph{rejects} floats, strings and non-Gaussian-rational
coefficients, and four such rejections are counted among the passing checks;
the script states that it is not an expression parser. Thirteen groups,
including rational monomials $87$ and the finite product rule $81$. \\[0.3em]
08 & $150$ checks passed &
Python with SymPy, deterministic, with no numerical tolerance --- an
assertion requires exact simplification to zero. Covers monomial derivatives
and primitives at nine rational exponents
$\{-\tfrac52,-1,-\tfrac13,0,\tfrac14,1,\tfrac43,2,\tfrac72\}$; the
logarithmic primitive; sixteen product-rule checks on four sparse
expressions with complex coefficients; twenty-four resolvent-truncation
residuals for $\lambda\in\{1,2,i,1+i\}$ and $N=1..6$, each required to equal
exactly the first omitted derivative with the alternating sign; sixteen
factorial residuals for $N=1..16$; the Riccati coefficients $a_{1}..a_{10}$;
constant-coefficient examples; the Cayley phase identity; and the oscillator
identities in the adjoined extension. \\[0.3em]
09 & $259$ checks passed &
Python with SymPy over formal symbols, with
$\der(\cdot)=\mathrm{expand}(-t^{2}\,\partial/\partial t)$; a matrix identity
or a whole checked coefficient prefix counts as \emph{one} check. Groups:
derivation identities $41$ (including $\der(t^{n})$ for $n=-12..12$ and
$\der^{n}t$ for $n=0..12$ against $(-1)^{n}n!\,t^{n+1}$); primitive
identities $15$; formal small-series identities $24$ at orders
$4,6,8,10$; coherent total-derivative examples $8$, including the moving
coefficient value $1-2/\omega^{2}-2/\omega^{3}$ and a moving-pole identity;
general resolvent residuals $81$ for $\lambda\in\{i,2,1+i\}$ and $N=0..8$;
factorial-series residuals $13$; whole-prefix coefficient recurrences $9$
through $t^{10}$ with lower bounds $1,-2,-5$; constant-coefficient operator
examples $33$; matrix identities $35$. \\[0.3em]
10 & $62$ checks passed &
Python with SymPy, exact symbolic and rational arithmetic throughout, with no
numerical sampling. Groups: the nonnormal counterexample $3$ --- unimodular
gauge, $\der G=AG$, and characteristic polynomial $\zeta^{2}+1$, asserted
\emph{as} a counterexample (\cref{diff:ex:matnonnormal}); the coupled unitary
example $4$ (orthogonality, skew-Hermiticity, the exact phase gauge, the
instantaneous determinant, \cref{diff:ex:matcoupled}); the Jordan metric $2$
(\cref{diff:warn:matlyapunov}); the three-phase system $6$, including its
invariant Hermitian metric, its gauge identity, three polynomial first
integrals verified in the \emph{original} coupled coordinates, and the
quadratic relation of \cref{diff:eq:mattoric}; the noncommuting recurrence
$12$ with $12$ unitarity-coefficient checks, one naive-exponential defect
identity and one assertion that the commutator defect is nonzero
(\cref{diff:ex:matnoncommuting}); the Airy recurrence $16$ with the first
amplitude correction and the exact truncation residual at $N=16$
(\cref{diff:eq:matairyrec,diff:eq:matairyres}); and three Pinney checks
(determinant one, invariant metric, symplectic amplitude reduction) carried out
on arbitrary symbolic jets under the substitution for the second derivative
(\cref{diff:eq:matpinneymetric,diff:eq:matcanonicalgauge}). The record also
carries the degree-$8$ amplitude jet of \cref{diff:eq:matairyjet} and the
$3\times3$ coefficient matrix of \cref{diff:eq:matthreematrix} as strings, both
matching the displayed forms. Its own trailing scope note reads: finite exact
symbolic checks; not a proof of general or infinite-support theorems. \\[0.3em]
11 & $10$ checks passed &
Python 3.13.5 with SymPy 1.14.0, exact symbolic arithmetic in a total-degree
quotient of $\Q[C][X,Y]$, at the default truncation \texttt{--degree 4} (the
option accepts $2$ to $7$, and the number of checks depends on it). Checks: the
logistic solution, with $e^{-\omega}$ modelled by a symbol $q$ with derivative
$-q$; the repeated-root counterexample $(\Dn y)^{2}=y^{3}$; the degree-two
counterexample, with $e^{\omega}$ modelled by $q$ with derivative $q$; the time
primitive of $-t^{2}(1+t)$; the implicit equation of
\cref{diff:aut:ex:logcorrection} through total degrees $1$, $2$, $3$ and $4$; the
printed homogeneous corrections $F_{1},F_{2}$; and the differential residual
through $X^{6}$, with the derivation modelled as
$-X^{2}\partial_{X}+X\partial_{L}$ after $X=1/\omega$, $Y=XL$, $L=\log\omega$. The
record also prints $F_{1},\dots,F_{4}$ and reports $7$ as the first possibly
nonzero degree of the truncation residual. Its trailing note: these tests do not
verify the general proofs, formalization or novelty. \\[0.3em]
12 & $8$ checks passed &
Python with SymPy 1.14.0 (no Python version recorded), exact rational and
Gaussian-rational arithmetic, a finite-cutoff normalizer for finite rational
supports with explicitly supplied generalized-eigenvalue blocks, and a residual
check of $\derT G-AG+GA^{\natural}$ made coefficientwise and independently of the
recursion. Checks: the interaction-generated resonance
$A^{\natural}_{2}=2E_{13}$ of \cref{diff:rs:ex:interaction}, with residual through
weight $6$; the logarithmic fundamental matrix, $\derT Y=AY$ identically and
$\det Y=t^{3}$, with $x=t^{1/2}$ and $\derT$ modelled as
$\tfrac{x}{2}\partial_{x}+\partial_{\tau}$; the nilpotent residual block of
\cref{diff:rs:ex:nonsemisimple} through weight $8$; the mixed-frequency system of
\cref{diff:rs:ex:mixed} through weight $7$; tail invariance after adding a weight
$\tfrac73$ coefficient, compared through weight $4$; the accumulating identity
$((1-\tfrac1j)-1)(-j)=1$ for $2\leq j\leq100$; the forced nilpotent example of
\cref{diff:rs:ex:forced} with a failed constant-image test; and twelve seeded
random cases (seed $20260922$) through weight $4$, with and without Jordan blocks.
Its scope field reads: finite algebraic checks, not a formal verification of the
theorems. \\[0.3em]
13 & $70$ checks passed &
Python 3.13.5 with SymPy 1.14.0, exact symbolic expressions and no
floating-point evaluation; a nonzero residual raises and gives a nonzero exit
status. Suites: the finite hierarchy $39$ --- for lengths $n=0,\dots,12$, with
generators $a_{j}$ and the derivation $a_{j}\mapsto-a_{j}\sum_{k\leq j}a_{k}$
modelling $\ell_{j}^{\dagger}$, the triangular identity \cref{diff:cp:eq:Ader},
the critical potential identity of \cref{diff:cp:thm:critical} and the Riccati
residual \cref{diff:cp:eq:riccati}; Euler identities $8$ --- real-power
eigenvalue, the Euler polynomial, the distinct-root basis with Wronskian $-2d$
for roots $\frac12\pm d$ and the repeated-root basis with Wronskian $1$, in an
ordinary positive variable $T$ with $d/dT$; the gauge identity
\cref{diff:cp:eq:gauge} with its normalization, first-derivative cancellation
and Wronskian $4$; the Schwarzian $4$, including that the ratio of two solutions
has Schwarzian $2Q$; and the Picard--Vessiot suite $15$ --- $\det Y=1$, the four
entries of the fundamental-matrix equation, derivative compatibility of the
automorphism on both generators, the four entries of the Borel action and the
four entries of the group law. Its trailing note: finite identities only, not a
proof-assistant verification; transfinite summability, the imported derivation,
exponent-field membership and novelty need the mathematical arguments. \\
\bottomrule
\end{longtable}

The first seven runs were recorded under Python 3.13.5, six of them under
SymPy 1.14.0 and the seventh using the standard library only. The eighth
records SymPy 1.14.0 and a status of \textsc{pass} with $62$ checks, but
records no Python version; its package pins \path{sympy==1.14.0} and its README
states Python 3.10 or later. The ninth record (member~11) states Python 3.13.5
and SymPy 1.14.0; the tenth (member~12) states SymPy 1.14.0, its README says
Python~3, and its program requires Python 3.10 or later; the eleventh
(member~13) states Python 3.13.5 and SymPy 1.14.0, and its program requires
Python 3.10 or later. All eleven record zero
failures. The consolidated scope limits are item~N48 to item~N52, item~N63 to
item~N64, items N78--N79, item~N94 and items N105--N106 of
\cref{diff:sec:nonclaims}.
Artifact-level discrepancies are recorded there and repeated here for
completeness: one package's PowerShell build script was supplied but not
executed on Windows in its preparation environment; one package's README lists a
\path{SHA256SUMS.txt} that is not present among its delivered files; the eighth
package's program rewrites its own delivered \path{data/verification.json},
computing the package root as the parent of its own directory, so the record
reproduced above is the delivered one and a rerun would overwrite it; member~12's
program writes its output file, by default \path{verification.json} in the
current directory, and its \texttt{verify} target also overwrites the text record;
member~11's build script and README redirect the program's output onto the
delivered text record; neither of those two packages carries a checksum
manifest; member~12's record labels the mixed-frequency ambient dimension
$1$ without computing it; and member~13's program writes a record only when
given \texttt{--output}, while its package README suggests
\texttt{--output verification.txt}, and its manifest lists three package files
--- the article source, the $21$-page PDF and the README --- that are not
shipped here.

\paragraph{Maintained review rerun.}
The eight earlier programs (members 01, 02 and 05--10), unmodified, were
subsequently run on separate temporary copies under Python 3.13.14 and SymPy 1.14.0 (member 01 uses only the standard
library). All exited successfully and reproduced the respective check totals
$1{,}395$, $115$, $149$, $258$, $291$, $150$, $259$ and $62$, with zero
failures. Their generated outputs stayed outside the repository; the delivered
\path{code/} and \path{data/} files were preserved. A separate exact
matrix-coefficient check confirmed the added $[t^4]$ cocycle counterexample
in \cref{diff:ex:matnoncommuting}; it is not included in those historical
suite counts. The later programs of members 11 and 12 were not rerun in this
review. Member~13's program, unmodified, was run once on a separate temporary
copy during its merge, under Python 3.14.4 and SymPy 1.14.0: it exited
successfully with $70$ passed and $0$ failed, and its output agreed with the
delivered record line for line apart from the recorded Python version. These
finite checks do not prove the general theorems.

\section{Building this report}
\label{diff:app:build}

The source \path{article.tex} is standalone: the bibliography is internal, and
no external bibliography database, graphics file, font file, shell escape or
network access is required. A standard TeX installation providing the packages
named in the preamble suffices. From this directory,
\begin{verbatim}
latexmk -pdf -interaction=nonstopmode article.tex
\end{verbatim}
and, if \texttt{latexmk} is unavailable, run \texttt{pdflatex} with
\texttt{-interaction=nonstopmode} three times so that the table of contents,
the cross-references and the \texttt{cleveref} labels settle. The directory
also contains \path{README.md}, member~11's review log
\path{11-autonomous-dynamics-sources.md}, and the unmodified \path{code/} and
\path{data/} directories of the eleven source packages. To re-run a source
suite, run it \emph{on a copy}: several of the delivered programs rewrite their own evidence files beside themselves, the
eighth writes to a \path{data/} directory it creates one level \emph{above} its
own, member~12's program writes \path{verification.json} into the current
directory, the build scripts of members~11 and~12 redirect output onto their
recorded transcripts, and member~13's program writes whatever file its
\texttt{--output} option names, all still under the packages' original file
names; so a
byte-identity check performed after a run compares two equally modified copies
and tells you nothing.

\clearpage
\phantomsection
\addcontentsline{toc}{section}{References}
\begin{thebibliography}{99}\small\raggedright

\bibitem{BM}
Alessandro Berarducci and Vincenzo Mantova,
\emph{Surreal numbers, derivations and transseries},
Journal of the European Mathematical Society \textbf{20} (2018), 339--390.
\href{https://doi.org/10.4171/JEMS/769}{doi:10.4171/JEMS/769}.
Preprint: \href{https://arxiv.org/abs/1503.00315v3}{arXiv:1503.00315v3}.
Theorems A and B of the introduction, with the normalization stated between
them, and the underlying Theorems 6.30 and 7.7; normal forms and summability
in Sections 2--3.

\bibitem{BMGerms}
Alessandro Berarducci and Vincenzo Mantova,
\emph{Transseries as germs of surreal functions},
Transactions of the American Mathematical Society \textbf{371} (2019),
3549--3592.
\href{https://doi.org/10.1090/tran/7428}{doi:10.1090/tran/7428}.
Preprint: \href{https://arxiv.org/abs/1703.01995}{arXiv:1703.01995v3}.
The composition framework and Taylor theorem (Corollaries 7.6--7.7, Theorems
7.14 and 8.2), the summability preliminaries of \S3, and the global
composition-incompatibility theorem (Theorem 8.4).

\bibitem{ADH}
Matthias Aschenbrenner, Lou van den Dries and Joris van der Hoeven,
\emph{The surreal numbers as a universal $H$-field},
Journal of the European Mathematical Society \textbf{21} (2019), 1179--1199.
\href{https://doi.org/10.4171/JEMS/858}{doi:10.4171/JEMS/858}.
Preprint: \href{https://arxiv.org/abs/1512.02267v3}{arXiv:1512.02267v3}.
Theorem 1 concerns an elementary embedding in the ordered differential
language. Section~2 --- Proposition~2.5, Lemmas~2.6--2.8, the pre-derivation
formula on p.~10 and the real-linear exponent formula in the proof of
Lemma~2.3, numbered as in the arXiv version~3 consulted by member~13 --- supplies
the ordinal log-atomic formulas imported in \cref{diff:cp:prop:tower}, and
p.~17 the solvability predicate $\Omega$ cited in \cref{diff:cp:rem:inputs}.

\bibitem{ADHNorm}
Matthias Aschenbrenner, Lou van den Dries and Joris van der Hoeven,
\emph{Normalizing Asymptotic Differential Equations},
\href{https://arxiv.org/abs/2403.19732v5}{arXiv:2403.19732v5} (2026).
Earlier versions carried the title \emph{A Normalization Theorem in Asymptotic
Differential Algebra}. Used for the splitting of monic scalar operators into
first-order factors and the surjectivity of nonzero scalar operators
(\S2.4, Lemma 2.4.19), which \cref{diff:thm:matsurjectivity} imports and
transfers, and for the universal exponential extensions and the
Schwarz-closed-field discussion cited in
\cref{diff:sub:matpinney,diff:sub:mattorus}.

\bibitem{ADHSecond}
Matthias Aschenbrenner, Lou van den Dries and Joris van der Hoeven,
\emph{Revisiting Second-Order Linear Differential Equations over Hardy Fields},
\href{https://arxiv.org/abs/2603.02013v1}{arXiv:2603.02013v1} (2026).
Theorems A and B and the amplitude--phase formulations of the introduction are
the antecedent of
\cref{diff:thm:matpinney,diff:thm:matamplitude}. No solution of the
Boshernitzan conjecture named there is claimed here.

\bibitem{Pinney}
Edmund Pinney,
\emph{The nonlinear differential equation $y''+p(x)y+cy^{-3}=0$},
Proceedings of the American Mathematical Society \textbf{1} (1950), no.~5,
681.
\href{https://doi.org/10.1090/S0002-9939-1950-0037979-4}{doi:10.1090/S0002-9939-1950-0037979-4}.
The sign of the parameter is convention dependent; the equation used in
\cref{diff:sub:matpinney} is $\der^{2}\rho+q\rho=\rho^{-3}$.

\bibitem{Gonshor}
Harry Gonshor,
\emph{An Introduction to the Theory of Surreal Numbers},
London Mathematical Society Lecture Note Series, vol.~110,
Cambridge University Press, 1986.
Standard reference for normal forms and the surreal exponential.

\bibitem{EK}
Philip Ehrlich and Elliot Kaplan,
\emph{Surreal ordered exponential fields},
Journal of Symbolic Logic \textbf{86}, no.~3 (2021), 1066--1115.
\href{https://doi.org/10.1017/jsl.2021.59}{doi:10.1017/jsl.2021.59}.
Preprint: \href{https://arxiv.org/abs/2002.07739}{arXiv:2002.07739v3}.
See \S11 for the global trigonometric and surcomplex exponential
constructions.

\bibitem{BG}
Olivier Bournez and Quentin Guilmant,
\emph{Surreal fields stable under exponential, logarithmic, derivative and
anti-derivative functions},
\href{https://arxiv.org/abs/2211.08396}{arXiv:2211.08396} (2022).
Cited for structural workspace results; no theorem of it is used here.

\bibitem{Neumann}
Bernhard H. Neumann,
\emph{On ordered division rings},
Transactions of the American Mathematical Society \textbf{66} (1949),
202--252.
\href{https://doi.org/10.1090/S0002-9947-1949-0032593-5}{doi:10.1090/S0002-9947-1949-0032593-5}.
The positive-support summability lemma, used in its standard
additive-exponent form and also proved in \cref{diff:app:support}.

\bibitem{Higman}
Graham Higman,
\emph{Ordering by divisibility in abstract algebras},
Proceedings of the London Mathematical Society (3) \textbf{2} (1952),
326--336.
\href{https://doi.org/10.1112/plms/s3-2.1.326}{doi:10.1112/plms/s3-2.1.326}.
The finite-word lemma.

\bibitem{Teschl}
Gerald Teschl,
\emph{Ordinary Differential Equations and Dynamical Systems},
Graduate Studies in Mathematics, vol.~140,
American Mathematical Society, 2012.
Chapters 3--4 concern differential equations in the complex domain, including
complex local existence.

\bibitem{PS}
Marius van der Put and Michael F. Singer,
\emph{Galois Theory of Linear Differential Equations},
Grundlehren der mathematischen Wissenschaften, vol.~328, Springer, 2003.
\href{https://doi.org/10.1007/978-3-642-55750-7}{doi:10.1007/978-3-642-55750-7}.
Chapters 1--2 supply the Picard--Vessiot and differential-module terminology.

\bibitem{Singer}
Michael F. Singer,
\emph{Introduction to the Galois theory of linear differential equations},
in M.~A.~H. MacCallum and A.~V. Mikhailov (eds.),
\emph{Algebraic Theory of Differential Equations},
London Mathematical Society Lecture Note Series, vol.~357,
Cambridge University Press, 2009, pp.~1--82.
Preprint: \href{https://arxiv.org/abs/0712.4124}{arXiv:0712.4124v2}.

\bibitem{Levelt}
A.~H.~M. Levelt,
\emph{Jordan decomposition for a class of singular differential operators},
Arkiv f\"or Matematik \textbf{13} (1975), 1--27.
\href{https://doi.org/10.1007/BF02386195}{doi:10.1007/BF02386195}.
Cited by member~12 as a direct antecedent of the regular-singular reduction of
\cref{diff:rs:part}.

\bibitem{KW}
Masoud Kamgarpour and Samuel Weatherhog,
\emph{Jordan decomposition for formal $G$-connections},
\href{https://arxiv.org/abs/1702.03608}{arXiv:1702.03608v2} (2019).
Member~12 cites in particular the discussion of unipotent differential operators
and regular-singular systems in its Section~3.5 as a close classical comparison.

\bibitem{vdH}
Joris van der Hoeven,
\emph{Operators on generalized power series},
Illinois Journal of Mathematics \textbf{45} (2001), 1161--1190.
\href{https://doi.org/10.1215/ijm/1258138061}{doi:10.1215/ijm/1258138061}.
Generalized-series operator methods, strong linearity and support-controlled
implicit constructions; the recursion of \cref{diff:rs:part} is proved directly
and not attributed as a new operator principle.

\bibitem{NPT}
Marc Paul Noordman, Marius van der Put and Jaap Top,
\emph{Autonomous first order differential equations},
Transactions of the American Mathematical Society \textbf{375} (2022).
Preprint: \href{https://arxiv.org/abs/1904.08152v1}{arXiv:1904.08152v1}. The
geometric formalism of \cref{diff:aut:part}, the local formal normal forms
(Lemma~6.1) and Rosenlicht's algebraic-independence theorem (Lemma~7.2,
Proposition~7.3); the numbered statements are those of the arXiv version~1
consulted by member~11, not assumed to match the journal numbering.

\bibitem{Milne}
J.~S. Milne,
\emph{Abelian Varieties}, course notes, version 2.00, 16 March 2008.
\url{https://www.jmilne.org/math/CourseNotes/AV.pdf}.
The standard algebraic-group and abelian-variety setting only; the formal
logarithm argument of \cref{diff:aut:lem:formallog} is proved in full.

\bibitem{Stacks}
The Stacks Project Authors,
\emph{The Stacks Project}, Section 29.43, Valuative criteria,
especially Lemma 29.43.1 (Tag~0BX5).
\url{https://stacks.math.columbia.edu/tag/0BX5}.
The valuative reading of the projective reduction of
\cref{diff:aut:lem:projective}, which is proved there by elementary means.

\bibitem{RepoMap}
\Repo, \emph{The surreal and surcomplex reports}: documentation catalogue and
manifest, \path{docs/README.md} and \path{docs/manifest.tex}, snapshot at
commit \Commit.
\url{https://github.com/VladimirReshetnikov/Surreal/blob/e260237db9b71da8b74a0c13c8e6355119091100/docs/README.md}.
The remaining repository references use this same snapshot.

\bibitem{RepoCatalogue}
\Repo, \emph{The surreal and surcomplex reports}: the report descriptions,
dependency guidance, coefficient-ring distinctions and workspace discussion
of the same catalogue, same snapshot.

\bibitem{RepoTrig}
\Repo, \emph{Trigonometry on the Surcomplex Plane: Canonical Phases,
Arbitrary-Scale Geometry, and Infinitesimal Degeneration},
\path{docs/surcomplex/trigonometry/}, same snapshot.
\url{https://github.com/VladimirReshetnikov/Surreal/blob/e260237db9b71da8b74a0c13c8e6355119091100/docs/surcomplex/trigonometry/README.md}.
The finite-angle decomposition, the global phase-extension discussion, and the
two explicit boundaries: it chooses no derivation on $\No$ and its
classification of phase extensions is not a classification of
differential-equation solutions.

\bibitem{RepoAnalysis}
\Repo, \emph{Surcomplex Analysis: holomorphic functions on $\No[i]$; germs,
coherent Hahn sections, and infinitesimal zero geometry},
\path{docs/surcomplex/analysis/}, same snapshot.
\url{https://github.com/VladimirReshetnikov/Surreal/blob/e260237db9b71da8b74a0c13c8e6355119091100/docs/surcomplex/analysis/README.md}.
In particular its derivative distinctions and its derivation-obstruction
theorem for one chosen canonical global surcomplex exponential.

\bibitem{RepoCAS}
\Repo, \emph{Surreal and Surcomplex Numbers in Symbolic Computer Algebra},
\path{docs/foundations-and-computation/computer-algebra/}, same snapshot.
\url{https://github.com/VladimirReshetnikov/Surreal/blob/e260237db9b71da8b74a0c13c8e6355119091100/docs/foundations-and-computation/computer-algebra/README.md}.
The capability-tier architecture, the derivative distinctions, and the
subsection ``Differential equations must name their category''.

\bibitem{RepoFound}
\Repo, \emph{Foundations for surreal and surcomplex mathematics},
\path{docs/foundations-and-computation/foundations/}, same snapshot.
\url{https://github.com/VladimirReshetnikov/Surreal/blob/e260237db9b71da8b74a0c13c8e6355119091100/docs/foundations-and-computation/foundations/README.md}.
Sets, classes, set-sized workspaces, the four-interface discipline, and the
limits of verification claims.

\bibitem{RepoSpectral}
\Repo, \emph{Surcomplex Spectral Theory: Singular Values, Infinitesimal Rank,
and Multiscale Stability},
\path{docs/surcomplex/spectral-theory/}, at revision \CommitB, the snapshot to
which the eighth source manuscript is pinned.
\url{https://github.com/VladimirReshetnikov/Surreal/blob/aa846271b4dcae2c055b216126a87210292ec19b/docs/surcomplex/spectral-theory/article.tex}.
Its ``spectrum'' is the algebraic eigenvalue and singular-value list of a
finite matrix, which is \emph{not} the differential phase spectrum of
\cref{diff:conv:phases}(e); see \cref{diff:warn:spectrum}.

\bibitem{RepoGaps}
\Repo, \emph{A countable genetic cut across the finite--infinite gap},
\path{docs/surreal/genetic-gaps-and-primitives/}, same snapshot.
\url{https://github.com/VladimirReshetnikov/Surreal/blob/e260237db9b71da8b74a0c13c8e6355119091100/docs/surreal/genetic-gaps-and-primitives/README.md}.
Its nonconstant functions with vanishing pointwise derivative concern the
\emph{fine} order-field derivative of a function, not a scalar derivation, and
neither its primitive non-uniqueness nor the analytic reports' local functions
contradicts the constant-field theorem used here.

\bibitem{RepoTailSpans}
\Repo, \emph{Tail-Spans and Differential Transcendence in Surreal and
Surcomplex Hahn Fields},
\path{docs/surreal/tail-spans-and-differential-transcendence/}.
It cites \cref{diff:eq:HahnQ} and \cref{diff:prop:hahnderiv} rather than
reproving them, and uses them over the real-algebraic coefficient subfield;
see \cref{diff:rem:tailspans}. Its own results are differential
transcendence statements over set-sized Hahn bases, with a targeted
literature search and no refereeing or proof-assistant verification
claimed. Its Euler derivation is the operator written $\derT$ in
\cref{diff:rs:part}.

\bibitem{RepoTate}
\Repo, \emph{Hahn--Tate Uniformization at Arbitrary Valuation Rank},
\path{docs/surcomplex/hahn-tate-uniformization/}. Its limitations section records
that curves with nonnegative valuation of $j$ need formal groups around a
good-reduction model, and it makes no assertion about the Berarducci--Mantova
derivation; see \cref{diff:aut:rem:tate}.

\bibitem{RepoSnapshotC}
\Repo, root \path{README.md}, documentation catalogue \path{docs/README.md} and
\path{docs/surcomplex/differential-equations/README.md} at commit \CommitC, the
snapshot to which members~11 and~12 are pinned and against which they compared
themselves (items N77 and N93).

\bibitem{RepoSnapshotD}
\Repo, repository tree, root \path{README.md}, catalogue \path{docs/README.md},
the listing of \path{docs/new/}, and
\path{docs/surcomplex/differential-equations/article.tex} at commit \CommitD, the
snapshot to which member~13 is pinned and against which it compared itself
(item~N104). At that commit this article was byte-identical to the one into
which member~13 was merged.

\bibitem{KT}
Helge Kr\"uger and Gerald Teschl,
\emph{Effective Pr\"ufer angles and relative oscillation criteria},
preprint \href{https://arxiv.org/abs/0709.0127}{arXiv:0709.0127v2} (2007).
Section~2: the Kneser--Weber--Hartman--Hille hierarchy and the positive-solution
reduction behind it, cited by member~13 as the antecedent of the finite critical
hierarchy of \cref{diff:cp:part}. The sign of their Schr\"odinger potential is
opposite to that of $Q$ in $\der^{2}y+Qy=0$.

\bibitem{ADHBook}
Matthias Aschenbrenner, Lou van den Dries and Joris van der Hoeven,
\emph{Asymptotic Differential Algebra and Model Theory of Transseries},
Annals of Mathematics Studies 195, Princeton University Press, 2017;
corrected preprint \href{https://arxiv.org/abs/1509.02588}{arXiv:1509.02588v7}
(2025). The general theory of transseries and $H$-fields; its existence theory
is neither replaced nor claimed in \cref{diff:cp:part}.

\bibitem{Mantova2026}
Vincenzo Mantova,
\emph{Monotonicity and a Taylor approximation theorem for transseries},
preprint \href{https://arxiv.org/abs/2601.07747}{arXiv:2601.07747v2} (2026).
Composition results cited by member~13 as context; not used or superseded.

\end{thebibliography}
\end{document}
